\documentclass[10pt,oneside]{book}

\usepackage[T1]{fontenc}
\usepackage[utf8]{inputenc}
\usepackage[french]{babel}
\usepackage{amsmath,amssymb,mathrsfs}
\usepackage[all]{xy}
\usepackage{tikz-cd}
\usepackage[shortlabels]{enumitem}
\usepackage[a4paper,margin=22mm]{geometry}
\usepackage[
  colorlinks=true,
  linkcolor=blue,
  hypertexnames=false,
  unicode=true,
  pdfencoding=auto
]{hyperref}
\hypersetup{
  pdftitle={Éléments de géométrie algébrique IV, troisième partie -- édition française},
  pdfauthor={Alexander Grothendieck et Jean Dieudonné},
  pdfsubject={Transcription diplomatique française depuis l'autorité NUMDAM}
}

\renewcommand{\thesection}{\arabic{section}}
\renewcommand{\thesubsection}{\thesection.\arabic{subsection}}
\newcommand{\oldpage}[2][]{%
  \setcounter{footnote}{0}%
  \marginpar{\scriptsize\textbf{#1}~|~#2}\ignorespaces}
\newenvironment{env}[1][]{\par\medskip\noindent\textbf{(#1)}\ }{\par}
\newenvironment{definition}[1][]{%
  \par\medskip\noindent\itshape Définition (#1). ---\ }{\par}
\newenvironment{proposition}[1][]{%
  \par\medskip\noindent\itshape Proposition (#1). ---\ }{\par}
\newenvironment{theorem}[1][]{%
  \par\medskip\noindent\itshape Théorème (#1). ---\ }{\par}
\newenvironment{corollary}[1][]{%
  \par\medskip\noindent\itshape Corollaire (#1). ---\ }{\par}
\newenvironment{remark}[1][]{%
  \par\medskip\noindent\itshape Remarque (#1). ---\ }{\par}
\newenvironment{remarks}[1][]{%
  \par\medskip\noindent\itshape Remarques (#1). ---\ }{\par}
\newenvironment{lemma}[1][]{%
  \par\medskip\noindent\itshape Lemme (#1). ---\ }{\par}
\newenvironment{proof}{%
  \par\medskip\noindent\textit{Démonstration. ---\ }}{\par}

\begin{document}
\clearpage
\thispagestyle{plain}
\oldpage[IV-3]{5}
\phantomsection
\label{chapter:IV3-fr}

\begin{center}
{\large CHAPITRE IV (suite)\par}
\vspace{1.5em}
{\Large\bfseries ÉTUDE LOCALE DES SCHÉMAS\\
ET DES MORPHISMES DE SCHÉMAS\par}
\end{center}

\setcounter{section}{7}
\section{Limites projectives de préschémas}
\label{section:IV.8-fr}

\setcounter{subsection}{0}
\subsection{Introduction}
\label{subsection:IV.8.1-fr}

\begin{env}[8.1.1]
\label{IV.8.1.1-fr}
Dans ce paragraphe, nous allons étudier systématiquement la situation
suivante. Soient $I$ un ensemble préordonné filtrant croissant,
$(A_\lambda,\varphi_{\mu\lambda})$ un système inductif d'anneaux ayant $I$
pour ensemble d'indices, $A=\varinjlim A_\lambda$ sa limite inductive. Pour
tout $\alpha\in I$ et tout $A_\alpha$-préschéma $X_\alpha$, considérons les
$A_\lambda$-préschémas
$X_\lambda=X_\alpha\otimes_{A_\alpha}A_\lambda$ pour $\lambda\geq\alpha$,
et le $A$-préschéma $X=X_\alpha\otimes_{A_\alpha}A$~; il est clair que les
préschémas $X_\lambda$ (pour $\lambda\geq\alpha$) forment un \emph{système
projectif}, et on verra (8.2.5) que $X$ est une \emph{limite projective} de ce
système dans la catégorie des préschémas. Nous nous proposons de trouver des
conditions sur $X_\alpha$ ou sur les $A_\lambda$ permettant de démontrer des
propriétés du type suivant~: pour que $X$ possède une propriété $P$ (par
exemple, la propriété d'être propre sur $S=\operatorname{Spec}(A)$, ou
irréductible, ou connexe, etc.), il faut et il suffit qu'il existe un indice
$\mu\geq\alpha$ tel que, pour tout $\lambda\geq\mu$, $X_\lambda$ ait (par
rapport à $S_\lambda=\operatorname{Spec}(A_\lambda)$, le cas échéant), la
même propriété $P$. Nous obtiendrons des énoncés analogues pour des propriétés
de $\mathcal O_X$-Modules, de $A$-morphismes de $A$-préschémas, etc. Nous
montrerons également (8.9.1) que la donnée d'un $A$-préschéma de présentation
finie (1.6.1) équivaut essentiellement à la donnée d'un $A_\lambda$-préschéma
de présentation finie $X_\lambda$ pour un $\lambda$ assez grand, $X$ étant
alors \emph{isomorphe à} $X_\lambda\otimes_{A_\lambda}A$. On a aussi des
énoncés analogues pour les $\mathcal O_X$-Modules de présentation finie,
leurs homomorphismes, les $A$-morphismes de $A$-préschémas de présentation
finie, etc.
\end{env}

\begin{env}[8.1.2]
\label{IV.8.1.2-fr}
L'utilité de tels résultats apparaîtra par exemple dans les questions
suivantes~:
\begin{enumerate}
\item[\emph{a)}] Soient $Y$ un préschéma, $y$ un point de $Y$, $(U_\lambda)$
le système projectif filtrant décroissant des voisinages ouverts affines de
$y$ dans $Y$~; si $A_\lambda$ est l'anneau de $U_\lambda$, les $A_\lambda$
forment un système inductif filtrant croissant, dont la limite inductive $A$
est l'anneau local $\mathcal O_y$~; d'ailleurs, si on désigne par
$\mathfrak p_\lambda$ l'idéal premier $\mathfrak j_y$ dans l'anneau
$A_\lambda$, le système inductif $(A_\lambda)$ est cofinal à tout système
inductif $(A_\alpha)_f$, où $f$ parcourt
$A_\alpha-\mathfrak p_\alpha$ (pour un $\alpha$ \emph{fixé}), car les $D(f)$
forment un système fondamental de voisinages de $y$ dans $U_\alpha$,
\oldpage[IV-3]{6}
donc dans $Y$. Les résultats du présent paragraphe impliqueront que la
géométrie algébrique des $\mathcal O_y$-préschémas de présentation finie (et
la théorie des Modules de présentation finie sur ces préschémas) est
essentiellement équivalente à la géométrie algébrique des préschémas de
présentation finie sur des voisinages ouverts «~assez petits~» de $y$. Ainsi,
l'énoncé (8.10.5, (xii)) implique que si un morphisme $f:X\to Y$ est de
présentation finie, alors, pour que
$X\times_Y\operatorname{Spec}(\mathcal O_y)$ soit propre sur
$\operatorname{Spec}(\mathcal O_y)$, il faut et il suffit qu'il existe un
voisinage ouvert $U$ de $y$ dans $Y$ tel que $f^{-1}(U)$ soit propre sur $U$.

Un cas particulièrement important, et dans une certaine mesure classique,
est celui où $Y$ est intègre et où $y$ est son point générique, de sorte que
$\mathcal O_y$ n'est autre que \emph{le corps $R(Y)=K$ des fonctions
rationnelles sur $Y$}. Les résultats du présent paragraphe reviennent alors à
interpréter la géométrie algébrique sur $K$ en termes de la géométrie
algébrique au-dessus d'ouverts non vides «~assez petits~» de $Y$, c'est-à-dire,
intuitivement, en termes de «~familles~» d'objets géométriques indexées par les
points d'un tel ouvert. Ce point de vue a d'ailleurs été couramment utilisé
depuis longtemps, non seulement en Géométrie algébrique sur des corps
algébriquement clos, mais aussi dans l'étude arithmétique des variétés définies
sur un \emph{corps de nombres} $K$ (extension finie de $\mathbf Q$), en
considérant ce dernier comme le corps des fractions de son anneau des entiers
$A$ («~théorie de la \emph{réduction modulo} $\mathfrak p$~» ($\mathfrak p$
idéal premier de $A$)~; cf. (I, 3.7)). Les résultats des \S\S~8 et 9
fournissent donc entre autres des fondements du langage de la «~réduction
modulo $\mathfrak p$~» en arithmétique.

On notera que dans l'exemple envisagé ici, les morphismes
$S_\mu\to S_\lambda$ (pour $\lambda\leq\mu$) sont les \emph{immersions
ouvertes} canoniques $U_\mu\to U_\lambda$, et \emph{a fortiori} sont des
morphismes \emph{plats} (mais non fidèlement plats en général), ce qui
explique l'intérêt des énoncés qui font appel à une telle restriction.

\item[\emph{b)}] Supposons que les $A_\lambda$ soient des \emph{corps}, de
sorte que $A=\varinjlim A_\lambda$ est aussi un corps. Ce cas se présente
généralement lorsqu'on part de données géométriques au-dessus d'un corps
quelconque $K$, que l'on considère comme extension d'un corps $k$ (par exemple
le sous-corps premier de $K$). Il y a alors intérêt à considérer $K$ comme
limite inductive de ses sous-extensions qui sont \emph{de type fini sur $k$},
ce qui permet dans de nombreuses questions de se ramener au cas où $K$ est
une extension de type fini de $k$. Utilisant aussi la méthode esquissée dans
\emph{a)}, on peut alors se ramener généralement au cas d'un anneau de base
$A$ qui est une \emph{algèbre intègre de type fini sur $k$}.

On notera que dans cet exemple, les morphismes $S_\mu\to S_\lambda$ sont
\emph{fidèlement plats}.

\item[\emph{c)}] Supposons qu'on s'intéresse aux propriétés, locales sur
$Y$, des préschémas de présentation finie au-dessus d'un préschéma quelconque
$Y$, que l'on peut dès lors supposer affine d'anneau $A$. Il y a alors intérêt
à considérer $A$ comme limite inductive de ses sous-anneaux qui sont des
$\mathbf Z$-algèbres \emph{de type fini}, ce qui permet de ramener de
nombreuses questions au cas où $Y$ est le spectre d'une telle algèbre. C'est
là l'explicitation du «~point de vue kroneckerien~» suivant lequel la
Géométrie algébrique se réduit à la Géométrie algébrique des préschémas de
type fini sur $\mathbf Z$ (que l'on qualifie parfois de
\oldpage[IV-3]{7}
«~Géométrie algébrique absolue~»). Cet exemple nous montre en particulier que
dans la plupart des questions «~relatives~» sur un préschéma de base $Y$, on
peut se ramener au cas où $Y$ est \emph{noethérien}.

On notera que dans cet exemple, au contraire des précédents, les morphismes
$S_\mu\to S_\lambda$ n'ont en général aucune propriété particulière de
régularité.
\end{enumerate}

Par la suite, quand nous appliquerons les résultats qui vont suivre à l'une
quelconque des trois situations particulières qui viennent d'être décrites,
nous nous dispenserons de reprendre en détail le procédé qui permet de faire
ces applications, nous contentant de renvoyer à ce qui précède.
\end{env}

\begin{env}[8.1.3]
\label{IV.8.1.3-fr}
Dans l'exemple \emph{a)} de (8.1.2), nous avons vu que si $Y$ est un
préschéma intègre de point générique $y$, $f:X\to Y$ un morphisme de
présentation finie, alors, si la fibre générique $f^{-1}(y)$ est propre sur
$\boldsymbol k(y)$, il y a un voisinage ouvert $U$ de $y$ tel que $f^{-1}(U)$
soit propre sur $U$~; il en résulte \emph{a fortiori} que pour tout $s\in U$,
$f^{-1}(s)$ est propre sur $\boldsymbol k(s)$. Il arrive qu'on ait besoin
d'une réciproque, assurant que si $f^{-1}(s)$ est propre sur
$\boldsymbol k(s)$ pour «~suffisamment~» de points $s\in U$, alors
$f^{-1}(y)$ est propre sur $\boldsymbol k(y)$. Par exemple, supposons que $X$
et $Y$ soient des préschémas algébriques sur un corps $k$ algébriquement clos
(on peut prendre pour $k$ le corps $\mathbf C$ des nombres complexes, pour
fixer les idées)~; on a parfois besoin de savoir que si, pour tout $s\in Y$,
\emph{rationnel sur $k$}, la fibre $f^{-1}(s)$ est propre sur
$\boldsymbol k(s)$, alors $f^{-1}(y)$ est propre sur $\boldsymbol k(y)$, et
par suite $f^{-1}(U)$ est propre sur $U$ pour un voisinage $U$ de
$y$\footnote{On notera qu'un tel énoncé est en fin de compte purement
géométrique, en ce sens qu'il ne fait appel qu'à des points rationnels sur
$k$, et non à des points génériques~; par exemple, lorsque $k=\mathbf C$, cet
énoncé a une signification topologique évidente pour l'analyste, en
interprétant «~propre~» au sens topologique du terme, pour les espaces
sous-jacents aux espaces analytiques formés des points de $X$ et $Y$
rationnels sur $\mathbf C$.}. Or cet énoncé résultera aisément du suivant~:
l'ensemble $E$ des points $s\in Y$ tels que $f^{-1}(s)$ soit propre sur
$\boldsymbol k(s)$ est \emph{constructible} (et par conséquent identique à
$Y$ tout entier s'il contient les points fermés de $Y$, grâce au théorème des
zéros de Hilbert (10.4.8))~; cela équivaut encore à dire que si $f^{-1}(y)$
\emph{n'est pas propre} sur $\boldsymbol k(y)$, alors il existe un voisinage
ouvert $U$ de $y$ tel que $f^{-1}(s)$ ne soit pas propre sur
$\boldsymbol k(s)$ pour tout $s\in U$ (cf. (9.6.1, (iv))). Cet exemple
illustre l'intérêt qu'il y a à développer systématiquement des \emph{critères
de constructibilité} pour les notions les plus importantes~: c'est ce qui
sera fait au \S~9.
\end{env}

\subsection{Limites projectives de préschémas}
\label{subsection:IV.8.2-fr}

\begin{env}[8.2.1]
\label{IV.8.2.1-fr}
Soient $S_0$ un espace annelé, $L$ un ensemble préordonné filtrant croissant,
$(\mathscr A_\lambda,\varphi_{\mu\lambda})$ un système inductif de
$\mathcal O_{S_0}$-Algèbres (non nécessairement commutatives), ayant $L$ pour
ensemble d'indices. On sait que, considéré comme système inductif de
$\mathcal O_{S_0}$-Modules, $(\mathscr A_\lambda,\varphi_{\mu\lambda})$ admet
une limite inductive $\mathscr A$~; désignons par
$\varphi_\lambda:\mathscr A_\lambda\to\mathscr A$ l'homomorphisme canonique
(de $\mathcal O_{S_0}$-Modules). Soit
$m_\lambda:\mathscr A_\lambda\otimes\mathscr A_\lambda\to\mathscr A_\lambda$
l'homomorphisme de $\mathcal O_{S_0}$-Modules qui définit la multiplication
dans la $\mathcal O_{S_0}$-Algèbre $\mathscr A_\lambda$~; l'hypothèse sur les
$\varphi_{\mu\lambda}$ entraîne que les $m_\lambda$ forment un système
inductif d'homomorphismes,
\oldpage[IV-3]{8}
et comme le foncteur $\varinjlim$ commute au produit tensoriel,
$m=\varinjlim m_\lambda$ est un homomorphisme
$\mathscr A\otimes\mathscr A\to\mathscr A$ de
$\mathcal O_{S_0}$-Modules~; par passage à la limite sur les diagrammes
commutatifs exprimant l'associativité de $m_\lambda$ et l'existence de la
section unité dans $\mathscr A_\lambda$, on voit que $m$ définit sur
$\mathscr A$ une structure de $\mathcal O_{S_0}$-Algèbre et que
$\varphi_\lambda$ est un homomorphisme de $\mathcal O_{S_0}$-Algèbres pour
tout $\lambda\in L$. En outre $\mathscr A$ est limite inductive du système
$(\mathscr A_\lambda,\varphi_{\mu\lambda})$ \emph{dans la catégorie des
$\mathcal O_{S_0}$-Algèbres}, autrement dit, pour toute
$\mathcal O_{S_0}$-Algèbre $\mathscr B$, l'application canonique
\begin{equation}
\operatorname{Hom}_{\mathcal O_{S_0}\text{-Alg.}}(\mathscr A,\mathscr B)
\longrightarrow
\varprojlim_\lambda
\operatorname{Hom}_{\mathcal O_{S_0}\text{-Alg.}}
(\mathscr A_\lambda,\mathscr B)
\tag{8.2.1.1}\label{IV.8.2.1.1-fr}
\end{equation}
qui à tout homomorphisme $f:\mathscr A\to\mathscr B$ de
$\mathcal O_{S_0}$-Algèbres, fait correspondre la famille
$(f\circ\varphi_\lambda)$, est \emph{bijective}. En effet, on sait déjà
qu'elle est injective et identifie
$\operatorname{Hom}_{\mathcal O_{S_0}\text{-Alg.}}(\mathscr A,\mathscr B)$
à une partie de
$\varprojlim_\lambda
\operatorname{Hom}_{\mathcal O_{S_0}\text{-Mod.}}
(\mathscr A_\lambda,\mathscr B)$~; tout revient à voir que si $(f_\lambda)$
est un système inductif d'homomorphismes de $\mathcal O_{S_0}$-Algèbres,
$f_\lambda:\mathscr A_\lambda\to\mathscr B$, sa limite inductive
$f:\mathscr A\to\mathscr B$, qui est par définition un homomorphisme de
$\mathcal O_{S_0}$-\emph{Modules}, est aussi un homomorphisme de
$\mathcal O_{S_0}$-Algèbres~; mais cela résulte du passage à la limite
inductive dans le diagramme commutatif d'homomorphismes de
$\mathcal O_{S_0}$-Modules
\[
\xymatrix{
  \mathscr A_\lambda\otimes\mathscr A_\lambda
    \ar[r]^{f_\lambda\otimes f_\lambda} \ar[d]_{m_\lambda}
    & \mathscr B\otimes\mathscr B \ar[d] \\
  \mathscr A_\lambda \ar[r]_{f_\lambda} & \mathscr B
}
\]
et du fait que le foncteur $\varinjlim$ commute aux produits tensoriels.

On notera enfin que si les $\mathscr A_\lambda$ sont des
$\mathcal O_{S_0}$-Algèbres commutatives, il en est de même de $\mathscr A$.
\end{env}

\begin{env}[8.2.2]
\label{IV.8.2.2-fr}
Supposons maintenant que $S_0$ soit un \emph{préschéma}, et que les
$\mathscr A_\lambda$ soient des $\mathcal O_{S_0}$-Algèbres (commutatives)
\emph{quasi-cohérentes}~; on sait alors que
$\mathscr A=\varinjlim\mathscr A_\lambda$ est une
$\mathcal O_{S_0}$-Algèbre quasi-cohérente (I, 4.1.1). Désignons par
$S_\lambda$ (resp. $S$) le spectre de la $\mathcal O_{S_0}$-Algèbre
$\mathscr A_\lambda$ (resp. $\mathscr A$) (II, 1.3.1), et soient
$u_{\lambda\mu}:S_\mu\to S_\lambda$ (pour $\lambda\leq\mu$) et
$u_\lambda:S\to S_\lambda$ les $S_0$-morphismes correspondant aux
homomorphismes $\varphi_{\mu\lambda}$ et $\varphi_\lambda$ respectivement
(II, 1.2.7)~; il est clair que $(S_\lambda,u_{\lambda\mu})$ est un
\emph{système projectif} dans la catégorie des $S_0$-préschémas. On notera
que les $u_{\lambda\mu}$ et $u_\lambda$ sont des morphismes \emph{affines}
(II, 1.6.2), donc \emph{quasi-compacts} et \emph{séparés}.
\end{env}

\begin{proposition}[8.2.3]
\label{IV.8.2.3-fr}
Avec les notations de (8.2.2), les morphismes $u_\lambda:S\to S_\lambda$
font de $S$ une limite projective du système projectif
$(S_\lambda,u_{\lambda\mu})$ dans la catégorie des préschémas. En outre, si
$h:S_0\to T$ est un morphisme, faisant de tout $S_0$-préschéma un
$T$-préschéma, $S$ est aussi limite projective du système
$(S_\lambda,u_{\lambda\mu})$ dans la catégorie des $T$-préschémas.
\end{proposition}

Prouvons d'abord la seconde assertion de l'énoncé dans le cas $T=S_0$.
\oldpage[IV-3]{9}
Tout revient à démontrer que si $X$ est un $S_0$-préschéma quelconque,
l'application canonique
\begin{equation}
\operatorname{Hom}_{S_0}(X,S)\longrightarrow
\varprojlim_\lambda\operatorname{Hom}_{S_0}(X,S_\lambda)
\tag{8.2.3.1}\label{IV.8.2.3.1-fr}
\end{equation}
qui à tout $S_0$-morphisme $v:X\to S$ fait correspondre la famille
$(u_\lambda\circ v)$, est \emph{bijective}. Or, si $g:X\to S_0$ est le
morphisme structural et si on pose $\mathscr B=g_*(\mathcal O_X)$, qui est
une $\mathcal O_{S_0}$-Algèbre, l'application (8.2.3.1) s'identifie
canoniquement à (8.2.1.1) (II, 1.2.7), et la conclusion résulte donc de ce
qu'on a vu dans (8.2.1).

Les autres assertions de (8.2.3) sont conséquences du lemme général suivant~:

\begin{lemma}[8.2.4]
\label{IV.8.2.4-fr}
Soient $\mathbf C$ une catégorie, $T$ un objet de $\mathbf C$,
$\mathbf C_T$ la sous-catégorie des objets de $\mathbf C$ au-dessus de $T$.
Soit $(S_\lambda,u_{\lambda\mu})$ un système projectif dans $\mathbf C_T$~;
alors toute limite projective de ce système dans $\mathbf C_T$ est aussi
limite projective dans $\mathbf C$, et réciproquement.
\end{lemma}

Soit $f_\lambda:S_\lambda\to T$ le morphisme structural. Supposons que $S$
soit limite projective de $(S_\lambda,u_{\lambda\mu})$ dans $\mathbf C$, et
désignons par $u_\lambda:S\to S_\lambda$ les morphismes canoniques
correspondants. Considérons alors un système projectif de $T$-morphismes
$w_\lambda:Y\to S_\lambda$, où $Y\in\mathbf C_T$. Il existe par hypothèse un
morphisme unique (dans $\mathbf C$) $w:Y\to S$ tel que
$w_\lambda=u_\lambda\circ w$ pour tout $\lambda$. L'hypothèse que les
$u_\lambda$ sont des $T$-morphismes entraîne que les morphismes
$f_\lambda\circ u_\lambda:S\to T$ sont tous les mêmes, et ce morphisme $f$
fait donc de $S$ un $T$-objet. Si $g:Y\to T$ est le morphisme structural de
$Y$, on a alors
$f\circ w=f_\lambda\circ u_\lambda\circ w=f_\lambda\circ w_\lambda=g$ pour
tout $\lambda$, ce qui prouve que $w$ est un $T$-morphisme. Inversement,
supposons (avec les mêmes notations) que $S$ soit limite projective de
$(S_\lambda,u_{\lambda\mu})$ dans $\mathbf C_T$, et considérons maintenant un
système projectif de morphismes (de $\mathbf C$) $w_\lambda:Y\to S_\lambda$.
Les morphismes composés $f_\lambda\circ w_\lambda:Y\to T$ sont alors tous
les mêmes~: en effet, pour deux indices quelconques $\lambda$, $\mu$, il y a
un indice $\nu$ tel que $\lambda\leq\nu$ et $\mu\leq\nu$, d'où
$f_\nu=f_\lambda\circ u_{\lambda\nu}=f_\mu\circ u_{\mu\nu}$~; comme
$w_\lambda=u_{\lambda\nu}\circ w_\nu$ et
$w_\mu=u_{\mu\nu}\circ w_\nu$, on a
$f_\lambda\circ w_\lambda=f_\lambda\circ u_{\lambda\nu}\circ w_\nu
=f_\nu\circ w_\nu$, et on voit de même que
$f_\mu\circ w_\mu=f_\nu\circ w_\nu$. Si $g:Y\to T$ est l'unique morphisme
ainsi défini, $g$ fait de $Y$ un $T$-objet, et les $w_\lambda$ sont alors des
$T$-morphismes~; ils ont par suite une limite projective $w:Y\to S$ qui est
un $T$-morphisme, et \emph{a fortiori} un morphisme de $\mathbf C$~;
d'ailleurs la première partie du raisonnement montre que toute limite
projective $w'$ (dans $\mathbf C$) du système projectif $(w_\lambda)$ est
nécessairement aussi un $T$-morphisme, donc égale à $w$, ce qui achève de
prouver le lemme.

\begin{proposition}[8.2.5]
\label{IV.8.2.5-fr}
Sous les conditions de (8.2.2), soient $\alpha$ un élément de $L$,
$X_\alpha$ un $S_\alpha$-préschéma. Pour tout $\lambda\geq\alpha$, posons
$X_\lambda=X_\alpha\times_{S_\alpha}S_\lambda$, et pour
$\alpha\leq\lambda\leq\mu$, posons
$v_{\lambda\mu}=1_{X_\alpha}\times u_{\lambda\mu}$, de sorte que
$(X_\lambda,v_{\lambda\mu})$ est un système projectif de
$X_\alpha$-préschémas, dont l'ensemble d'indices est formé des
$\lambda\geq\alpha$ dans $L$. Posons de même
$X=X_\alpha\times_{S_\alpha}S$ et $v_\lambda=1_{X_\alpha}\times u_\lambda$.
Alors les $X_\alpha$-morphismes $v_\lambda:X\to X_\lambda$ font de $X$ une
limite projective du système projectif $(X_\lambda,v_{\lambda\mu})$ dans la
catégorie des $X_\alpha$-préschémas, ou dans la catégorie de tous les
préschémas.
\end{proposition}

Cela résultera encore du lemme général suivant~:

\begin{lemma}[8.2.6]
\label{IV.8.2.6-fr}
Soient $\mathbf C$ une catégorie dans laquelle les produits fibrés existent,
$q:T'\to T$ un morphisme de $\mathbf C$, $\mathbf C_T$ (resp.
$\mathbf C_{T'}$) la catégorie des objets de $\mathbf C$ au-dessus
\oldpage[IV-3]{10}
de $T$ (resp. $T'$). Soit $(S_\lambda,u_{\lambda\mu})$ un système projectif
(non nécessairement filtrant) dans $\mathbf C_T$, et posons
$S'_\lambda=S_\lambda\times_T T'$,
$u'_{\lambda\mu}=u_{\lambda\mu}\times 1_{T'}$, de sorte que
$(S'_\lambda,u'_{\lambda\mu})$ est un système projectif dans
$\mathbf C_{T'}$. Alors, si $(S,u_\lambda)$ est une limite projective de
$(S_\lambda,u_{\lambda\mu})$ dans $\mathbf C_T$,
$(S\times_T T',u_\lambda\times 1_{T'})$ est une limite projective de
$(S'_\lambda,u'_{\lambda\mu})$ dans $\mathbf C_{T'}$.
\end{lemma}

On a par hypothèse, pour tout $\lambda$, un diagramme commutatif
\[
\xymatrix{
  S' \ar[r]^{u'_\lambda} \ar[d]_p
    & S'_\lambda \ar[r]^{h'_\lambda} \ar[d]_{p_\lambda}
    & T' \ar[d]^q \\
  S \ar[r]_{u_\lambda} & S_\lambda \ar[r]_{h_\lambda} & T
}
\]
où l'on a posé $S'=S\times_T T'$, $u'_\lambda=u_\lambda\times 1_{T'}$,
$h'_\lambda=h_\lambda\times 1_{T'}$. Soient $Y$ un $T'$-objet,
$g':Y\to T'$ le morphisme correspondant, et considérons un système
projectif de $T'$-morphismes $w'_\lambda:Y\to S'_\lambda$. Alors $Y$ est un
$T$-objet pour le morphisme $g=q\circ g'$, et les
$w_\lambda=p_\lambda\circ w'_\lambda$ sont des $T$-morphismes puisque
$h_\lambda\circ w_\lambda=h_\lambda\circ p_\lambda\circ w'_\lambda
=q\circ h'_\lambda\circ w'_\lambda=q\circ g'$ par hypothèse. En outre, on
vérifie aussitôt que $(w_\lambda)$ est un système projectif. Il existe donc
par hypothèse un $T$-morphisme $w:Y\to S$ et un seul tel que
$u_\lambda\circ w=w_\lambda$ pour tout $\lambda$. Par définition du produit
fibré, il y a un unique $T'$-morphisme $w':Y\to S'$ tel que
$p\circ w'=w$. On a alors
$u_\lambda\circ p\circ w'=u_\lambda\circ w=w_\lambda
=p_\lambda\circ w'_\lambda$, ce qui s'écrit aussi
$p_\lambda\circ u'_\lambda\circ w'=p_\lambda\circ w'_\lambda$~; d'autre
part, en écrivant que $w'$ est un $T'$-morphisme, il vient
$h'_\lambda\circ u'_\lambda\circ w'=g'=h'_\lambda\circ w'_\lambda$. La
définition de $S'_\lambda$ comme produit fibré $S_\lambda\times_T T'$ donne
donc $u'_\lambda\circ w'=w'_\lambda$, et il est immédiat que $w'$ est
l'unique $T'$-morphisme vérifiant ces relations, d'où le lemme.

\begin{remark}[8.2.7]
\label{IV.8.2.7-fr}
Étant donné un espace annelé quelconque $S$, les limites inductives par
rapport à un ensemble préordonné \emph{quelconque} $L$ (non nécessairement
filtrant) existent dans la catégorie des $\mathcal O_S$-Algèbres
commutatives, puisque la limite inductive filtrante existe d'après (8.2.1) et
d'autre part, pour deux homomorphismes de $\mathcal O_S$-Algèbres
$\mathscr A\to\mathscr B$, $\mathscr A\to\mathscr C$, le produit tensoriel
$\mathscr B\otimes_{\mathscr A}\mathscr C$ est la «~somme amalgamée~»
correspondante dans cette catégorie.

Lorsque $S$ est un préschéma, on sait que le produit tensoriel
$\mathscr B\otimes_{\mathscr A}\mathscr C$ est une $\mathcal O_S$-Algèbre
quasi-cohérente lorsqu'il en est ainsi de $\mathscr A$, $\mathscr B$ et
$\mathscr C$ (I, 1.3.13)~; on en conclut que, dans la catégorie des
$\mathcal O_S$-Algèbres \emph{quasi-cohérentes}, les limites inductives pour
un ensemble préordonné quelconque d'indices existent toujours. Cela permet de
généraliser la définition d'une limite projective de préschémas et les prop.
(8.2.3) et (8.2.5) au cas où l'ensemble préordonné $L$ n'est pas
nécessairement filtrant.
\end{remark}

\begin{env}[8.2.8]
\label{IV.8.2.8-fr}
Les notations étant celles de (8.2.2), posons
$u_{\lambda\mu}=(\psi_{\lambda\mu},\theta_{\lambda\mu})$ et
$u_\lambda=(\psi_\lambda,\theta_\lambda)$~; $(S_\lambda,\psi_{\lambda\mu})$
est donc un système projectif d'espaces topologiques et $(\psi_\lambda)$ un
système projectif d'applications continues $S\to S_\lambda$, des espaces
sous-jacents respectivement aux préschémas $S$ et $S_\lambda$.
\end{env}

\begin{proposition}[8.2.9]
\label{IV.8.2.9-fr}
Avec les notations de (8.2.8), la limite projective du système projectif
$(\psi_\lambda)$ d'applications continues est un homéomorphisme de l'espace
sous-jacent à $S$ sur la limite projective du système projectif
$(S_\lambda,\psi_{\lambda\mu})$ d'espaces topologiques.
\end{proposition}

\oldpage[IV-3]{11}
Soit $T$ l'espace topologique limite du système
$(S_\lambda,\psi_{\lambda\mu})$ et posons
$\psi=\varprojlim\psi_\lambda:S\to T$. On peut se borner au cas où
$S_0=S_\alpha$ pour un $\alpha\in L$, et $\lambda\geq\alpha$.

\medskip
\noindent
(i) Montrons en premier lieu que la topologie de $S$ est image réciproque par
$\psi$ de la topologie de $T$~; autrement dit, si
$\pi_\lambda:T\to S_\lambda$ est l'application canonique, il faut montrer
que tout ouvert de $S$ est réunion d'ouverts de la forme
$\psi^{-1}(\pi_\lambda^{-1}(U_\lambda))$, où $U_\lambda$ est ouvert dans
$S_\lambda$. Or tout ouvert de $S$ est par définition réunion d'ouverts
obtenus de la façon suivante~: on considère un ouvert affine $U_0$ de $S_0$,
d'anneau $A_0$, de sorte que $\psi_\alpha^{-1}(U_0)$ est l'ouvert affine de
$S$ d'anneau $A=\Gamma(U_0,\mathscr A)$, puis on prend un élément $f\in A$ et
on considère dans $\operatorname{Spec}(A)$, identifié à
$\psi_\alpha^{-1}(U_0)$, l'ouvert $D(f)$. Ce sont ces ouverts $D(f)$ qui
forment une base de la topologie de $S$ (II, 1.3.1). Or, si on pose
$A_\lambda=\Gamma(U_0,\mathscr A_\lambda)$, on a
$A=\varinjlim A_\lambda$ (I, 1.3.9), donc il existe un indice $\lambda$ tel
que $f$ soit l'image canonique d'un élément $f_\lambda\in A_\lambda$~; on a
alors $D(f)=\psi_\lambda^{-1}(D(f_\lambda))$ (I, 1.2.2), et comme
$\psi_\lambda=\pi_\lambda\circ\psi$, notre assertion est démontrée.

\medskip
\noindent
(ii) Prouvons maintenant que $\psi$ est \emph{bijective}, ce qui achèvera la
démonstration. Comme $S$ est un espace de Kolmogoroff, il résulte déjà de (i)
que $\psi$ est injective, et il reste donc à montrer que $\psi$ est
surjective. On peut évidemment remplacer pour cela $T$ par un ouvert
$\pi_\alpha^{-1}(U_0)$, où $U_0$ est ouvert affine dans $S_\alpha=S_0$, donc
on peut se limiter au cas où les $S_\lambda$ et $S$ sont affines, autrement
dit $\mathscr A_\lambda$ est le faisceau associé à une $A_0$-algèbre
$A_\lambda$, et $\mathscr A$ le faisceau d'algèbres associé à
$A=\varinjlim A_\lambda$~; nous noterons encore
$\varphi_{\mu\lambda}:A_\lambda\to A_\mu$ et
$\varphi_\lambda:A_\lambda\to A$ les homomorphismes canoniques. Par
définition, un élément de $T$ est une famille
$(\mathfrak p_\lambda)_{\lambda\in L}$, où $\mathfrak p_\lambda$ est un
idéal premier de $A_\lambda$ et où l'on a
$\mathfrak p_\lambda=\varphi_{\mu\lambda}^{-1}(\mathfrak p_\mu)$ pour
$\lambda\leq\mu$. On sait alors (5.13.3 et 5.13.1) qu'il y a un idéal premier
$\mathfrak p$ de $A$ tel que
$\mathfrak p_\lambda=\varphi_\lambda^{-1}(\mathfrak p)$ pour tout
$\lambda\in L$, ce qui achève la démonstration.

En particulier, on a ainsi démontré le

\begin{corollary}[8.2.10]
\label{IV.8.2.10-fr}
Soit $(A_\lambda)_{\lambda\in L}$ un système inductif filtrant d'anneaux, et
soient $A=\varinjlim A_\lambda$, $\varphi_\lambda:A_\lambda\to A$ les
homomorphismes canoniques. L'application canonique
$\mathfrak p\mapsto(\varphi_\lambda^{-1}(\mathfrak p))$ est un
homéomorphisme de $\operatorname{Spec}(A)$ sur l'espace topologique
$\varprojlim\operatorname{Spec}(A_\lambda)$.
\end{corollary}

\begin{corollary}[8.2.11]
\label{IV.8.2.11-fr}
Avec les notations de (8.2.8), pour tout ouvert quasi-compact $U$ de $S$, il
existe un indice $\lambda$ et un ouvert quasi-compact $U_\lambda$ de
$S_\lambda$ tel que $U=\psi_\lambda^{-1}(U_\lambda)$.
\end{corollary}

Cela résulte de ce que, par définition de la topologie limite projective, les
$\psi_\lambda^{-1}(U_\lambda)$ ($U_\lambda$ ouvert quasi-compact de
$S_\lambda$) forment une base de la topologie de $S$, et de ce que l'ensemble
d'indices $L$ est filtrant.

\begin{corollary}[8.2.12]
\label{IV.8.2.12-fr}
Avec les notations de (8.2.8), la limite inductive du système inductif
d'homomorphismes
$\theta_\lambda^\sharp:\psi_\lambda^*(\mathcal O_{S_\lambda})\to
\mathcal O_S$ de faisceaux d'anneaux sur $S$ est un isomorphisme
\begin{equation}
\varinjlim_\lambda\psi_\lambda^*(\mathcal O_{S_\lambda})
\xrightarrow{\sim}\mathcal O_S.
\tag{8.2.12.1}\label{IV.8.2.12.1-fr}
\end{equation}
\end{corollary}

On peut évidemment supposer les $S_\lambda$ affines~; avec les notations de
la démonstration de (8.2.9), tout revient à voir que la limite inductive du
système d'applications canoniques
$(A_\lambda)_{\mathfrak p_\lambda}\to A_{\mathfrak p}$ est un isomorphisme,
ce qui n'est autre que (5.13.3, (ii)).

\oldpage[IV-3]{12}
\begin{proposition}[8.2.13]
\label{IV.8.2.13-fr}
Supposons que les morphismes $u_{\lambda\mu}:S_\mu\to S_\lambda$ soient des
immersions ouvertes, de sorte que $S_\mu$ s'identifie à un sous-préschéma
induit sur un ouvert de $S_\lambda$ pour $\lambda\leq\mu$. Alors, pour tout
$\alpha\in L$, l'espace sous-jacent au préschéma $S$ s'identifie au
sous-espace de $S_\alpha$ intersection des $S_\lambda$ pour
$\lambda\geq\alpha$, et le faisceau structural $\mathcal O_S$ au faisceau
induit $(\mathrm G,\mathrm{II},1.5)$ par $\mathcal O_{S_\alpha}$ sur cette
intersection~; plus généralement, pour tout $\mathcal O_{S_\alpha}$-Module
$\mathscr F_\alpha$, $u_\alpha^*(\mathscr F_\alpha)$ s'identifie au
$\mathcal O_S$-Module induit par $\mathscr F_\alpha$ sur $S$.
\end{proposition}

La première assertion résulte de (8.2.9), vu la définition d'une limite
projective d'espaces topologiques~; en outre tous les
$\psi_\lambda^*(\mathcal O_{S_\lambda})$ sont égaux au faisceau induit par
$\mathcal O_{S_\alpha}$ sur $S$ par définition $(0_{\mathrm I},3.7.1)$ et,
avec les notations de la démonstration de (8.2.9), les homomorphismes
$(A_\lambda)_{\mathfrak p_\lambda}\to
(A_\mu)_{\mathfrak p_\mu}$ sont bijectifs pour un système
$(\mathfrak p_\lambda)$ d'idéaux premiers correspondant à un même point de
$S$~; l'assertion relative à $\mathcal O_S$ se déduit donc de (8.2.12). La
dernière assertion résulte alors de la définition de
$u_\alpha^*(\mathscr F_\alpha)$ $(0_{\mathrm I},4.3.1)$.

\begin{remark}[8.2.14]
\label{IV.8.2.14-fr}
Les résultats de (8.2.9) et (8.2.12) montrent que $S$ est limite projective du
système projectif $(S_\lambda,u_{\lambda\mu})$ dans la catégorie de
\emph{tous} les espaces annelés (ou de tous les espaces annelés en anneaux
locaux). En effet, soit $Y$ un espace annelé, et considérons un système
projectif de morphismes d'espaces annelés $w_\lambda:Y\to S_\lambda$. Si
l'on pose $w_\lambda=(\rho_\lambda,\omega_\lambda)$, les $\rho_\lambda$
forment un système projectif d'applications continues et, en vertu de
(8.2.9), leur limite projective $\rho$ s'identifie à une application continue
$Y\to S$ telle que $\rho_\lambda=\psi_\lambda\circ\rho$. D'autre part, les
\[
\omega_\lambda^\sharp:
\rho_\lambda^*(\mathcal O_{S_\lambda})\to\mathcal O_Y
\]
forment un système inductif d'homomorphismes de faisceaux d'anneaux~; comme
on peut écrire
$\rho_\lambda^*(\mathcal O_{S_\lambda})=
\rho^*(\psi_\lambda^*(\mathcal O_{S_\lambda}))$ et que le foncteur $\rho^*$
est exact, la limite inductive des
$\rho_\lambda^*(\mathcal O_{S_\lambda})$ est $\rho^*(\mathcal O_S)$ en
vertu de (8.2.12), et il y a donc un homomorphisme unique
$\omega^\sharp:\rho^*(\mathcal O_S)\to\mathcal O_Y$ tel que
$\omega_\lambda^\sharp=\omega^\sharp\circ
\rho^*(\theta_\lambda^\sharp)$, ce qui démontre notre assertion.
\end{remark}

\subsection{Parties constructibles dans une limite projective de préschémas}
\label{subsection:IV.8.3-fr}

\begin{env}[8.3.1]
\label{IV.8.3.1-fr}
Dans toute la suite de ce paragraphe, nous supposons remplies les conditions
de (8.2.2), dont nous conservons les notations.
\end{env}

\begin{theorem}[8.3.2]
\label{IV.8.3.2-fr}
Pour tout $\lambda$, soient $E_\lambda$, $F_\lambda$ deux parties de
$S_\lambda$. Posons
\begin{equation}
E=\bigcap_\lambda u_\lambda^{-1}(E_\lambda),\qquad
F=\bigcup_\lambda u_\lambda^{-1}(F_\lambda).
\tag{8.3.2.1}\label{IV.8.3.2.1-fr}
\end{equation}
On suppose vérifiées les conditions suivantes~:
\begin{enumerate}
\item[(i)] Pour tout $\lambda$, $E_\lambda$ est pro-constructible et
$F_\lambda$ est ind-constructible (1.9.4).
\item[(ii)] Pour $\lambda\leq\mu$, on a
\[
E_\mu\subset u_{\lambda\mu}^{-1}(E_\lambda),\qquad
F_\mu\supset u_{\lambda\mu}^{-1}(F_\lambda).
\]
\item[(iii)] Il existe $\alpha\in L$ tel que $S_\alpha$ soit quasi-compact
(ce qui entraîne que $S_\lambda$ est quasi-compact pour
$\lambda\geq\alpha$).
\end{enumerate}
Alors, les propriétés suivantes sont équivalentes~:
\begin{enumerate}
\item[\emph{a)}] $E\subset F$.
\oldpage[IV-3]{13}
\item[\emph{b)}] Il existe $\lambda\geq\alpha$ tel que
$u_\lambda^{-1}(E_\lambda)\subset u_\lambda^{-1}(F_\lambda)$ (et alors on a
$u_\mu^{-1}(E_\mu)\subset u_\mu^{-1}(F_\mu)$ pour $\mu\geq\lambda$).
\item[\emph{c)}] Il existe $\lambda\geq\alpha$ tel que
$E_\lambda\subset F_\lambda$ (et alors on a $E_\mu\subset F_\mu$ pour
$\mu\geq\lambda$).
\end{enumerate}
\end{theorem}

Les remarques entre parenthèses dans \emph{b)} et \emph{c)} résultent de (ii).
Posons
\[
G_\lambda=E_\lambda\cap(S_\lambda-F_\lambda),\qquad
G=E\cap(S-F).
\]
Alors $G_\lambda$ est une partie pro-constructible de $S_\lambda$ (1.9.5,
(i)), et en vertu de (8.3.2.1) et de (ii), on a
\[
G_\mu\subset u_{\lambda\mu}^{-1}(G_\lambda)
\quad\text{pour }\lambda\leq\mu,
\qquad
G=\bigcap_\lambda u_\lambda^{-1}(G_\lambda).
\]
On est donc ramené à prouver le cas particulier de (8.3.2) correspondant à
$F_\lambda=\emptyset$ pour tout $\lambda$~:

\begin{corollary}[8.3.3]
\label{IV.8.3.3-fr}
Pour tout $\lambda$, soit $E_\lambda$ une partie pro-constructible de
$S_\lambda$ telle que, pour $\lambda\leq\mu$, on ait
$E_\mu\subset u_{\lambda\mu}^{-1}(E_\lambda)$. Supposons qu'il existe
$\alpha\in L$ tel que $S_\alpha$ soit quasi-compact. Alors les conditions
suivantes sont équivalentes~:
\begin{enumerate}
\item[\emph{a)}]
$E=\bigcap_\lambda u_\lambda^{-1}(E_\lambda)=\emptyset$.
\item[\emph{b)}] Il existe un $\lambda$ tel que
$u_\lambda^{-1}(E_\lambda)=\emptyset$ (et alors
$u_\mu^{-1}(E_\mu)=\emptyset$ pour $\mu\geq\lambda$).
\item[\emph{c)}] Il existe $\lambda$ tel que $E_\lambda=\emptyset$ (et alors
$E_\mu=\emptyset$ pour $\mu\geq\lambda$).
\end{enumerate}
\end{corollary}

Il est clair que \emph{c)} implique \emph{a)}. Prouvons que \emph{a)} entraîne
\emph{b)}~: puisque $S_\lambda$ est quasi-compact, il en est de même de $S$
(8.2.2)~; $E_\lambda$ étant pro-constructible, il en est de même de
$u_\lambda^{-1}(E_\lambda)$ (1.9.5, (vi))~; alors la famille filtrante
décroissante d'ensembles pro-constructibles $u_\lambda^{-1}(E_\lambda)$ a
une intersection vide, donc (1.9.9) l'un d'eux est vide.

Enfin, montrons que \emph{b)} entraîne \emph{c)}. Comme $S_\alpha$ est
quasi-compact et $L$ filtrant, on peut remplacer $S_\alpha$ par un ouvert
affine, donc supposer (8.2.2) que $S$ et les $S_\lambda$ pour
$\lambda\geq\alpha$ sont affines~; on a alors (1.9.2.1), pour
$\lambda\geq\alpha$,
\[
u_\lambda(S)=\bigcap_{\mu\geq\lambda}u_{\lambda\mu}(S_\mu),
\]
d'où
\[
E_\lambda\cap u_\lambda(S)=
\bigcap_{\mu\geq\lambda}
\bigl(E_\lambda\cap u_{\lambda\mu}(S_\mu)\bigr).
\]
Or, comme $u_\lambda$ et les $u_{\lambda\mu}$ sont quasi-compacts,
$u_\lambda(S)$ et $u_{\lambda\mu}(S_\mu)$ sont pro-constructibles dans
$S_\lambda$ (1.9.5, (vii)), donc les ensembles
$E_\lambda\cap u_{\lambda\mu}(S_\mu)$ pour $\mu\geq\lambda$ forment une
famille filtrante décroissante de parties pro-constructibles de $S_\lambda$.
Comme $S_\lambda$ est quasi-compact, l'hypothèse \emph{b)} entraîne que
l'intersection de cette famille est vide, donc (1.9.9) un des ensembles
$E_\lambda\cap u_{\lambda\mu}(S_\mu)$ est vide, donc
$E_\mu\subset u_{\lambda\mu}^{-1}(E_\lambda)$ est vide. C.Q.F.D.

\begin{corollary}[8.3.4]
\label{IV.8.3.4-fr}
Pour tout $\lambda$, soit $F_\lambda$ une partie ind-constructible de
$S_\lambda$ telle que pour $\lambda\leq\mu$ on ait
$u_{\lambda\mu}^{-1}(F_\lambda)\subset F_\mu$. Supposons qu'il existe
$\alpha\in L$ tel que $S_\alpha$ soit quasi-compact. Alors les conditions
suivantes sont équivalentes~:
\begin{enumerate}
\item[\emph{a)}] L'ensemble
$F=\bigcup_\lambda u_\lambda^{-1}(F_\lambda)$ est égal à $S$.
\item[\emph{b)}] Il existe un $\lambda$ tel que
$u_\lambda^{-1}(F_\lambda)=S$ (et alors
$u_\mu^{-1}(F_\mu)=S$ pour $\mu\geq\lambda$).
\oldpage[IV-3]{14}
\item[\emph{c)}] Il existe un $\lambda$ tel que $F_\lambda=S_\lambda$ (et
alors $F_\mu=S_\mu$ pour $\mu\geq\lambda$).
\end{enumerate}
\end{corollary}

Cela se déduit aussitôt de (8.3.3) par passage aux complémentaires.

\begin{corollary}[8.3.5]
\label{IV.8.3.5-fr}
Pour tout $\lambda$, soient $E_\lambda$, $F_\lambda$ deux parties
constructibles de $S_\lambda$ telles que, pour $\mu\geq\lambda$, on ait
$E_\mu\subset u_{\lambda\mu}^{-1}(E_\lambda)$ et
$F_\mu\supset u_{\lambda\mu}^{-1}(F_\lambda)$. Supposons qu'il existe un
indice $\alpha$ tel que $S_\alpha$ soit quasi-compact. Alors, pour que
$E\subset F$ (resp. $E=F$), il faut et il suffit qu'il existe $\lambda$ tel
que $E_\lambda\subset F_\lambda$ (resp. $E_\lambda=F_\lambda$), auquel cas
on a aussi $E_\mu\subset F_\mu$ (resp. $E_\mu=F_\mu$) pour
$\mu\geq\lambda$.
\end{corollary}

C'est un cas particulier de (8.3.2).

En particulier~:

\begin{corollary}[8.3.6]
\label{IV.8.3.6-fr}
Supposons qu'il existe un $\alpha$ tel que $S_\alpha$ soit quasi-compact.
Pour que $S=\emptyset$, il faut et il suffit qu'il existe un $\lambda$ tel
que $S_\lambda=\emptyset$.
\end{corollary}

\begin{corollary}[8.3.7]
\label{IV.8.3.7-fr}
On a, pour tout $\lambda$,
\begin{equation}
u_\lambda(S)=\bigcap_{\mu\geq\lambda}u_{\lambda\mu}(S_\mu).
\tag{8.3.7.1}\label{IV.8.3.7.1-fr}
\end{equation}
\end{corollary}

Il est clair que le premier membre de (8.3.7.1) est contenu dans le second.
Soit $s$ un point de $S_\lambda$ et posons
$X_\lambda=\operatorname{Spec}(k(s))$~; considérons le système projectif
$(X_\mu,v_{\mu\nu})$ où
$X_\mu=X_\lambda\times_{S_\lambda}S_\mu$ et
$v_{\mu\nu}=1\times u_{\mu\nu}$ pour
$\lambda\leq\mu\leq\nu$~; sa limite projective est
$X=X_\lambda\times_{S_\lambda}S$ et $v_\lambda=1\times u_\lambda$ est
l'application canonique $X\to X_\lambda$ (8.2.5). Si
$s\in\bigcap_{\mu\geq\lambda}u_{\lambda\mu}(S_\mu)$, cela entraîne que
$X_\mu\neq\emptyset$ pour tout $\mu\geq\lambda$ (I, 3.4.8)~; il résulte
alors de (8.3.6) que $X\neq\emptyset$, donc $s\in u_\lambda(S)$ par
(I, 3.4.8).

\begin{proposition}[8.3.8]
\label{IV.8.3.8-fr}
\begin{enumerate}
\item[(i)] Pour que le morphisme $u_\lambda:S\to S_\lambda$ soit dominant
(resp. surjectif), il faut et il suffit que pour $\mu\geq\lambda$ le
morphisme $u_{\lambda\mu}:S_\mu\to S_\lambda$ soit dominant (resp.
surjectif).
\item[(ii)] Si, pour un indice $\lambda$, les morphismes
$u_{\lambda\mu}:S_\mu\to S_\lambda$ sont plats (resp. fidèlement plats)
pour tout $\mu\geq\lambda$, alors le morphisme
$u_\lambda:S\to S_\lambda$ est plat (resp. fidèlement plat).
\item[(iii)] Supposons que les morphismes $u_\mu:S\to S_\mu$ soient
surjectifs pour $\mu\geq\lambda$. Pour que $u_\lambda$ soit un morphisme
ouvert (resp. universellement ouvert), il faut et il suffit que, pour tout
$\mu\geq\lambda$, $u_{\lambda\mu}$ soit un morphisme ouvert (resp.
universellement ouvert).
\end{enumerate}
\end{proposition}

\noindent
(i) Comme $u_\lambda(S)\subset u_{\lambda\mu}(S_\mu)$ pour
$\mu\geq\lambda$, la nécessité des conditions est triviale, et il résulte
aussitôt de (8.3.7.1) que si les $u_{\lambda\mu}$ sont surjectifs, il en est
de même de $u_\lambda$. Supposons maintenant les $u_{\lambda\mu}$ dominants
pour $\mu\geq\lambda$, et considérons dans $S_\lambda$ un ouvert
quasi-compact non vide $U_\lambda$~; alors les
$U_\mu=u_{\lambda\mu}^{-1}(U_\lambda)$ pour $\mu\geq\lambda$ forment un
système projectif dont la limite projective est
$U=u_\lambda^{-1}(U_\lambda)$ (8.2.5). Par hypothèse les $U_\mu$ sont tous
non vides, donc il en est de même de $U$ par (8.3.6), ce qui prouve que
$u_\lambda$ est dominant.

\medskip
\noindent
(ii) En vertu de (i), il suffit de considérer le cas où les
$u_{\lambda\mu}$ sont plats~; on peut alors se borner au cas où $S_\lambda$
est affine, donc aussi les $S_\mu$ pour $\mu\geq\lambda$ et $S$, et
l'assertion résulte alors de (2.1.2) et $(0_{\mathrm I},6.2.3)$.

\medskip
\noindent
(iii) En vertu de (8.2.5) et de (I, 3.5.2), il suffit de traiter le cas des
morphismes ouverts. Comme $u_\lambda=u_{\lambda\mu}\circ u_\mu$ et que
$u_\mu$ est surjectif, on sait que si $u_\lambda$ est ouvert il en est de
même de $u_{\lambda\mu}$ pour $\mu\geq\lambda$ (Bourbaki,
\emph{Top. gén.}, chap. $\mathrm I^{\mathrm{er}}$, $3^{\mathrm e}$ éd.,
\S~5, n\textsuperscript{o} 1, prop. 1).

\oldpage[IV-3]{15}
Inversement, pour montrer que $u_\lambda$ est ouvert lorsque tous les
$u_{\lambda\mu}$ le sont pour $\mu\geq\lambda$, il suffit de voir que pour
tout ouvert quasi-compact $V$ de $S$, $u_\lambda(V)$ est ouvert dans
$S_\lambda$~; mais il existe alors $\mu\geq\lambda$ tel que
$V=u_\mu^{-1}(V_\mu)$, où $V_\mu$ est ouvert dans $S_\mu$ (8.2.11)~; comme
$u_\mu$ est surjectif, on a $V_\mu=u_\mu(V)$ et
$u_\lambda(V)=u_{\lambda\mu}(u_\mu(V))$ est donc ouvert par hypothèse.

On notera qu'il peut se faire que tous les $u_{\lambda\mu}$ soient ouverts
sans qu'il en soit ainsi de $u_\lambda$ si les $u_\mu$ ne sont pas
surjectifs. On en a un exemple en considérant un anneau intègre $A$ qui n'est
pas un corps, et son corps des fractions $K$, qui est limite inductive des
$A_f$, où $f$ parcourt $A-\{0\}$~; si l'on pose
$S_1=\operatorname{Spec}(A)$, $S_f=\operatorname{Spec}(A_f)$, on a
$S=\varprojlim S_f=\operatorname{Spec}(K)$, et le morphisme $S\to S_1$ n'est
pas ouvert, bien que les morphismes $S_f\to S_1$ le soient.

\begin{env}[8.3.9]
\label{IV.8.3.9-fr}
Pour tout préschéma $X$, nous noterons comme d'ordinaire $\mathfrak P(X)$
l'ensemble des parties de l'ensemble sous-jacent à $X$, par $\mathfrak C(X)$
(resp. $\mathfrak D_c(X)$, $\mathfrak F_c(X)$, $\mathfrak{LF}_c(X)$)
l'ensemble des parties constructibles (resp. constructibles et ouvertes,
resp. constructibles et fermées, resp. constructibles et localement fermées)
de $X$. Il est clair que
$(\mathfrak P(S_\lambda),u_{\lambda\mu}^{-1})$ est un \emph{système inductif
d'ensembles} et que les applications
$u_\lambda^{-1}:\mathfrak P(S_\lambda)\to\mathfrak P(S)$ forment un système
inductif d'applications, d'où, en passant à la limite inductive, une
application canonique
\begin{equation}
u_{\mathfrak P}:\varinjlim\mathfrak P(S_\lambda)\to\mathfrak P(S).
\tag{8.3.9.1}\label{IV.8.3.9.1-fr}
\end{equation}
En outre, il résulte de (1.8.2) que $u_{\lambda\mu}^{-1}$ applique
$\mathfrak C(S_\lambda)$ (resp. $\mathfrak D_c(S_\lambda)$,
$\mathfrak F_c(S_\lambda)$, $\mathfrak{LF}_c(S_\lambda)$) dans
$\mathfrak C(S_\mu)$ (resp. $\mathfrak D_c(S_\mu)$,
$\mathfrak F_c(S_\mu)$, $\mathfrak{LF}_c(S_\mu)$) et que
$u_\lambda^{-1}$ applique $\mathfrak C(S_\lambda)$ (resp.
$\mathfrak D_c(S_\lambda)$, $\mathfrak F_c(S_\lambda)$,
$\mathfrak{LF}_c(S_\lambda)$) dans $\mathfrak C(S)$ (resp.
$\mathfrak D_c(S)$, $\mathfrak F_c(S)$, $\mathfrak{LF}_c(S)$). On a donc par
restriction de (8.3.9.1), des applications canoniques
\begin{align}
\varinjlim\mathfrak C(S_\lambda)&\longrightarrow\mathfrak C(S)
\tag{8.3.9.2}\label{IV.8.3.9.2-fr}\\
\varinjlim\mathfrak D_c(S_\lambda)&\longrightarrow\mathfrak D_c(S)
\tag{8.3.9.3}\label{IV.8.3.9.3-fr}\\
\varinjlim\mathfrak F_c(S_\lambda)&\longrightarrow\mathfrak F_c(S)
\tag{8.3.9.4}\label{IV.8.3.9.4-fr}\\
\varinjlim\mathfrak{LF}_c(S_\lambda)&\longrightarrow\mathfrak{LF}_c(S).
\tag{8.3.9.5}\label{IV.8.3.9.5-fr}
\end{align}
\end{env}

\begin{env}[8.3.10]
\label{IV.8.3.10-fr}
Soit $g_\alpha:X_\alpha\to S_\alpha$ un morphisme~; avec les notations de
(8.2.5) on a comme ci-dessus une application canonique
$v_{\mathfrak P}:\varinjlim\mathfrak P(X_\lambda)\to\mathfrak P(X)$~;
d'autre part, on a des morphismes de projection
$g_\lambda:X_\lambda\to S_\lambda$ pour tout $\lambda\geq\alpha$ et un
morphisme de projection $g:X\to S$. Il est clair que les
$g_\lambda^{-1}:\mathfrak P(S_\lambda)\to\mathfrak P(X_\lambda)$ forment un
système inductif d'applications, et que les diagrammes
\[
\xymatrix{
  \mathfrak P(S_\lambda) \ar[r]^{u_\lambda^{-1}} \ar[d]_{g_\lambda^{-1}}
    & \mathfrak P(S) \ar[d]^{g^{-1}} \\
  \mathfrak P(X_\lambda) \ar[r]_{v_\lambda^{-1}} & \mathfrak P(X)
}
\]
\oldpage[IV-3]{16}
sont commutatifs~; on en déduit donc par passage à la limite inductive un
diagramme commutatif
\begin{equation}
\xymatrix{
  \varinjlim\mathfrak P(S_\lambda) \ar[r]^{u_{\mathfrak P}}
    \ar[d]_{\varinjlim g_\lambda^{-1}}
    & \mathfrak P(S) \ar[d]^{g^{-1}} \\
  \varinjlim\mathfrak P(X_\lambda) \ar[r]_{v_{\mathfrak P}}
    & \mathfrak P(X)
}
\tag{8.3.10.1}\label{IV.8.3.10.1-fr}
\end{equation}
et il résulte de (1.8.2) que l'on a des diagrammes analogues en remplaçant
$\mathfrak P$ par $\mathfrak C$, $\mathfrak D_c$, $\mathfrak F_c$ ou
$\mathfrak{LF}_c$.
\end{env}

Il résulte de (8.3.5) que sous l'hypothèse que pour un $\alpha\in L$,
$S_\alpha$ est \emph{quasi-compact}, l'application canonique (8.3.9.2) est
\emph{injective}. En outre~:

\begin{theorem}[8.3.11]
\label{IV.8.3.11-fr}
Supposons qu'il existe un $\alpha\in L$ tel que $S_\alpha$ soit quasi-compact
et quasi-séparé. Alors les applications canoniques (8.3.9.2), (8.3.9.3),
(8.3.9.4) et (8.3.9.5) sont \emph{bijectives}.
\end{theorem}

En vertu de la remarque précédente, il reste à prouver que ces applications
sont \emph{surjectives}~; comme toute partie constructible de $S$ est réunion
finie d'ensembles de la forme $U\cap\complement V$, où $U$ et $V$ sont
ouverts et constructibles, il suffira de prouver que (8.3.9.3) est surjective
pour qu'il en soit de même de (8.3.9.2) (et aussi de (8.3.9.4), par passage
aux complémentaires). Or, comme les morphismes $S_\lambda\to S_\alpha$ et
$S\to S_\alpha$ sont affines, donc séparés, les $S_\lambda$ pour
$\lambda\geq\alpha$ et $S$ sont quasi-compacts et quasi-séparés (1.2.2), et
l'on sait que les parties ouvertes constructibles dans un tel préschéma ne
sont autres que les parties ouvertes quasi-compactes (1.8.1). La conclusion
résulte donc de (8.2.11) sauf pour l'application (8.3.9.5). Pour prouver que
cette dernière est surjective, considérons une partie $Z$ localement fermée
et constructible dans $S$~; $Z$ est donc quasi-compact
$(0_{\mathrm{III}},9.1.4)$. Comme tout point $x\in Z$ admet par hypothèse un
voisinage ouvert quasi-compact $V_x$ dans $S$ tel que $Z\cap V_x$ soit fermé
dans $V_x$, on peut couvrir $Z$ par un nombre fini des $V_x$~; autrement dit,
il y a un ouvert quasi-compact $U$ contenant $Z$ et tel que $Z$ soit fermé
dans $U$~; comme $Z$ est constructible dans $S$, il l'est aussi dans $U$
$(0_{\mathrm{III}},9.1.8)$. On sait (8.2.11) qu'il y a un indice $\lambda$ et
un ouvert quasi-compact $U_\lambda$ dans $S_\lambda$ tels que
$U=u_\lambda^{-1}(U_\lambda)$. Appliquant à $U$ (qui est limite projective
des $U_\mu=u_{\lambda\mu}^{-1}(U_\lambda)$ pour $\mu\geq\lambda$) le fait
que l'application (8.3.9.4) est surjective, on voit qu'il existe un
$\mu\geq\lambda$ et un ensemble fermé constructible $Z_\mu$ dans $U_\mu$
tels que $Z=u_\mu^{-1}(Z_\mu)$. Mais comme l'immersion canonique
$U_\mu\to S_\mu$ est quasi-compacte par hypothèse (1.2.7), elle est de
présentation finie (1.6.2, (i)), et $Z_\mu$ est aussi une partie constructible
\emph{de} $S_\mu$ en vertu de (1.8.4) et (1.8.1)~; comme $Z_\mu$ est
évidemment localement fermée dans $S_\mu$, cela achève la démonstration.

\begin{corollary}[8.3.12]
\label{IV.8.3.12-fr}
Supposons qu'il existe un $\alpha$ tel que $S_\alpha$ soit quasi-compact, et
soit, pour tout $\lambda$, $Z_\lambda$ une partie constructible de
$S_\lambda$ telle que
$Z_\mu=u_{\lambda\mu}^{-1}(Z_\lambda)$ pour $\mu\geq\lambda$. Si
$Z=u_\lambda^{-1}(Z_\lambda)$,
\oldpage[IV-3]{17}
alors, pour que $Z$ soit ouverte (resp. fermée, resp. localement fermée) dans
$S$, il faut et il suffit qu'il existe $\lambda\geq\alpha$ tel que
$Z_\lambda$ le soit dans $S_\lambda$.
\end{corollary}

Couvrons $S_\alpha$ par un nombre fini d'ouverts affines
$U_\alpha^{(i)}$~; alors les $U^{(i)}=u_\alpha^{-1}(U_\alpha^{(i)})$ forment
un recouvrement ouvert affine de $S$, et pour que $Z$ soit ouvert (resp.
fermé, resp. localement fermé) dans $S$, il faut et il suffit que chacun des
$Z\cap U^{(i)}$ le soit dans $U^{(i)}$. Comme $L$ est filtrant et que chacun
des $Z\cap U^{(i)}$ est constructible dans $U^{(i)}$
$(0_{\mathrm{III}},9.1.8)$, on peut se borner à démontrer le corollaire
lorsque $S_\alpha$ est affine, donc quasi-compact et quasi-séparé~; mais alors
il résulte de (8.3.11).

\begin{proposition}[8.3.13]
\label{IV.8.3.13-fr}
Supposons que les morphismes $u_{\lambda\mu}:S_\mu\to S_\lambda$ soient
plats pour $\lambda\leq\mu$, et qu'il existe $\alpha$ tel que $S_\alpha$ soit
quasi-compact. Pour tout $\lambda$, soient $Z'_\lambda$, $Z''_\lambda$ deux
parties pro-constructibles de $S_\lambda$, telles que
$Z'_\mu=u_{\lambda\mu}^{-1}(Z'_\lambda)$ et
$Z''_\mu=u_{\lambda\mu}^{-1}(Z''_\lambda)$ pour $\mu\geq\lambda$~; on
suppose en outre que $\overline{Z'_\alpha}$ est constructible dans
$S_\alpha$. Soient $Z'=u_\lambda^{-1}(Z'_\lambda)$,
$Z''=u_\lambda^{-1}(Z''_\lambda)$~; pour que
$Z''\subset\overline{Z'}$, il faut et il suffit qu'il existe
$\lambda\geq\alpha$ tel que
$Z''_\lambda\subset\overline{Z'_\lambda}$.
\end{proposition}

En effet, on sait que $u_\lambda:S\to S_\lambda$ est aussi un morphisme plat
pour tout $\lambda$ (8.3.8)~; comme $Z'_\lambda$ est pro-constructible,
l'adhérence de $Z'_\mu$ pour $\mu\geq\lambda$ (resp. de $Z'$) dans $S_\mu$
(resp. $S$) est égale à
$u_{\lambda\mu}^{-1}(\overline{Z'_\lambda})$ (resp.
$u_\lambda^{-1}(\overline{Z'_\lambda})$) (2.3.10). Comme les
$u_{\lambda\mu}^{-1}(\overline{Z'_\lambda})$ et
$u_\lambda^{-1}(\overline{Z'_\lambda})$ sont constructibles (1.8.2), la
conclusion résulte de (8.3.2).

\subsection{Critères d'irréductibilité et de connexion pour les limites
projectives de préschémas}
\label{subsection:IV.8.4-fr}

\begin{proposition}[8.4.1]
\label{IV.8.4.1-fr}
Supposons qu'il existe un indice $\alpha$ tel que $S_\alpha$ soit
quasi-compact.
\begin{enumerate}
\item[(i)] Si $S$ n'est pas irréductible et si de plus l'espace sous-jacent à
$S$ est noethérien et $S_\alpha$ quasi-séparé, il existe
$\lambda\geq\alpha$ tel que, pour $\mu\geq\lambda$, $S_\mu$ ne soit pas
irréductible.
\item[(ii)] Si $S$ n'est pas connexe, il existe $\lambda\geq\alpha$ tel que,
pour $\mu\geq\lambda$, $S_\mu$ ne soit pas connexe.
\end{enumerate}
\end{proposition}

Supposons que $S$ soit réunion de deux ensembles fermés $S'$, $S''$ distincts
de $S$ (resp. fermés disjoints et non vides). Dans le cas (i), $S'$ et $S''$
sont constructibles puisque l'espace $S$ est noethérien. En vertu de
(8.3.11), il existe donc un $\lambda\geq\alpha$ et deux ensembles fermés
constructibles $S'_\lambda$, $S''_\lambda$ de $S_\lambda$ tels que
$S'=u_\lambda^{-1}(S'_\lambda)$,
$S''=u_\lambda^{-1}(S''_\lambda)$~; comme $S=S'\cup S''$, il résulte aussi
de (8.3.11) que l'on peut supposer que
$S_\lambda=S'_\lambda\cup S''_\lambda$~; comme $S'_\lambda$ et
$S''_\lambda$ sont distincts de $S_\lambda$, cela prouve que $S_\lambda$
n'est pas irréductible.

Dans le cas (ii), $S'$ et $S''$ sont ouverts quasi-compacts, donc, en vertu
de (8.2.11), il existe un $\lambda\geq\alpha$ et deux ensembles ouverts
quasi-compacts $S'_\lambda$, $S''_\lambda$ de $S_\lambda$ tels que
$S'=u_\lambda^{-1}(S'_\lambda)$,
$S''=u_\lambda^{-1}(S''_\lambda)$. D'ailleurs, comme $S'$ et $S''$ sont
ouverts et fermés dans $S$, ils sont à la fois pro-constructibles et
ind-constructibles (1.9.6), donc constructibles (1.9.11), et il résulte donc
de (8.3.5) que l'on peut supposer $\lambda$ pris tel que
$S_\lambda=S'_\lambda\cup S''_\lambda$ et
$S'_\lambda\cap S''_\lambda=\emptyset$, ce qui montre que $S_\lambda$
n'est pas connexe.

\begin{proposition}[8.4.2]
\label{IV.8.4.2-fr}
Supposons que l'espace sous-jacent à $S$ soit noethérien et que l'une des deux
conditions suivantes soit vérifiée~:
\begin{enumerate}
\item[\emph{a)}] Pour $\lambda\leq\mu$,
$u_{\lambda\mu}:S_\mu\to S_\lambda$ est dominant, et il existe $\alpha$ tel
que $S_\alpha$ soit quasi-compact.
\oldpage[IV-3]{18}
\item[\emph{b)}] Il existe $\alpha$ tel que l'espace sous-jacent à
$S_\alpha$ soit noethérien, et pour $\mu\geq\lambda$, $u_{\lambda\mu}$ est
un homéomorphisme de $S_\mu$ sur un sous-espace de $S_\lambda$.
\end{enumerate}
Sous ces conditions, il existe un $\lambda$ tel que, pour tout
$\mu\geq\lambda$~:
\begin{enumerate}
\item[(i)] Pour toute composante irréductible $Y_i$ de $S$
$(1\leq i\leq m)$, $\overline{u_\mu(Y_i)}$ est une composante irréductible
de $S_\mu$, et l'application $Y_i\mapsto\overline{u_\mu(Y_i)}$ est une
bijection de l'ensemble des composantes irréductibles de $S$ sur l'ensemble
des composantes irréductibles de $S_\mu$.
\item[(ii)] Pour toute composante connexe $C_j$ de $S$
$(1\leq j\leq n)$, $\overline{u_\mu(C_j)}$ est une composante connexe de
$S_\mu$, et l'application $C_j\mapsto\overline{u_\mu(C_j)}$ est une
bijection de l'ensemble des composantes connexes de $S$ sur l'ensemble des
composantes connexes de $S_\mu$.
\end{enumerate}
\end{proposition}

Nous établirons d'abord le

\begin{lemma}[8.4.2.1]
\label{IV.8.4.2.1-fr}
Sous la condition \emph{a)} ou \emph{b)} de (8.4.2), il existe $\lambda$ tel
que, pour $\mu\geq\lambda$, $u_\mu:S\to S_\mu$ soit dominant.
\end{lemma}

Dans le cas \emph{a)}, cela a déjà été prouvé sans supposer l'espace $S$
noethérien (8.3.8, (i)). Dans le cas \emph{b)}, posons
$Z_\alpha=\overline{u_\alpha(S)}$~; en tant que partie fermée de l'espace
noethérien $S_\alpha$, $Z_\alpha$ est constructible, et comme
$u_\alpha^{-1}(Z_\alpha)=S$, il résulte de (8.3.5) qu'il existe
$\lambda\geq\alpha$ tel que, pour $\mu\geq\lambda$, on ait
$Z_\mu=u_{\alpha\mu}^{-1}(Z_\alpha)=S_\mu$. Mais comme $u_{\alpha\mu}$ est
un homéomorphisme de $S_\mu$ sur un sous-espace de $S_\alpha$, et comme
l'application composée $S\to Z_\mu\to Z_\alpha$ est dominante, il en est de
même de $S\to Z_\mu=S_\mu$.

Ce lemme étant démontré, on peut supposer que pour tout $\lambda$,
$u_\lambda$ est un morphisme dominant.

\medskip
\noindent
(i) Chacun des $S_\lambda$ est réunion des
$\overline{u_\lambda(Y_i)}$, qui sont irréductibles. D'autre part, si $U_i$
est l'ouvert de $S$ complémentaire de la réunion des $Y_j$ d'indice $j\neq i$
$(1\leq i\leq m)$, les $U_i$ sont deux à deux disjoints et
$Y_i=\overline{U_i}$ $(0_{\mathrm I},2.1.6)$. Comme l'espace sous-jacent à
$S$ est noethérien, les $U_i$ sont quasi-compacts, donc il existe un indice
$\lambda$ et des ouverts $U_{i\lambda}$ de $S_\lambda$ tels que
$U_i=u_\lambda^{-1}(U_{i\lambda})$ pour $1\leq i\leq m$ (8.2.11). On en
conclut que si l'on pose
$U_{i\mu}=u_{\lambda\mu}^{-1}(U_{i\lambda})$ pour $\lambda\leq\mu$, les
$U_{i\mu}$ sont deux à deux disjoints, car les
$U_i=u_\mu^{-1}(U_{i\mu})$ le sont et $u_\mu$ est dominant. Par suite,
aucune des adhérences $\overline{U_{i\mu}}$ n'est contenue dans une autre et
$u_\mu(U_i)$ est dense dans $U_{i\mu}$ puisque $u_\mu$ est dominant~; on a
donc
$\overline{U_{i\mu}}=\overline{u_\mu(Y_i)}$, ce qui prouve que les
$\overline{U_{i\mu}}$ sont les composantes irréductibles de $S_\mu$
$(0_{\mathrm I},2.1.7)$ et achève la démonstration.

\medskip
\noindent
(ii) Comme l'espace $S$ est noethérien, les $C_j$ sont ouverts et fermés dans
$S$ et quasi-compacts~; le même raisonnement que dans (i) montre donc qu'il
existe un $\lambda$ et des ouverts $V_{j\lambda}$ de $S_\lambda$ tels que
$C_j=u_\lambda^{-1}(V_{j\lambda})$ pour $1\leq j\leq n$. On voit aussi,
comme dans (i), que si l'on pose
$V_{j\mu}=u_{\lambda\mu}^{-1}(V_{j\lambda})$ pour $\mu\geq\lambda$, les
$V_{j\mu}$ sont deux à deux disjoints, et $u_\mu(C_j)$ est dense dans
$V_{j\mu}$~; cela entraîne que $V_{j\mu}$ est connexe. En outre, il résulte
de (8.3.4) que pour $\mu$ assez grand, la réunion des $V_{j\mu}$ est
$S_\mu$, puisque tout ouvert dans un préschéma est ind-constructible (1.9.6).
Les $V_{j\mu}$ sont donc les composantes connexes de $S_\mu$, ce qui achève
la démonstration.

On notera que si les morphismes $u_{\lambda\mu}$ sont des \emph{immersions},
ils vérifieront en particulier la condition \emph{b)} de (8.4.2).

\oldpage[IV-3]{19}
\begin{corollary}[8.4.3]
\label{IV.8.4.3-fr}
Supposons vérifiée l'une des conditions \emph{a)}, \emph{b)} de (8.4.2),
l'espace sous-jacent à $S$ étant noethérien~; alors, pour que $S$ soit
irréductible, il faut et il suffit qu'il existe $\lambda$ tel que $S_\mu$ le
soit pour tout $\mu\geq\lambda$.
\end{corollary}

\begin{proposition}[8.4.4]
\label{IV.8.4.4-fr}
Supposons qu'il existe un $\alpha$ tel que $S_\alpha$ soit quasi-compact et
que, pour $\lambda\leq\mu$, $u_{\lambda\mu}:S_\mu\to S_\lambda$ soit dominant.
Alors, pour que $S$ soit connexe, il faut et il suffit qu'il existe $\lambda$
tel que $S_\mu$ le soit pour tout $\mu\geq\lambda$.
\end{proposition}

\begin{proof}
La condition est suffisante en vertu de (8.4.1)~; d'autre part, on a vu
(8.3.8, (i)) que $u_\lambda:S\to S_\lambda$ est dominant pour $\lambda$ assez
grand, donc, si $S$ est connexe, il en est de même de $S_\lambda$, puisque
$u_\lambda(S)$ est dense dans $S_\lambda$ et connexe.
\end{proof}

\begin{corollary}[8.4.5]
\label{IV.8.4.5-fr}
Soient $k$ un corps, $X$ un $k$-préschéma quasi-compact. Pour que $X$ soit
géométriquement connexe (4.5.2), il faut et il suffit que, pour toute extension
finie et séparable $K$ de $k$, $X\otimes_k K$ soit connexe.
\end{corollary}

\begin{proof}
La condition est trivialement nécessaire. Pour voir qu'elle est suffisante,
nous devons prouver que si $\Omega$ est une clôture algébrique de $k$,
$X\otimes_k\Omega$ est connexe (4.5.1). Or, $\Omega$ est limite inductive
filtrante des sous-extensions finies $k'$ de $k$, et pour
$k\subset k'\subset k''\subset\Omega$, le morphisme
$X\otimes_k k''\to X\otimes_k k'$ est surjectif. On est donc ramené, en vertu
de (8.4.4), à prouver que $X\otimes_k k'$ est connexe pour toute extension
finie $k'$ de $k$. Mais si $K$ est la plus grande extension séparable contenue
dans $k'$, le morphisme $X\otimes_k k'\to X\otimes_k K$ est fini, surjectif et
radiciel, donc (2.4.5) un homéomorphisme, et comme $X\otimes_k K$ est connexe
par hypothèse, il en est de même de $X\otimes_k k'$.
\end{proof}

\begin{remark}[8.4.6]
\label{IV.8.4.6-fr}
\begin{enumerate}[label=(\roman*)]
  \item La démonstration de (8.4.2) montre que la conclusion de cette
  proposition est valable si l'on suppose que l'espace sous-jacent à $S$ est
  noethérien, qu'il existe $\alpha$ tel que $S_\alpha$ soit quasi-compact, et
  enfin que les $u_\lambda:S\to S_\lambda$ sont dominants.

  \item Par contre, la conclusion de (8.4.2) peut être en défaut lorsque les
  $u_\lambda$ ne sont pas dominants pour $\lambda$ assez grand, même lorsque
  les $S_\lambda$ et $S$ sont noethériens, ainsi que le montre l'exemple
  suivant. Prenons pour ensemble d'indices $\mathbf N$, tous les $S_n$ étant
  égaux à
  $\operatorname{Spec}(A\times K)=\operatorname{Spec}(A)\amalg
  \operatorname{Spec}(K)$, où $K$ est un corps, $A$ une $K$-algèbre
  quelconque, et tous les morphismes $u_{n,n+1}$ étant égaux au même morphisme
  correspondant à l'homomorphisme $(x,y)\mapsto(j(y),y)$ de $A\times K$ dans
  lui-même, où $j:K\to A$ est l'homomorphisme canonique. On vérifie aisément
  que la limite inductive de ce système d'anneaux est $K$, l'homomorphisme
  canonique $u_n$ correspondant à la seconde projection $A\times K\to K$. On
  voit donc que $S=\operatorname{Spec}(K)$ est irréductible bien qu'aucun des
  $S_n$ ne soit connexe.
\end{enumerate}
\end{remark}

\subsection{Modules de présentation finie sur une limite projective de préschémas}
\label{subsection:IV.8.5-fr}

\begin{env}[8.5.1]
\label{IV.8.5.1-fr}
Nous conservons toujours les notations de (8.2.2)~; nous nous bornerons en
outre au cas où $S_0$ est l'un des $S_\lambda$, auquel on peut toujours se
ramener.

\emph{Lorsque, dans cette section, nous considérerons une famille
$(\mathcal F_\lambda)$, où, pour tout $\lambda\in L$, $\mathcal F_\lambda$ est
un $\mathcal O_{S_\lambda}$-Module, il sera sous-entendu que cette famille
vérifie la condition}
\begin{equation}
  \mathcal F_\mu=u_{\lambda\mu}^*(\mathcal F_\lambda)
  \qquad\text{pour }\lambda\leq\mu.
  \tag{8.5.1.1}\label{IV.8.5.1.1-fr}
\end{equation}

\oldpage[IV-3]{20}
Nous poserons alors
\begin{equation}
  \mathcal F=u_\lambda^*(\mathcal F_\lambda)
  \tag{8.5.1.2}\label{IV.8.5.1.2-fr}
\end{equation}
qui est un $\mathcal O_S$-Module ne dépendant pas de l'indice $\lambda\in L$,
en vertu de l'hypothèse (8.5.1.1).

Soient maintenant $(\mathcal F_\lambda)$, $(\mathcal G_\lambda)$ deux telles
familles de $\mathcal O_{S_\lambda}$-Modules. Il est clair que les applications
$f_\lambda\mapsto u_{\lambda\mu}^*(f_\lambda)$ de
$\operatorname{Hom}_{S_\lambda}(\mathcal F_\lambda,\mathcal G_\lambda)$ dans
$\operatorname{Hom}_{S_\mu}(\mathcal F_\mu,\mathcal G_\mu)$ définissent alors
un \emph{système inductif} de groupes abéliens
$(\operatorname{Hom}_{S_\lambda}(\mathcal F_\lambda,\mathcal G_\lambda),
u_{\lambda\mu}^*)$, et que les applications
$f_\lambda\mapsto u_\lambda^*(f_\lambda)$ forment un \emph{système inductif}
d'homomorphismes de groupes abéliens
\[
  u_\lambda^*:\operatorname{Hom}_{S_\lambda}
  (\mathcal F_\lambda,\mathcal G_\lambda)
  \longrightarrow\operatorname{Hom}_S(\mathcal F,\mathcal G)~;
\]
d'où, en passant à la limite inductive, un homomorphisme canonique de groupes
abéliens
\begin{equation}
  u_{\mathcal F,\mathcal G}^*:
  \varinjlim\operatorname{Hom}_{S_\lambda}
  (\mathcal F_\lambda,\mathcal G_\lambda)
  \longrightarrow\operatorname{Hom}_S(\mathcal F,\mathcal G).
  \tag{8.5.1.3}\label{IV.8.5.1.3-fr}
\end{equation}

Notons que lorsque $\mathcal F_\lambda=\mathcal O_{S_\lambda}$, la condition
(8.5.1.1) est vérifiée, et que l'on a $\mathcal F=\mathcal O_S$~;
l'homomorphisme (8.5.1.3) donne donc $(0_{\mathrm I},5.1.1)$ un homomorphisme
canonique de groupes abéliens
\begin{equation}
  u_{\mathcal G}:\varinjlim\Gamma(S_\lambda,\mathcal G_\lambda)
  \longrightarrow\Gamma(S,\mathcal G).
  \tag{8.5.1.4}\label{IV.8.5.1.4-fr}
\end{equation}
\end{env}

\begin{theorem}[8.5.2]
\label{IV.8.5.2-fr}
\begin{enumerate}[label=(\roman*)]
  \item Supposons $S_0$ quasi-compact (resp. quasi-compact et quasi-séparé) et
  que, pour un $\lambda\in L$, $\mathcal F_\lambda$ soit quasi-cohérent et de
  type fini (resp. de présentation finie) et $\mathcal G_\lambda$
  quasi-cohérent. Alors l'homomorphisme $u_{\mathcal F,\mathcal G}^*$ est
  injectif (resp. bijectif).

  \item Supposons $S_0$ quasi-compact et quasi-séparé. Pour tout
  $\mathcal O_S$-Module quasi-cohérent $\mathcal F$ de présentation finie, il
  existe un $\lambda\in L$ et un $\mathcal O_{S_\lambda}$-Module
  quasi-cohérent $\mathcal F_\lambda$ de présentation finie tels que
  $\mathcal F$ soit isomorphe à $u_\lambda^*(\mathcal F_\lambda)$.
\end{enumerate}
\end{theorem}

\begin{proof}
\begin{enumerate}[label=(\roman*)]
  \item On peut évidemment se borner au cas où $S_0=S_\lambda$ puisque les
  morphismes $u_{0\lambda}:S_\lambda\to S_0$ sont affines, donc quasi-compacts
  et séparés. Considérons d'abord le cas où
  $S_0=\operatorname{Spec}(A_0)$ est \emph{affine}. Alors l'assertion (i)
  équivaut au

  \begin{lemma}[8.5.2.1]
  \label{IV.8.5.2.1-fr}
  Soient $A_0$ un anneau, $(A_\lambda)_{\lambda\in L}$ un système inductif de
  $A_0$-algèbres, $A=\varinjlim A_\lambda$~; soient $M_0$, $N_0$ deux
  $A_0$-modules, et posons
  $M_\lambda=M_0\otimes_{A_0}A_\lambda$,
  $N_\lambda=N_0\otimes_{A_0}A_\lambda$,
  $M=M_0\otimes_{A_0}A=\varinjlim M_\lambda$,
  $N=N_0\otimes_{A_0}A=\varinjlim N_\lambda$. Si $M_0$ est de type fini
  (resp. de présentation finie), l'homomorphisme canonique
  \begin{equation}
    \varinjlim\operatorname{Hom}_{A_\lambda}(M_\lambda,N_\lambda)
    \longrightarrow\operatorname{Hom}_A(M,N)
    \tag{8.5.2.2}\label{IV.8.5.2.2-fr}
  \end{equation}
  est injectif (resp. bijectif).
  \end{lemma}

  On sait en effet (Bourbaki, \emph{Alg.}, chap. II, 3\textsuperscript{e} éd.,
  \S~5, n\textsuperscript{o}~1) que l'on a des isomorphismes canoniques
  fonctoriels
  \[
    \operatorname{Hom}_{A_\lambda}(M_\lambda,N_\lambda)
    \xrightarrow{\sim}\operatorname{Hom}_{A_0}(M_0,N_\lambda),
    \qquad
    \operatorname{Hom}_A(M,N)
    \xrightarrow{\sim}\operatorname{Hom}_{A_0}(M_0,N)
  \]

  \oldpage[IV-3]{21}
  si bien que l'homomorphisme (8.5.2.2) n'est autre, à des isomorphismes
  canoniques près, que l'homomorphisme canonique
  \begin{equation}
    \varinjlim\operatorname{Hom}_{A_0}(M_0,N_\lambda)
    \longrightarrow
    \operatorname{Hom}_{A_0}(M_0,\varinjlim N_\lambda),
    \tag{8.5.2.3}\label{IV.8.5.2.3-fr}
  \end{equation}
  qui, à tout système inductif d'homomorphismes de $A_0$-modules
  $\theta_\lambda:M_0\to N_\lambda$, fait correspondre sa limite inductive.

  Or, si $M_0$ est de type fini (resp. de présentation finie), on a une suite
  exacte $A_0^m\to M_0\to0$ (resp. $A_0^n\to A_0^m\to M_0\to0$)~; comme il est
  clair que (8.5.2.3) est bijectif lorsque $M_0$ est de la forme $A_0^q$, il
  suffit d'utiliser l'exactitude à gauche du foncteur
  $M_0\mapsto\operatorname{Hom}_{A_0}(M_0,P)$ et l'exactitude du foncteur
  $\varinjlim$ (dans la catégorie des groupes abéliens) pour conclure.

  Passons au cas où $S_0$ est quasi-compact, et soit $(U_i)$ un recouvrement
  fini de $S_0$ par des ouverts affines~; pour tout $\lambda$, les
  $U_{i\lambda}=u_{0\lambda}^{-1}(U_i)$ forment un recouvrement ouvert affine de
  $S_\lambda$, et les $V_i=u_0^{-1}(U_i)$ un recouvrement ouvert affine de $S$.
  Pour voir que $u_{\mathcal F,\mathcal G}^*$ est injectif, il faut prouver que
  si $f_\lambda:\mathcal F_\lambda\to\mathcal G_\lambda$ est tel que
  $f=u_\lambda^*(f_\lambda)=0$, alors il existe $\mu\geq\lambda$ tel que
  $f_\mu=u_{\lambda\mu}^*(f_\lambda)=0$. En vertu du lemme (8.5.2.1), pour
  chaque $i$ il existe un $\lambda_i$ tel que
  $f_\mu|U_{i\mu}=0$ pour $\mu\geq\lambda_i$. Il suffit donc de prendre $\mu$
  plus grand que tous les $\lambda_i$.

  Supposons en outre $S_0$ quasi-séparé et $\mathcal F_0$ de présentation finie,
  et soit $f:\mathcal F\to\mathcal G$ un homomorphisme de
  $\mathcal O_S$-Modules. En vertu du lemme (8.5.2.1), pour tout $i$, il existe
  un indice $\lambda_i$ et un homomorphisme
  $f_{\lambda_i}^{(i)}:\mathcal F_{\lambda_i}|U_{i,\lambda_i}
  \to\mathcal G_{\lambda_i}|U_{i,\lambda_i}$ tels que
  $u_{\lambda_i}^*(f_{\lambda_i}^{(i)})=f|V_i$. Comme $L$ est filtrant, on peut
  en outre supposer tous les $\lambda_i$ égaux à un même $\lambda$. Notons
  maintenant que $S_\lambda$ est quasi-séparé (1.2.3) et
  $\mathcal F_\lambda$ est un $\mathcal O_{S_\lambda}$-Module quasi-cohérent de
  présentation finie $(0_{\mathrm I},5.2.5)$~; comme, pour tout couple d'indices
  $i,j$, $U_{ij\lambda}=U_{i\lambda}\cap U_{j\lambda}$ est quasi-compact et que
  l'on a
  \[
    u_\lambda^*(f_\lambda^{(i)}|U_{ij\lambda})
    =u_\lambda^*(f_\lambda^{(j)}|U_{ij\lambda})
    =f|(V_i\cap V_j)
  \]
  par définition, il résulte de ce qu'on a vu plus haut qu'il existe un indice
  $\lambda_{ij}$ tel que
  \[
    u_{\lambda\mu}^*(f_\lambda^{(i)}|U_{ij\lambda})
    =u_{\lambda\mu}^*(f_\lambda^{(j)}|U_{ij\lambda})
    \quad\text{pour }\mu\geq\lambda_{ij}~;
  \]
  prenant $\mu$ plus grand que tous les $\lambda_{ij}$, on voit donc que
  $u_{\lambda\mu}^*(f_\lambda^{(i)})$ et
  $u_{\lambda\mu}^*(f_\lambda^{(j)})$ coïncident dans
  $U_{i\mu}\cap U_{j\mu}$ pour tout couple $(i,j)$, et par suite définissent
  un homomorphisme $f_\mu:\mathcal F_\mu\to\mathcal G_\mu$ tel que
  $f=u_\mu^*(f_\mu)$.
\end{enumerate}
\end{proof}

Avant de passer à la démonstration de (ii), notons les corollaires suivants de
(i)~:

\begin{corollary}[8.5.2.4]
\label{IV.8.5.2.4-fr}
Supposons $S_0$ quasi-compact, $\mathcal F_\lambda$ quasi-cohérent de type fini,
$\mathcal G_\lambda$ quasi-cohérent de présentation finie. Soit
$f_\lambda:\mathcal F_\lambda\to\mathcal G_\lambda$ un homomorphisme. Pour que
$f=u_\lambda^*(f_\lambda):\mathcal F\to\mathcal G$ soit un isomorphisme, il faut
et il suffit qu'il existe $\mu\geq\lambda$ tel que
$f_\mu=u_{\lambda\mu}^*(f_\lambda):\mathcal F_\mu\to\mathcal G_\mu$ soit un
isomorphisme.
\end{corollary}

\begin{proof}
On peut toujours supposer $S_\lambda=S_0$~; la question étant locale sur $S_0$,
on peut en outre ($S_0$ étant quasi-compact et $L$ filtrant) se ramener au cas
où $S_0$ est affine, donc quasi-séparé. La condition étant trivialement
suffisante, il reste à montrer qu'elle est nécessaire~: or, il y a par
hypothèse un $\mathcal O_S$-homomorphisme $g:\mathcal G\to\mathcal F$ tel que
$g\circ f=1_{\mathcal F}$ et $f\circ g=1_{\mathcal G}$. Comme $\mathcal G$ est
de présentation finie, il existe $\nu\geq\lambda$ et un homomorphisme
$g_\nu:\mathcal G_\nu\to\mathcal F_\nu$ tels que $g=u_\nu^*(g_\nu)$ en vertu de
(8.5.2, (i))~; on a par suite
$u_\nu^*(g_\nu\circ f_\nu)=1_{\mathcal F}$

\oldpage[IV-3]{22}
et $u_\nu^*(f_\nu\circ g_\nu)=1_{\mathcal G}$~; compte tenu de ce que
$\mathcal F_\nu$ et $\mathcal G_\nu$ sont de type fini, on en conclut par
(8.5.2, (i)) qu'il existe $\mu\geq\nu$ tel que l'on ait
$g_\mu\circ f_\mu=1_{\mathcal F_\mu}$ et
$f_\mu\circ g_\mu=1_{\mathcal G_\mu}$, d'où le corollaire.
\end{proof}

\begin{corollary}[8.5.2.5]
\label{IV.8.5.2.5-fr}
Supposons $S_0$ quasi-compact et quasi-séparé. Supposons que
$\mathcal F_\lambda$, $\mathcal G_\lambda$ soient des
$\mathcal O_{S_\lambda}$-Modules quasi-cohérents de présentation finie. Pour
que $\mathcal F$ et $\mathcal G$ soient isomorphes, il faut et il suffit qu'il
existe $\mu\geq\lambda$ tel que $\mathcal F_\mu$ et $\mathcal G_\mu$ soient
isomorphes. De plus, pour tout isomorphisme
$f:\mathcal F\xrightarrow{\sim}\mathcal G$, il existe un $\nu\geq\mu$ et un
isomorphisme $f_\nu:\mathcal F_\nu\xrightarrow{\sim}\mathcal G_\nu$ tels que
$f=u_\nu^*(f_\nu)$.
\end{corollary}

\begin{proof}
Cela résulte de (8.5.2.4) et de (8.5.2, (i)) puisque tout homomorphisme
$f:\mathcal F\to\mathcal G$ est de la forme $u_\mu^*(f_\mu)$ pour un
$\mu\geq\lambda$ et un homomorphisme
$f_\mu:\mathcal F_\mu\to\mathcal G_\mu$.
\end{proof}

\par\medskip\noindent\textup{(ii)} Considérons encore en premier lieu le cas où
$S_0=\operatorname{Spec}(A_0)$ est
affine. Alors l'assertion équivaut au lemme (5.13.7.1).

Dans le cas général, en partant d'un recouvrement ouvert affine fini $(U_i)$
de $S_0$, on déduit de (5.13.7.1) que pour tout $i$, il existe un indice
$\lambda(i)$ et un $\mathcal O_{U_{i,\lambda(i)}}$-Module quasi-cohérent de
présentation finie $\mathcal F_i$ tels que
$u_{\lambda(i)}^*(\mathcal F_i)=\mathcal F|V_i$ (avec les notations de (i)). En
outre, comme $L$ est filtrant, on peut supposer que tous les $\lambda_i$ sont
égaux à un même $\lambda$. Comme
$U_{ij\lambda}=U_{i\lambda}\cap U_{j\lambda}$ est quasi-compact et
quasi-séparé (1.2.7), il résulte de (8.5.2.5) que pour tout couple $(i,j)$, il
existe un indice $\lambda_{ij}\geq\lambda$ et un isomorphisme
\[
  \theta_{ij}:u_{\lambda,\lambda_{ij}}^*(\mathcal F_i|U_{ij\lambda})
  \xrightarrow{\sim}
  u_{\lambda,\lambda_{ij}}^*(\mathcal F_j|U_{ij\lambda})
\]
tel que $u_{\lambda_{ij}}^*(\theta_{ij})$ soit l'automorphisme identique de
$\mathcal F|(V_i\cap V_j)$~; on peut encore supposer tous les
$\lambda_{ij}$ égaux à $\lambda$. Changeant les notations, on peut donc
supposer qu'il existe pour tout couple $(i,j)$ un isomorphisme
\[
  \theta_{ij}:\mathcal F_i|U_{ij\lambda}
  \xrightarrow{\sim}\mathcal F_j|U_{ij\lambda},
\]
tels que $u_\lambda^*(\theta_{ij})$ soit l'automorphisme identique de
$\mathcal F|(V_i\cap V_j)$. Enfin, pour trois indices quelconques $i,j,k$, si
l'on pose
$U_{ijk,\lambda}=U_{i\lambda}\cap U_{j\lambda}\cap U_{k\lambda}$,
$U_{ijk,\lambda}$ est quasi-compact, et si $\theta'_{ij}$, $\theta'_{jk}$ et
$\theta'_{ik}$ désignent les restrictions de $\theta_{ij}$, $\theta_{jk}$ et
$\theta_{ik}$ à $U_{ijk,\lambda}$, on a
\[
  u_\lambda^*(\theta'_{ij}\circ\theta'_{jk})=u_\lambda^*(\theta'_{ik}).
\]
Il y a donc, en vertu de (i), un indice $\mu\geq\lambda$ tel que l'on ait
\[
  u_{\lambda\mu}^*(\theta'_{ik})
  =u_{\lambda\mu}^*(\theta'_{ij}\circ\theta'_{jk}),
\]
donc les isomorphismes
\[
  u_{\lambda\mu}^*(\theta'_{ij}):
  u_{\lambda\mu}^*(\mathcal F_i)|U_{ij\mu}
  \xrightarrow{\sim}
  u_{\lambda\mu}^*(\mathcal F_j)|U_{ij\mu}
\]
vérifient la condition de recollement, et définissent par suite sur $S_\mu$ un
$\mathcal O_{S_\mu}$-Module quasi-cohérent $\mathcal F_\mu$ de présentation
finie $(0_{\mathrm I},3.3.1)$ tel que $\mathcal F$ soit isomorphe à
$u_\mu^*(\mathcal F_\mu)$.

\begin{env}[8.5.3]
\label{IV.8.5.3-fr}
Le résultat de (8.5.2) s'exprime encore en disant que si $S_0$ est
quasi-compact et quasi-séparé, la catégorie des $\mathcal O_S$-Modules
quasi-cohérents de présentation finie est déterminée à équivalence près par la
donnée des catégories des $\mathcal O_{S_\lambda}$-Modules quasi-cohérents de
présentation finie, des foncteurs $u_{\lambda\mu}^*$ entre ces catégories, et
des isomorphismes de transition
$u_{\mu\nu}^*\circ u_{\lambda\mu}^*\xrightarrow{\sim}u_{\lambda\nu}^*$. De façon
imagée, on peut dire que la donnée d'un $\mathcal O_S$-Module quasi-cohérent de
présentation finie $\mathcal F$ équivaut «~fonctoriellement~» à la donnée d'un
$\mathcal O_{S_\lambda}$-Module de présentation finie $\mathcal F'_\lambda$
pour $\lambda$ grand et que, si un $\mathcal O_{S_\mu}$-Module quasi-cohérent de
présentation finie $\mathcal F''_\mu$ a aussi $\mathcal F$ comme image
réciproque, alors $\mathcal F'_\lambda$ et $\mathcal F''_\mu$ ont même image
réciproque dans un $S_\nu$ convenable ($\nu\geq\lambda$, $\nu\geq\mu$).

Nous allons interpréter de ce point de vue diverses notions liées aux
$\mathcal O_S$-Modules quasi-cohérents.
\end{env}

\oldpage[IV-3]{23}
\begin{corollary}[8.5.4]
\label{IV.8.5.4-fr}
Supposons $S_0$ quasi-compact et quasi-séparé~; alors, pour tout
$\mathcal O_{S_\lambda}$-Module quasi-cohérent $\mathcal G_\lambda$,
l'homomorphisme canonique (8.5.1.4) est bijectif.
\end{corollary}

\begin{proof}
En effet, il suffit d'appliquer (8.5.2, (i)) au cas où
$\mathcal F_\lambda=\mathcal O_{S_\lambda}$, qui est de présentation finie.
\end{proof}

\begin{proposition}[8.5.5]
\label{IV.8.5.5-fr}
Supposons $S_0$ quasi-compact, et supposons que $\mathcal F_\lambda$ soit un
$\mathcal O_{S_\lambda}$-Module quasi-cohérent de présentation finie. Pour que
$\mathcal F$ soit localement libre (resp. localement libre de rang $n$), il
faut et il suffit qu'il existe $\mu\geq\lambda$ tel que $\mathcal F_\mu$ le
soit.
\end{proposition}

\begin{proof}
La condition étant trivialement suffisante, prouvons qu'elle est nécessaire.
Si $\mathcal F$ est localement libre (resp. localement libre de rang $n$), il
existe un recouvrement ouvert affine fini $(V_i)$ de $S$ tel que
$\mathcal F|V_i$ soit isomorphe à $\mathcal O_S^{n_i}|V_i$ (resp.
$\mathcal O_S^n|V_i$) pour tout $i$. En vertu de (8.2.11), il existe un
$\nu\geq\lambda$ et pour chaque $i$ un ouvert quasi-compact $U_{i\nu}$ de
$S_\nu$ tels que $V_i=u_\nu^{-1}(U_{i\nu})$. Comme $S_\nu$ est quasi-compact,
chaque $U_{i\nu}$ est réunion finie d'ouverts affines~; on est donc ramené au
cas où $S_0$ est affine et $\mathcal F=\mathcal O_S^n$~; on sait alors qu'il
existe $\mu\geq\lambda$ tel que $\mathcal F_\mu$ soit isomorphe à
$\mathcal O_{S_\mu}^n$ (8.5.2.5).
\end{proof}

\begin{proposition}[8.5.6]
\label{IV.8.5.6-fr}
Supposons $S_0$ quasi-compact, et considérons une suite
\[
  \mathcal F_\lambda\longrightarrow\mathcal G_\lambda
  \longrightarrow\mathcal H_\lambda\longrightarrow0
\]
d'homomorphismes de $\mathcal O_{S_\lambda}$-Modules quasi-cohérents, où
$\mathcal F_\lambda$ et $\mathcal G_\lambda$ sont de type fini et
$\mathcal H_\lambda$ de présentation finie. Pour que la suite correspondante
$\mathcal F\to\mathcal G\to\mathcal H\to0$ soit exacte, il faut et il suffit
qu'il existe $\mu\geq\lambda$ tel que la suite
$\mathcal F_\mu\to\mathcal G_\mu\to\mathcal H_\mu\to0$ le soit (auquel cas il
en est de même de la suite
$\mathcal F_\nu\to\mathcal G_\nu\to\mathcal H_\nu\to0$ pour
$\nu\geq\mu$).
\end{proposition}

\begin{proof}
Le fait que la condition soit suffisante et la dernière assertion résultent
de ce que le foncteur $u_\lambda^*$ (resp. $u_{\mu\nu}^*$) est exact à droite.
Pour démontrer que la condition est nécessaire, notons qu'il résulte de
l'hypothèse et de (8.5.2, (i)) qu'il existe $\nu\geq\lambda$ tel que le composé
$\mathcal F_\nu\to\mathcal G_\nu\to\mathcal H_\nu$ soit nul. Si l'on pose
$\mathcal H'_\nu=\operatorname{Coker}(\mathcal F_\nu\to\mathcal G_\nu)$, on a
donc un homomorphisme $f_\nu:\mathcal H'_\nu\to\mathcal H_\nu$~; par
hypothèse, $u_\nu^*(f_\nu)$ est un isomorphisme, et il résulte donc de
(8.5.2.4) qu'il existe $\mu\geq\nu$ tel que $u_{\nu\mu}^*(f_\nu)$ soit un
isomorphisme, ce qui achève la démonstration.
\end{proof}

\begin{corollary}[8.5.7]
\label{IV.8.5.7-fr}
Supposons $S_0$ quasi-compact, $\mathcal F_\lambda$ quasi-cohérent,
$\mathcal G_\lambda$ quasi-cohérent de type fini, et soit
$f_\lambda:\mathcal F_\lambda\to\mathcal G_\lambda$ un homomorphisme. Pour que
$f=u_\lambda^*(f_\lambda):\mathcal F\to\mathcal G$ soit surjectif, il faut et
il suffit qu'il existe $\mu\geq\lambda$ tel que
$f_\mu=u_{\lambda\mu}^*(f_\lambda):\mathcal F_\mu\to\mathcal G_\mu$ le soit.
\end{corollary}

\begin{proof}
C'est le cas particulier de (8.5.6) appliqué à la suite
$0\to\mathcal H_\lambda\to0\to0$, où
$\mathcal H_\lambda=\operatorname{Coker}(f_\lambda)$, qui est quasi-cohérent et
de type fini (compte tenu de ce que l'on a
$\mathcal H=\operatorname{Coker}(f)$ et
$\mathcal H_\mu=\operatorname{Coker}(f_\mu)$, en vertu de l'exactitude à droite
des foncteurs $u_\lambda^*$ et $u_{\lambda\mu}^*$).
\end{proof}

\begin{corollary}[8.5.8]
\label{IV.8.5.8-fr}
Supposons $S_0$ quasi-compact et les morphismes
$u_{\lambda\mu}:S_\mu\to S_\lambda$ plats. Alors~:
\begin{enumerate}[label=(\roman*)]
  \item Soit
  $\mathcal F_\lambda\xrightarrow{f_\lambda}\mathcal G_\lambda
  \xrightarrow{g_\lambda}\mathcal H_\lambda$ une suite d'homomorphismes de
  $\mathcal O_{S_\lambda}$-Modules quasi-cohérents, tels que
  $\operatorname{Im}f_\lambda$ et $\operatorname{Ker}g_\lambda$ soient de type
  fini. Pour que la suite correspondante
  $\mathcal F\xrightarrow{f}\mathcal G\xrightarrow{g}\mathcal H$ soit

  \oldpage[IV-3]{24}
  exacte, il faut et il suffit qu'il existe $\mu\geq\lambda$ tel que la suite
  $\mathcal F_\mu\xrightarrow{f_\mu}\mathcal G_\mu
  \xrightarrow{g_\mu}\mathcal H_\mu$ soit exacte.

  \item Soit $f_\lambda:\mathcal F_\lambda\to\mathcal G_\lambda$ un
  homomorphisme de $\mathcal O_{S_\lambda}$-Modules quasi-cohérents tel que
  $\operatorname{Ker}f_\lambda$ soit de type fini. Pour que
  $f=u_\lambda^*(f_\lambda):\mathcal F\to\mathcal G$ soit injectif, il faut et
  il suffit qu'il existe $\mu\geq\lambda$ tel que
  $f_\mu=u_{\lambda\mu}^*(f_\lambda):\mathcal F_\mu\to\mathcal G_\mu$ le soit.
\end{enumerate}
\end{corollary}

\begin{proof}
\begin{enumerate}[label=(\roman*)]
  \item Compte tenu de (8.3.8, (ii)), notons que, par platitude,
  $\operatorname{Im}f$ et $\operatorname{Ker}g$ (resp.
  $\operatorname{Im}f_\mu$ et $\operatorname{Ker}g_\mu$ pour
  $\mu\geq\lambda$) sont les images réciproques de
  $\operatorname{Im}f_\lambda$ et $\operatorname{Ker}g_\lambda$
  $(0_{\mathrm I},6.7.2)$. Supposons que la suite
  $\mathcal F\xrightarrow{f}\mathcal G\xrightarrow{g}\mathcal H$ soit exacte.
  Puisque $\operatorname{Im}f_\lambda$ est de type fini, il existe
  $\mu\geq\lambda$ tel que le composé
  $\mathcal F_\mu\xrightarrow{f_\mu}\mathcal G_\mu
  \xrightarrow{g_\mu}\mathcal H_\mu$ soit nul, en vertu de (8.5.2, (i)).
  Changeant les notations, on peut donc déjà supposer que
  $g_\lambda\circ f_\lambda=0$. Alors comme l'homomorphisme
  $\mathcal F\to\operatorname{Ker}g$ est surjectif et que
  $\operatorname{Ker}g_\lambda$ est de type fini, il résulte de (8.5.7) qu'il
  existe $\mu\geq\lambda$ tel que l'homomorphisme
  $\mathcal F_\mu\to\operatorname{Ker}g_\mu$ soit surjectif, ce qui achève de
  prouver (i).

  \item L'assertion est le cas particulier de (i) appliqué à la suite
  $0\to\mathcal F_\lambda\to\mathcal G_\lambda$.
\end{enumerate}
\end{proof}

\begin{lemma}[8.5.9]
\label{IV.8.5.9-fr}
Supposons $S_0$ quasi-compact, $\mathcal F_\lambda$ quasi-cohérent de type
fini~; soient $\mathcal G'_\lambda$ et $\mathcal G''_\lambda$ deux quotients
quasi-cohérents de $\mathcal F_\lambda$, $\mathcal G'_\lambda$ étant en outre
supposé de présentation finie. Pour que $\mathcal G''$ soit quotient de
$\mathcal G'$, il faut et il suffit qu'il existe $\mu\geq\lambda$ tel que
$\mathcal G''_\mu$ soit quotient de $\mathcal G'_\mu$.
\end{lemma}

\begin{proof}
Par hypothèse, il y a deux homomorphismes surjectifs
$p'_\lambda:\mathcal F_\lambda\to\mathcal G'_\lambda$,
$p''_\lambda:\mathcal F_\lambda\to\mathcal G''_\lambda$~; en vertu de
l'exactitude à droite de $u_\lambda^*$ et de $u_{\lambda\mu}^*$,
$p'=u_\lambda^*(p'_\lambda)$, $p''=u_\lambda^*(p''_\lambda)$,
$p'_\mu=u_{\lambda\mu}^*(p'_\lambda)$,
$p''_\mu=u_{\lambda\mu}^*(p''_\lambda)$ sont aussi surjectifs~; en outre, s'il
existe un homomorphisme $f:\mathcal G'\to\mathcal G''$ (resp.
$f_\mu:\mathcal G'_\mu\to\mathcal G''_\mu$) tel que $p''=f\circ p'$ (resp.
$p''_\mu=f_\mu\circ p'_\mu$), cet homomorphisme est nécessairement unique, ce
qui montre que la question est locale sur $S_\lambda$, et qu'on peut par suite
($S_0$ étant quasi-compact et $L$ filtrant) supposer $S_\lambda$ affine, donc
quasi-séparé. Il est clair que la condition de l'énoncé est suffisante.
Inversement, comme $\mathcal G'_\lambda$ est de présentation finie,
$S_\lambda$ quasi-compact et quasi-séparé, il résulte de (8.5.2, (i)) que s'il
existe un homomorphisme $f:\mathcal G'\to\mathcal G''$ tel que
$p''=f\circ p'$, il existe un $\mu\geq\lambda$ et un homomorphisme
$f_\mu:\mathcal G'_\mu\to\mathcal G''_\mu$ tels que
$f=u_\mu^*(f_\mu)$ et $p''_\mu=f_\mu\circ p'_\mu$, d'où le lemme.
\end{proof}

\begin{env}[8.5.10]
\label{IV.8.5.10-fr}
Dans la fin de ce numéro, pour tout Module quasi-cohérent $\mathcal F$ sur un
préschéma, notons $\mathfrak Q(\mathcal F)$ l'ensemble des Modules quotients de
$\mathcal F$ qui sont de présentation finie. Si $\mathcal F_\lambda$ est
quasi-cohérent et $\mathcal G_\lambda\in\mathfrak Q(\mathcal F_\lambda)$, il
résulte du fait que $u_{\lambda\mu}^*$ et $u_\lambda^*$ sont exacts à droite,
que l'on a $\mathcal G_\mu\in\mathfrak Q(\mathcal F_\mu)$ pour
$\mu\geq\lambda$ et $\mathcal G\in\mathfrak Q(\mathcal F)$~; il est clair que
$(\mathfrak Q(\mathcal F_\lambda),u_{\lambda\mu}^*)$ est un système inductif
d'ensembles, et que les
$u_\lambda^*:\mathfrak Q(\mathcal F_\lambda)\to\mathfrak Q(\mathcal F)$
forment un système inductif d'applications, d'où, en passant à la limite
inductive, une application canonique
\begin{equation}
  u_{\mathfrak Q}:\varinjlim\mathfrak Q(\mathcal F_\lambda)
  \longrightarrow\mathfrak Q(\mathcal F).
  \tag{8.5.10.1}\label{IV.8.5.10.1-fr}
\end{equation}

En outre, si $(\mathcal F'_\lambda)$ est une seconde famille de
$\mathcal O_{S_\lambda}$-Modules quasi-cohérents et si, pour tout $\lambda$,
$\mathcal F'_\lambda$ est un quotient de $\mathcal F_\lambda$, alors
$\mathcal F'$ est un quotient de $\mathcal F$ et l'on a un diagramme
commutatif

\oldpage[IV-3]{25}
\[
\begin{tikzcd}
\varinjlim\mathfrak Q(\mathcal F_\lambda)
  \ar[r,"u_{\mathfrak Q}"] \ar[d]
  & \mathfrak Q(\mathcal F) \ar[d] \\
\varinjlim\mathfrak Q(\mathcal F'_\lambda)
  \ar[r,"u_{\mathfrak Q}"']
  & \mathfrak Q(\mathcal F')
\end{tikzcd}
\tag{8.5.10.2}\label{IV.8.5.10.2-fr}
\]
\end{env}

\begin{proposition}[8.5.11]
\label{IV.8.5.11-fr}
Supposons $S_0$ quasi-compact (resp. quasi-compact et quasi-séparé). Supposons
$\mathcal F_\lambda$ quasi-cohérent de type fini (resp. de présentation finie)
pour tout $\lambda$~; alors l'application canonique (8.5.10.1) est injective
(resp. bijective).
\end{proposition}

\begin{proof}
La première assertion résulte du lemme plus précis (8.5.9). Pour démontrer la
seconde, considérons un $\mathcal O_S$-Module quotient $\mathcal G$ de
$\mathcal F$ qui soit de présentation finie. Il résulte de (8.5.2, (ii)) qu'il
existe un $\lambda'\in L$ et un $\mathcal O_{S_{\lambda'}}$-Module
$\mathcal G_{\lambda'}$ quasi-cohérent de présentation finie tels que
$\mathcal G=u_{\lambda'}^*(\mathcal G_{\lambda'})$~; comme $L$ est filtrant, on
peut supposer $\lambda'=\lambda$ (en remplaçant $\lambda$ et $\lambda'$ par un
indice qui les majore). Considérons alors l'homomorphisme canonique
$p:\mathcal F\to\mathcal G$~; il résulte de (8.5.2, (i)) qu'il existe un
$\mu\geq\lambda$ et un homomorphisme
$p_\mu:\mathcal F_\mu\to\mathcal G_\mu$ tels que $p=u_\mu^*(p_\mu)$. En outre,
en vertu de (8.5.7), on peut supposer $\mu$ pris assez grand pour que $p_\mu$
soit surjectif, ce qui termine la démonstration.
\end{proof}

\subsection{Sous-préschémas de présentation finie d'une limite projective de préschémas}
\label{subsection:IV.8.6-fr}

\begin{env}[8.6.1]
\label{IV.8.6.1-fr}
Étant donné un préschéma $Y$, désignons dans ce numéro par
$\mathfrak{Spr}(Y)$ l'ensemble ordonné (I, 4.1.10) des sous-préschémas de $Y$
qui sont de présentation finie sur $Y$ (1.6.1), par $\mathfrak{Spro}(Y)$
(resp. $\mathfrak{Sprf}(Y)$) la partie de $\mathfrak{Spr}(Y)$ formée des
sous-préschémas induits sur des ouverts (resp. des sous-préschémas fermés) de
$Y$, de présentation finie sur $Y$~; il revient au même de dire qu'un
sous-préschéma $Z$ de $Y$ appartient à $\mathfrak{Spro}(Y)$ (resp.
$\mathfrak{Sprf}(Y)$) ou qu'il est induit sur un ouvert et que l'espace
sous-jacent $Z$ est rétrocompact dans $Y$ (resp. qu'il est fermé et que l'idéal
$\mathcal J$ de $\mathcal O_Y$ qui le définit est de type fini, ce qui signifie
aussi que $j_*(\mathcal O_Z)\in\mathfrak Q(\mathcal O_Y)$ si $j:Z\to Y$ est
l'injection canonique) (1.6.1 et 1.4.5).
\end{env}

\begin{env}[8.6.2]
\label{IV.8.6.2-fr}
Nous conservons toujours les notations de (8.2.2) et supposons que $S_0$ est
un des $S_\lambda$. Soit $Y_\lambda$ un sous-préschéma de $S_\lambda$~; alors
$Y_\mu=u_{\lambda\mu}^{-1}(Y_\lambda)$ (resp.
$Y=u_\lambda^{-1}(Y_\lambda)$) est un sous-préschéma de $S_\mu$ pour
$\mu\geq\lambda$ (resp. de $S$)~; il est induit sur un ouvert (resp. il est
fermé) si $Y_\lambda$ l'est (I, 4.3.2) et de présentation finie sur $S_\mu$
(resp. $S$) si $Y_\lambda$ est de présentation finie sur $S_\lambda$
(1.6.2, (iii)). Par suite
$(\mathfrak{Spr}(S_\lambda),u_{\lambda\mu}^{-1})$ (resp.
$(\mathfrak{Spro}(S_\lambda),u_{\lambda\mu}^{-1})$,
$(\mathfrak{Sprf}(S_\lambda),u_{\lambda\mu}^{-1})$) est un système inductif et
les applications
$u_\lambda^{-1}:\mathfrak{Spr}(S_\lambda)\to\mathfrak{Spr}(S)$ (resp.
$\mathfrak{Spro}(S_\lambda)\to\mathfrak{Spro}(S)$,
$\mathfrak{Sprf}(S_\lambda)\to\mathfrak{Sprf}(S)$) forment un système inductif
d'applications~; d'où, en passant à la limite inductive, des applications
canoniques
\begin{align}
  \varinjlim\mathfrak{Spr}(S_\lambda)
    &\longrightarrow\mathfrak{Spr}(S)
    \tag{8.6.2.1}\label{IV.8.6.2.1-fr}\\
  \varinjlim\mathfrak{Spro}(S_\lambda)
    &\longrightarrow\mathfrak{Spro}(S)
    \tag{8.6.2.2}\label{IV.8.6.2.2-fr}\\
  \varinjlim\mathfrak{Sprf}(S_\lambda)
    &\longrightarrow\mathfrak{Sprf}(S).
    \tag{8.6.2.3}\label{IV.8.6.2.3-fr}
\end{align}
\end{env}

\oldpage[IV-3]{26}

Rappelons (I, 4.1.10) que $\mathfrak{Spr}(X)$, pour tout préschéma $X$, est un
ensemble \emph{ordonné} par la relation «~$Z$ est un sous-préschéma du
sous-préschéma $Y$~» qui s'écrit $Z\leq Y$. Les applications
$u_{\lambda\mu}^{-1}$ et $u_\lambda^{-1}$ sont \emph{croissantes} pour les
relations d'ordre correspondantes dans $\mathfrak{Spr}(S_\lambda)$,
$\mathfrak{Spr}(S_\mu)$, $\mathfrak{Spr}(S)$. En outre, on définit une relation
d'ordre dans l'ensemble $\varinjlim\mathfrak{Spr}(S_\lambda)$ en écrivant que
$\eta\leq\eta'$ pour deux éléments de cet ensemble lorsqu'il existe un
$\lambda$ et deux éléments $Y_\lambda$, $Y'_\lambda$ de
$\mathfrak{Spr}(S_\lambda)$, dont $\eta$ et $\eta'$ sont les images canoniques,
et qui sont tels que $Y_\lambda\leq Y'_\lambda$~; on vérifie aisément que cela
ne dépend pas des représentants $Y_\lambda$, $Y'_\lambda$ considérés, et que
l'on a bien ainsi une \emph{relation d'ordre}. Cela étant, le fait que les
$u_\lambda^{-1}$ sont croissantes entraîne aussitôt que l'application
canonique (8.6.2.1) est \emph{croissante}~; il en est de même évidemment de
(8.6.2.2) et (8.6.2.3).

\begin{proposition}[8.6.3]
\label{IV.8.6.3-fr}
Supposons $S_0$ quasi-compact (resp. quasi-compact et quasi-séparé). Alors les
applications (8.6.2.1), (8.6.2.2), (8.6.2.3) sont injectives (resp.
bijectives).
\end{proposition}

\begin{proof}
Compte tenu des remarques de (8.6.1), les assertions relatives à (8.6.2.3)
découlent de (8.5.11) appliqué à $\mathcal F_\lambda=\mathcal O_{S_\lambda}$~;
de même, les assertions relatives à (8.6.2.2) sont des cas particuliers de
(8.3.5) et (8.3.11), compte tenu de ce que les $S_\lambda$ et $S$ sont
quasi-compacts. Reste à considérer l'application (8.6.2.1). Prouvons d'abord
qu'elle est surjective lorsque $S_0$ est quasi-compact et quasi-séparé. Soit
$Z$ un sous-préschéma de $S$, de présentation finie sur $S$~; comme $S$ est
quasi-compact, il en est de même de $Z$, donc il existe un ouvert
quasi-compact $U$ de $S$ tel que $Z$ soit un sous-préschéma fermé de $U$, de
présentation finie sur $U$. Il existe alors un indice $\lambda$ et un ouvert
quasi-compact $U_\lambda$ de $S_\lambda$ tels que
$U=u_\lambda^{-1}(U_\lambda)$ (8.2.11)~; comme $S_\lambda$ est quasi-séparé, il
en est de même de $U_\lambda$ (I, 1.2.7), et par suite on peut se borner au cas
où $U=S$~; mais dans ce cas, on est ramené au fait que (8.6.2.3) est
surjective.

Enfin, pour voir que (8.6.2.1) est injective lorsque $S_0$ est quasi-compact,
il suffira de prouver le résultat plus précis suivant~:
\end{proof}

\begin{corollary}[8.6.3.1]
\label{IV.8.6.3.1-fr}
Supposons $S_0$ quasi-compact et soient $Z'_\lambda$, $Z''_\lambda$ deux
sous-préschémas de $S_\lambda$, de présentation finie sur $S_\lambda$. Pour
que $Z'=u_\lambda^{-1}(Z'_\lambda)$ soit majoré par
$Z''=u_\lambda^{-1}(Z''_\lambda)$ (I, 4.1.10), il faut et il suffit qu'il
existe $\mu\geq\lambda$ tel que
$Z'_\mu=u_{\lambda\mu}^{-1}(Z'_\lambda)$ soit majoré par
$Z''_\mu=u_{\lambda\mu}^{-1}(Z''_\lambda)$.
\end{corollary}

\begin{proof}
Il est trivial que la condition est suffisante. Pour voir qu'elle est
nécessaire, notons d'abord que les ensembles sous-jacents $Z'_\lambda$ et
$Z''_\lambda$ sont localement constructibles dans $S_\lambda$ par hypothèse
(I, 1.8.4), donc l'hypothèse entraîne l'existence d'un $\nu\geq\lambda$ tel
que $Z'_\nu\subset Z''_\nu$ (8.3.5)~; remplaçant $\lambda$ par $\nu$, on peut
donc déjà supposer que l'on a $Z'_\lambda\subset Z''_\lambda$. En outre, par
hypothèse (I, 1.6.1), les sous-espaces $Z'_\lambda$ et $Z''_\lambda$ de
$S_\lambda$ sont quasi-compacts~; pour tout point $x\in Z'_\lambda$, il y a un
voisinage ouvert quasi-compact $V(x)$ dans $S_\lambda$ tels que
$V(x)\cap Z'_\lambda$ et $V(x)\cap Z''_\lambda$ soient fermés dans $V(x)$. En
recouvrant $S_\lambda$ par un nombre fini de voisinages $V(x_i)$ on voit donc
qu'il y a un ouvert quasi-compact $U_\lambda$ de $S_\lambda$ contenant
$Z'_\lambda$ et tel que $Z'_\lambda$ et $Z''_\lambda\cap U_\lambda$ soient
fermés dans $U_\lambda$. Si l'on désigne par $Y_\lambda$ le sous-préschéma
induit par $Z''_\lambda$ sur $U_\lambda\cap Z''_\lambda$, il est clair qu'avec
les notations habituelles, $Y_\mu$ (resp. $Y$) est induit par $Z''_\mu$ sur
$U_\mu\cap Z''_\mu$ pour $\mu\geq\lambda$ (resp. par $Z''$ sur
$U\cap Z''$)~; en outre $Z'$ est majoré par $Y$ (I, 4.4.1), et comme il suffit
de prouver que $Z'_\mu$ est majoré par $Y_\mu$ pour $\mu$
\oldpage[IV-3]{27}
assez grand, on voit finalement qu'on est ramené (en remplaçant $S_\lambda$
par $U_\lambda$) au cas où $Z'_\lambda$ et $Z''_\lambda$ sont des
sous-préschémas fermés de $S_\lambda$. Mais alors, cela a déjà été démontré
puisque (8.6.2.3) est croissante et injective.
\end{proof}

\begin{corollary}[8.6.4]
\label{IV.8.6.4-fr}
Supposons $S_0$ quasi-compact, et soit $Z_\lambda$ un sous-préschéma de
$S_\lambda$, de présentation finie sur $S_\lambda$. Pour que
$Z=u_\lambda^{-1}(Z_\lambda)$ soit un sous-préschéma induit sur un ouvert
(resp. un sous-préschéma fermé) de $S$, il faut et il suffit qu'il existe
$\mu\geq\lambda$ tel que $Z_\mu=u_{\lambda\mu}^{-1}(Z_\lambda)$ soit induit
sur un ouvert (resp. un sous-préschéma fermé) dans $S_\mu$.
\end{corollary}

\begin{proof}
Soit $(U_\lambda^{(i)})$ un recouvrement ouvert affine fini de $S_\lambda$, et
posons $U_\mu^{(i)}=u_{\lambda\mu}^{-1}(U_\lambda^{(i)})$ pour
$\mu\geq\lambda$ et $U^{(i)}=u_\lambda^{-1}(U_\lambda^{(i)})$. Si $Z$ est
ouvert (resp. fermé) dans $S$, $Z\cap U^{(i)}$ l'est dans $U^{(i)}$, et
inversement si chacun des $Z_\mu\cap U_\mu^{(i)}$ est ouvert (resp. fermé)
dans $U_\mu^{(i)}$, $Z_\mu$ l'est dans $S_\mu$. Puisque $L$ est filtrant, il
suffit de démontrer le corollaire lorsque $S_\lambda$ est affine, donc
quasi-séparé. Mais alors le résultat découle de ce que les applications
(8.6.2.1), (8.6.2.2) et (8.6.2.3) sont bijectives.
\end{proof}

\subsection{Critères pour qu'une limite projective de préschémas soit un préschéma réduit (resp. intègre)}
\label{subsection:IV.8.7-fr}

Nous conservons les hypothèses et notations de (8.2.2) et supposons toujours
que $S_0$ est un des $S_\lambda$.

\begin{proposition}[8.7.1]
\label{IV.8.7.1-fr}
Supposons que $S$ soit non réduit. Alors il existe $\lambda_0$ tel que pour
$\lambda\geq\lambda_0$, $S_\lambda$ soit non réduit.
\end{proposition}

\begin{proof}
La question étant locale sur $S_0$, on peut supposer
$S_0=\operatorname{Spec}(A_0)$ affine, d'où $S=\operatorname{Spec}(A)$, où
$A=\varinjlim A_\lambda$ est la limite inductive d'un système inductif de
$A_0$-algèbres $(A_\lambda)$. On sait alors (5.13.2) que le nilradical de $A$
est la limite inductive de ceux des $A_\lambda$~; s'il n'est pas nul, un des
$A_\lambda$ contient donc un élément nilpotent $a_\lambda\neq0$ dont l'image
dans $A$ est un élément nilpotent et $\neq0$, et l'image de $a_\lambda$ dans
les $A_\mu$ pour $\mu\geq\lambda$ est par suite un élément nilpotent et
$\neq0$.
\end{proof}

\begin{proposition}[8.7.2]
\label{IV.8.7.2-fr}
Supposons vérifiée l'une des hypothèses suivantes~:
\begin{enumerate}
\item[\emph{a)}] $S_0$ est quasi-compact, le nilradical $\mathcal N_0$ de
$\mathcal O_{S_0}$ est un Idéal de type fini (ce qui sera par exemple vérifié
lorsque $S_0$ est noethérien), et les morphismes
$u_{\lambda\mu}:S_\mu\to S_\lambda$ sont des immersions ouvertes.
\item[\emph{b)}] Les morphismes $u_{\lambda\mu}:S_\mu\to S_\lambda$ sont
fidèlement plats.
\end{enumerate}
Sous ces conditions, pour que $S$ soit réduit, il faut et il suffit qu'il
existe $\lambda_0$ tel que $S_\lambda$ soit réduit pour
$\lambda\geq\lambda_0$.

En outre, dans le cas \emph{b)}, si $S$ est réduit, tous les $S_\lambda$ le
sont.
\end{proposition}

\begin{proof}
La dernière assertion résulte de ce que le morphisme $S\to S_\lambda$ est
alors fidèlement plat pour tout $\lambda$ (8.3.8) et de (2.1.13). D'autre part,
(8.7.1) prouve que la condition de l'énoncé est suffisante (sans hypothèse sur
$S_0$ ni sur les $u_{\lambda\mu}$). Il reste donc à voir que la condition est
nécessaire dans l'hypothèse \emph{a)}~; alors (8.2.13), l'espace sous-jacent à
$S$ s'identifie à l'intersection des espaces sous-jacents aux $S_\lambda$ (les
$S_\lambda$ étant identifiés à des sous-préschémas induits sur des ouverts de
$S_0$), et le faisceau structural $\mathcal O_S$ s'identifie au
\oldpage[IV-3]{28}
faisceau induit sur $S$ par tous les $\mathcal O_{S_\lambda}$~; en particulier
pour tout $s\in S$, l'anneau local $\mathcal O_s$ est \emph{le même} pour $S$
et pour tous les $S_\lambda$. Si $\mathcal N_\lambda$ est le Nilradical de
$\mathcal O_{S_\lambda}$, le Nilradical $\mathcal N$ de $\mathcal O_S$ a donc
en chaque point de $S$ la même fibre $\mathcal N_s$ (nilradical de
$\mathcal O_s$) que $u_\lambda^*(\mathcal N_\lambda)$ (induit sur $S$ par
$\mathcal N_\lambda$). L'hypothèse que $S$ est réduit entraîne par suite
$u_\lambda^*(\mathcal N_\lambda)=0$~; comme $\mathcal N_0$ est supposé de
type fini, il en est de même des $\mathcal N_\lambda$, et la conclusion résulte
donc de (8.5.7).
\end{proof}

\begin{corollary}[8.7.3]
\label{IV.8.7.3-fr}
Supposons vérifiée l'une des hypothèses suivantes~:
\begin{enumerate}
\item[\emph{a)}] $S_0$ est un préschéma noethérien et les morphismes
$u_{\lambda\mu}:S_\mu\to S_\lambda$ sont des immersions ouvertes.
\item[\emph{b)}] Les morphismes $u_{\lambda\mu}:S_\mu\to S_\lambda$ sont
fidèlement plats.
\end{enumerate}
Alors, pour que $S$ soit intègre, il faut et il suffit qu'il existe
$\lambda_0$ tel que $S_\lambda$ soit intègre pour $\lambda\geq\lambda_0$.
\end{corollary}

\begin{proof}
Dire qu'un préschéma est intègre signifie qu'il est à la fois réduit et
irréductible~; le corollaire résulte donc de (8.7.2) et de (8.4.3).
\end{proof}

\begin{remark}[8.7.4]
\label{IV.8.7.4-fr}
Si l'on ne fait pas d'hypothèse sur les $u_{\lambda\mu}$, il peut se faire que
$S$ soit intègre bien que tous les $S_\lambda$ soient non réduits et non
connexes, comme le montre l'exemple (8.4.4), où on prend l'anneau $A$ non
réduit.
\end{remark}

\subsection{Préschémas de présentation finie sur une limite projective de préschémas}
\label{subsection:IV.8.8-fr}

\begin{env}[8.8.1]
\label{IV.8.8.1-fr}
Conservant toujours les notations et hypothèses de (8.2.2), nous supposerons
donnés dans cette section deux $S_\alpha$-préschémas $X_\alpha$, $Y_\alpha$,
ce qui définit (8.2.5) deux systèmes projectifs de préschémas
$(X_\lambda,v_{\lambda\mu})$ et $(Y_\lambda,w_{\lambda\mu})$ en posant
$X_\lambda=X_\alpha\times_{S_\alpha}S_\lambda$,
$Y_\lambda=Y_\alpha\times_{S_\alpha}S_\lambda$,
$v_{\lambda\mu}=1_{X_\alpha}\times u_{\lambda\mu}$,
$w_{\lambda\mu}=1_{Y_\alpha}\times u_{\lambda\mu}$ (pour
$\alpha\leq\lambda\leq\mu$), dont les limites projectives sont respectivement
$X=X_\alpha\times_{S_\alpha}S$, $Y=Y_\alpha\times_{S_\alpha}S$, les
morphismes canoniques $v_\lambda:X\to X_\lambda$ et
$w_\lambda:Y\to Y_\lambda$ étant respectivement égaux à
$1_{X_\alpha}\times u_\lambda$ et $1_{Y_\alpha}\times u_\lambda$. Pour
$\alpha\leq\lambda\leq\mu$, on a une application canonique
$e_{\mu\lambda}:\operatorname{Hom}_{S_\lambda}(X_\lambda,Y_\lambda)\to
\operatorname{Hom}_{S_\mu}(X_\mu,Y_\mu)$, qui à tout $S_\lambda$-morphisme
$f_\lambda:X_\lambda\to Y_\lambda$ fait correspondre
$f_\mu=f_\lambda\times1_{S_\mu}:X_\lambda\times_{S_\lambda}S_\mu\to
Y_\lambda\times_{S_\lambda}S_\mu$, et il est clair que
$(\operatorname{Hom}_{S_\lambda}(X_\lambda,Y_\lambda),e_{\mu\lambda})$ est un
\emph{système inductif} d'ensembles. De même, on a une application canonique
$e_\lambda:\operatorname{Hom}_{S_\lambda}(X_\lambda,Y_\lambda)\to
\operatorname{Hom}_S(X,Y)$ qui à $f_\lambda$ fait correspondre
$f=f_\lambda\times1_S:X_\lambda\times_{S_\lambda}S\to
Y_\lambda\times_{S_\lambda}S$ et $(e_\lambda)$ est un système inductif
d'applications~; d'où, en passant à la limite inductive, une application
canonique, fonctorielle en $S_\alpha$, $X_\alpha$ et $Y_\alpha$~:
\begin{equation}
e:\varinjlim\operatorname{Hom}_{S_\lambda}(X_\lambda,Y_\lambda)
\longrightarrow\operatorname{Hom}_S(X,Y).
\tag{8.8.1.1}\label{IV.8.8.1.1-fr}
\end{equation}
\end{env}

\begin{theorem}[8.8.2]
\label{IV.8.8.2-fr}
\begin{enumerate}
\item[(i)] Supposons $X_\alpha$ quasi-compact (resp. quasi-compact et
quasi-séparé), et $Y_\alpha$ localement de type fini (resp. localement de
présentation finie) sur $S_\alpha$. Alors l'application (8.8.1.1) est
injective (resp. bijective).
\item[(ii)] Supposons $S_0$ quasi-compact et quasi-séparé. Pour tout préschéma
$X$ de présentation finie sur $S$, il existe un $\lambda\in L$, un préschéma
$X_\lambda$ de présentation finie sur $S_\lambda$, et un $S$-isomorphisme
$X\xrightarrow{\sim}X_\lambda\times_{S_\lambda}S$.
\end{enumerate}
\end{theorem}

\begin{proof}
\oldpage[IV-3]{29}
(i) Considérons d'abord le cas où $S_0=\operatorname{Spec}(A_0)$,
$X_\alpha=\operatorname{Spec}(B_\alpha)$ et
$Y_\alpha=\operatorname{Spec}(C_\alpha)$ sont affines~; alors les
$S_\lambda=\operatorname{Spec}(A_\lambda)$ et $S=\operatorname{Spec}(A)$ sont
aussi affines, avec $A=\varinjlim A_\lambda$, et les assertions de (i) sont
équivalentes au
\end{proof}

\begin{lemma}[8.8.2.1]
\label{IV.8.8.2.1-fr}
Soient $A_0$ un anneau, $(A_\lambda)_{\lambda\in L}$ un système inductif de
$A_0$-algèbres, $A=\varinjlim A_\lambda$~; soient $B_\alpha$ une
$A_\alpha$-algèbre, $C_\alpha$ une $A_\alpha$-algèbre de type fini (resp. de
présentation finie). Alors l'homomorphisme canonique
\begin{equation}
\varinjlim_\lambda
\operatorname{Hom}_{A_\lambda\text{-alg.}}
  (C_\alpha\otimes_{A_\alpha}A_\lambda,
   B_\alpha\otimes_{A_\alpha}A_\lambda)
\longrightarrow
\operatorname{Hom}_{A\text{-alg.}}
  (C_\alpha\otimes_{A_\alpha}A,
   B_\alpha\otimes_{A_\alpha}A)
\tag{8.8.2.2}\label{IV.8.8.2.2-fr}
\end{equation}
est injectif (resp. bijectif).
\end{lemma}

\begin{proof}
On sait que l'on a des isomorphismes canoniques fonctoriels
\[
\operatorname{Hom}_{A_\lambda\text{-alg.}}
 (C_\alpha\otimes_{A_\alpha}A_\lambda,
  B_\alpha\otimes_{A_\alpha}A_\lambda)
\xrightarrow{\sim}
\operatorname{Hom}_{A_\alpha\text{-alg.}}
 (C_\alpha,B_\alpha\otimes_{A_\alpha}A_\lambda)
\]
\[
\operatorname{Hom}_{A\text{-alg.}}
 (C_\alpha\otimes_{A_\alpha}A,
  B_\alpha\otimes_{A_\alpha}A)
\xrightarrow{\sim}
\operatorname{Hom}_{A_\alpha\text{-alg.}}
 (C_\alpha,B_\alpha\otimes_{A_\alpha}A)
\]
en vertu de la propriété universelle du produit tensoriel de deux algèbres. Il
suffit donc de prouver le
\end{proof}

\begin{lemma}[8.8.2.3]
\label{IV.8.8.2.3-fr}
Soient $E$ un anneau, $G$ une $E$-algèbre de type fini (resp. de présentation
finie), $(F_\lambda)$ un système inductif de $E$-algèbres. Alors
l'homomorphisme canonique
\[
\varinjlim\operatorname{Hom}_{E\text{-alg.}}(G,F_\lambda)
\longrightarrow
\operatorname{Hom}_{E\text{-alg.}}(G,\varinjlim F_\lambda)
\]
qui, à tout système inductif d'homomorphismes
$\theta_\lambda:G\to F_\lambda$ de $E$-algèbres, fait correspondre sa limite
inductive, est injectif (resp. bijectif).
\end{lemma}

\begin{proof}
Supposons d'abord la $E$-algèbre $G$ de type fini, et soit
$(t_i)_{1\leq i\leq n}$ un système de générateurs de cette $E$-algèbre~;
montrons que si $(\theta'_\lambda)$, $(\theta''_\lambda)$ sont deux systèmes
inductifs d'homomorphismes $G\to F_\lambda$ tels que
$\varinjlim\theta'_\lambda=\varinjlim\theta''_\lambda$, il existe un
$\mu\geq\lambda$ tel que $\theta'_\mu=\theta''_\mu$. En effet, si
$\varphi_{\mu\lambda}:F_\lambda\to F_\mu$ et
$\varphi_\lambda:F_\lambda\to F=\varinjlim F_\lambda$ sont les
homomorphismes canoniques du système inductif $(F_\lambda)$, par hypothèse,
pour chaque indice $i$, il existe un indice $\lambda_i$ tel que
$\varphi_{\lambda_i}(\theta'_{\lambda_i}(t_i))=
\varphi_{\lambda_i}(\theta''_{\lambda_i}(t_i))$, et l'on peut supposer tous
les $\lambda_i$ égaux à un même $\lambda$~; il en résulte de la même manière
l'existence d'un $\mu\geq\lambda$ tel que
$\varphi_{\mu\lambda}(\theta'_\lambda(t_i))=
\varphi_{\mu\lambda}(\theta''_\lambda(t_i))$ pour $1\leq i\leq n$,
c'est-à-dire $\theta'_\mu(t_i)=\theta''_\mu(t_i)$ pour $1\leq i\leq n$,
d'où $\theta'_\mu=\theta''_\mu$.

Supposons en second lieu $G$ de présentation finie, de sorte que l'on a
\[
G=E[T_1,\ldots,T_n]/\mathfrak J,
\]
où $\mathfrak J$ est un idéal de type fini, $t_i$ étant la classe de $T_i$
mod. $\mathfrak J$. Soit $(P_j)_{1\leq j\leq m}$ un système de générateurs de
$\mathfrak J$. Supposons donné un homomorphisme de $E$-algèbres
$\theta:G\to F$~; posons $b_i=\theta(t_i)$~; par définition, on a donc
$P_j(b_1,b_2,\ldots,b_n)=\theta(P_j(t_1,\ldots,t_n))=0$ pour
$1\leq j\leq m$. Or, il existe un $\lambda$ et des éléments
$x_1,\ldots,x_n$ dans $F_\lambda$ tels que $b_i=\varphi_\lambda(x_i)$ pour
$1\leq i\leq n$~; on a donc
$\varphi_\lambda(P_j(x_1,\ldots,x_n))=P_j(b_1,\ldots,b_n)=0$, et par suite il
existe $\mu\geq\lambda$ tel que
$\varphi_{\mu\lambda}(P_j(x_1,\ldots,x_n))=
P_j(\varphi_{\mu\lambda}(x_1),\ldots,\varphi_{\mu\lambda}(x_n))=0$ pour
$1\leq j\leq m$~; on en conclut qu'il existe un homomorphisme de
$E$-algèbres $\theta_\mu:G\to F_\mu$ tel que
$\theta_\mu(t_i)=\varphi_{\mu\lambda}(x_i)$
\oldpage[IV-3]{30}
pour $1\leq i\leq n$~; on en déduit pour tout $\nu\geq\mu$ un
homomorphisme de $E$-algèbres
$\theta_\nu=\varphi_{\nu\mu}\circ\theta_\mu$ de $G$ dans $F_\nu$, et il est
clair que $\theta$ est la limite inductive de ce système inductif
d'homomorphismes, ce qui achève de démontrer le lemme.

Passons maintenant au cas où $X_\alpha$ est quasi-compact et $Y_\alpha$
localement de type fini sur $S_\alpha$. Posons
$Z_\alpha=X_\alpha\times_{S_\alpha}Y_\alpha$ et introduisons le système
projectif correspondant des
$Z_\lambda=Z_\alpha\times_{S_\alpha}S_\lambda=
X_\lambda\times_{S_\lambda}Y_\lambda$ et sa limite
$Z=Z_\alpha\times_{S_\alpha}S=X\times_S Y$~; les bijections canoniques
(I, 3.3.14) donnent des diagrammes commutatifs
\[
\begin{tikzcd}[column sep=large,row sep=large]
\operatorname{Hom}_{S_\lambda}(X_\lambda,Y_\lambda)
  \ar[r,"e_{\mu\lambda}"] \ar[d,"\wr"] &
\operatorname{Hom}_{S_\mu}(X_\mu,Y_\mu)
  \ar[r,"e_\mu"] \ar[d,"\wr"] &
\operatorname{Hom}_S(X,Y) \ar[d,"\wr"] \\
\operatorname{Hom}_{X_\lambda}(X_\lambda,Z_\lambda)
  \ar[r,"e_{\mu\lambda}"'] &
\operatorname{Hom}_{X_\mu}(X_\mu,Z_\mu)
  \ar[r,"e_\mu"'] &
\operatorname{Hom}_X(X,Z)
\end{tikzcd}
\]
et par suite on est ramené à prouver que (8.8.1.1) est injective dans le cas
particulier où $S_\alpha=X_\alpha$ (compte tenu de (I, 3.4)). De plus, comme
$X_\alpha$ est quasi-compact, donc réunion finie d'ouverts affines, on peut
supposer $X_\alpha$ affine ($L$ étant filtrant). Supposons donc donnés deux
$X_\alpha$-morphismes $f'_\alpha:X_\alpha\to Y_\alpha$,
$f''_\alpha:X_\alpha\to Y_\alpha$ tels que $f'$ et $f''$ soient des
$X$-morphismes égaux de $X$ dans $Y$~; il s'agit de prouver que
$f'_\mu=f''_\mu$ pour $\mu\geq\alpha$ assez grand. Comme $X_\alpha$ est
quasi-compact, il en est de même de
$f'_\alpha(X_\alpha)\cup f''_\alpha(X_\alpha)$, et comme $Y_\alpha$ est
localement de type fini sur $X_\alpha$,
$f'_\alpha(X_\alpha)\cup f''_\alpha(X_\alpha)$ est contenu dans une réunion
finie d'ouverts affines $V_{i\alpha}$ de $Y_\alpha$, de type fini sur
$X_\alpha$. Posons
$U'_{i\alpha}=f_\alpha^{'-1}(V_{i\alpha})$,
$U''_{i\alpha}=f_\alpha^{''-1}(V_{i\alpha})$,
$U_{i\alpha}=U'_{i\alpha}\cap U''_{i\alpha}$,
$U_\alpha=\bigcup_iU_{i\alpha}$. L'hypothèse $f'=f''$ entraîne
$v_\alpha^{-1}(U'_{i\alpha})=v_\alpha^{-1}(U''_{i\alpha})$, ces deux ensembles
étant respectivement égaux à $f^{'-1}(w_\alpha^{-1}(V_{i\alpha}))$ et à
$f^{''-1}(w_\alpha^{-1}(V_{i\alpha}))$. Comme les $V_{i\alpha}$ recouvrent
$Y_\alpha$, on a $v_\alpha^{-1}(U_\alpha)=f^{'-1}(Y)=X$. Comme $X_\alpha$
est quasi-compact et que toute partie ouverte de $X_\lambda$ est
ind-constructible, il résulte de (8.3.4) qu'il y a un indice
$\lambda\geq\alpha$ tel que les
$U_{i\mu}=v_{\alpha\mu}^{-1}(U_{i\alpha})$ forment un \emph{recouvrement} de
$X_\mu$. Remplaçant $\alpha$ par $\mu$, on peut donc supposer que les
$U_{i\alpha}$ recouvrent $X_\alpha$~; cela entraîne que pour tout
$x\in X_\alpha$, il y a un voisinage ouvert affine $W(x)$ contenu dans un des
$U_{i\alpha}$, autrement dit tel que les restrictions de $f'_\alpha$ et
$f''_\alpha$ à $W(x)$ appliquent $W(x)$ dans un \emph{même} ouvert affine
$V_{i\alpha}$. Comme $X_\alpha$ est quasi-compact, il est recouvert par un
nombre fini d'ouverts affines $W(x_k)$~; en vertu du lemme (8.8.2.1) et du fait
que $L$ est filtrant, il existe un $\lambda\geq\alpha$ tel que les
restrictions de $f'_\lambda$ et $f''_\lambda$ à chacun des ouverts
$v_{\alpha\lambda}^{-1}(W(x_k))$ soient égales, d'où
$f'_\lambda=f''_\lambda$.

Supposons maintenant $X_\alpha$ quasi-compact et quasi-séparé et $Y_\alpha$
localement de présentation finie sur $S_\alpha$, et prouvons que (8.8.1.1) est
surjective. Supposons donc donné un $X$-morphisme $f:X\to Y$. Comme $X$ est
quasi-compact, il en est de même
\oldpage[IV-3]{31}
de $f(X)$, et par suite il y a un ouvert quasi-compact $Y'$ dans $Y$ qui
contient $f(X)$~; il existe par suite un $\lambda\geq\alpha$ et un ouvert
quasi-compact $Y'_\lambda$ dans $Y_\lambda$ tels que
$Y'=w_\lambda^{-1}(Y'_\lambda)$ (8.2.11). Remplaçant au besoin $\alpha$ par
$\lambda$ et $Y_\lambda$ par $Y'_\lambda$, on peut donc se borner au cas où
$Y_\alpha$ est quasi-compact, donc aussi $Y$ et les $Y_\lambda$. Comme $Y$ est
quasi-compact, il est réunion finie d'ouverts affines $V_i$, et par suite $X$
est réunion des ouverts $f^{-1}(V_i)$. Comme tout point de $X$ a un voisinage
ouvert quasi-compact contenu dans un des $f^{-1}(V_i)$ et que $X$ est
quasi-compact, on peut, en tenant compte de (8.2.11) et remplaçant au besoin
$\alpha$ par un indice $\lambda\geq\alpha$, supposer que $Y$ est réunion finie
d'ouverts $V_i=w_\alpha^{-1}(V_{i\alpha})$, où les $V_{i\alpha}$ sont des
ouverts affines de $Y_\alpha$~; par suite $X$ est réunion des ouverts
$f^{-1}(V_i)$. Comme tout point de $X$ a un voisinage ouvert quasi-compact
contenu dans l'un des $f^{-1}(V_i)$ et que $X$ est quasi-compact, on peut
recouvrir $X$ par un nombre fini de tels voisinages, et (en répétant au besoin
certains des $V_i$) supposer que l'on a un recouvrement $(U_i)$ de $X$ par des
ouverts quasi-compacts ayant même ensemble d'indices que $(V_i)$ et tel que
$f(U_i)\subset V_i$ pour tout $i$. En outre, à l'aide de (8.2.11) (et en
remplaçant au besoin $\alpha$ par un indice $\lambda\geq\alpha$), on peut
supposer que l'on a $U_i=v_\alpha^{-1}(U_{i\alpha})$ où les $U_{i\alpha}$ sont
des ouverts quasi-compacts dans $X_\alpha$~; de plus, utilisant (8.3.4) comme
ci-dessus, on peut supposer que $(U_{i\alpha})$ est un \emph{recouvrement} de
$X_\alpha$.

Cela étant, il suffira de montrer qu'il existe un $\lambda\geq\alpha$ et pour
chaque $i$ un morphisme $f_{i\lambda}:U_{i\lambda}\to V_{i\lambda}$ (avec
$U_{i\lambda}=v_{\alpha\lambda}^{-1}(U_{i\alpha})$ et
$V_{i\lambda}=w_{\alpha\lambda}^{-1}(V_{i\alpha})$) tels que le morphisme
correspondant $f_i=e_\lambda(f_{i\lambda}):U_i\to V_i$ soit égal à la
\emph{restriction} de $f$ à $U_i$. En effet, s'il en est ainsi, comme
$X_\lambda$ est quasi-séparé (I, 1.2.3), les
$U_{i\lambda}\cap U_{j\lambda}$ sont \emph{quasi-compacts} et le résultat
d'unicité démontré plus haut (qui s'applique puisque $Y_\lambda$ est
localement de type fini sur $S_\lambda$) prouve qu'il existe
$\mu\geq\lambda$ tel que $f_{i\mu}$ et $f_{j\mu}$ coïncident dans
$U_{i\mu}\cap U_{j\mu}$ pour tout couple $(i,j)$, et par suite définissent un
$S_\mu$-morphisme $f_\mu:X_\mu\to Y_\mu$ tel que $e_\mu(f_\mu)=f$.

On est ainsi ramené au cas où $Y_\alpha$ est \emph{affine}, et comme en outre
on peut supposer que les $V_{i\alpha}$ ont des images dans $S_\alpha$
contenues dans des ouverts affines, on peut aussi supposer que $S_\alpha$ est
affine, soient donc $S_\alpha=\operatorname{Spec}(A_\alpha)$,
$Y_\alpha=\operatorname{Spec}(C_\alpha)$, $C_\alpha$ étant une
$A_\alpha$-algèbre de présentation finie, $S=\operatorname{Spec}(A)$,
$Y=\operatorname{Spec}(C)$, avec $A=\varinjlim A_\lambda$,
$C=\varinjlim(C_\alpha\otimes_{A_\alpha}A_\lambda)=
C_\alpha\otimes_{A_\alpha}A$. On a alors
\[
\operatorname{Hom}_S(X,Y)=
\operatorname{Hom}_{A\text{-alg.}}(C,\Gamma(X,\mathcal O_X))=
\operatorname{Hom}_{A_\alpha\text{-alg.}}(C_\alpha,\Gamma(X,\mathcal O_X))
\]
(I, 2.2.4) et de même
\[
\operatorname{Hom}_{S_\lambda}(X_\lambda,Y_\lambda)=
\operatorname{Hom}_{A_\lambda\text{-alg.}}
(C_\alpha\otimes_{A_\alpha}A_\lambda,\Gamma(X_\lambda,\mathcal O_{X_\lambda}))=
\operatorname{Hom}_{A_\alpha\text{-alg.}}
(C_\alpha,\Gamma(X_\lambda,\mathcal O_{X_\lambda})).
\]
Mais comme $X_\alpha$ est quasi-compact et quasi-séparé, on sait (8.5.4) que
l'on a
$\varinjlim\Gamma(X_\lambda,\mathcal O_{X_\lambda})=
\Gamma(X,\mathcal O_X)$~; puisque $C_\alpha$ est une $A_\alpha$-algèbre de
présentation finie, le fait que (8.8.1.1) est bijectif résulte alors de
(8.8.2.3).

Avant de passer à la démonstration de (ii), notons les corollaires suivants de
(i)~:
\end{proof}

\begin{corollary}[8.8.2.4]
\label{IV.8.8.2.4-fr}
Supposons $S_0$ quasi-compact, $X_\alpha$ de présentation finie sur $S_\alpha$,
$Y_\alpha$ de type fini sur $S_\alpha$ et quasi-séparé sur $S_\alpha$ (ce qui
sera par exemple le cas si $Y_\alpha$ est également
\oldpage[IV-3]{32}
de présentation finie sur $S_\alpha$). Soit
$f_\alpha:X_\alpha\to Y_\alpha$ un $S_\alpha$-morphisme. Pour que
$f:X\to Y$ soit un isomorphisme, il faut et il suffit qu'il existe
$\lambda\geq\alpha$ tel que $f_\lambda:X_\lambda\to Y_\lambda$ soit un
isomorphisme.
\end{corollary}

\begin{proof}
La condition est évidemment suffisante. Pour prouver qu'elle est nécessaire,
notons d'abord que la question étant locale sur $S_0$ (puisque $L$ est
filtrant), on peut toujours supposer $S_0$ affine, donc quasi-séparé. Il y a
par hypothèse un $S$-morphisme $g:Y\to X$ tel que $g\circ f=1_X$ et
$f\circ g=1_Y$. Comme $X_\alpha$ est de présentation finie sur $S_\alpha$,
et $Y_\alpha$ quasi-compact et quasi-séparé (puisque $S_\alpha$ est
quasi-compact et quasi-séparé), il existe un $\mu\geq\alpha$ et un
$S_\mu$-morphisme $g_\mu:Y_\mu\to X_\mu$ tels que
$g=g_\mu\times1_S$ en vertu de (8.8.2, (i)). D'autre part, il résulte aussi de
(8.8.2, (i)) et des relations $g\circ f=1_X$ et $f\circ g=1_Y$ qu'il existe
$\nu\geq\mu$ tel que l'on ait $g_\nu\circ f_\nu=1_{X_\nu}$ et
$f_\nu\circ g_\nu=1_{Y_\nu}$, puisque $X_\alpha$ et $Y_\alpha$ sont de type
fini sur $S_\alpha$ et quasi-compacts. Cela signifie que
$f_\nu:X_\nu\to Y_\nu$ est un isomorphisme, d'où le corollaire.
\end{proof}

\begin{corollary}[8.8.2.5]
\label{IV.8.8.2.5-fr}
Supposons $S_0$ quasi-compact et quasi-séparé, $X_\alpha$ et $Y_\alpha$ de
présentation finie sur $S_\alpha$. Pour que $X$ et $Y$ soient
$S$-isomorphes, il faut et il suffit qu'il existe $\lambda\geq\alpha$ tels que
$X_\lambda$ et $Y_\lambda$ soient $S_\lambda$-isomorphes. De plus, pour tout
$S$-isomorphisme $f:X\xrightarrow{\sim}Y$, il existe $\mu\geq\lambda$ et un
$S_\mu$-isomorphisme $f_\mu:X_\mu\to Y_\mu$ tels que
$f=f_\mu\times1_S$.
\end{corollary}

\begin{proof}
La condition est évidemment suffisante~; inversement, s'il existe un
$S$-isomorphisme $f:X\to Y$, il résulte de (8.8.2, (i)) que $f$ est de la
forme $f_\mu\times1_S$ pour un $\mu\geq\alpha$ et un homomorphisme
$f_\mu:X_\mu\to Y_\mu$~; mais comme $f$ est un isomorphisme, il résulte de
(8.8.2.4) qu'il existe $\nu\geq\mu$ tel que
$f_\nu:X_\nu\to Y_\nu$ soit un isomorphisme.
\end{proof}

\begin{proof}[Démonstration de (8.8.2), (ii)]
Considérons encore en premier lieu le cas où $S_0=\operatorname{Spec}(A_0)$ et
$X=\operatorname{Spec}(B)$ sont affines~; alors l'assertion de (ii) équivaut
au lemme (I, 8.4.2).

Pour démontrer (ii) dans le cas général, notons que $S$ est quasi-compact et
quasi-séparé~; comme $X$ est de présentation finie sur $S$ et $S$ affine sur
$S_0$, il existe donc un recouvrement fini $(U_i)$ de $S_0$ et, si $(W_i)$ est
le recouvrement ouvert affine de $S$ formé des images réciproques des $U_i$
par le morphisme $S\to S_0$, un recouvrement fini $(X_r)$ de $X$ par des
ouverts affines tel que l'image par $X\to S$ de chaque $X_r$ soit contenue
dans un $W_{i(r)}$~; l'anneau $A(X_r)$ est alors une algèbre de présentation
finie sur l'anneau $A(W_{i(r)})$ (I, 1.4.6). En vertu du lemme (I, 8.4.2) et du
fait que $L$ est filtrant, il existe un indice $\lambda\in L$ et, pour chaque
indice $r$, un schéma affine $Z_{r\lambda}$, de présentation finie sur l'image
réciproque $W_{i(r),\lambda}$ de $U_{i(r)}$ dans $S_\lambda$, et un
$S$-isomorphisme
$g_r:Z_{r\lambda}\times_{S_\lambda}S\xrightarrow{\sim}X_r$. Soit $Z_{rs}$
l'image réciproque par $g_r$ du sous-préschéma induit sur l'ouvert
$X_r\cap X_s$ de $X_r$, qui est quasi-compact puisque $X$ est quasi-séparé,
et désignons par $g'_{rs}$ la restriction
$Z_{rs}\xrightarrow{\sim}X_r\cap X_s$ de l'isomorphisme $g_r$. En vertu de
(8.8.2.5), il existe un $\mu\geq\lambda$ et, pour tout couple $(r,s)$, un
ouvert quasi-compact $Z_{rs\mu}$ de
$Z_{r\mu}=v_{\lambda\mu}^{-1}(Z_{r\lambda})$ tels que $Z_{rs}$ soit l'image
réciproque de $Z_{rs\mu}$~; d'ailleurs, comme $S_\mu$ est quasi-séparé, et
$W_{i(r),\mu}$ ouvert quasi-compact dans $S_\mu$, chacun des $Z_{rs\mu}$ est
de présentation finie sur $S_\mu$ (I, 1.6.2). Considérons alors, pour tout
couple $(r,s)$, l'isomorphisme
$h_{sr}=g_{sr}^{'-1}\circ g'_{rs}$ de $Z_{rs}$ sur $Z_{sr}$~; il résulte de
(8.8.2.4) qu'il existe un $\nu\geq\mu$ et, pour tout couple $(r,s)$, un
isomorphisme $h_{sr\nu}:Z_{rs\nu}\to Z_{sr\nu}$ tels que
$h_{sr}=h_{sr\nu}\times1_S$. Enfin, pour tout triplet $(r,s,t)$ désignons par
$h'_{sr}$ la restriction de $h_{sr}$ à
\oldpage[IV-3]{33}
$Z_{rs}\cap Z_{rt}$~; il résulte aussitôt des définitions que $h'_{sr}$ est un
isomorphisme de $Z_{rs}\cap Z_{rt}$ sur $Z_{sr}\cap Z_{st}$, et que l'on a la
relation $h'_{ts}\circ h'_{sr}=h'_{tr}$. En vertu de (8.8.2, (i)), il existe
donc $\rho\geq\nu$ tel que, pour tout triplet $(r,s,t)$, on ait
$h'_{ts\rho}\circ h'_{sr\rho}=h'_{tr\rho}$. On en conclut que les
isomorphismes $h_{sr\rho}$ vérifient la condition de recollement
$(0_{\mathrm I}, 4.1.7)$ et définissent donc un préschéma $X_\rho$ tel que $X$
soit isomorphe à $X_\rho\times_{S_\rho}S$. En outre, les $Z_{r\rho}$ sont de
présentation finie sur $S_\rho$ et, si on les identifie à des ouverts de
$X_\rho$, les intersections $Z_{r\rho}\cap Z_{s\rho}$, isomorphes aux
$Z_{rs\rho}$, sont quasi-compactes, donc (I, 1.6.3) $X_\rho$ est de
présentation finie sur $S_\rho$.
\end{proof}

\begin{env}[8.8.3]
\label{IV.8.8.3-fr}
Le contenu essentiel de (8.8.2) s'exprime encore en disant que si $S_0$ est
quasi-compact et quasi-séparé, la catégorie des $S$-préschémas de présentation
finie est déterminée à équivalence près par la donnée des catégories des
$S_\lambda$-préschémas de présentation finie, des foncteurs
$\rho_{\mu\lambda}:X_\lambda\rightsquigarrow
X_\lambda\times_{S_\lambda}S_\mu$ (pour $\lambda\leq\mu$) entre ces catégories
et des isomorphismes de transitivité
$\rho_{\nu\mu}\circ\rho_{\mu\lambda}\xrightarrow{\sim}\rho_{\nu\lambda}$. De
façon imagée, on peut dire que la donnée d'un $S$-préschéma de présentation
finie $X$ équivaut «~fonctoriellement~» à la donnée d'un
$S_\lambda$-préschéma de présentation finie $X'_\lambda$~; si un
$S_\mu$-préschéma de présentation finie $X''_\mu$ est tel que $X$ soit
$S$-isomorphe à $X''_\mu\times_{S_\mu}S$, il existe un $\nu$ tel que
$\lambda\leq\nu$, $\mu\leq\nu$ et tel que
$X'_\lambda\times_{S_\lambda}S_\nu$ et
$X''_\mu\times_{S_\mu}S_\nu$ soient $S_\nu$-isomorphes. Le fait que $L$ est
filtrant entraîne en outre que si on se donne une famille \emph{finie}
$(X^{(i)})_{i\in I}$ de $S$-préschémas de présentation finie, et une famille
\emph{finie} $(f^{(j)})_{j\in J}$ de $S$-morphismes entre ces préschémas
($f^{(j)}$ étant donc de la forme
$X^{(\sigma(j))}\to X^{(\tau(j))}$ où $\sigma$ et $\tau$ sont deux
applications de $J$ dans $I$), alors il y a un indice $\lambda\in L$, une
famille $(X_\lambda^{(i)})_{i\in I}$ de $S_\lambda$-préschémas de présentation
finie et une famille $(f_\lambda^{(j)})_{j\in J}$ de $S_\lambda$-morphismes
tels que $X^{(i)}$ s'identifie à
$X_\lambda^{(i)}\times_{S_\lambda}S$ et $f^{(j)}$ à
$f_\lambda^{(j)}\times1_S$ quels que soient $i$ et $j$~; en outre, si on a
une relation $f^{(j)}=f^{(k)}$, on peut supposer $\lambda$ choisi de telle
sorte que $f_\lambda^{(j)}=f_\lambda^{(k)}$.

Considérons en particulier trois $S$-préschémas de présentation finie $X$,
$Y$, $Z$ et deux $S$-morphismes $f:X\to Z$, $g:Y\to Z$, de sorte que le
produit $X\times_ZY$ relatif à ces morphismes est encore un $S$-préschéma de
présentation finie (I, 1.6.2). Alors il résulte de ce qui précède et de
(I, 3.3.11) que l'on peut écrire
$X\times_ZY=(X_\lambda\times_{Z_\lambda}Y_\lambda)\times_{S_\lambda}S$ pour
un $\lambda$ convenable, $X_\lambda$, $Y_\lambda$, $Z_\lambda$ étant des
$S_\lambda$-préschémas de présentation finie~; on peut donc dire que la
détermination des $S$-préschémas de présentation finie par la donnée des
$S_\lambda$-préschémas de présentation finie est «~\emph{compatible avec les
produits fibrés}~». On a vu d'autre part (8.6.3) que si $g$ est une immersion,
on peut supposer qu'il en est de même de $g_\lambda:Y_\lambda\to Z_\lambda$~;
identifiant $Y$ (resp. $Y_\lambda$) avec un sous-préschéma de $Z$ (resp.
$Z_\lambda$), on voit donc que l'on peut écrire, pour un $\lambda$ convenable,
$f^{-1}(Y)=f_\lambda^{-1}(Y_\lambda)\times_{S_\lambda}S$ (I, 4.4.1)~; il y a
donc aussi «~\emph{compatibilité avec la formation des images réciproques de
sous-préschémas}~». Plus particulièrement, si $f_1$, $f_2$ sont deux
$S$-morphismes de $X$ dans $Y$, on appelle \emph{noyau} de ce couple de
morphismes l'image réciproque $N$ du sous-préschéma diagonal de $Y\times_SY$
par le $S$-morphisme $(f_1,f_2)_S$~; on déduit de ce qui précède que l'on aura
pour $\lambda$ convenable $N=N_\lambda\times_{S_\lambda}S$, où $N_\lambda$
est le noyau du couple de $S_\lambda$-morphismes
$(f_{1\lambda},f_{2\lambda})$ de $X_\lambda$ dans $Y_\lambda$. Ces remarques
s'étendent à des produits finis quelconques et aux «~noyaux~» de systèmes
finis quelconques de $S$-morphismes d'un $S$-préschéma de présentation finie
\oldpage[IV-3]{34}
dans un autre~; on en conclut que d'une façon générale la formation des
$S$-préschémas de présentation finie par la donnée des
$S_\lambda$-préschémas de présentation finie est «~\emph{compatible avec les
limites projectives finies}~», une telle limite étant par définition le noyau
d'un système fini de morphismes d'un même $S$-préschéma dans un produit de
$S$-préschémas (T, 1.8).

Nous nous permettrons couramment par la suite de faire les traductions
impliquées par les propriétés précédentes (ainsi que par (8.3.11), (8.5.2) et
(8.6.3)) sans toujours nous astreindre à les justifier pas à pas comme
ci-dessus. Par exemple, la donnée d'un «~\emph{préschéma en groupes}~» $G$ de
présentation finie sur $S$ équivaut à la donnée d'un préschéma en groupes
$G_\lambda$ de présentation finie sur un $S_\lambda$ pour un $\lambda$ assez
grand~; écrire en effet la condition d'associativité pour le $S$-morphisme
«~loi de composition~» $G\times_SG\to G$ du préschéma en groupes revient à
écrire que le noyau de deux $S$-morphismes de la forme
$G\times_SG\times_SG\to G$ est égal à $G\times_SG\times_SG$ (II, 8.3.9), et
les autres conditions intervenant dans la définition d'un préschéma en
groupes s'interprètent de même.
\end{env}

\subsection{Premières applications à l'élimination des hypothèses noethériennes}
\label{subsection:IV.8.9-fr}

\begin{proposition}[8.9.1]
\label{IV.8.9.1-fr}
Soient $A$ un anneau, $X$ un $A$-préschéma.
\begin{enumerate}
\item[(i)] Les conditions suivantes sont équivalentes~:
\begin{enumerate}
\item[\emph{a)}] $X$ est de présentation finie sur $A$.
\item[\emph{b)}] Il existe un anneau noethérien $A_0$, un préschéma $X_0$ de
type fini sur $A_0$, un homomorphisme d'anneaux $A_0\to A$, et un
$A$-isomorphisme $X_0\otimes_{A_0}A\xrightarrow{\sim}X$.
\item[\emph{c)}] Il existe un sous-anneau $A_0$ de $A$, qui est une
$\mathbf Z$-algèbre de type fini, un préschéma $X_0$ de type fini sur $A_0$ et
un $A$-isomorphisme $X_0\otimes_{A_0}A\xrightarrow{\sim}X$.
\end{enumerate}
\item[(ii)] Si $\mathcal F$ est un $\mathcal O_X$-Module quasi-cohérent de
présentation finie, on peut supposer le sous-anneau $A_0$ tel qu'il existe un
$\mathcal O_{X_0}$-Module cohérent $\mathcal F_0$ tel que $\mathcal F$ soit
isomorphe à $\mathcal F_0\otimes_{A_0}A$~;
$\operatorname{Supp}(\mathcal F)$ est constructible et fermé dans $X$, et il
y a un sous-préschéma fermé $Z$ de $X$, ayant
$\operatorname{Supp}(\mathcal F)$ comme espace sous-jacent, tel que
l'immersion canonique $Z\to X$ soit de présentation finie.
\item[(iii)] Si $Y$ est un second $A$-préschéma de présentation finie, et
$f:X\to Y$ un $A$-morphisme, on peut supposer le sous-anneau $A_0$ de $A$ tel
qu'il existe un préschéma $Y_0$ de type fini sur $A_0$, un $A$-isomorphisme
$Y_0\otimes_{A_0}A\xrightarrow{\sim}Y$ et un $A_0$-morphisme
$f_0:X_0\to Y_0$ (nécessairement de type fini) tels que $f$ s'identifie à
$f_0\otimes1_A$.
\end{enumerate}
\end{proposition}

\begin{proof}
(i) Comme $A$ est limite inductive de ses sous-anneaux qui sont de type fini
sur $\mathbf Z$, \emph{a)} implique \emph{c)} en vertu de (8.8.2, (ii))~;
\emph{c)} implique \emph{b)} car une $\mathbf Z$-algèbre de type fini est un
anneau noethérien~; enfin, si $A_0$ est noethérien, un $A_0$-préschéma de type
fini est de présentation finie (I, 1.6.1), donc \emph{b)} implique \emph{a)}
en vertu de (I, 1.6.2, (iii)).

L'assertion (ii) se déduit immédiatement de (8.5.2, (ii)), (8.3.11) et
(I, 1.8.2) et l'assertion (iii) de (8.8.2, (i)).
\end{proof}

\begin{env}[8.9.2]
\label{IV.8.9.2-fr}
La proposition (8.9.1) et les résultats de (8.5.2) et (8.8.2) permettent
d'étendre dans de nombreux cas à des morphismes de présentation finie
$X\to Y$ des résultats prouvés par des techniques noethériennes en supposant
$Y$ \emph{localement noethérien}.
\oldpage[IV-3]{35}
Nous en verrons de nombreux exemples par la suite et nous nous bornerons ici
à donner quelques résultats de ce type.
\end{env}

\begin{proposition}[8.9.3]
\label{IV.8.9.3-fr}
Soient $A$ un anneau, $M$ un $A$-module de présentation finie. Alors tout
$A$-endomorphisme surjectif de $M$ est bijectif.
\end{proposition}

\begin{proof}
En effet, considérons $A$ comme limite inductive de ses sous-$\mathbf Z$-algèbres
de type fini. Il résulte de (8.5.2.6) qu'il existe une de ces sous-algèbres
$A_0$ et un $A_0$-module $M_0$ de présentation finie tels que $M$ soit
$A$-isomorphe à $M_0\otimes_{A_0}A$~; en outre, si $f:M\to M$ est un
$A$-endomorphisme surjectif, on peut supposer (8.5.2, (i)) qu'il existe un
$A_0$-endomorphisme $f_0:M_0\to M_0$ tel que $f=f_0\otimes1_A$~; enfin
(8.5.7) on peut supposer que $f_0$ est surjectif. Mais comme $A_0$ est
noethérien et $M_0$ un $A_0$-module de type fini, $M_0$ est un $A_0$-module
noethérien, donc (Bourbaki, \emph{Alg.}, chap. VIII, \S~2, n\textsuperscript{o}~2,
lemme 3) $f_0$ est bijectif, et il en est par suite de même de $f$.
\end{proof}

\begin{proposition}[8.9.4]
\label{IV.8.9.4-fr}
(«~théorème de platitude générique~»). --- Soient $Y$ un préschéma intègre,
$u:X\to Y$ un morphisme de type fini et localement de présentation finie,
$\mathcal F$ un $\mathcal O_X$-Module quasi-cohérent de présentation finie. Il
existe alors un ouvert non vide $U$ de $Y$ tel que
$\mathcal F|u^{-1}(U)$ soit plat sur $U$.
\end{proposition}

\begin{proof}
Le raisonnement du début de (6.9.1) ramène à prouver le

\begin{lemma}[8.9.4.1]
\label{IV.8.9.4.1-fr}
Soient $A$ un anneau intègre, $B$ une $A$-algèbre de présentation finie,
$M$ un $B$-module de présentation finie. Alors il existe un $f\neq0$ dans
$A$ tel que $M_f$ soit un $A_f$-module libre.
\end{lemma}

En effet, d'après (8.9.1) il y a une sous-$\mathbf Z$-algèbre de type fini
$A_0$ de $A$, une $A_0$-algèbre de type fini $B_0$ et un $B_0$-module de type
fini $M_0$ tels que $B$ soit $A$-isomorphe à $B_0\otimes_{A_0}A$ et que $M$
soit $B$-isomorphe à $M_0\otimes_{A_0}A$~; d'après (6.9.2), il existe
$f_0\neq0$ dans $A_0$ tel que $(M_0)_{f_0}$ soit un
$(A_0)_{f_0}$-module libre. Comme
$M_{f_0}=(M_0)_{f_0}\otimes_{A_0}A$ et
$A_{f_0}=(A_0)_{f_0}\otimes_{A_0}A$, $M_{f_0}$ est un $A_{f_0}$-module libre.
\end{proof}

\begin{corollary}[8.9.5]
\label{IV.8.9.5-fr}
Soient $S$ un préschéma quasi-compact et quasi-séparé, $u:X\to S$ un
morphisme de présentation finie, $\mathcal F$ un $\mathcal O_X$-Module
quasi-cohérent de présentation finie. Il existe alors une partition
$(S_i)_{1\leq i\leq n}$ de $S$ en un nombre fini d'ensembles localement fermés
dans $S$ telle que, pour $1\leq i\leq n$, il existe un sous-préschéma $S'_i$
de $S$, ayant pour espace sous-jacent $S_i$, de présentation finie sur $S$,
tel que si l'on pose $X_i=X\times_SS'_i$, le $\mathcal O_{X_i}$-Module
$\mathcal F_i:=\mathcal F\otimes_{\mathcal O_S}\mathcal O_{S'_i}$ soit plat
sur $S'_i$.
\end{corollary}

\begin{proof}
Considérons un recouvrement fini $(U_j)_{1\leq j\leq m}$ de $S$ par des
ouverts affines et définissons par récurrence $T_1=U_1$,
$T_j=U_j-\bigcup_{h<j}(U_j\cap U_h)$ pour $j\geq2$~; chaque $T_j$ est fermé
dans l'ouvert affine $U_j$, et les $T_j$ sont deux à deux sans point commun~;
en outre les $U_j\cap U_h$ sont quasi-compacts puisque $S$ est quasi-séparé,
et par suite ($S$ étant aussi quasi-compact), sont constructibles dans $S$
(I, 1.8.1), donc il en est de même des $T_j$. Il suffira évidemment de
démontrer le corollaire pour un sous-préschéma convenable $T'_j$ de $S$ ayant
$T_j$ pour espace sous-jacent, de présentation finie sur $S$ et pour le
morphisme et le Module déduits respectivement de $u$ et $\mathcal F$ par le
changement de base $T'_j\to S$. Notons pour cela que si l'on prend sur $U_j$
la structure de préschéma induite par celle de $S$, l'immersion ouverte
$U_j\to S$ est quasi-compacte puisque $S$ est quasi-séparé (I, 1.2.7), donc
elle est de
\oldpage[IV-3]{36}
présentation finie (I, 1.6.2, (i)). Il suffira donc que $T'_j$ soit de
présentation finie sur $U_j$, autrement dit, on peut déjà se borner au cas où
$U_j=S$ et $T_j=T$ est fermé constructible dans $S$.

Soit $S=\operatorname{Spec}(A)$, et considérons $A$ comme limite inductive de
ses sous-$\mathbf Z$-algèbres de type fini, de sorte que
$S=\varprojlim S_\lambda$, où les $S_\lambda$ sont affines et noethériens. En
vertu de (8.3.11), il existe un $\lambda$ et une partie fermée (nécessairement
constructible) $T_\lambda$ de $S_\lambda$ tels que
$T=u_\lambda^{-1}(T_\lambda)$ ($u_\lambda:S\to S_\lambda$ étant le morphisme
canonique). On munira $T_\lambda$ d'une structure de sous-préschéma de
$S_\lambda$ et on prendra $T'=T_\lambda\times_{S_\lambda}S$~; comme
l'immersion canonique $T_\lambda\to S_\lambda$ est de présentation finie
(I, 1.6.1), il en est de même de l'immersion $T'\to S$. Comme $T'$ est affine,
on voit finalement qu'on peut se borner au cas où $T'=S$ est affine.

Alors (8.9.1), avec les mêmes notations, il existe un $\lambda$, un morphisme
de type fini $u_\lambda:X_\lambda\to S_\lambda$ et un
$\mathcal O_{X_\lambda}$-Module cohérent $\mathcal F_\lambda$ tels que $X$
soit isomorphe à $X_\lambda\times_{S_\lambda}S$, $\mathcal F$ à
$\mathcal F_\lambda\otimes_{\mathcal O_{X_\lambda}}\mathcal O_X$. On peut
alors appliquer à $S_\lambda$, $X_\lambda$ et $\mathcal F_\lambda$ le
résultat de (6.9.3) et il y a des sous-préschémas $S'_{\lambda i}$ de
$S_\lambda$, dont les ensembles sous-jacents forment une partition de
$S_\lambda$ et qui sont tels que si l'on pose
$X_{\lambda i}=X_\lambda\times_{S_\lambda}S'_{\lambda i}$,
$\mathcal F_{\lambda i}=\mathcal F_\lambda\otimes_{\mathcal O_{S_\lambda}}
\mathcal O_{S'_{\lambda i}}$ soit plat sur $S'_{\lambda i}$. Les
$S'_i=S'_{\lambda i}\times_{S_\lambda}S$ sont alors des sous-préschémas de
$S$ répondant à la question en vertu de (2.1.4).
\end{proof}

\subsection{Propriétés de permanence des morphismes par passage à la limite projective}
\label{subsection:IV.8.10-fr}

Dans cette section, nous conservons les hypothèses générales et les notations
de (8.8.1).

\begin{proposition}[8.10.1]
\label{IV.8.10.1-fr}
S'il existe $\lambda$ tel que, pour $\mu\geq\lambda$, $f_\mu$ soit un morphisme
ouvert (resp. universellement ouvert), alors $f$ est un morphisme ouvert
(resp. universellement ouvert).
\end{proposition}

\begin{proof}
Comme $X=X_\lambda\times_{Y_\lambda}Y$, l'assertion relative aux morphismes
universellement ouverts est un cas particulier de (2.4.3, (iii)). Supposons
donc $f_\mu$ ouvert pour $\mu\geq\lambda$~; il suffit de voir que pour tout
ouvert quasi-compact $U$ de $X$, $f(U)$ est ouvert dans $Y$. Or, il existe
$\mu\geq\lambda$ tel que $U=v_\mu^{-1}(U_\mu)$, où $U_\mu$ est un ouvert
quasi-compact dans $X_\mu$ (2.3.11)~; on a alors
$f(v_\mu^{-1}(U_\mu))=w_\mu^{-1}(f_\mu(U_\mu))$ (I, 3.4.8), donc $f(U)$ est
ouvert dans $Y$.
\end{proof}

\begin{corollary}[8.10.2]
\label{IV.8.10.2-fr}
Soit $f:X\to Y$ un morphisme. Pour que $f$ soit universellement ouvert, il
suffit que, pour tout entier $n>0$, si l'on pose
$Y_n=Y\otimes_{\mathbf Z}\mathbf Z[T_1,\ldots,T_n](=\mathbf V_Y^n)$, et
$X_n=X\times_YY_n$, la projection canonique
$f_n=f\times1_{Y_n}:X_n\to Y_n$ soit un morphisme ouvert.
\end{corollary}

\begin{proof}
Pour démontrer que $f$ est universellement ouvert, il suffit de prouver qu'il
en est ainsi de la restriction $f^{-1}(U)\to U$ de $f$ pour tout ouvert affine
$U$ de $Y$~; comme, par hypothèse, si
$U_n=U\otimes_{\mathbf Z}\mathbf Z[T_1,\ldots,T_n]$ est l'image réciproque de
$U$ dans $Y_n$, le morphisme $f_n^{-1}(U_n)\to U_n$, restriction de $f_n$, est
ouvert, on voit qu'on peut se borner au cas où $Y=\operatorname{Spec}(A)$ est
affine. En outre, il suffit aussi évidemment de montrer que pour tout
morphisme $Y'\to Y$, où $Y'=\operatorname{Spec}(A')$ est lui aussi affine,
$f'=f_{(Y')}$ est ouvert. Supposons d'abord que $A'$ soit une $A$-algèbre de
type fini, donc quotient d'une algèbre de polynômes $A[T_1,\ldots,T_n]$~;
alors $Y'$ est un sous-préschéma fermé de $Y_n$ et $f'$ la
\oldpage[IV-3]{37}
restriction de $f_n$ à $f_n^{-1}(Y')$~; mais pour tout ouvert $V$ de $X_n$,
on a $f_n(V\cap f_n^{-1}(Y'))=f_n(V)\cap Y'$, et comme par hypothèse $f_n(V)$
est ouvert dans $Y_n$, cela montre que $f'$ est aussi un morphisme ouvert.
Lorsque $A'$ est quelconque, il peut être considéré comme limite inductive de
ses sous-$A$-algèbres de type fini $A'_\lambda$, et le fait que $f'$ soit un
morphisme ouvert résulte de ce qui précède et de (8.10.1).
\end{proof}

\begin{proposition}[8.10.3]
\label{IV.8.10.3-fr}
Supposons qu'il existe $\alpha$ tel que~: 1\textsuperscript{o} $S_\alpha$ soit
quasi-compact~; 2\textsuperscript{o} les morphismes $X_\alpha\to S_\alpha$,
$Y_\alpha\to S_\alpha$ soient quasi-compacts et le morphisme
$Y_\alpha\to S_\alpha$ quasi-séparé~; 3\textsuperscript{o} pour
$\alpha\leq\lambda\leq\mu$, les morphismes
$u_{\lambda\mu}:S_\mu\to S_\lambda$ soient plats~; 4\textsuperscript{o}
$\overline{f_\alpha(X_\alpha)}$ soit constructible dans $Y_\alpha$. Alors,
pour que $f$ soit dominant, il faut et il suffit qu'il existe
$\lambda\geq\alpha$ tel que $f_\lambda$ soit dominant.
\end{proposition}

\begin{proof}
Les hypothèses entraînent que $Y_\alpha$ est quasi-compact et que le morphisme
$f_\alpha$ est quasi-compact (I, 1.2.4)~; par suite
$\overline{f_\alpha(X_\alpha)}=Z_\alpha$ est pro-constructible (I, 1.9.5,
(vii)) dans $Y_\alpha$. Si l'on pose
$Z_\lambda=w_{\alpha\lambda}^{-1}(Z_\alpha)$ pour $\lambda\geq\alpha$ et
$Z=w_\alpha^{-1}(Z_\alpha)$, on a donc
$Z_\lambda=\overline{f_\lambda(X_\lambda)}$ et
$Z=\overline{f(X)}$ (I, 3.4.8), et $Z_\lambda$ est pro-constructible dans
$Y_\lambda$ (I, 1.9.5, (vii)). Il suffit alors d'appliquer (8.3.13) en
remplaçant $S_\lambda$, $Z'_\lambda$ et $Z''_\lambda$ par $Y_\lambda$,
$Y_\lambda$ et $Z_\lambda$ respectivement.
\end{proof}

\begin{proposition}[8.10.4]
\label{IV.8.10.4-fr}
Supposons qu'il existe $\alpha$ tel que $Y_\alpha$ soit quasi-compact et
$f_\alpha$ de type fini et quasi-séparé. Pour que le morphisme $f$ soit
séparé, il faut et il suffit qu'il existe $\lambda\geq\alpha$ tel que
$f_\lambda$ soit séparé.
\end{proposition}

\begin{proof}
La question étant locale sur $Y_\alpha$ (puisque $Y_\alpha$ est quasi-compact
et $L$ filtrant), on peut se borner au cas où $Y_\alpha$ est affine, donc
quasi-séparé, et l'hypothèse entraîne que $X_\alpha$ (donc les $X_\lambda$ et
$X$) sont quasi-compacts et quasi-séparés. Posons
$X'_\lambda=X_\lambda\times_{Y_\lambda}X_\lambda$ pour
$\lambda\geq\alpha$ et $X'=X\times_YX$~; on a
$X'_\lambda=X'_\alpha\times_{Y_\alpha}Y_\lambda$ et
$X'=X'_\alpha\times_{Y_\alpha}Y$~; le morphisme première projection
$X'_\alpha\to X_\alpha$ est quasi-compact et quasi-séparé par hypothèse
(I, 1.2.3, (iii)), donc $X'_\lambda$ est quasi-compact et quasi-séparé.
Notons maintenant que si l'on désigne par $\Delta_\lambda$ (resp. $\Delta$)
la diagonale de $X_\lambda\times_{Y_\lambda}X_\lambda$ (resp. de
$X\times_YX$), il résulte de (I, 5.3.4 et 3.4.8) que $\Delta_\mu$ (resp.
$\Delta$) est l'image réciproque de $\Delta_\lambda$ par l'application
$v'_{\lambda\mu}:X'_\mu\to X'_\lambda$ (resp.
$v'_\lambda:X'\to X'_\lambda$). D'autre part $\Delta_\lambda$ est
constructible dans $X'_\lambda$~: en effet, puisque $f_\lambda$ est
quasi-séparé, l'immersion diagonale $X_\lambda\to X'_\lambda$ est
quasi-compacte, et localement de présentation finie puisque $f_\lambda$ est
de type fini (I, 4.3 et I, 5.4, (iii))~; il résulte alors de (1.8.4.1) que
$\Delta_\lambda$ est localement constructible, donc constructible puisque
$X'_\lambda$ est quasi-compact. On peut maintenant appliquer (8.3.12) en
remplaçant $S_\lambda$ et $Z_\lambda$ par $X'_\lambda$ et $\Delta_\lambda$
respectivement.
\end{proof}

\begin{theorem}[8.10.5]
\label{IV.8.10.5-fr}
Supposons $S_0$ quasi-compact, $X_\alpha$ et $Y_\alpha$ de présentation finie
sur $S_\alpha$, et soit $f_\alpha:X_\alpha\to Y_\alpha$ un
$S_\alpha$-morphisme. Considérons, pour un morphisme, la propriété d'être~:
\begin{enumerate}
\item[(i)] un isomorphisme~;
\item[(i \emph{bis})] un monomorphisme~;
\item[(ii)] une immersion~;
\item[(iii)] une immersion ouverte~;
\item[(iv)] une immersion fermée~;
\item[(v)] séparé~;
\item[(vi)] surjectif~;
\oldpage[IV-3]{38}
\item[(vii)] radiciel~;
\item[(viii)] affine~;
\item[(ix)] quasi-affine~;
\item[(x)] fini~;
\item[(xi)] quasi-fini~;
\item[(xii)] propre.
\end{enumerate}
Alors, si $\mathbf P$ désigne une des propriétés précédentes, pour que $f$
possède la propriété $\mathbf P$, il faut et il suffit qu'il existe
$\lambda\geq\alpha$ tel que $f_\lambda$ possède la propriété $\mathbf P$
(auquel cas $f_\mu$ la possède aussi pour $\mu\geq\lambda$).

Si $S_0$ est en outre supposé quasi-séparé, la même conclusion est valable
lorsque $\mathbf P$ est la propriété d'être~:
\begin{enumerate}
\item[(xiii)] projectif~;
\item[(xiv)] quasi-projectif.
\end{enumerate}
\end{theorem}

\begin{proof}
Le cas où $\mathbf P$ est l'une des propriétés (i) ou (v) n'a été inséré dans
l'énoncé que pour mémoire~; dans ces cas, le théorème découle de ce qui a été
démontré respectivement dans (8.8.2.4) et (8.10.4). D'ailleurs, compte tenu de
(I, 5.4.1 et 5.3.4), (v) résulte aussi de (iv). Le cas (i \emph{bis}) se
déduit aussitôt de (i), en utilisant (I, 5.3.8) et en notant (comme on l'a
déjà utilisé dans (8.10.4)) que la diagonale $\Delta_f$ se déduit de
$\Delta_{f_\lambda}$ par le changement de base $S\to S_\lambda$.

On notera d'autre part que (vi), (vii) et (xi) sont en fait des conditions
sur les fibres $f^{-1}(y)$ des morphismes envisagés, compte tenu de la
transitivité des fibres par changement de base (I, 3.6.4)~: la condition (vi)
signifie en effet que toutes les fibres doivent être non vides, la condition
(vii) qu'elles doivent être radicielles (I, 3.5.8) et la condition (xi)
qu'elles doivent être finies (III, 6.2.2 et 6.2.3 et I, 6.4.4, compte tenu de
ce que $f$ et les $f_\lambda$ sont des morphismes de type fini par (I, 5.4,
(v)). Le théorème dans ces trois cas sera donc encore conséquence d'un
résultat général sur ce type de propriétés ne concernant que les fibres, qui
sera établi dans (9.3.3)~; nous repoussons donc jusqu'à ce moment-là la
démonstration du théorème dans le cas (xi) (bien entendu, le lecteur pourra
vérifier que, sauf aux n\textsuperscript{os}~8.11 et 8.12, nous ne ferons pas
usage du théorème dans ce cas jusqu'à (9.3.3) et que (8.11) et (8.12) ne
seront pas utilisés avant (9.3.3)).

Pour les cas qui restent à démontrer, on peut se borner à prouver que la
condition de l'énoncé est nécessaire, toutes les propriétés $\mathbf P$
envisagées étant invariantes par changement de base (voir chap. I\textsuperscript{er}
et II dans les numéros concernant chacune de ces propriétés). On peut en
outre supposer que $S_\alpha=S_0$ et que $Y_\alpha=S_\alpha$, donc
$Y_\lambda=S_\lambda$ pour tout $\lambda\geq\alpha$. Enfin, les propriétés
(i) à (xii) sont locales sur $S_0$, donc, comme $S_0$ est réunion finie
d'ouverts affines et $L$ filtrant, on peut se borner pour les démontrer au cas
où $S_0=\operatorname{Spec}(A_0)$ est affine (donc quasi-séparé). On désignera
par $A_\lambda$ (resp. $A$) l'anneau de $S_\lambda$ (resp. $S$).

\emph{Cas (ii), (iii), (iv)}~: Supposons que $f$ soit une immersion (resp.
une immersion ouverte, resp. une immersion fermée), et soit $X'$ le
sous-préschéma (resp. induit sur un ouvert, resp. fermé) de $S$ associé à $f$,
qui est donc un $S$-préschéma de présentation finie.
\oldpage[IV-3]{39}
En vertu de (8.6.3), il existe donc un $\lambda\geq\alpha$ et un
sous-préschéma (resp. induit sur un ouvert, resp. fermé) $X'_\lambda$ de
$S_\lambda$, de présentation finie sur $S_\lambda$, tels que $X'$ soit
isomorphe à $X'_\lambda\times_{S_\lambda}S$. Pour tout $\mu\geq\lambda$,
$X'_\mu=X'_\lambda\times_{S_\lambda}S_\mu$ est donc un sous-préschéma (resp.
induit sur un ouvert, resp. fermé) de $S_\mu$, de présentation finie sur
$S_\mu$, et il résulte donc de (8.8.2.4) et (8.8.2.5) qu'il existe un
$\mu\geq\lambda$ et un isomorphisme $g_\mu:X_\mu\to X'_\mu$ tel que
$g=g_\mu\times1_S$ soit l'isomorphisme $X\to X'$ associé à $f$~; d'où la
conclusion.

\emph{Cas (vi) et (vii)}~: On sait (1.8.4.1) que
$Z_\alpha=f_\alpha(X_\alpha)$ est constructible dans $S_\alpha$~; si l'on
pose $Z_\lambda=u_{\alpha\lambda}^{-1}(Z_\alpha)$ pour $\lambda\geq\alpha$
et $Z=u_\alpha^{-1}(Z_\alpha)$, on a $Z_\lambda=f_\lambda(X_\lambda)$ et
$Z=f(X)$ (I, 3.4.8). Comme en vertu de (8.3.11) l'application canonique
$\varinjlim\mathfrak C(S_\lambda)\to\mathfrak C(S)$ est injective, la
relation $f(X)=S$ implique l'existence d'un $\lambda\geq\alpha$ tel que
$f_\lambda(X_\lambda)=S_\lambda$, ce qui prouve le théorème dans le cas (vi).
Pour le prouver dans le cas (vii), il suffit de noter que le morphisme
structural $X_\alpha\times_{S_\alpha}X_\alpha\to S_\alpha$ est de présentation
finie puisqu'il en est ainsi de la première projection
$X_\alpha\times_{S_\alpha}X_\alpha\to X_\alpha$ (1.6.2)~; il suffit donc en
vertu de (1.8.7.1) d'appliquer le cas (vi) du théorème au morphisme diagonal
$\Delta_{f_\alpha}:X_\alpha\to X_\alpha\times_{S_\alpha}X_\alpha$ en notant
que l'on a $\Delta_{f_\lambda}=\Delta_{f_\alpha}\times1_{S_\lambda}$ et
$\Delta_f=\Delta_{f_\alpha}\times1_S$ (I, 5.3.4 et 3.3.11).

\emph{Cas (viii) et (ix)}~: Comme $S=\operatorname{Spec}(A)$ est affine, dire
que $f:X\to S$ est affine (resp. quasi-affine) signifie qu'il existe un entier
$r$ et une immersion fermée (resp. une immersion)
$j:X\to\mathbf V_S^r=\operatorname{Spec}(A[T_1,\ldots,T_r])$ de
$S$-préschémas, puisque $f$ est de type fini et $S$ quasi-compact
(II, 5.1.9). Comme
$\mathbf V_S^r=\mathbf V_{S_0}^r\times_{S_0}S$, et que
$\mathbf V_{S_0}^r$ est un $S_0$-préschéma de présentation finie, il résulte
de (8.8.2, (i)) appliqué au $S$-morphisme $j$, qu'il existe un $\lambda$ et un
$S_\lambda$-morphisme $j_\lambda:X_\lambda\to\mathbf V_{S_\lambda}^r$ tel que
$j=j_\lambda\times1_S$~; appliquant alors (iv) (resp. (ii)) à $j$, on en
déduit qu'il existe $\mu\geq\lambda$ tel que $j_\mu$ soit une immersion
fermée (resp. une immersion)~; par suite $f_\mu$ est affine (resp.
quasi-affine).

\emph{Cas (x)}~: Par hypothèse, on a $X=\operatorname{Spec}(B)$, où $B$ est
une $A$-algèbre qui est un $A$-module de présentation finie (1.4.7)~; il
résulte donc de (8.5.2, (ii)) qu'il y a un $\lambda$ et un $A_\lambda$-module
de présentation finie $B_\lambda$ tels que le $A$-module $B$ soit isomorphe à
$B_\lambda\otimes_{A_\lambda}A$. La structure de $A$-algèbre de $B$ est
définie par un $A$-homomorphisme $m:B\otimes_AB\to B$~; comme on a
$B\otimes_AB=(B_\lambda\otimes_{A_\lambda}B_\lambda)\otimes_AA$, il existe
d'après (8.5.2, (i)) un $\mu\geq\lambda$ et un $A_\mu$-homomorphisme
$m_\mu:B_\mu\otimes_{A_\mu}B_\mu\to B_\mu$ tel que $m=m_\mu\otimes1$. En
considérant les diagrammes usuels exprimant l'associativité et la
commutativité de $m$, on voit en appliquant encore (8.5.2, (i)) qu'il existe
$\nu\geq\mu$ tel que $m_\nu$ définisse sur $B_\nu$ une multiplication
associative et commutative~; de la même manière on voit qu'on peut supposer
$\nu$ pris assez grand pour que l'anneau $B_\nu$ ainsi défini admette un
élément unité. Si $X'_\nu=\operatorname{Spec}(B_\nu)$, il est clair alors que
$X$ est $S$-isomorphe à $X'_\nu\times_{S_\nu}S$, donc, en vertu de (i), il
existe $\rho\geq\nu$ tel que $X_\rho$ et $X'_\rho$ soient
$S_\rho$-isomorphes, ce qui achève la démonstration dans ce cas.

Pour démontrer le théorème dans le cas (xii), nous démontrerons d'abord la
proposition suivante~:
\end{proof}

\begin{proposition}[8.10.5.1]
\label{IV.8.10.5.1-fr}
(lemme de Chow pour les morphismes de présentation finie). --- Soient $A$ un
anneau, $X$, $Y$ deux $A$-préschémas de présentation finie,
$f:X\to Y$
\oldpage[IV-3]{40}
un $A$-morphisme séparé. Alors il existe deux $A$-préschémas $X'$, $P$ de
présentation finie, et des $A$-morphismes $p:P\to Y$, $j:X'\to P$,
$g:X'\to X$, tels que le diagramme
\[
\begin{tikzcd}[column sep=large,row sep=large]
X' \arrow[r,"j"] \arrow[d,"g"'] & P \arrow[d,"p"] \\
X \arrow[r,"f"'] & Y
\end{tikzcd}
\]
soit commutatif, et que~: 1\textsuperscript{o} $p$ soit projectif~;
2\textsuperscript{o} $g$ soit projectif et surjectif~;
3\textsuperscript{o} $j$ soit une immersion ouverte.
\end{proposition}

\begin{proof}
En effet, soient $A_0\subset A$, $X_0$, $Y_0$ et $f_0$ déterminés comme dans
(8.9.1) de sorte que $Y_0$ est noethérien et $f_0$ de type fini~; on peut en
outre supposer $f_0$ séparé d'après (8.10.4). Le lemme de Chow (II, 5.6.1)
montre alors l'existence de trois morphismes $p_0:P_0\to Y_0$,
$g_0:X'_0\to X_0$ et $j_0:X'_0\to P_0$, de type fini, tels que le diagramme
\[
\begin{tikzcd}[column sep=large,row sep=large]
X'_0 \arrow[r,"j_0"] \arrow[d,"g_0"'] & P_0 \arrow[d,"p_0"] \\
X_0 \arrow[r,"f_0"'] & Y_0
\end{tikzcd}
\]
soit commutatif, et que $p_0$ soit projectif, $g_0$ projectif et surjectif et
$j_0$ une immersion ouverte. Les propriétés de l'énoncé résultent alors de
l'invariance des propriétés précédentes par changement de base (II, 5.5.5,
(iii) et I, 3.5.2 et 4.3.2).
\end{proof}

\begin{proof}[Suite de la démonstration de (8.10.5)]
\emph{Cas (xii)}~: Appliquons au morphisme $f_0:X_0\to S_0$ la prop.
(8.10.5.1)~: on a donc un diagramme commutatif
\[
\begin{tikzcd}[column sep=large,row sep=large]
X'_0 \arrow[r,"j_0"] \arrow[d,"g_0"'] & P_0 \arrow[d,"p_0"] \\
X_0 \arrow[r,"f_0"'] & S_0
\end{tikzcd}
\]
où $p_0$ est projectif, $g_0$ projectif et surjectif, et $j_0$ une immersion
ouverte~; on en déduit pour chaque $\lambda$ un diagramme analogue où les
morphismes $p_\lambda=p_0\times1_{S_\lambda}$,
$g_\lambda=g_0\times1_{S_\lambda}$ et $j_\lambda=j_0\times1_{S_\lambda}$ ont
respectivement les mêmes propriétés, et de même pour les morphismes limites
projectives $p=p_0\times1_S$, $g=g_0\times1_S$, $j=j_0\times1_S$. Comme $g$
est propre (II, 5.5.3), il en est de même de $p\circ j=f\circ g$ (II, 5.4.2)
et comme $p$ est séparé, $j$ est propre, donc une immersion fermée~;
appliquant le cas (iv) au morphisme $j$ (en notant que $X'_0$ et $P_0$ sont
des $S_0$-préschémas de présentation finie (8.10.5.1 et 1.6.2)), on voit
qu'il existe $\lambda$ tel que $j_\lambda$ soit une immersion fermée, donc
soit propre (II, 5.4.2). Mais alors
$f_\lambda\circ g_\lambda=j_\lambda\circ p_\lambda$ est propre (II, 5.5.3 et
5.4.2) et comme $g_\lambda$ est surjectif, et qu'on peut supposer $f_\lambda$
séparé en vertu de l'hypothèse sur $f$ et de (v), il résulte de (II, 5.4.3)
que $f_\lambda$ est propre.

\emph{Cas (xiii) et (xiv)}~: En vertu de (xii) et de (II, 5.5.3) (qui est
applicable puisque les $S_\lambda$ sont quasi-compacts et quasi-séparés,
compte tenu de (1.7.19)), il suffit de
\oldpage[IV-3]{41}
considérer le cas où $f$ est quasi-projectif. Supposons donc qu'il existe un
$\mathcal O_X$-Module inversible $\mathcal L$ qui soit $f$-ample~; comme $S_0$
est quasi-compact et quasi-séparé il en est de même de $X_0$ (1.2.3), et il y
a donc un $\lambda$ et un $\mathcal O_{X_\lambda}$-Module quasi-cohérent
$\mathcal L_\lambda$ de présentation finie tel que
$\mathcal L=v_\lambda^*(\mathcal L_\lambda)$ (8.5.2, (ii))~; en outre, en
vertu de (8.5.5) on peut supposer $\mathcal L_\lambda$ inversible. Le théorème
dans ce cas est alors conséquence du lemme plus précis suivant.

\begin{lemma}[8.10.5.2]
\label{IV.8.10.5.2-fr}
Supposons $S_0$ quasi-compact, et soit $\mathcal L_\lambda$ un
$\mathcal O_{X_\lambda}$-Module inversible. Pour que $\mathcal L$ soit un
$\mathcal O_X$-Module ample pour $f$ (resp. très ample pour $f$), il faut et
il suffit qu'il existe $\mu\geq\lambda$ tel que $\mathcal L_\mu$ soit ample
pour $f_\mu$ (resp. très ample pour $f_\mu$).
\end{lemma}

La condition étant évidemment suffisante (II, 4.4.10 et 4.6.13), montrons
qu'elle est nécessaire~; les $S_\lambda$ étant quasi-compacts et les
$f_\lambda$ de type fini, on peut, en remplaçant $\mathcal L$ par une puissance
tensorielle convenable, se borner à considérer le cas où $\mathcal L$ est très
ample (II, 4.6.11). En outre, la question étant ici locale sur $S_0$ (vu
(II, 4.4.5) et le fait que $L$ est filtrant), on peut se borner au cas où
$S_0$ (et par suite $S$) est affine. Alors, en vertu de (II, 4.4.1, (ii) et
II, 4.1.2), il existe une $S$-immersion
$j:X\to\mathbf P_S^r=P$ telle que $\mathcal L$ soit isomorphe à
$j^*(\mathcal O_P(1))$. Compte tenu de (8.8.2, (i)), de (ii) et de
(II, 4.1.3), il existe donc un $\mu\geq\lambda$ et une immersion
$j_\mu:X_\mu\to\mathbf P_{S_\mu}^r=P_\mu$ tels que
$j=j_\mu\times1_S$~; utilisant ensuite (II, 4.1.3.2) et (8.5.2.5), on voit
qu'il existe $\nu\geq\mu$ tel que $\mathcal L_\nu$ soit isomorphe à
$j_\nu^*(\mathcal O_{P_\nu}(1))$, ce qui montre que $\mathcal L_\nu$ est très
ample pour $f_\nu$ (II, 4.4.2).
\end{proof}

\subsection{Application aux morphismes quasi-finis}
\label{subsection:IV.8.11-fr}
\label{IV.8.11-fr}

Nous nous proposons dans cette section de démontrer les deux théorèmes
suivants~:

\begin{theorem}[8.11.1]
\label{IV.8.11.1-fr}
Soit $f:X\to Y$ un morphisme propre, localement de présentation finie, et
quasi-fini. Alors le morphisme $f$ est fini.
\end{theorem}

\begin{theorem}[8.11.2]
\label{IV.8.11.2-fr}
Soit $f:X\to Y$ un morphisme localement de présentation finie, quasi-fini et
séparé. Alors le morphisme $f$ est quasi-affine, et \emph{a fortiori}
quasi-projectif.
\end{theorem}

\begin{remark}[8.11.3]
\label{IV.8.11.3-fr}
On va voir ci-dessous que l'on peut se ramener, pour la démonstration de
(8.11.1) et (8.11.2), au cas où $Y$ est localement noethérien~; on notera que
dans ce cas on obtiendra ainsi une autre démonstration du théorème de
Chevalley (III, 4.4.2).
\end{remark}

\begin{env}[8.11.4]
\label{IV.8.11.4-fr}
Les hypothèses et les conclusions de (8.11.1) et (8.11.2) sont toutes locales
sur $Y$ (I, 6.1, I, 2.6, (II, 5.1.1), (II, 5.4.1) et (II, 6.2.2)), donc on
peut supposer $Y=\operatorname{Spec}(A)$ affine. On sait qu'il existe alors un
sous-anneau $A_0$ de $A$, qui est une $\mathbf Z$-algèbre de type fini, et un
morphisme de type fini $f_0:X_0\to\operatorname{Spec}(A_0)$ tel que $X$
s'identifie à $X_0\otimes_{A_0}A$ et $f$ à $f_0\times1$ (8.9.1). En outre,
$A$ peut être considéré comme limite inductive de ses sous-anneaux contenant
$A_0$ et qui sont des $\mathbf Z$-algèbres de type fini~; utilisant la méthode
de (8.1.2, c)) ainsi que (8.10.5, (v), (xi) et (xii)), on voit qu'il suffit de
démontrer les théorèmes (8.11.1) et (8.11.2) pour $f_0$. Supposons donc
désormais $Y$ noethérien~; utilisant maintenant la méthode de (8.1.2, a))
ainsi que (8.10.5, (v), (ix), (x), (xi) et (xii)) on peut remplacer $Y$ par
$\operatorname{Spec}(\mathcal O_y)$, où $y$ est un
\oldpage[IV-3]{42}
point quelconque de $Y$, donc on voit finalement que l'on peut supposer que
$Y=\operatorname{Spec}(A)$, où $A$ est un \emph{anneau local noethérien}.
Soient $\mathfrak m$ l'idéal maximal de $A$, $\widehat A$ le complété de $A$
pour la topologie $\mathfrak m$-préadique~; on sait que $\widehat A$ est un
anneau local noethérien et que le morphisme
$\operatorname{Spec}(\widehat A)\to\operatorname{Spec}(A)$ est fidèlement
plat et quasi-compact (0\textsubscript{I}, 7.3.5)~; quant à (2.7.1, (i),
(vii), (xiv), (xv) et (xvi)), on voit en outre qu'on peut se borner au cas où
$A$ est \emph{complet}. Il résulte alors de (II, 6.2.6) que
$X=X'\amalg X''$, où $X'$ est un $Y$-schéma fini et $X''$ un $Y$-schéma
quasi-fini tel que $X''\cap f^{-1}(y)=\varnothing$.
\end{env}

Plaçons-nous d'abord dans les hypothèses de (8.11.1)~; comme $f$ est propre,
$X''$, qui est fermé dans $X$, est propre sur $Y$ (II, 5.4.10), donc
$f(X'')$ est fermé dans $Y$~; mais $y$ n'est pas contenu dans $f(X'')$, et
par ailleurs est adhérent à tout point de $Y$, donc $f(X'')=\varnothing$, et
par suite $X''$ est vide et $f$ est \emph{fini}.

Plaçons-nous maintenant dans les hypothèses de (8.11.2) et, en nous ramenant
encore (comme on peut le faire d'après ce qui précède) au cas où
$Y=\operatorname{Spec}(A)$ est affine et noethérien de dimension finie,
raisonnons en outre par récurrence sur la dimension de $Y$. En se ramenant
comme ci-dessus au cas où $A$ est en outre local et complet, on a
$\dim(\mathcal O_y)=\dim(A)=\dim(Y)$ et pour tout $z\neq y$,
$\dim(\mathcal O_z)<\dim(\mathcal O_y)$, donc
$\dim(Y-\{y\})<\dim(Y)$. Or, par hypothèse on a
$f(X'')\subset Y-\{y\}$ et la restriction de $f$ à $X''$ est évidemment un
morphisme quasi-fini et séparé~; appliquant à $Y-\{y\}$ et à $X''$
l'hypothèse de récurrence, on voit que $X''$ est quasi-affine sur
$Y-\{y\}$~; mais $Y-\{y\}$ étant quasi-affine sur $Y$ puisque $Y$ est
noethérien, $X''$ est aussi quasi-affine sur $Y$ (II, 5.1.10, (ii))~; comme
par ailleurs $X'$ est fini (et \emph{a fortiori} affine) sur $Y$, $X$ est
quasi-affine sur $Y$ (II, 4.6.17 et 5.1.2, c')).

\begin{proposition}[8.11.5]
\label{IV.8.11.5-fr}
Soit $f:X\to Y$ un morphisme de présentation finie. Les conditions suivantes
sont équivalentes~:
\begin{enumerate}
\item[\emph{a)}] $f$ est une immersion fermée.
\item[\emph{b)}] $f$ est un monomorphisme propre.
\item[\emph{c)}] $f$ est propre et pour tout $y\in Y$, $f^{-1}(y)$ est
radiciel et géométriquement réduit sur $k(y)$ (c'est-à-dire vide ou
$k(y)$-isomorphe à $\operatorname{Spec}(k(y))$).
\end{enumerate}
\end{proposition}

Il est clair que \emph{a)} entraîne \emph{b)}. Pour voir que \emph{b)}
entraîne \emph{c)}, notons (I, 5.3.8) que pour tout $y\in Y$, le morphisme
$f^{-1}(y)\to\operatorname{Spec}(k(y))$ déduit de $f$ par extension de la base
est un monomorphisme, donc est injectif, et par suite $f^{-1}(y)$ est vide ou
réduit à un point, et en tout cas affine. De plus, si $A$ est l'anneau de
$f^{-1}(y)$, l'homomorphisme canonique $A\otimes_kA\to A$ est bijectif
(I, 5.3.8). Cela entraîne évidemment que $A=k$, sans quoi il y aurait un
élément $a\in A$ non dans $k$ et les images de $a\otimes1$ et $1\otimes a$
dans $A$ seraient toutes deux égales à $a$, alors que $a\otimes1\neq1\otimes a$
puisque $1$ et $a$ forment un système libre sur $k$.

Reste à prouver que \emph{c)} entraîne \emph{a)}. Il résulte tout d'abord de
(8.11.1) que $f$ est un morphisme \emph{fini}, donc
$X=\operatorname{Spec}(\mathcal A)$, où $\mathcal A$ est une
$\mathcal O_Y$-Algèbre \emph{finie}. Il suffit donc de prouver que
l'homomorphisme canonique $\mathcal O_Y\to\mathcal A$ est \emph{surjectif}
(II, 1.4.10), ou encore que pour tout $y\in Y$, l'homomorphisme
$\mathcal O_y\to\mathcal A_y$ est surjectif. Mais par hypothèse
\oldpage[IV-3]{43}
$f^{-1}(y)=\operatorname{Spec}(\mathcal A_y/\mathfrak m_y\mathcal A_y)$
(II, 1.5.5) est tel que l'homomorphisme correspondant
$k(y)=\mathcal O_y/\mathfrak m_y\to
\mathcal A_y/\mathfrak m_y\mathcal A_y$ soit bijectif~; comme $\mathcal A_y$
est un $\mathcal O_y$-module de type fini, le lemme de Nakayama montre que
$\mathcal O_y\to\mathcal A_y$ est surjectif, ce qui achève la démonstration.

\begin{remark}[8.11.5.1]
\label{IV.8.11.5.1-fr}
On notera que le raisonnement précédent prouve que si $f$ est un
\emph{monomorphisme}, alors, pour tout $y\in Y$, $f^{-1}(y)$ est vide ou
$k(y)$-isomorphe à $\operatorname{Spec}(k(y))$.
\end{remark}

\begin{proposition}[8.11.6]
\label{IV.8.11.6-fr}
Si un morphisme $f:X\to Y$ de présentation finie est un homéomorphisme
universel, il est fini, surjectif et radiciel (la réciproque étant vraie
d'après (2.4.5, (i))).
\end{proposition}

\begin{proof}
En effet, $f$ étant de type fini, universellement fermé, et séparé en vertu de
(2.4.4), est propre par définition (II, 5.4.1). Comme il est évidemment
quasi-fini (II, 6.2.3), il est fini d'après (8.11.1). On sait par ailleurs
qu'il est radiciel (2.4.4), et évidemment surjectif.
\end{proof}

\subsection{Nouvelle démonstration et généralisation du «~Main Theorem~» de Zariski}
\label{subsection:IV.8.12-fr}
\label{IV.8.12-fr}

\begin{lemma}[8.12.1]
\label{IV.8.12.1-fr}
Soient $f:X\to Y$ un morphisme quasi-compact et quasi-séparé, $\mathcal C$ une
$\mathcal O_Y$-Algèbre quasi-cohérente, $Z=\operatorname{Spec}(\mathcal C)$,
qui est un $Y$-préschéma affine au-dessus de $Y$. Soient $g:X\to Z$ un
$Y$-morphisme, $\varphi=\mathcal A(g):\mathcal C\to f_*(\mathcal O_X)$ le
$\mathcal O_Y$-homomorphisme correspondant de $\mathcal O_Y$-Algèbres
(II, 1.2.7). Supposons que $g$ soit une immersion. Alors, pour que l'image
fermée de $X$ par $g$ (I, 9.5.3) soit égale à $Z$, il faut et il suffit que
$\varphi$ soit injectif~; $g$ est alors une immersion ouverte.
\end{lemma}

\begin{proof}
L'hypothèse entraîne que $f_*(\mathcal O_X)$ est une $\mathcal O_Y$-Algèbre
quasi-cohérente (1.7.5)~; en outre, comme le morphisme canonique $h:Z\to Y$
est affine, donc quasi-compact et quasi-séparé (1.2.2 et 1.1.2), $g$ est un
morphisme quasi-compact et quasi-séparé, donc $g_*(\mathcal O_X)$ est une
$\mathcal O_Z$-Algèbre quasi-cohérente (1.7.5). Cela étant, dire que l'image
fermée de $X$ par $g$ est égale à $Z$ signifie (I, 9.5.2) que
l'homomorphisme canonique $\theta:\mathcal O_Z\to g_*(\mathcal O_X)$ est
injectif. Mais on a $h_*(\mathcal O_Z)=\mathcal C$ par définition de $Z$
(II, 1.3.1), et $h_*(g_*(\mathcal O_X))=f_*(\mathcal O_X)$. Comme $Z$ est
affine au-dessus de $Y$, il revient au même de dire que l'homomorphisme
$\theta:\mathcal O_Z\to g_*(\mathcal O_X)$ est injectif ou que
l'homomorphisme correspondant
$\varphi=h_*(\theta):\mathcal C\to f_*(\mathcal O_X)$ est injectif
(I, 1.3.9). Le fait que $g$ est alors une immersion ouverte résulte de
(I, 9.5.10) et de l'hypothèse que $g$ est une immersion.
\end{proof}

\begin{lemma}[8.12.2]
\label{IV.8.12.2-fr}
Soient $Y$ un préschéma quasi-compact et quasi-séparé, $f:X\to Y$ un
morphisme quasi-séparé de type fini, $\mathcal C$ une $\mathcal O_Y$-Algèbre
quasi-cohérente, $Z=\operatorname{Spec}(\mathcal C)$. Soient $g:X\to Z$ un
$Y$-morphisme, $\varphi:\mathcal C\to f_*(\mathcal O_X)$ le
$\mathcal O_Y$-homomorphisme correspondant de $\mathcal O_Y$-Algèbres. Soit
$(\mathcal C_\lambda)$ la famille filtrante croissante des sous-$\mathcal O_Y$-Algèbres
quasi-cohérentes de type fini de $\mathcal C$ (dont $\mathcal C$ est la
réunion ((I, 9.6.6) et (1.7.9)))~; posons
$Z_\lambda=\operatorname{Spec}(\mathcal C_\lambda)$ et soit $g_\lambda$ le
morphisme composé $X\xrightarrow{g}Z\to Z_\lambda$. Alors les conditions
suivantes sont équivalentes~:
\begin{enumerate}
\item[\emph{a)}] $g$ est une immersion.
\item[\emph{b)}] Il existe $\lambda$ tel que $g_\lambda$ soit une immersion.
\end{enumerate}
\oldpage[IV-3]{44}
De plus, lorsque $g_\lambda$ est une immersion, il en est de même de $g_\mu$
pour $\mu\geq\lambda$.

Il suffit d'appliquer (II, 3.8.4) en remplaçant $\mathcal L$ par
$\mathcal O_X$ et $\mathcal S$ par $\mathcal C[\mathrm T]$, et en tenant
compte de (II, 3.1.7).
\end{lemma}

\begin{proposition}[8.12.3]
\label{IV.8.12.3-fr}
Soient $Y$ un préschéma quasi-compact et quasi-séparé, $f:X\to Y$ un
morphisme séparé de type fini. Soit $\mathcal B=f_*(\mathcal O_X)$, qui est
une $\mathcal O_Y$-Algèbre quasi-cohérente (I, 9.2.2)~; soit $\mathcal C$ la
fermeture intégrale de $\mathcal O_Y$ dans $\mathcal B$, qui est une
$\mathcal O_Y$-Algèbre quasi-cohérente (II, 6.3.4)~; posons
$Z=\operatorname{Spec}(\mathcal C)$, et soit $g:X\to Z$ le $Y$-morphisme
correspondant à l'injection
$\varphi:\mathcal C\to\mathcal B=f_*(\mathcal O_X)$ (II, 1.2.7). Soit
$(\mathcal C_\lambda)$ la famille filtrante croissante des
sous-$\mathcal O_Y$-Algèbres quasi-cohérentes de type fini de $\mathcal C$
(dont $\mathcal C$ est la réunion ((I, 9.6.6) et (1.7.9))), et, pour tout
$\lambda$, posons $Z_\lambda=\operatorname{Spec}(\mathcal C_\lambda)$, et
soit $g_\lambda:X\to Z_\lambda$ le $Y$-morphisme correspondant à l'injection
$\mathcal C_\lambda\to\mathcal B=f_*(\mathcal O_X)$~; les conditions
suivantes sont équivalentes~:
\begin{enumerate}
\item[\emph{a)}] Il existe une factorisation de $f$ en
\[
X\xrightarrow{f'}Y'\xrightarrow{u}Y,
\]
où $f'$ est une immersion et $u$ un morphisme fini.
\item[\emph{a$'$)}] Il existe une factorisation
$X\xrightarrow{f'}Y'\xrightarrow{u}Y$ de $f$, où $f'$ est une immersion
ouverte et $u$ un morphisme fini.
\item[\emph{b)}] Le morphisme $g:X\to Z$ est une immersion.
\item[\emph{c)}] Il existe $\lambda$ tel que
$g_\lambda:X\to Z_\lambda$ soit une immersion.
\end{enumerate}
De plus, lorsqu'il en est ainsi, $g$ est une immersion ouverte, $g(X)$ est
dense dans $Z$, et il existe $\lambda$ tel que, pour $\mu\geq\lambda$,
$g_\mu$ soit une immersion ouverte.
\end{proposition}

Comme l'homomorphisme $\varphi:\mathcal C\to f_*(\mathcal O_X)$ est
injectif, il résulte de (8.12.1) que si $g$ est une immersion, c'est une
immersion ouverte et $g(X)$ est dense dans $Z$, et il en est de même pour
$g_\lambda$. Le fait que \emph{a)} entraîne \emph{a$'$)} résulte aussi de
(8.12.1), appliqué en remplaçant $Z$ par $Y'$ et $g$ par $f'$ ($Y'$ étant
fini, donc affine au-dessus de $Y$)~: en effet si $Y''$ est l'image fermée de
$X$ par $f'$, $Y''$ est fini au-dessus de $Y$ et $f'$ se factorise en
$X\xrightarrow{f''}Y''\xrightarrow{j}Y'$, où $j$ est l'injection canonique,
et $f''$ est une immersion (I, 4.1.10)~; il résulte donc de (8.12.1) que
$f''$ est une immersion ouverte.

L'équivalence de \emph{b)} et \emph{c)} résulte de (8.12.2) ainsi que le fait
que $g_\lambda$ est alors une immersion pour $\lambda$ assez grand. Il est
clair que \emph{c)} implique \emph{a)}, puisque $Z_\lambda$ est fini
au-dessus de $Y$ (II, 6.3.4 et 6.1.2). Montrons enfin que \emph{a)} implique
\emph{c)}. On a vu ci-dessus qu'on peut supposer que $Y'$ est l'image fermée
de $X$ par $f'$, et alors il résulte de (8.12.1) que si l'on pose
$\mathcal B'=u_*(\mathcal O_{Y'})$, de sorte que $Y'$ s'identifie à
$\operatorname{Spec}(\mathcal B')$, l'homomorphisme
$\varphi':\mathcal B'\to\mathcal B=f_*(\mathcal O_X)$ est \emph{injectif}.
Mais comme par hypothèse $\mathcal B'$ est une $\mathcal O_Y$-Algèbre finie,
elle s'identifie par définition de $\mathcal B$ à une des
sous-$\mathcal O_Y$-Algèbres $\mathcal C_\lambda$, ce qui démontre
\emph{c)}.

Nous dirons qu'un morphisme $f:X\to Y$, où $Y$ est quasi-compact et
quasi-séparé, est \emph{pseudo-fini}, s'il est de type fini et vérifie la
condition \emph{a)} de (8.12.3) (auquel cas il est nécessairement séparé).

\begin{corollary}[8.12.4]
\label{IV.8.12.4-fr}
Soient $Y$ un préschéma quasi-compact et quasi-séparé, $f:X\to Y$ un
morphisme.
\oldpage[IV-3]{45}
\begin{enumerate}
\item[(i)] Supposons $f$ pseudo-fini. Alors, pour tout morphisme $Y'\to Y$,
où $Y'$ est quasi-compact et quasi-séparé,
$f_{(Y')}:X'=X_{(Y')}\to Y'$ est pseudo-fini.
\item[(ii)] Soit $(U_\lambda)$ un recouvrement de $Y$ formé d'ouverts
quasi-compacts. Pour que $f$ soit pseudo-fini, il faut et il suffit que pour
tout $\lambda$, la restriction
$f_\lambda:f^{-1}(U_\lambda)\to U_\lambda$ de $f$ soit un morphisme
pseudo-fini.
\item[(iii)] Supposons en outre $Y$ noethérien, et $f$ de type fini. Alors,
pour que $f$ soit pseudo-fini, il faut et il suffit que, pour tout $y\in Y$,
le morphisme
$f_y=f\times 1:X_y=X\times_Y\operatorname{Spec}(\mathcal O_y)
\to\operatorname{Spec}(\mathcal O_y)$ le soit.
\end{enumerate}
\end{corollary}

\begin{proof}
(i) Il est clair que $f_{(Y')}$ est de type fini (1.5.4)~; en outre, une
factorisation
$X\xrightarrow{g}Z\xrightarrow{u}Y$, où $g$ est une immersion et $u$ est
fini, donne une factorisation
$X'\xrightarrow{g'}Z'\xrightarrow{u'}Y'$ de $f_{(Y')}$, où
$Z'=Z_{(Y')}$, $g'=g_{(Y')}$ et $u'=u_{(Y')}$~; $g'$ est une immersion
(I, 4.3.2) et $u'$ est fini (II, 6.1.5)~; donc $f_{(Y')}$ est pseudo-fini.

(ii) La condition est nécessaire en vertu de (i), les $U_\lambda$ étant
quasi-séparés puisque $Y$ l'est. Pour voir qu'elle est suffisante, observons
(avec les notations de (8.12.3)) que si l'on pose
$X_\lambda=f^{-1}(U_\lambda)$, on a
$\mathcal B|U_\lambda=(f_\lambda)_*(\mathcal O_{X_\lambda})$,
$\mathcal C|U_\lambda$ est la fermeture intégrale de $\mathcal O_{U_\lambda}$
dans $\mathcal B|U_\lambda$ et par suite, si $h:Z\to Y$ est le morphisme
canonique, $Z_\lambda=\operatorname{Spec}(\mathcal C|U_\lambda)$ s'identifie à
$h^{-1}(U_\lambda)$. Or, pour que $g:X\to Z$ soit une immersion, il faut et
il suffit que pour tout $\lambda$, la restriction
$g_\lambda:f^{-1}(U_\lambda)\to h^{-1}(U_\lambda)$ de $g$ le soit
(I, 4.2.4). Cela entraîne la conclusion en vertu de (8.12.3).

(iii) Il suffit, en vertu de (ii), de prouver que $y$ admet un voisinage $U$
tel que la restriction $f^{-1}(U)\to U$ de $f$ soit un morphisme pseudo-fini.
Désignons par $(U_\lambda)$ le système projectif filtrant décroissant des
voisinages ouverts affines de $y$, et appliquons la méthode de (8.1.2,
\emph{a)}). Puisque $Y$ est noethérien, les restrictions
$f_\lambda:f^{-1}(U_\lambda)\to U_\lambda$ de $f$ sont de présentation
finie et il en est de même de $f_y$. Par hypothèse $f_y$ se factorise en
$X_y\xrightarrow{g_y}Z_y\xrightarrow{u_y}\operatorname{Spec}(\mathcal O_y)$,
où $u_y$ est fini et $g_y$ est une immersion. Comme $\mathcal O_y$ est
noethérien, il en est de même de $Z_y$, et comme $Z_y$ est de présentation
finie sur $\operatorname{Spec}(\mathcal O_y)$, il existe un $\lambda$ et un
morphisme de présentation finie $u_\lambda:Z_\lambda\to U_\lambda$ tels que
$Z_y$ s'identifie à
$Z_\lambda\times_{U_\lambda}\operatorname{Spec}(\mathcal O_y)$ et $u_y$ à
$u_\lambda\times 1$ (8.8.2, (ii))~; en outre, il existe un morphisme
$g_\lambda:X_\lambda\to Z_\lambda$ tel que $g_y=g_\lambda\times 1$ et que
$f_\lambda=u_\lambda\circ g_\lambda$ (8.8.2, (i)). De plus, on peut supposer
$\lambda$ pris de sorte que $g_\lambda$ soit une immersion et $u_\lambda$ un
morphisme fini (8.10.5, (ii) et (x)), ce qui prouve que $f_\lambda$ est
pseudo-fini.
\end{proof}

\begin{env}[8.12.5]
\label{IV.8.12.5-fr}
Nous pouvons maintenant donner du «~Main theorem~» de Zariski (III, 4.4.3),
une démonstration n'utilisant pas les résultats cohomologiques de nature
«~globale~» du chap.~III, mais faisant appel par contre aux propriétés plus
fines des anneaux locaux noethériens~; nous généraliserons en outre l'énoncé
du théorème en le débarrassant des hypothèses noethériennes~:
\end{env}

\begin{theorem}[8.12.6]
\label{IV.8.12.6-fr}
(«~Main Theorem~» de Zariski). --- Soit $Y$ un préschéma quasi-compact et
quasi-séparé. Si un morphisme $f:X\to Y$ est quasi-fini, séparé et de
présentation finie, il existe une factorisation de $f$
\begin{equation}
\label{IV.8.12.6.1-fr}
X\xrightarrow{f'}Y'\xrightarrow{u}Y. \tag{8.12.6.1}
\end{equation}
où $f'$ est une immersion ouverte et $u$ un morphisme fini.
\end{theorem}

\oldpage[IV-3]{46}
\begin{proof}
En vertu de (8.12.4, (ii)) et du caractère local (sur $Y$) des notions de
morphisme quasi-fini, séparé et de présentation finie, on peut se borner au
cas où $Y=\operatorname{Spec}(A)$ est affine. Appliquant (8.9.1), on peut
supposer qu'il y a un sous-anneau $A_0$ de $A$, qui est une
$\mathbf Z$-algèbre de type fini, et un $A$-isomorphisme
$X_0\otimes_{A_0}A\xrightarrow{\sim}X$, $f$ s'identifiant par cet
isomorphisme à $f_0\times 1$, où
$f_0:X_0\to\operatorname{Spec}(A_0)$ est un morphisme \emph{de type fini}~;
en outre (8.10.5, (v) et (xii)) on peut supposer que $f_0$ est séparé et
quasi-fini~; si l'on prouve que $f_0$ est pseudo-fini, il en sera de même de
$f$ d'après (8.12.4, (i)). Comme $A_0$ est alors noethérien et que les
notions de morphisme de type fini, séparé et quasi-fini se conservent par
changement de base, il résulte de (8.12.4, (iii)) que l'on peut même supposer
que $A$ est un anneau local, \emph{essentiellement de type fini sur}
$\mathbf Z$ (1.3.8). Posons $n=\dim(A)$, et procédons par récurrence sur $n$~;
pour $n=0$, le théorème est évident, $A$ étant un corps et le morphisme $f$
étant déjà fini (II, 6.2.2). Posons $B=\Gamma(X,\mathcal O_X)$~; désignons
par $C$ la fermeture intégrale de $A$ dans $B$, posons
$Z=\operatorname{Spec}(C)$ et soit $g:X\to Z$ le $Y$-morphisme correspondant
au $A$-homomorphisme canonique injectif $C\to B$~; en vertu de (8.12.3), il
s'agit de montrer que $g$ est une immersion ouverte. Soit $a$ le point fermé
de $Y$, et soit $U=Y-\{a\}$~; $U$ est noethérien et tous ses anneaux locaux
sont essentiellement de type fini sur $\mathbf Z$ et de dimension $<n$~;
compte tenu de l'hypothèse de récurrence, et de (8.12.4, (iii)), on voit que
la restriction $f^{-1}(U)\to U$ de $f$ est un morphisme pseudo-fini. On en
conclut (8.12.3) que, si $h:Z\to Y$ est le morphisme structural, la
restriction $f^{-1}(U)\to h^{-1}(U)$ de $g$ est une immersion ouverte.
Posons $A'=\widehat A$, $Y'=\operatorname{Spec}(A')$,
$X'=X_{(Y')}$, $f'=f_{(Y')}:X'\to Y'$. Comme le morphisme canonique
$u:Y'\to Y$ est plat, il résulte de (2.3.1) que
$B'=\Gamma(X',\mathcal O_{X'})$ s'identifie à la $A'$-algèbre
$B\otimes_A A'$.

D'autre part, comme $A$ est un anneau local \emph{excellent} (7.8.3), le
morphisme $u:Y'\to Y$ est régulier, et \emph{a fortiori} normal, et par suite
(6.14.4) \emph{la fermeture intégrale de $A'$ dans $B'$ est égale à}
$C'=C\otimes_A A'$. On voit donc que $Z'=\operatorname{Spec}(C')$ est égal à
$Z_{(Y')}$ et le morphisme $g':X'\to Z'$ provenant de l'injection $C'\to B'$
est égal à $g_{(Y')}$. Comme $u:Y'\to Y$ est fidèlement plat et
quasi-compact, pour prouver que $g$ est une immersion ouverte, il suffit de
prouver que $g'$ est une immersion ouverte (2.7.1, (x)). Notons maintenant
que $u^{-1}(a)$ est réduit au point fermé $a'$ de $Y'$ et par suite
$U'=Y'-\{a'\}=u^{-1}(U)$. Si $h':Z'\to Y'$ est le morphisme canonique, le
fait que la restriction $f^{-1}(U)\to h^{-1}(U)$ de $g$ soit une immersion
ouverte entraîne qu'il en est de même de la restriction
$f'^{-1}(U')\to h'^{-1}(U')$ de $g'$. Notons maintenant que $f'$ est un
morphisme séparé et quasi-fini (II, 6.2.4)~; comme $A'$ est \emph{complet},
on déduit de (II, 6.2.6) que $X'$ est $Y'$-isomorphe à une somme
$X'_1\amalg X'_2$, où la restriction
$f'|X'_1=f'_1:X'_1\to Y'$ est un morphisme \emph{fini}, et
$X'_2\subset f'^{-1}(U')$. Il en résulte que $B'$ est composée directe des
deux $A'$-algèbres
$\Gamma(X'_1,\mathcal O_{X'_1})=B'_1$ et
$\Gamma(X'_2,\mathcal O_{X'_2})=B'_2$~; on en conclut aussitôt que la
fermeture intégrale $C'$ de $A'$ dans $B'$ est composée directe des
fermetures intégrales $C'_1$, $C'_2$ de $A'$ dans $B'_1$, $B'_2$
respectivement, d'où $Z'=Z'_1\amalg Z'_2$, où
$Z'_i=\operatorname{Spec}(C'_i)$ ($i=1,2$)~; et le morphisme canonique
$g':X'\to Z'$ est tel que $g'|X'_i$ soit le morphisme canonique
$g'_i:X'_i\to Z'_i$ ($i=1,2$). Mais comme $B'_1$ est déjà une
$A'$-algèbre finie, on a $C'_1=B'_1$, et $g'_1$ est donc un isomorphisme.
D'autre part, comme $X'_2\subset f'^{-1}(U')$ et est ouvert dans
$f'^{-1}(U')$, on sait
\oldpage[IV-3]{47}
déjà que $g'_2$ est une immersion ouverte. On conclut bien que $g'$ est une
immersion ouverte.
\end{proof}

\begin{remark}[8.12.7]
\label{IV.8.12.7-fr}
Lorsque, dans (8.12.6), on suppose que $Y$ est un \emph{schéma affine}, la
démonstration par réduction au cas noethérien montre que, dans la
factorisation (8.12.6.1), les morphismes $f'$ et $u$ sont aussi des
morphismes \emph{de présentation finie} (1.6.2).
\end{remark}

\begin{corollary}[8.12.8]
\label{IV.8.12.8-fr}
Soient $Y$ un schéma quasi-compact tel qu'il existe un $\mathcal O_Y$-Module
ample (II, 4.5.3), $f:X\to Y$ un morphisme quasi-fini et quasi-projectif.
Alors il existe une factorisation de $f$ en
\[
X\xrightarrow{f'}Y'\xrightarrow{u}Y
\]
où $f'$ est une immersion ouverte et $u$ un morphisme fini.
\end{corollary}

\begin{proof}
L'hypothèse entraîne que $X$ s'identifie à un sous-$Y$-schéma quasi-compact
d'un $Y$-schéma de la forme $Z=\mathbf P_Y^r$ (II, 5.3.3). Il y a par suite
un voisinage ouvert quasi-compact $U$ de $X$ dans $Z$ tel que $X$ soit fermé
dans $U$~; comme $Z$ est un schéma, l'injection canonique $U\to Z$ est un
morphisme de présentation finie ((1.2.7) et (1.6.2)), donc le morphisme
composé $g:U\to Z\to Y$ est aussi un morphisme de présentation finie (le
fait que $\mathbf P_Y^r$ est de présentation finie sur $Y$ résultant
aussitôt de la définition (II, 4.1.1)). Soit $\mathcal J$ l'Idéal
quasi-cohérent de $\mathcal O_U$ définissant le sous-préschéma fermé $X$~;
comme $U$ est un schéma quasi-compact, $\mathcal J$ est limite inductive
filtrante de ses sous-Idéaux quasi-cohérents de type fini $\mathcal J_\lambda$
(I, 9.4.9). Si $X_\lambda$ est le sous-préschéma fermé de $U$ défini par
$\mathcal J_\lambda$, on a par suite $X=\bigcap_\lambda X_\lambda$. Pour tout
$y\in Y$, on a donc
$f^{-1}(y)=\bigcap_\lambda(X_\lambda\cap g^{-1}(y))$, et comme les ensembles
$X_\lambda\cap g^{-1}(y)$ sont fermés dans l'espace noethérien
$g^{-1}(y)$, il existe pour tout $y$ un indice $\lambda(y)$ tel que
$f^{-1}(y)=X_{\lambda(y)}\cap g^{-1}(y)$. Désignons par $E_\lambda$
l'ensemble des $y\in Y$ tels que la fibre $X_\lambda\cap g^{-1}(y)$ de la
restriction de $g$ à $X_\lambda$ soit un $k(y)$-préschéma fini. L'hypothèse
que $f$ est quasi-fini entraîne, en vertu de ce qui précède, que
$Y=\bigcup_\lambda E_\lambda$. Or, chacun des $X_\lambda$ est, par
définition, \emph{de présentation finie} sur $Y$~; il résulte donc de
(9.2.3) et (9.2.6)\footnote{Le lecteur vérifiera que les corollaires (8.12.8)
à (8.12.11) ne sont pas utilisés dans le \S~9.} que les $E_\lambda$ sont des
ensembles \emph{constructibles} dans le schéma $Y$~; comme ils forment un
ensemble filtrant croissant, il existe un indice $\lambda$ tel que
$E_\lambda=Y$ (1.9.9), et pour cet indice $\lambda$, le morphisme
$f_\lambda:X_\lambda\to Y$, restriction de $g$ à $X_\lambda$, est donc
quasi-fini. Comme il est de présentation finie et séparé, on peut lui
appliquer (8.12.6), et $f_\lambda$ se factorise donc en
\[
X_\lambda\xrightarrow{j_\lambda}Y'_\lambda
\xrightarrow{u_\lambda}Y
\]
où $j_\lambda$ est une immersion et $u_\lambda$ un morphisme fini. Comme
$X$ est un sous-préschéma fermé de $X_\lambda$, on a ainsi prouvé que $f$
possède la propriété (8.12.3, \emph{a)}), d'où le corollaire en vertu de
l'équivalence de (8.12.3, \emph{a)}) et (8.12.3, \emph{a$'$)}).
\end{proof}

\oldpage[IV-3]{48}
\begin{corollary}[8.12.9]
\label{IV.8.12.9-fr}
Soit $f:X\to Y$ un morphisme localement quasi-fini
$(\mathbf{Err}_{\mathrm{III}}, 20)$. Pour tout $x\in X$ il existe un
voisinage ouvert $U$ de $x$ dans $X$, un voisinage ouvert $V$ de
$y=f(x)$ dans $Y$, tels que $f(U)\subset V$ et une factorisation
\[
U\xrightarrow{f'}V'\xrightarrow{u}V
\]
de la restriction de $f$ à $U$, où $f'$ est une immersion ouverte et $u$ un
morphisme fini.
\end{corollary}

Il suffit évidemment de prendre pour $V$ un voisinage affine de $y$ dans
$Y$, pour $U$ un voisinage affine de $x$ dans $X$ contenu dans $f^{-1}(U)$
et tel que $f|U$ soit quasi-fini. Le morphisme $U\to V$ restriction de $f$
étant alors affine (donc quasi-projectif), on peut lui appliquer (8.12.8).

\begin{corollary}[8.12.10]
\label{IV.8.12.10-fr}
Soient $Y$ un préschéma intègre et normal, $X$ un préschéma intègre,
$f:X\to Y$ un morphisme birationnel et localement quasi-fini
$(\mathbf{Err}_{\mathrm{III}}, 20)$. Alors $f$ est un isomorphisme local~;
pour que $f$ soit une immersion ouverte, il faut et il suffit que $f$ soit
en outre séparé.
\end{corollary}

La seconde assertion résulte aussitôt de la première et de (I, 8.2.8). Pour
prouver la première assertion, on peut supposer $X$ et $Y$ affines et $f$
quasi-fini~; considérons la factorisation $f=u\circ f'$ de (8.12.8), qui
permet d'identifier $X$ par $f'$ à un sous-préschéma induit sur un ouvert de
$Y'$. Comme $X$ est intègre, on peut, en vertu de (I, 5.2.3), remplacer
$Y'$ par le sous-préschéma réduit de $Y'$ ayant pour espace sous-jacent
$\overline X$, donc on peut supposer aussi que $Y'$ est \emph{intègre}. En
outre, puisque $f$ est birationnel, $u$ l'est aussi. La conclusion résulte
alors du lemme suivant~:

\begin{lemma}[8.12.10.1]
\label{IV.8.12.10.1-fr}
Soient $Y'$ un préschéma intègre, $Y$ un préschéma intègre et normal~; alors
un morphisme $u:Y'\to Y$ fini et birationnel est un isomorphisme.
\end{lemma}

Posons en effet $\mathcal A=u_*(\mathcal O_{Y'})$, de sorte que $\mathcal A$
est une $\mathcal O_Y$-Algèbre finie, $Y'$ s'identifiant à
$\operatorname{Spec}(\mathcal A)$ (II, 1.3.6). Si $R(Y)$ est le corps des
fonctions rationnelles de $Y$, on a donc, pour tout $y\in Y$,
$\mathcal O_{Y,y}\subset\mathcal A_y\subset R(Y)$~; mais comme l'anneau
$\mathcal O_{Y,y}$ est par hypothèse intégralement clos et a $R(Y)$ pour
corps des fractions, on a nécessairement
$\mathcal A_y=\mathcal O_{Y,y}$, d'où le lemme.

\begin{corollary}[8.12.11]
\label{IV.8.12.11-fr}
Soient $Y$ un préschéma intègre, $X$ un préschéma intègre et normal,
$f:X\to Y$ un morphisme dominant et localement quasi-fini. Soient $K$ et
$L$ (extension de $K$) les corps de fonctions rationnelles de $Y$ et $X$
respectivement~; et soit $Y'$ la fermeture intégrale de $Y$ relativement à
$L$ (II, 6.3.4)~; alors $f$ se factorise d'une seule manière en
$f:X\xrightarrow{f'}Y'\xrightarrow{u}Y$, où $f'$ est birationnel, et
correspond à l'automorphisme identique de $L$~; $f'$ est alors un
isomorphisme local, et pour que $f'$ soit une immersion ouverte, il faut et
il suffit que $f$ soit séparé.
\end{corollary}

L'existence et l'unicité de la factorisation de $f$ résultent de (II, 6.3.9).
Il résulte de (II, 6.2.4, (v)), en se ramenant au cas affine, que $f'$ est
localement quasi-fini~; en outre, il résulte de (I, 5.5.1) que, pour que
$f'$ soit séparé, il faut et il suffit que $f$ le soit, puisque $u$ est
affine, donc séparé~; les deux dernières assertions sont donc des
conséquences de (8.12.10) appliqué à $f'$.

\oldpage[IV-3]{49}
\subsection{Traduction en termes de pro-objets}
\label{subsection:IV.8.13-fr}

La proposition suivante est essentiellement équivalente à (8.8.2, (i))~:

\begin{proposition}[8.13.1]
\label{IV.8.13.1-fr}
Soient $S$ un préschéma, $(X_\lambda,v_{\lambda\mu})$ un système projectif
filtrant de $S$-préschémas~; on suppose qu'il existe $\alpha$ tel que
$v_{\alpha\lambda}$ soit un morphisme affine pour tout $\lambda\geqslant
\alpha$ (ce qui entraîne (II, 1.6.2) que $v_{\lambda\mu}$ soit affine pour
$\alpha\leqslant\lambda\leqslant\mu$), de sorte que la limite projective
$X=\varprojlim_\lambda X_\lambda$ existe dans la catégorie des $S$-préschémas
(8.2.3). Soit $Y$ un $S$-préschéma, et, pour tout $\lambda\geqslant\alpha$,
soit
\[
e_\lambda:\operatorname{Hom}_S(X_\lambda,Y)
\longrightarrow\operatorname{Hom}_S(X,Y)
\]
l'application qui, à tout $S$-morphisme $f_\lambda:X_\lambda\to Y$, fait
correspondre $f=f_\lambda\circ v_\lambda$, où $v_\lambda:X\to X_\lambda$
est le morphisme canonique. La famille $(e_\lambda)$ est un système inductif
d'applications, qui définit donc une application canonique
\begin{equation}
\label{IV.8.13.1.1-fr}
\varinjlim_\lambda\operatorname{Hom}_S(X_\lambda,Y)
\longrightarrow\operatorname{Hom}_S(X,Y). \tag{8.13.1.1}
\end{equation}
Supposons $X_\alpha$ quasi-compact (resp. quasi-compact et quasi-séparé), et le
morphisme structural $Y\to S$ localement de type fini (resp. localement de
présentation finie). Alors l'application (8.13.1.1) est injective (resp.
bijective).
\end{proposition}

Posons en effet, pour $\lambda\geqslant\alpha$,
$Z_\lambda=Y\times_S X_\lambda$, de sorte que l'on a
$Z_\lambda=Z_\alpha\times_{X_\alpha}X_\lambda$. Posons de même
$Z=Y\times_S X=Z_\alpha\times_{X_\alpha}X$~; on sait alors (8.2.5) que, si
l'on pose aussi
\[
w_{\lambda\mu}=1\times v_{\lambda\mu}:Z_\mu\to Z_\lambda
\quad(\alpha\leqslant\lambda\leqslant\mu),
\qquad
w_\lambda=1\times v_\lambda:Z\to Z_\lambda
\quad(\alpha\leqslant\lambda),
\]
$Z$ est limite projective du système projectif $(Z_\lambda,w_{\lambda\mu})$
et les $w_\lambda$ sont les morphismes canoniques correspondants. Notons
d'autre part que le morphisme $Z_\alpha\to X_\alpha$ est localement de type
fini (resp. localement de présentation finie) (1.3.4 et 1.4.3). Enfin, on
sait que l'on a
\[
\operatorname{Hom}_S(X_\lambda,Y)
=\operatorname{Hom}_{X_\lambda}(X_\lambda,Z_\lambda)
\quad\hbox{et}\quad
\operatorname{Hom}_S(X,Y)=\operatorname{Hom}_X(X,Z)
\]
(I, 3.3.14). Il suffit maintenant d'appliquer (8.8.2, (i)) en y faisant
$X_\lambda=S_\lambda$ et en y remplaçant $Y_\lambda$ par $Z_\lambda$.

\begin{corollary}[8.13.2]
\label{IV.8.13.2-fr}
Avec les notations de (8.13.1), on suppose $X_\alpha$ quasi-compact et
quasi-séparé, et les $v_{\alpha\lambda}$ affines pour
$\alpha\leqslant\lambda$~; on suppose en outre que
$Y=\varprojlim_\rho Y_\rho$, où $(Y_\rho,t_{\rho\sigma})$ est un système
projectif filtrant de $S$-préschémas tel que, pour chaque $\rho$, le morphisme
structural $Y_\rho\to S$ soit localement de présentation finie. On a alors une
bijection canonique
\begin{equation}
\label{IV.8.13.2.1-fr}
\operatorname{Hom}_S(X,Y)
\xrightarrow{\sim}
\varprojlim_\rho\left(\varinjlim_\lambda
\operatorname{Hom}_S(X_\lambda,Y_\rho)\right). \tag{8.13.2.1}
\end{equation}
\end{corollary}

En effet, le fait que $Y$ soit limite projective des $Y_\rho$ entraîne en
particulier que l'application canonique
\[
\operatorname{Hom}_S(X,Y)\longrightarrow
\varprojlim_\rho\operatorname{Hom}_S(X,Y_\rho)
\]
est bijective~; et d'autre part, les hypothèses entraînent, pour chaque
$\rho$, l'existence d'une bijection canonique
\[
\operatorname{Hom}_S(X,Y_\rho)\xrightarrow{\sim}
\varinjlim_\lambda\operatorname{Hom}_S(S_\lambda,Y_\rho)
\]
en vertu de (8.13.1)~; d'où la conclusion.

\begin{env}[8.13.3]
\label{IV.8.13.3-fr}
Les résultats précédents permettent d'interpréter dans la théorie des
préschémas les notions de «~pro-variété~» ou de «~pro-schéma~» qui
s'introduisent dans certaines applications (par exemple dans la théorie du
corps de classes local suivant les idées de Serre [39] ou dans la théorie de
Néron sur la réduction des variétés
\oldpage[IV-3]{50}
abéliennes [32]). Rappelons rapidement ici la notion de \emph{pro-objet} d'une
catégorie, renvoyant au chap.~V pour de plus amples développements (nous
n'utiliserons d'ailleurs pas avant le chap.~V l'interprétation qui va suivre,
et le lecteur peut donc omettre jusque-là la lecture de la fin de ce numéro).
Étant donnée une catégorie $\mathcal C$, la catégorie
$\mathbf{Pro}(\mathcal C)$ des \emph{pro-objets} de $\mathcal C$ a pour objets
les \emph{systèmes projectifs} (dans l'univers où l'on se place)
$\mathbf X=(X_\mu)_{\mu\in M}$ d'objets de $\mathcal C$ dont les ensembles
d'indices (dépendant du système projectif considéré) sont supposés préordonnés
filtrants. Étant donnés deux tels pro-objets
$\mathbf X=(X_\mu)_{\mu\in M}$,
$\mathbf X'=(X'_{\mu'})_{\mu'\in M'}$, les \emph{morphismes} de $\mathbf X$
dans $\mathbf X'$ sont par définition les éléments de l'ensemble
\[
\varprojlim_{\mu'}\left(\varinjlim_\mu
\operatorname{Hom}(X_\mu,X'_{\mu'})\right)~;
\]
la vérification du fait que l'on peut prendre ces ensembles pour ensembles de
morphismes est immédiate, la composition de systèmes de morphismes
$u_{\mu'}^\mu:X_\mu\to X'_{\mu'}$,
$u_{\mu''}^{\prime\mu'}:X'_{\mu'}\to X''_{\mu''}$ qui sont inductifs en
l'indice supérieur et projectifs en l'indice inférieur, se faisant
«~argument par argument~», autrement dit en considérant le système des
$u_{\mu''}^{\prime\mu'}\circ u_{\mu'}^\mu$.
\end{env}

\begin{env}[8.13.4]
\label{IV.8.13.4-fr}
Considérons alors un préschéma quasi-compact et quasi-séparé $S$, et désignons
par $\mathcal C$ la sous-catégorie pleine de la catégorie
$(\mathbf{Sch})_{/S}$ des $S$-préschémas, formée des $S$-préschémas $X$ ayant
la propriété suivante~: le morphisme structural $X\to S$ se factorise en
\[
X\xrightarrow{g}X_0\xrightarrow{f}S,
\]
où $g:X\to X_0$ est affine et $f:X_0\to S$ de présentation finie~; nous
dirons pour abréger que les préschémas de $\mathcal C$ sont
\emph{essentiellement affines au-dessus de $S$}.

Considérons d'autre part la sous-catégorie pleine $\mathcal C'_0$ de
$(\mathbf{Sch})_{/S}$ formée des $S$-préschémas \emph{de présentation finie},
et la catégorie $\mathbf{Pro}(\mathcal C'_0)$ des pro-objets de
$\mathcal C'_0$. Nous dirons qu'un objet
$\mathbf X=(X_\mu)_{\mu\in M}$ de $\mathbf{Pro}(\mathcal C'_0)$ est
\emph{essentiellement affine} s'il existe un $\gamma\in M$ tel que pour tout
$\mu\geqslant\gamma$, le morphisme de transition
$v_{\gamma\mu}:X_\mu\to X_\gamma$ soit affine (ce qui entraîne que pour
$\gamma\leqslant\mu\leqslant\nu$, $v_{\mu\nu}:X_\nu\to X_\mu$ est affine).
On notera qu'un objet de $\mathbf{Pro}(\mathcal C'_0)$ \emph{isomorphe} à un
objet essentiellement affine \emph{n'est pas} nécessairement essentiellement
affine lui-même.

Nous désignerons par $\mathcal C'$ la sous-catégorie pleine de
$\mathbf{Pro}(\mathcal C'_0)$ formée des pro-objets de $\mathcal C'_0$
essentiellement affines.

Cela étant, il résulte de (8.2.2) et (8.2.3) que pour tout objet
$\mathbf X=(X_\mu)_{\mu\in M}$ de $\mathcal C'$, le $S$-préschéma
$X=\varprojlim_\mu X_\mu$ existe~; en outre, comme, pour $\mu$ assez grand,
le morphisme $X\to X_\mu$ est affine (8.2.2), $X$ est
\emph{essentiellement affine au-dessus de $S$} par définition. Posons
$X=L(\mathbf X)$~; montrons qu'on a ainsi défini un \emph{foncteur canonique}
\begin{equation}
\label{IV.8.13.4.1-fr}
L:\mathcal C'\longrightarrow\mathcal C. \tag{8.13.4.1}
\end{equation}

On a en effet, pour deux objets $\mathbf X=(X_\mu)$,
$\mathbf X'=(X'_{\mu'})$ de $\mathcal C'$, une application canonique pour
chaque $\mu'$
\[
\varinjlim_\mu\operatorname{Hom}_S(X_\mu,X'_{\mu'})
\longrightarrow
\operatorname{Hom}_S(\varprojlim_\mu X_\mu,X'_{\mu'})
\]
définie dans (8.13.1.1), et d'autre part, par définition de la limite
projective, une bijection canonique
\[
\varprojlim_{\mu'}
\operatorname{Hom}_S(\varprojlim_\mu X_\mu,X'_{\mu'})
\xrightarrow{\sim}
\operatorname{Hom}_S(\varprojlim_\mu X_\mu,
\varprojlim_{\mu'}X'_{\mu'}),
\]
\oldpage[IV-3]{51}
d'où une application canonique
\begin{equation}
\label{IV.8.13.4.2-fr}
\varprojlim_{\mu'}\left(\varinjlim_\mu
\operatorname{Hom}_S(X_\mu,X'_{\mu'})\right)
\longrightarrow
\operatorname{Hom}_S(\varprojlim_\mu X_\mu,
\varprojlim_{\mu'}X'_{\mu'}), \tag{8.13.4.2}
\end{equation}
évidemment fonctorielle en $\mathbf X$ et $\mathbf X'$, et qui complète la
définition du foncteur $L$.
\end{env}

\begin{proposition}[8.13.5]
\label{IV.8.13.5-fr}
Les hypothèses et notations étant celles de (8.13.4), le foncteur $L$ est
pleinement fidèle. Si de plus $S$ est un préschéma noethérien (ce qui implique
déjà que $S$ est quasi-compact et quasi-séparé (1.2.8)), $L$ est une
\emph{équivalence de catégories}.
\end{proposition}

Dire que $L$ est pleinement fidèle signifie que l'application (8.13.4.2) est
bijective quels que soient $\mathbf X$, $\mathbf X'$ dans $\mathcal C'$, ce
qui est un cas particulier de (8.13.2)~: en effet, les morphismes structuraux
$X_\mu\to S$ étant de présentation finie, sont en particulier quasi-compacts
et quasi-séparés, donc les $X_\mu$ sont quasi-compacts et quasi-séparés.

Pour montrer que lorsque $S$ est noethérien $L$ est une équivalence de
catégories, il suffit, puisque l'on sait déjà que $L$ est pleinement fidèle,
de prouver que tout $S$-préschéma essentiellement affine $X$ est
$S$-\emph{isomorphe} à un objet de la forme $L(\mathbf X)$ où
$\mathbf X\in\mathcal C'$ $(\mathbf 0_{\mathrm{III}},8.1.5)$.

Or, il y a par hypothèse une factorisation
$X\xrightarrow{g}X_0\xrightarrow{f}S$ du morphisme structural, $f$ étant de
présentation finie et $g$ affine. On peut donc écrire
$X=\operatorname{Spec}(\mathcal A)$, où $\mathcal A$ est une
$\mathcal O_{X_0}$-Algèbre quasi-cohérente (II, 1.3.1). Or, comme $X_0$ est
noethérien (puisqu'il en est ainsi de $S$ et que $f$ est de type fini),
$\mathcal A$ est limite inductive filtrante de la famille
$(\mathcal A_\lambda)$ de ses sous-$\mathcal O_{X_0}$-Algèbres quasi-cohérentes
de type fini (I, 9.6.6). Posons
$X_\lambda=\operatorname{Spec}(\mathcal A_\lambda)$~; les morphismes
$X_\lambda\to X_0$ sont de type fini, donc de présentation finie puisque
$X_0$ est noethérien, et par suite il en est de même des morphismes composés
$X_\lambda\to X_0\to S$ (1.6.2)~; autrement dit, les $X_\lambda$
appartiennent à $\mathcal C'_0$ et comme les morphismes $X_\lambda\to X_0$
sont affines, $\mathbf X=(X_\lambda)$ est un objet de $\mathcal C'$ dont la
limite projective existe et est $S$-isomorphe à $X$ en vertu de (8.2.2). Ceci
achève la démonstration.

\begin{remark}[8.13.6]
\label{IV.8.13.6-fr}
Il résulte de (1.6.2) et de (II, 1.6.2) que si $X$ et $Y$ sont
essentiellement affines au-dessus de $S$, alors il en est de même de
$X\times_S Y$. On en conclut par exemple $(\mathbf 0_{\mathrm{III}},8.2.5)$
qu'un $\mathcal C$-groupe n'est autre qu'un $(\mathbf{Sch})_{/S}$-groupe qui
est un préschéma essentiellement affine au-dessus de $S$. D'autre part, les
produits finis existent dans la catégorie $\mathcal C'$~: en effet, si
$\mathbf X=(X_\mu)_{\mu\in M}$, $\mathbf Y=(Y_\rho)_{\rho\in R}$ sont deux
objets de $\mathcal C'$, les produits $X_\mu\times_S Y_\rho$ sont des
$S$-préschémas de présentation finie, et en prenant pour morphismes de
transition
\[
X_\nu\times_S Y_\sigma\longrightarrow X_\mu\times_S Y_\rho
\]
les produits des morphismes de transition $X_\nu\to X_\mu$ et
$Y_\sigma\to Y_\rho$, on voit aussitôt que
$(X_\mu\times_S Y_\rho)$ est produit de $\mathbf X$ et de $\mathbf Y$ dans
$\mathbf{Pro}(\mathcal C'_0)$~; en outre (II, 1.6.2) les morphismes de
transition ainsi définis sont affines pour $\mu$ et $\rho$ assez grands, donc
le produit $\mathbf X\times\mathbf Y$ ainsi défini appartient bien à
$\mathcal C'$. On conclut alors comme plus haut qu'un $\mathcal C'$-groupe est
un $\mathbf{Pro}(\mathcal C'_0)$-groupe qui est essentiellement affine. On
déduit donc de (8.13.5) que les catégories des $\mathcal C'$-groupes et des
$\mathcal C$-groupes sont \emph{équivalentes} lorsque $S$ est
\emph{noethérien}. Il semble plausible que lorsque $S$ est le spectre d'un
corps $k$, la catégorie des $\mathcal C$-groupes soit \emph{équivalente} à
celle des $\boldsymbol K$-groupes, où $\boldsymbol K$ est la catégorie des
$S$-préschémas \emph{quasi-compacts}~; autrement dit, tout préschéma en
groupes sur $k$ qui est quasi-compact serait essentiellement affine. D'autre
part, si l'on désigne
\oldpage[IV-3]{52}
par $\mathcal C'_0$-$\boldsymbol{gr}$ la catégorie des
$\mathcal C'_0$-groupes, il est vraisemblable que la catégorie des
$\mathcal C'$-groupes est équivalente à la sous-catégorie pleine de
$\mathbf{Pro}(\mathcal C'_0\text{-}\boldsymbol{gr})$ formée des «~groupes
proalgébriques essentiellement affines~», c'est-à-dire des systèmes
projectifs $\mathbf G=(G_\mu)_{\mu\in M}$, où les $G_\mu$ sont des groupes
algébriques sur $k$ et les morphismes de transition $G_\nu\to G_\mu$ sont
affines pour $\mu$ assez grand (ce que l'on peut exprimer aussi en disant que
$\mathbf G$ est extension d'un groupe algébrique par un pro-groupe algébrique
affine). La conjonction de ces deux conjectures équivaut d'ailleurs à la
suivante~: tout préschéma en groupes quasi-compacts sur $k$ est «~extension~»
d'un «~groupe algébrique~» (i.e. un préschéma en groupes de type fini sur $k$)
par un préschéma en groupes \emph{affine} sur $k$.

Les seuls pro-groupes algébriques rencontrés en pratique jusqu'à présent étant
en fait essentiellement affines, il y aura donc sans doute avantage à
substituer à l'étude des groupes pro-algébriques généraux (introduits et
étudiés par Serre [40]) celle des schémas en groupes quasi-compacts sur $k$,
dont la définition est conceptuellement plus simple.
\end{remark}

\subsection{Caractérisation d'un préschéma localement de présentation finie sur un autre, en termes du foncteur qu'il représente}
\label{subsection:IV.8.14-fr}

\begin{env}[8.14.1]
\label{IV.8.14.1-fr}
Étant donné un préschéma $S$, nous dirons encore, comme en (8.13.4), qu'un
système projectif filtrant $(X_\lambda,v_{\lambda\mu})$ de $S$-préschémas est
\emph{essentiellement affine} s'il existe $\alpha$ tel que
$v_{\alpha\lambda}$ soit un morphisme affine pour $\lambda\geqslant\alpha$.

L'énoncé suivant, qui sera surtout utile au chap.~V, précise (8.8.2, (i)) en
lui fournissant une réciproque~:
\end{env}

\begin{proposition}[8.14.2]
\label{IV.8.14.2-fr}
Soient $S$ un préschéma, $f:X\to S$ un morphisme. Pour tout $S$-préschéma
$T$, posons
\[
h_X(T)=\operatorname{Hom}_S(T,X)
\]
de sorte que $h_X$ est un foncteur contravariant de la catégorie
$(\mathbf{Sch})_{/S}$ des $S$-préschémas dans la catégorie
$\mathbf{Ens}$ des ensembles $(\mathbf 0_{\mathrm{III}},8.1.1)$, et $X$ un
objet représentant ce foncteur $(\mathbf 0_{\mathrm{III}},8.1.8)$. Les
conditions suivantes sont équivalentes~:
\begin{enumerate}[label=\emph{\alph*)}]
  \item $f$ est localement de présentation finie.
  \item Pour tout système projectif filtrant $(Z_\lambda)$ de
  $S$-préschémas, essentiellement affine (8.13.4) et formé de préschémas
  quasi-compacts et quasi-séparés, l'application canonique (8.13.1.1)
\begin{equation}
\label{IV.8.14.2.1-fr}
\varinjlim_\lambda h_X(Z_\lambda)
\longrightarrow h_X(\varprojlim_\lambda Z_\lambda) \tag{8.14.2.1}
\end{equation}
  est bijective.
  \item Pour tout système projectif filtrant $(Z_\lambda)$ de
  $S$-préschémas, tel que les $Z_\lambda$ soient des schémas affines,
  l'application (8.14.2.1) est bijective.
  \item[\rm c$'$)] Pour tout ouvert affine $U$ de $S$ et tout système
  projectif filtrant $(Z_\lambda)$ de $U$-préschémas, tel que les $Z_\lambda$
  soient des schémas affines, l'application (8.14.2.1) est bijective.
\end{enumerate}
\end{proposition}

Le fait que \emph{a)} implique \emph{b)} n'est autre que (8.13.1)~; il est
trivial que \emph{b)} implique \emph{c)} et que \emph{c)} implique
\emph{c$'$)}. Reste à voir que \emph{c$'$)} entraîne \emph{a)}, et comme la
propriété \emph{a)} est locale sur $S$, on peut se borner au cas où $S$ est
\emph{affine}.

\oldpage[IV-3]{53}
Supposons d'abord que $X$ soit aussi affine~; l'assertion à prouver est alors
équivalente au

\begin{corollary}[8.14.2.2]
\label{IV.8.14.2.2-fr}
Soient $A$ un anneau, $B$ une $A$-algèbre. Afin que, pour tout système
inductif filtrant $(C_\lambda)$ de $A$-algèbres, l'application canonique
\begin{equation}
\label{IV.8.14.2.3-fr}
\varinjlim_\lambda\operatorname{Hom}_{A\text{-}\mathrm{alg.}}(B,C_\lambda)
\longrightarrow
\operatorname{Hom}_{A\text{-}\mathrm{alg.}}
(B,\varinjlim_\lambda C_\lambda) \tag{8.14.2.3}
\end{equation}
soit bijective, il faut et il suffit que $B$ soit une $A$-algèbre de
présentation finie.
\end{corollary}

Il reste seulement à montrer que la condition est nécessaire. Prenons d'abord
pour $(C_\lambda)$ le système inductif filtrant des sous-$A$-algèbres de type
fini de $B$, de sorte que $\varinjlim_\lambda C_\lambda=B$. Le fait que (8.14.2.3)
soit bijective entraîne en particulier que l'application identique
$1_B$ se factorise en $B\to C_\lambda\to B$ pour un $\lambda$ convenable,
ce qui entraîne $C_\lambda=B$, donc $B$ est une $A$-algèbre \emph{de type
fini}. Posons alors $B=C/\mathfrak J$, où
$C=A[T_1,\ldots,T_n]$ et $\mathfrak J$ est un idéal de $C$. Alors
$\mathfrak J$ est limite inductive filtrante des idéaux de type fini
$\mathfrak J_\lambda\subset\mathfrak J$ de $C$~; posant
$C_\lambda=C/\mathfrak J_\lambda$, et utilisant l'exactitude du foncteur
$\varinjlim$, on voit que $B$ est encore isomorphe à la limite inductive du
système inductif filtrant $(C_\lambda)$. Il existe donc un $\lambda$ et un
$A$-homomorphisme $u:B\to C_\lambda$ tels que le composé
\[
B\xrightarrow{u}C_\lambda\xrightarrow{p_\lambda}B
\]
(où $p_\lambda$ est l'homomorphisme canonique) soit l'identité. Soit
$q_\lambda:C\to C_\lambda$ l'homomorphisme canonique, et posons
$t_i=p_\lambda(q_\lambda(T_i))$~; on a donc
$p_\lambda(u(t_i))=p_\lambda(q_\lambda(T_i))$, autrement dit
\[
u(t_i)-q_\lambda(T_i)\in\mathfrak J/\mathfrak J_\lambda.
\]
Il existe donc un $\mu\geqslant\lambda$ tel que les $n$ éléments
$u(t_i)-q_\lambda(T_i)$ appartiennent à
$\mathfrak J_\mu/\mathfrak J_\lambda$ $(1\leqslant i\leqslant n)$~; si
$p_{\mu\lambda}:C_\lambda\to C_\mu$ est l'homomorphisme canonique, on a par
suite
\[
p_{\mu\lambda}(u(t_i))
=p_{\mu\lambda}(q_\lambda(T_i))=q_\mu(T_i).
\]
Remplaçant $\lambda$ par $\mu$ et $u$ par $p_{\mu\lambda}\circ u$, on peut
donc supposer que $u(t_i)=q_\lambda(T_i)$ pour tout $i$, et si
$r=p_\lambda\circ q_\lambda$ est l'homomorphisme canonique
$C\to C/\mathfrak J=B$, on peut donc écrire
$u(r(T_i))=q_\lambda(T_i)$ pour tout $i$, d'où
$q_\lambda=u\circ r$. Mais cela entraîne nécessairement que
$\mathfrak J=\mathfrak J_\lambda$, car si $z\in\mathfrak J$, on a
$r(z)=0$~; donc on a $B=C_\lambda$.

Passons maintenant au cas où $S$ est affine et $X$ quelconque~; tout revient à
prouver qu'un ouvert affine $V$ de $X$ est de présentation finie sur $S$, et
en vertu de ce qui vient d'être démontré, il suffit de prouver que pour tout
système projectif filtrant $(Z_\lambda)$ de $S$-préschémas affines,
l'application
\begin{equation}
\label{IV.8.14.2.4-fr}
\varinjlim_\lambda\operatorname{Hom}_S(Z_\lambda,V)
\longrightarrow
\operatorname{Hom}_S(\varprojlim_\lambda Z_\lambda,V) \tag{8.14.2.4}
\end{equation}
est bijective. Il est immédiat que cette application est injective, car si
$(v_\lambda)$, $(v'_\lambda)$ sont deux systèmes inductifs de
$S$-homomorphismes $v_\lambda:Z_\lambda\to V$,
$v'_\lambda:Z_\lambda\to V$ tels que les morphismes correspondants
\[
Z\xrightarrow{u_\lambda}Z_\lambda\xrightarrow{v_\lambda}V,
\qquad
Z\xrightarrow{u_\lambda}Z_\lambda\xrightarrow{v'_\lambda}V
\]
soient égaux ($u_\lambda$ étant le morphisme canonique), alors les morphismes
\[
Z\xrightarrow{u_\lambda}Z_\lambda\xrightarrow{v_\lambda}V
\xrightarrow{j}X,
\qquad
Z\xrightarrow{u_\lambda}Z_\lambda\xrightarrow{v'_\lambda}V
\xrightarrow{j}X
\]
(où $j$ est l'injection canonique) sont égaux, ce qui entraîne
$j\circ v_\lambda=j\circ v'_\lambda$ par hypothèse pour un $\lambda$
convenable, donc $v_\lambda=v'_\lambda$.

\oldpage[IV-3]{54}
Reste à prouver que (8.14.2.4) est surjective. Soit donc $v:Z\to V$ un
$S$-morphisme~; par hypothèse il existe un $\lambda$ et un $S$-morphisme
$w_\lambda:Z_\lambda\to X$ tels que $j\circ v$ se factorise en
\[
Z\xrightarrow{u_\lambda}Z_\lambda\xrightarrow{w_\lambda}X,
\]
et tout revient à prouver qu'il existe $\mu\geqslant\lambda$ tel que le
morphisme
\[
Z_\mu\xrightarrow{u_{\lambda\mu}}Z_\lambda
\xrightarrow{w_\lambda}X
\]
(où $u_{\lambda\mu}$ est le morphisme de transition) se factorise en
\[
Z_\mu\xrightarrow{v_\mu}V\xrightarrow{j}X.
\]
Posons, pour tout $\mu\geqslant\lambda$,
$U_\mu=u_{\lambda\mu}^{-1}(w_\lambda^{-1}(V))$. On a
\[
u_\mu^{-1}(U_\mu)=u_\lambda^{-1}(U_\lambda)
=u_\lambda^{-1}(w_\lambda^{-1}(V))=v^{-1}(V)=Z.
\]
Comme les $Z_\lambda$ sont quasi-compacts et que les $U_\mu$, étant ouverts,
sont ind-constructibles (1.9.6), on déduit de (8.3.4) qu'il existe
$\mu\geqslant\lambda$ tel que $U_\mu=Z_\mu$. C.Q.F.D.

\begin{remark}[8.14.3]
\label{IV.8.14.3-fr}
Le fait que l'application (8.14.2.1) soit injective lorsque $f$ est localement
de type fini (8.8.2, (i)) amène naturellement à se demander si ce résultat
admet aussi une réciproque. Il n'en est rien, même lorsque $S$ et $X$ sont
affines, car il existe des monomorphismes $X\to S$ qui ne sont pas de type fini
(I, 2.4.2), et qui mettent donc en défaut cette conjecture.
\end{remark}

\section{Propriétés constructibles}
\label{section:IV.9-fr}

Soient $S$ un préschéma, $f:X\to S$ un morphisme de
\emph{présentation finie} (1.6.1), $\mathcal F$ un $\mathcal O_X$-Module
quasi-cohérent de \emph{présentation finie}. On se propose dans ce paragraphe
de donner des critères assurant, par exemple, que l'ensemble des $s\in S$ tels
que le $k(s)$-préschéma
$X_s=f^{-1}(s)=X\times_S\operatorname{Spec}(k(s))$ a une certaine propriété,
ou tels que le $\mathcal O_{X_s}$-Module
$\mathcal F_s=\mathcal F\otimes_{\mathcal O_X}k(s)$ a une certaine propriété,
est \emph{localement constructible}, ou tout au moins
\emph{ind-constructible} (1.9.4)~; on verra que c'est le cas pour la plupart
des propriétés qui s'introduisent en Géométrie algébrique. En supposant
seulement $f$ \emph{localement de présentation finie}, on donnera aussi (9.9)
des critères pour que l'ensemble des points $x\in X$ où la fibre $X_{f(x)}$
(ou le $\mathcal O_{X_{f(x)}}$-Module $\mathcal F_{f(x)}$) a une certaine
propriété, soit \emph{localement constructible}. On verra au \S~12 que ces
résultats, joints à l'hypothèse supplémentaire que $f$ est \emph{plat}
(resp. \emph{propre et plat}), permettent de prouver que les ensembles
considérés dans $X$ (resp. dans $S$) sont même \emph{ouverts} dans de nombreux
cas.

\subsection{Le principe de l'extension finie}
\label{subsection:IV.9.1-fr}

\begin{proposition}[9.1.1]
\label{IV.9.1.1-fr}
\textnormal{(Principe de l'extension finie).} Soient $k$ un corps,
$\mathcal C$ un ensemble d'extensions de $k$. On suppose vérifiées les
conditions suivantes~:
\begin{enumerate}[label=(\roman*)]
  \item Si $K\in\mathcal C$ et s'il existe un $k$-homomorphisme $K\to K'$
  (où $K'$ appartient à l'univers où l'on se place), alors
  $K'\in\mathcal C$.
  \item Si $K\in\mathcal C$, il existe une sous-extension $K'$ de $K$, de
  type fini sur $k$, telle que $K'\in\mathcal C$.
  \item Si $K\in\mathcal C$ est le corps de fractions d'une $k$-algèbre
  intègre de type fini $A$, il existe $f\in A-\{0\}$ tel que pour tout idéal
  maximal $\mathfrak m$ de $A$ on ait $A_f/\mathfrak m\in\mathcal C$.
\end{enumerate}
\oldpage[IV-3]{55}
Soit alors $\Omega$ une extension algébriquement close de $k$ (dans l'univers
considéré). Les conditions suivantes sont équivalentes~:
\begin{enumerate}[label=\emph{\alph*)}]
  \item $\mathcal C$ est non vide.
  \item Il existe $K'\in\mathcal C$ qui est une extension finie de $k$.
  \item On a $\Omega\in\mathcal C$.
\end{enumerate}
\end{proposition}

La condition (i) implique évidemment que \emph{b)} entraîne \emph{c)}, et
\emph{c)} entraîne trivialement \emph{a)}~; prouvons que \emph{a)} entraîne
\emph{b)}. En vertu de (ii) et (iii) il existe une extension
$K'\in\mathcal C$ qui est de la forme $A/\mathfrak m$, où $A$ est une
$k$-algèbre de type fini sur $k$ et $\mathfrak m$ un idéal maximal de $A$.
On sait, par le théorème des zéros de Hilbert (Bourbaki, \emph{Alg. comm.},
chap.~V, \S~3, n\textsuperscript{o}~1, cor.~2 du th.~1) que $K'$ est une
extension finie de $k$.

\begin{corollary}[9.1.2]
\label{IV.9.1.2-fr}
Sous les hypothèses (i), (ii), (iii) de (9.1.1), si $k$ est algébriquement
clos et si $\mathcal C$ n'est pas vide, $\mathcal C$ contient toutes les
extensions de $k$ dans l'univers envisagé.
\end{corollary}

En effet, comme \emph{a)} entraîne \emph{c)}, on a $k\in\mathcal C$, et la
conclusion résulte de l'hypothèse (i).

\begin{remark}[9.1.3]
\label{IV.9.1.3-fr}
En pratique, on vérifiera la condition (ii) de (9.1.1) en notant que $K$ est
limite inductive de ses sous-extensions de type fini $K_\alpha$, et en
appliquant les résultats du \S~8, compte tenu éventuellement du fait que pour
$K_\alpha\subset K_\beta$, $K_\beta$ est fidèlement plat sur $K_\alpha$.
Fréquemment l'ensemble $\mathcal C$ est formé des \emph{corps} appartenant à
un ensemble $\mathcal C'$ de $k$-algèbres qui vérifie la condition suivante~:
\begin{enumerate}[label=(i bis)]
  \item Si $A\in\mathcal C'$ et s'il existe un homomorphisme de $k$-algèbres
  $A\to A'$ (où $A'$ appartient à l'univers où l'on se place), alors
  $A'\in\mathcal C'$.
\end{enumerate}
Lorsqu'il en est ainsi, la condition (i) est trivialement satisfaite, et on
vérifiera généralement la condition (iii) de (9.1.1) en notant que le corps
des fractions $K$ de $A$ est la limite inductive des algèbres $A_f$ (pour la
relation $D(f)\supset D(g)$), et en appliquant les résultats du \S~8, compte
tenu éventuellement du fait que les morphismes $D(g)\to D(f)$ sont des
immersions ouvertes.

Par contre, lorsque (i bis) n'est pas vérifiée, la démonstration de (iii) est
souvent plus délicate, et est liée à des critères de constructibilité qui
seront développés plus loin.
\end{remark}

Voici des exemples typiques d'application du principe de l'extension finie~:

\begin{proposition}[9.1.4]
\label{IV.9.1.4-fr}
Soient $k$ un corps, $\Omega$ une extension algébriquement close de $k$,
$X$ et $Y$ deux préschémas de type fini sur $k$. Les conditions suivantes
sont équivalentes~:
\begin{enumerate}[label=\emph{\alph*)}]
  \item Il existe un $\Omega$-morphisme $X_{(\Omega)}\to Y_{(\Omega)}$
  (resp. un $\Omega$-morphisme possédant l'une (déterminée) des propriétés
  (i) à (xiv) de (8.10.5)).
  \item Il existe une extension finie $k'$ de $k$ et un $k'$-morphisme
  $X_{(k')}\to Y_{(k')}$ (resp. un $k'$-morphisme possédant la propriété
  considérée).
  \item Il existe une extension $K$ de $k$ et un $K$-morphisme
  $X_{(K)}\to Y_{(K)}$ (resp. un $K$-morphisme possédant la propriété
  considérée).
\end{enumerate}
\end{proposition}

On applique la remarque (9.1.3) en prenant pour $\mathcal C'$ l'ensemble de
toutes les $k$-algèbres $A$ (de l'univers où on est placé) telles qu'il existe
un $A$-morphisme $X\otimes_k A\to Y\otimes_k A$ (resp. un morphisme ayant
celle des propriétés de (8.10.5) que l'on considère). La condition (i bis) de
(9.1.3) est alors vérifiée grâce au fait que les

\oldpage[IV-3]{56}
propriétés envisagées dans (8.10.5) sont toutes stables par changement de base.
La condition (iii) de (9.1.1) est donc satisfaite en vertu de (8.8.2, (i))
(resp. (8.10.5)), puisque $\operatorname{Spec}(k)$ est quasi-compact et
quasi-séparé. Reste à vérifier la condition (ii) de (9.1.1) qui découle encore
de (8.8.2, (i)) et de (8.10.5), en considérant $K$ comme limite inductive de
ses sous-extensions de type fini. On conclut donc par (9.1.1).

En particulier, s'il existe une extension $K$ de $k$ et un $K$-isomorphisme
$X_{(K)}\xrightarrow{\sim}Y_{(K)}$, on dit que $X$ et $Y$ sont
\emph{géométriquement isomorphes}.

Le corollaire suivant généralise (II, 6.6.5)~:

\begin{corollary}[9.1.5]
\label{IV.9.1.5-fr}
Soient $k$ un corps, $X$ un $k$-préschéma. S'il existe une extension $K$ de
$k$ telle que $X_{(K)}$ soit projectif (resp. quasi-projectif) sur $K$, alors
$X$ est projectif (resp. quasi-projectif) sur $k$.
\end{corollary}

Le morphisme $\operatorname{Spec}(K)\to\operatorname{Spec}(k)$ étant
fidèlement plat et quasi-compact, il résulte déjà de (2.7.1, (v)) que $X$ est
de type fini sur $k$. L'hypothèse signifie qu'il existe une immersion fermée
(resp. une immersion) $X_{(K)}\to\mathbf P_K^r=\mathbf P_k^r\otimes_k K$
(II, 5.5.4, (ii) et 5.5.3)~; appliquant (9.1.4) pour la propriété (iv)
(resp. (ii)) de (8.10.5), on en déduit qu'il y a une extension finie $k'$ de
$k$ et une immersion fermée (resp. une immersion)
$X_{(k')}\to\mathbf P_{k'}^r$, autrement dit $X_{(k')}$ est projectif
(resp. quasi-projectif) sur $k'$. On conclut alors par (II, 6.6.5).

\begin{proposition}[9.1.6]
\label{IV.9.1.6-fr}
Soient $k$ un corps, $\Omega$ une extension algébriquement close de $k$,
$X$ un préschéma de type fini sur $k$, $\mathcal F$, $\mathcal G$ deux
$\mathcal O_X$-Modules cohérents. Supposons qu'il existe un isomorphisme
$\mathcal F\otimes_k\Omega\xrightarrow{\sim}\mathcal G\otimes_k\Omega$.
Alors il existe une extension finie $k'$ de $k$ et un isomorphisme
$\mathcal F\otimes_k k'\xrightarrow{\sim}\mathcal G\otimes_k k'$.
\end{proposition}

Le raisonnement est le même que dans (9.1.4), en appliquant (8.5.2, (i))
(on utilise ici, dans la démonstration de la propriété (iii) de (9.1.1), le
fait que les morphismes $D(g)\to D(f)$ (avec les notations de (9.1.3)) sont
des immersions ouvertes, et \emph{a fortiori} des morphismes plats).

\subsection{Propriétés constructibles et ind-constructibles}
\label{subsection:IV.9.2-fr}

\begin{definition}[9.2.1]
\label{IV.9.2.1-fr}
Soit $\mathbf P(X,\mathcal F,k)$ une relation. On dit que $\mathbf P$ est une
propriété \emph{constructible} (resp. \emph{ind-constructible}) de préschémas
algébriques si les deux conditions suivantes sont vérifiées~:
\begin{enumerate}[label=(\roman*)]
  \item Si $k$ est un corps, $X$ un préschéma algébrique sur $k$,
  $\mathcal F$ un $\mathcal O_X$-Module cohérent, $k'$ une extension de $k$,
  alors, pour que $\mathbf P(X,\mathcal F,k)$ soit vraie, il faut et il suffit
  que $\mathbf P(X_{(k')},\mathcal F\otimes_k k',k')$ soit vraie.
  \item Soient $S$ un préschéma intègre noethérien, de point générique $\eta$,
  $u:X\to S$ un morphisme de type fini, $\mathcal F$ un $\mathcal O_X$-Module
  cohérent. Pour tout $s\in S$, posons
  $X_s=u^{-1}(s)=X\times_S\operatorname{Spec}(k(s))$,
  $\mathcal F_s=\mathcal F\otimes_{\mathcal O_S}k(s)$. Soit $E$ l'ensemble
  des $s\in S$ tels que $\mathbf P(X_s,\mathcal F_s,k(s))$ soit vraie.
  Alors l'un des ensembles $E$, $S-E$ (resp. l'ensemble $E$) contient un
  ouvert non vide (et par suite est un voisinage de $\eta$) (resp. contient
  un ouvert non vide s'il contient $\eta$).
\end{enumerate}
\end{definition}

\oldpage[IV-3]{57}
\begin{remarks}[9.2.2]
\label{IV.9.2.2-fr}
\begin{enumerate}[label=(\roman*)]
  \item Il s'agit là bien entendu d'une convention de langage de nature
  métamathématique et non d'une définition mathématique proprement dite. On a
  des «~définitions~» analogues pour des relations entre $k$, un ou plusieurs
  $k$-préschémas algébriques, des $k$-morphismes entre ces préschémas, des
  Modules cohérents sur ces préschémas ou des homomorphismes entre ces Modules.

  \item Nous aurons aussi à considérer des relations où figurent des
  \emph{parties constructibles} de préschémas. Par exemple, soit
  $\mathbf P(X,Z,k)$ une relation~; nous dirons (par abus de langage) que
  $\mathbf P$ est une propriété constructible (resp. ind-constructible) de la
  partie constructible $Z$ de $X$ si les deux conditions suivantes sont
  vérifiées~:
  \begin{enumerate}[label=$\arabic*^{\mathrm o}$]
    \item Si $k$ est un corps, $X$ un préschéma algébrique sur $k$, $Z$ une
    \emph{partie constructible de $X$}, $k'$ une extension de $k$, alors, pour
    que $\mathbf P(X,Z,k)$ soit vraie, il faut et il suffit que
    $\mathbf P(X_{(k')},p^{-1}(Z),k')$ soit vraie
    ($p:X_{(k')}\to X$ étant la projection canonique).
    \item Soient $S$ un préschéma intègre noethérien, de point générique
    $\eta$, $u:X\to S$ un morphisme de type fini, $Z$ une
    \emph{partie constructible de $X$}. Pour tout $s\in S$, posons
    $X_s=u^{-1}(s)$, $Z_s=Z\cap X_s$. Soit $E$ l'ensemble des $s\in S$ tels
    que $\mathbf P(X_s,Z_s,k(s))$ soit vraie. Alors l'un des ensembles $E$,
    $S-E$ (resp. l'ensemble $E$) contient un ouvert non vide (resp. contient
    un ouvert non vide s'il contient $\eta$).
  \end{enumerate}

  On notera que dans la condition $2^{\mathrm o}$ il faut supposer que $Z$ est
  une partie constructible de $X$, et non seulement que $Z_s$ est une partie
  constructible de $X_s$ pour tout $s$~; la première de ces deux propriétés
  entraîne la seconde (1.8.2), mais non réciproquement.

  \item Si $\mathbf P$ est une propriété constructible, il est clair qu'il en
  est de même de «~non $\mathbf P$~». Si $\mathbf P$, $\mathbf Q$ sont deux
  propriétés constructibles (resp. ind-constructibles), il en est de même des
  propriétés «~$\mathbf P$ ou $\mathbf Q$~» et «~$\mathbf P$ et
  $\mathbf Q$~»~; en effet, si, sous les hypothèses de (9.2.1, (ii)), $E$,
  $E'$ sont deux parties de $S$ et si $E$ contient un ouvert non vide, il en
  est de même de $E\cup E'$, et si $S-E$ et $S-E'$ contiennent chacun un
  ouvert non vide, il en est de même de
  $S-(E\cup E')=(S-E)\cap(S-E')$.

  \item Soit $\mathbf P(X,\mathcal F,k)$ une relation vérifiant la condition
  (9.2.1, (i))~; soient $S$ un préschéma, $u:X\to S$ un morphisme de type
  fini~; avec les notations de (9.2.1, (ii)), soit $E$ l'ensemble des $s\in S$
  tels que $\mathbf P(X_s,\mathcal F_s,k(s))$ soit vraie. Soit d'autre part
  $g:S'\to S$ un morphisme quelconque, et posons
  $X'=X_{(S')}$, $\mathcal F'=\mathcal F\otimes_{\mathcal O_S}\mathcal O_{S'}$~;
  alors il résulte de la transitivité des fibres (I, 3.6.4) et de la condition
  (9.2.1, (i)) que \emph{l'ensemble des $s'\in S'$ tels que
  $\mathbf P(X'_{s'},\mathcal F'_{s'},k(s'))$ soit vraie est égal à
  $g^{-1}(E)$}. Cela s'étend aussitôt au cas où il figure dans $\mathbf P$
  plusieurs préschémas, Modules, morphismes de préschémas ou homomorphismes de
  Modules, ainsi qu'aux propriétés du type considéré dans (ii).

  \item Comme nous le verrons dans la suite de ce paragraphe et dans le courant
  du reste du chap.~IV, la plupart des propriétés $\mathbf P$ vérifiant la
  condition (9.2.1, (i)) vérifient aussi (9.2.1, (ii)). Comme exceptions
  possibles, notons la propriété d'être \emph{projectif}, ou
  \emph{quasi-projectif}, ou \emph{affine}, ou \emph{quasi-affine} sur le
  corps de base (pour un préschéma algébrique)~; nous verrons (9.6.2) que ces
  propriétés sont \emph{ind-constructibles}, mais nous donnerons plus tard un
  exemple où $S$ est une partie ouverte non vide de $\operatorname{Spec}(\mathbf Z)$
  (ou une partie ouverte d'une courbe elliptique sur un corps fini) et où
  toutes les fibres $X_s$

  \oldpage[IV-3]{58}
  sauf la fibre générique $X_\eta$ sont projectives sur $k(s)$ (tous les
  préschémas $X_s$ étant de dimension 2).
\end{enumerate}
\end{remarks}

\begin{proposition}[9.2.3]
\label{IV.9.2.3-fr}
Soient $\mathbf P$ une propriété constructible (resp. ind-constructible) de
préschémas algébriques, $S$ un préschéma, $X$ un préschéma de présentation
finie sur $S$, $\mathcal F$ un $\mathcal O_X$-Module quasi-cohérent de
présentation finie. Alors l'ensemble $E$ des $s\in S$ tels que
$\mathbf P(X_s,\mathcal F_s,k(s))$ soit vraie est localement constructible
(resp. ind-constructible). En outre, si $S$ est irréductible de point générique
$\eta$, un des deux ensembles $E$, $S-E$ est un voisinage de $\eta$ dans $S$
(resp. $E$ est un voisinage de $\eta$ s'il contient ce point).
\end{proposition}

Pour démontrer ces assertions, on peut se borner au cas où
$S=\operatorname{Spec}(A)$ est affine. On sait alors qu'il existe un
sous-anneau $A_0$ de $A$ qui est une $\mathbf Z$-algèbre de type fini, un
$A_0$-préschéma de type fini $X_0$ et un $\mathcal O_{X_0}$-Module cohérent
$\mathcal F_0$ tels que $X$ soit isomorphe à $X_0\otimes_{A_0}A$ et
$\mathcal F$ à $\mathcal F_0\otimes_{A_0}A$ (8.9.1). Soit
$p:S\to S_0=\operatorname{Spec}(A_0)$ le morphisme correspondant à
l'injection $A_0\to A$, et soit $E_0$ l'ensemble des $s_0\in S_0$ tels que
\[
\mathbf P((X_0)_{s_0},(\mathcal F_0)_{s_0},k(s_0))
\]
soit vraie~; alors, en vertu de (9.2.2, (iv)), on a $E=p^{-1}(E_0)$~; on peut
par suite (1.8.2) se borner au cas où $S$ est le spectre d'une
$\mathbf Z$-algèbre de type fini, donc un schéma noethérien. Utilisons le
critère de constructibilité $(\mathbf 0_{\mathrm{III}},9.2.3)$ (resp. le
critère d'ind-constructibilité (1.9.10))~; on est alors ramené, en utilisant
comme ci-dessus (9.2.2, (iv)) et en remplaçant $S$ par un sous-schéma fermé
intègre de $S$, au cas où $S$ est noethérien et intègre, et où il faut prouver
que $E$ est rare dans $S$ ou contient un ouvert non vide de $S$ (resp. que $E$
contient un ouvert non vide de $S$ s'il contient le point générique)~; mais
cela est garanti en vertu de la condition (9.2.1, (ii)).

On notera que l'on n'a utilisé (9.2.1, (ii)) que lorsque $S$ est le spectre
d'un anneau intègre de type fini sur $\mathbf Z$. Il est clair d'autre part
que l'énoncé de (9.2.3) s'applique aussi lorsqu'il figure dans $\mathbf P$
plusieurs préschémas, Modules sur ces préschémas, morphismes de préschémas ou
homomorphismes de Modules. Il s'applique encore lorsqu'il y figure des parties
(en nombre fini) des préschémas considérés, pourvu qu'on impose à ces parties
la condition d'être \emph{localement constructibles}. En effet, la restriction
au cas où $S$ est affine montre qu'on peut se borner au cas où ces parties sont
\emph{constructibles}~: on applique alors (8.3.11) qui montre (avec les
notations précédentes) qu'une partie constructible de $X$ est l'image
réciproque d'une partie constructible de $X_0$ pour un choix convenable de
$A_0$.

\begin{corollary}[9.2.4]
\label{IV.9.2.4-fr}
Soient $\mathbf P$ une propriété constructible (resp. ind-constructible) de
préschémas algébriques, $X$, $Y$ deux $S$-préschémas de présentation finie,
$f:X\to Y$ un $S$-morphisme. Pour tout $s\in S$, posons
$X_s=X\times_S\operatorname{Spec}(k(s))$,
$Y_s=Y\times_S\operatorname{Spec}(k(s))$,
$f_s=f\times1:X_s\to Y_s$. Alors l'ensemble $E$ des $s\in S$ tels que pour
tout $y\in Y_s$, la propriété $\mathbf P(f_s^{-1}(y),k(y))$ soit vraie, est
localement constructible (resp. ind-constructible).
\end{corollary}

En effet, soit $Z$ l'ensemble des $y\in Y$ tels que
$\mathbf P(f^{-1}(y),k(y))$ soit vraie. Comme les fibres $f^{-1}(y)$ et
$f_s^{-1}(y)$ sont isomorphes, on voit que $E$ est l'ensemble des $s\in S$
tels que $Y_s\subset Z$~; si $g:Y\to S$ est le morphisme structural, on a
donc $E=S-g(Y-Z)$.

\oldpage[IV-3]{59}
Or, $f$ est de présentation finie (1.6.2, (v)), donc il résulte de (9.2.3) que
$Z$ est localement constructible (resp. ind-constructible) dans $Y$, donc
$Y-Z$ est localement constructible (resp. pro-constructible) dans $Y$. Comme
$g$ est de présentation finie, $g(Y-Z)$ est localement constructible
(resp. pro-constructible) dans $S$, en vertu du théorème de Chevalley (1.8.4)
(resp. de (1.9.5, (vii)))~; donc $E$ est localement constructible
(resp. ind-constructible) dans $S$.

\begin{remark}[9.2.5]
\label{IV.9.2.5-fr}
On notera que si $\mathbf P$ est une propriété de préschémas algébriques pour
laquelle la prop. (9.2.3) est vraie, alors $\mathbf P$ satisfait aussi à la
condition (9.2.1, (ii))~: cela résulte en effet du fait que dans un espace
irréductible noethérien, un ensemble constructible est rare ou contient un
ouvert non vide $(\mathbf 0_{\mathrm{III}},9.2.3)$.
\end{remark}

\begin{proposition}[9.2.6]
\label{IV.9.2.6-fr}
Désignons par $\mathbf P$ une des propriétés suivantes d'un $k$-préschéma
algébrique $X$~:
\begin{enumerate}[label=(\roman*)]
  \item $X$ est vide.
  \item $X$ est fini sur $k$.
  \item $X$ est radiciel sur $k$.
  \item $\dim(X)$ appartient à une partie donnée $\Phi$ de l'ensemble
  $\overline{\mathbf Z}=\mathbf Z\cup\{-\infty\}$.
\end{enumerate}
Alors $\mathbf P$ est constructible.
\end{proposition}

Il est clair que (i) et (ii) sont des cas particuliers de (iv), en prenant
respectivement pour $\Phi$ l'ensemble $\{-\infty\}$ et l'ensemble
$\{-\infty,0\}$. On n'a donc qu'à démontrer (iii) et (iv). Dans chacun de ces
deux cas la condition (i) de (9.2.1) est remplie en vertu de (2.7.1, (xv)) et
(4.1.4). D'autre part, dans le cas (iii), la propriété $\mathbf P$ vérifie la
conclusion de (9.2.3) en vertu de (1.8.7)~; il reste donc à voir qu'il en est de
même dans le cas (iv). Cela va résulter de la proposition plus précise
suivante~:

\begin{proposition}[9.2.6.1]
\label{IV.9.2.6.1-fr}
Si $f:X\to S$ est un morphisme de présentation finie, la fonction
$s\mapsto\dim(f^{-1}(s))$ est localement constructible.
\end{proposition}

La question est locale sur $S$, donc on peut supposer que
$S=\operatorname{Spec}(A)$ est affine et prouver que pour tout $n$, l'ensemble
$E$ des $s\in S$ tels que $\dim(X_s)=n$ est constructible. Le même
raisonnement que dans (9.2.3) ramène au cas où $A$ est noethérien et intègre,
et il suffit alors de prouver le

\begin{corollary}[9.2.6.2]
\label{IV.9.2.6.2-fr}
Soient $S$ un préschéma intègre noethérien de point générique $\eta$,
$f:X\to S$ un morphisme de type fini. Alors il existe un voisinage de $\eta$
dans $S$ tel que la fonction $s\mapsto\dim(X_s)$ soit constante dans ce
voisinage.
\end{corollary}

Les images par $f$ des composantes irréductibles (en nombre fini) de $X$ qui
ne rencontrent pas $X_\eta$, sont contenues dans des parties fermées de $S$ ne
contenant pas $\eta$ (puisque $S$ est intègre
$(\mathbf 0_{\mathrm I},2.1.5)$), donc (en remplaçant $S$ par un voisinage
ouvert de $\eta$) on peut se borner au cas où toutes les composantes
irréductibles $X_i$ de $X$ rencontrent $X_\eta$~; désignons encore par $X_i$
le sous-préschéma fermé réduit de $X$ ayant $X_i$ pour espace sous-jacent~;
comme $\dim(X_s)=\sup_i(\dim((X_i)_s))$ (4.1.1), on peut se borner au cas où
$X$ est irréductible. Il existe alors un recouvrement fini $(U_j)$ de $X$ par
des

\oldpage[IV-3]{60}
ouverts affines partout denses, et les nombres $\dim((U_j)_\eta)$ sont tous
égaux à $n=\dim(X_\eta)$ (4.1.1.3)~; on peut par suite se borner au cas où $X$
est affine, donc aussi $X_\eta$. Il existe alors, en vertu de (4.1.2), un
ouvert non vide $W$ de $X$ tel que $W\supset X_\eta$, et un
$k(\eta)$-morphisme fini surjectif
$h:W_\eta\to\mathbf V_{k(\eta)}^n
(=\operatorname{Spec}(k(\eta)[T_1,\ldots,T_n]))$~; en appliquant
(8.8.2, (i)) et la méthode de (8.1.2, \emph{a)}), on en déduit (en remplaçant
au besoin $S$ par un voisinage de $\eta$) que l'on a $h=g_\eta$, où
$g:W\to\mathbf V_A^n(=\operatorname{Spec}(A[T_1,\ldots,T_n]))$ est un
$S$-morphisme, et on peut supposer ce morphisme fini et surjectif en vertu de
(8.10.5, (vi) et (x)). On en conclut que pour tout $s\in S$, le morphisme
$g_s:W_s\to\mathbf V_{k(s)}^n$ est fini et surjectif, donc que
$\dim(W_s)=n$ (4.1.2).

\subsection{Propriétés constructibles de morphismes de préschémas algébriques}

\begin{proposition}[9.3.1]
\label{IV.9.3.1-fr}
Soit $\mathbf P(X,k)$ une propriété constructible (resp. ind-constructible)
de préschémas algébriques. Désignons par $\mathbf P'(f,X,Y,k)$ la relation
suivante~: $f:X\to Y$ est un $k$-morphisme de $k$-préschémas algébriques tel
que pour tout $y\in Y$, on ait la propriété
$\mathbf P(f^{-1}(y),k(y))$. Alors $\mathbf P'$ est une propriété
constructible (resp. ind-constructible).
\end{proposition}

En effet, comme $\mathbf P$ vérifie la condition (9.2.1, (i)), il en est de
même de $\mathbf P'$ en vertu de la transitivité des fibres (I, 3.6.4)~; par
ailleurs le fait que $\mathbf P'$ vérifie la condition (9.2.1, (ii)) résulte
de (9.2.4), en raison de la remarque (9.2.5).

\begin{proposition}[9.3.2]
\label{IV.9.3.2-fr}
Désignons par $\mathbf P$ une des propriétés suivantes d'un $k$-morphisme
$f:X\to Y$ de $k$-préschémas algébriques~:
\begin{enumerate}[label=(\roman*)]
  \item $f$ est surjectif.
  \item $f$ est quasi-fini.
  \item $f$ est radiciel.
  \item Pour tout $y\in Y$, $\dim(f^{-1}(y))$ appartient à $\Phi$ (notation
  de (9.2.6)).
\end{enumerate}
Alors $\mathbf P$ est constructible.
\end{proposition}

Cela résulte aussitôt de (9.3.1) et (9.2.6) si l'on tient compte de ce que $f$
est de type fini (1.5.4, (v)), de la caractérisation des morphismes radiciels
(I, 3.5.8), et de celle des morphismes quasi-finis (II, 6.2.2).

\begin{proposition}[9.3.3]
\label{IV.9.3.3-fr}
Supposons vérifiées les hypothèses de (8.8.1) dont nous gardons les
notations~; supposons en outre que $S_\alpha$ soit quasi-compact, $X_\alpha$
et $Y_\alpha$ de présentation finie sur $S_\alpha$, et soit
$f_\alpha:X_\alpha\to Y_\alpha$ un $S_\alpha$-morphisme. Soit $\mathbf P$ une
propriété ind-constructible de morphismes de préschémas algébriques. Pour tout
$s\in S$ (resp. $s_\lambda\in S_\lambda$) posons
$X_s=X\times_S\operatorname{Spec}(k(s))$,
$Y_s=Y\times_S\operatorname{Spec}(k(s))$,
$f_s=f\times1:X_s\to Y_s$ (resp.
$X_{\lambda,s_\lambda}=X_\lambda\times_{S_\lambda}
\operatorname{Spec}(k(s_\lambda))$,
$Y_{\lambda,s_\lambda}=Y_\lambda\times_{S_\lambda}
\operatorname{Spec}(k(s_\lambda))$,
$f_{\lambda,s_\lambda}=f_\lambda\times1:
X_{\lambda,s_\lambda}\to Y_{\lambda,s_\lambda}$). Alors, afin que pour
tout $s\in S$ on ait la propriété $\mathbf P(f_s)$, il faut et il suffit
qu'il existe $\lambda\geqslant\alpha$ tel que pour tout
$s_\lambda\in S_\lambda$, on ait $\mathbf P(f_{\lambda,s_\lambda})$.
\end{proposition}

En effet, soit $E$ (resp. $E_\lambda$) l'ensemble des $s\in S$ (resp.
$s_\lambda\in S_\lambda$) tels que la propriété $\mathbf P(f_s)$ (resp.
$\mathbf P(f_{\lambda,s_\lambda})$) soit vraie~; il résulte de
(9.2.2, (iv)) que l'on a
$E_\mu=u_{\lambda\mu}^{-1}(E_\lambda)$ pour $\lambda\leqslant\mu$, et
$E=u_\lambda^{-1}(E_\lambda)$~; en outre, en vertu de (9.2.3), $E$ (resp.
$E_\lambda$) est ind-constructible dans $S$ (resp. $S_\lambda$)~; la
proposition résulte donc de (8.3.4) appliqué à la partie ind-constructible
$E$ de $S$.

\oldpage[IV-3]{61}
On généralise sans peine ce résultat à des propriétés $\mathbf P$ du type
considéré dans (9.2.3, (i) et (ii)).

\begin{remark}[9.3.4]
\label{IV.9.3.4-fr}
La conjonction de (9.3.3) et de (9.3.2, (ii)) démontre l'assertion
(8.10.5, (xi)).
\end{remark}

\begin{proposition}[9.3.5]
\label{IV.9.3.5-fr}
Soit $\mathbf P(X,Y,k)$ la propriété~: «~$X$ et $Y$ sont deux préschémas de
type fini sur le corps $k$, et il existe une extension $k'$ de $k$ et un
$k'$-morphisme $g:X_{(k')}\to Y_{(k')}$ vérifiant $\mathbf Q(g)$~», où
$\mathbf Q$ est une des propriétés (i) à (xiv) de (8.10.5). Alors $\mathbf P$
est une propriété ind-constructible.
\end{proposition}

La définition de $\mathbf P$ montre en effet que la condition (9.2.1, (i))
est satisfaite, compte tenu de ce que la propriété $\mathbf Q(g)$ est stable
par changement du corps de base, et de ce que deux extensions de $k$ peuvent
toujours être considérées comme des sous-extensions d'une troisième extension
de $k$. Pour vérifier (9.2.1, (ii)), on peut se borner au cas où
$S=\operatorname{Spec}(A)$ est affine~; si $K=k(\eta)$, corps des fractions de
$A$, il existe par hypothèse et en vertu de (9.1.4) une extension finie $K'$ de
$K$ et un $K'$-morphisme
$g':(X_\eta)_{(K')}\to(Y_\eta)_{(K')}$, vérifiant $\mathbf Q(g')$, et $K'$ est
évidemment le corps des fractions d'une $A$-algèbre intègre finie $A'$. Si
l'on pose $S'=\operatorname{Spec}(A')$, il résulte alors de (8.10.5) qu'il y
a un voisinage $U'$ du point générique $\eta'$ de $S'$ tel que, si l'on pose
$X'=X\otimes_A A'$, $Y'=Y\otimes_A A'$, il existe, pour tout $s'\in U'$, un
morphisme $X'_{s'}\to Y'_{s'}$ ayant la propriété $\mathbf Q$. Mais le
morphisme $h:S'\to S$ est fini, donc fermé, et comme
$h^{-1}(\eta)=\{\eta'\}$, $h(U')$ contient un voisinage ouvert $U$ de $\eta$
dans $S$~; pour tout $s\in U$, il y a donc un $s'\in U'$ tel que $h(s')=s$,
et comme
$X'_{s'}=X_s\otimes_{k(s)}k(s')$,
$Y'_{s'}=Y_s\otimes_{k(s)}k(s')$, la propriété
$\mathbf P(X_s,Y_s,k(s))$ est vraie pour tout $s\in U$.

\begin{env}[9.3.6]
\label{IV.9.3.6-fr}
\textbf{Exemple.} Prenons par exemple pour $\mathbf Q$ la propriété d'être un
isomorphisme. Alors, en conjuguant (9.3.5) et (9.3.3), on a la propriété
suivante~: les notations et hypothèses étant celles de (8.8.1), $S_\alpha$
étant quasi-compact, $X_\alpha$ et $Y_\alpha$ de présentation finie sur
$S_\alpha$, afin que, pour tout $s\in S$, $X_s$ et $Y_s$ soient
géométriquement isomorphes (9.1.4), il faut et il suffit qu'il existe
$\lambda\geqslant\alpha$ tel que, pour tout $s_\lambda\in S_\lambda$,
$X_{\lambda,s_\lambda}$ et $Y_{\lambda,s_\lambda}$ soient géométriquement
isomorphes.
\end{env}

On a un résultat analogue lorsque les préschémas que l'on considère sont munis
de «~lois de composition~» d'une certaine espèce
$(\mathbf 0_{\mathrm{III}},8.2.1)$, par exemple des «~préschémas en groupes~»,
«~préschémas en anneaux~», etc. $(\mathbf 0_{\mathrm{III}},8.2.3)$. Alors
l'énoncé précédent est encore valable lorsque par «~isomorphisme~» on entend
des isomorphismes de préschémas qui sont des homomorphismes pour les lois de
composition considérées $(\mathbf 0_{\mathrm{III}},8.2.2)$~; il suffit ici
d'utiliser non seulement (8.10.5) mais aussi (8.8.2, (i)) en remarquant que la
notion d'homomorphisme pour une loi de composition s'exprime en écrivant que
des diagrammes de morphismes de préschémas sont commutatifs (il faut bien
entendu que les morphismes de transition $X_\mu\to X_\lambda$ et
$Y_\mu\to Y_\lambda$ pour $\lambda\leqslant\mu$ soient des homomorphismes
pour les lois de composition envisagées).

On peut aussi, au lieu de considérer des morphismes de préschémas comme dans
(9.3.5), considérer des homomorphismes de Modules, en utilisant (9.1.6) au
lieu de (9.1.5).

\oldpage[IV-3]{62}

\subsection{Constructibilité de certaines propriétés des modules}
\label{subsection:IV.9.4-fr}

\begin{env}[9.4.1]
\label{IV.9.4.1-fr}
\textbf{Notations.} Dans ce numéro et les suivants jusqu'à la fin du \S~9,
nous utiliserons systématiquement les notations suivantes~: étant donné un
morphisme $f:X\to S$, nous poserons, pour tout $s\in S$,
$X_s=f^{-1}(s)=X\times_S\operatorname{Spec}(k(s))$~; pour tout
$\mathcal O_X$-Module quasi-cohérent $\mathcal F$, $\mathcal F_s$ désignera
le $\mathcal O_{X_s}$-Module
$\mathcal F\otimes_{\mathcal O_S}k(s)$, et pour tout homomorphisme
$u:\mathcal F\to\mathcal G$ de $\mathcal F$ dans un $\mathcal O_X$-Module
quasi-cohérent $\mathcal G$, $u_s:\mathcal F_s\to\mathcal G_s$ sera le
morphisme $p^*(u)$, en désignant par $p$ la projection canonique $X_s\to X$.
Pour toute section $g$ de $\mathcal F$ au-dessus de $X$, on désignera par
$g_s$ l'image de $g$ par l'homomorphisme canonique
$\Gamma(X,\mathcal F)\to\Gamma(X_s,\mathcal F_s)$. Pour toute partie $Z$ de
$X$, on notera $Z_s$ l'image réciproque
$p^{-1}(Z)=Z\cap X_s$ (\textbf{I}, 3.6.1). Enfin, si $Y$ est un second
$S$-préschéma et $h:X\to Y$ un $S$-morphisme, on notera $h_s$ le morphisme
$h\times1:X_s\to Y_s$.
\end{env}

\begin{proposition}[9.4.2]
\label{IV.9.4.2-fr}
Soient $S$ un préschéma intègre de point générique $\eta$, $f:X\to S$ un
morphisme de présentation finie, $\mathcal F$, $\mathcal G$, $\mathcal H$
trois $\mathcal O_X$-Modules quasi-cohérents de présentation finie. Soient
$u:\mathcal F\to\mathcal G$, $v:\mathcal G\to\mathcal H$ deux
homomorphismes de $\mathcal O_X$-Modules, et supposons que la suite
$\mathcal F_\eta\xrightarrow{u_\eta}\mathcal G_\eta
\xrightarrow{v_\eta}\mathcal H_\eta$ soit exacte. Alors il existe un
voisinage ouvert $U$ de $\eta$ dans $S$ tel que, pour tout $s\in U$, la suite
$\mathcal F_s\xrightarrow{u_s}\mathcal G_s
\xrightarrow{v_s}\mathcal H_s$ soit exacte.
\end{proposition}

Avec les notations générales de (9.2.1), il s'agit ici de la relation
$\mathbf P(X,\mathcal F,\mathcal G,\mathcal H,u,v,k)$~: «~$X$ est un
préschéma algébrique sur le corps $k$,
$\mathcal F\xrightarrow{u}\mathcal G\xrightarrow{v}\mathcal H$ une suite
exacte de $\mathcal O_X$-Modules quasi-cohérents~». Comme, pour toute extension
$k'$ de $k$, la projection canonique $X_{(k')}\to X$ est un morphisme
fidèlement plat (2.2.13), la condition (9.2.1, (i)) est satisfaite (2.2.7). En
vertu de (9.2.3), on peut se borner au cas où $S$ est intègre et
\emph{noethérien}, auquel cas $X$ est noethérien, et $\mathcal F$,
$\mathcal G$, $\mathcal H$ sont des $\mathcal O_X$-Modules cohérents.
L'hypothèse implique qu'il existe un voisinage ouvert $U$ de $\eta$ dans $S$
tel que la suite
$\mathcal F|f^{-1}(U)\to\mathcal G|f^{-1}(U)\to\mathcal H|f^{-1}(U)$ soit
\emph{exacte}, en vertu de (8.5.8, (ii)) appliqué suivant la méthode générale
de (8.1.2, \emph{a)}), et l'on peut donc déjà supposer que la suite
$\mathcal F\xrightarrow{u}\mathcal G\xrightarrow{v}\mathcal H$ est exacte~;
il suffit évidemment de prouver que l'on a
$\operatorname{Ker}(v_s)=(\operatorname{Ker}v)_s$ et
$\operatorname{Im}(u_s)=(\operatorname{Im}u)_s$ pour tout $s$ voisin de
$\eta$ dans $S$~; par suite (tenant compte de ce que les
$\mathcal O_X$-Modules $\mathcal F$, $\mathcal G$, $\mathcal H$ sont
\emph{cohérents} et de $(\mathbf 0_{\mathrm I},5.3.4)$) on est ramené à
prouver la proposition dans le cas particulier où la suite
\[
0\to\mathcal F\xrightarrow{u}\mathcal G\xrightarrow{v}\mathcal H\to0
\]
est exacte. Or, il y a alors un ouvert $U$ dans $S$ contenant $\eta$ et tel
que $\mathcal H|f^{-1}(U)$ soit \emph{plat} sur $U$ (6.9.1)~; il résulte donc
de (2.1.8) que pour tout $s\in U$, la suite
$0\to\mathcal F_s\to\mathcal G_s\to\mathcal H_s\to0$ est exacte, ce qui
achève la démonstration.

\begin{corollary}[9.4.3]
\label{IV.9.4.3-fr}
Soient $S$ un préschéma intègre, de point générique $\eta$, $f:X\to S$ un
morphisme de présentation finie. Soit
$\mathcal L^\bullet=(\mathcal L^i)_{i\in\mathbf Z}$ un complexe de
$\mathcal O_X$-Modules quasi-cohérents de présentation finie. Pour tout $i$,
il existe un voisinage ouvert $U$ de $\eta$ dans $S$ tel que les homomorphismes
canoniques
\begin{equation}
\label{IV.9.4.3.1-fr}
(\mathcal H^i(\mathcal L^\bullet))_s\to
\mathcal H^i(\mathcal L_s^\bullet) \tag{9.4.3.1}
\end{equation}
soient bijectifs pour tout $s\in U$.
\end{corollary}

\oldpage[IV-3]{63}

On peut évidemment se borner à un complexe à trois termes de degrés
$-1$, $0$, $+1$~:
\[
\mathcal M\xrightarrow{u}\mathcal N\xrightarrow{v}\mathcal P
\]
et à $i=0$~; l'homomorphisme à considérer est alors l'homomorphisme canonique
$(\operatorname{Ker}v/\operatorname{Im}u)_s\to
\operatorname{Ker}(v_s)/\operatorname{Im}(u_s)$. Utilisant (8.9.1) et
(8.5.2, (i)), on voit qu'on peut se ramener au cas où $S$ (donc aussi $X$)
est \emph{noethérien}, et par suite $\mathcal M$, $\mathcal N$, $\mathcal P$
des $\mathcal O_X$-Modules \emph{cohérents}~; alors
$\operatorname{Im}(u)$ et $\operatorname{Ker}(v)$ sont aussi cohérents
$(\mathbf 0_{\mathrm I},5.3.4)$ et en outre il existe un voisinage $U$ de
$\eta$ tel que pour $s\in U$, on ait
$\operatorname{Ker}(v_s)=(\operatorname{Ker}(v))_s$ et
$\operatorname{Im}(u_s)=(\operatorname{Im}(u))_s$ (9.4.2)~; la conclusion
résulte alors de (9.4.2) appliqué à la suite exacte
\[
0\to\operatorname{Im}(u)\to\operatorname{Ker}(v)\to
\operatorname{Ker}(v)/\operatorname{Im}(u)\to0,
\]
compte tenu de ce que $\mathcal O_\eta=k(\eta)$ (puisque $S$ est intègre) et
par conséquent la suite
\[
0\to(\operatorname{Im}u)_\eta\to(\operatorname{Ker}v)_\eta\to
(\operatorname{Ker}v/\operatorname{Im}u)_\eta\to0
\]
est exacte.

\begin{proposition}[9.4.4]
\label{IV.9.4.4-fr}
Soient $f:X\to S$ un morphisme de présentation finie, $\mathcal F$,
$\mathcal G$, $\mathcal H$ trois $\mathcal O_X$-Modules quasi-cohérents de
présentation finie. Soient $u:\mathcal F\to\mathcal G$,
$v:\mathcal G\to\mathcal H$ deux homomorphismes de $\mathcal O_X$-Modules.
Alors l'ensemble $E$ des $s\in S$ tels que la suite
$\mathcal F_s\xrightarrow{u_s}\mathcal G_s
\xrightarrow{v_s}\mathcal H_s$ soit exacte est localement constructible.
\end{proposition}

Compte tenu de (9.2.3), il faut établir que la propriété $\mathbf P$ considérée
dans (9.4.2) est constructible. On a déjà remarqué dans la démonstration de
(9.4.2) que la condition (9.2.1, (i)) est satisfaite pour cette propriété, et
il reste à vérifier la condition (9.2.1, (ii)). Supposons donc $S$ intègre
noethérien, de point générique $\eta$, et prouvons que $E$ ou $S-E$ est un
voisinage de $\eta$. Si $\eta\in E$, notre assertion résulte de (9.4.2), et
l'on peut donc se borner au cas où $\eta\notin E$, autrement dit, la suite
$\mathcal F_\eta\xrightarrow{u_\eta}\mathcal G_\eta
\xrightarrow{v_\eta}\mathcal H_\eta$ n'est pas exacte. Distinguons alors
deux cas~:

$1^\circ$ Posons $w=v\circ u$, et supposons d'abord que
$w_\eta=v_\eta\circ u_\eta\neq0$. Comme $\mathcal F$, $\mathcal G$,
$\mathcal H$ sont cohérents, il en est de même de
$\mathcal N=\operatorname{Ker}(w)$ $(\mathbf 0_{\mathrm I},5.3.4)$~; il
résulte donc de (9.4.2) appliqué à la suite exacte
$0\to\mathcal N\to\mathcal F\xrightarrow{w}\mathcal H$ qu'il y a un
voisinage $U$ de $\eta$ dans $S$ tel que, pour $s\in S$,
$\operatorname{Ker}(w_s)=\mathcal N_s$~; en restreignant $S$, on peut donc
supposer cette relation vérifiée pour tout $s\in S$. Soit $j$ l'injection
canonique $\mathcal N\to\mathcal F$, et posons
$\mathcal M=\operatorname{Coker}(j)$~; l'exactitude à droite du foncteur
$\mathcal F\mapsto\mathcal F_s$ entraîne que
$\mathcal M_s=\operatorname{Coker}(j_s)$ pour tout $s\in S$. L'hypothèse
$w_\eta\neq0$ signifie que $\mathcal M_\eta\neq0$~; comme $\mathcal M$ est
cohérent $(\mathbf 0_{\mathrm I},5.3.4)$, il résulte de (1.8.6) qu'il y a un
voisinage ouvert $U$ de $\eta$ dans $S$ tel que $\mathcal M_s\neq0$ pour
$s\in U$, donc $w_s\neq0$ pour $s\in U$, et \emph{a fortiori} $S-E$ est un
voisinage de $\eta$.

$2^\circ$ Supposons que $w_\eta=0$~; en vertu de (8.5.2, (i)), appliqué
suivant la méthode générale de (8.1.2, \emph{a)}), il existe un voisinage
ouvert $U$ de $\eta$ tel que $w|f^{-1}(U)=0$~; remplaçant $S$ par $U$, on peut
déjà supposer $w=0$ dans $X$. Alors
$\mathcal F\xrightarrow{u}\mathcal G\xrightarrow{v}\mathcal H$ est un
complexe à trois termes $\mathcal L^\bullet$, auquel on peut appliquer
(9.4.3)~; par hypothèse on a
$(\mathcal H^0(\mathcal L^\bullet))_\eta=
\mathcal H^0(\mathcal L^\bullet_\eta)\neq0$, et
$\mathcal H^0(\mathcal L)$ est cohérent
$(\mathbf 0_{\mathrm I},5.3.4\text{ et }5.3.3)$, donc il résulte de (1.8.6)
qu'il y a un voisinage ouvert $U$ de $\eta$ tel que
$(\mathcal H^0(\mathcal L^\bullet))_s\neq0$ pour tout $s\in U$~; mais comme
on peut supposer que
$\mathcal H^0(\mathcal L_s^\bullet)=(\mathcal H^0(\mathcal L^\bullet))_s$
pour $s\in U$ par (9.4.3), on voit de nouveau que $S-E$ est un voisinage de
$\eta$ dans $S$.

\oldpage[IV-3]{64}

\begin{corollary}[9.4.5]
\label{IV.9.4.5-fr}
Soient $f:X\to S$ un morphisme de présentation finie, $\mathcal F$,
$\mathcal G$ deux $\mathcal O_X$-Modules quasi-cohérents de présentation finie,
$u:\mathcal F\to\mathcal G$ un homomorphisme de $\mathcal O_X$-Modules.
Alors l'ensemble des $s\in S$ tels que $u_s$ soit injectif
\textnormal{(}resp. surjectif, bijectif, nul\textnormal{)} est localement
constructible.
\end{corollary}

Il suffit d'appliquer (9.4.4) aux suites
$0\to\mathcal F\xrightarrow{u}\mathcal G$,
$\mathcal F\xrightarrow{u}\mathcal G\to0$,
$0\to\mathcal F\xrightarrow{u}\mathcal G\to0$,
$\mathcal F\xrightarrow{u}\mathcal G\xrightarrow{1_\mathcal G}\mathcal G$.

\begin{corollary}[9.4.6]
\label{IV.9.4.6-fr}
Soient $f:X\to S$ un morphisme de présentation finie, $\mathcal F$ un
$\mathcal O_X$-Module quasi-cohérent de présentation finie. Soit $h$ une
section de $\mathcal F$ au-dessus de $X$~; pour tout $s\in S$, soit $h_s$ la
section correspondante de $\mathcal F_s$ au-dessus de $X_s$ \textnormal{(}pour
le morphisme projection $X_s\to X$\textnormal{)}. Alors, l'ensemble des
$s\in S$ tels que $h_s=0$ est localement constructible.
\end{corollary}

Il suffit de remarquer que $h$ correspond à un homomorphisme
$u:\mathcal O_X\to\mathcal F$ $(\mathbf 0_{\mathrm I},5.1.1)$ et $h_s$ à
l'homomorphisme $u_s$.

\begin{proposition}[9.4.7]
\label{IV.9.4.7-fr}
Soient $f:X\to S$ un morphisme de présentation finie, $\mathcal F$ un
$\mathcal O_X$-Module quasi-cohérent de présentation finie. L'ensemble $E$
\textnormal{(}resp. $E'$\textnormal{)} des $s\in S$ tels que $\mathcal F_s$
soit un $\mathcal O_{X_s}$-Module localement libre \textnormal{(}resp.
localement libre de rang $n$\textnormal{)} est localement constructible.
\end{proposition}

Si $X$ est un préschéma algébrique sur un corps $k$, $\mathcal F$ un
$\mathcal O_X$-Module cohérent, $k'$ une extension de $k$, alors, pour que
$\mathcal F$ soit localement libre (resp. localement libre de rang $n$), il
faut et il suffit qu'il en soit de même de $\mathcal F\otimes_k k'$, puisque la
projection $X_{(k')}\to X$ est un morphisme fidèlement plat (2.2.7). Autrement
dit, la condition (9.2.1, (i)) est vérifiée pour les propriétés dont on veut
démontrer la constructibilité, et il reste à vérifier (9.2.1, (ii))~; on peut
donc encore supposer que $S$ est affine, noethérien et intègre. Il y a de
nouveau quatre cas à envisager~:

$1^\circ$ $\eta\in E$. Il résulte de (8.5.5), appliqué suivant la méthode
générale de (8.1.2, \emph{a)}), qu'il existe un voisinage ouvert $U$ de $\eta$
dans $S$ tel que $\mathcal F|f^{-1}(U)$ soit localement libre~;
\emph{a fortiori} $\mathcal F_s$ est localement libre pour tout $s\in U$.

$2^\circ$ $\eta\in E'$. Même raisonnement que dans le $1^\circ$.

$3^\circ$ $\eta\in S-E$. Comme $\mathcal F_\eta$ est un
$\mathcal O_{X_\eta}$-Module cohérent, dire qu'il n'est pas localement libre
équivaut à dire qu'il \emph{n'est pas plat} sur $\mathcal O_{X_\eta}$
(Bourbaki, \emph{Alg. comm.}, chap.~II, \S~5, n\textsuperscript{o}~2,
cor.~2 du th.~1). Le fait que $S-E$ est un voisinage de $\eta$ résultera donc
du lemme plus général suivant (appliqué au cas où $g$ est l'identité)~:

\begin{lemma}[9.4.7.1]
\label{IV.9.4.7.1-fr}
Soient $S$ un préschéma intègre noethérien, de point générique $\eta$, $X$,
$Y$ deux $S$-préschémas de type fini sur $S$, $g:X\to Y$ un $S$-morphisme,
$\mathcal F$ un $\mathcal O_X$-Module cohérent. Si $\mathcal F_\eta$ n'est
pas $g_\eta$-plat, alors il existe un voisinage ouvert $U$ de $\eta$ dans $S$
tel que pour tout $s\in U$, $\mathcal F_s$ ne soit pas $g_s$-plat.
\end{lemma}

Compte tenu de (2.1.2) et de Bourbaki, \emph{Alg. comm.}, chap.~I\textsuperscript{er},
\S~2, n\textsuperscript{o}~3, Remarque~1, l'hypothèse signifie qu'il existe un
ouvert non vide $V$ de $Y_\eta$, et un homomorphisme injectif
$v:\mathcal M\to\mathcal N$ de $\mathcal O_V$-Modules \emph{cohérents}, tel
que l'homomorphisme
$1\otimes v:\mathcal F_\eta\otimes_{\mathcal O_V}\mathcal M\to
\mathcal F_\eta\otimes_{\mathcal O_V}\mathcal N$ ne soit pas injectif. On a
$V=Y_\eta\cap W$, où $W$ est ouvert dans $Y$ (\textbf{I}, 3.6.1), et il
résulte de (8.5.2, (i) et (ii)), appliqué suivant la méthode de
(8.1.2, \emph{a)}), qu'il existe un voisinage ouvert $U_0$ de $\eta$ dans $S$,
deux $\mathcal O_Z$-Modules

\oldpage[IV-3]{65}

cohérents $\mathcal G$, $\mathcal H$ (où $Z=W\cap h^{-1}(U_0)$,
$h:Y\to S$ étant le morphisme structural) et un $\mathcal O_Z$-homomorphisme
$u:\mathcal G\to\mathcal H$ tel que $\mathcal M=\mathcal G_\eta$,
$\mathcal N=\mathcal H_\eta$ et $v=u_\eta$~; on peut donc supposer $U_0$ pris
tel que, pour $s\in U_0$, $u_s:\mathcal G_s\to\mathcal H_s$ soit injectif
(9.4.5). Mais pour tout $s\in U_0$, l'homomorphisme
$1\otimes u_s:\mathcal F_s\otimes_{\mathcal O_{Z_s}}\mathcal G_s\to
\mathcal F_s\otimes_{\mathcal O_{Z_s}}\mathcal H_s$ n'est autre que
$(1\otimes u)_s$~; l'hypothèse que $(1\otimes u)_\eta$ est non injectif
entraîne donc (9.4.5) l'existence d'un ouvert non vide $U\subset U_0$ tel que
pour tout $s\in U$, $(1\otimes u)_s$ soit non injectif, et par suite
$\mathcal F_s$ n'est pas $g_s$-plat pour tout $s\in U$.

$4^\circ$ $\eta\in S-E'$. Il est clair que $S-E\subset S-E'$, et si
$\eta\in S-E$, $S-E'$ est \emph{a fortiori} un voisinage de $\eta$ d'après le
$3^\circ$. Supposons donc $\eta\in E$, donc $\mathcal F_\eta$ localement
libre~; ces hypothèses entraînent que $X_\eta$ est non connexe, et que les
rangs du $\mathcal O_{X_\eta}$-Module localement libre $\mathcal F_\eta$ ne
sont pas les mêmes sur les diverses composantes connexes de $X_\eta$. Or, il
résulte de (8.4.2), appliqué suivant la méthode de (8.1.2, \emph{a)}), que l'on
peut supposer (en remplaçant $S$ par un voisinage ouvert de $\eta$) que $X$ et
$X_\eta$ ont le même nombre de composantes connexes, les composantes connexes
de $X_\eta$ étant les intersections de $X_\eta$ et des composantes connexes de
$X$. La conclusion résulte alors du raisonnement fait dans le $2^\circ$,
appliqué à chacune des composantes connexes de $X$ (qui sont en nombre fini).

\begin{remark}[9.4.7.2]
\label{IV.9.4.7.2-fr}
Le lemme (9.4.7.1) sera plus tard généralisé et débarrassé d'hypothèses
noethériennes (11.2.8).
\end{remark}

\begin{proposition}[9.4.8]
\label{IV.9.4.8-fr}
Soient $S$ un préschéma localement noethérien, $f:X\to S$ un morphisme de type
fini, $\mathcal F$ un $\mathcal O_X$-Module cohérent. On suppose que pour tout
$s\in S$, $X_s$ soit un préschéma localement intègre. Alors l'ensemble $E$
\textnormal{(}resp. $E'$\textnormal{)} des $s\in S$ tels que $\mathcal F_s$
soit un $\mathcal O_{X_s}$-Module de torsion \textnormal{(}resp. un
$\mathcal O_{X_s}$-Module sans torsion\textnormal{)} est localement
constructible.
\end{proposition}

On peut évidemment supposer $S=\operatorname{Spec}(A)$ affine et noethérien et
prouver que $E$ (resp. $E'$) est alors constructible en utilisant le critère
$(\mathbf 0_{\mathrm{III}},9.2.3)$~; remplaçant $S$ par le sous-préschéma fermé
réduit de $S$ ayant pour espace sous-jacent une partie fermée irréductible $Y$
de $S$, et notant (\textbf{I}, 3.6.4) que pour $s\in Y$, la fibre
$(X_{(Y)})_s$ s'identifie canoniquement à $X_s$ et le faisceau
$(\mathcal F\otimes_{\mathcal O_S}\mathcal O_Y)_s$ à $\mathcal F_s$, on voit
qu'on est ramené au cas où $S$ est intègre de point générique $\eta$, et à
prouver que $E$ ou $S-E$ (resp. $E'$ ou $S-E'$) est un voisinage de $\eta$
dans $S$. Notons en outre que $X$ est réunion finie d'ouverts affines $U_i$,
de type fini sur $S$, et chacun des $(U_i)_s$ est induit sur un ouvert de
$X_s$, donc localement intègre~; en outre, si $(U_i)_\eta$ est vide (resp. non
vide), on sait que $(U_i)_s$ est aussi vide (resp. non vide) dans un voisinage
de $\eta$ (9.2.6). On peut donc supposer tous les $(U_i)_\eta$ non vides et
intègres, et dire que $\mathcal F_s$ est de torsion (resp. sans torsion)
équivaut à dire que chacun des $\mathcal F_s|(U_i)_s$ est de torsion (resp.
sans torsion). On est donc ramené au cas où $X=\operatorname{Spec}(B)$ est
affine, et $\mathcal F=\widetilde M$, où $M$ est un $B$-module de type fini~;
on posera $B_s=B\otimes_A k(s)$, $M_s=M\otimes_A k(s)$ et l'on peut supposer
$B_\eta$ intègre. Nous avons quatre cas à envisager~:

$1^\circ$ $\eta\in E$~; $M_\eta$ est donc un $B_\eta$-module de torsion de
type fini, et il y a par suite un $h\neq0$ dans $B_\eta$ tel que
$hM_\eta=0$~; en vertu de (8.5.2, (i)), appliqué suivant la méthode

\oldpage[IV-3]{66}

de (8.1.2, \emph{a)}), on peut (en remplaçant au besoin $S$ par un voisinage
de $\eta$) supposer que $h\in\Gamma(X_\eta,\mathcal O_{X_\eta})$ est de la
forme $g_\eta$, où $g\in\Gamma(X,\mathcal O_X)$~; soit
$u:\mathcal F\to\mathcal F$ l'endomorphisme de $\mathcal F$ défini par la
multiplication par $g$~; par hypothèse, on a $u_\eta=0$, donc (9.4.5)
l'endomorphisme $u_s:\mathcal F_s\to\mathcal F_s$ défini par la multiplication
par $g_s$, est nul dans un voisinage de $\eta$. D'autre part, soit
$v:\mathcal O_X\to\mathcal O_X$ l'endomorphisme défini par la multiplication
par $g$~; comme $B_\eta$ est intègre et $h=g_\eta\neq0$, $v_\eta$ est injectif,
et il résulte donc de (9.4.5) que $v_s$ est un endomorphisme injectif de
$\mathcal O_{X_s}$ pour $s$ voisin de $\eta$, autrement dit, $g_s$ est un
élément $\mathcal O_{X_s}$-régulier pour ces valeurs de $s$~; donc
$\mathcal F_s$ est de torsion dans un voisinage de $\eta$.

$2^\circ$ $\eta\in S-E$. Dire qu'un $B_\eta$-module de type fini $M_\eta$
n'est pas un module de torsion signifie que son quotient $M_\eta/T$ par son
sous-module de torsion est $\neq0$, et comme c'est un $B_\eta$-module sans
torsion de type fini, il est isomorphe à un sous-module d'un $B_\eta$-module
$B_\eta^n$~; il existe par suite un homomorphisme
$w:M_\eta\to B_\eta$ qui est $\neq0$. Appliquant (8.5.2, (i)) suivant la
méthode de (8.1.2, \emph{a)}), on en déduit (en remplaçant au besoin $S$ par un
voisinage de $\eta$) qu'il existe un homomorphisme
$v:\mathcal F\to\mathcal O_X$ tel que $v_\eta=\widetilde w$. L'hypothèse
$v_\eta\neq0$ entraîne donc (9.4.5) que $v_s\neq0$ dans un voisinage de
$\eta$, et comme $X_s$ est localement intègre, $\mathcal F_s$ n'est pas de
torsion pour ces valeurs de $s$.

$3^\circ$ $\eta\in E'$. Comme $M_\eta$ est un $B_\eta$-module de type fini
sans torsion, il existe un homomorphisme injectif
$w:M_\eta\to B_\eta^n$. Utilisant (8.5.2, (i)) et (9.4.5) comme dans le
$2^\circ$ (en restreignant au besoin $S$), on en déduit qu'il existe un
homomorphisme $v:\mathcal F\to\mathcal O_X^n$ tel que $v_\eta=\widetilde w$,
et que pour $s$ voisin de $\eta$,
$v_s:\mathcal F_s\to\mathcal O_{X_s}^n$ est injectif~; pour ces valeurs de
$s$, $\mathcal F_s$ est donc sans torsion.

$4^\circ$ $\eta\in S-E'$. Soit $T$ le sous-module de torsion de $M_\eta$~;
par hypothèse $T\neq0$, et $T$ est de type fini puisque $M_\eta$ est
noethérien. Utilisant cette fois (8.5.2, (i) et (ii)) on voit (en restreignant
au besoin $S$) qu'il existe un $\mathcal O_X$-Module cohérent $\mathcal G$ et
un homomorphisme injectif $u:\mathcal G\to\mathcal F$ tels que
$\mathcal G_\eta=\widetilde T$ et que $u_\eta$ soit l'injection canonique
$\widetilde T\to\mathcal F_\eta$. Il résulte alors du $1^\circ$ et de (1.8.6)
que dans un voisinage de $\eta$, $\mathcal G_s$ est un
$\mathcal O_{X_s}$-Module de torsion $\neq0$, et d'autre part il résulte de
(9.4.5) que dans un voisinage de $\eta$, $u_s$ est injectif. On en conclut que
dans un voisinage de $\eta$, le sous-Module de torsion de $\mathcal F_s$ n'est
pas nul. C.Q.F.D.

\begin{remark}[9.4.9]
\label{IV.9.4.9-fr}
La propriété «~$X$ est un $k$-préschéma algébrique localement intègre~» ne
vérifie pas la condition (9.2.1, (i)), et il n'est donc pas certain que
l'énoncé (9.4.8) soit encore valable quand on ne fait aucune hypothèse sur $S$
et qu'on suppose seulement que $f$ est un morphisme de présentation finie et
$\mathcal F$ un $\mathcal O_X$-Module de présentation finie. Considérons
toutefois le cas particulier suivant~: $S_0$ étant un préschéma localement
noethérien, soient $f_0:X_0\to S_0$ un morphisme de type fini, tel que les
fibres $(X_0)_{s_0}$ soient localement intègres (pour tout $s_0\in S_0$), et
$\mathcal F_0$ un $\mathcal O_{X_0}$-Module cohérent~; soit $g:S\to S_0$ un
morphisme quelconque, posons
$X=X_0\times_{S_0}S$, $\mathcal F=\mathcal F_0\otimes_{\mathcal O_{S_0}}
\mathcal O_S$, $f=f_0\times1_S:X\to S$, et \emph{supposons} que pour tout
$s\in S$, la fibre $X_s$ soit encore localement intègre. Alors l'ensemble $E$
(resp. $E'$) des $s\in S$ tels que $\mathcal F_s$ soit de torsion (resp. sans

\oldpage[IV-3]{67}

torsion) est encore \emph{localement constructible}. En effet, soit $s\in S$
et soit $s_0=g(s)$~; il suffira (compte tenu de (1.8.2)) de prouver que, pour
que $\mathcal F_s$ soit de torsion (resp. sans torsion), il faut et il suffit
que $(\mathcal F_0)_{s_0}$ le soit. Or,
$\operatorname{Supp}(\mathcal F_s)$ est l'image réciproque de
$\operatorname{Supp}((\mathcal F_0)_{s_0})$ par la projection
$p:X_s\to(X_0)_{s_0}$ (\textbf{I}, 9.1.13)~; comme $p$ est fidèlement plat et
quasi-compact, dire que $\operatorname{Supp}(\mathcal F_s)$ contient un point
maximal de $X_s$ équivaut à dire que
$\operatorname{Supp}((\mathcal F_0)_{s_0})$ contient un point maximal de
$(X_0)_{s_0}$ (1.1.5 et 2.3.4)~; d'où notre assertion en ce qui concerne
l'ensemble $E$ (\textbf{I}, 7.4.6). Si le sous-Module de torsion $\mathcal T$
de $(\mathcal F_0)_{s_0}$ n'est pas nul,
$\mathcal T\otimes_{k(s_0)}k(s)$ (qui est de torsion d'après ce qui précède)
n'est pas nul et s'identifie à un sous-Module de $\mathcal F_s$ (2.2.7), donc
le sous-Module de torsion de $\mathcal F_s$ n'est pas nul. Enfin, si
$(\mathcal F_0)_{s_0}$ est sans torsion, on peut supposer (en considérant un
ouvert affine de $(X_0)_{s_0}$) que $(\mathcal F_0)_{s_0}$ est isomorphe à un
sous-Module d'un $\mathcal O_{(X_0)_{s_0}}^n$, donc $\mathcal F_s$ est
isomorphe à un sous-Module d'un $\mathcal O_{X_s}^n$ (2.2.7), et ceci établit
notre assertion concernant $E'$.
\end{remark}

\subsection{Constructibilité de propriétés topologiques}
\label{subsection:IV.9.5-fr}

\begin{proposition}[9.5.1]
\label{IV.9.5.1-fr}
Soient $f:X\to S$ un morphisme de présentation finie, $Z$ une partie
localement constructible de $X$. Alors l'ensemble $E$ des $s\in S$ tels que
$Z_s\neq\varnothing$ est localement constructible.
\end{proposition}

En effet, on a $E=f(Z)$, et il suffit d'appliquer le th. de Chevalley (1.8.4).

\begin{corollary}[9.5.2]
\label{IV.9.5.2-fr}
Si $Z'$, $Z''$ sont deux parties localement constructibles de $X$, l'ensemble
des $s\in S$ tels que $Z'_s\subset Z''_s$ \textnormal{(}resp.
$Z'_s=Z''_s$\textnormal{)} est localement constructible.
\end{corollary}

En effet, la relation $Z'_s\subset Z''_s$ équivaut à
$(Z'\cap(\complement Z''))_s=\varnothing$ et
$Z'\cap\complement Z''$ est localement constructible.

\begin{proposition}[9.5.3]
\label{IV.9.5.3-fr}
Soient $f:X\to S$ un morphisme de présentation finie, $Z$, $Z'$ deux parties
localement constructibles de $X$ telles que $Z\subset Z'$. Alors l'ensemble
$E$ des $s\in S$ tels que $Z_s$ soit dense dans $Z'_s$ est localement
constructible dans $S$.
\end{proposition}

Il faut vérifier les deux conditions de (9.2.2, (iii)). En ce qui concerne la
première, considérons un préschéma $X$ algébrique sur un corps $k$, deux parties
constructibles $Z$, $Z'$ de $X$ telles que $Z\subset Z'$, et une extension
$k'$ de $k$. Alors la projection canonique $p:X_{(k')}\to X$ est un morphisme
fidèlement plat et quasi-compact, et l'on a donc
$p^{-1}(\overline Z)=\overline{p^{-1}(Z)}$ et
$p^{-1}(\overline{Z'})=\overline{p^{-1}(Z')}$ en vertu de (2.3.10)~; comme
$p$ est surjectif, la relation $\overline Z=\overline{Z'}$ est équivalente à
$\overline{p^{-1}(Z)}=\overline{p^{-1}(Z')}$.

Vérifions maintenant la seconde condition, et supposons donc $S$ affine,
noethérien et intègre, de point générique $\eta$. Distinguons deux cas~:

$1^\circ$ $\eta\in S-E$, autrement dit, $Z_\eta$ n'est pas dense dans
$Z'_\eta$~; il existe donc dans $X$ un ouvert $V$ tel que
$V\cap Z_\eta=\varnothing$ et $V\cap Z'_\eta\neq\varnothing$. Comme $X$ est
noethérien, $V$ est localement constructible, donc il en est de même de
$V\cap Z$, et en vertu de (9.5.1), il y a un voisinage $U$ de $\eta$ dans $S$
tel que pour tout $s\in U$, on ait $(V\cap Z)_s=\varnothing$ et
$(V\cap Z')_s\neq\varnothing$~; cela entraîne que $Z_s$ n'est pas dense dans
$Z'_s$ pour $s\in U$, autrement dit $U\subset S-E$.

\oldpage[IV-3]{68}

$2^\circ$ $\eta\in E$, donc $Z_\eta$ est dense dans $Z'_\eta$. Montrons
d'abord que l'on peut supposer $Z'$ \emph{fermé}. En effet, $Z_\eta$ est dense
dans $\overline{Z'_\eta}$ (adhérence prise dans $X_\eta$)~; posons
$V_\eta=X_\eta-\overline{Z_\eta}$, qui est ouvert dans $X_\eta$ et ne
rencontre pas $Z'_\eta$~; on peut supposer $V_\eta$ de la forme $V\cap X_\eta$,
où $V$ est ouvert (donc constructible) dans $X$, et l'hypothèse
$V_\eta\cap Z'_\eta=\varnothing$ entraîne alors
$V_s\cap Z'_s=\varnothing$ pour tout $s$ voisin de $\eta$ en vertu de (9.5.1).
Remplaçant $S$ par un voisinage ouvert de $\eta$, on peut donc supposer que
$V\cap Z'=\varnothing$, donc $V\cap\overline{Z'}=\varnothing$ (adhérence prise
dans $X$), et par suite $(\overline{Z'})_\eta=\overline{Z'_\eta}$, d'où notre
assertion. L'ensemble $Z'$ est alors réunion de ses composantes irréductibles
en nombre fini, et en restreignant encore $S$ à un voisinage de $\eta$, on
peut supposer que toutes les composantes irréductibles $Z'_i$ de $Z'$
rencontrent $X_\eta$, d'où résulte que $X_\eta$ contient les points génériques
des $Z'_i$ $(\mathbf 0_{\mathrm I},2.1.8)$. Dire que $Z_s$ est dense dans
$Z'_s$ équivaut alors à dire que chacun des $(Z\cap Z'_i)_s$ est dense dans
$(Z'_i)_s$, et l'on est ainsi ramené au cas où $Z'$ est irréductible.
Remplaçant alors au besoin $X$ par le sous-préschéma réduit ayant $Z'$ pour
espace sous-jacent, on voit qu'on peut supposer que $Z'=X$ et que $X$ est
\emph{intègre} et domine $S$. Enfin, en recouvrant $X$ par un nombre fini
d'ouverts affines $W_j$ et remplaçant $Z$ par $Z\cap W_j$, on peut supposer
que $X=\operatorname{Spec}(A)$, où $A$ est un anneau noethérien \emph{intègre}.
Comme $X_\eta$ est noethérien intègre et que $Z_\eta$ est constructible dans
$X_\eta$ et dense dans $X_\eta$, $Z_\eta$ contient un ouvert non vide de
$X_\eta$ $(\mathbf 0_{\mathrm{III}},9.2.2)$, que l'on peut supposer de la forme
$(D(t))_\eta$, où $t$ est un élément $\neq0$ de $A$. En remplaçant au besoin
$S$ par un voisinage de $\eta$, on peut en outre, en vertu de la relation
$(D(t))_\eta\subset Z_\eta$, supposer que $D(t)\subset Z$ (9.5.2). Enfin,
comme l'homothétie de rapport $t_\eta$ dans $\mathcal O_{X_\eta}$ est
injective, il résulte de (9.4.5) que pour $s$ voisin de $\eta$, $t_s$ est
$\mathcal O_{X_s}$-régulier, donc $(X_s)_{t_s}$ est dense dans $X_s$, et
\emph{a fortiori} il en est de même de $Z_s$ qui contient $(X_s)_{t_s}$.

\begin{corollary}[9.5.4]
\label{IV.9.5.4-fr}
Soient $f:X\to S$ un morphisme de présentation finie, $Z$ une partie
localement constructible de $X$. L'ensemble $E$ des $s\in S$ tels que $Z_s$
soit fermé \textnormal{(}resp. ouvert, resp. localement fermé\textnormal{)}
dans $X_s$ est localement constructible dans $S$.
\end{corollary}

Dire que $Z_s$ est ouvert dans $X_s$ signifie que
$(X-Z)_s=X_s-Z_s$ est fermé dans $X_s$, et comme $X-Z$ est localement
constructible, on peut se borner à considérer l'ensemble des $s\in S$ tels que
$Z_s$ soit fermé et l'ensemble des $s\in S$ tels que $Z_s$ soit localement
fermé.

Vérifions encore dans chaque cas les deux conditions de (9.2.2, (ii)). La
première résulte du fait que $p:X_{(k')}\to X$ est fidèlement plat et
quasi-compact, et de (2.3.12) et (2.3.14). Vérifions donc la seconde condition,
$S$ étant supposé affine, noethérien et intègre, de point générique $\eta$.
Posons $Z'=\overline Z$~; le même raisonnement que dans (9.5.3) montre que
$Z'_\eta$ est égal à l'adhérence de $Z_\eta$ dans $X_\eta$~; en vertu de
(9.5.3), il y a donc un voisinage $U$ de $\eta$ tel que pour $s\in U$, $Z_s$
soit dense dans $Z'_s$, ce dernier étant fermé dans $X_s$. Dire que $Z_s$ est
fermé dans $X_s$ signifie alors que $Z''_s=\varnothing$, où $Z''=Z'-Z$~; il
résulte donc de (9.5.1) que l'ensemble $E$ des $s\in S$ où
$Z''_s=\varnothing$ est tel que $E$ ou $S-E$ contienne un voisinage de
$\eta$. Dire que $Z_s$ est localement fermé dans $X_s$ signifie que $Z''_s$
est fermé dans $X_s$~; il suffit donc d'appliquer le résultat précédent en
remplaçant $Z$ par $Z''$, qui est localement constructible dans $X$.

\oldpage[IV-3]{69}

\begin{proposition}[9.5.5]
\label{IV.9.5.5-fr}
Soient $f:X\to S$ un morphisme de présentation finie, $Z$ une partie
localement constructible de $X$ telle que, pour tout $s\in S$, $Z_s$ soit
localement fermé dans $X_s$. Pour tout $s\in S$, soit
$D(s)\subset\mathbf Z\cup\{-\infty\}$ l'ensemble des dimensions des
composantes irréductibles de $Z_s$. Alors la fonction $s\mapsto D(s)$ est
localement constructible dans $S$.
\end{proposition}

Soit $\Phi$ une partie finie de $\mathbf Z\cup\{-\infty\}$~; il s'agit de
montrer que l'ensemble des $s\in S$ tels que $D(s)=\Phi$ est localement
constructible~; compte tenu de (9.2.3), nous avons encore à vérifier les deux
conditions de (9.2.2, (iii)).

En ce qui concerne la première, il s'agit de voir que si $X$ est un préschéma
algébrique sur un corps $k$, $Z$ une partie localement fermée de $X$, $k'$ une
extension de $k$, $p:X_{(k')}\to X$ la projection canonique, alors l'ensemble
des dimensions des composantes irréductibles de $Z$ est le même que celui des
dimensions des composantes irréductibles de $p^{-1}(Z)$~; compte tenu de
l'existence d'un sous-préschéma de $X$ ayant $Z$ pour espace sous-jacent
(\textbf{I}, 5.2.1), cela résulte de (4.2.8).

Pour la seconde condition de (9.2.2, (iii)), on est dans le cas où $S$ est
noethérien et intègre de point générique $\eta$, et $f:X\to S$ est un
morphisme de type fini, de sorte que $X$ est noethérien. Le sous-espace
$Z_\eta$ de l'espace noethérien $X_\eta$ est par hypothèse localement fermé,
donc a un nombre fini de composantes irréductibles $Z_{i\eta}$, qui sont
localement fermées dans $X_\eta$. Il existe par suite pour chaque indice $i$
une partie localement fermée $Z_i$ de $X$ telle que $(Z_i)_\eta=Z_{i\eta}$,
donc si $Z'=\bigcup_i Z_i$, on a $Z'_\eta=Z_\eta$. Mais comme $Z$ et $Z'$ sont
localement constructibles, on peut, en remplaçant $S$ par un voisinage de
$\eta$, supposer que $Z'=Z$ (9.5.1). En outre, pour $i\neq j$,
$Z_{i\eta}\cap Z_{j\eta}$ est rare et \emph{fermé} dans $Z_{i\eta}$~; donc
(9.5.3 et 9.5.4), on peut encore supposer, en restreignant $S$, que pour
$j\neq i$, $(Z_i)_s\cap(Z_j)_s$ est rare et fermé dans $(Z_i)_s$. Comme $Z_i$
est localement fermé dans $X$, il y a un sous-préschéma de $X$ ayant $Z_i$
pour espace sous-jacent (et encore noté $Z_i$), qui est de type fini sur
$S$ (\textbf{I}, 6.3.5). Posons
$U_i=Z_i-\bigcup_{j\neq i}(Z_i\cap Z_j)$, qui est ouvert dans $Z_i$ et tel que,
pour tout $s\in S$, $(U_i)_s$ soit ouvert partout dense dans $(Z_i)_s$~; en
outre, par construction, les $U_i$ sont deux à deux sans point commun. Comme
$(Z_i)_s$ est un $k(s)$-préschéma algébrique, l'ensemble des dimensions des
composantes irréductibles de $Z_s=\bigcup_i(Z_i)_s$ est le même que l'ensemble
des dimensions des composantes irréductibles de la réunion des $(U_i)_s$
(4.1.1.3), chacune de ces composantes étant déjà une composante irréductible
de l'un des $(U_i)_s$. On peut donc se borner au cas où $Z=U_i$~; en outre,
comme $Z_\eta$ est alors irréductible, il n'y a qu'une seule composante
irréductible de $Z$ qui rencontre $X_\eta$ $(\mathbf 0_{\mathrm I},2.1.8)$ et
l'on peut évidemment, en restreignant $Z$, supposer $Z$ irréductible. On est
finalement ramené à prouver le cas particulier suivant de (9.5.5)~:

\begin{corollary}[9.5.6]
\label{IV.9.5.6-fr}
Soient $S$ un préschéma noethérien et irréductible de point générique $\eta$,
$X$ un préschéma irréductible, $f:X\to S$ un morphisme dominant de type fini.
Alors il existe un voisinage $U$ de $\eta$ dans $S$ tel que, pour tout
$s\in U$, toutes les composantes irréductibles de $X_s$ soient de dimension
$n=\dim(X_\eta)$.
\end{corollary}

On peut évidemment se borner au cas où $S=\operatorname{Spec}(A)$ est affine,
$A$ étant donc noethérien~; remplaçant $f$ par $f_{\mathrm{red}}$, qui est de
type fini (1.5.4, (vii)), on peut

\oldpage[IV-3]{70}

supposer $A$ intègre et $X$ réduit, donc intègre puisqu'il est irréductible. On
sait (4.1.2) qu'il existe un ouvert $W$ dense dans $X_\eta$ et un
$k(\eta)$-morphisme \emph{fini et surjectif}
$h:W\to\mathbf V((k(\eta))^n)=\operatorname{Spec}(B')$, où
$B'=k(\eta)[T_1,\ldots,T_n]$ ($T_i$ indéterminées). Si $V$ est un ouvert de
$X$ tel que $V\cap X_\eta=W$, on sait (9.5.3) que pour $s$ voisin de $\eta$
dans $S$, $V_s$ est un ouvert dense dans $X_s$, et l'on peut par suite (4.1.1.3)
se borner au cas où $V=X$, $W=X_\eta$. Posons
$B=A[T_1,\ldots,T_n]$ et $Y=\operatorname{Spec}(B)=\mathbf V_A^n$, de sorte
que $\operatorname{Spec}(B')=Y_\eta$~; il résulte de (8.8.2, (i)) et
(8.10.5, (vi) et (x)), appliqués suivant la méthode de (8.1.2, \emph{a)}),
qu'en remplaçant au besoin $S$ par un voisinage de $\eta$, on peut supposer que
$h=g_\eta$, où $g:X\to Y$ est un morphisme fini surjectif~; autrement dit, on
a $X=\operatorname{Spec}(C)$, où $C$ est une $B$-algèbre finie et
l'homomorphisme $B\to C$ correspondant à $g$ est \emph{injectif}~; comme $C$
est un anneau intègre, $C$ est donc un $B$-module de type fini \emph{sans
torsion}, et $C_\eta=C\otimes_A k(\eta)$ est donc un module de type fini sans
torsion sur $B_\eta=B'=B\otimes_A k(\eta)=k(\eta)[T_1,\ldots,T_n]$ (étant un
module de fractions dont les dénominateurs sont dans
$A-\{0\}\subset B-\{0\}$). Il résulte donc de (9.4.8) qu'il y a un voisinage
$U$ de $\eta$ dans $S$ tel que pour $s\in S$,
$C_s=C\otimes_A k(s)$ soit un module de type fini \emph{sans torsion} sur
$B_s=B\otimes_A k(s)=k(s)[T_1,\ldots,T_n]$, et en particulier l'homomorphisme
$g_s:B_s\to C_s$ est injectif. Comme aucun élément de $B_s$ n'est diviseur de
0 dans $C_s$, pour tout idéal premier minimal $\mathfrak p_i$ de $C_s$ (dont
les éléments sont diviseurs de 0 dans $C_s$), on a nécessairement
$\mathfrak p_i\cap B_s=0$, donc l'homomorphisme canonique
$B_s\to C_s/\mathfrak p_i$ est injectif. On en déduit que pour chaque
composante irréductible $Z_i=\operatorname{Spec}(C_s/\mathfrak p_i)$ de $X_s$,
la restriction à $Z_i$ de $g$ est un morphisme fini et dominant $Z_i\to Y_s$,
donc \emph{surjectif} (\textbf{II}, 6.1.10). On conclut par suite de (4.1.2)
que $\dim Z_i=n$, ce qui achève la démonstration.

\begin{remark}[9.5.7]
\label{IV.9.5.7-fr}
On aura soin de noter que sous les hypothèses de (9.5.6) il peut se faire que
$X_s$ \emph{ne soit irréductible pour aucun} $s\neq\eta$ dans un voisinage de
$\eta$, autrement dit, la propriété «~$X$ est un $k$-préschéma algébrique
irréductible~» \emph{n'est pas constructible}. Prenons par exemple
$S=\operatorname{Spec}(k[T])$, où $k$ est un corps algébriquement clos, $T$ une
indéterminée~; on a donc $k(\eta)=K=k(T)$. Soit $L$ une extension finie
séparable de $K$ de degré $>1$, et soit $X$ la \emph{fermeture intégrale} de
$S$ dans $L$ (\textbf{II}, 6.3.4)~; on a donc $X=\operatorname{Spec}(B)$, où
$B$ est la fermeture intégrale de $k[T]$ dans $L$. On sait que $B$ est un
anneau de Dedekind, et que tous les idéaux maximaux de $k[T]$, sauf un nombre
fini, sont non ramifiés dans $B$~; comme en outre le corps résiduel de tout
idéal maximal de $B$ est nécessairement $k$ (puisque c'est une extension finie
de $k$), on voit que pour presque tous les idéaux maximaux $\mathfrak j_s$ de
$k[T]$, $B_s=B/\mathfrak j_sB$ est composé direct de $[L:K]$ corps isomorphes
à $k$, autrement dit $X_s$ n'est pas irréductible, bien que
$X_\eta=\operatorname{Spec}(L)$ le soit. Le même exemple montre que la
propriété «~$X$ est un $k$-préschéma algébrique intègre~» n'est pas
constructible. Enfin il en est de même de la propriété «~$X$ est un
$k$-préschéma algébrique réduit~». Il suffit pour le voir de prendre encore
$S=\operatorname{Spec}(k[T])$, où cette fois $k$ est un corps algébriquement
clos de caractéristique $p>0$, et pour $X$ la fermeture intégrale de $S$ dans
$L=K^{1/p}$ (où $K=k(T)$), de sorte que $X=\operatorname{Spec}(B)$ avec
$B=k[T^{1/p}]$ ($k$ étant parfait)~; tout idéal maximal de $k[T]$ est de la
forme $(T-\alpha)$ avec $\alpha\in k$~; l'unique idéal de $B$ au-dessus de
l'idéal $(T-\alpha)$ est l'idéal principal
$(T^{1/p}-\alpha^{1/p})$ et il est immédiat

\oldpage[IV-3]{71}

que l'anneau quotient $B/(T-\alpha)B$ admet par suite des éléments nilpotents~;
autrement dit, $X_s$ n'est réduit pour aucun $s\neq\eta$, tandis que
$X_\eta=\operatorname{Spec}(L)$ est intègre.

Nous verrons un peu plus loin (9.7) que l'on obtient par contre des propriétés
constructibles lorsqu'on considère les notions «~géométriques~»
correspondantes aux notions de préschéma irréductible, réduit ou intègre (4.5
et 4.6).
\end{remark}

\subsection{Constructibilité de certaines propriétés des morphismes}
\label{subsection:IV.9.6-fr}

\begin{proposition}[9.6.1]
\label{IV.9.6.1-fr}
Soient $X$, $Y$ deux $S$-préschémas de présentation finie, $f:X\to Y$ un
$S$-morphisme. Soit $E$ l'ensemble des $s\in S$ pour lesquels $f_s$ a l'une
des propriétés suivantes~: être~:
\begin{enumerate}[label=(\roman*),widest=xii]
  \item surjectif~;
  \item dominant~;
  \item séparé~;
  \item propre~;
  \item radiciel~;
  \item fini~;
  \item quasi-fini~;
  \item une immersion~;
  \item une immersion fermée~;
  \item une immersion ouverte~;
  \item un isomorphisme~;
  \item un monomorphisme.
\end{enumerate}
Alors $E$ est localement constructible dans $S$.
\end{proposition}

Les assertions de (i), (v) et (vii) ne sont insérées que pour mémoire, ayant
déjà été établies dans (9.3.2).

(ii)~: Notons d'abord que $f$ est de présentation finie (1.6.2, (v)), donc, en
vertu du théorème de Chevalley (1.8.4), $f(X)$ est localement constructible dans
$Y$. Par ailleurs, on a $f_s(X_s)=(f(X))_s$ (\textbf{I}, 3.6.1)~; l'ensemble
des $s$ tels que $f_s$ soit dominant est l'ensemble des $s$ tels que
$(f(X))_s$ soit dense dans $Y_s$~; la conclusion résulte donc dans ce cas de
(9.5.3).

(iii)~: Comme $f:X\to Y$ est de présentation finie, l'immersion diagonale
$\Delta_f:X\to X\times_YX$ est de présentation finie (1.6.2, (iv) et (v))~; il
résulte donc de (1.8.4.1) que $\Delta_f(X)=\Delta_X$ est une partie localement
constructible de $X\times_YX$. Dire que $f_s$ est séparé signifie (compte tenu
de (\textbf{I}, 5.3.4)) que $(\Delta_X)_s$ est fermé dans
$X_s\times_{Y_s}X_s=(X\times_YX)_s$~; la conclusion résulte donc ici de
(9.5.4).

(iv)~: Montrons que la propriété «~$X$ et $Y$ sont des préschémas algébriques
sur un corps $k$, et $f:X\to Y$ un $k$-morphisme propre~» est constructible. La
condition (9.2.1, (i)) est vérifiée en vertu de (2.7.1, (vii)). On peut donc se
borner au cas où $S$ est affine, noethérien et intègre, de point générique
$\eta$, et on doit prouver que $E$ ou $S-E$ est un voisinage de $\eta$.
Supposons d'abord que $\eta\in E$~; il résulte de (8.10.5, (xii)) appliqué
suivant la méthode de (8.1.2, \emph{a)}) qu'en remplaçant $S$ par un voisinage
de $\eta$, on peut

\oldpage[IV-3]{72}

supposer que $f$ est lui-même un morphisme propre~; on sait alors qu'il en est
de même de $f_s$ pour tout $s\in S$ (\textbf{II}, 5.4.2, (iii)). Supposons au
contraire que $\eta\in S-E$, et distinguons deux cas~:

\begin{enumerate}[label=\upshape{$\arabic{enumi}^\circ$}]
  \item Supposons que $f_\eta$ soit non séparé~; alors il résulte de (iii) que
  $f_s$ est non séparé (et \emph{a fortiori} non propre) dans un voisinage de
  $\eta$.

  \item Supposons $f_\eta$ séparé~; $Y$ est réunion finie d'ouverts affines
  $V_i$, et pour que $f_s$ soit propre, il faut et il suffit que chacune de ses
  restrictions $f_s^{-1}((V_i)_s)\to(V_i)_s$ le soit (\textbf{II}, 5.4.1)~;
  on peut donc se borner au cas où $Y$ est \emph{affine}, donc un schéma. Dire
  que $f_\eta$ n'est pas propre signifie (\textbf{II}, 5.6.3) qu'il existe un
  morphisme \emph{de type fini} $h:Z\to Y_\eta$ tel que le morphisme
  $(f_\eta)_{(Z)}:X_\eta\times_{Y_\eta}Z\to Z$ ne soit pas fermé. Comme $Y$ est
  un schéma, on déduit de (8.8.2, (i) et (ii)) (en restreignant au besoin $S$ à
  un voisinage de $\eta$) qu'il existe un morphisme de type fini $g:Y'\to Y$
  tel que $Z$ soit isomorphe à $Y'_\eta$, et que $g_\eta=h$~; si l'on pose
  $X'=X\times_YY'$, $f'=f_{(Y')}:X'\to Y'$, on a
  $(f_\eta)_{(Z)}=f'_\eta$, et par hypothèse il existe donc une partie fermée
  $M'$ de $X'_\eta$ telle que $f'_\eta(M')$ ne soit pas fermé dans $Y'_\eta$.
  Or $M'$ est la trace sur $X'_\eta$ d'une partie fermée $N'$ de $X'$~; comme
  $X'$ est noethérien et $f'$ de type fini, $f'(N')$ est constructible dans
  $Y'$ (1.8.4) et par hypothèse
  $(f'(N'))_\eta=f'_\eta(N'_\eta)=f'_\eta(M')$ n'est pas fermé dans
  $Y'_\eta$. On conclut alors de (9.5.4) qu'il existe un voisinage $U$ de
  $\eta$ dans $S$ tel que pour $s\in U$,
  $(f'(N'))_s=f'_s(N'_s)$ ne soit pas fermé dans $Y'_s$~; autrement dit, le
  morphisme $f'_s$ n'est pas fermé, et \emph{a fortiori} le morphisme $f_s$
  n'est pas propre.
\end{enumerate}

(vi) La propriété est conjonction des propriétés (iv) et (vii) (8.11.1), donc
la proposition résulte dans ce cas de ce qui a déjà été prouvé.

(ix) La vérification de la condition (9.2.1, (i)) résulte de (2.7.1, (xi)). On
peut donc se borner au cas où $S$ est affine, noethérien et intègre de point
générique $\eta$, et à prouver que $E$ ou $S-E$ est un voisinage de $\eta$. Si
$\eta\in E$, on peut (en remplaçant $S$ par un voisinage de $\eta$) supposer
que $f$ est une immersion fermée en vertu de (8.10.5, (iv)), et alors il est
clair que $f_s$ est une immersion fermée pour tout $s\in S$ (\textbf{I},
4.3.2). Supposons donc que $\eta\in S-E$ et distinguons deux cas~:

\begin{enumerate}[label=\upshape{$\arabic{enumi}^\circ$}]
  \item $f_\eta$ n'est pas un morphisme fini. Alors il résulte de (vi) que dans
  un voisinage de $\eta$, $f_s$ n'est pas fini, ni \emph{a fortiori} une
  immersion fermée.

  \item $f_\eta$ est fini~; alors, en vertu de (8.10.5, (x)), on peut supposer
  (en restreignant au besoin $S$) que $f$ lui-même est un morphisme \emph{fini}.
  Dans ce cas $\mathcal A=\mathcal A(X)=f_*(\mathcal O_X)$ est un
  $\mathcal O_Y$-Module cohérent (\textbf{II}, 6.1.3) et
  $f=\mathcal A(u)$, où $u:\mathcal O_Y\to\mathcal A$ est un homomorphisme de
  $\mathcal O_Y$-Algèbres (\textbf{II}, 1.1.2)~; comme $f_\eta$ n'est pas une
  immersion fermée par hypothèse, $u_\eta$ n'est pas surjectif (\textbf{II},
  1.4.10), donc (9.4.5) il existe un voisinage de $\eta$ dans lequel $u_s$
  n'est pas surjectif, et par suite (\textbf{I}, 4.2.3) $f_s$ n'est pas une
  immersion fermée.
\end{enumerate}

(viii) La vérification de la condition (9.2.1, (i)) se fait comme dans (ix), en
utilisant cette fois (2.7.1, (x)) et le fait que toute immersion d'un préschéma
noethérien dans un autre est quasi-compacte. On est donc ramené au cas où $S$
est affine, noethérien et intègre de point générique $\eta$, et à prouver que
$E$ ou $S-E$ est un voisinage de $\eta$. Si $\eta\in E$, on conclut comme dans
(ix) à l'aide de (8.10.5, (ii)). Si $\eta\in S-E$, on distingue de nouveau
deux cas~:

\begin{enumerate}[label=\upshape{$\arabic{enumi}^\circ$}]
  \item $f_\eta(X_\eta)$ n'est pas une partie localement fermée de $Y_\eta$.
  Comme $f(X)$ est constructible dans $Y$ (1.8.4) et que
  $f_s(X_s)=(f(X))_s$, on déduit de (9.5.4) que pour $s$ voisin de $\eta$,
  $f_s(X_s)$ n'est pas localement fermé dans $Y_s$, et \emph{a fortiori}
  $f_s$ n'est pas une immersion.

  \item\oldpage[IV-3]{73} $f_\eta(X_\eta)$ est localement fermé dans $Y_\eta$.
  Comme $f(X)$ est constructible dans $Y$ (1.8.4) et qu'il en est de même de
  $\overline{f(X)}$ puisque $Y$ est noethérien, il résulte de (8.3.11) qu'en
  restreignant au besoin $S$, on peut supposer que $f(X)$ est localement fermé
  dans $Y$. Il y a alors un ouvert $V$ de $Y$ contenant $f(X)$ et dans lequel
  $f(X)$ est fermé. Comme $Y$ est noethérien, $V$ est de type fini sur $S$, et
  en remplaçant $Y$ par $V$, on peut donc se ramener au cas où $f(X)$ est
  \emph{fermé} dans $Y$. Mais alors $f_s(X_s)=(f(X))_s$ est fermé dans $Y_s$
  pour tout $s\in S$, et dire que $f_s$ est une immersion équivaut à dire que
  $f_s$ est une immersion fermée~; on est donc ramené à ce qui a été démontré
  dans (ix).
\end{enumerate}

(x) Utilisant cette fois (2.7.1, (ix)) et (8.10.5, (iii)), on est ramené au cas
où $S$ est affine, noethérien, intègre de point générique $\eta$, et où
$\eta\in S-E$. Distinguons trois cas~:

\begin{enumerate}[label=\upshape{$\arabic{enumi}^\circ$}]
  \item $f_\eta(X_\eta)$ n'est pas ouvert dans $Y_\eta$. Comme $f(X)$ est
  constructible dans $Y$ (1.8.4), on déduit de (9.5.4) que $f_s(X_s)$ n'est
  pas ouvert dans $Y_s$ pour $s$ voisin de $\eta$, et \emph{a fortiori} $f_s$
  n'est pas une immersion ouverte.

  \item $f_\eta(X_\eta)$ est ouvert dans $Y_\eta$, mais $f_\eta$ n'est pas une
  immersion. Il résulte alors de (viii) que pour $s$ voisin de $\eta$ dans $S$,
  $f_s$ n'est pas une immersion, ni \emph{a fortiori} une immersion ouverte.

  \item $f_\eta(X_\eta)$ est ouvert dans $Y_\eta$ et $f_\eta$ est une
  immersion. Comme $f(X)$ est constructible dans $Y$, il résulte de (8.3.11)
  qu'en restreignant au besoin $S$, on peut déjà supposer que $f(X)$ est ouvert
  dans $Y$. Comme $Y$ est noethérien, le sous-préschéma induit sur $f(X)$ est
  de type fini sur $S$, donc on peut se ramener au cas où $f$ est \emph{surjectif}
  en remplaçant $Y$ par $f(X)$. Par hypothèse, $f_\eta$ est une immersion
  fermée, donc on peut, comme dans (ix), supposer que $f$ est une immersion
  \emph{fermée}, et par suite que $X$ est un sous-préschéma fermé de $Y$ défini
  par un Idéal cohérent $\mathcal J$ de $\mathcal O_Y$. Par hypothèse
  $f_\eta$ n'est pas un isomorphisme, donc $\mathcal J_\eta\neq0$~; on en
  conclut (9.4.5) que dans un voisinage de $\eta$, on a $\mathcal J_s\neq0$,
  et par suite l'immersion fermée surjective $f_s$ n'est pas ouverte.
\end{enumerate}

(xi) La propriété est conjonction des propriétés (i) et (x) et résulte donc de
ce qui a été démontré.

(xii) En vertu de (\textbf{I}, 5.3.8), dire que $f_s$ est un monomorphisme
signifie que
$\Delta_{f_s}=(\Delta_f)_s:X_s\to X_s\times_{Y_s}X_s=(X\times_YX)_s$ est un
isomorphisme, et comme $X\times_YX$ est un $S$-préschéma de présentation finie
(1.6.2, (iv)), la conclusion résulte de (xi).

\begin{proposition}[9.6.2]
\label{IV.9.6.2-fr}
Soient $X$, $Y$ deux $S$-préschémas de présentation finie, $f:X\to Y$ un
$S$-morphisme.
\begin{enumerate}[label=\upshape{I)}]
  \item Soit $E$ l'ensemble des $s\in S$ pour lesquels $f_s$ a l'une des
  propriétés suivantes~: être~:
  \begin{enumerate}[label=(\roman*)]
    \item affine~;
    \item quasi-affine~;
    \item projectif~;
    \item quasi-projectif.
  \end{enumerate}
  Alors $E$ est ind-constructible dans $S$.
\end{enumerate}

\oldpage[IV-3]{74}

\begin{enumerate}[label=\upshape{II)}]
  \item Soit $\mathcal L$ un $\mathcal O_X$-Module inversible. Alors l'ensemble
  $E'$ des $s\in S$ tels que $\mathcal L_s$ soit un $\mathcal O_{X_s}$-Module
  ample (resp. très ample) relativement à $f_s$, est ind-constructible dans
  $S$.
\end{enumerate}
\end{proposition}

I) Vérifions les conditions (i) et (ii) de (9.2.1). En ce qui concerne la
condition (9.2.1, (i)), elle résulte, pour les propriétés (i) et (ii), de
(2.7.1, (xiii) et (xiv))~; pour les propriétés (iii) et (iv), elle résulte de
(9.1.5). Vérifions ensuite la condition (9.2.1, (ii)), en supposant donc $S$
noethérien, intègre et de point générique $\eta$, et que
$f_\eta:X_\eta\to Y_\eta$ ait l'une des propriétés (i) à (iv) de l'énoncé.
Appliquant (8.10.5, (viii), (ix), (xiii) et (xiv)) suivant la méthode de
(8.1.2, \emph{a)}), on voit d'abord qu'il existe un voisinage ouvert $U$ de
$\eta$ tel que, si $V$ et $W$ sont les images réciproques de $U$ dans $X$ et
$Y$ respectivement par les morphismes structuraux, la restriction $V\to W$ de
$f$ a celle des propriétés (i) à (iv) que l'on considère. La conclusion
résulte alors de ce que ces propriétés sont toutes stables par changement de
base.

II) On procède de la même façon. La condition (9.2.1, (i)) résulte cette fois
de (2.7.2). Pour la condition (9.2.1, (ii)), avec les mêmes notations que dans
I), il résulte de (8.10.5.2) que pour un voisinage $U$ de $\eta$ dans $S$, la
restriction $\mathcal L|V$ est ample (resp. très ample) relativement à la
restriction $V\to W$ de $f$. La conclusion résulte encore de la stabilité des
propriétés considérées par changement de base (\textbf{II}, 4.6.13 et
4.4.10).

On peut améliorer (9.6.2, II) sous certaines conditions~:

\begin{proposition}[9.6.3]
\label{IV.9.6.3-fr}
Soient $X$, $Y$ deux $S$-préschémas de présentation finie, $f:X\to Y$ un
$S$-morphisme propre, $\mathcal L$ un $\mathcal O_X$-Module inversible. Alors
l'ensemble $E$ (resp. $E'$) des $s\in S$ tels que $\mathcal L_s$ soit un
$\mathcal O_{X_s}$-Module ample (resp. très ample) relativement à
$f_s:X_s\to Y_s$ est localement constructible dans $S$.
\end{proposition}

Soient $k$ un corps, $Z$, $T$ deux préschémas algébriques sur $k$, $\mathcal M$
un $\mathcal O_Z$-Module inversible, $g:Z\to T$ un $k$-morphisme de type fini~;
alors, si $\mathbf P(Z,T,\mathcal M,g,k)$ désigne la relation~: «~$\mathcal M$ est ample
(resp. très ample) relativement à $g$~», on a déjà remarqué dans (9.6.2) que
$\mathbf P(Z,T,\mathcal M,g,k)$ vérifie la condition (9.2.1, (i)) en vertu de (2.7.2).
On sait déjà d'autre part que $E$ et $E'$ sont ind-constructibles. Il reste donc
à voir que si $S$ est noethérien, intègre, de point générique $\eta$ et si
$\eta\in S-E$ (resp. $\eta\in S-E'$), alors $S-E$ (resp. $S-E'$) contient un
voisinage de $\eta$. Nous considérerons séparément le cas de $E'$ et celui de
$E$.

I) \emph{Cas de $E'$.} --- Notons que puisque $f$ est séparé et $Y$
quasi-compact, il existe un entier $h$ tel que, pour $q>h$, on ait
$R^qf_*(\mathcal L)=0$ (\textbf{III}, 1.4.12)~; d'autre part, puisque $f$ est
propre et $Y$ noethérien, les $R^qf_*(\mathcal L)$ sont tous des
$\mathcal O_Y$-Modules cohérents (\textbf{III}, 3.2.1)~; comme ils sont nuls
sauf pour un nombre fini de valeurs de $q$, le théorème de platitude générique
(6.9.1) montre qu'en restreignant $S$ à un voisinage de $\eta$, on peut
supposer que $\mathcal L$ et les $R^qf_*(\mathcal L)$ sont tous $S$-plats. On
conclut alors de (\textbf{III}, 6.9.9) que l'homomorphisme canonique
\begin{equation}
\label{IV.9.6.3.1-fr}
 f_*(\mathcal L)\otimes_{\mathcal O_S} k(s)
 \longrightarrow (f_s)_*(\mathcal L_s)
 \tag{9.6.3.1}
\end{equation}
est un \emph{isomorphisme}.

\oldpage[IV-3]{75}

Cela étant, il résulte de (\textbf{II}, 4.4.4) que dire que $\mathcal L_s$
n'est pas très ample relativement à $f_s$ signifie que~: \emph{ou bien}
l'homomorphisme canonique
$(f_s)^*((f_s)_*(\mathcal L_s))\to\mathcal L_s$ n'est pas surjectif~; \emph{ou
bien} l'homomorphisme précédent est surjectif et le morphisme canonique
$r:X_s\to\mathbf P((f_s)_*(\mathcal L_s))$ n'est pas une immersion. Compte
tenu de l'isomorphisme (9.6.3.1), ces conditions s'écrivent respectivement sous
la forme~: $1^\circ$ l'homomorphisme canonique
$(f^*(f_*(\mathcal L)))_s\to\mathcal L_s$ n'est pas surjectif~; $2^\circ$
l'homomorphisme précédent est surjectif et le morphisme canonique
$X_s\to\mathbf P((f_*(\mathcal L))_s)$ n'est pas une immersion.

Supposons d'abord que l'homomorphisme canonique
$(f^*(f_*(\mathcal L)))_\eta\to\mathcal L_\eta$ ne soit pas surjectif. Comme
$f_*(\mathcal L)$ est cohérent, il en est de même de $f^*(f_*(\mathcal L))$,
et alors il résulte de (9.4.5) que pour tout $s$ dans un voisinage de $\eta$,
l'homomorphisme $(f^*(f_*(\mathcal L)))_s\to\mathcal L_s$ n'est pas
surjectif, ce qui prouve dans ce cas que $S-E'$ est un voisinage de $\eta$.

Supposons en second lieu que l'homomorphisme canonique
$(f^*(f_*(\mathcal L)))_\eta\to\mathcal L_\eta$ soit surjectif mais que le
morphisme $X_\eta\to\mathbf P((f_*(\mathcal L))_\eta)$ ne soit pas une
immersion. Alors le même raisonnement que ci-dessus montre d'abord que pour
tout $s$ assez voisin de $\eta$, l'homomorphisme
$(f^*(f_*(\mathcal L)))_s\to\mathcal L_s$ est surjectif~; d'autre part, en
vertu de (9.6.1, (viii)), pour $s$ assez voisin de $\eta$, le morphisme
$X_s\to\mathbf P((f_*(\mathcal L))_s)$ n'est pas une immersion, ce qui achève
la démonstration dans le cas de $E'$.

II) \emph{Cas de $E$.} --- Considérons d'abord le cas particulier où $Y=S$~:

\begin{corollary}[9.6.4]
\label{IV.9.6.4-fr}
Soit $f:X\to S$ un morphisme propre de présentation finie, $\mathcal L$ un
$\mathcal O_X$-Module inversible. Alors, l'ensemble $E$ des $s\in S$ tels que
$\mathcal L_s$ soit ample (relativement à $f_s$) est ouvert dans $S$, et
$\mathcal L|f^{-1}(E)$ est ample relativement à la restriction
$f^{-1}(E)\to E$ de $f$.
\end{corollary}

Comme la condition (9.2.1, (i)) est vérifiée par la propriété $\mathbf P$ définie plus
haut, le résultat de (9.2.2, (iv)) et le raisonnement de (9.2.3) montrent qu'on
peut se borner au cas où $S$ est \emph{noethérien}~; mais alors le résultat
découle de (\textbf{III}, 4.7.1) et de la stabilité de l'amplitude par
changement de base (\textbf{II}, 4.6.13).

\begin{corollary}[9.6.5]
\label{IV.9.6.5-fr}
Sous les hypothèses de (9.6.4), pour que $\mathcal L$ soit ample relativement
à $f$, il faut et il suffit que, pour tout $s\in S$, $\mathcal L_s$ soit ample
relativement à $f_s$.
\end{corollary}

\phantomsection\label{IV.9.6.6-fr}%
\textbf{(9.6.6) Fin de la démonstration de (9.6.3).} --- Revenons au cas
général, $S$ étant noethérien, intègre et de point générique
$\eta\in S-E$. Comme $f$ est propre, il résulte de (9.6.4) que l'ensemble $V$
des $y\in Y$ tels que $\mathcal L_y$ soit un $\mathcal O_{X_y}$-Module ample,
relativement au morphisme $f_y:X_y\to\operatorname{Spec}(k(y))$ est ouvert,
donc $F=Y-V$ est fermé dans $Y$. Cela étant, comme $f_s$ est propre et que,
pour tout $y\in Y$ au-dessus de $s\in S$, on a
$\mathcal L_y=(\mathcal L_s)_y$, il résulte de (9.6.5) que pour que $s$
appartienne à $S-E$, il faut et il suffit que l'on ait $F_s\neq\varnothing$.
Mais comme $F$ est fermé dans $Y$ et que $F_\eta\neq\varnothing$, il résulte
de (9.5.1) que l'ensemble des $s\in S$ tels que $F_s\neq\varnothing$ est un
voisinage de $\eta$ dans $S$. C.Q.F.D.

\begin{remarks}[9.6.7]
\label{IV.9.6.7-fr}
(i) Pour chacune des propriétés $\mathbf P$ considérées dans (9.6.1) la proposition
(9.3.3) est applicable, et ces propriétés (pour les morphismes $f_s$) sont donc
«~stables~» par passage d'une limite projective d'un système projectif
essentiellement affine (8.13.4) de préschémas $X_\lambda$ à l'un convenable
d'entre eux.

(ii) Soit $Z$ une partie \emph{localement constructible} de $X$ telle que, pour
tout $s\in S$, $Z\cap X_s$ soit \emph{ouvert} dans $X_s$, et désignons par
$Z_s$ le sous-préschéma de $X_s$ induit sur l'ouvert $Z\cap X_s$.

\oldpage[IV-3]{76}

Alors, dans les propositions (9.6.1) et (9.6.2, (I)), on peut partout remplacer
$f_s$ par sa restriction $f_s|Z_s:Z_s\to Y_s$ sans changer les conclusions. En
effet, la vérification de (9.2.1, (i)) se fait comme dans (9.6.1) et (9.6.2).
D'autre part, dans la réduction au cas où $S$ est noethérien, faite dans
(9.2.3), si $Z=q^{-1}(Z_0)$, où $q:X\to X_0$ est la projection canonique et
$Z_0$ une partie constructible de $X_0$ (8.3.11), le fait que
$(Z_0)_{s_0}$ soit ouvert dans $(X_0)_{s_0}$, pour $s_0=p(s)$, découle de
(2.4.10) et de ce que la projection $X_s\to(X_0)_{s_0}$ est surjective. On est
donc ramené à vérifier (9.2.1, (ii)) sous les nouvelles hypothèses. Or, comme
$Z_\eta$ est ouvert dans $X_\eta$, il existe un ouvert $Z'\subset X$ tel que
$Z_\eta=Z'\cap X_\eta$~; comme $X$ est alors noethérien, $Z'$ est
constructible, et il en est de même de $Z$ par hypothèse~; on conclut donc de
(9.5.2) et de ($\mathbf 0_{\mathrm{III}}$, 9.2.2) qu'il y a un voisinage $U$
de $\eta$ dans $S$ tel que $Z_s=Z'_s$ pour $s\in U$. Remplaçant $S$ par $U$,
on peut donc se borner au cas où $Z$ est ouvert dans $X$, et alors on est
ramené à ce qui a été prouvé dans (9.6.1) et (9.6.2).
\end{remarks}

\subsection{Constructibilité des propriétés de séparabilité, d'irréductibilité
géométrique et de connexité géométrique}
\label{subsection:IV.9.7-fr}

\begin{lemma}[9.7.1]
\label{IV.9.7.1-fr}
Soient $S$ un préschéma irréductible de point générique $\eta$, $f:X\to S$
un morphisme de présentation finie. Si $X_\eta$ a $n$ composantes irréductibles
(resp. $n'$ composantes connexes), il existe un voisinage ouvert $U$ de $\eta$
dans $S$ tel que pour tout $s\in U$, $X_s$ ait au moins $n$ composantes
irréductibles (resp. $n'$ composantes connexes).
\end{lemma}

On peut se borner au cas où $S$ est affine, donc $X$ quasi-compact et
quasi-séparé. Soient $V_i$ ($1\leqslant i\leqslant n$) les intérieurs des
composantes irréductibles de $X_\eta$, qui sont deux à deux disjoints et
quasi-compacts, et dont la réunion est dense dans $X_\eta$~; en vertu de
(8.2.11), appliqué suivant la méthode de (8.1.2, \emph{a)}), il existe pour
chaque $i$ un ouvert quasi-compact $W_i$ dans $X$ tel que
$W_i\cap X_\eta=V_i$~; comme $X$ est quasi-séparé, les intersections
$W_i\cap W_j$ sont quasi-compactes (1.2.7), donc, en remplaçant $S$
par un voisinage de $\eta$, on peut supposer que $W_i\cap W_j=\varnothing$ pour
$i\neq j$ en vertu de (8.3.3). En outre, comme les $W_i$ sont constructibles,
il y a un voisinage $U$ de $\eta$ dans $S$ tel que pour tout $s\in U$ les
$(W_i)_s$ soient non vides (9.5.1 et 9.2.3) et que la réunion des $(W_i)_s$ soit
\emph{dense} dans $X_s$ (9.5.3 et 9.2.3). Cela étant, les composantes
irréductibles (en nombre fini) des $(W_i)_s$ sont aussi les composantes
irréductibles de la réunion des $(W_i)_s$ $(\mathbf 0_{\mathrm I},2.1.7)$, donc
les adhérences dans $X_s$ de ces composantes sont les composantes
irréductibles de $X_s$ $(\mathbf 0_{\mathrm I},2.1.6)$ et leur nombre est
évidemment $\geqslant n$.

Soient maintenant $C_j$ ($1\leqslant j\leqslant n'$) les composantes connexes
de $X_\eta$, qui sont des ouverts quasi-compacts de $X_\eta$, deux à deux
disjoints. Si on remplace les $V_i$ par les $C_j$ dans le raisonnement
précédent, on voit (utilisant deux fois (8.3.3)) qu'on peut supposer $X$ réunion
de $n'$ ouverts quasi-compacts $M_j$ ($1\leqslant j\leqslant n'$) deux à deux
disjoints et que, dans un voisinage de $\eta$, les $(M_j)_s$ soient non vides.
Comme la réunion des $(M_j)_s$ est $X_s$, les composantes connexes des $(M_j)_s$
sont les composantes connexes de $X_s$, donc leur nombre est $\geqslant n'$.

\oldpage[IV-3]{77}

\begin{lemma}[9.7.2]
\label{IV.9.7.2-fr}
Soient $S$ un préschéma irréductible de point générique $\eta$, $f:X\to S$
un morphisme de présentation finie. Si $X_\eta$ n'est pas réduit, il existe un
voisinage ouvert $U$ de $\eta$ dans $S$ tel que, pour tout $s\in U$, $X_s$ ne
soit pas réduit.
\end{lemma}

En effet, soit $\mathcal N$ le Nilradical de $\mathcal O_X$~; il résulte de
(8.2.13), appliqué suivant la méthode de (8.1.2, \emph{a)}) que $\mathcal N_\eta$
est le nilradical de $\mathcal O_{X_\eta}$, et l'hypothèse est que
$\mathcal N_\eta\neq0$. On conclut de (9.4.5) et (9.2.3) qu'il y a un voisinage
$U$ de $\eta$ tel que, pour tout $s\in U$, $\mathcal N_s$ s'identifie à un Idéal
de $\mathcal O_{X_s}$ et $\mathcal N_s\neq0$~; comme $\mathcal N_s$ est
évidemment contenu dans le Nilradical de $\mathcal O_{X_s}$, on voit que $X_s$
n'est pas réduit pour $s\in U$.

\phantomsection\label{IV.9.7.3-fr}%
\textbf{(9.7.3)} Étant donné un polynôme
$F\in A[\mathrm T_1,\ldots,\mathrm T_n]$, où $A$ est un anneau et les
$\mathrm T_i$ des indéterminées, pour tout homomorphisme d'anneaux
$\rho:A\to B$, nous désignerons par $F_{(\rho)}$ ou $F_{(B)}$ le polynôme de
$B[\mathrm T_1,\ldots,\mathrm T_n]$ obtenu en remplaçant chaque coefficient de
$F$ par son image par $\rho$. Si $k$ est un corps,
$F\in k[\mathrm T_1,\ldots,\mathrm T_n]$ un polynôme non constant et
$X=\operatorname{Spec}(k[\mathrm T_1,\ldots,\mathrm T_n]/(F))$, dire que $X$
est intègre (ou que l'idéal $(F)$ est premier) signifie que $F$ est
\emph{irréductible} (c'est-à-dire que dans toute factorisation $F=F_1F_2$ en
polynômes de $k[\mathrm T_1,\ldots,\mathrm T_n]$, $F_1$ ou $F_2$ est de degré
$0$)~; cela résulte de ce que l'anneau $k[\mathrm T_1,\ldots,\mathrm T_n]$ est
factoriel. On en déduit aussitôt ((4.6.2) et (4.5.2))~:

\begin{lemma}[9.7.4]
\label{IV.9.7.4-fr}
Soient $k$ un corps, $\Omega$ une extension algébriquement close de $k$, $F$ un
polynôme non constant de $k[\mathrm T_1,\ldots,\mathrm T_n]$. Les conditions
suivantes sont équivalentes~:
\begin{enumerate}[label=\emph{\alph*)}]
  \item $X=\operatorname{Spec}(k[\mathrm T_1,\ldots,\mathrm T_n]/(F))$ est
  géométriquement intègre.
  \item $F_{(K)}$ est irréductible pour toute extension $K$ de $k$.
  \item $F_{(\Omega)}$ est irréductible.
\end{enumerate}
Nous dirons dans ce cas que $F$ est \emph{géométriquement irréductible}.
\end{lemma}

\begin{lemma}[9.7.5]
\label{IV.9.7.5-fr}
Soient $A$ un anneau intègre, $K$ son corps des fractions, $F$ un polynôme non
constant de $A[\mathrm T_1,\ldots,\mathrm T_n]$ de degré $d$ tel que $F_{(K)}$
soit géométriquement irréductible. Alors il existe $f\neq0$ dans $A$ tel que
pour tout $x\in D(f)$, $F_{(\boldsymbol{k}(x))}$ soit géométriquement
irréductible.
\end{lemma}

Posons $F=\sum_\alpha c_\alpha\mathrm T^\alpha$ avec comme d'ordinaire
$c_\alpha=c_{\alpha_1\alpha_2\ldots\alpha_n}$,
$\mathrm T^\alpha=\mathrm T_1^{\alpha_1}\mathrm T_2^{\alpha_2}\cdots
\mathrm T_n^{\alpha_n}$,
$|\alpha|=\alpha_1+\alpha_2+\cdots+\alpha_n\leqslant d$ (un au moins des
$c_\alpha$ tels que $|\alpha|=d$ n'étant pas nul). Comme les $c_\alpha$ non
nuls sont inversibles dans $K$, on peut supposer, en remplaçant au besoin $A$
par un anneau $A_g$ ($g\neq0$ dans $A$) que les $c_\alpha$ non nuls sont
inversibles dans $A$. Il en résulte que pour tout $x\in\operatorname{Spec}(A)$,
$F_{(\boldsymbol{k}(x))}$ est \emph{de degré $d$}.

Nous démontrerons d'abord un lemme préliminaire.

\begin{lemma}[9.7.5.1]
\label{IV.9.7.5.1-fr}
Soient $A$ un anneau intègre, $S=\operatorname{Spec}(A)$, $\eta$ le point
générique de $S$, $B$ l'anneau de polynômes
$A[\mathrm T_1,\ldots,\mathrm T_m]$, $(P_i)_{i\in I}$ une famille finie
d'éléments de $B$, $\mathfrak a$ l'idéal de $B$ engendré par les $P_i$. Pour
tout $s\in S$, soit $\mathfrak a_s$ l'idéal de
$\boldsymbol{k}(s)[\mathrm T_1,\ldots,\mathrm T_m]$ engendré par les polynômes
$(P_i)_{(\boldsymbol{k}(s))}$ ($i\in I$). Alors, si
$V(\mathfrak a_\eta)=\varnothing$, il existe un voisinage $U$ de $\eta$ dans
$S$ tel que $V(\mathfrak a_s)=\varnothing$ pour tout $s\in U$.
\end{lemma}

En effet, soit $Y=\operatorname{Spec}(B)$, et soit $Z$ la partie fermée
$V(\mathfrak a)$ de $Y$~; comme $\mathfrak a$ est un idéal de type fini, $Z$
est constructible dans $Y$ $(\mathbf 0_{\mathrm{III}},9.1.5)$ et avec les
notations introduites en (9.4.1), on a $V(\mathfrak a_s)=Z_s$ pour tout
$s\in S$~; comme le morphisme structural $Y\to S$ est de présentation finie,
la conclusion du lemme résulte de (9.5.1) et (9.2.3).

\oldpage[IV-3]{78}

Soit $(p,q)$ un couple d'entiers $>0$ tels que $p+q=d$~; introduisons des
indéterminées $\mathrm T'_\beta$, $\mathrm T''_\gamma$ pour tous les systèmes
d'entiers $\beta$, $\gamma$ tels que $|\beta|\leqslant p$ et
$|\gamma|\leqslant q$~; pour tout système d'entiers $\alpha$ tel que
$|\alpha|\leqslant d$, considérons le polynôme de
\[
B=A[\mathrm T'_\beta,\mathrm T''_\gamma]_{|\beta|\leqslant p,
  |\gamma|\leqslant q}:
\qquad
P_\alpha(\mathrm T'_\beta,\mathrm T''_\gamma)
=\sum_{\beta+\gamma=\alpha}\mathrm T'_\beta\mathrm T''_\gamma-c_\alpha.
\]
Soit $\Omega$ une clôture algébrique de $K$~; dire qu'il existe deux polynômes
$F_1$, $F_2$ de $\Omega[\mathrm T_1,\ldots,\mathrm T_n]$, de degrés respectifs
$p$ et $q$, tels que $F_1F_2=F_{(\Omega)}$, signifie que le système d'équations
$P_\alpha(t'_\beta,t''_\gamma)=0$ ($|\alpha|\leqslant d$) admet une solution
$(t'_\beta,t''_\gamma)$ ($|\beta|\leqslant p$, $|\gamma|\leqslant q$) formée
d'éléments de $\Omega$. Soit $\mathfrak a$ l'idéal de $B$ engendré par les
$P_\alpha$~; l'interprétation précédente, et le théorème des zéros de Hilbert,
montrent que l'hypothèse sur $F_{(K)}$ implique que
$V(\mathfrak a_\eta)=\varnothing$, en désignant par $\eta$ le point générique de
$\operatorname{Spec}(A)$~; le lemme (9.7.5.1) prouve donc que dans un voisinage
de $\eta$, on a $V(\mathfrak a_x)=0$, et par suite pour ces valeurs de $x$,
$F_{(\boldsymbol{k}(x))}$ n'admet pas de factorisation
$F_{(\boldsymbol{k}(x))}=G_1G_2$, où $G_1$, $G_2$ sont des polynômes de degrés
respectifs $p$ et $q$, dont les coefficients sont dans une clôture algébrique
de $\boldsymbol{k}(x)$. Il suffit d'appliquer ce résultat à tous les couples
d'entiers $(p,q)$ tels que $p>0$, $q>0$ et $p+q=d$, pour obtenir la conclusion
du lemme (9.7.5).

\begin{proposition}[9.7.6]
\label{IV.9.7.6-fr}
Soient $S$ un préschéma noethérien intègre de point générique $\eta$, $f:X\to S$
un morphisme de type fini, $\mathcal F$ un $\mathcal O_X$-Module cohérent. Si
$\mathcal F_\eta$ est sans cycle premier associé immergé, il existe un voisinage
$U$ de $\eta$ dans $S$ tel que, pour tout $s\in U$, $\mathcal F_s$ soit sans cycle
premier associé immergé.
\end{proposition}

Comme $X_\eta$ est noethérien, il existe un homomorphisme injectif
$v:\mathcal F_\eta\to\bigoplus_{i=1}^m\mathcal G_i$, où chacun des $\mathcal G_i$ est
irrédondant et $\operatorname{Ass}(\mathcal F_\eta)$ est l'ensemble des $x_i$, où
$\{x_i\}=\operatorname{Ass}(\mathcal G_i)$ (3.2.6)~; si $Z_i$ est l'adhérence de
$\{x_i\}$ dans $X_\eta$, on a $Z_i=\operatorname{Supp}(\mathcal G_i)$ (3.1.4), et
par hypothèse les $Z_i$ sont les composantes irréductibles de
$Z=\operatorname{Supp}(\mathcal F_\eta)$. Il résulte de (8.5.2, (i) et (ii)) (en
restreignant au besoin $S$ à un voisinage de $\eta$) et (8.5.8) qu'il existe des
$\mathcal O_X$-Modules cohérents $\mathcal F_i$ et un homomorphisme injectif
$u:\mathcal F\to\bigoplus_{i=1}^m\mathcal F_i$ tels que $(\mathcal F_i)_\eta=\mathcal G_i$ pour
tout $i$ et $v=u_\eta$. Si $Y_i=\operatorname{Supp}(\mathcal F_i)$, on a
$(Y_i)_s=\operatorname{Supp}((\mathcal F_i)_s)$ pour tout $s\in S$
(\textbf{I}, 9.1.13) et en particulier $(Y_i)_\eta=Z_i$ pour tout $i$.
L'hypothèse entraîne que pour $i\neq j$, $Z_i-(Z_i\cap Z_j)$ est ouvert dense
dans $Z_i$~; on déduit donc de (9.5.3) et (9.5.4) que dans un voisinage de
$\eta$, $(Y_i)_s-((Y_i)_s\cap(Y_j)_s)$ est ouvert et dense dans $(Y_i)_s$.
Supposons que l'on ait prouvé que chacun des $(\mathcal F_i)_s$ est sans cycle
premier associé pour $s\in U$. Alors les éléments de
$\operatorname{Ass}((\mathcal F_i)_s)$ sont les points maximaux de $(Y_i)_s$~;
aucun d'eux ne peut donc appartenir à un $(Y_j)_s$ pour $i\neq j$, et la
proposition sera démontrée. On peut donc supposer que $\mathcal F_\eta$ est
\emph{irrédondant}~; en outre, on peut se borner au cas où
$S=\operatorname{Spec}(A)$ est affine~; $X$ est alors réunion d'un nombre fini
d'ouverts affines $V_j$, et si $V_j\cap\operatorname{Supp}(\mathcal F_\eta)
=\varnothing$, on aura aussi $V_j\cap\operatorname{Supp}(\mathcal F_s)=\varnothing$
pour $s$ voisin de $\eta$ (9.5.1), donc on peut se borner au cas où
$X=\operatorname{Spec}(B)$ est aussi affine et où le morphisme $f:X\to S$ est
dominant~; $A$ est donc un anneau noethérien intègre, \emph{sous-anneau} d'une
$A$-algèbre de type fini $B$, et $\mathcal F=\widetilde M$, où $M$ est un
$B$-module de type fini~; par hypothèse, si $K$ est le corps des fractions de
$A$, le $B_{(K)}$-module $M_{(K)}$ est irrédondant. Soit $\mathfrak q$ l'unique

\oldpage[IV-3]{79}

élément de $\operatorname{Ass}(M_{(K)})$, et soit $\mathfrak p$ l'idéal premier
de $B$ image réciproque de $\mathfrak q$ par l'application canonique
$B\to B_{(K)}$. On sait (5.11.1.1) qu'il existe une filtration finie
$(N_h)_{0\leqslant h\leqslant m}$ de $M_{(K)}$ telle que $N_0=M_{(K)}$,
$N_m=0$, et que $N_h/N_{h+1}$ soit isomorphe à un sous-$B_{(K)}$-module
$\neq0$ de $B_{(K)}/\mathfrak q$. Soit $M_h$ l'image réciproque de $N_h$ par
l'application canonique $M\to M_{(K)}$. Il résulte de (9.4.4) que pour $s$
assez voisin de $\eta$, $(M_h)_{(\boldsymbol{k}(s))}$ s'identifie à un
sous-$B_{(\boldsymbol{k}(s))}$-module de $M_{(\boldsymbol{k}(s))}$, et le
quotient
$(M_h)_{(\boldsymbol{k}(s))}/(M_{h+1})_{(\boldsymbol{k}(s))}$ à un
sous-$B_{(\boldsymbol{k}(s))}$-module $\neq0$ de
$B_{(\boldsymbol{k}(s))}/\mathfrak p_{(\boldsymbol{k}(s))}
=(B/\mathfrak p)_{(\boldsymbol{k}(s))}$. Compte tenu de (3.1.7), on voit qu'on
est ramené à prouver, avec les mêmes notations, que si $B$ est \emph{intègre},
alors $B_{(\boldsymbol{k}(s))}$ n'a pas d'idéaux premiers associés immergés
pour $s$ voisin de $\eta$. Or, en remplaçant au besoin $A$ par $A_g$ et $B$ par
$B_g$ (où $g$ est un élément $\neq0$ de $A$), on peut supposer que $B$ contient
une $A$-algèbre de polynômes $C=A[\mathrm T_1,\ldots,\mathrm T_n]$ telle que
$B$ soit un $C$-module de type fini (Bourbaki, \emph{Alg. comm.}, chap. V,
\S~3, n$^\circ$~1, cor. 1 du th. 1). Comme $B_{(K)}$ est un $C_{(K)}$-module
sans torsion, on peut appliquer de nouveau le raisonnement fait plus haut en
remplaçant $B$, $M$ et $\mathfrak q$ par $C$, $B$ et $(0)$ respectivement, et
il suffit donc de voir que pour $s$ voisin de $\eta$, $C_{(\boldsymbol{k}(s))}$
n'a pas d'idéaux premiers associés immergés. Mais cela est évident puisque
$C_{(\boldsymbol{k}(s))}=\boldsymbol{k}(s)[\mathrm T_1,\ldots,\mathrm T_n]$
est un anneau intègre. C.Q.F.D.

\begin{theorem}[9.7.7]
\label{IV.9.7.7-fr}
Soit $f:X\to S$ un morphisme de présentation finie, et soit $E$ l'ensemble des
$s\in S$ pour lesquels $X_s$ a l'une des propriétés suivantes~: être~:
\begin{enumerate}[label=(\roman*)]
  \item géométriquement irréductible~;
  \item géométriquement connexe~;
  \item géométriquement réduit~;
  \item géométriquement intègre.
\end{enumerate}
Alors $E$ est localement constructible dans $S$.
\end{theorem}

Pour voir que les propriétés envisagées sont constructibles, remarquons d'abord
qu'elles vérifient trivialement la condition (9.2.1, (i)), en vertu de
(4.5.6, (i)) et (4.6.5, (i)). Il reste donc à vérifier (9.2.1, (ii)), donc on
peut supposer $S=\operatorname{Spec}(A)$ affine noethérien et intègre de point
générique $\eta$. En vertu de (4.6.8), il existe une extension finie $K'$ de
$K=\boldsymbol{k}(\eta)$ telle que $(X_\eta)_{(K')}$ soit tel que ses
composantes irréductibles (resp. connexes) soient géométriquement irréductibles
(resp. géométriquement connexes) et que $((X_\eta)_{(K')})_{\mathrm{red}}$ soit
géométriquement réduit. Comme il existe une base de $K'$ sur $K$ formée
d'éléments entiers sur $A$, l'anneau intègre $A'$ engendré par ces éléments est
fini sur $A$ et a $K'$ pour corps des fractions. Posons
$S'=\operatorname{Spec}(A')$~; le morphisme $g:S'\to S$ est fini et dominant,
donc surjectif (\textbf{II}, 6.1.10). Posons $X'=X_{(S')}$, de sorte que si
$\eta'$ est le point générique de $S'$, on a
$X'_{\eta'}=(X_\eta)_{(K')}$~; l'ensemble $E'$ des $s'\in S'$ tels que
$X'_{s'}$ ait l'une des propriétés (i), (ii), (iii) ou (iv) est égal à
$g^{-1}(E)$ (9.2.2, (iv)) ($E$ correspondant bien entendu à la même propriété)~;
$g$ étant surjectif, on a $E=g(E')$ et $S-E=g(S'-E')$~; en outre, $g$ est fermé
et $g^{-1}(\eta)=\{\eta'\}$ puisque $A'$ est intègre et fini sur $A$ (Bourbaki,
\emph{Alg. comm.}, chap. V, \S~2, n$^\circ$~1, cor. 1 de la prop. 1), donc
l'image par $g$ de tout voisinage de $\eta'$ est un voisinage de $\eta$. Le
théorème sera donc démontré si l'on prouve que $E'$ ou $S'-E'$ est un voisinage
de $\eta'$. Autrement dit, on peut désormais supposer que les composantes
irréductibles (resp. connexes) de $X_\eta$ sont géométriquement

\oldpage[IV-3]{80}

irréductibles (resp. géométriquement connexes) et que $(X_\eta)_{\mathrm{red}}$
est géométriquement réduit.

Supposons d'abord que $\eta\in S-E$ pour une des propriétés (i) à (iv). Si
$X_\eta$ n'est pas géométriquement irréductible (resp. géométriquement connexe),
il n'est pas irréductible (resp. connexe) en vertu de l'hypothèse précédente,
donc il en est de même de $X_s$ pour $s$ dans un voisinage de $\eta$ (9.7.1), et
\emph{a fortiori} dans ce voisinage $X_s$ n'est pas géométriquement irréductible
(resp. géométriquement connexe). D'autre part, si $X_\eta$ n'est pas
géométriquement réduit, il n'est pas réduit (sans quoi il serait égal à
$(X_\eta)_{\mathrm{red}}$ qui est géométriquement réduit par hypothèse)~; donc,
dans un voisinage de $\eta$, $X_s$ n'est pas réduit (9.7.2) et \emph{a fortiori}
n'est pas géométriquement réduit. Enfin, si $X_\eta$ n'est pas géométriquement
intègre, ou bien il n'est pas réduit, et alors on vient de voir que $X_s$ n'est
pas réduit (ni \emph{a fortiori} intègre) dans un voisinage de $\eta$~; ou bien
$X_\eta$ est réduit (donc géométriquement réduit par hypothèse), et alors il
n'est pas géométriquement irréductible, et on a vu ci-dessus qu'il en est alors
de même de $X_s$ pour $s$ voisin de $\eta$~; \emph{a fortiori} $X_s$ n'est pas
géométriquement intègre pour ces valeurs de $s$.

Nous allons donc désormais supposer que $\eta\in E$ et examiner séparément
chacune des propriétés considérées.

1$^\circ$ Supposons que $X_\eta$ soit géométriquement intègre. Soit $L$ le corps
des fonctions rationnelles sur $X_\eta$~; l'hypothèse sur $X_\eta$ implique que
$L$ est extension séparable de $K$ (4.6.3), donc extension séparable finie d'une
extension pure $K(\mathrm T_1,\ldots,\mathrm T_n)$ ($\mathrm T_i$
indéterminées)~; il y a par suite un élément $x\in L$, entier sur l'anneau
$K[\mathrm T_1,\ldots,\mathrm T_n]$, et tel que
$L=K(\mathrm T_1,\ldots,\mathrm T_n)(x)$~; soit
$G\in K[\mathrm T_1,\ldots,\mathrm T_n,\mathrm T_{n+1}]$ son polynôme minimal.
Il existe un élément $g\neq0$ de $A$ tel que tous les coefficients $\neq0$ de
$G$ (qui appartiennent à $K$) soient dans l'anneau $A_g$~; remplaçant $A$ par
$A_g$ (ce qui revient à remplacer $S$ par un voisinage de $\eta$), on peut donc
supposer que $G$ a ses coefficients dans $A$~; désignant par $F$ le polynôme
$G$ considéré comme élément de
$A[\mathrm T_1,\ldots,\mathrm T_n,\mathrm T_{n+1}]$, on a donc
$G=F_{(\boldsymbol{k}(\eta))}$. Posons
$Y=\operatorname{Spec}(A[\mathrm T_1,\ldots,\mathrm T_{n+1}]/(F))$, de sorte
que
$Y_\eta=\operatorname{Spec}(K[\mathrm T_1,\ldots,\mathrm T_{n+1}]/(G))$~;
$Y_\eta$ est un schéma intègre ayant $L$ pour corps des fonctions rationnelles.
Comme $X_\eta$ et $Y_\eta$ sont noethériens, il existe un ouvert non vide
$V\subset Y_\eta$ et une immersion ouverte $v:V\to X_\eta$ (nécessairement
dominante) (\textbf{I}, 6.5.1, (ii) et 6.5.4, (ii)). Soit $W$ un ouvert de $Y$
tel que $W\cap Y_\eta=V$~; appliquant (8.8.2, (i)) et (8.10.5, (iii)) suivant
la méthode de (8.1.2, \emph{a)}), on voit qu'en remplaçant au besoin $S$ par un
voisinage de $\eta$, on peut supposer que $v=w_\eta$, où $w:W\to X$ est une
immersion ouverte.

Cela étant, on a vu (4.6.3) que le critère pour qu'un préschéma algébrique
intègre soit géométriquement intègre ne dépend que de son corps de fonctions
rationnelles~; comme $X_\eta$ est géométriquement intègre par hypothèse, il en
est de même de $Y_\eta$, et la définition de $G$ entraîne donc que ce polynôme
est géométriquement irréductible (9.7.4). Appliquant (9.7.5), on voit qu'il y a
un voisinage $U$ de $\eta$ dans $S$ tel que $F_{(\boldsymbol{k}(s))}$ soit
géométriquement irréductible pour tout $s\in U$, donc $Y_s$ est géométriquement
intègre pour $s\in U$ (9.7.4)~; en outre on peut supposer que pour $s\in U$,
$W_s$ n'est pas vide (9.5.1), et par suite est géométriquement intègre (4.6.3)~;
enfin, on peut supposer aussi que

\oldpage[IV-3]{81}

pour $s\in U$, $w_s:W_s\to X_s$ est une immersion ouverte (9.6.1, (x)), et que
son image dans $X_s$ est partout dense (9.5.3). Autrement dit, pour $s\in U$,
il y a dans $X_s$ un ouvert partout dense $W'_s$ qui est géométriquement
intègre~; le critère (4.5.9, \emph{c)}) montre donc déjà que $X_s$ est
géométriquement irréductible pour $s\in U$. Enfin, comme $X_\eta$ est réduit et
par suite n'a pas de cycle premier associé immergé (3.2.1), on peut aussi
supposer que pour $s\in U$, $X_s$ n'a pas de cycle premier associé immergé
(9.7.6)~; soit alors (pour un $s\in U$ fixé) $\Omega$ une extension
algébriquement close de $\boldsymbol{k}(s)$, et soit
$p:(X_s)_{(\Omega)}\to X_s$ la projection canonique~; $p^{-1}(W'_s)$ est un
ouvert dense dans $(X_s)_{(\Omega)}$ (2.3.10) et est intègre par hypothèse~; en
outre $(X_s)_{(\Omega)}$ n'a pas de cycle premier associé immergé (4.2.7), donc
on conclut de (3.2.1) que $(X_s)_{(\Omega)}$ est réduit~; cela achève de prouver
que $X_s$ est géométriquement intègre (4.6.1).

2$^\circ$ Supposons que $X_\eta$ soit géométriquement irréductible~; comme
$X_{\mathrm{red}}$ est aussi de type fini sur $S$ (1.5.4, (vi)) on peut, en
tenant compte de (\textbf{I}, 5.1.8), remplacer $X$ par $X_{\mathrm{red}}$~;
alors $X_\eta$ est aussi intègre, et comme par hypothèse $X_\eta$ est
géométriquement réduit, il est géométriquement intègre. On est alors dans les
conditions de 1$^\circ$ et on en conclut (revenant aux hypothèses initiales) que
$X_s$ est géométriquement irréductible pour $s$ voisin de $\eta$.

3$^\circ$ Supposons $X_\eta$ géométriquement connexe, et soient $Z_i$
($1\leqslant i\leqslant n$) les composantes irréductibles de $X_\eta$~; il
existe (en vertu de (5.10.8.1) appliqué à $\mathfrak T=\{\varnothing\}$) une
application \emph{surjective} $j\mapsto\nu(j)$ d'un intervalle $[1,m]$ de
$\mathbf N$ sur $[1,n]$ telle que
$Z_{\nu(j)}\cap Z_{\nu(j+1)}\neq\varnothing$ pour $1\leqslant j\leqslant m$.
Pour tout $i$, soit $X_i$ l'adhérence de $Z_i$ dans $X$, et soit $Y$ la réunion
des $X_i$~; comme $Y_\eta=X_\eta$ par définition, on peut supposer, en vertu de
(9.5.1), que $Y_s=X_s$ pour tout $s\in S$, donc que $X_s$ est réunion des
$(X_i)_s$. Or, en considérant les sous-préschémas fermés réduits de $X$ ayant
les $X_i$ pour espaces sous-jacents, on voit par 2$^\circ$ qu'il existe un
voisinage $U$ de $\eta$ dans $S$ tel que pour $s\in U$ les $(X_i)_s$ soient
\emph{géométriquement irréductibles} (puisque les $(X_i)_\eta$ peuvent être
supposés géométriquement irréductibles, comme on l'a vu au début). En outre, on
peut aussi supposer que pour $s\in U$, on a
$(X_{\nu(j)})_s\cap(X_{\nu(j+1)})_s\neq\varnothing$ (9.5.1) pour
$1\leqslant j\leqslant m$~; on en conclut aussitôt que $X_s$ est connexe, donc
(4.5.13.1) géométriquement connexe pour $s\in U$.

4$^\circ$ Supposons $X_\eta$ géométriquement réduit~; soient $Z_i$ les
composantes irréductibles de $X_\eta$, $W_i$ l'intérieur de $Z_i$ dans
$X_\eta$~; il y a pour chaque $i$ un ouvert $V_i$ de $X$ tel que
$W_i=V_i\cap X_\eta$ pour tout $i$~; comme les $W_i$ sont ouverts et deux à deux
disjoints et que leur réunion est dense dans $X_\eta$, on peut (9.5.1, 9.5.3 et
9.5.4) supposer que pour $s$ voisin de $\eta$, les $(V_i)_s$ sont des ouverts
deux à deux disjoints dans $X_s$ et que leur réunion est dense dans $X_s$. En
outre, comme les $W_i$ sont géométriquement réduits et ont été supposés au
début géométriquement irréductibles, il résulte du 1$^\circ$ que pour $s$ voisin
de $\eta$, les $(V_i)_s$ sont géométriquement intègres, et \emph{a fortiori}
réduits. D'autre part, on tire de (9.7.6) que pour $s$ voisin de $\eta$, $X_s$
n'a pas de cycle premier associé immergé, puisqu'il en est ainsi de $X_\eta$
qui est réduit (3.2.1)~; on conclut donc de (3.2.1) que $X_s$ est réduit, et de
(4.6.1) qu'il est géométriquement réduit.

\oldpage[IV-3]{82}

Les parties de l'énoncé (9.7.7) concernant les propriétés (i) et (ii) se
généralisent comme suit~:

\begin{proposition}[9.7.8]
\label{IV.9.7.8-fr}
Soient $S$ un préschéma irréductible de point générique $\eta$, et $f:X\to S$
un morphisme de présentation finie. Soit $n$ (resp. $n'$) le nombre géométrique
de composantes irréductibles (resp. connexes) de $X_\eta$ (4.5.2). Alors il
existe un voisinage $U$ de $\eta$ dans $S$ tel que pour tout $s\in U$ le nombre
géométrique de composantes irréductibles (resp. connexes) de $X_s$ soit égal à
$n$ (resp. $n'$).
\end{proposition}

Compte tenu de ce que le nombre géométrique de composantes irréductibles (resp.
connexes) d'un préschéma algébrique est invariant par extension du corps de
base (4.5.6), on voit par la méthode de (9.2.3) que l'on peut se ramener au cas
où $S$ est affine, noethérien et intègre. En outre, en raisonnant comme au début
de la démonstration de (9.7.7), on voit qu'on peut supposer que les composantes
irréductibles (resp. connexes) de $X_\eta$ soient \emph{géométriquement}
irréductibles (resp. \emph{géométriquement connexes}). On sait déjà (9.7.1) que
pour tout $s$ voisin de $\eta$, $X_s$ a au moins $n$ composantes irréductibles
et $n'$ composantes connexes. D'autre part, si $Z_i$ (resp. $Z'_j$) sont les
composantes irréductibles (resp. connexes) de $X_\eta$, $X_i$ (resp. $X'_j$) le
sous-préschéma fermé réduit de $X$ ayant pour espace sous-jacent l'adhérence de
$Z_i$ (resp. $Z'_j$) dans $X$, il résulte de (9.5.1) que dans un voisinage de
$\eta$, les $(X_i)_s$ (resp. $(X'_j)_s$) forment un recouvrement de $X_s$ et que
les $(X'_j)_s$ sont deux à deux disjoints~; comme les $(X_i)_s$ (resp.
$(X'_j)_s$) sont géométriquement irréductibles (resp. géométriquement connexes)
en vertu de (9.7.7) pour $s$ voisin de $\eta$, on voit que $X_s$ a au plus $n$
composantes irréductibles et au plus $n'$ composantes connexes, d'où la
proposition.

\begin{corollary}[9.7.9]
\label{IV.9.7.9-fr}
Soit $f:X\to S$ un morphisme de présentation finie. Pour tout $s\in S$, soit
$n(s)$ (resp. $n'(s)$) le nombre géométrique de composantes irréductibles
(resp. connexes) de $f^{-1}(s)$ (4.5.2). Alors la fonction $s\mapsto n(s)$
(resp. $s\mapsto n'(s)$) est localement constructible dans $S$.
\end{corollary}

Il s'agit de montrer que la propriété «~$X$ est un préschéma algébrique sur un
corps $k$ et le nombre géométrique des composantes irréductibles (resp.
connexes) de $X$ est égal à $n$ (resp. $n'$)~» est constructible. Il résulte de
(4.5.6) que cette propriété vérifie la condition (9.2.1, (i)), et on est donc
ramené au cas où $S$ est affine, noethérien et intègre de point générique
$\eta$~; la conclusion résulte alors de (9.7.8).

\begin{corollary}[9.7.10]
\label{IV.9.7.10-fr}
Soit $X$ un préschéma localement noethérien tel que, si $X'$ est le normalisé de
$X_{\mathrm{red}}$, le morphisme canonique $f:X'\to X$ soit fini. Alors
l'ensemble des points $x\in X$ tels que $X$ soit géométriquement unibranche au
point $x$ est localement constructible dans $X$.
\end{corollary}

En effet, cet ensemble est par définition l'ensemble des points $x\in X$ tels
que le nombre de points géométriques de $f^{-1}(x)$ soit égal à 1. Mais comme
$f$ est fini, ce nombre est aussi le nombre géométrique des composantes
irréductibles de l'espace discret $f^{-1}(x)$ (compte tenu de la définition du
normalisé (\textbf{II}, 6.3.8) et de (4.5.11))~; la conclusion résulte donc de
(9.7.9).

\begin{remark}[9.7.11]
\label{IV.9.7.11-fr}
Soit $Z$ une partie localement constructible de $X$ telle que, pour tout
$s\in S$, $Z\cap X_s$ soit ouvert dans $X_s$, et désignons par $Z_s$ le
sous-préschéma de $X_s$ induit

\oldpage[IV-3]{83}

sur l'ouvert $Z\cap X_s$. Alors, dans le théorème (9.7.7), on peut remplacer
$X_s$ par $Z_s$ sans changer la conclusion~: on le voit en répétant le
raisonnement fait dans (9.6.7, (ii)).
\end{remark}

\begin{proposition}[9.7.12]
\label{IV.9.7.12-fr}
Soient $f:X\to S$ un morphisme de présentation finie, et $g:S\to X$ une
$S$-section de $X$ (\textbf{I}, 2.5.5). Pour tout $s\in S$, soit $X_s^0$ la
composante connexe de $f^{-1}(s)$ qui contient $g(s)$~; alors, la réunion $X^0$
des $X_s^0$ pour $s\in S$ est une partie localement constructible de $X$.
\end{proposition}

Montrons d'abord qu'on peut se ramener au cas où $S$ est affine et
\emph{noethérien}. On peut toujours supposer $S=\operatorname{Spec}(A)$ affine~;
avec les notations de (9.2.3), on a $f=(f_0)_{(S)}$, où $f_0:X_0\to S_0$ est un
morphisme de type fini, et on peut en outre supposer qu'il existe une
$S_0$-section $g_0:S_0\to X_0$ telle que $g=(g_0)_{(S)}$ (8.9.1). Remarquons
maintenant que si $p$ est le morphisme $S\to S_0$, alors, pour tout $s_0\in S_0$,
la composante connexe $(X_0)^0_{s_0}$ de $f_0^{-1}(s_0)$ qui contient $g_0(s_0)$
est \emph{géométriquement connexe} (4.5.13), et par suite, si $s_0=p(s)$, on a
$X_s^0=q^{-1}((X_0)^0_{s_0})$ où $q:X_s\to(X_0)_{s_0}$ est la projection
canonique ((4.5.8) et (4.4.1))~; notre assertion résulte donc de (1.8.2).

Utilisons alors le critère de constructibilité $(\mathbf 0_{\mathrm{III}},9.2.3)$~:
soient $x$ un point de $X$, $Z$ le sous-préschéma réduit de $X$ ayant
$\overline{\{x\}}$ comme espace sous-jacent, $Y$ le sous-préschéma réduit de $S$
ayant $\overline{f(Z)}$ comme espace sous-jacent~; comme la restriction de $f$
à $Z$ se factorise en $Z\to Y\to S$ (\textbf{I}, 5.2.2), on peut remplacer $S$
par $Y$, autrement dit supposer que $f(x)=\eta$, point générique du préschéma
intègre $S$. Par hypothèse, $X_\eta$ est somme de deux sous-préschémas
$X_\eta^0$, $X_\eta^1$, induits sur des ouverts complémentaires de $X_\eta$. En
vertu de (9.5.4) et (9.5.1), on peut donc, en remplaçant au besoin $S$ par un
voisinage de $\eta$, supposer que $X$ est réunion de deux ouverts disjoints
$X^0$, $X^1$ tels que $(X^0)_\eta=X_\eta^0$ et $(X^1)_\eta=X_\eta^1$. Comme
$g:S\to X$ est continu et injectif, $S$ est réunion des deux ouverts disjoints
$g^{-1}(X^0)$ et $g^{-1}(X^1)$~; mais comme $S$ est irréductible, et
\emph{a fortiori} connexe, un de ces deux ouverts est vide, et comme
$g(\eta)\in X_\eta^0$ par définition, on a $g^{-1}(X^1)=\varnothing$. En
d'autres termes, $g$ est une $S$-section de $X^0$~; d'autre part, comme
$(X^0)_\eta$ est géométriquement connexe (4.5.13), il en est de même de
$(X^0)_s$ pour tout $s$ voisin de $\eta$ (9.7.7)~; comme
$g(s)\in(X^0)_s$, on a bien $(X^0)_s=X_s^0$. C.Q.F.D.

\subsection{Décomposition primaire au voisinage d'une fibre générique}
\label{subsection:IV.9.8-fr}

\begin{proposition}[9.8.1]
\label{IV.9.8.1-fr}
Soient $f:X\to S$ un morphisme de présentation finie, $\mathcal F$ un
$\mathcal O_X$-Module quasi-cohérent de présentation finie. Alors l'ensemble
$E$ des $s\in S$ tels que $\mathcal F_s$ soit un $\mathcal O_{X_s}$-Module sans
cycle premier associé immergé est localement constructible.
\end{proposition}

On sait que pour les Modules quasi-cohérents sur les préschémas algébriques, la
propriété d'être sans cycle premier associé immergé est invariante par
changement du corps de base (4.2.7)~; on peut donc se borner au cas où $S$ est
affine, noethérien et intègre de point générique $\eta$, et à prouver que dans
ce cas $E$ ou $S-E$ est un voisinage de $\eta$. On a vu dans (9.7.6) que si
$\eta\in E$, $E$ est un voisinage de $\eta$~; reste donc à considérer le cas où
$\mathcal F_\eta$ a des cycles premiers associés immergés. On peut se borner au
cas où $X$ est affine, et où il y a un sous-$\mathcal O_{X_\eta}$-Module
cohérent $\mathcal H$ de $\mathcal F_\eta$ dont le support $Z'$ est non vide et
\emph{rare} par rapport au support $Z$ de $\mathcal F$ (3.1.3)~; alors, en vertu
de (8.5.2, (i) et (ii))

\oldpage[IV-3]{84}

et (8.5.8), appliqués suivant la méthode de (8.1.2, \emph{a)}), il existe un
sous-$\mathcal O_X$-Module cohérent $\mathcal G$ de $\mathcal F$ tel que
$\mathcal G_\eta=\mathcal H$~; si $Y=\operatorname{Supp}(\mathcal F)$,
$Y'=\operatorname{Supp}(\mathcal G)$, on a par suite
$Y_s=\operatorname{Supp}(\mathcal F_s)$, $Y'_s=\operatorname{Supp}(\mathcal G_s)$
pour tout $s\in S$ (\textbf{I}, 9.1.13) et en particulier $Z=Y_\eta$,
$Z'=Y'_\eta$~; comme $Z-Z'$ est dense dans $Z$ et $Z'\neq\varnothing$, il y a
un voisinage $U$ de $\eta$ tel que pour $s\in U$, $Y_s-Y'_s$ soit dense dans
$Y_s$, et $Y'_s\neq\varnothing$ (9.5.1 et 9.5.3)~; en considérant un point
générique d'une composante irréductible de $Y'_s$ et un voisinage assez petit de
ce point dans $X_s$, on déduit aussitôt de (3.1.3) que $\mathcal F_s$ admet des
cycles premiers associés immergés.

\phantomsection\label{IV.9.8.2-fr}%
\textbf{(9.8.2)} Soient $S$ un préschéma noethérien \emph{intègre}, de point
générique $\eta$, $f:X\to S$ un morphisme de type fini, $\mathcal F$ un
$\mathcal O_X$-Module cohérent. Considérons une décomposition irredondante
réduite $(\mathcal G_i)_{i\in I}$ de $\mathcal F_\eta$ (3.2.6)~; les
$\mathcal G_i$ sont donc des quotients de $\mathcal F_\eta$ et il y a un
homomorphisme injectif
$h:\mathcal F_\eta\to\bigoplus_i\mathcal G_i$~; en outre
$\operatorname{Ass}(\mathcal G_i)$ est réduit à un point $x_i$. Soit $Z_i$
l'adhérence de $\{x_i\}$ dans $X$, de sorte que
$(Z_i)_\eta=\operatorname{Supp}(\mathcal G_i)$. Désignons par $\mathcal J_i$
l'Idéal cohérent de $\mathcal O_X$ définissant le sous-préschéma fermé réduit de
$X$ d'espace sous-jacent $Z_i$ (sous-préschéma encore noté $Z_i$). En vertu de
(8.5.2) et (8.5.8), appliqués suivant la méthode de (8.1.2, \emph{a)}), on peut
(en restreignant au besoin $S$ à un voisinage de $\eta$) supposer qu'il existe
des quotients $\mathcal F_i$ de $\mathcal F$ ($i\in I$) tels que
$\mathcal G_i=(\mathcal F_i)_\eta$ pour tout $i$, et un homomorphisme
$g:\mathcal F\to\bigoplus_i\mathcal F_i$ tel que $h=g_\eta$. En outre
(\textbf{I}, 9.3.5), il existe un entier $m$ tel que
$(\mathcal J_i^m\mathcal F_i)_\eta=(\mathcal J_i)_\eta^m\mathcal G_i=0$ et en
restreignant de nouveau $S$, on peut donc supposer que l'on a aussi
$\mathcal J_i^m\mathcal F_i=0$ (8.5.2.5), donc que le support de $\mathcal F_i$
est contenu dans $Z_i$~; mais comme il est fermé et contient $x_i$, il est
nécessairement égal à $Z_i$.

\begin{proposition}[9.8.3]
\label{IV.9.8.3-fr}
Sous les conditions de (9.8.2), pour tout $s\in S$, et tout $i\in I$, désignons
par $x_{is\alpha}$ ($\alpha\in J_{s,i}$) les points maximaux de $(Z_i)_s$. Il
existe un voisinage $U$ de $\eta$ dans $S$ tel que, pour tout $s\in U$, les
$x_{is\alpha}$ (pour $i\in I$ et, pour chaque $i$, $\alpha\in J_{s,i}$) soient
deux à deux distincts et que $\operatorname{Ass}(\mathcal F_s)$ soit l'ensemble
des $x_{is\alpha}$ (autrement dit, les cycles premiers associés à $\mathcal F_s$
sont les composantes irréductibles des $(Z_i)_s$). En outre, on peut prendre
$U$ tel que, pour que l'adhérence de $\{x_{is\alpha}\}$ dans $X_s$ soit un cycle
premier associé maximal de $\mathcal F_s$, il faut et il suffit que $(Z_i)_\eta$
(adhérence de $\{x_i\}$ dans $X_\eta$) soit un cycle premier associé maximal de
$\mathcal F_\eta$.
\end{proposition}

Il résulte de (3.1.3, \emph{c')}) que pour chaque $i$, il existe un ouvert $W_i$
dans $X_\eta$ tel que $W_i\cap(Z_i)_\eta$ soit non vide, et un
$\mathcal O_{X_\eta}$-Module cohérent $\mathcal H_i$, de support
$W_i\cap(Z_i)_\eta$, tels qu'il y ait un homomorphisme injectif
$v_i:\mathcal H_i\to\mathcal F_\eta|W_i$. Soit $V_i$ un ouvert de $X$ tel que
$V_i\cap X_\eta=W_i$~; appliquant, comme dans (9.8.2), les résultats de (8.5.2)
et (8.5.8), on peut (en restreignant $S$) supposer qu'il existe un Module
cohérent $\mathcal F'_i$, de support $V_i\cap Z_i$ et un homomorphisme
$u_i:\mathcal F'_i\to\mathcal F|V_i$ tels que
$\mathcal H_i=(\mathcal F'_i)_\eta$ et $v_i=(u_i)_\eta$. Nous allons démontrer
qu'il y a un voisinage $U$ de $\eta$ tel que pour $s\in U$, les propriétés
suivantes soient vérifiées~:
\begin{enumerate}[label=\emph{(\roman*)}]
  \item Les $(\mathcal F_i)_s$ sont sans cycle premier associé immergé.
  \item L'homomorphisme
  $g_s:\mathcal F_s\to\bigoplus_i(\mathcal F_i)_s$ est injectif.
  \item $(Z_i)_s\cap(V_i)_s$ est dense dans $(Z_i)_s$,
  $(\mathcal F'_i)_s$ a pour support $(Z_i)_s\cap(V_i)_s$ et
  $(u_i)_s:(\mathcal F'_i)_s\to\mathcal F_s|(V_i)_s$ est injectif.
\end{enumerate}

\oldpage[IV-3]{85}

\begin{enumerate}[label=\emph{(\roman*)}]
  \item[(iv)] Pour $i\neq j$, toute composante irréductible de $(Z_i)_s$ est
  distincte de toute composante irréductible de $(Z_j)_s$.
\end{enumerate}
Or, (i) a déjà été vu (9.7.6)~; (ii) est un cas particulier de (9.4.5)~; (iii)
résulte de même de (9.5.3) et (9.4.5). Enfin, si $i\neq j$,
$(Z_i)_\eta\cap(Z_j)_\eta=(Z_i\cap Z_j)_\eta$ est rare dans $(Z_i)_\eta$ ou
dans $(Z_j)_\eta$~; supposons par exemple que $(Z_i\cap Z_j)_\eta$ soit rare
dans $(Z_j)_\eta$~; alors il résulte de (9.5.3) et (9.5.4) que pour $s$ voisin
de $\eta$, $(Z_i\cap Z_j)_s=(Z_i)_s\cap(Z_j)_s$ est rare dans $(Z_j)_s$, ce qui
montre qu'aucune composante irréductible de $(Z_j)_s$ ne peut être contenue
dans une composante irréductible de $(Z_i)_s$, ni \emph{a fortiori} lui être
égale.

Cela étant, il résulte de (ii) et de (3.1.7) que pour $s\in U$, on a
\[
\operatorname{Ass}(\mathcal F_s)\subset
\bigcup_i\operatorname{Ass}((\mathcal F_i)_s),
\]
de (i) que $\operatorname{Ass}((\mathcal F_i)_s)$ est l'ensemble des
$x_{is\alpha}$ ($\alpha\in J_{s,i}$). D'autre part, en vertu de (iii) et du
critère (3.1.3, \emph{c')}), chacun des $x_{is\alpha}$ ($\alpha\in J_{s,i}$,
$i\in I$) appartient à $\operatorname{Ass}(\mathcal F_s)$. Enfin (iv) signifie
que les $x_{is\alpha}$ sont deux à deux distincts.

Reste à prouver la dernière assertion de l'énoncé. Or, il résulte de (i) que
pour $s$ et pour un $i$ donnés, aucun des ensembles $\{x_{is\alpha}\}$ ne peut
être contenu dans un autre pour $\alpha\in J_{s,i}$. D'autre part, si
$(Z_j)_\eta\subset(Z_i)_\eta$, on peut supposer $U$ pris tel que
$(Z_j)_s\subset(Z_i)_s$ pour $s\in U$ (9.5.1), donc chaque $x_{js\beta}$
appartient à l'adhérence d'un $x_{is\alpha}$~; au contraire, si
$(Z_j)_\eta\cap(Z_i)_\eta$ est rare dans $(Z_j)_\eta$, on a vu en prouvant (iv)
que $(Z_j)_s\cap(Z_i)_s$ est rare dans $(Z_j)_s$, donc aucun des $x_{js\beta}$
n'est adhérent à un $x_{is\alpha}$. En particulier, si $(Z_j)_\eta$ est
maximal, ce qui équivaut à dire que $(Z_j)_\eta\cap(Z_i)_\eta$ est rare dans
$(Z_j)_\eta$ pour \emph{tout} $i\neq j$, on en conclut que
$(Z_j)_s\cap(Z_i)_s$ est rare dans $(Z_j)_s$ pour tout $i\neq j$, donc que tout
$x_{js\alpha}$ ($\alpha\in J_{s,j}$) est \emph{maximal}. C.Q.F.D.

\begin{corollary}[9.8.4]
\label{IV.9.8.4-fr}
Les notations et hypothèses étant celles de (9.8.2), il existe un voisinage $U$
de $\eta$ dans $S$ tel que, pour tout $s\in U$, chacun des $(\mathcal F_i)_s$
soit sans cycle premier associé immergé~; en outre, si
$(\mathcal F_{is\alpha})_{\alpha\in J_{s,i}}$ est l'unique décomposition
irredondante réduite de $(\mathcal F_i)_s$, alors, pour tout $s\in U$, la
famille $(\mathcal F_{is\alpha})_{i\in I,\,\alpha\in J_{s,i}}$ est une
décomposition irredondante réduite de $\mathcal F_s$.
\end{corollary}

La première assertion résulte de (9.8.1) et de la définition des
$\mathcal F_i$. D'autre part, on a vu dans (9.8.3) que l'homomorphisme
$\mathcal F_s\to\bigoplus_i(\mathcal F_i)_s$ est injectif et par définition il
en est de même de chacun des homomorphismes
$(\mathcal F_i)_s\to\bigoplus_\alpha\mathcal F_{is\alpha}$, donc
l'homomorphisme $\mathcal F_s\to\bigoplus_{i,\alpha}\mathcal F_{is\alpha}$ est
injectif. Comme on peut supposer (9.4.5) que pour $s\in U$ chacun des
$(\mathcal F_i)_s$ est un quotient de $\mathcal F_s$, les
$\mathcal F_{is\alpha}$ sont des quotients de $\mathcal F_s$, et il reste à
vérifier (3.2.5) que les $x_{is\alpha}$ sont deux à deux distincts et
appartiennent à $\operatorname{Ass}(\mathcal F_s)$, ce qui a été prouvé dans
(9.8.3).

\begin{proposition}[9.8.5]
\label{IV.9.8.5-fr}
Soient $f:X\to S$ un morphisme de présentation finie, $\mathcal F$ un
$\mathcal O_X$-Module quasi-cohérent de présentation finie. Pour tout $s\in S$,
soit $E(s)$ (resp. $E'(s)$) l'ensemble fini (partie de
$\mathbf Z\cup\{-\infty\}$) des dimensions des cycles premiers associés à
$\mathcal F_s$ (resp. des cycles premiers maximaux associés à $\mathcal F_s$,
c'est-à-dire les composantes irréductibles de $\operatorname{Supp}(\mathcal F_s)$).
Alors les fonctions $s\mapsto E(s)$ et $s\mapsto E'(s)$ sont localement
constructibles dans $S$.
\end{proposition}

\oldpage[IV-3]{86}

Il résulte de (4.2.7) et (4.2.8) que la condition (9.2.1, (i)) est satisfaite
pour les propriétés dont nous voulons démontrer qu'elles sont constructibles.
On est donc ramené au cas où $S$ est affine, noethérien et intègre, et à prouver
que $E$ et $E'$ sont constantes au voisinage du point générique $\eta$ de $S$~;
mais cela résulte de (9.8.3) et (9.5.5).

\begin{proposition}[9.8.6]
\label{IV.9.8.6-fr}
Avec les hypothèses et notations de (9.8.2), soit $i\in I$ tel que $(Z_i)_\eta$
soit maximal (autrement dit, une composante irréductible de
$\operatorname{Supp}(\mathcal F_\eta)$). Alors, il existe un voisinage $U$ de
$\eta$ dans $S$ tel que pour tout $s\in U$ et tout $\alpha\in J_{s,i}$, la
longueur géométrique de $\mathcal F_s$ en $x_{is\alpha}$ (relative à
$\boldsymbol{k}(s)$) (4.7.5) est égale à la longueur géométrique de
$\mathcal F_\eta$ en $x_i$ (relative à $\boldsymbol{k}(\eta)$).
\end{proposition}

On peut évidemment se borner au cas où $S=\operatorname{Spec}(A)$ est affine~;
montrons d'abord qu'on peut se ramener au cas où le sous-préschéma $(Z_i)_\eta$
de $X_\eta$, qui est réduit, est \emph{géométriquement intègre}. Il y a en effet
une extension finie $K'$ de $K=\boldsymbol{k}(\eta)$ telle que
$(((Z_i)_\eta)_{(K')})_{\mathrm{red}}$ soit géométriquement réduit et que les
composantes irréductibles de $((Z_i)_\eta)_{(K')}$ soient géométriquement
irréductibles (4.6.8). Procédons comme dans la démonstration de (9.7.7) en
considérant une sous-$A$-algèbre $A'$ de $K'$ ayant $K'$ pour corps des
fractions et finie sur $A$. On pose $S'=\operatorname{Spec}(A')$ et on considère
le morphisme fini surjectif $g:S'\to S$~; soient ensuite $X'=X_{(S')}$ et
$\mathcal F'=\mathcal F\otimes_{\mathcal O_S}\mathcal O_{S'}$, et soit $\eta'$
le point générique de $S'$. Pour tout $s'\in S'$, soit $s=g(s')$~; si $T$ est
une composante irréductible de $\operatorname{Supp}(\mathcal F_s)$, les
composantes irréductibles $T'_j$ de $T_{(\boldsymbol{k}(s'))}$ sont des
composantes irréductibles de $\operatorname{Supp}(\mathcal F'_{s'})$ et dominent
$T$ (4.2.7) et les multiplicités radicielles de $T$ par rapport à $\mathcal F_s$
et de chaque $T'_j$ par rapport à $\mathcal F'_{s'}$ sont les mêmes (4.7.9). Le
raisonnement de la première partie de (9.7.7) montre donc qu'on peut se borner à
démontrer la proposition pour $X'$ et $\mathcal F'$~; et en vertu du choix de
$K'$, les sous-préschémas réduits ayant pour espaces sous-jacents les composantes
irréductibles de $((Z_i)_\eta)_{(K')}$ sont géométriquement intègres (4.6.1).

Supposons donc désormais $(Z_i)_\eta$ géométriquement intègre~; alors (9.7.7),
il en est de même de $(Z_i)_s$ pour $s$ voisin de $\eta$~; la définition (4.7.5)
montre qu'il suffira donc de prouver que la \emph{longueur} du
$\mathcal O_{X_\eta,x_i}$-module $(\mathcal F_\eta)_{x_i}$ est égale à celle du
$\mathcal O_{X_s,x_{is}}$-module $(\mathcal F_s)_{x_{is}}$ (on a ici supprimé
l'indice $\alpha$, inutile par hypothèse). La question étant évidemment locale
sur $X$, on peut supposer (en se restreignant à un voisinage de $x_i$) que
$\mathcal F=\mathcal F_i$, donc que $\mathcal F_\eta$ est \emph{irredondant} sur
$X=\operatorname{Spec}(B)$ affine, et on écrira $Z$ au lieu de $Z_i$, et $x$
pour le point générique de $Z$ (et de $Z_\eta$). L'anneau noethérien $B$
contient donc $A$ comme sous-anneau, et $\mathcal F=\widetilde M$, où $M$ est un
$B$-module de type fini~; en outre, si $K$ est le corps des fractions de $A$,
le $B_{(K)}$-module $M_{(K)}$ est $\neq0$ et irredondant. Soit $q$ l'unique
élément de $\operatorname{Ass}(M_{(K)})$ et soit $\mathfrak p$ l'idéal premier
de $B$, image réciproque de $q$. En utilisant (5.11.1.1) comme dans la
démonstration de (9.7.6), on se ramène au cas où $B$ est \emph{intègre} et $M$
un sous-module $\neq0$ de $B$~; alors $\mathcal F_\eta$ est un
sous-$\mathcal O_{X_\eta}$-Module non nul de $\mathcal O_{Z_\eta}$, et en vertu
de (9.4.5), pour $s$ voisin de $\eta$, $\mathcal F_s$ est isomorphe à un
sous-$\mathcal O_{X_s}$-Module non nul de $\mathcal O_{Z_s}$~; comme $Z_s$ est
géométriquement intègre, les longueurs de $(\mathcal F_\eta)_x$ et de
$(\mathcal F_s)_{x_s}$ sont toutes deux égales à $1$, ce qui achève la
démonstration.

\oldpage[IV-3]{87}

\begin{corollary}[9.8.7]
\label{IV.9.8.7-fr}
Soient $f:X\to S$ un morphisme de présentation finie, $\mathcal F$ un
$\mathcal O_X$-Module quasi-cohérent de présentation finie. Pour tout $s\in S$,
soit $M(s)$ l'ensemble des couples $(d,m)$ tels qu'il existe une composante
irréductible de $\operatorname{Supp}(\mathcal F_s)$ de dimension $d$ et de
multiplicité radicielle $m$ pour $\mathcal F_s$ (4.7.8). Alors la fonction
$s\mapsto M(s)$ est localement constructible dans $S$.
\end{corollary}

Il résulte de (4.2.7) et (4.7.9) que la condition (9.2.1, (i)) est satisfaite
pour la propriété dont nous voulons démontrer qu'elle est constructible. On est
donc ramené au cas où $S$ est affine, noethérien et intègre et à prouver que $M$
est constante dans un voisinage du point générique de $S$~; mais alors la
proposition résulte de (9.8.3) et (9.8.6).

\begin{proposition}[9.8.8]
\label{IV.9.8.8-fr}
Soient $f:X\to S$ un morphisme de présentation finie, $\mathcal F$ un
$\mathcal O_X$-Module quasi-cohérent de présentation finie. Pour tout $s\in S$,
soit $l(s)$ la somme des multiplicités totales pour $\mathcal F_s$ (relatives à
$\boldsymbol{k}(s)$) des points génériques des composantes irréductibles de
$\operatorname{Supp}(\mathcal F_s)$ (4.7.12). Alors la fonction $s\mapsto l(s)$
est localement constructible dans $S$.
\end{proposition}

Compte tenu de (4.7.12), la condition (9.2.1, (i)) est satisfaite pour la
propriété dont nous voulons démontrer qu'elle est constructible, et on est donc
ramené au cas où $S$ est affine, noethérien et intègre de point générique
$\eta$, et à montrer que $l$ est constante dans un voisinage de $\eta$.
Utilisant les notations de (9.8.2), cela résulte de la définition (4.7.12), du
fait que le nombre géométrique des composantes irréductibles de chaque
$(Z_i)_s$ est constant dans un voisinage de $\eta$ (9.7.8), que la longueur
géométrique de $\mathcal F_s$ en $x_{is\alpha}$ est également à celle de
$\mathcal F_\eta$ en $x_i$ pour chaque $i$ tel que $(Z_i)_\eta$ soit maximal
(9.8.6) et enfin de ce que l'adhérence de $\{x_{is\alpha}\}$ est un cycle
premier associé maximal de $\mathcal F_s$ si et seulement si $(Z_i)_\eta$ est
un cycle premier associé maximal de $\mathcal F_\eta$ (9.8.3).

\begin{remark}[9.8.9]
\label{IV.9.8.9-fr}
On peut préciser de diverses façons les propositions précédentes~;
bornons-nous à un énoncé à titre d'exemple. Disons qu'une partie finie $P$ d'un
$k$-préschéma algébrique $X$ est \emph{saturée} si, pour tout couple de points
$x$, $y$ de $P$, les points génériques des composantes irréductibles de
$\overline{\{x\}}\cap\overline{\{y\}}$ appartiennent aussi à $P$~; pour toute
partie finie $Q$ de $X$, il existe une plus petite partie finie $P$ de $X$
contenant $Q$ et saturée~; nous dirons que $P$ est le \emph{saturé} de $Q$. Pour
tout $\mathcal O_X$-Module cohérent $\mathcal F$, nous appellerons
\emph{squelette primaire} de $\mathcal F$ le système $(P,Q,\omega,d,m)$ où
$Q=\operatorname{Ass}(\mathcal F)$, $P$ est le saturé de $Q$, $\omega$ la
relation d'ordre $\overline{\{x\}}\subset\overline{\{y\}}$ sur $P$, $d$ la
fonction $x\mapsto\dim\overline{\{x\}}$ sur $P$, $m$ la fonction
$x\mapsto\operatorname{long}_{\mathcal O_x}\mathcal F_x$ définie sur l'ensemble
des éléments de $Q$ maximaux pour la relation $\omega$. Appelons d'autre part
\emph{squelette virtuel} tout système $(P,Q,\omega,d,m)$, où $P$ est un
ensemble, $Q$ une partie de $P$, $\omega$ une relation d'ordre sur $P$, $d$ une
application de $P$ dans $\mathbf N$, $m$ une application dans $\mathbf N$ de
l'ensemble des éléments maximaux de $Q$~; on définit de façon évidente la notion
d'\emph{isomorphisme} de deux squelettes virtuels. Enfin, avec les notations
précédentes, appelons \emph{type primaire} de $\mathcal F$ la classe (pour la
relation d'isomorphie des squelettes virtuels) du squelette primaire de
$\mathcal F$. Il résulte de (4.2.6), (4.2.7), (4.2.8), (4.5.1) et (4.7.9) que
si, pour une extension algébriquement close $k'$ de $k$, on pose
$X'=X\otimes_k k'$ et $\mathcal F'=\mathcal F\otimes_k k'$, le
\emph{type primaire} de $\mathcal F'$ est indépendant de

\oldpage[IV-3]{88}

l'extension algébriquement close $k'$ de $k$ considérée~; nous dirons que c'est
le \emph{type primaire géométrique} de $\mathcal F$. Ces définitions posées,
l'énoncé que nous avons en vue est le suivant~:

\phantomsection\label{IV.9.8.9.1-fr}%
\textbf{(9.8.9.1)} Soient $f:X\to S$ un morphisme de présentation finie,
$\mathcal F$ un $\mathcal O_X$-Module quasi-cohérent de présentation finie.
Pour tout $s\in S$, soit $T(s)$ le type primaire géométrique de $\mathcal F_s$.
Alors la fonction $s\mapsto T(s)$ est localement constructible dans $S$.

Compte tenu des remarques qui précèdent, on est ramené comme d'ordinaire à
prouver que si $S$ est affine, noethérien et intègre de point générique $\eta$,
$T(s)$ est constante au voisinage de $\eta$. Raisonnant comme au début de
(9.7.7), on peut supposer que toutes les parties irréductibles de $X_\eta$ qui
interviennent sont géométriquement irréductibles, et alors la proposition
résulte de (9.5.1), (9.5.5), (9.8.3) et (9.8.6)~; nous laissons les détails au
lecteur. On pourrait généraliser en considérant plusieurs Modules cohérents et
en définissant leur «~squelette primaire simultané~», etc. La conclusion
générale de ce qui a été vu depuis le début de ce paragraphe est que pour toutes
les propriétés du type envisagé (et pour un $S$ irréductible) \emph{les
propriétés valables sur la «~fibre générique~» le sont encore sur toutes les
fibres voisines}.
\end{remark}

\subsection{Constructibilité des propriétés locales des fibres}
\label{subsection:IV.9.9-fr}

\begin{proposition}[9.9.1]
\label{IV.9.9.1-fr}
Soient $f:X\to S$ un morphisme localement de présentation finie, $Z$ une
partie localement constructible de $X$ telle que pour tout $s\in S$, $Z_s$ soit
fermé dans $X_s$, $\Phi$ une partie finie de $\mathbf Z\cup\{\pm\infty\}$.
Alors les parties suivantes de $X$ sont localement constructibles~:
\begin{enumerate}
\item[(i)] L'ensemble des $x\in X$ tels que
$\dim_x(Z_{f(x)})$ appartienne à $\Phi$.
\item[(ii)] L'ensemble des $x\in X$ tels que
$\operatorname{codim}_x(Z_{f(x)},X_{f(x)})$ appartienne à $\Phi$.
\item[(iii)] L'ensemble des $x\in X$ tels que l'anneau local
$\mathcal O_{X_{f(x)},x}$ soit équidimensionnel.
\end{enumerate}
\end{proposition}

On notera que les propriétés (i) et (ii) s'expriment encore en disant que les
fonctions $x\mapsto\dim_x(Z_{f(x)})$ et
$x\mapsto\operatorname{codim}_x(Z_{f(x)},X_{f(x)})$ sont
\emph{localement constructibles} dans $X$ ($\mathbf 0_{\mathrm{III}},9.3.1$).

Les questions étant locales sur $X$, on peut se borner au cas où
$S=\operatorname{Spec}(A)$ et $X=\operatorname{Spec}(B)$ sont affines et où
$f$ est un morphisme de présentation finie~; il existe alors un sous-anneau
$A_0$ de $A$ qui est une $\mathbf Z$-algèbre de type fini, un
$A_0$-préschéma de type fini $X_0$ et une partie constructible $Z_0$ de $X_0$
tels que $X=X_0\otimes_{A_0}A$ et $Z=h^{-1}(Z_0)$, où $h:X\to X_0$ est la
projection canonique ((8.9.1) et (8.3.11)). En outre, pour tout $s\in S$, si
$s_0$ est la projection de $s$ dans $S_0=\operatorname{Spec}(A_0)$, on a
\[
X_s=(X_0)_{s_0}\otimes_{k(s_0)}k(s),
\]
et $h_s$ est la projection $X_s\to(X_0)_{s_0}$, on a
$Z_s=h_s^{-1}((Z_0)_{s_0})$. Comme le morphisme $h_s$ est fidèlement plat et
quasi-compact, l'hypothèse que $Z_s$ est fermé dans $X_s$ entraîne que
$(Z_0)_{s_0}$ est fermé dans $(X_0)_{s_0}$ (2.3.12).

Cela étant, la transitivité des fibres (\textbf{I}, 3.6.4) et la prop. (4.2.7)
entraînent que l'ensemble des dimensions des composantes irréductibles de
$Z_s$ contenant $x$ est le même que l'ensemble des dimensions des composantes
irréductibles de $(Z_0)_{s_0}$ contenant $x_0=h_s(x)$. En particulier, on a
$\dim_x(Z_s)=\dim_{x_0}((Z_0)_{s_0})$. D'autre part, si $Z_s^{(\beta)}$ sont les
composantes irréductibles de $Z_s$ contenant $x$ et $X_s^{(\alpha)}$ les
composantes irréductibles de $X_s$

\oldpage[IV-3]{89}

contenant $x$, on a
\[
\operatorname{codim}_x(Z_s,X_s)=
\inf_\beta\left(\sup_\alpha
\bigl(\operatorname{codim}(Z_s^{(\beta)},X_s^{(\alpha)})\bigr)\right),
\]
$(\alpha,\beta)$ variant dans l'ensemble des couples tels que
$x\in Z_s^{(\beta)}\subset X_s^{(\alpha)}$ ($\mathbf 0$, 14.2.6). Comme les
préschémas algébriques irréductibles sont biéquidimensionnels (5.2.1), on peut
écrire, en vertu de ($\mathbf 0$, 14.3.3.1)
\begin{equation}
\tag{9.9.1.1}
\label{IV.9.9.1.1-fr}
\operatorname{codim}_x(Z_s,X_s)=
\inf_\beta\left(\sup_\alpha
\bigl(\dim(X_s^{(\alpha)})-\dim(Z_s^{(\beta)})\bigr)\right)
\end{equation}
avec le même choix des couples $(\alpha,\beta)$. Comme $h_s$ est fidèlement
plat et quasi-compact, pour tout couple formé d'une composante irréductible
$(X_0)_{s_0}^{(\lambda)}$ de $(X_0)_{s_0}$ contenant $x_0$ et d'une composante
irréductible $(Z_0)_{s_0}^{(\mu)}$ de $(Z_0)_{s_0}$ contenant $x_0$ et contenue
dans $(X_0)_{s_0}^{(\lambda)}$, il existe un couple
$(X_s^{(\alpha)},Z_s^{(\beta)})$ du type décrit plus haut et tel que
$Z_s^{(\beta)}$ domine $(Z_0)_{s_0}^{(\mu)}$ et que $X_s^{(\alpha)}$ domine
$(X_0)_{s_0}^{(\lambda)}$ (2.3.5). La formule (9.9.1.1) (et la formule analogue
appliquée à $(X_0)_{s_0}$) montrent alors, en vertu de (4.2.7), que l'on a
\[
\operatorname{codim}_x(Z_s,X_s)=
\operatorname{codim}_{x_0}((Z_0)_{s_0},(X_0)_{s_0}).
\]

On voit ainsi que si $E$ (resp. $E_0$) est l'ensemble des $x\in X$ (resp. des
$x_0\in X_0$) vérifiant l'une des conditions (i), (ii), (iii) de l'énoncé
(resp. la même condition), on a $E=h^{-1}(E_0)$, et en vertu de (1.8.2), on
voit qu'on peut se borner au cas où $A$ est \emph{noethérien}, donc aussi $B$.
Compte tenu de ($\mathbf 0_{\mathrm{III}},9.2.3$), ainsi que de (9.9.1.1), on
est ramené à voir que pour tout $x\in X$, il y a un voisinage $V$ de $x$ dans
$\overline{\{x\}}$ tel que, pour tout $x'\in V$, l'ensemble des dimensions des
composantes irréductibles de $X_{f(x')}$ (resp. $Z_{f(x')}$) contenant $x'$ est
\emph{le même}, et en outre qu'il en est de même de l'ensemble des couples
\[
\bigl(\dim(Z_{f(x')}^{(\beta)}),\dim(X_{f(x')}^{(\alpha)})\bigr)
\]
pour les couples formés d'une composante irréductible $X_{f(x')}^{(\alpha)}$
de $X_{f(x')}$ et d'une composante irréductible $Z_{f(x')}^{(\beta)}$ de
$Z_{f(x')}$ contenue dans $X_{f(x')}^{(\alpha)}$ et contenant $x'$. On peut
évidemment pour cela remplacer $S$ par le sous-préschéma réduit $S'$ de $S$
ayant $\overline{\{f(x)\}}$ pour espace sous-jacent, et $X$ par
$X'=f^{-1}(S')$, les fibres de $X$ et $X'$ aux points de $S'$ étant les mêmes.
Autrement dit, on peut se borner au cas où $S$ est \emph{intègre} et où
$\eta=f(x)$ est son point générique.

Par hypothèse $Z_\eta$ est fermé dans $X_\eta$~; comme $Z$ est constructible,
il résulte de (8.3.11), appliqué suivant la méthode de (8.1.2, a)), que l'on
peut, en remplaçant au besoin $S$ par un voisinage ouvert de $\eta$, supposer
que $Z$ est \emph{fermé} dans $X$. Soient $X_i$ (resp. $Z_j$) les composantes
irréductibles de $X$ (resp. $Z$) contenant $x$~; en vertu de
($\mathbf 0_{\mathrm I},2.1.8$), les $X_i\cap X_\eta$ (resp.
$Z_j\cap X_\eta$) sont les composantes irréductibles de $X_\eta$ (resp.
$Z_\eta$) contenant $x$~; en vertu de (9.5.1), on peut supposer en outre, en
restreignant au besoin $S$ à un voisinage de $\eta$, que les $X_i$ (resp.
$Z_j$) sont exactement les composantes irréductibles de $X$ (resp. $Z$)
rencontrant $\overline{\{x\}}$ et que les
$(X_i)_s\cap\overline{\{x\}}$ et $(Z_j)_s\cap\overline{\{x\}}$ sont non vides
pour tout $s\in S$. Cela étant, il résulte encore de (9.7.1) que l'on peut
supposer, en restreignant $S$, que les couples $(i,j)$ tels que
$(Z_j)_s\subset(X_i)_s$ sont les mêmes pour tout $s\in S$. La conclusion
résulte alors de (9.5.6)~: pour tout $s$ assez voisin de $\eta$, toutes les
composantes irréductibles de $(X_i)_s$ (resp. $(Z_j)_s$) ont une même dimension
égale à celle de $(X_i)_\eta$ (resp. $(Z_j)_\eta$). En outre, si $(i,j)$ est un
couple tel que $Z_j\not\subset X_i$, $(X_i)_\eta$ ne contient pas le

\oldpage[IV-3]{90}

point générique de $(Z_j)_\eta$, donc
$\dim((X_i\cap Z_j)_\eta)<\dim((Z_j)_\eta)$~; par suite, la dimension commune
des composantes irréductibles de $(X_i)_s\cap(Z_j)_s$ est, pour tout $s\in S$,
strictement inférieure à la dimension commune des composantes irréductibles de
$(Z_j)_s$, ce qui prouve qu'aucune des composantes irréductibles de $(Z_j)_s$
n'est contenue dans une composante irréductible de $(X_i)_s$ pour $s\in S$.
C.Q.F.D.

\begin{proposition}[9.9.2]
\label{IV.9.9.2-fr}
Soient $f:X\to S$ un morphisme localement de présentation finie, $\mathcal F$
un $\mathcal O_X$-Module quasi-cohérent de présentation finie, $\Phi$ une
partie finie de $\mathbf Z\cup\{\pm\infty\}$. Les parties suivantes de $X$
sont localement constructibles~:
\begin{enumerate}
\item[(i)] L'ensemble des points $x\in X$ tels que $\mathcal F_{f(x)}$ soit
équidimensionnel au point $x$ (5.1.12).
\item[(ii)] L'ensemble des $x\in X$ tels que les longueurs géométriques de
$\mathcal F_{f(x)}$ relatives à $k(f(x))$, aux points génériques des composantes
irréductibles de $\operatorname{Supp}(\mathcal F_{f(x)})$ contenant $x$
(4.7.5), appartiennent à $\Phi$.
\item[(iii)] L'ensemble des $x\in X$ tels que les dimensions des cycles
premiers associés à $\mathcal F_{f(x)}$ et contenant $x$ appartiennent à
$\Phi$.
\item[(iv)] L'ensemble des $x\in X$ tels que $\mathcal F_{f(x)}$ soit
géométriquement réduit au point $x$ (4.6.17).
\item[(v)] L'ensemble des $x\in X$ tels que $\mathcal F_{f(x)}$ soit
géométriquement intègre au point $x$ (4.6.22).
\item[(vi)] L'ensemble des $x\in X$ tels que
$\operatorname{dim.proj}((\mathcal F_{f(x)})_x)\in\Phi$.
\item[(vii)] L'ensemble des $x\in X$ tels que
$\operatorname{coprof}((\mathcal F_{f(x)})_x)\in\Phi$.
\item[(viii)] L'ensemble des $x\in X$ tels que $\mathcal F_{f(x)}$ possède la
propriété $(S_k)$ au point $x$ (5.7.2).
\item[(ix)] L'ensemble des $x\in X$ tels que $\mathcal F_{f(x)}$ soit un
$\mathcal O_{X_{f(x)}}$-Module de Cohen-Macaulay au point $x$ (5.7.1).
\end{enumerate}
\end{proposition}

(i) Le support $Z$ de $\mathcal F$ est localement constructible et fermé dans
$X$ (8.9.1), et en considérant un sous-préschéma de $X$ ayant $Z$ pour espace
sous-jacent et qui soit de présentation finie sur $S$ (8.9.1), on voit que
l'assertion (i) est un cas particulier de (9.9.1, (iii)).

(ii) Toutes les propriétés envisagées sont locales sur $X$, et on se bornera
donc au cas où $X=\operatorname{Spec}(B)$ et $S=\operatorname{Spec}(A)$ sont
affines et $f$ un morphisme de présentation finie. On garde les notations du
début de la démonstration de (9.9.1), et on suppose en outre $A_0$ choisi de
sorte qu'il existe un $\mathcal O_{X_0}$-Module cohérent $\mathcal F_0$ tel que
$\mathcal F$ soit isomorphe à $\mathcal F_0\otimes_{A_0}A$. Alors ((4.2.7) et
(4.7.9)) l'ensemble des longueurs géométriques de
$(\mathcal F_0)_{f(x_0)}$ aux points génériques des composantes irréductibles de
$\operatorname{Supp}((\mathcal F_0)_{f(x_0)})$ qui contiennent $x_0$ est le
même que l'ensemble analogue pour $x$ et $\mathcal F_{f(x)}$~; autrement dit,
si $E$ (resp. $E_0$) est l'ensemble des $x\in X$ (resp. des $x_0\in X_0$)
vérifiant la condition (ii) de l'énoncé, on a $E=h^{-1}(E_0)$, et en vertu de
(1.8.2), on voit qu'on peut se borner à considérer le cas où $A$ est
\emph{noethérien}. Comme dans la démonstration de (9.9.1), on voit qu'on est
ramené à montrer que, pour tout $x\in X$, il existe un voisinage $V$ de $x$
dans $\overline{\{x\}}$ tel que, pour tout $x'\in V$, l'ensemble des longueurs
géométriques de $\mathcal F_{f(x')}$ aux points génériques des composantes
irréductibles de son support contenant $x'$ soit \emph{le même}. En outre, si
$S'$ est le sous-préschéma réduit de $S$ ayant $\overline{\{f(x)\}}$ pour espace
sous-jacent, et si $X'=f^{-1}(S')$, les fibres de $X$ et de $X'$ aux points de
$S'$ sont les mêmes,

\oldpage[IV-3]{91}

donc on peut remplacer $S$ par $S'$ et $X$ par $X'$, autrement dit supposer que
$S$ est \emph{intègre} et que $\eta=f(x)$ est son point générique.

Or, si l'on pose $Z=\operatorname{Supp}(\mathcal F)$, le même raisonnement que
dans la démonstration de (9.9.1) montre que, si $Z_i$ sont les composantes
irréductibles de $Z$ contenant $x$, on peut supposer que ce sont exactement les
composantes irréductibles de $Z$ rencontrant $\overline{\{x\}}$ et que
$(Z_i)_s\cap\overline{\{x\}}\neq\varnothing$ pour tout $s\in S$. La conclusion
résulte alors de (9.8.3) et (9.8.6).

(iii) On se ramène comme dans (ii) au cas où $S$ est noethérien et intègre et
$\eta=f(x)$ son point générique, en utilisant (4.2.7). Comme dans la
démonstration de (9.9.1), on est ramené à montrer qu'il existe un voisinage $V$
de $x$ dans $\overline{\{x\}}$ tel que, pour tout $x'\in V$, l'ensemble des
dimensions des cycles premiers associés à $\mathcal F_{f(x')}$ qui contiennent
$x'$ soit \emph{le même}. Or, si $Z'_i$ sont les adhérences dans $X$ des cycles
premiers associés à $\mathcal F_\eta$ qui contiennent $x$ (cf. (9.8.2)), il
résulte de (9.8.3) et (9.5.1) que pour tout $s$ assez voisin de $\eta$, les
cycles premiers associés à $\mathcal F_s$ qui rencontrent
$\overline{\{x\}}$ sont des composantes irréductibles des $(Z'_i)_s$ et, en
vertu de (9.5.6), toutes ces composantes irréductibles ont même dimension égale
à $\dim((Z'_i)_\eta)$, d'où la conclusion.

(iv) On raisonne comme dans (iii), en utilisant cette fois (4.7.11). Avec les
mêmes notations que dans (iii), les cycles premiers associés à $\mathcal F_s$
qui sont des composantes irréductibles de $(Z'_i)_s$ sont immergés si et
seulement si $(Z'_i)_\eta$ est un cycle premier associé immergé de
$\mathcal F_\eta$. On conclut déjà que si $x$ appartient à un (resp.
n'appartient à aucun) cycle premier associé immergé de $\mathcal F_\eta$,
l'ensemble des $x'\in\overline{\{x\}}$ tels que $x'$ appartienne à un (resp.
n'appartienne à aucun) cycle premier associé immergé de $\mathcal F_{f(x')}$
est un voisinage de $x$ dans $\overline{\{x\}}$. La conclusion résulte de cette
remarque, de la caractérisation des points où un Module est géométriquement
réduit (4.7.10), et de (ii).

(v) Compte tenu de (4.7.11), on se ramène encore au cas où $S$ est noethérien
et intègre et où $\eta=f(x)$ est son point générique. Soit $Y$ un
sous-préschéma fermé de $X$ ayant $\operatorname{Supp}(\mathcal F)$ pour espace
sous-jacent. Raisonnant comme au début de la démonstration de (9.7.7), on voit
qu'il existe une $A$-algèbre intègre finie $A'$ (de sorte que si
$S'=\operatorname{Spec}(A')$, le morphisme $S'\to S$ est fini et surjectif)
telle que si $Y'=Y_{(S')}$ et si $\eta'$ est le point générique de $S'$, les
composantes irréductibles de $Y'_{\eta'}$ soient géométriquement irréductibles.
D'autre part, si $X'=X_{(S')}$, le morphisme projection $h:X'\to X$ est fini
et surjectif, donc fermé~; par suite, si $x'$ est un point de $X'$ au-dessus de
$x$, on a $h(\overline{\{x'\}})=\overline{\{x\}}$ et l'image par $h$ d'un
voisinage de $x'$ dans $\overline{\{x'\}}$ est un voisinage de $x$ dans
$\overline{\{x\}}$. Compte tenu de (4.7.11) et posant
$\mathcal F'=\mathcal F_{(S')}$, on est donc ramené à prouver que si
$\mathcal F'_{\eta'}$ est (resp. n'est pas) géométriquement intègre au point
$x'$, l'ensemble des $y'\in\overline{\{x'\}}$ tels que
$\mathcal F'_{f'(y')}$ (où $f'=f_{(S')}:X'\to S'$) soit (resp. ne soit pas)
géométriquement intègre au point $y'$ est un voisinage de $x'$ dans
$\overline{\{x'\}}$. On peut donc se borner au cas où $X'=X$ et où les
composantes irréductibles de $Y_\eta$ sont géométriquement irréductibles~; si
l'on désigne par $Z_i$ des parties fermées de $X$ telles que les
$Z_i\cap X_\eta$ soient les

\oldpage[IV-3]{92}

composantes irréductibles de $Y_\eta$, il résulte de (9.7.7), (9.7.8) et
(9.5.3) que pour tout $s$ voisin de $\eta$, les $(Z_i)_s$ sont les composantes
irréductibles de $Y_s$ et qu'elles sont \emph{géométriquement irréductibles}.
Dire qu'en un point $y\in X_s$, $\mathcal F_s$ est géométriquement intègre
signifie alors que $\mathcal F_s$ est géométriquement réduit en ce point et en
outre que $y$ n'appartient qu'à \emph{un seul} des $(Z_i)_s$ (4.6.22). La
conclusion résulte donc d'une part de (iv) et d'autre part de (9.5.1) appliqué à
l'intersection de $\overline{\{x\}}$ et de chaque $Z_i$.

(vi) Gardant les mêmes notations que dans (ii), il résulte de (6.2.1) que l'on
a
\[
\operatorname{dim.proj}((\mathcal F_s)_x)=
\operatorname{dim.proj}(((\mathcal F_0)_{s_0})_{x_0})~;
\]
on peut donc encore se borner au cas où $A$ est \emph{noethérien}. En outre, on
se ramène encore à montrer que, pour tout $x\in X$, il existe un voisinage $V$
de $x$ dans $\overline{\{x\}}$ tel que, pour tout $x'\in V$, on ait
\[
\operatorname{dim.proj}((\mathcal F_{f(x')})_{x'})=
\operatorname{dim.proj}((\mathcal F_{f(x)})_x)~;
\]
et comme ci-dessus, on peut supposer que $S$ est \emph{intègre} et que
$\eta=f(x)$ est son point générique, de sorte que l'on a
$(\mathcal F_\eta)_x=\mathcal F_x$. En vertu du théorème de platitude générique
(6.9.1), on peut, en remplaçant au besoin $S$ par un voisinage ouvert de
$\eta$, supposer que le morphisme $f$ est \emph{plat} et que $\mathcal F$ est
$f$-plat~; on a alors
$\operatorname{dim.proj}((\mathcal F_{f(x')})_{x'})=
\operatorname{dim.proj}(\mathcal F_{x'})$ pour tout $x'\in X$ en vertu de
(6.2.3). Cela étant, en vertu de (6.11.1), on peut (en remplaçant au besoin $X$
par un voisinage ouvert de $x$) supposer que
$\operatorname{dim.proj}(\mathcal F_{x'})\leq
\operatorname{dim.proj}(\mathcal F_x)$ pour tout $x'\in X$. D'autre part, si
$\operatorname{dim.proj}(\mathcal F_x)=n$, il y a par hypothèse un
$\mathcal O_x$-module de type fini $M$ tel que
$\operatorname{Ext}_{\mathcal O_x}^n(\mathcal F_x,M)\neq0$
($\mathbf 0$, 17.2.4). Or, il existe un $\mathcal O_X$-Module cohérent
$\mathcal G$ tel que $M=\mathcal G_x$ (en remplaçant au besoin $X$ par un
voisinage ouvert de $x$ ($\mathbf 0_{\mathrm I}$, 5.3.8))~; en vertu de
(T, 4.2.2), on a donc
$(\mathcal{E}xt_{\mathcal O_X}^n(\mathcal F,\mathcal G))_x\neq0$. Mais
$\mathcal{E}xt_{\mathcal O_X}^n(\mathcal F,\mathcal G)$ est un
$\mathcal O_X$-Module cohérent ($\mathbf 0_{\mathrm{III}}$, 12.3.3), donc son
support $Y$ est fermé ($\mathbf 0_{\mathrm I}$, 5.2.2)~; comme il contient $x$,
il contient aussi $\overline{\{x\}}$, d'où l'on conclut (en appliquant encore
(T, 4.2.2)) que l'on a
$\operatorname{dim.proj}(\mathcal F_{x'})\geq n$ pour tout
$x'\in\overline{\{x\}}$, ce qui achève de prouver l'assertion dans le cas (vi).

(vii) Comme $B$ est une $A$-algèbre de type fini, $X$ est $S$-isomorphe à un
\emph{sous-schéma fermé} d'un $S$-schéma de la forme
$Y=\operatorname{Spec}(A[T_1,\ldots,T_r])$~; si $j:X\to Y$ est l'injection
canonique, et $\mathcal G=j_*(\mathcal F)$, on a
$\mathcal G_s=(j_s)_*(\mathcal F_s)$ pour tout $s\in S$, et, en vertu de
($\mathbf 0$, 16.4.11), on peut se borner à démontrer l'assertion pour $Y$ et
$\mathcal G$. Autrement dit, on peut supposer que $B=A[T_1,\ldots,T_r]$, de
sorte que chacun des schémas $X_s=\operatorname{Spec}(k(s)[T_1,\ldots,T_r])$
est \emph{régulier} ($\mathbf 0$, 17.3.7). Soit alors
$W=\operatorname{Supp}(\mathcal F)$, de sorte que
$W_s=\operatorname{Supp}(\mathcal F_s)$ (\textbf{I}, 9.1.13)~; on a, par
(6.11.2.1)
\begin{equation}
\tag{9.9.2.1}
\label{IV.9.9.2.1-fr}
\operatorname{coprof}((\mathcal F_{f(x)})_x)=
\operatorname{dim.proj}((\mathcal F_{f(x)})_x)-
\operatorname{codim}_x(W_{f(x)},X_{f(x)}).
\end{equation}
Mais comme $W$ est constructible (8.9.1) et chaque $W_s$ fermé, il résulte de
(vi) et de (9.9.1, (ii)) que les deux fonctions du second membre de (9.9.2.1)
sont constructibles~; il en est donc de même de leur différence ce qui achève
de prouver la proposition dans le cas (vii).

(viii) Soit $U_n$ l'ensemble des $x\in X$ tels que
$\operatorname{coprof}((\mathcal F_{f(x)})_x)\leq n$, et posons
$Z_n=X-U_n$~; il résulte de (vii) que les $Z_n$ sont constructibles~; en outre,
comme

\oldpage[IV-3]{93}

la fonction $x\mapsto\dim_x(W_{f(x)})$ est constructible en vertu de
(9.9.1, (i)), elle ne prend qu'un nombre fini de valeurs, donc les nombres
$\dim(W_{f(x)})$ ont une borne supérieure finie $m$ lorsque $x$ parcourt $X$~;
comme
$\operatorname{coprof}((\mathcal F_{f(x)})_x)\leq
\dim((\mathcal F_{f(x)})_x)\leq\dim(W_{f(x)})$, on voit que l'on a
$Z_n=\varnothing$ pour $n\geq m$. Enfin, il résulte de (6.11.2, (i)) que pour
tout $n$ et tout $s\in S$, $(Z_n)_s$ est \emph{fermé} dans $X_s$. D'après
(5.7.4), l'ensemble des $x\in X$ où $(\mathcal F_{f(x)})_x$ possède la
propriété $(S_k)$ est l'ensemble des $x\in X$ vérifiant toutes les relations
\begin{equation}
\tag{9.9.2.2}
\label{IV.9.9.2.2-fr}
\operatorname{codim}_x((Z_n)_{f(x)},W_{f(x)})>n+k
\end{equation}
pour tout $n\geq0$~; comme cette relation est automatiquement vérifiée pour
$n\geq m$, on n'a à considérer les relations (9.9.2.2) que pour $0\leq n<m$.
Mais en vertu de (9.9.1, (ii)), l'ensemble $V_{n,k}$ des $x$ vérifiant (9.9.2.2)
est constructible, et il en est de même de l'intersection de ces ensembles pour
$0\leq n<m$.

(ix) L'assertion résulte ici aussitôt de (vii), l'ensemble des $x\in X$ où
$(\mathcal F_{f(x)})_x$ est un module de Cohen-Macaulay étant défini par la
relation $\operatorname{coprof}((\mathcal F_{f(x)})_x)=0$.

\begin{corollary}[9.9.3]
\label{IV.9.9.3-fr}
Soient $f:X\to S$ un morphisme de présentation finie, $\mathcal F$ un
$\mathcal O_X$-Module de présentation finie, $\mathbf P$ l'une des propriétés
(i) à (ix) de (9.9.2). Alors l'ensemble des $s\in S$ tels que la propriété
$\mathbf P$ soit vraie en tous les points $x\in X_s$, est localement
constructible dans $S$.
\end{corollary}

En effet, son complémentaire est l'image par $f$ du complémentaire de
l'ensemble $E$ des points où $\mathbf P$ est vraie. Comme $E$ est localement
constructible dans $X$, il en est de même de $X-E$, donc $f(X-E)$ est
localement constructible dans $S$ en vertu du théorème de Chevalley (1.8.4).

\begin{proposition}[9.9.4]
\label{IV.9.9.4-fr}
Soit $f:X\to S$ un morphisme localement de présentation finie. L'ensemble des
$x\in X$ tels que $X_{f(x)}$ ait au point $x$ l'une (déterminée) des propriétés
suivantes~:
\begin{enumerate}
\item[(i)] être géométriquement régulier~;
\item[(ii)] posséder la propriété $(R_k)$ géométrique~;
\item[(iii)] être géométriquement normal~;
\item[(iv)] être géométriquement réduit (i.e. séparable)~;
\item[(v)] être géométriquement ponctuellement intègre~;
\end{enumerate}
est une partie localement constructible de $X$.
\end{proposition}

Pour les propriétés (iv) et (v), cela résulte de (9.9.2, (iv) et (v)) appliqué à
$\mathcal F=\mathcal O_X$. Pour les autres propriétés, compte tenu de (6.7.8),
on se ramène, comme au début de (9.9.2) au cas où $S$ est noethérien et intègre
et de point générique $\eta=f(x)$. En outre, en vertu du théorème de platitude
générique (6.9.1), on peut, en remplaçant $S$ par un voisinage ouvert de
$\eta$, supposer que $f$ est un morphisme \emph{plat}. Dire que $X_{f(y)}$ est
géométriquement régulier au point $y$ signifie alors que $f$ est
\emph{régulier} au point $y$ (6.8.1), et on sait que l'ensemble $L$ de ces
points est \emph{ouvert} dans $X$ (6.8.7), ce qui prouve la proposition dans le
cas (i).

Pour prouver le cas (ii), posons $F=X-L$, qui est \emph{fermé} dans $X$. Dire
que $X_s$

\oldpage[IV-3]{94}

vérifie la propriété géométrique $(R_k)$ au point $y\in f^{-1}(s)$ signifie, ou
bien que l'on a $y\in L$, ou bien que les points génériques $z_i$ des
composantes irréductibles de l'ensemble fermé $F_s$ qui contiennent $y$
vérifient la relation $\dim(\mathcal O_{X_s,z_i})\geq k+1$ (compte tenu de
(4.2.7) et (5.2.3))~; autrement dit, les points $y\in F_s$ où $X_s$ vérifie la
propriété $(R_k)$ géométrique sont ceux tels que
$\operatorname{codim}_y(F_s,X_s)\geq k+1$ (5.1.2). La conclusion résulte donc
de (i) et de (9.9.1, (ii)).

Enfin, dans le cas (iii), la conclusion résulte de (ii), de (9.9.2, (viii)), du
fait que la propriété $(S_k)$ est stable par extension du corps de base (6.7.1)
et enfin du critère de Serre (5.8.6).

\begin{corollary}[9.9.5]
\label{IV.9.9.5-fr}
Soit $f:X\to S$ un morphisme de présentation finie, et notons $\mathbf P$
l'une des propriétés (i) à (v) de (9.9.4). Alors l'ensemble des $s\in S$ tels
que la propriété $\mathbf P$ soit vraie en tous les points $x\in X_s$, est
localement constructible dans $S$.
\end{corollary}

La démonstration à partir de (9.9.4) est la même que celle de (9.9.3) à partir
de (9.9.2).

\begin{proposition}[9.9.6]
\label{IV.9.9.6-fr}
Soient $f:X\to S$ un morphisme localement de présentation finie,
$\mathcal L_\bullet$ un complexe formé de $\mathcal O_X$-Modules
quasi-cohérents de présentation finie~; pour tout $s\in S$, soit
$(\mathcal L_\bullet)_s$ le complexe
$\mathcal L_\bullet\otimes_{\mathcal O_S}k(s)$ de
$\mathcal O_{X_s}$-Modules cohérents de type fini. Alors, pour un entier $n$
donné, l'ensemble des $x\in X$ tels que
$(\mathcal H_n((\mathcal L_\bullet)_{f(x)}))_x=0$ est localement constructible
dans $X$.
\end{proposition}

On peut se borner au cas où $\mathcal L_i=0$ sauf pour $i=0$, 1 ou 2, et où
$n=1$. En outre, la question étant locale sur $X$, on peut se borner au cas où
$S=\operatorname{Spec}(A)$ et $X=\operatorname{Spec}(B)$ sont affines, $B$
étant une $A$-algèbre de présentation finie. Il existe alors un sous-anneau
noethérien $A_0$ de $A$, un $A_0$-préschéma de type fini $X_0$ et un complexe
$\mathcal L_\bullet^{(0)}$ de $\mathcal O_{X_0}$-Modules cohérents, nuls sauf en
dimensions 0, 1 et 2 et tels que
$X=X_0\otimes_{A_0}A$ et
$\mathcal L_\bullet=\mathcal L_\bullet^{(0)}\otimes_{A_0}A$. Pour tout
$s\in S$, si $s_0$ est la projection de $s$ dans
$S_0=\operatorname{Spec}(A_0)$, on a
$X_s=(X_0)_{s_0}\otimes_{k(s_0)}k(s)$, et le morphisme projection
$X_s\to(X_0)_{s_0}$ est fidèlement plat~; on en conclut que l'on a
\[
\mathcal H_n((\mathcal L_\bullet)_s)=
\mathcal H_n((\mathcal L_\bullet^{(0)})_{s_0})
\otimes_{k(s_0)}k(s),
\]
et par suite, si $E$ (resp. $E_0$) est l'ensemble des $x\in X$ (resp.
$x_0\in X_0$) tels que
\[
(\mathcal H_n((\mathcal L_\bullet)_{f(x)}))_x=0
\quad\text{(resp.}\quad
(\mathcal H_n((\mathcal L_\bullet^{(0)})_{f_0(x_0)}))_{x_0}=0\text{)},
\]
on a $E=h^{-1}(E_0)$, où $h:X\to X_0$ est la projection canonique. En vertu
de (1.8.2), on peut donc se borner au cas où $A$ est \emph{noethérien}~; il
s'agit de voir ($\mathbf 0_{\mathrm{III}},9.2.3$) que si $x\in X$ est tel que
$(\mathcal H_n((\mathcal L_\bullet)_{f(x)}))_x=0$ (resp. $\neq0$), il existe un
voisinage ouvert $V$ de $x$ dans $\overline{\{x\}}$ tel que, pour tout
$x'\in V$, on ait
$(\mathcal H_n((\mathcal L_\bullet)_{f(x')}))_{x'}=0$ (resp. $\neq0$).
Remplaçant $S$ par le sous-préschéma réduit de $S$ ayant
$\overline{\{f(x)\}}$ pour espace sous-jacent, on peut supposer que $S$ est
intègre et que $f(x)=\eta$ est son point générique. Alors, en restreignant $S$
à un voisinage ouvert de $\eta$, on peut supposer que pour tout $s\in S$, on a
$(\mathcal H_n(\mathcal L_\bullet))_s=
\mathcal H_n((\mathcal L_\bullet)_s)$ (9.4.3), et par suite, si $Z$ est le
support de $\mathcal H_n(\mathcal L_\bullet)$, le support de
$\mathcal H_n((\mathcal L_\bullet)_s)$ est $Z_s=Z\cap X_s$
(\textbf{I}, 9.1.13.1). L'hypothèse est que
$Z_\eta\cap\overline{\{x\}}=\varnothing$ (resp. $\neq\varnothing$). Comme
$Z\cap\overline{\{x\}}$ est fermé dans l'espace noethérien $X$, on conclut de
(9.5.1) qu'il y a un voisinage $U$ de $\eta$ dans $S$ tel que, pour tout
$s\in U$, on ait $Z_s\cap\overline{\{x\}}=\varnothing$ (resp.
$\neq\varnothing$). L'ensemble $V=f^{-1}(U)\cap\overline{\{x\}}$ répond donc à
la question.

\oldpage[IV-3]{95}
\section{Préschémas de Jacobson}
\label{section:IV.10-fr}

On a déjà eu l'occasion d'observer (5.2.5) que même les préschémas
\emph{excellents} (7.8.5) ne se comportent pas toujours comme les
«~variétés~» de la Géométrie algébrique classique, notamment en ce qui
concerne les questions de dimension~; c'est ainsi que si $X$ est le spectre
d'un anneau de valuation discrète complet, l'ensemble des points fermés de
$X$ (réduit à un seul élément) n'est pas partout dense dans $X$, et que son
complémentaire (réduit à un seul élément) est un ouvert partout dense dans
$X$ mais dont la dimension est nulle, donc $<\dim(X)$. On examine dans ce
paragraphe des conditions générales moyennant lesquelles de tels phénomènes
ne se présentent pas~; il en résulte un comportement plus satisfaisant à
certains égards pour les relations entre dimension, codimension, profondeur
et coprofondeur dans ces préschémas (10.6 et 10.8). En outre, dans les
préschémas considérés, le fait que l'ensemble des points fermés soit partout
dense (et même «~très dense~» (10.1.3)) permet de ne considérer que ces points
dans beaucoup de démonstrations~; on rejoint ainsi le point de vue classique
des «~variétés algébriques~» qui, de notre point de vue, sont les ensembles
de points fermés des \emph{préschémas algébriques sur un corps}, et on fait le
lien entre le langage des schémas et celui des «~variétés~» ou «~espaces
algébriques~» de Serre (10.9 et 10.10).

\subsection{Parties très denses d'un espace topologique}
\label{subsection:IV.10.1-fr}

\begin{env}[10.1.1]
\label{IV.10.1.1-fr}
On dit qu'une partie d'un espace topologique $X$ est
\emph{quasi-constructible} si elle est réunion finie de parties localement
fermées de $X$. On dit qu'une partie $T$ de $X$ est \emph{localement
quasi-constructible} si pour tout $x\in X$, il existe un voisinage ouvert $V$
de $x$ tel que $T\cap V$ soit quasi-constructible dans $V$. Il est clair que
toute partie quasi-constructible de $X$ est localement quasi-constructible~:
la réciproque est vraie si $X$ est quasi-compact. Le raisonnement de
($\mathbf 0_{\mathrm{III}}$, 9.1.3), où l'on supprime le mot
«~rétrocompact~», montre que l'ensemble des parties quasi-constructibles
(resp. localement quasi-constructibles) de $X$, que nous désignerons par
$\mathfrak{Qc}(X)$ (resp. $\mathfrak{Lqc}(X)$) est stable par réunion finie,
intersection finie et passage aux complémentaires. Si $f:X\to Y$ est une
application continue, il résulte aussitôt de ces définitions que, pour toute
partie quasi-constructible (resp. localement quasi-constructible) $Z$ de $Y$,
$f^{-1}(Z)$ est quasi-constructible (resp. localement quasi-constructible)
dans $X$.

Les parties constructibles (resp. localement constructibles) de $X$ sont
évidemment quasi-constructibles (resp. localement quasi-constructibles)~; la
réciproque est vraie lorsque $X$ est noethérien (resp. localement
noethérien).

Dans ce qui suit, nous désignerons respectivement par $\mathfrak O(X)$,
$\mathfrak F(X)$, $\mathfrak{Lf}(X)$, $\mathfrak C(X)$, $\mathfrak{Lc}(X)$
l'ensemble des parties de $X$ qui sont respectivement ouvertes, fermées,
localement fermées, constructibles, localement constructibles.
\end{env}

\oldpage[IV-3]{96}
\begin{proposition}[10.1.2]
\label{IV.10.1.2-fr}
Soient $X$ un espace topologique, $X_0$ une partie de $X$. Les conditions
suivantes sont équivalentes~:
\begin{enumerate}[label=\emph{\alph*)}]
  \item Pour toute partie localement fermée $Z\neq\varnothing$ de $X$, on a
  $Z\cap X_0\neq\varnothing$.
  \item[\emph{a$'$)}] Pour toute partie fermée $Z$ de $X$, on a
  $Z=\overline{Z\cap X_0}$.
  \item Pour toute partie $Z\neq\varnothing$ de $X$ localement
  quasi-constructible, on a $Z\cap X_0\neq\varnothing$.
  \item[\emph{b$'$)}] Pour toute partie localement quasi-constructible $Z$ de
  $X$, on a $Z\subset\overline{Z\cap X_0}$ (autrement dit $Z\cap X_0$ est
  \emph{dense} dans $Z$).
  \item L'application $U\mapsto U\cap X_0$ de $\mathfrak O(X)$ dans
  $\mathfrak O(X_0)$ est injective (donc bijective).
  \item[\emph{c$'$)}] L'application $F\mapsto F\cap X_0$ de $\mathfrak F(X)$
  dans $\mathfrak F(X_0)$ est injective (donc bijective).
  \item[\emph{c$''$)}] L'application $Z\mapsto Z\cap X_0$ de
  $\mathfrak{Qc}(X)$ dans $\mathfrak{Qc}(X_0)$ est injective (donc bijective).
  \item[\emph{c$'''$)}] L'application $Z\mapsto Z\cap X_0$ de
  $\mathfrak{Lqc}(X)$ dans $\mathfrak{Lqc}(X_0)$ est injective.
\end{enumerate}
En outre, lorsque ces conditions sont vérifiées, l'application
$Z\mapsto Z\cap X_0$ de $\mathfrak{Lqc}(X)$ dans $\mathfrak{Lqc}(X_0)$ est
bijective.
\end{proposition}

Notons que les assertions de surjectivité dans \emph{c)} et \emph{c$'$)} sont
triviales~: elles entraînent que toute partie localement fermée dans $X_0$ est
trace sur $X_0$ d'une partie localement fermée dans $X$, donc l'application
$\mathfrak{Qc}(X)\to\mathfrak{Qc}(X_0)$ définie dans \emph{c$''$)} est aussi
surjective.

Nous prouverons les implications
\[
c''')\Rightarrow c'')\Rightarrow c')\Rightarrow c)\Rightarrow b)
\Rightarrow b')\Rightarrow a')\Rightarrow a)\Rightarrow c''').
\]
Les trois premières sont triviales. Pour voir que \emph{c)} implique
\emph{b)}, notons d'abord que \emph{c)} entraîne que $X_0$ est \emph{dense}
dans $X$. Remplaçant $X$ et $X_0$ par un ouvert convenable $U$ de $X$ et par
$U\cap X_0$ respectivement, on peut donc se borner au cas où $Z$ est
localement fermée dans $X$, donc $Z=V\cap\complement W$, où $V$ et $W$ sont
ouverts dans $X$~; l'hypothèse $Z\neq\varnothing$ signifie que
$V\not\subset W$, ou encore $V\cup W\neq W$. En vertu de \emph{c)}, on a donc
$(V\cup W)\cap X_0\neq W\cap X_0$, donc $V\cap X_0\not\subset W$, et par
suite $(V\cap\complement W)\cap X_0\neq\varnothing$.

Pour voir que \emph{b)} entraîne \emph{b$'$)}, il suffit d'appliquer
\emph{b)} à $Z\cap U$, où $U$ est un voisinage ouvert arbitraire d'un point de
$Z$. Comme $Z\cap\overline{X_0}\subset\overline{Z\cap X_0}$, il est trivial
que \emph{b$'$)} implique \emph{a$'$)}. Pour montrer que \emph{a$'$)} entraîne
\emph{a)}, remarquons que si $Z$ est localement fermée dans $X$, on peut
écrire $Z=F-F'$ où $F$ et $F'$ sont fermés dans $X$ et $F'\subset F$~; d'où
$Z\cap X_0=(F\cap X_0)-(F'\cap X_0)$. Si l'on avait
$Z\cap X_0=\varnothing$, on en déduirait $F\cap X_0=F'\cap X_0$, donc
$F=F'$ en vertu de \emph{a$'$)}, c'est-à-dire $Z=\varnothing$.

Pour voir que \emph{a)} entraîne \emph{c$'''$)}, il suffit de montrer que si
$Z\neq\varnothing$ est localement quasi-constructible, on a
$Z\cap X_0\neq\varnothing$~: en effet, la relation
$Z'\cap X_0=Z''\cap X_0$ est équivalente à
$((Z'\cup Z'')-(Z'\cap Z''))\cap X_0=\varnothing$. Autrement dit, il suffit de
prouver que \emph{a)} entraîne \emph{b)}~; en outre, remplaçant $X$ et $X_0$
par un ouvert $U$ de $X$ et par $U\cap X_0$ respectivement, on est bien
ramené au cas où $Z$ est localement fermée dans $X$, d'où la conclusion.

Reste à montrer que l'application
$\mathfrak{Lqc}(X)\to\mathfrak{Lqc}(X_0)$ est surjective. Soit donc $Z_0$ une
partie localement quasi-constructible de $X_0$~: il y a un recouvrement
$(V_\alpha)$ de $X_0$ par des ouverts de $X_0$ tel que $Z_0\cap V_\alpha$ soit
quasi-constructible dans $V_\alpha$ (et
\oldpage[IV-3]{97}
donc aussi dans $X_0$). En vertu de \emph{c)}, il y a pour tout $\alpha$ un
seul ouvert $U_\alpha$ de $X$ tel que $X_0\cap U_\alpha=V_\alpha$, et en
vertu de \emph{c$''$)} un seul ensemble quasi-constructible $Z_\alpha$ dans
$X$ tel que $Z_0\cap V_\alpha=X_0\cap Z_\alpha$. Si $\alpha$ et $\beta$ sont
deux indices quelconques, on a
\[
Z_\beta\cap U_\alpha\cap X_0
=U_\alpha\cap Z_0\cap V_\beta
=V_\alpha\cap Z_0\cap V_\beta
=Z_\alpha\cap X_0\cap V_\beta
\subset Z_\alpha\cap X_0~;
\]
comme $Z_\beta\cap U_\alpha$ et $Z_\alpha$ sont quasi-constructibles dans
$X$, il résulte de \emph{c$''$)} que
$Z_\beta\cap U_\alpha\subset Z_\alpha$. Si l'on pose
$Z=\bigcup_\alpha Z_\alpha$, on a donc $Z\cap U_\alpha=Z_\alpha$ pour tout
$\alpha$~; d'ailleurs, comme les $V_\alpha$ recouvrent $X_0$, il résulte de
\emph{c)} que les $U_\alpha$ recouvrent $X$, et l'on voit donc que $Z$ est
localement quasi-constructible dans $X$ et $Z_0=Z\cap X_0$.

\begin{definition}[10.1.3]
\label{IV.10.1.3-fr}
Lorsqu'une partie $X_0$ d'un espace topologique vérifie les conditions
équivalentes de (10.1.2), on dit que $X_0$ est \emph{très dense} dans $X$.
\end{definition}

On a déjà vu au cours de la démonstration de (10.1.2) que $X_0$ est alors
\emph{dense} dans $X$.

\begin{corollary}[10.1.4]
\label{IV.10.1.4-fr}
Si $X_0$ est très dense dans $X$ et $U$ une partie ouverte de $X$,
$U\cap X_0$ est très dense dans $U$. Inversement, si $(U_\alpha)$ est un
recouvrement ouvert de $X$ tel que $U_\alpha\cap X_0$ soit très dense dans
$U_\alpha$ pour tout $\alpha$, alors $X_0$ est très dense dans $X$.
\end{corollary}

Comme toute partie localement fermée dans $U$ est localement fermée dans $X$,
la première assertion résulte du critère \emph{a)} de (10.1.2)~; il en est de
même de la seconde, car si $Z\neq\varnothing$ est localement fermée dans $X$,
$Z\cap U_\alpha$ est localement fermée dans $U_\alpha$ pour tout $\alpha$,
et $Z\cap U_\alpha\neq\varnothing$ pour un $\alpha$ au moins.

\subsection{Quasi-homéomorphismes}
\label{subsection:IV.10.2-fr}

\begin{proposition}[10.2.1]
\label{IV.10.2.1-fr}
Soient $X_0$, $X$ deux espaces topologiques, $f:X_0\to X$ une application
continue. Les conditions suivantes sont équivalentes~:
\begin{enumerate}[label=\emph{\alph*)}]
  \item L'application $U\mapsto f^{-1}(U)$ de $\mathfrak O(X)$ dans
  $\mathfrak O(X_0)$ est bijective.
  \item[\emph{a$'$)}] L'application $F\mapsto f^{-1}(F)$ de $\mathfrak F(X)$
  dans $\mathfrak F(X_0)$ est bijective.
  \item La topologie de $X_0$ est image réciproque par $f$ de celle de $X$, et
  $f(X_0)$ est très dense (10.1.3) dans $X$.
  \item Le foncteur $\mathcal F\mapsto f^*(\mathcal F)$ de la catégorie des
  faisceaux d'ensembles (resp. des faisceaux de groupes abéliens) sur $X$ dans
  la catégorie des faisceaux d'ensembles (resp. des faisceaux de groupes
  abéliens) sur $X_0$ est une équivalence de catégories.
\end{enumerate}
\end{proposition}

Il est clair que \emph{a)} et \emph{a$'$)} sont équivalentes, et \emph{a)}
implique que la topologie de $X_0$ est image réciproque par $f$ de celle de
$X$. D'autre part, si $f(X_0)$ n'est pas très dense dans $X$, il y a deux
ouverts distincts $U_1$, $U_2$ de $X$ tels que
$U_1\cap f(X_0)=U_2\cap f(X_0)$, et par suite
$f^{-1}(U_1)=f^{-1}(U_2)$, ce qui achève de montrer que \emph{a)} entraîne
\emph{b)}. Inversement, la condition \emph{b)} implique que les applications
$U\mapsto U\cap f(X_0)$ de $\mathfrak O(X)$ dans $\mathfrak O(f(X_0))$ et
$V\mapsto f^{-1}(V)$ de $\mathfrak O(f(X_0))$ dans $\mathfrak O(X_0)$ sont
bijectives, donc il en est de même de leur composée $U\mapsto f^{-1}(U)$.

Pour voir que \emph{a)} entraîne \emph{c)}, il suffit d'appliquer la définition
de $f^*(\mathcal F)$ ($\mathbf 0_{\mathrm I}$, 3.7.1) et les axiomes des
faisceaux~: \emph{a)} entraîne que pour tout ouvert $U$ de $X$, l'application
canonique
$\Gamma(U,\mathcal F)\to\Gamma(f^{-1}(U),f^*(\mathcal F))$ est une bijection
fonctorielle en $\mathcal F$, d'où \emph{c)}. Il reste à montrer que
\emph{c)} entraîne \emph{b)}.

\oldpage[IV-3]{98}
Supposons d'abord que $f(X_0)$ ne soit pas très dense dans $X$~; il existe
alors deux parties fermées distinctes $Y\supset Y'$ de $X$ telles que
$Y\cap f(X_0)=Y'\cap f(X_0)$. Soit $\mathcal F$ (resp. $\mathcal F'$) le
faisceau de groupes abéliens sur $X$ image directe par l'injection canonique
$Y\to X$ (resp. $Y'\to X$) du faisceau simple associé au préfaisceau constant
$\mathbf Z$ sur $Y$ (resp. $Y'$)~; la définition de $f^*$
($\mathbf 0_{\mathrm I}$, 3.7.1) montre que $f^*(\mathcal F)$ est isomorphe à
$f^*(\mathcal F')$ mais $\mathcal F$ n'est pas isomorphe à $\mathcal F'$,
donc la condition \emph{c)} n'est pas vérifiée. (On notera que le foncteur
$f^*$ n'est même pas alors \emph{fidèle}, car si $u$ et $v$ sont
l'automorphisme identique et l'endomorphisme nul de $\mathcal F/\mathcal F'$,
$f^*(u)$ et $f^*(v)$ sont égaux.)

Montrons maintenant que \emph{c)} entraîne que la topologie de $X_0$ est
l'image réciproque de celle de $X$ par $f$. Notons que si la condition
\emph{c)} est vérifiée pour la catégorie des faisceaux d'ensembles, elle l'est
aussi pour la catégorie des faisceaux de groupes abéliens, car il résulte
aussitôt des définitions ($\mathbf 0_{\mathrm{III}}$, 8.2.5) que cette dernière
n'est autre que la catégorie des \emph{objets en groupes abéliens} dans la
catégorie des faisceaux d'ensembles. Il suffira donc de prouver notre assertion
en supposant \emph{c)} vérifiée pour les catégories de faisceaux de groupes
abéliens sur $X$ et $X_0$. Or, si $\mathbf Z_X$ désigne le faisceau simple sur
$X$ associé au préfaisceau constant égal à $\mathbf Z$, on a un isomorphisme
canonique
\[
\Gamma(X,\mathcal F)\xrightarrow{\sim}\operatorname{Hom}(\mathbf Z_X,\mathcal F)
\]
fonctoriel en $\mathcal F$ (même raisonnement que dans
($\mathbf 0_{\mathrm I}$, 5.1.1))~; comme il est clair par définition
($\mathbf 0_{\mathrm I}$, 3.7.1) que $f^*(\mathbf Z_X)=\mathbf Z_{X_0}$,
l'hypothèse que $f^*$ est une équivalence entraîne que l'homomorphisme
canonique $\Gamma(X,\mathcal F)\to\Gamma(X_0,f^*(\mathcal F))$ est
\emph{bijectif} pour tout faisceau de groupes abéliens $\mathcal F$ sur $X$.
Or, soit $Y_0$ une partie fermée de $X_0$~; soit $\mathcal F_0$ le faisceau
sur $X$ image directe par l'injection canonique $Y_0\to X_0$ du faisceau
simple associé au préfaisceau constant $\mathbf Z$ sur $Y_0$. Comme $f^*$ est
une équivalence, $\mathcal F_0$ est isomorphe à un faisceau de la forme
$f^*(\mathcal F)$, où $\mathcal F$ est un faisceau de groupes abéliens sur
$X$. Considérons alors la section $s_0$ de $\mathcal F_0$ au-dessus de $X_0$
telle que $(s_0)_{x_0}=1_{x_0}$ si $x_0\in Y_0$,
$(s_0)_{x_0}=0_{x_0}$ si $x_0\notin Y_0$. Les remarques précédentes montrent
qu'il existe une section $s$ de $\mathcal F$ au-dessus de $X$ telle que
$(s_0)_{x_0}=s_{f(x_0)}$ pour tout $x_0\in X_0$ (les fibres
$(\mathcal F_0)_{x_0}$ et $\mathcal F_{f(x_0)}$ étant canoniquement
identifiées ($\mathbf 0_{\mathrm I}$, 3.7.1))~; cela entraîne que $Y_0$ est
l'image réciproque par $f$ de l'ensemble $Y$ des $x\in X$ tels que
$s_x\neq0_x$, et $y$ est fermé dans $X$. C.Q.F.D.

\begin{definition}[10.2.2]
\label{IV.10.2.2-fr}
Lorsqu'une application continue $f:X_0\to X$ vérifie les conditions
équivalentes de (10.2.1), on dit que $f$ est un \emph{quasi-homéomorphisme} de
$X_0$ dans $X$.
\end{definition}

En vertu de (10.2.1, \emph{b)}), dire qu'une partie $X_0$ d'un espace
topologique $X$ est \emph{très dense} dans $X$ signifie donc que l'injection
canonique $X_0\to X$ est un \emph{quasi-homéomorphisme}.

\begin{corollary}[10.2.3]
\label{IV.10.2.3-fr}
Le composé de deux quasi-homéomorphismes est un quasi-homéomorphisme.
\end{corollary}

Cela résulte aussitôt de (10.2.1, \emph{a)}).

\begin{corollary}[10.2.4]
\label{IV.10.2.4-fr}
Soient $f:X\to Y$ un quasi-homéomorphisme, $Y'$ une partie localement
quasi-constructible de $Y$, $X'=f^{-1}(Y')$~; alors la restriction
$f'=f|X':X'\to Y'$ est un quasi-homéomorphisme.
\end{corollary}

Il est clair en vertu de (10.2.1, \emph{b)}) que la topologie induite sur $X'$
par celle de $X$ est l'image réciproque par $f'$ de la topologie induite sur
$Y'$ par celle de $Y$. D'autre part,
\oldpage[IV-3]{99}
soit $Z\neq\varnothing$ une partie fermée \emph{dans} $Y'$, et $z\in Z$~; il
y a un voisinage ouvert $U$ de $z$ dans $Y$ tel que $U\cap Y'$ soit réunion
d'un nombre fini de parties $Y'_j$ fermées dans $U$~; si $j$ est un indice tel
que $z\in Y'_j$, $Z\cap Y'_j$ est donc fermé dans $U$. Puisque $f(X)\cap U$
est très dense dans $U$ (10.1.4), $Z\cap Y'_j\cap f(X)$ n'est pas vide
(10.1.2, \emph{a)}), ni \emph{a fortiori} $Z\cap f(X)$~; mais comme
$Z\subset Y'$, on a $Z\cap f(X)=Z\cap f'(X')$~; le critère
(10.1.2, \emph{a)}) montre donc que $f'(X')$ est très dense dans $Y'$, et
l'on conclut à l'aide de (10.2.1, \emph{b)}).

\begin{corollary}[10.2.5]
\label{IV.10.2.5-fr}
Soient $f:X\to Y$ une application continue, $(V_\alpha)$ un recouvrement
ouvert de $Y$. Si, pour tout $\alpha$, la restriction
$f^{-1}(V_\alpha)\to V_\alpha$ de $f$ est un quasi-homéomorphisme, alors $f$
est un quasi-homéomorphisme.
\end{corollary}

Cela résulte aussitôt du critère (10.2.1, \emph{b)}) et de (10.1.4).

\begin{corollary}[10.2.6]
\label{IV.10.2.6-fr}
Soient $f:X\to Y$ un quasi-homéomorphisme, $Y'$ une partie localement
quasi-constructible de $Y$, $X'=f^{-1}(Y')$. Pour que $Y'$ soit quasi-compact
(resp. noethérien, resp. rétrocompact dans $Y$), il faut et il suffit que $X'$
soit quasi-compact (resp. noethérien, resp. rétrocompact dans $X$).
\end{corollary}

Démontrons d'abord les deux premières assertions~; en vertu de (10.2.4), on
peut se borner au cas où $Y'=Y$. Dire que $X$ est quasi-compact (resp.
noethérien) signifie que pour toute famille filtrante croissante $(U_\alpha)$
dans $\mathfrak O(X)$ ayant $X$ pour plus grand élément (resp. pour toute
famille filtrante croissante $(U_\alpha)$ dans $\mathfrak O(X)$), il existe
$\gamma$ tel que $U_\alpha=U_\gamma$ pour $\alpha\geq\gamma$. Comme
$U\mapsto f^{-1}(U)$ est une bijection de $\mathfrak O(Y)$ sur
$\mathfrak O(X)$, notre assertion résulte aussitôt de la remarque précédente.

Les ouverts quasi-compacts de $X$ sont donc les ensembles de la forme
$f^{-1}(U)$ où $U$ est ouvert quasi-compact dans $Y$, en vertu de
(10.2.1, \emph{a)}) et de ce qui précède. Pour que $X'$ soit rétrocompact dans
$X$, il faut et il suffit alors que pour tout ouvert quasi-compact $U$ dans
$Y$, $f^{-1}(U)\cap X'=f^{-1}(U\cap Y')$ soit quasi-compact
($\mathbf 0_{\mathrm{III}}$, 9.1.1)~; la première partie de la démonstration
montre que cela équivaut à dire que $U\cap Y'$ est quasi-compact pour tout
ouvert quasi-compact $U$, c'est-à-dire que $Y'$ est rétrocompact dans $Y$.

\begin{proposition}[10.2.7]
\label{IV.10.2.7-fr}
Soit $f:X\to Y$ un quasi-homéomorphisme. Alors l'application
$Z\mapsto f^{-1}(Z)$ de $\mathfrak P(Y)$ dans $\mathfrak P(X)$ définit par
restriction les bijections suivantes (cf. (10.1.1) pour les notations)~:
\[
\begin{gathered}
\mathfrak O(Y)\xrightarrow{\sim}\mathfrak O(X)\\
\mathfrak F(Y)\xrightarrow{\sim}\mathfrak F(X)\\
\mathfrak{Lf}(Y)\xrightarrow{\sim}\mathfrak{Lf}(X)\\
\mathfrak{Qc}(Y)\xrightarrow{\sim}\mathfrak{Qc}(X)\\
\mathfrak{Lqc}(Y)\xrightarrow{\sim}\mathfrak{Lqc}(X)\\
\mathfrak C(Y)\xrightarrow{\sim}\mathfrak C(X)\\
\mathfrak{Lc}(Y)\xrightarrow{\sim}\mathfrak{Lc}(X).
\end{gathered}
\]
\end{proposition}

\oldpage[IV-3]{100}
Pour les deux premières, ce n'est autre que la définition (10.2.2)~; comme la
topologie de $X$ est l'image réciproque par $f$ de celle de $f(X)$ les cinq
dernières applications, où l'on remplace $Y$ par $f(X)$, sont bijectives. On
peut donc (par (10.2.1, \emph{b)})) se borner au cas où $X$ est un sous-espace
très dense de $Y$, et le fait que les applications
$\mathfrak{Lf}(Y)\to\mathfrak{Lf}(X)$,
$\mathfrak{Qc}(Y)\to\mathfrak{Qc}(X)$,
$\mathfrak{Lqc}(Y)\to\mathfrak{Lqc}(X)$ sont bijectives a déjà été prouvé
(10.1.2). Toute partie localement constructible étant localement
quasi-constructible, les applications $\mathfrak C(Y)\to\mathfrak C(X)$ et
$\mathfrak{Lc}(Y)\to\mathfrak{Lc}(X)$ sont injectives~; en outre, pour tout
ouvert $U\subset Y$, la restriction $f^{-1}(U)\to U$ de $f$ est un
quasi-homéomorphisme (10.2.4), donc si l'on montre que
$\mathfrak C(Y)\to\mathfrak C(X)$ est surjective, il en sera de même de
$\mathfrak{Lc}(Y)\to\mathfrak{Lc}(X)$. Mais en vertu de (10.2.6), toute partie
ouverte rétrocompacte $Z$ dans $X$ est de la forme $f^{-1}(T)$, où $T$ est
ouvert rétrocompact dans $Y$~; cela prouve évidemment la surjectivité de
$\mathfrak C(Y)\to\mathfrak C(X)$.

\begin{remarks}[10.2.8]
\label{IV.10.2.8-fr}
\upshape (i) On a vu dans la démonstration de (10.2.1) que si $f$ est un
quasi-homéomorphisme, l'application canonique
\begin{equation}
\label{IV.10.2.8.1-fr}
\tag{10.2.8.1}
\Gamma(Y,\mathcal F)\to\Gamma(X,f^*(\mathcal F))
\end{equation}
est un isomorphisme de groupes abéliens fonctoriel en $\mathcal F$ dans la
catégorie des faisceaux de groupes abéliens sur $Y$. Comme $f^*$ est exact
dans cette catégorie ($\mathbf 0_{\mathrm I}$, 3.7.2), cela implique que les
homomorphismes canoniques (T, 3.2.2)
\[
\mathrm H^i(Y,\mathcal F)\xrightarrow{\sim}
\mathrm H^i(X,f^*(\mathcal F))
\]
sont \emph{bijectifs} pour tout $i$.

(ii) Si $f:X\to Y$ est une application continue et $\mathcal F$,
$\mathcal G$ deux faisceaux de groupes abéliens sur $Y$, on a
\[
f^*(\mathcal F\otimes_{\mathbf Z_Y}\mathcal G)
=f^*(\mathcal F)\otimes_{\mathbf Z_X}f^*(\mathcal G)
\]
à un isomorphisme canonique près ($\mathbf 0_{\mathrm I}$, 4.3.3). On en
conclut que si $f$ est un quasi-homéomorphisme, $f^*$ est encore une
\emph{équivalence} de la catégorie des faisceaux d'anneaux sur $Y$, et de la
catégorie des faisceaux d'anneaux sur $X$~; la donnée d'une structure d'espace
annelé sur $Y$ est donc équivalente à celle d'une structure d'espace annelé
sur $X$.

Étant donnés deux espaces annelés $(X,\mathcal O_X)$, $(Y,\mathcal O_Y)$, nous
dirons qu'un morphisme d'espaces annelés
$f=(\psi,\theta):(X,\mathcal O_X)\to(Y,\mathcal O_Y)$ est un
\emph{quasi-isomorphisme} si $\psi$ est un quasi-homéomorphisme de $X$ dans
$Y$ et $\theta^\sharp:\psi^*(\mathcal O_Y)\to\mathcal O_X$ un isomorphisme de
faisceaux d'anneaux~; lorsqu'il en est ainsi, l'espace annelé
$(X,\mathcal O_X)$ est entièrement déterminé, à isomorphisme près, par
$(Y,\mathcal O_Y)$, l'espace $X$ et le quasi-homéomorphisme $\psi$ (que l'on
peut prendre arbitraire). Lorsque $f$ est un quasi-isomorphisme d'espaces
annelés, le foncteur
\[
\mathcal F\mapsto f^*(\mathcal F)
\]
est une \emph{équivalence} de la catégorie des $\mathcal O_Y$-Modules et de
celle des $\mathcal O_X$-Modules, puisque $f^*(\mathcal F)$ s'identifie
canoniquement ici à $\psi^*(\mathcal F)$. On en conclut par exemple des
isomorphismes de bi-$\partial$-foncteurs
\[
\operatorname{Ext}^i_{\mathcal O_Y}(\mathcal F,\mathcal G)
\xrightarrow{\sim}
\operatorname{Ext}^i_{\mathcal O_X}(f^*(\mathcal F),f^*(\mathcal G)).
\]
De façon générale, on peut dire que les constructions usuelles de la théorie
des faisceaux et de l'algèbre homologique, faites sur l'espace annelé $Y$ ou
l'espace annelé $X$, sont équivalentes.
\end{remarks}

\oldpage[IV-3]{101}
\subsection{Espaces de Jacobson}
\label{subsection:IV.10.3-fr}

\begin{definition}[10.3.1]
\label{IV.10.3.1-fr}
On dit qu'un espace topologique $X$ est un \emph{espace de Jacobson} si
l'ensemble $X_0$ des points fermés de $X$ est très dense dans $X$ (autrement
dit, si l'injection canonique $X_0\to X$ est un
quasi-homéomorphisme).
\end{definition}

Cela signifie donc (10.1.2) que toute partie localement fermée (ou seulement
localement quasi-constructible) $Z\neq\varnothing$ de $X$ contient un point
fermé de $X$, ou encore que \emph{toute partie fermée de $X$ est l'adhérence
de l'ensemble de ses points fermés}.

\begin{proposition}[10.3.2]
\label{IV.10.3.2-fr}
Soient $X$ un espace de Jacobson, $Z$ une partie localement
quasi-constructible de $X$~; alors le sous-espace $Z$ de $X$ est un espace de
Jacobson, et pour qu'un point de $Z$ soit fermé dans $Z$, il faut et il suffit
qu'il soit fermé dans $X$.
\end{proposition}

Si $X_0$ est l'ensemble des points fermés de $X$, $Z\cap X_0$ est très dense
dans $Z$ en vertu de (10.2.4) appliqué à l'injection $i:X_0\to X$~; comme
l'ensemble $Z_0$ des points fermés de $Z$ contient évidemment $Z\cap X_0$,
il est \emph{a fortiori} très dense dans $Z$, donc $Z$ est un espace de
Jacobson. Montrons maintenant qu'en fait on a $Z_0=Z\cap X_0$~; soit $x\in Z$
un point fermé dans $Z$~; soit $\overline{\{x\}}$ son adhérence dans $X$~;
$\overline{\{x\}}\cap Z=\{x\}$ est donc une partie localement
quasi-constructible de $X$, et comme son intersection avec $X_0$ est non vide
(10.1.2), on a $x\in X_0$.

\begin{proposition}[10.3.3]
\label{IV.10.3.3-fr}
Soient $X$ un espace topologique, $(U_\alpha)$ un recouvrement ouvert de $X$.
Pour que $X$ soit un espace de Jacobson, il faut et il suffit que chacun des
sous-espaces $U_\alpha$ le soit.
\end{proposition}

La condition est nécessaire en vertu de (10.3.2). Inversement, montrons
d'abord que l'hypothèse que les $U_\alpha$ sont des espaces de Jacobson
entraîne que pour qu'un point $x\in U_\alpha$ soit fermé dans $X$, il suffit
qu'il soit fermé dans $U_\alpha$. Il suffit en effet de voir que cette
condition entraîne que $x$ est aussi fermé dans chacun des $U_\beta$ qui le
contiennent~; mais $U_\alpha\cap U_\beta$ est ouvert dans $U_\alpha$, donc $x$
est fermé dans $U_\alpha\cap U_\beta$, et en vertu de (10.3.2), $x$ est aussi
fermé dans $U_\beta$, ce qui achève la démonstration.

\subsection{Préschémas de Jacobson et anneaux de Jacobson}
\label{subsection:IV.10.4-fr}

\begin{definition}[10.4.1]
\label{IV.10.4.1-fr}
On dit qu'un préschéma $X$ est un \emph{préschéma de Jacobson} si l'espace
topologique $X$ sous-jacent est un espace de Jacobson. On dit qu'un anneau
$A$ est un \emph{anneau de Jacobson} si $\operatorname{Spec}(A)$ est un
préschéma de Jacobson.
\end{definition}

Toute partie fermée de $X=\operatorname{Spec}(A)$ est de la forme
$Z=V(\mathfrak a)$, où $\mathfrak a$ est un idéal égal à sa racine, et
l'ensemble $X_0$ des points fermés de $X$ est l'ensemble des idéaux maximaux
de $A$~; dire que $Z\cap X_0$ est dense dans $Z$ signifie donc que
$\mathfrak a$ est \emph{intersection d'idéaux maximaux} (\textbf{I}, 1.1.4)~;
comme $\mathfrak a$ est intersection d'idéaux premiers, il revient au même de
dire que tout idéal premier de $A$ est intersection d'idéaux maximaux~; en
vertu de (10.3.1) et (10.1.2), la définition usuelle des anneaux de Jacobson
(Bourbaki, \emph{Alg. comm.}, chap.~V, \S~3, n\textsuperscript{o}~4, déf.~1)
coïncide donc avec la définition (10.4.1).

\oldpage[IV-3]{102}
\begin{proposition}[10.4.2]
\label{IV.10.4.2-fr}
Soient $X$ un préschéma, $(U_\alpha)$ un recouvrement de $X$ formé d'ouverts
affines. Pour que $X$ soit un préschéma de Jacobson, il faut et il suffit que
les anneaux $\Gamma(U_\alpha,\mathcal O_X)$ des $U_\alpha$ soient des anneaux
de Jacobson.
\end{proposition}

Cela résulte de (10.3.3) et de la définition (10.4.1).

\begin{env}[10.4.3]
\label{IV.10.4.3-fr}
Un préschéma discret est un préschéma de Jacobson~; un anneau artinien est donc
un anneau de Jacobson. Tout anneau principal ayant une infinité d'idéaux
maximaux (par exemple $\mathbf Z$) est un anneau de Jacobson~; un anneau local
noethérien $A$ est un anneau de Jacobson si et seulement si son idéal maximal
est son seul idéal premier, c'est-à-dire si $A$ est artinien. Tout
sous-préschéma d'un préschéma de Jacobson est un préschéma de Jacobson en vertu
de (10.3.2).
\end{env}

\begin{proposition}[10.4.4]
\label{IV.10.4.4-fr}
Soit $B$ un anneau intègre. Les conditions suivantes sont équivalentes~:
\begin{enumerate}[label=\emph{\alph*)}]
  \item Il existe $f\neq0$ dans $B$ tel que $B_f$ soit un corps.
  \item Le corps des fractions de $B$ est une $B$-algèbre de type fini.
  \item Il existe un corps $K$ contenant $B$ et qui est une $B$-algèbre de
  type fini.
  \item Le point générique de $\operatorname{Spec}(B)$ est isolé dans
  $\operatorname{Spec}(B)$.
\end{enumerate}
\end{proposition}

Il est clair que \emph{d)} est équivalent à \emph{a)}, puisque \emph{d)}
signifie qu'il existe $f\neq0$ dans $B$ tel que $D(f)$ soit réduit au point
générique de $\operatorname{Spec}(B)$. Il est trivial que \emph{a)} entraîne
\emph{b)} et que \emph{b)} entraîne \emph{c)}. Enfin, \emph{c)} entraîne
\emph{a)}, en vertu de Bourbaki, \emph{Alg. comm.}, chap.~V, \S~3,
n\textsuperscript{o}~1, cor.~2 du th.~1.

\begin{proposition}[10.4.5]
\label{IV.10.4.5-fr}
Soit $A$ un anneau. Les conditions suivantes sont équivalentes~:
\begin{enumerate}[label=\emph{\alph*)}]
  \item $A$ est un anneau de Jacobson.
  \item Pour tout idéal premier non maximal $\mathfrak p$ de $A$ et tout
  $f\neq0$ dans $B=A/\mathfrak p$, $B_f$ n'est pas un corps.
  \item[\emph{b$'$)}] Toute $A$-algèbre de type fini $K$ qui est un corps est
  une $A$-algèbre finie.
\end{enumerate}
\end{proposition}

On sait que \emph{a)} entraîne \emph{b$'$)} (Bourbaki,
\emph{Alg. comm.}, chap.~V, \S~3, n\textsuperscript{o}~4, cor.~3 du th.~3).
En outre, le noyau de l'homomorphisme $A\to K$ est alors un idéal
\emph{maximal} $\mathfrak m$ de $A$, et $K$ est une extension \emph{finie} de
$A/\mathfrak m$ (\emph{loc. cit.}). Il est trivial que \emph{b$'$)} entraîne
\emph{b)}, puisque $A/\mathfrak p=B$ n'est pas un corps, toute $B$-algèbre de
type fini est une $A$-algèbre de type fini et $B_f$ est une $B$-algèbre de
type fini. Reste à voir que \emph{b)} implique \emph{a)}. Nous utiliserons le
lemme suivant~:

\begin{lemma}[10.4.5.1]
\label{IV.10.4.5.1-fr}
Soit $X$ un espace topologique ayant la propriété suivante~: pour toute partie
localement fermée $Z\neq\varnothing$ de $X$, il existe une partie $Z'$ de $Z$,
localement fermée dans $X$ (ou dans $Z$, ce qui revient au même), et un point
$x\in Z'$, fermé dans $Z'$. Alors, pour que $X$ soit un espace de Jacobson, il
faut et il suffit qu'aucun point non fermé $x$ de $X$ ne soit isolé dans
$\overline{\{x\}}$.
\end{lemma}

Si $X$ est un espace de Jacobson, un point non fermé $x$ de $X$ ne peut être
isolé dans $\overline{\{x\}}$, car cela signifierait qu'il existe un ouvert
$U$ de $X$ tel que $U\cap\overline{\{x\}}=\{x\}$~; or
$U\cap\overline{\{x\}}$ est localement fermé dans $X$ et ne contiendrait aucun
point fermé dans $X$, ce qui est contraire à l'hypothèse que $X$ est un espace
de Jacobson. Inversement, supposons vérifiée la condition de l'énoncé~; alors
l'ensemble des points fermés de $X$ est identique à l'ensemble $X_0$ des
$x\in X$ qui sont isolés dans $\overline{\{x\}}$~; mais il résulte de
(5.1.10.1) que

\oldpage[IV-3]{103}
cet ensemble est \emph{très dense} dans $X$, donc $X$ est un espace de
Jacobson par définition.

Rappelons (5.1.10) que l'hypothèse faite sur $X$ dans (10.4.5.1) est toujours
vérifiée lorsque $X$ est l'espace sous-jacent à un \emph{préschéma}.

Revenant alors à la démonstration de (10.4.5), la condition \emph{b)} entraîne
que pour tout point $x$ non fermé de $\operatorname{Spec}(A)$, le point
générique $x$ de
$\operatorname{Spec}(A/\mathfrak j_x)=V(\mathfrak j_x)=\overline{\{x\}}$
n'est pas isolé dans $\overline{\{x\}}$, donc \emph{b)} entraîne \emph{a)} en
vertu du lemme (10.4.5.1) et de (10.4.4).

\begin{corollary}[10.4.6]
\label{IV.10.4.6-fr}
Toute algèbre de type fini $B$ sur un anneau de Jacobson $A$ est un anneau de
Jacobson, et l'image réciproque dans $A$ de tout idéal maximal de $B$ est un
idéal maximal de $A$. En particulier, toute algèbre de type fini sur un corps
ou sur $\mathbf Z$ est un anneau de Jacobson.
\end{corollary}

Une $B$-algèbre de type fini $K$ qui est un corps est aussi une $A$-algèbre de
type fini, donc un $A$-module de type fini et \emph{a fortiori} un $B$-module
de type fini, d'où la première assertion~; la seconde a été démontrée au cours
de la démonstration de (10.4.5), appliqué à $K=B/\mathfrak m$.

\begin{corollary}[10.4.7]
\label{IV.10.4.7-fr}
Si $X$ est un préschéma de Jacobson, $f:Y\to X$ un morphisme localement de type
fini, alors $Y$ est un préschéma de Jacobson et l'image par $f$ de tout point
fermé dans $Y$ est un point fermé dans $X$.
\end{corollary}

La question étant locale sur $X$ et sur $Y$, on est ramené au cas où $X$ et
$Y$ sont affines, et le corollaire résulte alors de (10.4.6).

\begin{corollary}[10.4.8]
\label{IV.10.4.8-fr}
Si $k$ est un corps algébriquement clos, $X$ un $k$-préschéma localement de
type fini, l'ensemble des points de $X$ rationnels sur $k$ est très dense dans
$X$.
\end{corollary}

En effet, $X$ est un préschéma de Jacobson (10.4.7) et les points fermés de
$X$ sont exactement les points rationnels sur $k$ (\textbf{I}, 6.4.2).

\begin{env}[10.4.9]
\label{IV.10.4.9-fr}
Le fait que les préschémas localement de type fini sur un corps ou sur
$\mathbf Z$ sont des préschémas de Jacobson est particulièrement important,
en raison de la possibilité de se ramener à ce cas dans de nombreuses
questions de Géométrie algébrique (8.1.2, \emph{c)}). Nous allons en donner
deux exemples~:
\end{env}

\begin{env}[10.4.10]
\label{IV.10.4.10-fr}
\textbf{Applications (10.4.10)~: I~: Démonstration de (6.15.9).} --- Soient
$k$ un corps séparablement clos, $X$ un $k$-préschéma localement de type fini
sur $k$ et \emph{unibranche}. On sait que la fermeture intégrale d'une
$k$-algèbre intègre de type fini $A$ dans une extension finie de son corps des
fractions est une $A$-algèbre finie (Bourbaki, \emph{Alg. comm.}, chap.~V,
\S~3, n\textsuperscript{o}~2, th.~2), donc toute $k$-algèbre de type fini est
un anneau universellement japonais~; on en conclut que l'ensemble des points
$x\in X$ où $X$ est géométriquement unibranche est localement constructible
(9.7.10). Mais l'hypothèse et le lemme (6.15.8) entraînent que cet ensemble
contient tous les points \emph{fermés} de $X$. La conclusion résulte donc de
(10.4.6), (10.3.1) et de la bijectivité de l'application canonique
$\mathfrak{Lc}(X)\to\mathfrak{Lc}(X_0)$ (où $X_0$ est l'ensemble des points
fermés de $X$) (10.2.7).
\end{env}

\begin{env}[10.4.11]
\label{IV.10.4.11-fr}
\textbf{Applications~: II~: Proposition (10.4.11).} --- Soient $S$ un
préschéma, $X$ un $S$-préschéma de type fini. Tout $S$-endomorphisme de $X$
qui est radiciel est surjectif (donc bijectif).
\end{env}

\oldpage[IV-3]{104}
Soient $f:X\to S$ le morphisme structural, $g:X\to X$ le $S$-endomorphisme
considéré et, pour tout $s\in S$, soit $g_s$ le morphisme déduit de $g$ par le
changement de base $\operatorname{Spec}(k(s))\to S$, qui est un
$k(s)$-endomorphisme de la fibre
$f^{-1}(s)=X\times_S\operatorname{Spec}(k(s))$.

Pour prouver que $g$ est surjectif, il suffit de prouver que $g_s$ est
surjectif pour tout $s\in S$, donc on peut (en vertu de (\textbf{I}, 3.5.7))
se borner au cas où $S=\operatorname{Spec}(k)$ est le spectre d'un corps $k$,
auquel cas $f$ est un morphisme de présentation finie, puisque $S$ est
noethérien. Appliquant (8.9.1) et (8.10.5, (vii)), on est ramené au cas où
$S=\operatorname{Spec}(A)$, où $A$ est une sous-$\mathbf Z$-algèbre de type
fini de $k$. Or $X$ est alors un préschéma de Jacobson (10.4.7), et $g(X)$ est
constructible dans $X$ (\textbf{I}, 1.8.5), donc, pour prouver que $g(X)=X$,
il suffit de montrer que $g(X)$ contient tous les points \emph{fermés} de $X$
(10.3.1).

\begin{lemma}[10.4.11.1]
\label{IV.10.4.11.1-fr}
Soit $Y$ un $\mathbf Z$-préschéma de type fini.
\begin{enumerate}[label=(\roman*)]
  \item Pour qu'un point $y\in Y$ soit fermé dans $Y$, il faut et il suffit
  que $k(y)$ soit un corps fini.
  \item Pour tout nombre premier $p$ et tout entier $d\geq1$, l'ensemble des
  points $y\in Y$ tels que $k(y)$ soit une extension de $\mathbf F_p$ dont le
  degré divise $d$ est fini.
\end{enumerate}
\end{lemma}

L'assertion (i) résulte de ce que l'image d'un point fermé $y\in Y$ dans
$\operatorname{Spec}(\mathbf Z)$ est un point fermé (10.4.7), autrement dit
un nombre premier, et de (\textbf{I}, 6.4.2). D'autre part, comme $Y$ est
réunion finie d'ouverts affines de type fini sur $\mathbf Z$, on peut se
borner pour prouver (ii) au cas où $Y=\operatorname{Spec}(C)$, où $C$ est une
$\mathbf Z$-algèbre de type fini. Or, les points $y\in Y$ tels que le degré de
$k(y)$ sur $\mathbf F_p$ divise $d$ correspondent biunivoquement aux
homomorphismes $C\to\mathbf F_{p^d}$~; mais si $(t_i)$ est un système de
générateurs de la $\mathbf Z$-algèbre $C$, tout homomorphisme de $C$ est
déterminé par ses valeurs aux éléments $t_i$, et par suite il n'y a qu'un
nombre fini d'homomorphismes de $C$ dans un corps fini.

Cela étant, $X$ est un $\mathbf Z$-préschéma de type fini~; soit $T_{p,d}$
l'ensemble des points fermés $z\in X$ tels que $k(z)$ soit une extension de
$\mathbf F_p$ de degré divisant $d$~; il résulte de (10.4.11.1) que l'ensemble
$T_{p,d}$ est fini et que l'ensemble des points fermés de $X$ est réunion des
$T_{p,d}$. En outre, si $z\in T_{p,d}$ et si $h$ est un endomorphisme
quelconque de $X$, $k(h(z))$ est isomorphe à un sous-corps de $k(z)$, donc
$h(z)\in T_{p,d}$, autrement dit $T_{p,d}$ est \emph{stable} par tout
endomorphisme de $X$. Comme $g$ est par hypothèse injectif, sa restriction à
$T_{p,d}$ est \emph{bijective} puisque $T_{p,d}$ est fini, ce qui achève de
démontrer la proposition.

Nous verrons plus loin (17.9.7) que lorsqu'on suppose de plus, d'une part que
$X$ est un $S$-préschéma de présentation finie, d'autre part que $g$ est un
monomorphisme, alors on peut affirmer que $g$ est un \emph{automorphisme} de
$X$.

\subsection{Préschémas de Jacobson noethériens}
\label{subsection:IV.10.5-fr}

\begin{proposition}[10.5.1]
\label{IV.10.5.1-fr}
Soit $B$ un anneau intègre noethérien. Les conditions équivalentes \emph{a)} à
\emph{d)} de (10.4.4) sont alors aussi équivalentes aux suivantes~:
\begin{enumerate}[label=\emph{\alph*)}]
  \setcounter{enumi}{4}
  \item $\operatorname{Spec}(B)$ est fini.
  \item $B$ est un anneau semi-local de dimension $\leq1$.
\end{enumerate}
\end{proposition}

Il résulte en effet du théorème d'Artin--Tate ($\mathbf 0$, 16.3.3) que les
conditions \emph{a)} et \emph{f)} sont équivalentes. La condition \emph{f)}
implique que $\operatorname{Spec}(B)$ est réunion de l'ensemble

\oldpage[IV-3]{105}
fini de ses points fermés et de son point générique, donc \emph{f)} implique
\emph{e)} sans supposer $A$ noethérien. Enfin \emph{e)} implique \emph{d)} sans
supposer $A$ noethérien, car le point générique $x$ de $X$ est le complémentaire
de la réunion des adhérences $\overline{\{y\}}$, où $y$ parcourt l'ensemble des
points $y\neq x$, et comme ces points sont en nombre fini, $\{x\}$ est
complémentaire d'un ensemble fermé dans $X$.

\begin{corollary}[10.5.2]
\label{IV.10.5.2-fr}
Soit $A$ un anneau noethérien. Pour que $A$ soit un anneau de Jacobson, il faut
et il suffit qu'il n'existe aucun idéal premier $\mathfrak p$ de $A$ tel que
$A/\mathfrak p$ soit un anneau semi-local de dimension $1$.
\end{corollary}

Cela résulte aussitôt de (10.5.1) et de la condition \emph{b)} de (10.4.5),
les idéaux premiers $\mathfrak p$ de $A$ tels que $A/\mathfrak p$ soit
semi-local de dimension $0$ étant les idéaux maximaux de $A$.

\begin{corollary}[10.5.3]
\label{IV.10.5.3-fr}
Soit $X$ un préschéma noethérien irréductible. Les conditions suivantes sont
équivalentes~:
\begin{enumerate}[label=\emph{\alph*)}]
  \item Le point générique de $X$ est isolé.
  \item $X$ est fini.
  \item $X$ est de dimension $\leq1$ et l'ensemble de ses points fermés est
  fini.
\end{enumerate}
\end{corollary}

Il existe par hypothèse un nombre fini d'ouverts affines irréductibles $U_i$
($1\leq i\leq n$) recouvrant $X$, et dont chacun contient donc le point
générique de $X$~; il suffit de prouver l'équivalence de \emph{a)}, \emph{b)}
et \emph{c)} pour chacun des $(U_i)_{\mathrm{red}}$ (compte tenu de
($\mathbf 0$, 14.1.7)). Mais cette équivalence résulte alors de (10.5.1).

\begin{remark}[10.5.4]
\label{IV.10.5.4-fr}
Un préschéma noethérien $X$ vérifiant les conditions équivalentes de (10.5.3)
n'est pas nécessairement un schéma affine~; en fait il peut même être
\emph{non séparé}. On en a un exemple en remplaçant, dans l'exemple
(\textbf{I}, 5.5.11) de «~droite affine avec point dédoublé~», $X_1$ et $X_2$
par le spectre de l'anneau de valuation discrète $(K[s])_{(s)}$, $U_{12}$ et
$U_{21}$ par l'ouvert réduit au point générique dans $X_1$ et $X_2$
respectivement~; le préschéma non séparé $X$ que l'on obtient a exactement
$3$ points.
\end{remark}

\begin{proposition}[10.5.5]
\label{IV.10.5.5-fr}
Soit $X$ un préschéma noethérien.
\begin{enumerate}[label=(\roman*)]
  \item L'ensemble $X_0$ des points $x\in X$ tels que $\overline{\{x\}}$ soit
  fini est très dense dans $X$.
  \item Pour que $X$ soit un préschéma de Jacobson, il faut et il suffit qu'il
  n'existe aucun sous-préschéma de $X$ isomorphe au spectre d'un anneau
  semi-local intègre de dimension $1$.
\end{enumerate}
\end{proposition}

La condition que $\overline{\{x\}}$ soit fini équivaut en effet ici (10.5.3)
au fait que $x$ soit isolé dans $\overline{\{x\}}$, $\overline{\{x\}}$ étant
espace sous-jacent d'un sous-préschéma (noethérien) de $X$~; l'assertion (i)
résulte donc de (5.1.10.1). De même, compte tenu de (10.4.5.1), pour démontrer
l'assertion (ii), remarquons que pour qu'un point non fermé $x$ de $X$
appartienne à $X_0$, il faut et il suffit que le sous-préschéma fermé intègre
$Y$ de $X$ ayant $\overline{\{x\}}$ pour espace sous-jacent soit de dimension
$1$ et fini, donc réunion finie de sous-préschémas ouverts (dans $Y$) affines
$U_i$ qui sont des spectres d'anneaux semi-locaux intègres de dimension $1$.
Inversement, s'il y a un sous-préschéma $Z$ de $X$ qui soit spectre d'un anneau
semi-local intègre de dimension $1$, $Z$ n'est pas un préschéma de Jacobson
(10.5.2), donc il en est de même de $X$ (10.3.2).

\oldpage[IV-3]{106}
\begin{remark}[10.5.6]
\label{IV.10.5.6-fr}
L'assertion (ii) de (10.5.5) est encore valable lorsque $X$ est
\emph{localement noethérien}~: en effet, si $(V_\alpha)$ est un recouvrement de
$X$ formé d'ouverts affines (noethériens), tout sous-préschéma d'un $V_\alpha$
est un sous-préschéma de $X$~; inversement, si un sous-préschéma $Z$ de $X$ est
isomorphe au spectre d'un anneau semi-local intègre de dimension $1$, il y a
un $\alpha$ tel que $V_\alpha\cap Z$ contienne un ouvert affine $U$ de $Z$ non
réduit au point générique de $Z$, et qui est donc aussi spectre d'un anneau
semi-local intègre de dimension $1$. On conclut à l'aide de (10.3.3).
\end{remark}

\begin{proposition}[10.5.7]
\label{IV.10.5.7-fr}
Soient $X$ un préschéma localement noethérien, $Y$ une partie fermée de $X$
telle que toute partie fermée non vide de $X$ rencontre $Y$. Alors le
préschéma induit sur l'ouvert $X-Y$ est un préschéma de Jacobson.
\end{proposition}

Appliquons le critère (10.5.5, (ii)), et supposons qu'il y ait un
sous-préschéma $Z$ de $X-Y$ qui soit spectre d'un anneau semi-local intègre de
dimension $1$, le point générique $z$ de $Z$ étant donc isolé dans $Z$ (ou
dans l'adhérence $\overline Z$ de $Z$ dans $X$), et $Z$ étant distinct de
$\{z\}$. Soit $y\neq z$ un point de $Z$~; comme il n'appartient pas à $Y$, il
n'est pas fermé dans $X$, et son adhérence $\overline{\{y\}}$ dans $X$
rencontre $Y$ en un point $x\neq y$ qui est donc spécialisation de $y$.
L'existence de la chaîne
$\{x\}\subset\overline{\{y\}}\subset\overline{\{z\}}$ montre alors que la
dimension de $\overline{\{z\}}$ serait $\geq2$, et il en serait de même de la
dimension de $U\cap\overline{\{z\}}$, où $U$ est un voisinage ouvert affine
(donc noethérien) de $x$ dans $X$. Or, cela contredit le fait que le point
générique $z$ de $U\cap\overline{\{z\}}$ est isolé dans
$U\cap\overline{\{z\}}$ (10.5.3).

\begin{corollary}[10.5.8]
\label{IV.10.5.8-fr}
Soit $A$ un anneau noethérien~; pour tout élément $f$ du radical $\mathfrak R$
de $A$, l'anneau $A_f$ est un anneau de Jacobson, et l'ouvert
$\operatorname{Spec}(A)-V(\mathfrak R)$ est un schéma de Jacobson.
\end{corollary}

Si $X=\operatorname{Spec}(A)$, $Y=V(\mathfrak J)$, où $\mathfrak J$ est un
idéal de $A$, dire que toute partie fermée non vide de $X$ rencontre $Y$
équivaut ($\mathbf 0_{\mathrm I}$, 2.1.3) à dire que $Y$ contient tous les
\emph{points fermés} de $X$, ou encore que $\mathfrak J$ est \emph{contenu
dans le radical} $\mathfrak R$ de $A$. Si $f\in\mathfrak R$, l'ouvert $D(f)$
ne rencontre donc pas $V(\mathfrak R)$, et est un espace de Jacobson en vertu
de (10.5.7).

\begin{corollary}[10.5.9]
\label{IV.10.5.9-fr}
Soient $A$ un anneau local noethérien, $\mathfrak m$ son idéal maximal,
$X=\operatorname{Spec}(A)$, $Y=X-\{\mathfrak m\}$~; alors $Y$ est un schéma de
Jacobson, dont les points fermés sont les idéaux premiers
$\mathfrak p\in\operatorname{Spec}(A)$ tels que $\dim(A/\mathfrak p)=1$.
\end{corollary}

La première assertion est un cas particulier de (10.5.8)~; d'autre part les
points fermés de $Y$ sont les idéaux premiers $\mathfrak p$ de $A$ qui sont
des éléments maximaux dans l'ensemble des idéaux premiers $\neq\mathfrak m$,
ce qui, par définition de la dimension, signifie que
$\dim(A/\mathfrak p)=1$.

\begin{proposition}[10.5.10]
\label{IV.10.5.10-fr}
Soit $A$ un anneau local noethérien réduit et complet et qui n'est pas un
corps. Alors les intersections finies des noyaux des homomorphismes locaux de
$A$ dans des anneaux de valuation discrète $V$, qui font de $V$ une
$A$-algèbre finie, forment une base de filtre tendant vers $0$ pour la
topologie adique de $A$.
\end{proposition}

Il suffit (Bourbaki, \emph{Alg. comm.}, chap.~III, \S~2,
n\textsuperscript{o}~7, prop.~8) de prouver que l'intersection des noyaux
envisagés dans l'énoncé est réduite à $0$. Supposons d'abord que $A$ soit
intègre et de dimension $1$~; en vertu du théorème de Nagata ($\mathbf 0$,
23.1.5 et 23.1.6), la clôture intégrale $A'$ de $A$ est un anneau local
complet intègre et

\oldpage[IV-3]{107}
intégralement clos et de dimension $1$ et une $A$-algèbre finie~; c'est donc
un anneau de valuation discrète, et la proposition en résulte aussitôt dans
ce cas.

Passons au cas général~; soit $x$ le point fermé de
$X=\operatorname{Spec}(A)$, et posons $Y=X-\{x\}$~; on sait (10.5.9) que $Y$
est un préschéma de Jacobson. Montrons que cela entraîne que l'intersection
des idéaux premiers $\mathfrak p$ de $A$ tels que
$\dim(A/\mathfrak p)=1$ est \emph{réduite à $0$}~; la proposition en résultera
puisque, pour chacun de ces idéaux $\mathfrak p$, l'intersection des noyaux des
homomorphismes locaux de $A/\mathfrak p$ dans des anneaux de valuation
discrète $V$, faisant de $V$ une $(A/\mathfrak p)$-algèbre finie, est réduite
à $0$. Mais dire que l'intersection de ces idéaux premiers est réduite à $0$
signifie que l'ensemble de ces idéaux est dense dans $X$, ou encore dans $Y$
(puisque $Y$ est dense dans $X$), et cela résulte aussitôt de (10.5.9).

\subsection{Dimension dans les préschémas de Jacobson}
\label{subsection:IV.10.6-fr}

Les résultats de ce numéro précisent dans certains cas et généralisent des résultats
du \S~5.

\begin{proposition}[10.6.1]
\label{IV.10.6.1-fr}
Soit $S$ un préschéma localement noethérien, vérifiant en outre les conditions
suivantes~: $1^\circ$ $S$ est un préschéma de Jacobson~; $2^\circ$ pour tout
$s\in S$, $\mathcal O_{S,s}$ est universellement caténaire (5.6.2)~; $3^\circ$ toute
composante irréductible $S'$ de $S$ est équicodimensionnelle (autrement dit, pour
tout point fermé $s$ de $S'$ et tout sous-préschéma de $S$ ayant $S'$ pour espace
sous-jacent, on a $\dim(S')=\dim(\mathcal O_{S',s})$). On a alors les propriétés
suivantes~:
\begin{enumerate}[label=(\roman*)]
  \item Pour tout morphisme localement de type fini $g:X\to S$, $X$ vérifie les
  conditions $1^\circ$, $2^\circ$ et $3^\circ$ précédentes. En particulier, si $X$
  est équidimensionnel (par exemple si $X$ est irréductible), $X$ est
  biéquidimensionnel (autrement dit, $X$ est caténaire et pour tout point fermé
  $x$ de $X$, on a $\dim(\mathcal O_{X,x})=\dim(X)$ ($\mathbf 0$, 14.3.3)).
  \item Soient $X$, $Y$ deux $S$-préschémas localement de type fini sur $S$,
  $f:X\to Y$ un $S$-morphisme~; on suppose $X$ irréductible et $f$ dominant. Si
  $\xi$ (resp. $\eta$) est le point générique de $X$ (resp. $Y$) et
  $e=\dim(f^{-1}(\eta))=\deg.\operatorname{tr}_{k(\eta)}k(\xi)$, on a
  \begin{equation}
  \label{IV.10.6.1.1-fr}
  \tag{10.6.1.1}
  \dim(X)=\dim(Y)+e.
  \end{equation}
  \item Soient $X$, $Y$ deux $S$-préschémas localement de type fini sur $S$,
  $f:X\to Y$ un $S$-morphisme, $n$ un entier $\geq0$ tel que l'on ait
  $\dim(f^{-1}(y))\geq n$ (resp. $\dim(f^{-1}(y))\leq n$) pour tout $y\in Y$.
  Alors on a
  \begin{equation}
  \label{IV.10.6.1.2-fr}
  \tag{10.6.1.2}
  \dim(X)\geq\dim(Y)+n
  \end{equation}
  (resp.
  \begin{equation}
  \label{IV.10.6.1.3-fr}
  \tag{10.6.1.3}
  \dim(X)\leq\dim(Y)+n).
  \end{equation}
\end{enumerate}
\end{proposition}

(i) La propriété $1^\circ$ pour $X$ résulte de (10.4.7). Pour tout $x\in X$,
$\mathcal O_{X,x}$ est anneau local d'une $\mathcal O_{S,g(x)}$-algèbre de type fini en
un idéal premier, et l'homomorphisme
$\mathcal O_{S,g(x)}\to\mathcal O_{X,x}$

\oldpage[IV-3]{108}
est local~; donc (5.6.3, (iv)) $\mathcal O_{X,x}$ est universellement caténaire. Pour
démontrer que $X$ vérifie la condition $3^\circ$, considérons plusieurs cas~:

\emph{a)} $X$ est un sous-préschéma fermé irréductible de $S$~; soient $S'$ une
composante irréductible de $S$ contenant $X$, $\xi$ le point générique de $X$,
$x$ un point fermé de $X$~; pour tout $s\in S'$, $\mathcal O_{S',s}$, quotient de
$\mathcal O_{S,s}$, est caténaire (5.6.1), donc les conditions $2^\circ$ et $3^\circ$
entraînent que $S'$ est biéquidimensionnel ($\mathbf 0$, 14.3.3). En vertu de
(5.1.2) et de ($\mathbf 0$, 14.3.3.2), on a donc
\[
  \dim(\mathcal O_{X,x})=\dim(\mathcal O_{S',x})-\dim(\mathcal O_{S',\xi})
  =\dim(S')-\dim(\mathcal O_{S',\xi})
\]
ce qui prouve que $\dim(\mathcal O_{X,x})$ ne dépend pas du point fermé $x$
considéré, d'où l'assertion dans ce cas (5.1.4).

\emph{b)} $X$ est irréductible et $g$ dominant. Alors, pour tout point fermé $x$
de $X$, $g(x)$ est fermé dans $S$ par (10.4.7)~; comme $\mathcal O_{S,g(x)}$ est
universellement caténaire, il résulte de (5.6.5.3) que l'on a
\[
  \dim(\mathcal O_{X,x})=\dim(\mathcal O_{S,g(x)})+e
\]
où $e=\dim(g^{-1}(\zeta))$, $\zeta$ étant le point générique de $S$. Comme
$\dim(S)=\dim(\mathcal O_{S,g(x)})$ en vertu de la condition $3^\circ$ pour $S$, on
a $\dim(\mathcal O_{X,x})=\dim(S)+e$ pour tout point fermé $x\in X$~; cela prouve
la condition $3^\circ$ pour $X$ (5.1.4), et en même temps la formule
\begin{equation}
\label{IV.10.6.1.4-fr}
\tag{10.6.1.4}
\dim(X)=\dim(S)+e.
\end{equation}

\emph{c)} Cas général. En considérant un sous-préschéma réduit $X'$ de $X$ ayant
pour espace sous-jacent une composante irréductible de $X$, on est ramené au cas
où $X$ est intègre~; utilisant (\textbf{I}, 5.2.2) et le cas \emph{a)} prouvé
ci-dessus, on peut ensuite remplacer $S$ par le sous-préschéma réduit ayant
$\overline{g(X)}$ pour espace sous-jacent~; on est alors ramené au cas \emph{b)},
et cela achève de démontrer (i).

(ii) Le morphisme $f$ étant localement de type fini (\textbf{I}, 1.3.4), on peut
appliquer les résultats de (i) en remplaçant $S$ par $Y$~; en outre, comme $X$
est irréductible et $f$ dominant, on peut aussi remplacer $S$ par $Y$ dans
(10.6.1.4), ce qui donne (10.6.1.1).

(iii) L'assertion relative au cas où $\dim(f^{-1}(y))\leq n$ pour tout $y\in Y$
a déjà été démontrée sous des hypothèses plus générales dans (5.6.7). Supposons
que $\dim(f^{-1}(y))\geq n$ pour tout $y\in Y$, et considérons un point générique
$\eta$ d'une composante irréductible $Y'$ de $Y$~; il existe au moins une
composante irréductible $Z$ de $f^{-1}(\eta)$ de dimension $\geq n$~; si $\xi$
est le point générique de $Z$, $\xi$ est aussi point générique d'une composante
irréductible $X'$ de $X$ telle que $f(X')$ soit dense dans $Y'$ ($\mathbf 0_{\mathrm I}$,
2.1.8). Considérons les sous-préschémas réduits de $X$, $Y$ ayant respectivement
pour espaces sous-jacents $X'$, $Y'$, et la restriction $X'\to Y'$ de $f$
(\textbf{I}, 5.2.2)~; il résulte alors de (ii) que
$\dim(X')\geq\dim(Y')+n$, et a fortiori $\dim(X)\geq\dim(Y')+n$~; ceci étant vrai
pour toute composante irréductible $Y'$ de $Y$, on en conclut l'inégalité
(10.6.1.2).

\begin{corollary}[10.6.2]
\label{IV.10.6.2-fr}
Supposons que $X$ vérifie les conditions $1^\circ$, $2^\circ$ et $3^\circ$ de
(10.6.1). Alors, pour tout ouvert $U$ dense dans $X$, on a
$\dim(U)=\dim(X)$.
\end{corollary}

\oldpage[IV-3]{109}
On sait en effet que $\dim(U)\leq\dim(X)$ ($\mathbf 0$, 14.1.4)~; en outre,
comme $X$ est un préschéma de Jacobson, $U$ contient un point fermé de toute
composante irréductible de $X$, donc $\dim(U)=\dim(X)$ en vertu de (10.6.1) et
($\mathbf 0$, 14.1.2.1).

\begin{proposition}[10.6.3]
\label{IV.10.6.3-fr}
Supposons que $X$ vérifie les conditions $1^\circ$, $2^\circ$ et $3^\circ$ de
(10.6.1), et soit $Y$ une partie fermée de $X$. Alors, pour tout $x\in Y$, et
tout voisinage ouvert $U$ de $x$ dans $X$ ne rencontrant pas les composantes
irréductibles de $Y$ qui ne contiennent pas $x$, on a
\begin{equation}
\label{IV.10.6.3.1-fr}
\tag{10.6.3.1}
\dim_x(Y)=\dim(U\cap Y)=\dim(\overline{\{x\}})
 +\operatorname{codim}(\overline{\{x\}},Y)=\sup_i\dim(Y_i)
\end{equation}
où $Y_i$ ($1\leq i\leq m$) sont les composantes irréductibles de $Y$ contenant
$x$.
\end{proposition}

En considérant un sous-préschéma fermé de $X$ ayant $Y$ pour espace sous-jacent,
on peut se borner, en vertu de (10.6.1, (i)), au cas où $Y=X$. En vertu du choix
de $U$, $U$ est réunion des $U\cap X_i$, donc
$\dim(U)=\sup_i\dim(U\cap X_i)$~; mais d'après (10.6.2) appliqué à un
sous-préschéma fermé de $X$ ayant $X_i$ pour espace sous-jacent, on a (compte
tenu de (10.6.1, (i))) $\dim(U\cap X_i)=\dim(X_i)$~; cela démontre que le second
et le quatrième terme de (10.6.3.1) sont égaux. Comme $U\cap X_i$ est
biéquidimensionnel (10.6.1, (i)) (puisque l'immersion $U\cap X_i\to X_i$ est de
type fini (\textbf{I}, 6.3.5)), on a
\begin{equation}
\label{IV.10.6.3.2-fr}
\tag{10.6.3.2}
\dim(X_i)=\dim(U\cap X_i)=\dim(U\cap\overline{\{x\}})
 +\operatorname{codim}(U\cap\overline{\{x\}},U\cap X_i)
\end{equation}
($\mathbf 0$, 14.3.5). D'après (10.6.2) appliqué à un sous-préschéma fermé de
$X$ ayant $\overline{\{x\}}$ pour espace sous-jacent,
$\dim(U\cap\overline{\{x\}})=\dim(\overline{\{x\}})$~; comme on a aussi, pour les
mêmes raisons,
\begin{equation}
\label{IV.10.6.3.3-fr}
\tag{10.6.3.3}
\dim(X_i)=\dim(\overline{\{x\}})
 +\operatorname{codim}(\overline{\{x\}},X_i)
\end{equation}
on obtient
\[
  \dim(U)=\dim(\overline{\{x\}})
  +\sup_i\operatorname{codim}(\overline{\{x\}},X_i)
  =\dim(\overline{\{x\}})+\operatorname{codim}(\overline{\{x\}},X)
\]
par définition de la codimension ($\mathbf 0$, 14.2.1). Ceci montre que $\dim(U)$
est indépendante du voisinage ouvert $U$ de $x$ vérifiant les conditions de
l'énoncé, donc est égal à $\dim_x(X)$, par ($\mathbf 0$, 14.1.4.1).

\begin{corollary}[10.6.4]
\label{IV.10.6.4-fr}
Sous les hypothèses de (10.6.3), soient $\mathcal F$ un $\mathcal O_X$-Module cohérent,
$Y$ son support. Pour tout $x\in Y$, on a
\begin{equation}
\label{IV.10.6.4.1-fr}
\tag{10.6.4.1}
\dim(\overline{\{x\}})+\dim(\mathcal F_x)=\dim_xY.
\end{equation}
\end{corollary}

En effet, cela résulte de (10.6.3.1) et de la formule
$\dim(\mathcal F_x)=\operatorname{codim}(\overline{\{x\}},Y)$ (5.1.12.2).

\subsection{Exemples et contre-exemples}
\label{subsection:IV.10.7-fr}

\begin{env}[10.7.1]
\label{IV.10.7.1-fr}
Soit $S$ un préschéma localement noethérien de dimension $\leq1$ et supposons
que $S$ soit un \emph{préschéma de Jacobson}~; lorsque $S$ est noethérien, il
revient au même de dire que les composantes irréductibles de $S$ de dimension
$1$ sont infinies, car tout $x\in S$ qui n'est pas fermé est le point générique
d'une telle composante (10.4.5 et 10.5.4). Alors $S$ vérifie aussi les conditions
$2^\circ$ et $3^\circ$ de (10.6.1)~: en effet, tout anneau local $\mathcal O_{S,s}$
est de dimension $0$ ou $1$, et par suite est universellement caténaire (7.2.9)~;
d'autre part, une composante irréductible $S'$ de $S$ est, soit réduite

\oldpage[IV-3]{110}
à un point, soit de dimension $1$, et pour tout point fermé $s\in S'$,
$\mathcal O_{S',s}$ est nécessairement de dimension $1$.

On déduit de ces remarques et de (10.6.1) que tout préschéma localement de type
fini sur $S$ vérifie aussi les propriétés $1^\circ$, $2^\circ$ et $3^\circ$ de
(10.6.1)~: il en est ainsi en particulier des préschémas \emph{localement de type
fini sur un corps ou sur $\mathbf Z$}.
\end{env}

\begin{env}[10.7.2]
\label{IV.10.7.2-fr}
Soit $A$ un anneau local noethérien \emph{universellement caténaire} et soit $S$
le \emph{complémentaire dans $X=\operatorname{Spec}(A)$ du point fermé} $a$.
Alors $S$ vérifie les conditions $1^\circ$, $2^\circ$ et $3^\circ$ de (10.6.1)~:
en effet, on a déjà vu que $S$ est un préschéma de Jacobson (10.5.9)~; comme $A$
est universellement caténaire, il en est de même des anneaux locaux $A_{\mathfrak p}$
aux idéaux premiers de $A$ (5.6.3). D'autre part, une composante irréductible $S'$
de $S$ est le complémentaire de $a$ dans une composante irréductible $X'$ de $X$~;
pour tout point fermé $x$ de $S$, l'adhérence de $x$ dans $X$ est donc
$\{x,a\}$, autrement dit $\dim(\overline{\{x\}})=1$ et par suite, puisque $X'$
est par hypothèse biéquidimensionnel ($\mathbf 0$, 14.3.3), on a
$\dim(\mathcal O_{X',x})=\dim(X')-1$ ($\mathbf 0$, 14.3.2), ce qui prouve que $S'$
est équicodimensionnel (5.1.4).
\end{env}

\begin{env}[10.7.3]
\label{IV.10.7.3-fr}
Soit $A$ un anneau de valuation discrète~; montrons que dans l'anneau
$B=A[T_1,\ldots,T_n]$ des polynômes à $n$ indéterminées, il existe deux idéaux
maximaux $\mathfrak m$, $\mathfrak n$ de \emph{hauteurs respectives} $n$ et
$n+1$. On l'a vu dans (5.2.5, (i)) pour $n=1$~; démontrons-le par récurrence sur
$n$. Comme $A[T_1,\ldots,T_{n+1}]$ est un $A[T_1,\ldots,T_n]$-module libre,
donc fidèlement plat, il y a dans $A[T_1,\ldots,T_{n+1}]$ deux idéaux maximaux
$\mathfrak m'$, $\mathfrak n'$ respectivement au-dessus de $\mathfrak m$ et
$\mathfrak n$ ($\mathbf 0_{\mathrm I}$, 6.5.1)~; en outre, d'après (5.5.3), ces
idéaux sont nécessairement de hauteurs respectives $n+1$ et $n+2$, d'où notre
assertion. Supposons $n\geq2$ dans ce qui suit. Soit
$\mathfrak J=\mathfrak m\cap\mathfrak n=\mathfrak m\mathfrak n$, et
$R=1+\mathfrak J$ qui est une partie multiplicative de $B$~; si l'on pose
$B'=R^{-1}B$, l'idéal $\mathfrak J'=R^{-1}\mathfrak J$ est contenu dans le
radical de $B'$ (Bourbaki, \emph{Alg. comm.}, chap.~III, \S~3,
n\textsuperscript{o}~5, prop.~12)~; on sait que $X'=\operatorname{Spec}(B')$
s'identifie en tant qu'espace topologique à un sous-espace de
$X=\operatorname{Spec}(B)$, et qu'aux points $x$ de $X'$, les anneaux locaux
$\mathcal O_{X,x}$ et $\mathcal O_{X',x}$ sont les mêmes (\textbf{I}, 1.6.2).
Considérons alors dans $X'$ l'ensemble fermé $Y'=V(\mathfrak J')$, et soit
$S=X'-Y'$~; on sait (10.5.7) que $S$ est un préschéma de Jacobson, évidemment
irréductible et noethérien~; en outre les anneaux locaux
$\mathcal O_{S,s}=\mathcal O_{X,s}$ sont universellement caténaires pour tout $s\in S$
en vertu de (5.6.3), puisque $A$ est universellement caténaire (5.6.4). Pourtant,
il y a deux points fermés $a$, $b$ de $S$ tels que $\mathcal O_{S,a}$ et
$\mathcal O_{S,b}$ \emph{n'aient pas même dimension}, autrement dit $S$ ne vérifie
pas la condition $3^\circ$ de (10.6.1). Pour le voir, considérons les deux idéaux
maximaux $\mathfrak m'=R^{-1}\mathfrak m$, $\mathfrak n'=R^{-1}\mathfrak n$ de
$B'$, qui sont de hauteurs $n$ et $n+1$ respectivement~; on a
$\mathfrak J'=\mathfrak m'\cap\mathfrak n'$, et $\mathfrak J'$ n'est donc
contenu dans aucun idéal premier de $B'$ distinct de $\mathfrak m'$ et
$\mathfrak n'$, qui sont par suite les seuls idéaux maximaux de $B'$. Soient
$a'$, $b'$ les seuls points fermés de $X'$, correspondant à $\mathfrak m'$ et
$\mathfrak n'$. Il existe dans $B'$ un idéal premier non maximal
$\mathfrak p'\subset\mathfrak m'$ qui n'est pas contenu dans $\mathfrak n'$~:
il suffit de montrer qu'il y a dans $B$ un idéal premier non maximal contenu dans
$\mathfrak m$ et non dans $\mathfrak n$~; pour cela, on pourra par exemple
considérer la fibre de $\mathfrak m$ pour le morphisme correspondant à l'injection
$A[T_1]\to B=A[T_1,\ldots,T_n]$, et appliquer (6.1.2). En considérant une chaîne
maximale d'idéaux premiers entre $\mathfrak p'$ et $\mathfrak m'$ et remplaçant
$\mathfrak p'$ par l'avant-dernier idéal de cette chaîne, on peut donc supposer
que le point $a$ de $X'$ correspondant à $\mathfrak p'$ est tel que son adhérence
dans $X'$ soit $\{a,a'\}$~; comme $B'_{\mathfrak m'}$ est biéquidimensionnel,
$\mathfrak p'$ est alors de hauteur $n-1$. On construit de la même manière un
idéal premier non maximal $\mathfrak q'$ de $B'$ de hauteur $n$, tel que si $b$
est le point correspondant de $X'$, l'adhérence de $b$ dans $X'$ soit
$\{b,b'\}$. Cela étant, $a$ et $b$ sont dans $S$, donc \emph{fermés dans $S$},
et répondent par suite à la question.
\end{env}

\subsection{Profondeur rectifiée}
\label{subsection:IV.10.8-fr}

\begin{definition}[10.8.1]
\label{IV.10.8.1-fr}
Soient $X$ un préschéma localement noethérien, $\mathcal F$ un $\mathcal O_X$-Module
cohérent. Pour tout $x\in X$, on appelle \emph{profondeur rectifiée de $\mathcal F$
au point $x$} et on note $\operatorname{prof}^{*}_{x}(\mathcal F)$ le nombre (entier
$\geq0$ ou $+\infty$) égal à
\begin{equation}
\label{IV.10.8.1.1-fr}
\tag{10.8.1.1}
\operatorname{prof}^{*}_{x}(\mathcal F)
=\operatorname{prof}(\mathcal F_x)+\dim(\overline{\{x\}})
\end{equation}
où $\overline{\{x\}}$ est l'adhérence du point $x$ dans $X$. Pour toute partie
$Z$ de $X$, on appelle \emph{profondeur rectifiée de $\mathcal F$ le long de $Z$}
et on note $\operatorname{prof}^{*}_{Z}(\mathcal F)$ le nombre
\begin{equation}
\label{IV.10.8.1.2-fr}
\tag{10.8.1.2}
\operatorname{prof}^{*}_{Z}(\mathcal F)
=\inf_{x\in Z}\operatorname{prof}^{*}_{x}(\mathcal F).
\end{equation}
\end{definition}

En d'autres termes pour tout entier $n$, la relation
$\operatorname{prof}^{*}_{Z}(\mathcal F)\geq n$ équivaut à
$\operatorname{prof}^{*}_{x}(\mathcal F)\geq n$ pour tout $x\in Z$. Si $Z=X$, on
écrit $\operatorname{prof}^{*}(\mathcal F)$ au lieu de
$\operatorname{prof}^{*}_{X}(\mathcal F)$.

\oldpage[IV-3]{111}
\begin{remarks}[10.8.2]
\label{IV.10.8.2-fr}
(i) En tout point fermé $x\in X$, la profondeur rectifiée est égale à la
profondeur.

(ii) Soient $Y$ un sous-préschéma fermé de $X$, $j:Y\to X$ l'injection canonique,
$\mathcal G$ un $\mathcal O_Y$-Module cohérent. On sait (5.7.3, (vi)) que l'on a
$\operatorname{prof}(\mathcal G_x)=\operatorname{prof}(j_*(\mathcal G)_x)$ pour tout
$x\in Y$~; on en déduit que l'on a aussi
$\operatorname{prof}^{*}_{x}(\mathcal G)
=\operatorname{prof}^{*}_{x}(j_*(\mathcal G))$ pour tout $x\in Y$.

(iii) La notion de profondeur rectifiée n'a d'intérêt que lorsqu'elle est de
caractère \emph{local}, c'est-à-dire qu'elle ne change pas lorsqu'on remplace $X$
par un voisinage ouvert arbitraire $U$ de $x$. Cela exige évidemment que $x$ ne
soit pas isolé dans $\overline{\{x\}}$ lorsque $x$ n'est pas fermé et par suite
que $X$ soit un \emph{préschéma de Jacobson} (10.4.5.1)~; le plus souvent, il sera
aussi nécessaire de savoir que $\dim(U)=\dim(X)$ pour tout ouvert dense $U$ dans
$X$, et il faudra donc supposer que $X$ vérifie aussi les conditions $2^\circ$ et
$3^\circ$ de (10.6.1).
\end{remarks}

\begin{lemma}[10.8.3]
\label{IV.10.8.3-fr}
Soient $X$ un préschéma régulier et biéquidimensionnel, $\mathcal F$ un
$\mathcal O_X$-Module cohérent. Alors on a, pour tout $x\in X$,
\begin{equation}
\label{IV.10.8.3.1-fr}
\tag{10.8.3.1}
\operatorname{prof}^{*}_{x}(\mathcal F)
=\dim(X)-\dim.\operatorname{proj}(\mathcal F_x).
\end{equation}
\end{lemma}

En effet, comme $X$ est biéquidimensionnel, on a ($\mathbf 0$, 14.3.5.1)
\[
  \dim(\overline{\{x\}})=\dim(X)-\operatorname{codim}(\overline{\{x\}},X)
  =\dim(X)-\dim(\mathcal O_{X,x})
\]
en vertu de (5.1.2). D'autre part, comme $X$ est régulier, on a par
($\mathbf 0$, 17.3.4)
\[
  \operatorname{prof}(\mathcal F_x)
  =\dim(\mathcal O_{X,x})-\dim.\operatorname{proj}(\mathcal F_x)
\]
d'où le lemme.

\begin{corollary}[10.8.4]
\label{IV.10.8.4-fr}
Sous les hypothèses de (10.8.3), la fonction
$x\mapsto\operatorname{prof}^{*}_{x}(\mathcal F)$ est semi-continue inférieurement.
\end{corollary}

Cela résulte de (10.8.3.1), puisque
$x\mapsto\dim.\operatorname{proj}(\mathcal F_x)$ est semi-continue supérieurement
(6.11.1).

\begin{proposition}[10.8.5]
\label{IV.10.8.5-fr}
Soient $S$ un préschéma localement noethérien, $X$ un préschéma localement de
type fini sur $S$, $\mathcal F$ un $\mathcal O_X$-Module cohérent. On suppose que $S$
vérifie les conditions suivantes~: $1^\circ$ $S$ est un préschéma de Jacobson~;
$2^\circ$ $S$ est régulier~; $3^\circ$ les composantes irréductibles de $S$ sont
équicodimensionnelles. Alors la fonction
$x\mapsto\operatorname{prof}^{*}_{x}(\mathcal F)$ est semi-continue inférieurement
dans $X$~; autrement dit, pour tout entier $n$, l'ensemble $U_n$ des $x\in X$
tels que $\operatorname{prof}^{*}_{x}(\mathcal F)\geq n$ est ouvert.
\end{proposition}

Comme les anneaux locaux $\mathcal O_{S,s}$ de $S$ sont réguliers, ils sont
universellement caténaires (5.6.4)~; autrement dit, $S$ vérifie les conditions
$1^\circ$, $2^\circ$ et $3^\circ$ de (10.6.1), donc il en est de même de $X$
(10.6.1, (i)). La notion de profondeur rectifiée étant alors de caractère local
(10.8.2, (iii)), on peut se borner au cas où $S=\operatorname{Spec}(A)$ et
$X=\operatorname{Spec}(B)$ sont affines, $A$ étant un anneau régulier et $B$ une
$A$-algèbre de type fini, donc quotient d'un anneau de polynômes
$C=A[T_1,\ldots,T_n]$, et ce dernier est régulier ($\mathbf 0$, 17.3.7). On peut
donc supposer que $X$ est un sous-préschéma fermé d'un préschéma régulier $Y$
vérifiant également les conditions $1^\circ$, $2^\circ$ et $3^\circ$ de (10.6.1)~;
compte tenu de la remarque (10.8.2, (ii)), on est ainsi ramené au cas où $X$ est
en outre régulier et noethérien. Mais comme les anneaux locaux de $X$ sont alors
intègres, les composantes irréductibles de $X$ sont

\oldpage[IV-3]{112}
ouvertes (\textbf{I}, 6.1.10), et on peut par suite supposer aussi $X$
irréductible. Alors, comme les anneaux locaux de $X$ sont caténaires
($\mathbf 0$, 16.5.12), l'hypothèse $3^\circ$ de (10.6.1) entraîne que $X$ est
biéquidimensionnel ((5.1.5) et ($\mathbf 0$, 14.3.3))~; il suffit donc d'appliquer
(10.8.4).

On notera que si $S$ est le \emph{spectre d'un corps ou de $\mathbf Z$}, il
vérifie les conditions de (10.8.5).

\begin{corollary}[10.8.6]
\label{IV.10.8.6-fr}
Les hypothèses étant celles de (10.8.5), pour tout $x\in X$, le nombre
$\operatorname{prof}^{*}_{x}(\mathcal F)$ est l'unique entier $n$ ayant la propriété
suivante~: il existe un voisinage ouvert $U$ de $x$ dans $X$ tel que pour tout
point $x'\in U\cap\overline{\{x\}}$, fermé dans $U$, on ait
$\operatorname{prof}(\mathcal F_{x'})=n$. En particulier, pour que
$\operatorname{prof}^{*}_{x}(\mathcal F)\geq m$ (resp.
$\operatorname{prof}^{*}_{x}(\mathcal F)\leq m$), il faut et il suffit qu'il existe
un voisinage ouvert $V$ de $x$ dans $X$ tel que, pour tout
$x'\in V\cap\overline{\{x\}}$, fermé dans $V$, on ait
$\operatorname{prof}(\mathcal F_{x'})\geq m$ (resp.
$\operatorname{prof}(\mathcal F_{x'})\leq m$).
\end{corollary}

En effet, on peut se borner au cas où $x$ n'est pas fermé~; si
$\operatorname{prof}^{*}_{x}(\mathcal F)=n$, l'ensemble des $y\in X$ tels que
$\operatorname{prof}^{*}_{y}(\mathcal F)\leq n$ est fermé en vertu de (10.8.5), donc
contient $\overline{\{x\}}$, et en vertu de la semi-continuité inférieure de
$y\mapsto\operatorname{prof}^{*}_{y}(\mathcal F)$, il existe un voisinage ouvert $U$
de $x$ tel que $\operatorname{prof}^{*}_{y}(\mathcal F)=n$ pour tout
$y\in U\cap\overline{\{x\}}$, donc $\operatorname{prof}(\mathcal F_{x'})=n$ si
$x'\in U\cap\overline{\{x\}}$ est fermé dans $U$ (puisque la notion de profondeur
rectifiée est locale).

Pour les préschémas vérifiant les hypothèses de (10.8.5), la notion de profondeur
rectifiée peut donc se définir à l'aide des valeurs de la profondeur aux
\emph{points fermés} de $X$ (ces derniers formant un ensemble \emph{très dense}
dans toute partie fermée de $X$).

\begin{proposition}[10.8.7]
\label{IV.10.8.7-fr}
Soit $S$ un préschéma localement noethérien vérifiant les conditions $1^\circ$,
$2^\circ$ et $3^\circ$ de (10.6.1). Soient $X$ un préschéma localement de type
fini sur $S$, $\mathcal F$ un $\mathcal O_X$-Module cohérent,
$Y=\operatorname{Supp}(\mathcal F)$. Alors, pour tout $x\in Y$, on a
\begin{equation}
\label{IV.10.8.7.1-fr}
\tag{10.8.7.1}
\operatorname{prof}^{*}_{x}(\mathcal F)
=\dim_x(Y)-\operatorname{coprof}(\mathcal F_x).
\end{equation}
\end{proposition}

En effet, par définition, on a
$\operatorname{coprof}(\mathcal F_x)
=\dim(\mathcal F_x)-\operatorname{prof}(\mathcal F_x)$, et il résulte de (10.6.4) que
l'on a $\dim_x(Y)=\dim(\overline{\{x\}})+\dim(\mathcal F_x)$~; d'où (10.8.7.1) par
définition de $\operatorname{prof}^{*}_{x}(\mathcal F)$.

\begin{corollary}[10.8.8]
\label{IV.10.8.8-fr}
Les hypothèses sur $f$ et $\mathcal F$ étant celles de (9.9.1), la fonction
$x\mapsto\operatorname{prof}^{*}_{x}(\mathcal F_{f(x)})$ est constructible.
\end{corollary}

Notons que pour tout $s\in S$, la fibre $X_s$ étant un préschéma localement de
type fini sur $k(s)$, est un préschéma de Jacobson (10.4.7)~; comme en outre
$\operatorname{Spec}(k(s))$ vérifie les conditions de (10.6.1), on a
\[
  \operatorname{prof}^{*}_{x}(\mathcal F_{f(x)})
  =\dim_x(\operatorname{Supp}(\mathcal F_{f(x)}))
  -\operatorname{coprof}((\mathcal F_{f(x)})_x)
\]
(10.8.7). Or, si $Z=\operatorname{Supp}(\mathcal F)$, on a
$Z_{f(x)}=\operatorname{Supp}(\mathcal F_{f(x)})$ (\textbf{I}, 9.1.13) et $Z$ est
localement constructible (8.9.1)~; donc les fonctions
$x\mapsto\dim_x(\operatorname{Supp}(\mathcal F_{f(x)}))$ et
$x\mapsto\operatorname{coprof}((\mathcal F_{f(x)})_x)$ sont localement
constructibles ((9.9.1) et (9.9.3)), ce qui prouve la proposition.

\subsection{Spectres maximaux et ultrapréschémas}
\label{subsection:IV.10.9-fr}

Les résultats de ce numéro ne seront pas utilisés par la suite.

\begin{env}[10.9.1]
\label{IV.10.9.1-fr}
Soit $X$ un préschéma de Jacobson, et soit $S(X)$ l'\emph{espace annelé} dont
l'espace sous-jacent est le sous-espace des points fermés de $X$, et le faisceau
d'anneaux le faisceau induit sur ce sous-espace par $\mathcal O_X$, autrement dit
le faisceau d'anneaux $\psi^*(\mathcal O_X)$, en désignant par
$\psi:S(X)\to X$ l'injection canonique. Comme $\psi$ est un quasi-
\oldpage[IV-3]{113}
homéomorphisme, on a vu (10.2.8, (ii)) que si
$\theta:\mathcal O_X\to\psi_*(\mathcal O_{S(X)})$ est l'homomorphisme de
faisceaux d'anneaux tel que
$\theta^\sharp:\psi^*(\mathcal O_X)\to\mathcal O_{S(X)}$ soit l'identité,
alors $j_X=(\psi,\theta)$ est un \emph{quasi-isomorphisme} d'espaces annelés,
et $\mathcal F\mapsto j_X^*(\mathcal F)$ une équivalence de la catégorie des
$\mathcal O_X$-Modules et de celle des $\mathcal O_{S(X)}$-Modules. Il est clair
que dans cette équivalence, aux $\mathcal O_X$-Modules localement libres (resp.
cohérents) correspondent les $\mathcal O_{S(X)}$-Modules localement libres
(resp. cohérents)~; en outre, si $\mathcal L$ est un $\mathcal O_X$-Module
localement libre et si $U$ est un ouvert de $X$ tel que $\mathcal L|U$ soit
isomorphe à $(\mathcal O_X|U)^n$, $j_X^*(\mathcal L)$ est tel que
$j_X^*(\mathcal L)|(U\cap S(X))$ soit isomorphe à
$(\mathcal O_{S(X)}|(U\cap S(X)))^n$.
\end{env}

\begin{env}[10.9.2]
\label{IV.10.9.2-fr}
Soient $X$, $Y$ deux préschémas de Jacobson et
$f=(\rho,\lambda):X\to Y$ un morphisme \emph{localement de type fini}. On a vu
(10.4.7) que l'on a $\rho(S(X))\subset S(Y)$ et par restriction de $\rho$ à
$S(X)$, on définit donc une application continue $S(\rho):S(X)\to S(Y)$.
D'autre part, pour tout ouvert $V$ de $Y$, on définit par composition un
homomorphisme d'anneaux
\[
  \Gamma(V\cap S(Y),\mathcal O_{S(Y)})\longrightarrow
  \Gamma(V,\mathcal O_Y)\xrightarrow{\lambda_V}
  \Gamma(f^{-1}(V),\mathcal O_X)\longrightarrow
  \Gamma(f^{-1}(V)\cap S(X),\mathcal O_{S(X)})
\]
où les deux isomorphismes extrêmes ont été définis dans (10.9.1)~; il est clair
que cela définit un homomorphisme de faisceaux d'anneaux
$S(\lambda):\mathcal O_{S(Y)}\to\rho_*(\mathcal O_{S(X)})$ (en se rappelant
que les ouverts de $S(X)$ (resp. $S(Y)$) correspondent biunivoquement à ceux de
$X$ (resp. $Y$) (10.2.1))~; on obtient ainsi un morphisme d'espaces annelés
\[
  S(f)=(S(\rho),S(\lambda)):(S(X),\mathcal O_{S(X)})\longrightarrow
  (S(Y),\mathcal O_{S(Y)})
\]
tel que le diagramme
\[
\begin{tikzcd}
S(X) \ar[r,"S(f)"] \ar[d,"j_X"'] & S(Y) \ar[d,"j_Y"] \\
X \ar[r,"f"'] & Y
\end{tikzcd}
\]
soit commutatif~; en outre, si $Z$ est un troisième préschéma de Jacobson et
$g:Y\to Z$ un morphisme localement de type fini, il est clair que
$S(g\circ f)=S(g)\circ S(f)$. On a ainsi défini un \emph{foncteur covariant}
$S:\mathbf C\to\mathbf C'$, où $\mathbf C'$ est la catégorie des \emph{espaces
annelés en anneaux locaux}, et $\mathbf C$ la catégorie dont les objets sont les
\emph{préschémas de Jacobson} et les morphismes sont les morphismes
\emph{localement de type fini} entre préschémas de Jacobson.
\end{env}

\begin{env}[10.9.3]
\label{IV.10.9.3-fr}
Proposons-nous de déterminer la sous-catégorie $\mathbf C''$ de $\mathbf C'$
formée des espaces annelés isomorphes aux $S(X)$ et dont les morphismes
proviennent des $S(f)$. Supposons d'abord que $X=\operatorname{Spec}(A)$, où
$A$ est un anneau de Jacobson~; alors $S(X)$ est l'ensemble des \emph{idéaux
maximaux} de $A$, muni~: $1^\circ$ de la topologie induite par celle de $X$, de
sorte qu'une base de cette topologie est formée des
$D^m(h)=D(h)\cap S(X)$, ensemble des idéaux maximaux $\mathfrak m$ de $A$ tels
que $h\notin\mathfrak m$, où $h$ parcourt $A$~; $2^\circ$ du faisceau d'anneaux
$\mathcal O_{S(X)}$ tel que
$\Gamma(D^m(h),\mathcal O_{S(X)})=A_h$. Nous dirons que cet espace annelé est
le \emph{spectre maximal} de l'anneau de Jacobson $A$ et nous le noterons
$\operatorname{Spm}(A)$.

On notera que si $j:D(h)\to X$ est l'injection canonique, l'espace annelé
induit sur $D^m(h)$ par $S(X)$ est $S(D(h))$ et l'injection canonique
$D^m(h)\to S(X)$ d'espaces annelés est égale à $S(j)$.

Soient $B$ un second anneau de Jacobson, $Y=\operatorname{Spec}(B)$,
$\varphi:B\to A$ un homomorphisme d'anneaux faisant de $A$ une $B$-algèbre de
type fini, $f=({}^a\varphi,\widetilde\varphi):X\to Y$ le morphisme
correspondant de préschémas, et
\[
  S(f):\operatorname{Spm}(A)\longrightarrow\operatorname{Spm}(B)
\]
le morphisme d'espaces annelés correspondant à $f$. Il est clair que
$S(f)=(\psi,\theta)$ est un morphisme d'espaces annelés en anneaux locaux,
c'est-à-dire ($\mathrm{Err}_{\mathrm{II}}$, 1.8.2) que pour tout
$x\in\operatorname{Spm}(A)$, $\theta_x^\sharp$ est un homomorphisme local.
Réciproquement~:
\end{env}

\begin{proposition}[10.9.4]
\label{IV.10.9.4-fr}
Soient $A$, $B$ deux anneaux de Jacobson. Si
$u=(\psi,\theta):\operatorname{Spm}(A)\to\operatorname{Spm}(B)$ est un
morphisme d'espaces annelés en anneaux locaux tel que
$\Gamma(\theta):B\to A$ fasse de $A$ une $B$-algèbre de type fini, il existe
un morphisme de préschémas
$f:\operatorname{Spec}(A)\to\operatorname{Spec}(B)$ et un seul tel que
$u=S(f)$.
\end{proposition}

L'unicité de $f$ est évidente, puisque si
$f=({}^a\varphi,\widetilde\varphi)$, on doit avoir
$\varphi=\Gamma(\theta)$~; il reste à voir que $S(f)$ est défini et qu'on a
bien $u=S(f)$. Or, la première assertion résulte de ce que $\varphi$ est
supposé faire de $A$ une $B$-algèbre de type fini, et par suite
$f(\operatorname{Spm}(A))\subset\operatorname{Spm}(B)$~; le fait que
$\theta_x^\sharp$ est un homomorphisme local pour tout
$x\in\operatorname{Spm}(A)$ permet alors de répéter le raisonnement de
(\textbf I, 1.7.3) en se bornant aux points $x$ de $\operatorname{Spm}(A)$~:
on montre ainsi successivement que $\psi(x)={}^a\varphi(x)$ pour tout
$x\in\operatorname{Spm}(A)$, puis que
$\theta_x^\sharp=\varphi_x:B_{\psi(x)}\to A_x$, ce qui achève de prouver que
$u=S(f)$.

\oldpage[IV-3]{114}

\begin{env}[10.9.5]
\label{IV.10.9.5-fr}
Considérons maintenant un espace annelé $(X,\mathcal O_X)$~; nous dirons qu'une
partie ouverte $U$ de $X$ est \emph{ultra-affine} si l'espace annelé induit
$(U,\mathcal O_X|U)$ est isomorphe à un spectre maximal
$\operatorname{Spm}(A)$, où $A$ est un anneau de Jacobson. Nous dirons que $X$
est un \emph{ultra-préschéma} si tout point de $X$ admet un voisinage ouvert
ultra-affine. On montre, comme dans (\textbf I, 2.1.3 et 2.1.4) que les ouverts
ultra-affines forment une \emph{base} de la topologie de $X$ et que $X$ est un
espace de Kolmogoroff. Si $Y$ est un second ultra-préschéma, nous dirons qu'un
morphisme d'espaces annelés $f:X\to Y$ est un \emph{morphisme
d'ultra-préschémas} s'il vérifie les conditions suivantes~: $1^\circ$ $f$ est
un morphisme d'espaces annelés en anneaux locaux~; $2^\circ$ pour tout $x\in X$,
il y a un voisinage ouvert ultra-affine $V$ de $f(x)$ dans $Y$ et un voisinage
ouvert ultra-affine $U$ de $x$ dans $X$ tels que $f(U)\subset V$ et que
l'homomorphisme
$\Gamma(V,\mathcal O_Y)\to\Gamma(U,\mathcal O_X)$ correspondant à $f$ fasse
de $\Gamma(U,\mathcal O_X)$ une $\Gamma(V,\mathcal O_Y)$-algèbre de type fini.

Il est immédiat qu'on définit bien ainsi des morphismes, le composé de deux
morphismes en étant un troisième grâce à la remarque finale de (10.9.3). Il
est clair que la catégorie $\mathbf C''_0$ ainsi définie est une sous-catégorie
de $\mathbf C'$ qui contient $\mathbf C''$~; nous nous proposons de montrer que
$\mathbf C''=\mathbf C''_0$, autrement dit~:
\end{env}

\begin{proposition}[10.9.6]
\label{IV.10.9.6-fr}
Le foncteur $X\mapsto S(X)$ de $\mathbf C$ dans $\mathbf C''_0$ est une
équivalence de catégories.
\end{proposition}

$1^\circ$ Montrons d'abord que le foncteur $X\mapsto S(X)$ est
\emph{pleinement fidèle}, autrement dit que pour $X$, $Y$ dans $\mathbf C$,
l'application canonique
\[
  \operatorname{Hom}_{\mathbf C}(X,Y)\longrightarrow
  \operatorname{Hom}_{\mathbf C''_0}(S(X),S(Y))
\]
est \emph{bijective}. En premier lieu elle est \emph{injective}~: soient en
effet $f$, $g$ deux morphismes localement de type fini de $X$ dans $Y$ et
supposons que $S(f)=S(g)$. Cela entraîne d'abord que pour tout ouvert $V$ de
$Y$, on a
$f^{-1}(V)\cap S(X)=g^{-1}(V)\cap S(X)$, donc (10.3.1 et 10.2.7)
$f^{-1}(V)=g^{-1}(V)$~; il suffit donc de prouver que pour tout ouvert affine
$V$ de $Y$, $f$ et $g$ coïncident dans
$f^{-1}(V)=g^{-1}(V)$, autrement dit, on est ramené au cas où
$Y=\operatorname{Spec}(B)$ est le spectre d'un anneau de Jacobson $B$. Il
suffit (\textbf I, 2.2.4) de montrer que les homomorphismes d'anneaux
$B\to\Gamma(X,\mathcal O_X)$ correspondant à $f$ et $g$ sont alors les mêmes.
Or, pour tout $s\in B$, les images de $s$ par ces homomorphismes sont deux
sections de $\mathcal O_X$ au-dessus de $X$ qui, par hypothèse, induisent la
même section au-dessus de $S(X)$~; on sait donc (10.2.8) que ces sections sont
identiques, d'où notre assertion. Prouvons en second lieu que tout morphisme
$h:S(X)\to S(Y)$ (pour la catégorie $\mathbf C''_0$) est de la forme $S(f)$,
où $f:X\to Y$ est un morphisme localement de type fini. Par hypothèse, il
existe un recouvrement ouvert ultra-affine $(U'_\lambda)$ (resp.
$(V'_\mu)$) de $S(X)$ (resp. $S(Y)$) tel que pour tout $\lambda$,
$h(U'_\lambda)$ soit contenu dans un $V'_\mu$ et qu'il corresponde au
morphisme $h_\lambda:U'_\lambda\to V'_\mu$, restriction de $h$, un
homomorphisme d'anneaux
$\Gamma(V'_\mu,\mathcal O_{S(Y)})\to
\Gamma(U'_\lambda,\mathcal O_{S(X)})$ faisant du second anneau une algèbre de
type fini sur le premier. On peut supposer que
$U'_\lambda=U_\lambda\cap S(X)$ et $V'_\mu=V_\mu\cap S(Y)$, où
$U_\lambda$ et $V_\mu$ sont des ouverts affines uniquement déterminés dans
$X$ et $Y$ respectivement~; si l'on prouve que, pour tout $\lambda$,
$h_\lambda=S(f_\lambda)$ où $f_\lambda:U_\lambda\to V_\mu$ est un morphisme
de type fini, alors il résulte de la première partie du raisonnement,
appliquée aux restrictions de $f_\alpha$ et $f_\beta$ à
$U_\alpha\cap U_\beta$, que les $f_\lambda$ sont les restrictions d'un même
morphisme $f:X\to Y$, et on aura évidemment $h=S(f)$. On est ainsi ramené au
cas où $X$ et $Y$ sont affines, et la conclusion résulte alors de (10.9.4).

$2^\circ$ Il reste à prouver que tout ultra-préschéma $X'$ est de la forme
$S(X)$ pour un préschéma de Jacobson $X$ (qui sera nécessairement unique à
isomorphisme près, en vertu de $1^\circ$). Il y a un recouvrement
$(U'_\alpha)$ de $X'$ par des ouverts ultra-affines, dont chacun est de la
forme $S(U_\alpha)$, $U_\alpha$ étant le spectre d'un anneau de Jacobson. Pour
tout couple d'indices $\alpha$, $\beta$, considérons l'unique ouvert
$U_{\alpha\beta}$ de $U_\alpha$ dont la trace sur $U'_\alpha$ est
$U'_\alpha\cap U'_\beta$~; en vertu de $1^\circ$, l'automorphisme identique de
$U'_\alpha\cap U'_\beta$ est de la forme $S(\theta_{\alpha\beta})$, où
$\theta_{\alpha\beta}:U_{\alpha\beta}\to U_{\beta\alpha}$ est un
isomorphisme de préschémas. On constate aussitôt (en vertu du $1^\circ$) que la
famille $(\theta_{\alpha\beta})$ vérifie la condition de recollement
($\mathbf 0_{\mathrm I}$, 4.1.7), et que cette famille définit donc un
préschéma $X$, tel que les $U_\alpha$ s'identifient à des ouverts affines de
$X$~; il est clair alors que l'on a $X'=S(X)$, ce qui achève la démonstration.

\subsection{Espaces algébriques de Serre}
\label{subsection:IV.10.10-fr}

\begin{env}[10.10.1]
\label{IV.10.10.1-fr}
Le langage introduit par Serre dans (FAC) est parfois commode, notamment dans
les questions où l'intérêt majeur s'attache aux points rationnels sur le corps
de base $k$ (algébriquement clos par hypothèse) des «~variétés algébriques~»
sur $k$ que l'on considère. Nous allons esquisser ici ce langage en le
rattachant aux considérations qui précèdent, pour permettre au lecteur de
traduire les énoncés de Serre en langage des schémas. Il est d'ailleurs
possible de développer le langage de Serre également pour des préschémas sur
un corps non algébriquement clos (et même sur un anneau artinien)~; mais cela
introduit des complications techniques considérables, et d'ailleurs, sur un
corps de base arbitraire, les avantages (surtout psychologiques) du point de
vue de Serre disparaissent~; aussi
\oldpage[IV-3]{115}
nous tiendrons-nous dans le cadre fixé par Serre. Le présent numéro, tout comme
le précédent, ne sera d'ailleurs pas utilisé dans la suite de ce Traité, et
nous nous bornerons donc à de brèves indications.
\end{env}

\begin{env}[10.10.2]
\label{IV.10.10.2-fr}
Étant donné un ultra-préschéma fixe $R$, on peut naturellement (comme dans
toute catégorie) définir la notion de $R$-\emph{ultra-préschéma}. Considérons
en particulier un corps algébriquement clos $k$~; $\operatorname{Spm}(k)$ est
alors identique à $\operatorname{Spec}(k)$~; nous dirons qu'un
$\operatorname{Spm}(k)$-ultra-préschéma $X'$ est un \emph{espace
préalgébrique sur $k$}~: c'est donc un espace $k$-annelé en anneaux locaux dont
tout point a un voisinage ouvert isomorphe au spectre maximal d'une $k$-algèbre
de type fini~; il revient au même de dire (par (10.9.6)) que $X'=S(X)$, où $X$
est un préschéma localement de type fini sur $k$. Si $X$, $Y$, $Z$ sont trois
$k$-préschémas localement de type fini, il en est de même de
$X\times_ZY$ (\textbf I, 3.4), donc en vertu de (10.9.6), la notion de produit
existe dans la catégorie des espaces préalgébriques sur $k$ (aussi bien le
produit «~sur $k$~» $X'\times_kY'$ que le «~produit fibré~»
$X'\times_{Z'}Y'$, où $X'$, $Y'$, $Z'$ sont trois espaces préalgébriques sur
$k$). On peut donc définir le morphisme diagonal
$\Delta_{X'}:X'\to X'\times_kX'$ (qui n'est autre d'ailleurs que
$S(\Delta_X)$)~; on dit que $X'$ est un \emph{espace algébrique sur $k$} si
$X$ est un schéma, et il revient au même de dire que l'image de $\Delta_{X'}$
est une partie \emph{fermée} de $X'\times_kX'$.
\end{env}

\begin{env}[10.10.3]
\label{IV.10.10.3-fr}
Les simplifications qui proviennent de l'hypothèse que $k$ est algébriquement
clos tiennent d'abord à ce que, pour un $k$-préschéma $X$, localement de type
fini, il y a correspondance biunivoque entre \emph{points fermés de $X$},
\emph{points de $X$ à valeurs dans $k$} (\textbf I, 3.4.4) et \emph{points de
$X$ rationnels sur $k$} (\textbf I, 3.4.5), en vertu de (\textbf I, 6.4.2).
Cela montre en particulier (en vertu de (\textbf I, 3.4.3.1)) que pour deux
espaces $k$-préalgébriques $X'$, $Y'$, l'ensemble sous-jacent au produit
$X'\times_kY'$ est identique à l'\emph{ensemble produit} $X'\times Y'$ des
ensembles sous-jacents (mais bien entendu la topologie de l'espace sous-jacent
à $X'\times_kY'$ n'est pas la \emph{topologie produit} des topologies des
espaces sous-jacents à $X'$ et $Y'$, elle est en général strictement plus fine
que cette dernière).

D'autre part, les anneaux locaux $\mathcal O_x$ aux points d'un espace
préalgébrique $X'$ sur $k$ sont des $k$-algèbres dont on vient de voir que le
corps résiduel est \emph{isomorphe à $k$}~; si $A$ et $B$ sont deux telles
$k$-algèbres locales, tout $k$-homomorphisme $\varphi:A\to B$ est
nécessairement \emph{local}~: en effet, si un élément $x$ de l'idéal maximal de
$A$ était tel que $\varphi(x)$ soit inversible, il existerait $\lambda\in k$
non nul tel que $\varphi(x-\lambda.1)$ appartienne à l'idéal maximal de $B$, ce
qui est absurde puisque $x-\lambda.1$ est inversible dans $A$. On conclut
aussitôt de là que si $X'=S(X)$, $Y'=S(Y)$ sont deux espaces préalgébriques sur
$k$, tout morphisme $X'\to Y'$ d'espaces $k$-annelés est aussi un morphisme
d'espaces $k$-annelés en anneaux locaux (\textbf I, 1.8.2~; cf.
$\mathrm{Err}_{\mathrm{II}}$). D'ailleurs, avec les notations précédentes, si
$A$ et $B$ sont des $k$-algèbres de type fini, $\varphi$ fait de $B$ une
$A$-algèbre de type fini~; donc tout morphisme $X'\to Y'$ de $k$-espaces
annelés est, en vertu de (10.9.6), de la forme $S(f)$, où $f:X\to Y$ est un
morphisme de $k$-préschémas.

Enfin, pour tout ouvert $U$ de $X'$, toute section
$s\in\Gamma(U,\mathcal O_{X'})$, et tout $x\in U$, $s(x)$
($\mathbf 0_{\mathrm I}$, 5.5.1) s'identifie à un élément de $k$, et l'on a
ainsi associé à $s$ une application $\overline s:x\mapsto s(x)$ de $U$ dans
$k$, autrement dit une section au-dessus de $U$ du faisceau $\mathcal A(X')$
des germes d'applications de $X'$ dans $k$~; comme l'application
$h_U:s\mapsto\overline s$ est évidemment un homomorphisme d'anneaux
$\Gamma(U,\mathcal O_{X'})\to\Gamma(U,\mathcal A(X'))$ et commute aux
restrictions à un ouvert $V\subset U$, les $h_U$ définissent un homomorphisme
de faisceaux d'anneaux $h:\mathcal O_{X'}\to\mathcal A(X')$. Si l'on prend
pour $U$ un ouvert ultra-affine $\operatorname{Spm}(A)$, où $A$ est un anneau
de Jacobson, dire que $\overline s=0$ signifie que pour tout idéal maximal
$\mathfrak m$ de $A$, $s$ appartient à $\mathfrak m$, ou encore que $s$ est
dans le \emph{radical} de $A$~; mais comme $A$ est un anneau de Jacobson, son
radical est égal à son nilradical~; pour que $h_U$ soit injective, il faut et
il suffit par suite que $A$ soit réduit.
\end{env}

\begin{env}[10.10.4]
\label{IV.10.10.4-fr}
On dit que l'espace $k$-préalgébrique $X'=S(X)$ est \emph{réduit} s'il en est
ainsi de $X$~; comme l'ensemble des points $x\in X$ où $X$ est réduit est
ouvert ($\mathbf 0_{\mathrm I}$, 5.2.2), son complémentaire contient au moins
un point fermé s'il est non vide (5.1.11), et il revient donc au même de dire
que $X'$ est réduit ou que chacun de ses anneaux locaux $\mathcal O_x$ (pour
$x\in X'$) est réduit. On vient de voir dans (10.10.3) que pour que
l'homomorphisme $h:\mathcal O_{X'}\to\mathcal A(X')$ soit \emph{injectif}, il
faut et il suffit que $X'$ soit réduit. Dans (FAC), Serre se borne en fait aux
espaces préalgébriques \emph{réduits}, ce qui lui permet de définir
$\mathcal O_{X'}$ comme un \emph{sous-faisceau} de $\mathcal A(X')$. On notera
que si $X'$ et $Y'$ sont des $k$-espaces préalgébriques réduits, il en est de
même de $X'\times_kY'$~: en effet, tout revient à voir que si $A$ et $B$ sont
deux $k$-algèbres de type fini réduites, il en est de même de $A\otimes_kB$~;
mais nous avons vu que les radicaux de $A$ et $B$ sont alors réduits à $0$, et
comme $k$ est algébriquement clos, $A$ et $B$ sont des algèbres «~séparables~»
sur $k$ au sens de Bourbaki (Bourbaki, \emph{Alg.}, chap.~VIII, \S~7,
n$^\circ$~5, prop.~5)~; donc $A\otimes_kB$ est sans radical (\emph{loc.
cit.}, n$^\circ$~6, cor.~3 du th.~3), et puisque c'est un anneau de Jacobson,
il est réduit. Toutefois, si $Z'$ est un troisième espace préalgébrique sur
$k$, le «~produit fibré~» $X'\times_{Z'}Y'$ de deux espaces préalgébriques
réduits sur $Z'$ n'est pas réduit en général, ce qui implique que la catégorie
de ces espaces est insuffisante dans de nombreuses questions (notamment dans
la théorie des groupes algébriques). Mais comme on l'a vu ci-dessus, on peut
garder le langage de Serre sans se limiter, comme ce dernier (qui en outre ne
considère que des espaces préalgébriques quasi-compacts), au cas des espaces
préalgébriques réduits.
\end{env}

\oldpage[IV-3]{116}

\begin{env}[10.10.5]
\label{IV.10.10.5-fr}
Enfin, on peut aussi considérer des ultra-préschémas sur un corps $k$
quelconque en conservant un langage qui reste proche de celui de Serre, et en
introduisant, comme chez Weil, une extension algébriquement close fixe $K$ de
$k$ (choisie assez grande, par exemple de degré de transcendance infini sur
$k$, pour avoir suffisamment de «~points génériques~» au sens de Weil). À tout
préschéma $X$ localement de type fini sur $k$, on associe alors l'ensemble
$S_K(X)=S(X\otimes_kK)$ des points de $X$ à valeurs dans $K$~; on a une
application canonique $j:S_K(X)\to X$ que l'on démontre être un
quasi-homéomorphisme quand on munit $S_K(X)$ de la topologie image réciproque
de celle de $X$ par $j$~; on munit $S_K(X)$ du faisceau de $k$-algèbres
$j^*(\mathcal O_X)$, et on obtient ainsi une sous-catégorie de la catégorie des
espaces $k$-annelés en anneaux locaux, que l'on pourrait appeler catégorie des
$(k,K)$-espaces préalgébriques. On peut montrer qu'on y peut encore définir des
produits et généraliser les résultats de (10.10.3) et (10.10.4)
($\mathcal A(X')$ étant ici remplacé par le faisceau $\mathcal A_K(X')$ des
germes d'applications de $X'$ dans $K$). Toutefois ce point de vue fait jouer
un rôle artificiel à un surcorps arbitrairement choisi de $k$, et nous ne le
signalons que pour le rejeter.
\end{env}

\oldpage[IV-3]{116}
\section{Propriétés topologiques des morphismes plats de présentation finie. Critères de platitude}
\label{section:IV.11-fr}

Alors que dans le \S~2 nous avons considéré les énoncés concernant la platitude
qui ne dépendent d'aucune hypothèse de finitude, et que le \S~6 étudie la notion
de platitude dans le cadre des préschémas localement noethériens (mais sans
hypothèse de finitude sur les morphismes), le présent paragraphe est consacré à
la notion de $f$-platitude dans le cas où le morphisme $f:X\to Y$ est
\emph{localement de présentation finie}. L'intérêt de la notion de morphisme
plat de présentation finie tient au fait que c'est elle qui semble exprimer
techniquement de la façon la plus adéquate la notion intuitive de «~famille de
préschémas algébriques paramétrée par un schéma $Y$~», dont l'étude est un des
objets principaux de la Géométrie algébrique. D'ailleurs, alors même qu'on ne
s'intéresserait au départ qu'au cas d'un schéma de base noethérien, il est
indispensable, pour certaines raisons techniques (par exemple pour certaines
applications de la théorie de la «~descente~», qui conduit à introduire des
schémas non nécessairement noethériens) de ne pas se limiter à ce cas, dès
qu'il s'agit de problèmes de nature essentiellement relative liés aux
morphismes localement de présentation finie. Nous suivrons systématiquement ce
principe, déjà étayé par les résultats des \S\S~8 et 9, dans toute la suite de ce
Chapitre, et même dans la suite de notre Traité, quitte à lui sacrifier à
l'occasion la simplicité de certaines démonstrations, que des hypothèses
noethériennes permettent parfois d'alléger\footnote{Ce principe s'inspire
également de la nécessité de donner droit de cité, comme «~espaces de
paramètres~» pour les familles de schémas algébriques, à des espaces annelés
(et même des «~topos~» annelés) quelconques, pour lesquels il ne peut plus être
question en général d'hypothèses noethériennes. Il semble assez clair que l'on
ne pourra plus éluder longtemps cette nouvelle extension de la Géométrie
algébrique, et il convient dès à présent de développer les notions et
techniques de nature «~relative~» de la théorie des schémas de sorte qu'elles
puissent s'adapter pratiquement telles quelles à ce cadre plus général.}. Dans
le présent paragraphe, cela nous conduit à reprendre, dans le contexte de la
«~présentation finie~» (notamment au n$^\circ$~3) certains énoncés de platitude,
déjà obtenus dans le contexte noethérien. L'outil technique essentiel pour
faire la réduction au cas noethérien est le théorème de compatibilité de la
platitude avec les limites projectives de préschémas (11.2.6), complétant les
résultats généraux du \S~8. Nous prouvons aussi en passant (11.3.1) un résultat
souvent utilisé

\oldpage[IV-3]{117}
par la suite, impliquant que l'ensemble des points de platitude d'un morphisme
localement de présentation finie est ouvert.

Dans les n$^\mathrm{os}$~4 à 8, nous étudions la question de la «~descente~» de
la platitude, consistant à trouver des conditions utiles sur un morphisme de
changement de base $Y'\to Y$ (non plat en général) pour pouvoir conclure que si
$X\times_YY'$ est plat sur $Y'$, $X$ est plat sur $Y$. Ces résultats, plus
techniques que ceux des n$^\mathrm{os}$~1 à 3, sont d'une utilisation moins
fréquente dans la suite~; ils joueront cependant un rôle important dans les
techniques de construction non projectives, au chapitre suivant. Le seul
résultat des n$^\mathrm{os}$~4 à 8 utilisé dans la suite du chap.~IV est le
critère valuatif de platitude (11.8), qu'on appliquera dans (15.2).

Enfin, les n$^\mathrm{os}$~9 et 10 sont consacrés à l'étude d'une notion qui
précise, en théorie des schémas, celle de densité au sens topologique, savoir la
notion de famille de sous-préschémas \emph{schématiquement dense} dans un
préschéma donné, et notamment l'étude du comportement de cette notion par
changement de base (plat ou quelconque). Cette notion est surtout utilisée,
pour l'instant, dans l'étude des schémas en groupes.

\subsection{Ensembles de platitude (cas noethérien)}
\label{subsection:IV.11.1-fr}

\begin{theorem}[11.1.1]
\label{IV.11.1.1-fr}
Soient $Y$ un préschéma localement noethérien, $f:X\to Y$ un morphisme
localement de type fini, $\mathcal F$ un $\mathcal O_X$-Module cohérent. Alors
l'ensemble $U$ des $x\in X$ tels que $\mathcal F$ soit $f$-plat au point $x$
est ouvert dans $X$.
\end{theorem}

La question étant évidemment locale sur $X$ et $Y$, on peut supposer
$Y=\operatorname{Spec}(A)$, $X=\operatorname{Spec}(B)$, $A$ étant noethérien et
$B$ une $A$-algèbre de type fini. On a alors $\mathcal F=\widetilde M$, où $M$
est un $B$-module de type fini. Appliquons le critère
($\mathbf 0_{\mathrm{III}}$, 9.2.6)~: il suffit donc de démontrer l'assertion
suivante~:

\begin{env}[11.1.1.1]
\label{IV.11.1.1.1-fr}
\emph{Soient $A$ un anneau noethérien, $B$ une $A$-algèbre de type fini, $M$ un
$B$-module de type fini, $\mathfrak q$ un idéal premier de $B$, $\mathfrak p$
son image réciproque dans $A$. On suppose que $M_{\mathfrak q}$ soit un
$A_{\mathfrak p}$-module plat. Alors il existe $g\in B-\mathfrak q$ tel que,
pour tout idéal premier $\mathfrak q'\supset\mathfrak q$ de $B$ tel que
$g\notin\mathfrak q'$, $M_{\mathfrak q'}$ soit un
$A_{\mathfrak p'}$-module plat, en désignant par $\mathfrak p'$ l'image
réciproque de $\mathfrak q'$ dans $A$ (il revient d'ailleurs au même
($\mathbf 0_{\mathrm I}$, 6.3.1) de dire que $M_{\mathfrak q'}$ est un
$A$-module plat).}
\end{env}

Considérons pour cela $B_{\mathfrak q'}$ comme une $A$-algèbre~; on a
évidemment
$\mathfrak pB_{\mathfrak q'}\subset\mathfrak q'B_{\mathfrak q'}$. On sait
alors ($\mathbf 0_{\mathrm{III}}$, 10.2.2) que pour que $M_{\mathfrak q'}$ soit
un $A$-module plat, il faut et il suffit que
$M_{\mathfrak q'}/\mathfrak pM_{\mathfrak q'}$ soit un
$(A/\mathfrak p)$-module plat et que l'on ait
$\operatorname{Tor}_1^A(M_{\mathfrak q'},A/\mathfrak p)=0$. Or, on a
$M_{\mathfrak q'}=M\otimes_BB_{\mathfrak q'}$~; $B_{\mathfrak q'}$ étant plat
sur $B$, on a
\[
  M_{\mathfrak q'}/\mathfrak pM_{\mathfrak q'}
  =(M/\mathfrak pM)_{\mathfrak q'}
\]
et
\[
  \operatorname{Tor}_1^A(M_{\mathfrak q'},A/\mathfrak p)
  =(\operatorname{Tor}_1^A(M,A/\mathfrak p))_{\mathfrak q'}
\]
(en définissant le Tor à l'aide d'une résolution projective de
$A/\mathfrak p$)~; pour la même raison, comme on doit avoir
$g\notin\mathfrak q'$, $B_{\mathfrak q'}$ est plat sur $B_g$, donc
\[
  (M/\mathfrak pM)_{\mathfrak q'}
  =((M/\mathfrak pM)_g)_{\mathfrak q'B_g}
\]
et
\[
  (\operatorname{Tor}_1^A(M,A/\mathfrak p))_{\mathfrak q'}
  =((\operatorname{Tor}_1^A(M,A/\mathfrak p))_g)_{\mathfrak q'B_g},
\]
où dans ces formules, $M/\mathfrak pM$ et
$\operatorname{Tor}_1^A(M,A/\mathfrak p)$ sont considérés comme des
$B$-modules. Compte tenu de ($\mathbf 0_{\mathrm I}$, 6.3.2), on voit qu'on
est ramené à prouver le

\begin{lemma}[11.1.1.2]
\label{IV.11.1.1.2-fr}
Sous les conditions de (11.1.1.1), il existe $g\in B-\mathfrak q$ tel que~:
\begin{enumerate}[label=(\roman*)]
  \item $(M/\mathfrak pM)_g$ soit un $(A/\mathfrak p)$-module plat~;
  \item $(\operatorname{Tor}_1^A(M,A/\mathfrak p))_g=0$.
\end{enumerate}
\end{lemma}

\oldpage[IV-3]{118}

En vertu du théorème de platitude générique (6.9.1) appliqué à l'anneau
intègre $A/\mathfrak p$, à la $(A/\mathfrak p)$-algèbre de type fini
$B/\mathfrak pB$ et au $(B/\mathfrak pB)$-module de type fini
$M/\mathfrak pM$, il existe un $h\in A-\mathfrak p$ tel que, si $\bar h$ est
son image canonique dans $A/\mathfrak p$,
$(M/\mathfrak pM)_{\bar h}$ soit un $(A/\mathfrak p)$-module plat. D'autre
part, comme $M_{\mathfrak q}$ est un $A_{\mathfrak p}$-module plat, et par
suite un $A$-module plat ($\mathbf 0_{\mathrm I}$, 6.3.1), on a
$\operatorname{Tor}_1^A(M_{\mathfrak q},A/\mathfrak p)=0$, ce qui s'écrit
aussi
$(\operatorname{Tor}_1^A(M,A/\mathfrak p))_{\mathfrak q}=0$. Mais comme $A$
et $B$ sont noethériens, $\operatorname{Tor}_1^A(M,A/\mathfrak p)$ est un
$B$-module de type fini, donc ($\mathbf 0_{\mathrm I}$, 5.2.2) il existe
$g\in B-\mathfrak q$ tel que
$(\operatorname{Tor}_1^A(M,A/\mathfrak p))_g=0$. D'ailleurs, on a
$(M/\mathfrak pM)_{\bar h}=(M/\mathfrak pM)_h$ ($M/\mathfrak pM$ étant
considéré dans le second membre comme un $A$-module)~; de plus,
$(M/\mathfrak pM)_h$ étant un $B$-module,
$(M/\mathfrak pM)_{hg}$ est encore un $(A/\mathfrak p)$-module plat, car il
s'écrit $(M/\mathfrak pM)_{\overline{hg}}$, où $\overline{hg}$ est l'image
canonique de $hg$ dans $B/\mathfrak pB$, et il suffit d'appliquer
($\mathbf 0_{\mathrm I}$, 6.3.2)~; enfin, on a
$(\operatorname{Tor}_1^A(M,A/\mathfrak p))_{hg}=0$ et
$hg\in B-\mathfrak q$ puisque $h\in A-\mathfrak p$. C.Q.F.D.

\begin{corollary}[11.1.2]
\label{IV.11.1.2-fr}
Soient $Y$ un préschéma localement noethérien, $f:X\to Y$ un morphisme
localement de type fini, $\mathcal F$, $\mathcal F'$ deux
$\mathcal O_X$-Modules cohérents,
$u:\mathcal F'\to\mathcal F$ un homomorphisme de $\mathcal O_X$-Modules. On
suppose que $\mathcal F$ est $f$-plat. Alors l'ensemble $U$ des $x\in X$ tels
que, en posant $y=f(x)$, l'homomorphisme
\[
  u_x\otimes1:\mathcal F'_x\otimes_{\mathcal O_y}k(y)\longrightarrow
  \mathcal F_x\otimes_{\mathcal O_y}k(y)
\]
soit injectif, est ouvert dans $X$.
\end{corollary}

En effet, soit $\mathcal N$ (resp. $\mathcal P$) le noyau (resp. le conoyau) de
$u$~; appliquons ($\mathbf 0_{\mathrm{III}}$, 10.2.4)\footnote{Voici une
démonstration de ($\mathbf 0_{\mathrm{III}}$, 10.2.4), qui n'a pas été publiée
dans l'\emph{Algèbre commutative} de N.~Bourbaki. Compte tenu de
($\mathbf 0_{\mathrm I}$, 6.1.2), il suffit de voir que b) entraîne a). Posons
$P=\operatorname{Im}(u)$, $Q=\operatorname{Coker}(u)$,
$R=\operatorname{Ker}(u)$. Le composé
$M\otimes k\to P\otimes k\to N\otimes k$ est injectif, et
$M\otimes k\to P\otimes k$ est surjectif, donc
$P\otimes k\to N\otimes k$ est injectif et
$M\otimes k\to P\otimes k$ est bijectif. La suite exacte
$0\to P\to N\to Q\to0$ donne la suite exacte
\[
0=\operatorname{Tor}_1^A(N,k)\to\operatorname{Tor}_1^A(Q,k)
\to P\otimes k\to N\otimes k\to Q\otimes k\to0
\]
et puisque $P\otimes k\to N\otimes k$ est injectif, on a
$\operatorname{Tor}_1^A(Q,k)=0$~; comme $Q$ est un $B$-module de type fini,
($\mathbf 0_{\mathrm{III}}$, 10.2.2) montre que $Q$ est un $A$-module plat~;
alors $P$ est aussi un $A$-module plat par ($\mathbf 0_{\mathrm I}$, 6.1.2).
La suite $0\to R\to M\to P\to0$ étant exacte, il en est de même de
$0\to R\otimes k\to M\otimes k\to P\otimes k\to0$ par
($\mathbf 0_{\mathrm I}$, 6.1.2)~; puisque $M\otimes k\to P\otimes k$ est
bijectif, on a $R\otimes k=0$~; mais comme $B$ est noethérien, $R$ est un
$B$-module de type fini, donc on a $R=0$ en vertu du lemme de Nakayama.} aux
anneaux locaux $\mathcal O_y$ et $\mathcal O_x$ et aux
$\mathcal O_x$-modules $\mathcal F'_x$ et $\mathcal F_x$~: dire que
$u_x\otimes1$ est injectif équivaut à dire que $\mathcal N_x=0$ et que
$\mathcal P$ est $f$-plat au point $x$. Or comme $\mathcal N$ et $\mathcal P$
sont cohérents ($\mathbf 0_{\mathrm I}$, 5.3.4), l'ensemble des $x$ où
$\mathcal N_x=0$ est ouvert ($\mathbf 0_{\mathrm I}$, 5.2.2), et l'ensemble
des $x$ où $\mathcal P$ est $f$-plat est ouvert par (11.1.1)~; d'où la
conclusion.

En particulier~:

\begin{corollary}[11.1.3]
\label{IV.11.1.3-fr}
Soient $Y$ un préschéma localement noethérien, $f:X\to Y$ un morphisme plat et
localement de type fini, $g$ une section de $\mathcal O_X$ au-dessus de $X$.
L'ensemble des $x\in X$ tels que $g_x\otimes1$ soit non diviseur de zéro dans
$\mathcal O_x\otimes_{\mathcal O_{f(x)}}k(f(x))$ est ouvert dans $X$.
\end{corollary}

Il suffit d'appliquer (11.1.2) à l'endomorphisme de $\mathcal O_X$ défini par
$g$.

\begin{corollary}[11.1.4]
\label{IV.11.1.4-fr}
Soient $Y$ un préschéma localement noethérien, $f:X\to Y$ un morphisme
localement de type fini, $\mathcal F$ un $\mathcal O_X$-Module cohérent et
$f$-plat. Soit $(g_i)_{1\leq i\leq n}$ une suite de sections de $\mathcal O_X$
au-dessus de $X$. Alors l'ensemble $U$ des $x\in X$ tels que la suite
$((g_i)_x\otimes1)$ soit
$(\mathcal F_x\otimes_{\mathcal O_{f(x)}}k(f(x)))$-régulière est ouvert dans
$X$.
\end{corollary}

\oldpage[IV-3]{119}

Comme $\mathcal F$ est $f$-plat, il résulte de ($\mathbf 0$, 15.1.16) que $U$
est aussi l'ensemble des $x\in X$ tels que la suite $((g_i)_x)$ soit
$\mathcal F_x$-régulière et que le $\mathcal O_x$-module $\mathcal G_x$ (où
$\mathcal G=\mathcal F/(\sum_{i=1}^n g_i\mathcal F)$) soit $f$-plat. Mais
$\mathcal F$ et $\mathcal G$ sont cohérents, donc le corollaire résulte de
(11.1.1) et de ($\mathbf 0$, 15.2.4).

\begin{corollary}[11.1.5]
\label{IV.11.1.5-fr}
Soient $X$, $Y$, $Z$ trois préschémas localement noethériens,
$f:X\to Y$, $g:Y\to Z$ deux morphismes de type fini, $\mathcal F$ un
$\mathcal O_X$-Module cohérent. Alors l'ensemble $U$ des $y\in Y$ tels que,
pour toute générisation $y'$ de $y$, $\mathcal F$ soit $(g\circ f)$-plat en
tous les points de $f^{-1}(y')$ (i.e. tels que
$\mathcal F\otimes_Y\operatorname{Spec}(\mathcal O_{Y,y})$ soit plat
relativement au morphisme
$X\times_Y\operatorname{Spec}(\mathcal O_{Y,y})\to Z$) est ouvert dans $Y$.
\end{corollary}

Si $V$ est l'ensemble des $x\in X$ où $\mathcal F$ est $(g\circ f)$-plat, $U$
est l'ensemble des $y\in Y$ tels que pour toute générisation $y'$ de $y$, on
ait $f^{-1}(y')\subset V$. Or $V$ est ouvert (11.1.1) donc localement
constructible dans $X$, et l'ensemble $W$ des $y\in Y$ tels que
$f^{-1}(y)\subset V$ est égal à $Y-f(X-V)$, donc est aussi localement
constructible dans $Y$ en vertu du théorème de Chevalley (1.8.4). Mais il
résulte alors de ($\mathbf 0_{\mathrm{III}}$, 9.2.5) que les points de $U$ sont
les points intérieurs à $W$, d'où la conclusion.

\begin{corollary}[11.1.6]
\label{IV.11.1.6-fr}
Soient $A$ un anneau noethérien, $B$ une $A$-algèbre de type fini, $C$ une
$B$-algèbre de type fini, $M$ un $C$-module de type fini. Alors l'ensemble des
$\mathfrak q\in\operatorname{Spec}(B)$ tels que $M_{\mathfrak q}$ soit un
$A$-module plat est ouvert dans $\operatorname{Spec}(B)$.
\end{corollary}

Compte tenu de (2.1.2), c'est une conséquence de (11.1.5) appliqué à
$Z=\operatorname{Spec}(A)$, $Y=\operatorname{Spec}(B)$,
$X=\operatorname{Spec}(C)$, $\mathcal F=\widetilde M$.

Les résultats de ce numéro seront débarrassés des hypothèses noethériennes
dans (11.3).

\subsection{Platitude d'une limite projective de préschémas}
\label{subsection:IV.11.2-fr}

\begin{env}[11.2.1]
\label{IV.11.2.1-fr}
Soient $A$ un anneau, $M$, $N$ deux $A$-modules, $A'$ une $A$-algèbre~;
posons $M'=M\otimes_AA'$, $N'=N\otimes_AA'$. Rappelons
($\mathbf{III}$, 6.3.8) que l'on définit, pour tout $i$, un homomorphisme
canonique de $A$-modules
\begin{equation}
\psi_i:\operatorname{Tor}_i^A(M,N)\to\operatorname{Tor}_i^{A'}(M',N')
\label{IV.11.2.1.1-fr}
\tag{11.2.1.1}
\end{equation}
de la façon suivante~: on considère une résolution gauche de $M$ par des
$A$-modules libres
\begin{equation}
\cdots\longrightarrow L_{i+1}\xrightarrow{f_i}L_i\longrightarrow\cdots
\longrightarrow L_0\xrightarrow{\varepsilon}M\longrightarrow0
\label{IV.11.2.1.2-fr}
\tag{11.2.1.2}
\end{equation}
d'où l'on déduit par tensorisation avec $A'$ un complexe de $A'$-modules
\begin{equation}
\cdots\longrightarrow L'_{i+1}\xrightarrow{f'_i}L'_i\longrightarrow\cdots
\longrightarrow L'_0\xrightarrow{\varepsilon'}M'\longrightarrow0
\label{IV.11.2.1.3-fr}
\tag{11.2.1.3}
\end{equation}
où l'on a posé $L'_i=L_i\otimes_AA'$, $f'_i=f_i\otimes1_{A'}$,
$\varepsilon'=\varepsilon\otimes1_{A'}$. Considérons d'autre part une
résolution gauche de $M'$ par des $A'$-modules libres
\begin{equation}
\cdots\longrightarrow L''_{i+1}\xrightarrow{f''_i}L''_i\longrightarrow\cdots
\longrightarrow L''_0\xrightarrow{\varepsilon''}M'\longrightarrow0.
\label{IV.11.2.1.4-fr}
\tag{11.2.1.4}
\end{equation}
\end{env}

\oldpage[IV-3]{120}

Comme les $L'_i$ sont des $A'$-modules libres, on sait (M, V, 1.1) qu'il y a
des $A'$-homomorphismes $u_i$ formant un diagramme commutatif
\[
\xymatrix{
\cdots\ar[r]&L'_{i+1}\ar[r]^{f'_i}\ar[d]_{u_{i+1}}&
L'_i\ar[r]\ar[d]_{u_i}&\cdots\ar[r]&
L'_0\ar[r]^{\varepsilon'}\ar[d]_{u_0}&M'\ar[r]\ar[d]_{1_{M'}}&0\\
\cdots\ar[r]&L''_{i+1}\ar[r]_{f''_i}&L''_i\ar[r]&\cdots\ar[r]&
L''_0\ar[r]_{\varepsilon''}&M'\ar[r]&0.
}
\]
Si l'on compose l'homomorphisme $u_\bullet:L'_\bullet\to L''_\bullet$ de
complexes ainsi défini avec l'homomorphisme canonique
$L_\bullet\to L'_\bullet$, on obtient un homomorphisme de complexes de
$A$-modules $L_\bullet\to L''_\bullet$~; en remarquant que l'on a
$(L_i\otimes_AN)\otimes_AA'=L'_i\otimes_{A'}N'$, on en déduit un
homomorphisme de complexes de $A$-modules
$L_\bullet\otimes_AN\to L''_\bullet\otimes_{A'}N'$, d'où, en passant à
l'homologie, les homomorphismes canoniques (11.2.1.1). Comme les $u_i$ sont
bien déterminés à homotopie près (M, V, 1.1), les homomorphismes (11.2.1.1)
ne dépendent pas du choix des $u_i$ ni du choix des résolutions libres
$L_\bullet$ et $L''_\bullet$.

\begin{sloppypar}
Comme les $\operatorname{Tor}_i^{A'}(M',N')$ sont des $A'$-modules, on déduit
canoniquement de (11.2.1.1) des $A'$-homomorphismes
\begin{equation}
\begin{gathered}
\varphi_i:\operatorname{Tor}_i^A(M,N)\otimes_AA'\\[-2pt]
\longrightarrow\operatorname{Tor}_i^{A'}(M',N').
\end{gathered}
\label{IV.11.2.1.5-fr}
\tag{11.2.1.5}
\end{equation}
\end{sloppypar}

Considérons maintenant deux homomorphismes d'anneaux
\[
\rho:A\to A^{(1)},\qquad \sigma:A^{(1)}\to A^{(2)},
\]
et leur composé $\sigma\circ\rho:A\to A^{(2)}$~; posons
$M^{(j)}=M\otimes_AA^{(j)}$, $N^{(j)}=N\otimes_AA^{(j)}$ pour $j=1,2$.
Alors l'homomorphisme canonique composé
\[
\operatorname{Tor}_i^A(M,N)\to
\operatorname{Tor}_i^{A^{(1)}}(M^{(1)},N^{(1)})\to
\operatorname{Tor}_i^{A^{(2)}}(M^{(2)},N^{(2)})
\]
est le même que l'homomorphisme canonique déduit de $\sigma\circ\rho$~;
cela résulte de ce que, si $L_\bullet^{(j)}$ est une résolution libre de
$M^{(j)}$, le diagramme
\[
\xymatrix{
L_\bullet\ar[r]&L_\bullet\otimes_AA^{(1)}\ar[r]\ar[d]&
L_\bullet^{(1)}\ar[d]\\
&L_\bullet\otimes_AA^{(2)}\ar[r]&
L_\bullet^{(1)}\otimes_{A^{(1)}}A^{(2)}\ar[r]&L_\bullet^{(2)}
}
\]
est commutatif.

\begin{env}[11.2.2]
\label{IV.11.2.2-fr}
Les notations étant celles de (11.2.1), considérons maintenant un
\emph{système inductif filtrant} de $A$-algèbres
$(A_\alpha,g_{\beta\alpha})$, et pour tout indice $\alpha$, posons
$M_\alpha=M\otimes_AA_\alpha$, $N_\alpha=N\otimes_AA_\alpha$~; il résulte
alors de (11.2.1) que pour chaque $i$,
$(\operatorname{Tor}_i^{A_\alpha}(M_\alpha,N_\alpha),
\psi_{i\beta\alpha})$, où $\psi_{i\beta\alpha}$ est l'homomorphisme
canonique (11.2.1.1) correspondant à $g_{\beta\alpha}$, est un \emph{système
inductif} de $A_\alpha$-modules. Posons
$A'=\varinjlim A_\alpha$, $M'=\varinjlim M_\alpha=M\otimes_AA'$,
$N'=\varinjlim N_\alpha=N\otimes_AA'$~; si on désigne par
$g_\alpha:A_\alpha\to A'$ l'homomorphisme canonique, on en déduit des
homomorphismes canoniques (11.2.1.1)
$\psi_{i\alpha}:\operatorname{Tor}_i^{A_\alpha}(M_\alpha,N_\alpha)
\to\operatorname{Tor}_i^{A'}(M',N')$ qui (toujours en vertu de (11.2.1))
forment un système inductif d'homomorphismes~; nous nous proposons de
compléter le résultat de (M, V, 9.4$^*$) en montrant que les
\begin{equation}
\psi_i=\varinjlim_\alpha\psi_{i\alpha}:
\varinjlim_\alpha\operatorname{Tor}_i^{A_\alpha}(M_\alpha,N_\alpha)
\to\operatorname{Tor}_i^{A'}(M',N')
\label{IV.11.2.2.1-fr}
\tag{11.2.2.1}
\end{equation}
sont des \emph{isomorphismes de $A'$-modules}. Pour cela, nous procédons
comme dans (M, V, 9.5$^*$), en associant à chaque $M_\alpha$ sa
\emph{résolution libre canonique}. Tout revient (compte tenu de l'exactitude
du foncteur $\varinjlim$) à prouver le
\end{env}

\oldpage[IV-3]{121}

\begin{lemma}[11.2.2.2]
\label{IV.11.2.2.2-fr}
Soient $(A_\alpha,g_{\beta\alpha})$ un système inductif filtrant d'anneaux,
$(M_\alpha,h_{\beta\alpha})$ un système inductif d'ensembles,
$A'=\varinjlim A_\alpha$, $M'=\varinjlim M_\alpha$,
$g_\alpha:A_\alpha\to A'$, $h_\alpha:M_\alpha\to M'$ les applications
canoniques. Pour tout $\alpha$, soit $F(M_\alpha)$ le $A_\alpha$-module des
combinaisons linéaires formelles d'éléments de $M_\alpha$~; soit de même
$F(M')$ le $A'$-module des combinaisons linéaires formelles d'éléments de
$M'$~; si $h^0_{\beta\alpha}:F(M_\alpha)\to F(M_\beta)$ (pour
$\alpha\leq\beta$) et $h'_\alpha:F(M_\alpha)\to F(M')$ sont les
$A_\alpha$-homomorphismes déduits de $h_{\beta\alpha}$ et $h_\alpha$
respectivement, $(F(M_\alpha),h^0_{\beta\alpha})$ est un système inductif
de $A_\alpha$-modules et $(h'_\alpha)$ un système inductif
d'homomorphismes. Alors
\[
h''=\varinjlim h'_\alpha:\varinjlim F(M_\alpha)\to
F(M')=F(\varinjlim M_\alpha)
\]
est un isomorphisme.
\end{lemma}

Pour la démonstration, voir Bourbaki, \emph{Alg.}, chap.~II, 3$^e$ éd.,
\S~6, n$^\circ$~6, cor. de la prop.~10.

\begin{env}[11.2.3]
\label{IV.11.2.3-fr}
Reprenons les notations de (11.2.1) et considérons particulièrement le cas
$i=1$~; posons $T=\operatorname{Tor}_1^A(M,N)$, $T'=T\otimes_AA'$,
$T''=\operatorname{Tor}_1^{A'}(M',N')$~; alors $\psi_1$ est
l'homomorphisme
\[
\operatorname{Ker}(f_0\otimes1_N)/\operatorname{Im}(f_1\otimes1_N)
\to
\operatorname{Ker}(f''_0\otimes1_{N'})/
\operatorname{Im}(f''_1\otimes1_{N'})
\]
qui se déduit par passage aux quotients de la restriction
$\operatorname{Ker}(f_0\otimes1_N)\to
\operatorname{Ker}(f''_0\otimes1_{N'})$ de
\[
u_1\otimes1:L_0\otimes_AN\to L''_0\otimes_{A'}N'=L''_0\otimes_AN.
\]
Posons $R=\operatorname{Ker}(\varepsilon)$,
$R''=\operatorname{Ker}(\varepsilon'')$, de sorte que l'on a les suites
exactes
\[
0\to R\xrightarrow{j}L_0\xrightarrow{\varepsilon}M\to0
\qquad\text{et}\qquad
0\to R''\xrightarrow{j''}L''_0\xrightarrow{\varepsilon''}M'\to0,
\]
d'où on déduit les suites exactes d'homologie
\begin{equation}
0=\operatorname{Tor}_1^A(L_0,N)\to T\xrightarrow{\partial}R\otimes_AN
\xrightarrow{j\otimes1}L_0\otimes_AN\xrightarrow{\varepsilon\otimes1}
M\otimes_AN\to0
\label{IV.11.2.3.1-fr}
\tag{11.2.3.1}
\end{equation}
\begin{equation}
0=\operatorname{Tor}_1^{A'}(L''_0,N')\to T''\xrightarrow{\partial''}
R''\otimes_{A'}N'\xrightarrow{j''\otimes1}L''_0\otimes_{A'}N'
\xrightarrow{\varepsilon''\otimes1}M'\otimes_{A'}N'\to0.
\label{IV.11.2.3.2-fr}
\tag{11.2.3.2}
\end{equation}
On a d'autre part un homomorphisme de $A'$-modules
\[
v:R'=\operatorname{Ker}(\varepsilon)\otimes_AA'
\to\operatorname{Ker}(\varepsilon')\xrightarrow{u_0}
\operatorname{Ker}(\varepsilon'')=R''.
\]
\end{env}

\oldpage[IV-3]{122}

Montrons que le diagramme
\begin{equation}
\xymatrix{
&T'\ar[r]^{\partial\otimes1}\ar[d]_{\varphi_1}&
R'\otimes_{A'}N'\ar[r]\ar[d]_{v\otimes1}&
L'_0\otimes_{A'}N'\ar[r]\ar[d]_{u_0\otimes1}&
M'\otimes_{A'}N'\ar[r]\ar@{=}[d]&0\\
0\ar[r]&T''\ar[r]_{\partial''}&R''\otimes_{A'}N'\ar[r]&
L''_0\otimes_{A'}N'\ar[r]&M'\otimes_{A'}N'\ar[r]&0
}
\label{IV.11.2.3.3-fr}
\tag{11.2.3.3}
\end{equation}
est \emph{commutatif}. Pour cela, on vérifie aussitôt (M, IV, 1) que
l'homomorphisme $\partial:T\to R\otimes_AN$ provient (dans le cas actuel)
par passage au quotient de l'homomorphisme
$w:\operatorname{Ker}(f_0\otimes1_N)\to R\otimes_AN$, restriction de
l'homomorphisme $g_0\otimes1_N:L_1\otimes_AN\to R\otimes_AN$, où
$g_0:L_1\to R$ est tel que $f_0=j\circ g_0$~; de même pour $\partial''$.
Il suffit donc de voir que le diagramme
\[
\xymatrix{
\operatorname{Ker}(f_0\otimes1_N)\ar[r]^w\ar[d]&R\otimes_AN\ar[r]&
R'\otimes_{A'}N'=(R\otimes_AN)\otimes_AA'\ar[d]_{v\otimes1}\\
\operatorname{Ker}(f''_0\otimes1_{N'})\ar[rr]&&R''\otimes_{A'}N'
}
\]
est commutatif, ce qui résulte de la commutativité du diagramme
\[
\xymatrix{
L_1\ar[r]\ar[d]&R\ar[d]\\
L''_1\ar[r]&R''.
}
\]

\begin{lemma}[11.2.4]
\label{IV.11.2.4-fr}
Soient $A$ un anneau, $C$ une $A$-algèbre, $M$ un $A$-module, $N$ un
$C$-module, $A'$ une $A$-algèbre. Posons $C'=C\otimes_AA'$,
$M'=M\otimes_AA'$, $N'=N\otimes_AA'=N\otimes_CC'$. Supposons que
$M\otimes_AN$ soit un $C$-module plat. Alors l'homomorphisme canonique
\[
\varphi_1:\operatorname{Tor}_1^A(M,N)\otimes_AA'
\to\operatorname{Tor}_1^{A'}(M',N')
\]
(cf. (11.2.1.5)) est surjectif.
\end{lemma}

Gardons les notations de (11.2.3)~; l'exactitude à droite du produit
tensoriel montre que la suite
$L'_1\xrightarrow{f'_0}L'_0\xrightarrow{\varepsilon'}M'\to0$ est exacte~;
comme $L'_0$ et $L'_1$ sont des $A'$-modules libres, \emph{on peut supposer
que l'on a pris} $L'_0=L''_0$, $L'_1=L''_1$, $u_0$ et $u_1$ étant les
\emph{applications identiques} et $f''_0=f'_0$. Comme
$R=\operatorname{Im}(f_0)$ et $R''=\operatorname{Im}(f''_0)=
\operatorname{Im}(f'_0)$, l'homomorphisme $v$ est surjectif, et il en est
donc de même de $v\otimes1$. La démonstration sera achevée si l'on prouve que
la première ligne de (11.2.3.3) est \emph{exacte}, $v\otimes1$ étant
surjective et $u_0\otimes1$ injectif (Bourbaki, \emph{Alg. comm.},
chap.~I$^{\mathrm{er}}$, \S~1, n$^\circ$~4, cor.~2 de la prop.~2).
Utilisons pour cela

\oldpage[IV-3]{123}

l'hypothèse que $M\otimes_AN$ est un $C$-module plat. Posant
$Q=\operatorname{Ker}(\varepsilon\otimes1)=\operatorname{Im}(j\otimes1)$
dans la suite exacte (11.2.3.1), on a les deux suites exactes
$0\to Q\to L_0\otimes_AN\to M\otimes_AN\to0$ et
$T\to R\otimes_AN\to Q\to0$, où les homomorphismes sont des homomorphismes
de $C$-modules~; utilisant l'hypothèse de platitude, on en déduit
($\mathbf 0_{\mathrm I}$, 6.1.2) la suite exacte
\[
0\to Q\otimes_CC'\to(L_0\otimes_AN)\otimes_CC'
\to(M\otimes_AN)\otimes_CC'\to0
\]
et d'autre part, le produit tensoriel étant exact à droite, on a la suite
exacte
\[
T\otimes_CC'\to(R\otimes_AN)\otimes_CC'\to Q\otimes_CC'\to0,
\]
d'où finalement la suite exacte
\[
T\otimes_CC'\to(R\otimes_AN)\otimes_CC'\to
(L_0\otimes_AN)\otimes_CC'\to(M\otimes_AN)\otimes_CC'\to0.
\]
Mais par définition, pour tout $C$-module $P$,
$P\otimes_CC'=P\otimes_AA'$, d'où la conclusion.

\begin{lemma}[11.2.5]
\label{IV.11.2.5-fr}
Soient $A$ un anneau, $\mathfrak J$ un idéal de $A$, $B$ une $A$-algèbre,
$M$ un $B$-module, $A\to A'$ un homomorphisme d'anneaux. On pose
$\mathfrak J'=\mathfrak JA'$, $B'=B\otimes_AA'$,
$M'=M\otimes_AA'=M\otimes_BB'$. Soit $\mathfrak p'$ un idéal premier de
$B'$ contenant $\mathfrak J'B'$. On suppose vérifiée l'une des hypothèses
suivantes~:
\begin{enumerate}[label=\emph{\alph*)}]
\item $\mathfrak J$ est nilpotent.
\item $A'$ est noethérien, $B'$ est une $A'$-algèbre de type fini, $M'$ un
$B'$-module de type fini.
\end{enumerate}
Dans ces conditions, supposons vérifiées les deux propriétés suivantes~:
\begin{enumerate}[label=\emph{(\roman*)}]
\item $M\otimes_A(A/\mathfrak J)$ est un $(A/\mathfrak J)$-module plat.
\item L'homomorphisme canonique composé
\[
\operatorname{Tor}_1^A(M,A/\mathfrak J)\xrightarrow{\psi_1}
\operatorname{Tor}_1^{A'}(M',A'/\mathfrak J')\xrightarrow{\theta}
(\operatorname{Tor}_1^{A'}(M',A'/\mathfrak J'))_{\mathfrak p'}
\]
(où $\psi_1$ est l'homomorphisme (11.2.1.1) et $\theta$ l'homomorphisme
canonique d'un $B'$-module dans son localisé en $\mathfrak p'$) est nul.
\end{enumerate}
Alors $M'_{\mathfrak p'}$ est un $A'$-module plat.
\end{lemma}

Notons que dans l'hypothèse b) $M'_{\mathfrak p'}$ est un
$B'_{\mathfrak p'}$-module de type fini, $B'_{\mathfrak p'}$ une
$A'$-algèbre noethérienne, et que $\mathfrak J'B'_{\mathfrak p'}$ est
contenu dans le radical $\mathfrak p'B'_{\mathfrak p'}$ de
$B'_{\mathfrak p'}$~; dans l'hypothèse a), $\mathfrak J'$ est nilpotent~;
on va donc pouvoir appliquer le critère de platitude
($\mathbf 0_{\mathrm{III}}$, 10.2.1) ou ($\mathbf 0_{\mathrm{III}}$, 10.2.2)
suivant que a) ou b) est vérifiée. En premier lieu, on a
\[
M'_{\mathfrak p'}\otimes_{A'}(A'/\mathfrak J')
=(M'\otimes_{A'}(A'/\mathfrak J'))_{\mathfrak p'}
=((M\otimes_A(A/\mathfrak J))\otimes_{A/\mathfrak J}
(A'/\mathfrak J'))_{\mathfrak p'}~;
\]
l'hypothèse (i) entraîne donc que
$M'_{\mathfrak p'}\otimes_{A'}(A'/\mathfrak J')$ est un
$(A'/\mathfrak J')$-module plat, compte tenu de
($\mathbf 0_{\mathrm I}$, 6.2.1 et 6.3.2). Reste donc à voir que
$\operatorname{Tor}_1^{A'}(M'_{\mathfrak p'},A'/\mathfrak J')=0$~; or ce
$B'_{\mathfrak p'}$-module est égal à
$(\operatorname{Tor}_1^{A'}(M',A'/\mathfrak J'))_{\mathfrak p'}$, en vertu
de la platitude de $B'_{\mathfrak p'}$ sur $B'$. Mais en vertu de
l'hypothèse (ii), l'homomorphisme composé
$\theta\circ\varphi_1:\operatorname{Tor}_1^A(M,A/\mathfrak J)\otimes_AA'
\to(\operatorname{Tor}_1^{A'}(M',A'/\mathfrak J'))_{\mathfrak p'}$ est
nul~; par ailleurs, (11.2.4) appliqué à $C=A/\mathfrak J$ et $N=C$ montre
(compte tenu de l'hypothèse (i)) que $\varphi_1$ est \emph{surjectif} (car
$A'/\mathfrak J'=C\otimes_AA'$)~; donc l'homomorphisme $\theta$ est
\emph{nul} et comme l'image par $\theta$ de
$\operatorname{Tor}_1^{A'}(M',A'/\mathfrak J')$ engendre le
$B'_{\mathfrak p'}$-module
$(\operatorname{Tor}_1^{A'}(M',A'/\mathfrak J'))_{\mathfrak p'}$, ce
dernier est nul. C.Q.F.D.

\begin{theorem}[11.2.6]
\label{IV.11.2.6-fr}
Les notations étant celles de (8.5.1) et (8.8.1), on suppose $S_\alpha$
quasi-compact, $X_\alpha$ et $Y_\alpha$ de présentation finie sur
$S_\alpha$~; soient $f_\alpha:X_\alpha\to Y_\alpha$ un
$S_\alpha$-morphisme, $\mathcal F_\alpha$ un
$\mathcal O_{X_\alpha}$-Module quasi-cohérent de présentation finie.

\oldpage[IV-3]{124}

\begin{enumerate}[label=\emph{(\roman*)}]
\item Soient $x$ un point de $X$, $x_\lambda$ sa projection canonique dans
$X_\lambda$. Pour que $\mathcal F$ soit $f$-plat au point $x$, il faut et il
suffit qu'il existe $\lambda\geq\alpha$ tel que $\mathcal F_\lambda$ soit
$f_\lambda$-plat au point $x_\lambda$.
\item Pour que $\mathcal F$ soit $f$-plat, il faut et il suffit qu'il existe
$\lambda\geq\alpha$ tel que $\mathcal F_\lambda$ soit $f_\lambda$-plat.
\end{enumerate}
\end{theorem}

On peut supposer que $S_\alpha=S_0$~; comme $Y_0$ est de présentation finie
sur $S_0$, $Y_0$ est quasi-compact, et $f_0:X_0\to Y_0$ est un morphisme de
présentation finie (1.6.2, (v)), donc on peut aussi se borner au cas où
$S_0=Y_0$. En outre, en vertu de la quasi-compacité de $S_0$ et du fait que
l'ensemble d'indices $L$ est filtrant, on peut se borner au cas où
$S_0=\operatorname{Spec}(A_0)$ est affine. De plus $X_0$ est quasi-compact,
donc le même raisonnement montre qu'on peut aussi supposer
$X_0=\operatorname{Spec}(B_0)$ affine~; on a alors
$\mathcal F_0=\widetilde M_0$, où $M_0$ est un $B_0$-module de présentation
finie, et l'énoncé (11.2.6) est dans ce cas équivalent au suivant (compte tenu
de ($\mathbf 0_{\mathrm I}$, 6.3.1))~:

\begin{corollary}[11.2.6.1]
\label{IV.11.2.6.1-fr}
Soient $A_0$ un anneau, $(A_\lambda)_{\lambda\in L}$ un système inductif
filtrant de $A_0$-algèbres, $B_0$ une $A_0$-algèbre de présentation finie,
$M_0$ un $B_0$-module de présentation finie. On pose
\[
\begin{gathered}
B_\lambda=B_0\otimes_{A_0}A_\lambda,\qquad
M_\lambda=M_0\otimes_{A_0}A_\lambda=M_0\otimes_{B_0}B_\lambda,\\
A=\varinjlim A_\lambda,\qquad
B=\varinjlim B_\lambda=B_0\otimes_{A_0}A,\\
M=\varinjlim M_\lambda=M_0\otimes_{A_0}A=M_0\otimes_{B_0}B.
\end{gathered}
\]
\begin{enumerate}[label=\emph{(\roman*)}]
\item Soit $\mathfrak p$ un idéal premier de $B$, et pour tout $\lambda$,
soit $\mathfrak p_\lambda$ son image réciproque dans $B_\lambda$. Pour que
$M_{\mathfrak p}$ soit un $A$-module plat, il faut et il suffit qu'il existe
$\lambda$ tel que $(M_\lambda)_{\mathfrak p_\lambda}$ soit un
$A_\lambda$-module plat.
\item Pour que $M$ soit un $A$-module plat, il faut et il suffit qu'il existe
$\lambda$ tel que $M_\lambda$ soit un $A_\lambda$-module plat.
\end{enumerate}
\end{corollary}

On doit seulement démontrer que les conditions sont nécessaires (2.1.4).
Nous procéderons en plusieurs étapes.

\medskip
\noindent\textbf{I)} \emph{Réduction au cas où les $A_\lambda$ sont
noethériens.} En vertu de (8.9.1), il existe un sous-anneau $A'_0$ de $A_0$,
qui est une $\mathbf Z$-algèbre de type fini, une $A'_0$-algèbre de type fini
$B'_0$ et un $B'_0$-module de type fini $M'_0$ tels que l'on ait
$B_0=B'_0\otimes_{A'_0}A_0$ et $M_0=M'_0\otimes_{A'_0}A_0$~; comme on a
$B_\lambda=B'_0\otimes_{A'_0}A_\lambda$ et
$M_\lambda=M'_0\otimes_{A'_0}A_\lambda$, on peut remplacer $A_0$, $B_0$,
$M_0$ par $A'_0$, $B'_0$, $M'_0$ dans l'énoncé de (11.2.6.1), en considérant
les $A_\lambda$ comme des $A'_0$-algèbres~; on peut donc déjà supposer que
$A_0$ est noethérien. Soit $H$ l'ensemble des couples
$(\lambda,C_\lambda)$, où $C_\lambda$ est une sous-$A_0$-algèbre de
\emph{type fini} de $A_\lambda$~; ordonnons $H$ en posant
$(\lambda,C_\lambda)\leq(\mu,D_\mu)$ si $\lambda\leq\mu$ et si
l'homomorphisme $\varphi_{\mu\lambda}:A_\lambda\to A_\mu$ est tel que
$\varphi_{\mu\lambda}(C_\lambda)\subset D_\mu$~; pour cet ordre $H$ est
filtrant croissant, car si $(\lambda,C_\lambda)$ et $(\mu,D_\mu)$ sont deux
éléments quelconques de $H$, on les majorera par un $(\nu,E_\nu)$ en prenant
$\nu\geq\lambda$, $\nu\geq\mu$ dans $L$, puis $E_\nu$ égal à la
sous-$A_0$-algèbre de $A_\nu$ engendrée par
$\varphi_{\nu\lambda}(C_\lambda)$ et $\varphi_{\nu\mu}(D_\mu)$. Pour un
élément $\xi=(\lambda,C_\lambda)$ de $H$, on posera $A_\xi=C_\lambda$, et
pour $\xi\leq\eta=(\mu,D_\mu)$ (donc $\lambda\leq\mu$ et
$\varphi_{\mu\lambda}(C_\lambda)\subset D_\mu$),
$\varphi_{\eta\xi}:A_\xi\to A_\eta$ sera la restriction à $C_\lambda$ de
$\varphi_{\mu\lambda}$, considéré comme homomorphisme dans $D_\mu$~; il est
clair que l'on obtient ainsi un système inductif filtrant de $A_0$-algèbres.
On pose $B_\xi=B_0\otimes_{A_0}A_\xi$,
$M_\xi=M_0\otimes_{A_0}A_\xi$~; cette fois les $A_\xi$ sont
noethériens~; en outre la formule de la double limite inductive (Bourbaki,
\emph{Alg.}, chap.~II, 3$^e$ éd., \S~6, n$^\circ$~4, prop.~7) prouve que
l'on a encore
$\varinjlim_HA_\xi=A$, $\varinjlim_HB_\xi=B$,
$\varinjlim_HM_\xi=M$. Supposons (11.2.6.1) prouvé pour le système inductif
$(A_\xi)_{\xi\in H}$~; soit $\mathfrak p$ un idéal premier de $B$,

\oldpage[IV-3]{125}

tel que $M_{\mathfrak p}$ soit un $A$-module plat~; il existe alors
$\xi\in H$ tel que, si $\mathfrak p_\xi$ est l'image réciproque de
$\mathfrak p$ dans $B_\xi$, $(M_\xi)_{\mathfrak p_\xi}$ soit un
$A_\xi$-module plat. Soit $\xi=(\lambda,C_\lambda)$, de sorte que
l'injection $A_\xi=C_\lambda\to A_\lambda$ donne un homomorphisme
$B_\xi=B_0\otimes_{A_0}C_\lambda\to B_\lambda=B_0\otimes_{A_0}A_\lambda$,
et si $\mathfrak p_\lambda$ est l'image réciproque de $\mathfrak p$ dans
$B_\lambda$, $\mathfrak p_\xi$ est l'image réciproque de
$\mathfrak p_\lambda$ dans $B_\xi$~; on a par suite
\[
(M_\lambda)_{\mathfrak p_\lambda}
=(M_\xi\otimes_{C_\lambda}A_\lambda)_{\mathfrak p_\lambda}
=((M_\xi)_{\mathfrak p_\xi}\otimes_{C_\lambda}A_\lambda)_{\mathfrak p_\lambda},
\]
donc $(M_\lambda)_{\mathfrak p_\lambda}$ est un $A_\lambda$-module plat
($\mathbf 0_{\mathrm I}$, 6.2.1 et 6.3.2). On traite de même le cas (ii) de
l'énoncé. Nous pouvons donc par la suite supposer les $A_\lambda$
noethériens pour $\lambda\in L$ (mais non nécessairement $A$ lui-même).

\medskip
\noindent\textbf{II)} \emph{Réduction de l'énoncé global (ii) à l'énoncé
ponctuel (i).} Supposons que $\mathcal F$ soit $f$-plat. Pour tout $\lambda$,
soit $U_\lambda$ l'ensemble des $x_\lambda\in X_\lambda$ tels que
$\mathcal F_\lambda$ soit $f_\lambda$-plat au point $x_\lambda$~; on sait
que $U_\lambda$ est \emph{ouvert} dans $X_\lambda$ puisque $S_\lambda$ est
noethérien et $f_\lambda$ de type fini (11.1.1)~; soit
$V_\lambda=v_\lambda^{-1}(U_\lambda)$ son image réciproque dans $X$. Comme,
par hypothèse, pour tout $x\in X$, il y a un $\lambda$ tel que
$\mathcal F_\lambda$ soit $f_\lambda$-plat au point $x_\lambda$, projection
de $x$ dans $X_\lambda$, cela signifie que $x\in V_\lambda$ pour un
$\lambda$~; autrement dit, $X$ est réunion des $V_\lambda$. D'ailleurs
(2.1.4), pour $\lambda\leq\mu$, on a $V_\lambda\subset V_\mu$, donc, comme
$X$ est quasi-compact, il existe un indice $\mu$ tel que $X=V_\mu$. Comme
les $X_\lambda$ sont quasi-compacts, il résulte de (8.3.4) qu'il existe un
indice $\nu\geq\mu$ tel que $v_{\mu\nu}^{-1}(U_\mu)=X_\nu$~; mais par
(2.1.4), cela entraîne que $\mathcal F_\nu$ est $f_\nu$-plat.

\medskip
\noindent\textbf{III)} \emph{Fin de la démonstration.} Il reste à prouver
(i) en supposant $S_0$ affine et noethérien~; si $y_0$ est la projection de
$x$ dans $S_0$, on peut en outre, en vertu de (2.1.4) et de
($\mathbf I$, 3.6.5), remplacer $S_0$ par
$\operatorname{Spec}(\mathcal O_{y_0})$, autrement dit on peut se borner au
cas où $A_0$ est un anneau \emph{local} noethérien, d'idéal maximal
$\mathfrak m_0=\mathfrak j_{y_0}$~; par définition, $\mathfrak m_0$ est
l'image réciproque de l'idéal premier $\mathfrak p=\mathfrak j_x$ de $B$,
et $M_{\mathfrak p}$ est supposé être un $A$-module plat~; on a donc en
particulier $\operatorname{Tor}_1^A(A/\mathfrak m_0A,M_{\mathfrak p})=0$,
ce qui s'écrit aussi, puisque les $\operatorname{Tor}_1^A(A/\mathfrak m_0A,M)$
sont des $B$-modules et que $B_{\mathfrak p}$ est plat sur $B$,
$(\operatorname{Tor}_1^A(A/\mathfrak m_0A,M))_{\mathfrak p}=0$.
Notons maintenant que $\operatorname{Tor}_1^{A_0}(A_0/\mathfrak m_0,M_0)$
est un $B_0$-module \emph{de type fini}, car on peut le définir en prenant
une résolution de $A_0/\mathfrak m_0$ par des $A_0$-modules libres
\emph{de type fini} ($A_0$ étant noethérien) et en tensorisant par $M_0$, ce
qui donne des $B_0$-modules de type fini~; comme $B_0$ est noethérien,
l'homologie du complexe ainsi obtenu est bien formée de $B_0$-modules de
type fini. Soit $(t_i^0)_{1\leq i\leq m}$ un système de générateurs du
$B_0$-module $\operatorname{Tor}_1^{A_0}(A_0/\mathfrak m_0,M_0)$ et soit
$t_i$ l'image canonique (11.2.1.1) de $t_i^0$ dans
$\operatorname{Tor}_1^A(A/\mathfrak m_0A,M)$. L'hypothèse entraîne qu'il
existe un $h\in B-\mathfrak p$ tel que $ht_i=0$ pour $1\leq i\leq m$. Or,
on a (11.2.2.1)
\[
\operatorname{Tor}_1^A(A/\mathfrak m_0A,M)=
\varinjlim\operatorname{Tor}_1^{A_\lambda}
(A_\lambda/\mathfrak m_0A_\lambda,M_\lambda)~;
\]
il existe donc un $\lambda$ tel que, si les $t_i^\lambda$ sont les images
des $t_i^0$ dans
$\operatorname{Tor}_1^{A_\lambda}(A_\lambda/\mathfrak m_0A_\lambda,
M_\lambda)$, il existe $h_\lambda\in B_\lambda$ d'image $h$, tel que
$h_\lambda t_i^\lambda=0$ pour $1\leq i\leq m$. Soit
$\mathfrak p_\lambda=\mathfrak j_{x_\lambda}$ l'idéal premier de
$B_\lambda$ image réciproque de $\mathfrak p$~; on a
$h_\lambda\in B_\lambda-\mathfrak p_\lambda$, donc les images canoniques
des $t_i^0$ dans
$(\operatorname{Tor}_1^{A_\lambda}(A_\lambda/\mathfrak m_0A_\lambda,
M_\lambda))_{\mathfrak p_\lambda}$ sont nulles et par suite
l'homomorphisme
$\operatorname{Tor}_1^{A_0}(A_0/\mathfrak m_0,M_0)\to
(\operatorname{Tor}_1^{A_\lambda}(A_\lambda/\mathfrak m_0A_\lambda,
M_\lambda))_{\mathfrak p_\lambda}$ est nul. Les conditions du lemme
(11.2.5) sont donc remplies ($A_0/\mathfrak m_0$ étant un corps), et
$(M_\lambda)_{\mathfrak p_\lambda}$ est un $A_\lambda$-module plat, ce qui
achève de démontrer le théorème.

\begin{corollary}[11.2.7]
\label{IV.11.2.7-fr}
Soient $S=\operatorname{Spec}(A)$ un schéma affine, $f:X\to S$ un morphisme,
$\mathcal F$ un $\mathcal O_X$-Module quasi-cohérent, $x$ un point de $X$.
Les conditions suivantes sont équivalentes~:
\begin{enumerate}[label=\emph{\alph*)}]
\item $f$ est un morphisme de présentation finie, $\mathcal F$ est un
$\mathcal O_X$-Module de présentation finie et $\mathcal F$ est $f$-plat au
point $x$ (resp. $f$-plat).

\oldpage[IV-3]{126}

\item Il existe un schéma affine noethérien $S_0=\operatorname{Spec}(A_0)$,
un morphisme de type fini $f_0:X_0\to S_0$, un
$\mathcal O_{X_0}$-Module cohérent $\mathcal F_0$, un morphisme $S\to S_0$
tels que le $S$-préschéma $X_0\otimes_{S_0}S$ soit $S$-isomorphe à $X$ et
que, si on identifie $X$ à $X_0\otimes_{S_0}S$, $\mathcal F$ soit isomorphe
à $\mathcal F_0\otimes_{\mathcal O_{S_0}}\mathcal O_S$, et $\mathcal F_0$
soit $f_0$-plat au point $x_0$ projection de $x$ dans $X_0$ (resp.
$f_0$-plat).
\item Les conditions de b) sont vérifiées et en outre $A_0$ est une
sous-$\mathbf Z$-algèbre de type fini de $A$, le morphisme $S\to S_0$
correspondant à l'injection canonique $A_0\to A$.
\end{enumerate}
\end{corollary}

Il est clair que c) implique b) et b) implique a) en vertu de (2.1.4).
D'autre part, on peut considérer $A$ comme limite inductive de ses
sous-$\mathbf Z$-algèbres de type fini, et on sait par (8.9.1) qu'il y a une
telle sous-algèbre $A_0$ et un morphisme de type fini $f_0:X_0\to S_0$ tels
que $X$ soit $S$-isomorphe à $X_0\times_{S_0}S$ et $\mathcal F$ isomorphe à
$\mathcal F_0\otimes_{\mathcal O_{S_0}}\mathcal O_S$~; on peut donc écrire
$X=\varprojlim X_\lambda$, où $X_\lambda=X_0\otimes_{A_0}A_\lambda$,
$A_\lambda$ parcourant l'ensemble des sous-$\mathbf Z$-algèbres de type
fini de $A$ contenant $A_0$, et $\mathcal F=v_\lambda^*(\mathcal F_\lambda)$,
où $v_\lambda:X\to X_\lambda$ est la projection canonique et
$\mathcal F_\lambda=\mathcal F_0\otimes_{A_0}A_\lambda$. Il résulte de
(11.2.6) qu'il existe $\lambda$ tel que $\mathcal F_\lambda$ soit
$f_\lambda$-plat au point $x_\lambda=v_\lambda(x)$ (resp.
$f_\lambda$-plat)~; alors le sous-anneau $A_\lambda$ de $A$ vérifie les
conditions de c).

\begin{proposition}[11.2.8]
\label{IV.11.2.8-fr}
Soient $g:X\to S$, $h:Y\to S$ deux morphismes de présentation finie,
$f:X\to Y$ un $S$-morphisme, $\mathcal F$ un
$\mathcal O_X$-Module quasi-cohérent de présentation finie~; pour tout
$s\in S$, soient $X_s=g^{-1}(s)=X\times_S\operatorname{Spec}(k(s))$,
$Y_s=h^{-1}(s)=Y\times_S\operatorname{Spec}(k(s))$, $f_s$ le morphisme
$f\times1:X_s\to Y_s$, $\mathcal F_s=\mathcal F\otimes_{\mathcal O_S}k(s)$.
Alors l'ensemble $E$ des $s\in S$ tels que $\mathcal F_s$ soit $f_s$-plat
est localement constructible.
\end{proposition}

La propriété dont nous voulons démontrer qu'elle est constructible vérifie la
condition (9.2.1, (i)), en vertu de (2.2.13) et (2.5.1). Compte tenu de
(9.2.3), on peut donc se limiter au cas où $S$ est affine, noethérien et
intègre de point générique $\eta$ et à prouver que $E$ ou $S-E$ est un
voisinage de $\eta$ dans $S$. Si $\eta\in S-E$, cela résulte aussitôt du
lemme (9.4.7.1). Si au contraire $\eta\in E$, il résulte tout d'abord de
(11.2.6), appliqué suivant la méthode de (8.1.2, a)), qu'il y a un voisinage
ouvert $U$ de $\eta$ dans $S$ tel que
$\mathcal F|g^{-1}(U)$ soit $h^{-1}(U)$-plat~; \emph{a fortiori} (2.1.4)
$\mathcal F_s$ est $f_s$-plat pour tout $s\in U$, ce qui achève la
démonstration.

Le théorème suivant généralise (11.2.6.1, (ii))~:

\begin{theorem}[11.2.9]
\label{IV.11.2.9-fr}
(Raynaud). --- Soient $A_0$ un anneau, $(A_\lambda)_{\lambda\in L}$ un système
inductif filtrant de $A_0$-algèbres, $B_0$ une $A_0$-algèbre de présentation
finie, $\mathfrak J_0$ un idéal de type fini de $B_0$, $M_0$ un $B_0$-module
de présentation finie. On pose
$B_\lambda=B_0\otimes_{A_0}A_\lambda$,
$\mathfrak J_\lambda=\mathfrak J_0B_\lambda$,
$M_\lambda=M_0\otimes_{A_0}A_\lambda=M_0\otimes_{B_0}B_\lambda$,
$A=\varinjlim A_\lambda$,
$B=\varinjlim B_\lambda=B_0\otimes_{A_0}A$,
$\mathfrak J=\mathfrak J_0B$,
$M=\varinjlim M_\lambda=M_0\otimes_{A_0}A=M_0\otimes_{B_0}B$. Pour que
$\operatorname{gr}_{\mathfrak J}^*(M)$ soit un $A$-module plat, il faut et il
suffit qu'il existe $\lambda$ tel que
$\operatorname{gr}_{\mathfrak J_\lambda}^*(M_\lambda)$ soit un
$A_\lambda$-module plat~; les homomorphismes canoniques
\begin{equation}
\left\{
\begin{aligned}
\operatorname{gr}_{\mathfrak J_\lambda}^*(M_\lambda)
  \otimes_{A_\lambda}A&\longrightarrow
  \operatorname{gr}_{\mathfrak J}^*(M),\\
\operatorname{gr}_{\mathfrak J_\lambda}^*(M_\lambda)
  \otimes_{A_\lambda}A_\mu&\longrightarrow
  \operatorname{gr}_{\mathfrak J_\mu}^*(M_\mu)
\end{aligned}
\right.
\qquad\text{pour }\lambda\leq\mu
\label{IV.11.2.9.1-fr}
\tag{11.2.9.1}
\end{equation}
sont alors bijectifs et
$\operatorname{gr}_{\mathfrak J_\mu}^*(M_\mu)$ est un $A_\lambda$-module plat
pour $\lambda\leq\mu$.
\end{theorem}

\begin{proof}
Montrons d'abord que les conditions sont suffisantes, ce qui revient à prouver
la bijectivité de (11.2.9.1). Cela résulte du lemme suivant~:

\begin{lemma}[11.2.9.2]
\label{IV.11.2.9.2-fr}
Soient $A$ un anneau, $B$ une $A$-algèbre, $M$ un $B$-module, $\mathfrak J$ un
idéal de $B$. Soient $A'$ une $A$-algèbre~; posons
$B'=B\otimes_AA'$, $M'=M\otimes_AA'=M\otimes_BB'$,
$\mathfrak J'=\mathfrak JB'$. Si
$\operatorname{gr}_{\mathfrak J}^*(M)$ est un $A$-module plat,
l'homomorphisme canonique
\[
(\operatorname{gr}_{\mathfrak J}^*(M))\otimes_AA'
\longrightarrow \operatorname{gr}_{\mathfrak J'}^*(M')
\]
est bijectif.
\end{lemma}

\oldpage[IV-3]{127}

En effet, par récurrence sur $k$, l'hypothèse que les
$\mathfrak J^kM/\mathfrak J^{k+1}M$ sont des $A$-modules plats pour $k\geq0$
entraîne d'abord, par ($\mathbf 0_{\mathrm I}$, 6.1.2), que
$M/\mathfrak J^{k+1}M$ est un $A$-module plat~; en outre
($\mathbf 0_{\mathrm I}$, 6.1.2), la suite
\[
0\longrightarrow(\mathfrak J^{k+1}M)\otimes_AA'
\longrightarrow M\otimes_AA'
\longrightarrow(M/\mathfrak J^{k+1}M)\otimes_AA'\longrightarrow0
\]
est alors exacte, autrement dit $(\mathfrak J^{k+1}M)\otimes_AA'$ s'identifie
à son image canonique $\mathfrak J'^{,k+1}M'$ dans $M'=M\otimes_AA'$.
D'autre part, toujours en vertu de ($\mathbf 0_{\mathrm I}$, 6.1.2), la suite
\[
0\longrightarrow(\mathfrak J^{k+1}M)\otimes_AA'
\longrightarrow(\mathfrak J^kM)\otimes_AA'
\longrightarrow(\mathfrak J^kM/\mathfrak J^{k+1}M)\otimes_AA'
\longrightarrow0
\]
est exacte, ce qui prouve le lemme.

Pour démontrer la nécessité des conditions de (11.2.9), nous procéderons en
plusieurs étapes.

\medskip
\noindent\textbf{I)} \emph{Réduction au cas où les $A_\lambda$ sont
noethériens.} --- On procède comme dans la réduction I) de (11.2.6.1), dont
nous gardons les notations~; il faut simplement commencer par remplacer
$A'_0$ par une sous-$A'_0$-algèbre de type fini $A''_0$ de $A_0$ telle que,
si l'on pose $B''_0=B'_0\otimes_{A'_0}A''_0$, il y ait dans $B''_0$ un idéal
de type fini $\mathfrak J''_0$ tel que $\mathfrak J''_0B_0=\mathfrak J_0$.
On considère pour cela les sous-$A'_0$-algèbres $A'_\alpha$ de type fini de
$A_0$, qui forment une famille filtrante, et l'on a
$B_0=\varinjlim B'_\alpha$, où
$B'_\alpha=B'_0\otimes_{A'_0}A'_\alpha$~; il y a donc un indice $\beta$ tel
qu'un système fini de générateurs de $\mathfrak J_0$ soit formé d'images dans
$B_0$ d'éléments de $B'_\beta$~; on prendra alors $A''_0=A'_\beta$,
$B''_0=B'_\beta$ et pour $\mathfrak J''_0$ l'idéal engendré par ces éléments.
On peut donc supposer que $A_0$ (donc aussi $B_0$) est noethérien. On définit
ensuite comme \emph{loc. cit.} l'ensemble filtrant $H$ et les $A_\xi$,
$B_\xi$, $M_\xi$ pour $\xi\in H$~; on posera aussi
$\mathfrak J_\xi=\mathfrak J_0B_\xi$ pour tout $\xi\in H$. Supposons alors
que l'on ait démontré qu'il existe un $\xi=(\lambda,C_\lambda)$ tel que
$\operatorname{gr}_{\mathfrak J_\xi}^*(M_\xi)$ soit un $C_\lambda$-module
plat~; comme $M_\lambda=M_\xi\otimes_{C_\lambda}A_\lambda$, il résulte de
(11.2.9.2) que $\operatorname{gr}_{\mathfrak J_\lambda}^*(M_\lambda)$ est un
$A_\lambda$-module plat.

\medskip
\noindent\textbf{II)} \emph{Lemmes préliminaires.}

\begin{lemma}[11.2.9.3]
\label{IV.11.2.9.3-fr}
Soient $A$ un anneau,
$B=\bigoplus_{i\geq0}B_i$ une $A$-algèbre graduée à degrés positifs,
$M=\bigoplus_{i\in\mathbf Z}M_i$ un $B$-module gradué. On suppose que $B_0$
est un anneau local, et que chacun des $M_i$ est un $B_0$-module de type fini.
Pour que $M$ soit un $B$-module de type fini, il faut et il suffit que, si
$\mathfrak q$ est l'image réciproque dans $A$ de l'idéal maximal $\mathfrak m$
de $B_0$, $M\otimes_Ak(\mathfrak q)$ soit un
$(B\otimes_Ak(\mathfrak q))$-module de type fini.
\end{lemma}

La condition étant évidemment nécessaire, prouvons qu'elle est suffisante.
Comme $B_0$ (et par suite $B$) est une $A_{\mathfrak q}$-algèbre, on peut
remplacer $A$ par $A_{\mathfrak q}$, autrement dit supposer que $A$ est aussi
un anneau local dont $\mathfrak q$ est l'idéal maximal. Par hypothèse, il existe
un entier $N$ tel que l'homomorphisme canonique
\[
\left(\bigoplus_{-N\leq i\leq N}(M_i\otimes_{B_0}B)\right)
\otimes_Ak(\mathfrak q)\longrightarrow M\otimes_Ak(\mathfrak q)
\]

\oldpage[IV-3]{128}

soit surjectif~; cela signifie aussi que, pour tout entier $n$,
l'homomorphisme canonique de $k(\mathfrak q)$-modules
\[
\left(\bigoplus_{-N\leq i\leq N}(M_i\otimes_{B_0}B_{n-i})\right)
\otimes_Ak(\mathfrak q)\longrightarrow M_n\otimes_Ak(\mathfrak q)
\]
est surjectif. Or, $M_n$ est un $B_0$-module de type fini, $B_0$ un anneau
local et $\mathfrak qB_0\subset\mathfrak m$~; donc le lemme de Nakayama prouve
que chacun des homomorphismes canoniques
\[
\bigoplus_{-N\leq i\leq N}(M_i\otimes_{B_0}B_{n-i})\longrightarrow M_n
\]
est surjectif, d'où la conclusion.

\begin{corollary}[11.2.9.4]
\label{IV.11.2.9.4-fr}
Sous les hypothèses de (11.2.9.3), on suppose en outre que chacun des $B_i$ et
chacun des $M_i$ soit un $B_0$-module de présentation finie, et que $M$ soit
un $A$-module plat. Pour que $M$ soit un $B$-module de présentation finie, il
faut et il suffit que $M\otimes_Ak(\mathfrak q)$ soit un
$(B\otimes_Ak(\mathfrak q))$-module de présentation finie.
\end{corollary}

En vertu de (11.2.9.3), si la condition de l'énoncé est vérifiée, $M$ est un
$B$-module de type fini, donc il existe un $B$-module libre \emph{gradué} de
type fini $L$ (ayant donc une base finie formée d'éléments homogènes) et un
homomorphisme gradué surjectif de degré 0, $u:L\to M$. Soit $R=\ker(u)$, qui
est un $B$-module gradué~; il y a donc un nombre fini d'entiers $m_j$
($1\leq j\leq r$) tels que pour chaque entier $i$, $R_i$ soit le noyau d'un
homomorphisme surjectif
\[
\bigoplus_{1\leq j\leq r}B_{i+m_j}\longrightarrow M_i~;
\]
on conclut alors de l'hypothèse sur les $M_i$ et les $B_i$ que $R_i$ est un
$B_0$-module de type fini (Bourbaki, \emph{Alg. comm.}, chap.~I, \S~2,
n$^\circ$~8, lemme 9). Pour prouver que $R$ est un $B$-module de type fini,
notons qu'en vertu de l'hypothèse de platitude et de
($\mathbf 0_{\mathrm I}$, 6.1.2) la suite
\[
0\longrightarrow R\otimes_Ak(\mathfrak q)
\longrightarrow L\otimes_Ak(\mathfrak q)
\longrightarrow M\otimes_Ak(\mathfrak q)\longrightarrow0
\]
est exacte, et l'hypothèse entraîne donc (Bourbaki, \emph{loc. cit.}) que
$R\otimes_Ak(\mathfrak q)$ est un $(B\otimes_Ak(\mathfrak q))$-module de type
fini~; il suffit donc d'appliquer à $R$ le lemme (11.2.9.3).

\medskip
\noindent\textbf{III)} \emph{Réduction au cas où les homomorphismes de
transitions $\varphi_{\mu\lambda}:A_\lambda\to A_\mu$ (pour
$\lambda\leq\mu$) sont injectifs.} --- Soit $A'_\lambda$ l'image de
$A_\lambda$ par l'homomorphisme canonique $\varphi_\lambda:A_\lambda\to A$~;
il est clair que les $A'_\lambda$ forment un système inductif de sous-anneaux
noethériens de $A$, dont la limite inductive est $A$, et les homomorphismes de
transition $A'_\lambda\to A'_\mu$ (pour $\lambda\leq\mu$) sont injectifs.
Posons $B'_\lambda=B_0\otimes_{A_0}A'_\lambda$,
$\mathfrak J'_\lambda=\mathfrak J_0B'_\lambda$,
$M'_\lambda=M_0\otimes_{A_0}A'_\lambda=M_\lambda\otimes_{A_\lambda}A'_\lambda$~;
on a encore $B=\varinjlim B'_\lambda$, $M=\varinjlim M'_\lambda$. Supposons
que l'on ait prouvé qu'il existe un $\lambda$ tel que, pour $\mu\geq\lambda$,
$\operatorname{gr}_{\mathfrak J'_\mu}^*(M'_\mu)$ soit un $A'_\mu$-module
plat. Soit $\mathfrak a_\lambda$ le noyau de l'homomorphisme $\varphi_\lambda$,
qui est donc un idéal de type fini de l'anneau noethérien $A_\lambda$. Il
résulte de la définition de la limite inductive qu'il existe un indice
$\mu\geq\lambda$ tel que $\varphi_{\mu\lambda}(\mathfrak a_\lambda)=0$,
autrement dit l'homomorphisme $\varphi_{\mu\lambda}$ se factorise en
$\varphi_{\mu\lambda}:A_\lambda\to A'_\lambda\to A_\mu$, et l'on peut écrire
$B_\mu=B'_\lambda\otimes_{A'_\lambda}A_\mu$,
$\mathfrak J_\mu=\mathfrak J'_\lambda B_\mu$, et
$M_\mu=M'_\lambda\otimes_{A'_\lambda}A_\mu$~; on déduit donc de (11.2.9.2)
que $\operatorname{gr}_{\mathfrak J_\mu}^*(M_\mu)$ est un $A_\mu$-module
plat.

\medskip
\noindent\textbf{IV)} \emph{Réduction au cas où
$B_0=A_0[T_1,\ldots,T_n]$ et
$\operatorname{gr}_{\mathfrak J_0}^*(B_0)=B_0$.} --- En vertu de l'hypothèse,
il y a un système $(t_i)_{1\leq i\leq n}$ de générateurs de la $A_0$-algèbre
de type fini $B_0$, tel que les $t_i$ pour $i\leq m$ engendrent l'idéal
$\mathfrak J_0$ de $B_0$. Posons $B'_0=A_0[T_1,\ldots,T_n]$,

\oldpage[IV-3]{129}

algèbre de polynômes (donc noethérienne), et soit $\mathfrak J'_0$ l'idéal de
$B'_0$ engendré par les $T_i$ d'indice $i\leq m$~; on a alors un
$A_0$-homomorphisme surjectif $u_0:B'_0\to B_0$ transformant chaque $T_i$ en
$t_i$ ($1\leq i\leq n$), donc tel que $u_0(\mathfrak J'_0)=\mathfrak J_0$~;
cet homomorphisme permet de considérer $M_0$ comme un $B'_0$-module de type
fini. On pose alors
\[
B'_\lambda=B'_0\otimes_{A_0}A_\lambda
=A_\lambda[T_1,\ldots,T_n],\qquad
\mathfrak J'_\lambda=\mathfrak J'_0B'_\lambda,
\]
\[
B'=B'_0\otimes_{A_0}A=A[T_1,\ldots,T_n],\qquad
\mathfrak J'=\mathfrak J'_0B',
\]
de sorte que $\mathfrak J'$ est l'idéal de $B'$ engendré par les $T_i$
d'indice $i\leq m$~; en outre, il est clair que pour tout entier $k\geq0$,
on a $\mathfrak J'^kM=\mathfrak J^kM$ et
$\mathfrak J'_\lambda{}^kM_\lambda=\mathfrak J_\lambda^kM_\lambda$ pour tout
$\lambda$~; donc
$\operatorname{gr}_{\mathfrak J'}^*(M)=\operatorname{gr}_{\mathfrak J}^*(M)$
en tant que $A$-module et
$\operatorname{gr}_{\mathfrak J'_\lambda}^*(M_\lambda)
=\operatorname{gr}_{\mathfrak J_\lambda}^*(M_\lambda)$ en tant que
$A_\lambda$-module~; on peut donc substituer $B'$, $B'_\lambda$,
$\mathfrak J'$, $\mathfrak J'_\lambda$ à $B$, $B_\lambda$, $\mathfrak J$,
$\mathfrak J_\lambda$ respectivement dans la démonstration~; on notera en
outre que par construction $\operatorname{gr}_{\mathfrak J'}^*(B')$
s'identifie à $B'$ et $\operatorname{gr}_{\mathfrak J'_\lambda}^*(B'_\lambda)$
à $B'_\lambda$.

\medskip
\noindent\textbf{V)} \emph{Preuve du fait que
$\operatorname{gr}_{\mathfrak J}^*(M)$ est un
$\operatorname{gr}_{\mathfrak J}^*(B)$-module de présentation finie.} --- Nous
supposons donc désormais $A_0$ et les $A_\lambda$ noethériens, les
homomorphismes de transition $A_\lambda\to A_\mu$ injectifs,
$B_0=A_0[T_1,\ldots,T_n]$, $\mathfrak J_0$ étant engendré par les $T_i$
d'indice $i\leq m$~; l'anneau
$C_0=\operatorname{gr}_{\mathfrak J_0}^*(B_0)$ s'identifie donc à $B_0$ et
$C=\operatorname{gr}_{\mathfrak J}^*(B)$ s'identifie à $B$~; nous utiliserons
seulement tout d'abord dans ce qui suit le fait que
$C\cong C_0\otimes_{A_0}A$.

Nous aurons besoin de la variante suivante de (6.9.3)~:

\begin{lemma}[11.2.9.5]
\label{IV.11.2.9.5-fr}
Soient $A_0$ un anneau noethérien, $B_0$ une $A_0$-algèbre de type fini,
$\mathfrak J_0$ un idéal de $B_0$, $M_0$ un $B_0$-module de type fini. Il
existe alors une suite $(S_{0i})_{1\leq i\leq m}$ de sous-schémas de
$S_0=\operatorname{Spec}(A_0)$ ayant les propriétés suivantes~:
\begin{enumerate}[label=\arabic{enumi}$^\circ$]
\item Les espaces sous-jacents aux $S_{0i}$ sont deux à deux disjoints et
forment un recouvrement de $S_0$.
\item Pour chaque $i$, l'ensemble $S_{0i}$ est ouvert dans
$\bigcup_{j\geq i}S_{0j}$.
\item Chaque schéma $S_{0i}$ est affine.
\item Si $A_{0i}$ est l'anneau de $S_{0i}$ et si l'on pose
$B_{0i}=B_0\otimes_{A_0}A_{0i}$, $\mathfrak J_{0i}=\mathfrak J_0B_{0i}$,
alors
$\operatorname{gr}_{\mathfrak J_{0i}}^*(M_0\otimes_{A_0}A_{0i})$ est un
$A_{0i}$-module plat.
\end{enumerate}
\end{lemma}

On procède comme dans (6.9.3) par récurrence noethérienne, en supposant le
lemme vrai lorsqu'on y remplace $A_0$ par $A'_0=A_0/\mathfrak a_0$, où
l'idéal $\mathfrak a_0$ est tel que $V(\mathfrak a_0)\neq S_0$, $B_0$ par
$B'_0=B_0\otimes_{A_0}A'_0$, $\mathfrak J_0$ par
$\mathfrak J'_0=\mathfrak J_0B'_0$ et $M_0$ par $M_0\otimes_{A_0}A'_0$. On
considère le nilradical $\mathfrak N_0$ de $A_0$, et il suffit évidemment de
prouver le lemme en remplaçant $A_0$ par $A_0/\mathfrak N_0$ et $B_0$,
$\mathfrak J_0$, $M_0$ par les objets correspondants comme ci-dessus.
Autrement dit, on peut supposer que $A_0$ est \emph{réduit}~; d'autre part,
$\operatorname{gr}_{\mathfrak J_0}^*(B_0)$ est une $A_0$-algèbre de type fini
et $\operatorname{gr}_{\mathfrak J_0}^*(M_0)$ un
$\operatorname{gr}_{\mathfrak J_0}^*(B_0)$-module de type fini~; en vertu du
théorème de platitude générique (6.9.1), il existe un ouvert
$D(t_0)\subset S_0$ tel que si l'on pose $A_{01}=(A_0)_{t_0}$,
$\operatorname{gr}_{\mathfrak J_0}^*(M_0)\otimes_{A_0}A_{01}$ soit un
$A_{01}$-module plat~; mais puisque $A_{01}$ est un $A_0$-module plat,
$\operatorname{gr}_{\mathfrak J_0}^*(M_0)\otimes_{A_0}A_{01}$ s'identifie à
$\operatorname{gr}_{\mathfrak J_{01}}^*(M_0\otimes_{A_0}A_{01})$, où
$B_{01}=B_0\otimes_{A_0}A_{01}$ et $\mathfrak J_{01}=\mathfrak J_0B_{01}$.
Le complémentaire de $D(t_0)$ dans $S_0$ est alors de la forme
$V(\mathfrak a_0)$, et on conclut en appliquant l'hypothèse de récurrence
noethérienne.

Appliquons ce lemme à la situation actuelle, en gardant les mêmes notations~;
posons $U_{0i}=\bigcup_{j\leq i}S_{0j}$, qui est un ouvert quasi-compact de
$S_0$. Il y a donc une famille
$(f_{ik})_{1\leq i\leq m,\,1\leq k\leq n}$ d'éléments de $A_0$ tels que pour
chaque $i$, $U_{0i}$ soit réunion des $D(f_{ik})$ ($1\leq k\leq n$).

\oldpage[IV-3]{130}

Le $C_0$-module $\operatorname{gr}_{\mathfrak J_0}^*(M_0)$ est engendré par
un nombre fini d'éléments homogènes de degré 1, de sorte qu'il y a un
épimorphisme de $C_0$-modules gradués
\[
u_0:C_0^r\longrightarrow\operatorname{gr}_{\mathfrak J_0}^*(M_0).
\]
Comme $C=C_0\otimes_{A_0}A$ et que l'homomorphisme canonique
$\operatorname{gr}_{\mathfrak J_0}^*(M_0)\otimes_{A_0}A
\to\operatorname{gr}_{\mathfrak J}^*(M)$ est surjectif, on déduit donc de
$u_0$ un épimorphisme de $C$-modules gradués
\[
u:C^r\xrightarrow{u_0\otimes1_A}
\operatorname{gr}_{\mathfrak J_0}^*(M_0)\otimes_{A_0}A
\longrightarrow\operatorname{gr}_{\mathfrak J}^*(M),
\]
et tout revient à voir que le $C$-module gradué $Q=\ker(u)$ est \emph{de type
fini}.

\begin{lemma}[11.2.9.6]
\label{IV.11.2.9.6-fr}
Sous les hypothèses précédentes, soit $\mathfrak R_0$ un idéal de $A_0$~;
posons
\[
A'_0=A_0/\mathfrak R_0,\qquad
B'_0=B_0\otimes_{A_0}A'_0=B_0/\mathfrak R_0B_0,\qquad
\mathfrak J'_0=\mathfrak J_0B'_0,
\]
et $\mathfrak R=\mathfrak R_0A$, $A'=A/\mathfrak R$, et supposons en outre
que $\operatorname{gr}_{\mathfrak J'_0}^*(M_0\otimes_{A_0}A'_0)$ soit un
$A'_0$-module plat. Alors il existe un nombre fini d'éléments $q_h$
($1\leq h\leq s$) de $Q$ tels que, pour tout idéal premier
$\mathfrak p\supset\mathfrak RB$ de $B$, les images canoniques des $q_h$ dans
$Q_\mathfrak p$ engendrent le $C_\mathfrak p$-module $Q_\mathfrak p$.
\end{lemma}

Supposons ce lemme démontré, et notons qu'on peut encore l'appliquer en
remplaçant $A_0$ par $(A_0)_{t_0}$, $B_0$, $\mathfrak J_0$, $M_0$,
$A_\lambda$ et $A$ par les objets correspondants $(B_0)_{t_0}$,
$(\mathfrak J_0)_{t_0}$, $(M_0)_{t_0}$, $(A_\lambda)_{t_0}$ et $A_{t_0}$ pour
tout $t_0\in A_0$, puisque $(A_0)_{t_0}$ est un $A_0$-module plat. On
appliquera alors ce lemme en remplaçant successivement $A_0$ par chacun des
anneaux $(A_0)_{f_{ik}}$, $\mathfrak R_0$ étant remplacé par l'idéal
$\mathfrak R_{ik}$ définissant le sous-schéma fermé de $D(f_{ik})$ induit par
$S_{0i}$ sur l'ouvert $S_{0i}\cap D(f_{ik})$ de $S_{0i}$. Comme l'hypothèse
de platitude de (11.2.9.6) est vérifiée dans chacun de ces cas en raison de
(11.2.9.5), on obtient ainsi pour chaque couple $(i,k)$ une famille finie
$(a_{ikh})_{1\leq h\leq l}$ d'éléments de $Q_{f_{ik}}$ dont les images dans
$Q_{\mathfrak p_{ik}}$ engendrent ce $C_{\mathfrak p_{ik}}$-module, pour tout
idéal premier $\mathfrak p_{ik}$ de $B_{f_{ik}}$ contenant
$\mathfrak R_{ik}B_{f_{ik}}$. On peut écrire, pour un entier $d>0$ convenable,
$a_{ikh}=b_{ikh}/f_{ik}^d$ pour tous les indices $i$, $k$, $h$, avec
$b_{ikh}\in Q$. Puisque les $S_{0i}$ recouvrent $S_0$, tout idéal premier
$\mathfrak p$ de $B$ est tel que son image dans $S_0$ appartienne à un
ensemble $S_{0i}\cap D(f_{ik})$, donc l'image $\mathfrak p_{ik}$ de
$\mathfrak p$ dans $B_{f_{ik}}$ contient $\mathfrak R_{ik}B_{f_{ik}}$. Comme
les images des $b_{ikh}$ ($1\leq h\leq l$) dans
$Q_{\mathfrak p_{ik}}=Q_\mathfrak p$ engendrent ce $C_\mathfrak p$-module,
on en conclut que les $b_{ikh}$
($1\leq i\leq m$, $1\leq k\leq n$, $1\leq h\leq l$) engendrent le
$C$-module $Q$ (Bourbaki, \emph{Alg. comm.}, chap.~II, \S~3, n$^\circ$~3,
th.~1).

Reste à démontrer le lemme (11.2.9.6). Posons
\[
C'=C\otimes_AA'=C/\mathfrak RC,\qquad
B'=B\otimes_AA',\qquad \mathfrak J'=\mathfrak JB',\qquad
M'=M\otimes_AA'.
\]
Par hypothèse $\operatorname{gr}_{\mathfrak J}^*(M)$ est un $A$-module plat,
donc $\operatorname{gr}_{\mathfrak J'}^*(M')$ s'identifie à
$\operatorname{gr}_{\mathfrak J}^*(M)\otimes_AA'$ (11.2.9.2), et on a la
suite exacte ($\mathbf 0_{\mathrm I}$, 6.1.2)
\begin{equation}
0\longrightarrow Q\otimes_AA'\longrightarrow C'^{,r}
\xrightarrow{u\otimes1_{A'}}\operatorname{gr}_{\mathfrak J'}^*(M')
\longrightarrow0.
\label{IV.11.2.9.7-fr}
\tag{11.2.9.7}
\end{equation}
Posons $C'_0=C_0\otimes_{A_0}A'_0=C_0/\mathfrak R_0C_0$ et considérons
l'épimorphisme de $C'_0$-modules gradués
\[
u'_0:C'_0{}^{,r}\xrightarrow{u_0\otimes1_{A'_0}}
\operatorname{gr}_{\mathfrak J_0}^*(M_0)\otimes_{A_0}A'_0
\longrightarrow
\operatorname{gr}_{\mathfrak J'_0}^*(M_0\otimes_{A_0}A'_0).
\]
Soit $Q_0=\ker(u'_0)$~; utilisant le fait que
$\operatorname{gr}_{\mathfrak J'_0}^*(M_0\otimes_{A_0}A'_0)$ est un
$A'_0$-module plat, on voit que son produit tensoriel sur $A'_0$ par $A'$
s'identifie à $\operatorname{gr}_{\mathfrak J'}^*(M')$ (11.2.9.2) et on a la
suite exacte ($\mathbf 0_{\mathrm I}$, 6.1.2)
\begin{equation}
0\longrightarrow Q_0\otimes_{A'_0}A'\longrightarrow C'^{,r}
\xrightarrow{u\otimes1_{A'}}\operatorname{gr}_{\mathfrak J'}^*(M')
\longrightarrow0.
\label{IV.11.2.9.8-fr}
\tag{11.2.9.8}
\end{equation}

\oldpage[IV-3]{131}

d'où, par comparaison avec (11.2.9.7), un isomorphisme de $C'$-modules
gradués
\[
Q\otimes_AA'\xrightarrow{\sim}Q_0\otimes_{A'_0}A'.
\]
Comme $C'_0$ est noethérien, $Q_0$ est un $C'_0$-module de type fini, donc
$Q_0\otimes_{A'_0}A'$ est un $C'$-module gradué de type fini, et il en est par
suite de même de $Q\otimes_AA'$~; soient $q_h$ ($1\leq h\leq s$) des éléments
homogènes de $Q$ dont les images dans $Q\otimes_AA'$ engendrent ce
$C'$-module.

Considérons maintenant un idéal premier quelconque
$\mathfrak p\supset\mathfrak RB$ de $B$~; on a
$C_\mathfrak p=\operatorname{gr}_{\mathfrak J_\mathfrak p}^*(B_\mathfrak p)$
par platitude, et
$\operatorname{gr}_{\mathfrak J_\mathfrak p}^0(B_\mathfrak p)
=B_\mathfrak p/\mathfrak J_\mathfrak p$ est réduit à 0 ou est un anneau
local~; pour prouver le lemme, on peut évidemment se borner au second cas. Il
est clair alors que chacun des $(B_\mathfrak p/\mathfrak J_\mathfrak p)$-modules
$\operatorname{gr}_{\mathfrak J_\mathfrak p}^i(B_\mathfrak p)$ est de
présentation finie. D'autre part, pour tout indice $i$,
$M/\mathfrak J^iM$ s'identifie à
$(M_0/\mathfrak J_0^iM_0)\otimes_{A_0}A$, et est donc un $B$-module de
présentation finie. Par récurrence sur $i$, l'hypothèse que
$\operatorname{gr}_{\mathfrak J}^*(M)$ est un $A$-module plat entraîne, en
vertu de ($\mathbf 0_{\mathrm I}$, 6.1.2), que $M/\mathfrak J^iM$ est un
$A$-module plat. L'application de (11.2.6.1) où l'on remplace $M_0$ par
$M_0/\mathfrak J_0^kM_0$ pour $k\leq i$ montre qu'il existe un indice
$\lambda_i$ tel que, pour $\mu\geq\lambda_i$, chacun des
$M_\mu/\mathfrak J_\mu^{k+1}M_\mu$ soit un $A_\mu$-module plat, et par suite
aussi ($\mathbf 0_{\mathrm I}$, 6.1.2) chacun des
$\operatorname{gr}_{\mathfrak J_\mu}^k(M_\mu)$ est un $A_\mu$-module plat
pour $k\leq i$. On déduit par suite de (11.2.9.2) que chacun des
$(B/\mathfrak J)$-modules $\operatorname{gr}_{\mathfrak J}^i(M)$ est de
présentation finie, donc
$\operatorname{gr}_{\mathfrak J_\mathfrak p}^i(M_\mathfrak p)$ est un
$(B_\mathfrak p/\mathfrak J_\mathfrak p)$-module de présentation finie. Par
ailleurs les images des $q_h$ dans
\[
Q_\mathfrak p\otimes_{B_\mathfrak p}k(\mathfrak p)
=(Q\otimes_AA')_\mathfrak p\otimes_{A'}k(\mathfrak p)
\]
engendrent ce module sur
$C_\mathfrak p\otimes_{B_\mathfrak p}k(\mathfrak p)$, puisque
\[
C_\mathfrak p\otimes_{B_\mathfrak p}k(\mathfrak p)
=C'_\mathfrak p\otimes_{A'}k(\mathfrak p)~;
\]
comme on a la suite exacte
\[
Q_\mathfrak p\otimes_{B_\mathfrak p}k(\mathfrak p)
\longrightarrow C_\mathfrak p^{,r}\otimes_{B_\mathfrak p}k(\mathfrak p)
\longrightarrow
\operatorname{gr}_{\mathfrak J_\mathfrak p}^*(M_\mathfrak p)
\otimes_{B_\mathfrak p}k(\mathfrak p)\longrightarrow0,
\]
on en conclut que
$\operatorname{gr}_{\mathfrak J_\mathfrak p}^*(M_\mathfrak p)
\otimes_{B_\mathfrak p}k(\mathfrak p)$ est un
$(C_\mathfrak p\otimes_{B_\mathfrak p}k(\mathfrak p))$-module de présentation
finie. Appliquant le lemme (11.2.9.4) où $A$, $B$, $M$ sont remplacés
respectivement par $B_\mathfrak p$, $C_\mathfrak p$ et
$\operatorname{gr}_{\mathfrak J_\mathfrak p}^*(M_\mathfrak p)$, on en conclut
que $Q_\mathfrak p$ est un $C_\mathfrak p$-module de type fini. Utilisant
maintenant le lemme de Nakayama (et le fait que $C=B$) on en déduit que les
images des $q_h$ dans $Q_\mathfrak p$ engendrent ce $C_\mathfrak p$-module.
Ceci achève de prouver le lemme (11.2.9.6) et le fait que
$\operatorname{gr}_{\mathfrak J}^*(M)$ est un $C$-module de présentation
finie.

\medskip
\noindent\textbf{VI)} \emph{Fin de la démonstration.} --- Posons pour abréger
$C_\lambda=C_0\otimes_{A_0}A_\lambda$ (égal en fait à $B_\lambda$),
$N_\lambda=\operatorname{gr}_{\mathfrak J_\lambda}^*(M_\lambda)$ et
$N=\operatorname{gr}_{\mathfrak J}^*(M)$. Notons d'abord que pour chaque
entier $k$, $\mathfrak J^kM$ s'identifie à
$\varinjlim\mathfrak J_\lambda^kM_\lambda$ par exactitude du foncteur
$\varinjlim$~; l'image de $\mathfrak J_\lambda^kM_\lambda$ dans $M_\mu$ pour
$\lambda\leq\mu$ (resp. dans $M$) engendrant $\mathfrak J_\mu^kM_\mu$
(resp. $\mathfrak J^kM$)~; utilisant encore l'exactitude de $\varinjlim$, on
en conclut que $N$ s'identifie canoniquement à $\varinjlim N_\lambda$. Faisant
cette identification, nous allons d'abord prouver que~:
\begin{equation}
\text{\emph{Pour $\lambda$ assez grand, l'homomorphisme canonique
$v_\lambda:N_\lambda\otimes_{A_\lambda}A\to N$ est bijectif.}}
\label{IV.11.2.9.9-fr}
\tag{11.2.9.9}
\end{equation}
Comme les $C_\lambda$ sont noethériens et $N_\lambda$ un $C_\lambda$-module
de type fini, les $C$-modules $N_\lambda\otimes_{A_\lambda}A$ sont de
présentation finie et forment un système inductif filtrant, dont la limite
inductive s'identifie canoniquement à $N$ en vertu du fait que $\varinjlim$
commute aux produits tensoriels. En outre, les homomorphismes de transition
$v_{\mu\lambda}:N_\lambda\otimes_{A_\lambda}A
\to N_\mu\otimes_{A_\mu}A$

\oldpage[IV-3]{132}

pour $\lambda\leq\mu$ et les homomorphismes $v_\lambda$ sont \emph{surjectifs}.
Pour un $\lambda$ fixé, considérons les sous-$C$-modules
$\ker(v_{\mu\lambda})=N'_\mu$ de $N_\lambda\otimes_{A_\lambda}A$~; par
définition de la limite inductive, ils forment une famille filtrante croissante
de sous-$C$-modules de $\ker(v_\lambda)$, dont la réunion est
$\ker(v_\lambda)$~; mais on a vu dans V) que $N$ est un $C$-module de
présentation finie, donc (Bourbaki, \emph{Alg. comm.}, chap.~I, \S~2,
n$^\circ$~8, lemme 9) $\ker(v_\lambda)$ est un $C$-module de type fini~; il
existe par suite un indice $\mu\geq\lambda$ tel que
$N'_\mu=\ker(v_\lambda)$, ce qui signifie (vu le fait que $v_{\mu\lambda}$ est
surjectif) que $v_\mu$ est \emph{injectif}~; comme il est surjectif, cela
prouve (11.2.9.9).

Quitte à remplacer $A_0$ et $M_0$ par $A_\lambda$ et $M_\lambda$ pour
$\lambda$ assez grand, on peut donc supposer que l'homomorphisme canonique
$v_\lambda$ est bijectif pour \emph{tout} $\lambda$. Posons alors
$P_\lambda=N_0\otimes_{C_0}C_\lambda=N_0\otimes_{A_0}A_\lambda$, de sorte que
$N=\varinjlim P_\lambda$. Comme $C_0$ est une $A_0$-algèbre de présentation
finie et $P_0$ un $C_0$-module de présentation finie, on peut appliquer à
$C$ et $N$ le résultat de (11.2.6.1) et on voit donc qu'il existe un indice
$\lambda$ tel que, pour $\mu\geq\lambda$, $P_\mu$ soit un $A_\mu$-module
plat. Or, pour $\mu\geq\lambda$, on a un diagramme commutatif
\[
\xymatrix{
P_\mu \ar[r] \ar[d]^{w_\mu} & N \ar@{=}[d] \\
N_\mu \ar[r] & N
}
\]
où il résulte des définitions que $w_\mu$ est \emph{surjectif}. Comme, en
vertu de III), les homomorphismes $A_\mu\to A$ sont injectifs et que $P_\mu$
est un $A_\mu$-module plat, l'homomorphisme canonique
$P_\mu\to N=P_\mu\otimes_{A_\mu}A$ est \emph{injectif}~; on conclut donc du
diagramme commutatif précédent que $w_\mu$ est aussi injectif, donc bijectif,
et par suite $N_\mu$ est un $A_\mu$-module plat pour $\mu\geq\lambda$, ce qui
achève de prouver (11.2.9).
\end{proof}

\begin{remark}[11.2.10]
\label{IV.11.2.10-fr}
Nous ignorons si la généralisation de (11.2.6, (i)) analogue au th. de Raynaud,
est valable.
\end{remark}

\subsection{Application à l'élimination d'hypothèses noethériennes}
\label{subsection:IV.11.3-fr}
\label{IV.11.3-fr}

\begin{theorem}[11.3.1]
\label{IV.11.3.1-fr}
Soient $f:X\to Y$ un morphisme localement de présentation finie,
$\mathcal F$ un $\mathcal O_X$-Module quasi-cohérent de présentation finie.
Alors l'ensemble $U$ des points $x\in X$ tels que $\mathcal F$ soit $f$-plat
au point $x$ est ouvert dans $X$. Si de plus $\operatorname{Supp}(\mathcal F)=X$,
$f|U$ est un morphisme ouvert.
\end{theorem}

\begin{proof}
La seconde assertion a déjà été démontrée (2.4.6) et n'a été insérée que pour
mémoire.

La question étant locale sur $X$ et $Y$, on peut supposer que $X$ et $Y$ sont
affines, et que $f$ est un morphisme de présentation finie. Soit $x\in X$ un
point tel que $\mathcal F$ soit $f$-plat au point $x$. Appliquant (11.2.7), on
peut supposer que
\[
X=X_0\times_{Y_0}Y,\qquad f=f_0\times1,\qquad
\mathcal F=\mathcal F_0\otimes_{\mathcal O_{Y_0}}\mathcal O_Y,
\]
où $Y_0$ est noethérien, $f_0:X_0\to Y_0$ un morphisme de type fini,
$\mathcal F_0$ un $\mathcal O_{X_0}$-Module cohérent~; en outre, si $x_0$ est la
projection canonique de $x$ dans $X_0$, on peut supposer que $\mathcal F_0$ est
$f_0$-plat au point $x_0$. Alors, en vertu de (11.1.1), l'ensemble $U_0$ des
points de $X_0$ où $\mathcal F_0$ est $f_0$-plat est un voisinage de $x_0$~;
donc $\mathcal F$ est $f$-plat aux points de l'image réciproque de $U_0$ dans
$X$ (2.1.4), et cela prouve que $U$ est un voisinage de $x$.
\end{proof}

\oldpage[IV-3]{133}

\begin{corollary}[11.3.2]
\label{IV.11.3.2-fr}
Soient $f:X\to Y$, $g:Y\to Z$ deux morphismes de présentation finie,
$\mathcal F$ un $\mathcal O_X$-Module quasi-cohérent de présentation finie.
Alors l'ensemble $U$ des $y\in Y$ tels que, pour toute générisation $y'$ de
$y$, $\mathcal F$ soit $(g\circ f)$-plat en tous les points de $f^{-1}(y')$
(i.e. tels que $\mathcal F\otimes_Y\operatorname{Spec}(\mathcal O_{Y,y})$
soit plat relativement au morphisme
$X\times_Y\operatorname{Spec}(\mathcal O_{Y,y})\to Z$) est ouvert dans $Y$.
\end{corollary}

Le même raisonnement que dans (11.1.5) montre que si $V$ est l'ensemble des
$x\in X$ où $\mathcal F$ est $(g\circ f)$-plat, $U$ est l'ensemble des $y\in Y$
dont toute générisation appartient à $E=Y-f(X-V)$. Comme $g\circ f$ est de
présentation finie, $V$ est ouvert dans $X$ en vertu de (11.3.1), donc
ind-constructible (1.9.6), et par suite $X-V$ est pro-constructible~; mais comme
$f$ est quasi-compact, $f(X-V)$ est pro-constructible dans $Y$ (1.9.5, (vii)),
et par suite $E$ est ind-constructible dans $Y$. Mais il résulte alors de
(1.10.1) que $U$ est l'intérieur de $E$, donc ouvert dans $Y$.

\begin{corollary}[11.3.3]
\label{IV.11.3.3-fr}
Soient $A$ un anneau, $B$ une $A$-algèbre de présentation finie, $C$ une
$B$-algèbre de présentation finie, $M$ un $C$-module de présentation finie.
Alors l'ensemble des $\mathfrak q\in\operatorname{Spec}(B)$ tels que
$M_\mathfrak q$ soit un $A$-module plat est ouvert dans
$\operatorname{Spec}(B)$.
\end{corollary}

\begin{proposition}[11.3.4]
\label{IV.11.3.4-fr}
Soient $f:X\to S$ un morphisme localement de présentation finie, $\mathcal J$
un Idéal quasi-cohérent de type fini de $\mathcal O_X$, $Y$ le sous-préschéma
fermé de $X$ défini par $\mathcal J$, $\mathcal F$ un $\mathcal O_X$-Module de
présentation finie. On suppose que $\operatorname{gr}_{\mathcal J}^*(\mathcal F)$
est $f$-plat. Dans ces conditions~:
\begin{enumerate}[label=\emph{(\roman*)}]
\item $\mathcal F$ est $f$-plat aux points de $Y$.
\item Si $\operatorname{gr}_{\mathcal J}^*(\mathcal O_X)$ est $f$-plat,
$\operatorname{gr}_{\mathcal J}^*(\mathcal O_X)$ est une
$(\mathcal O_X/\mathcal J)$-Algèbre de présentation finie et
$\operatorname{gr}_{\mathcal J}^*(\mathcal F)$ est un
$(\operatorname{gr}_{\mathcal J}^*(\mathcal O_X))$-Module de présentation
finie.
\item Supposons $\operatorname{gr}_{\mathcal J}^*(\mathcal O_X)$ $f$-plat
(ce qui entraîne que $\mathcal O_X/\mathcal J$ est $f$-plat). Alors l'ensemble
$W$ des $y\in Y$ tels que $(\operatorname{gr}_{\mathcal J}^*(\mathcal F))_y$
soit un $(\mathcal O_X/\mathcal J)_y$-module plat est ouvert dans $Y$.
\item Supposons $\operatorname{gr}_{\mathcal J}^*(\mathcal O_X)$ $f$-plat.
Soit $S'\to S$ un morphisme quelconque~; soient $X'=X\times_SS'$,
$p:X'\to X$ la projection canonique, $Y'=p^{-1}(Y)=Y\times_SS'$ le
sous-préschéma fermé de $X'$, défini par
$\mathcal J'=p^*(\mathcal J)\mathcal O_{X'}$, $\mathcal F'=p^*(\mathcal F)
=\mathcal F\otimes_{\mathcal O_S}\mathcal O_{S'}$~; alors, si
$W'=p^{-1}(W)$, on a
$\operatorname{gr}_{\mathcal J'}^*(\mathcal F')|W'
\simeq p^*(\operatorname{gr}_{\mathcal J}^*(\mathcal F)|W)$, et
$\operatorname{gr}_{\mathcal J'}^*(\mathcal F')$ est un
$(\mathcal O_{X'}/\mathcal J')$-Module plat aux points de $W'$.
\end{enumerate}
\end{proposition}

\begin{proof}
Les questions étant locales sur $X$ et $S$, on peut supposer que
$S=\operatorname{Spec}(A)$ et $X=\operatorname{Spec}(B)$ sont affines, $B$ étant
une $A$-algèbre de présentation finie, et $\mathcal F=\widetilde M$, où $M$ est
un $B$-module de présentation finie, $\mathcal J=\widetilde{\mathfrak J}$, où
$\mathfrak J$ est un idéal de type fini de $B$~; en vertu de (8.9.1) et
(8.5.11), il existe un sous-anneau noethérien $A_0$ de $A$, une $A_0$-algèbre
de type fini $B_0$ tels que $B=B_0\otimes_{A_0}A$, un idéal $\mathfrak J_0$ de
$B_0$ tel que $\mathfrak J=\mathfrak J_0B$, un $B_0$-module de type fini $M_0$
tel que $M=M_0\otimes_{A_0}A=M_0\otimes_{B_0}B$. En outre, $A$ est limite
inductive de ses sous-$A_0$-algèbres de type fini $A_\lambda$~; on pose
$B_\lambda=B_0\otimes_{A_0}A_\lambda$,
$M_\lambda=M_0\otimes_{A_0}A_\lambda=M_0\otimes_{B_0}B_\lambda$,
$\mathfrak J_\lambda=\mathfrak J_0B_\lambda$, de sorte que
$B=\varinjlim B_\lambda$, $M=\varinjlim M_\lambda$. Cela étant, par hypothèse
$\operatorname{gr}_{\mathfrak J}^*(M)$ est un $A$-module plat~; donc il résulte
de (11.2.9) qu'il existe un $\lambda$ tel que
$\operatorname{gr}_{\mathfrak J_\lambda}^*(M_\lambda)$ soit un
$A_\lambda$-module plat. Pour prouver que $\mathcal F$ est $f$-plat aux points
de $Y$, on peut donc se borner au cas où $S$ est noethérien~; mais alors
$(\mathbf 0_{\mathrm I},6.1.2)$, appliqué par récurrence sur $n$, prouve que les
$\mathcal F/\mathcal J^{n+1}\mathcal F$ sont $f$-plats, et on conclut par
$(\mathbf 0_{\mathrm{III}},10.2.6)$.

\oldpage[IV-3]{134}

La démonstration de (ii) se ramène de la même façon au cas où $S$ est
noethérien, en utilisant (11.2.9)~; la conclusion est alors évidente,
$\operatorname{gr}_{\mathfrak J}^*(B)$ étant dans ce cas une
$(B/\mathfrak J)$-algèbre de type fini et
$\operatorname{gr}_{\mathfrak J}^*(M)$ un
$\operatorname{gr}_{\mathfrak J}^*(B)$-module de type fini.

Pour prouver (iii), on se ramène comme dans (i) au cas où $S$ et $X$ sont
affines~; avec les mêmes notations, on peut supposer, en vertu de (11.2.9) que
$\operatorname{gr}_{\mathfrak J_\lambda}^*(B_\lambda)$ et
$\operatorname{gr}_{\mathfrak J_\lambda}^*(M_\lambda)$ sont des
$A_\lambda$-modules plats et que l'on a
\[
\operatorname{gr}_{\mathfrak J}^*(B)
=\operatorname{gr}_{\mathfrak J_\lambda}^*(B_\lambda)\otimes_{A_\lambda}A
\quad\text{et}\quad
\operatorname{gr}_{\mathfrak J}^*(M)
=\operatorname{gr}_{\mathfrak J_\lambda}^*(M_\lambda)\otimes_{A_\lambda}A.
\]
Par suite $\operatorname{gr}_{\mathfrak J}^*(B)$ est une
$(B/\mathfrak J)$-algèbre de présentation finie et
$\operatorname{gr}_{\mathfrak J}^*(M)$ un
$\operatorname{gr}_{\mathfrak J}^*(B)$-module de présentation finie. Si $W^m$
est l'ensemble des $\mathfrak p\in\operatorname{Spec}(B/\mathfrak J)$ tels que
$\operatorname{gr}_{\mathfrak J}^m(M)_\mathfrak p$ soit un
$(B/\mathfrak J)$-module plat, $W^m$ est ouvert dans
$\operatorname{Spec}(B/\mathfrak J)$ (11.3.1) et l'on a
$W=\bigcap_mW^m$, donc $W$ est stable par générisation. L'assertion (iii)
résulte alors de (11.3.2) appliqué en y prenant
$X=\operatorname{Spec}(\operatorname{gr}_{\mathfrak J}^*(B))$,
$Y=Z=\operatorname{Spec}(B/\mathfrak J)$ et
$\mathcal F=(\operatorname{gr}_{\mathfrak J}^*(M))^{\sim}$.

Enfin, (iv) découle aussitôt de (11.2.9.2).

Généralisant la définition de (6.10.1), on dit que pour un préschéma $X$, un
sous-préschéma fermé $Y$ de $X$ défini par un Idéal quasi-cohérent $\mathcal J$
de $\mathcal O_X$, et un $\mathcal O_X$-Module quasi-cohérent $\mathcal F$,
$\mathcal F$ est \emph{normalement plat le long de $Y$ en un point $y$} si
$(\operatorname{gr}_{\mathcal J}^*(\mathcal F))_y$ est un
$\mathcal O_{Y,y}$-module plat. On dit que $\mathcal F$ est
\emph{normalement plat le long de $Y$} s'il l'est en tous les points de $Y$.
\end{proof}

\begin{corollary}[11.3.5]
\label{IV.11.3.5-fr}
Sous les hypothèses générales de (11.3.4), on suppose que
$\operatorname{gr}_{\mathcal J}^*(\mathcal O_X)$ et
$\operatorname{gr}_{\mathcal J}^*(\mathcal F)$ soient $f$-plats. Alors
l'ensemble $W$ des $y\in Y$ tels que $\mathcal F$ soit normalement plat le long
de $Y$ au point $y$ est ouvert dans $Y$, et (avec les notations de
(11.3.4, (iv))) $\mathcal F'$ est normalement plat le long de $Y'$ en tous les
points de $W'=p^{-1}(W)$.
\end{corollary}

\begin{proposition}[11.3.6]
\label{IV.11.3.6-fr}
Les notations étant celles de (8.5.1) et (8.8.1), on suppose $S_\alpha$
quasi-compact, $X_\alpha$ de présentation finie sur $S_\alpha$, $Y_\alpha$ un
sous-préschéma fermé de $X_\alpha$ défini par un Idéal quasi-cohérent
$\mathcal J_\alpha$ de type fini de $\mathcal O_{X_\alpha}$ tel que
$\operatorname{gr}_{\mathcal J_\alpha}^*(\mathcal O_{X_\alpha})$ soit plat sur
$S_\alpha$~; enfin on suppose que $\mathcal F_\alpha$ est un
$\mathcal O_{X_\alpha}$-Module de présentation finie. Pour que $\mathcal F$
soit normalement plat le long de $Y$, il faut et il suffit qu'il existe
$\lambda$ tel que $\mathcal F_\lambda$ soit normalement plat le long de
$Y_\lambda$.
\end{proposition}

\begin{proof}
Notons que $Y$ (resp. $Y_\lambda$) est le sous-préschéma fermé de $X$ (resp.
$X_\lambda$) défini par $\mathcal J=\mathcal J_\alpha\mathcal O_X$ (resp.
$\mathcal J_\lambda=\mathcal J_\alpha\mathcal O_{X_\lambda}$)~; en vertu de
l'hypothèse et de (11.2.9.2), on a
\[
\operatorname{gr}_{\mathcal J}^*(\mathcal O_X)
=\operatorname{gr}_{\mathcal J_\alpha}^*(\mathcal O_{X_\alpha})
 \otimes_{\mathcal O_{S_\alpha}}\mathcal O_S
\]
et
\[
\operatorname{gr}_{\mathcal J_\lambda}^*(\mathcal O_{X_\lambda})
=\operatorname{gr}_{\mathcal J_\alpha}^*(\mathcal O_{X_\alpha})
 \otimes_{\mathcal O_{S_\alpha}}\mathcal O_{S_\lambda}
\quad\text{pour }\lambda\geq\alpha,
\]
et $\operatorname{gr}_{\mathcal J}^*(\mathcal O_X)$ (resp.
$\operatorname{gr}_{\mathcal J_\lambda}^*(\mathcal O_{X_\lambda})$) est plat
sur $S$ (resp. $S_\lambda$), ce qui entraîne que $Y$ est plat sur $S$ (resp.
$Y_\lambda$ plat sur $S_\lambda$). Si $\mathcal F_\lambda$ est normalement plat
le long de $Y_\lambda$, $\operatorname{gr}_{\mathcal J_\lambda}^*
(\mathcal F_\lambda)|Y_\lambda$ est plat sur $Y_\lambda$, donc aussi sur
$S_\lambda$, et comme
$\operatorname{gr}_{\mathcal J_\lambda}^*(\mathcal F_\lambda)|x_\lambda=0$
pour $x_\lambda\notin Y_\lambda$,
$\operatorname{gr}_{\mathcal J_\lambda}^*(\mathcal F_\lambda)$ est plat sur
$S_\lambda$. On en conclut (11.2.9.2) que
\[
\operatorname{gr}_{\mathcal J}^*(\mathcal F)
=\operatorname{gr}_{\mathcal J_\lambda}^*(\mathcal F_\lambda)
 \otimes_{\mathcal O_{S_\lambda}}\mathcal O_S,
\]
donc $\mathcal F$ est normalement plat le long de $Y$. Pour prouver la
réciproque, on peut supposer que $S_\alpha$ et $X_\alpha$ sont affines et
adopter les notations de (11.2.9)~; puisque
$\operatorname{gr}_{\mathfrak J}^*(B)$ est un $A$-module plat, il en est de même
de $\operatorname{gr}_{\mathfrak J}^*(M)$, donc, en vertu de (11.2.9), il
existe $\lambda\geq\alpha$ tel que
$\operatorname{gr}_{\mathfrak J_\lambda}^*(M_\lambda)$ soit un
$A_\lambda$-module plat, d'où
$\operatorname{gr}_{\mathfrak J}^*(M)
=\operatorname{gr}_{\mathfrak J_\lambda}^*(M_\lambda)\otimes_{A_\lambda}A$.
En outre (11.3.4, (ii)),
$\operatorname{gr}_{\mathfrak J_\lambda}^*(M_\lambda)$ est un
$\operatorname{gr}_{\mathfrak J_\lambda}^*(B_\lambda)$-module de présentation
finie et $\operatorname{gr}_{\mathfrak J_\lambda}^*(B_\lambda)$ une
$(B_\lambda/\mathfrak J_\lambda)$-algèbre de présentation finie. La conclusion
résulte alors de (11.2.6).
\end{proof}

\oldpage[IV-3]{135}

\begin{proposition}[11.3.7]
\label{IV.11.3.7-fr}
Soient $f:X\to Y$ un morphisme localement de présentation finie,
$\mathcal F$, $\mathcal F'$ deux $\mathcal O_X$-Modules quasi-cohérents de
présentation finie, $u:\mathcal F'\to\mathcal F$ un
$\mathcal O_X$-homomorphisme, $x$ un point de $X$ et $y=f(x)$~; on suppose que
$\mathcal F$ est $f$-plat au point $x$. Les conditions suivantes sont alors
équivalentes~:
\begin{enumerate}[label=\emph{\alph*)}]
\item On a $(\ker u)_x=0$ et $\operatorname{Coker}u$ est un
$\mathcal O_X$-Module $f$-plat au point $x$.
\item L'homomorphisme
\[
(u\otimes1)_x:(\mathcal F'\otimes_{\mathcal O_Y}k(y))_x
\longrightarrow(\mathcal F\otimes_{\mathcal O_Y}k(y))_x
\]
est injectif.
\end{enumerate}
Si de plus $\mathcal F$ est $f$-plat, l'ensemble des points $x$ vérifiant les
conditions équivalentes précédentes est ouvert dans $X$.
\end{proposition}

\begin{proof}
La condition a) entraîne b) en vertu de $(\mathbf 0_{\mathrm I},6.1.2)$, sans
hypothèse sur $\mathcal F$. Pour démontrer la réciproque, on peut se borner au
cas où $X$ et $Y$ sont affines, puis, en vertu de (8.9.1) et (11.2.7), supposer
que l'on a
\[
X=X_0\times_{Y_0}Y,\quad f=f_0\times1,\quad
\mathcal F=\mathcal F_0\otimes_{\mathcal O_{Y_0}}\mathcal O_Y,
\]
\[
\mathcal F'=\mathcal F'_0\otimes_{\mathcal O_{Y_0}}\mathcal O_Y,
\qquad u=u_0\otimes1_Y,
\]
où $Y_0$ est noethérien, $f_0:X_0\to Y_0$ un morphisme de type fini,
$\mathcal F_0$, $\mathcal F'_0$ deux $\mathcal O_{X_0}$-Modules cohérents,
$u_0:\mathcal F'_0\to\mathcal F_0$ un homomorphisme~; en outre, si $x_0$ est la
projection de $x$ dans $X_0$, on peut supposer que $\mathcal F_0$ est
$f_0$-plat au point $x_0$. Posons $y_0=f_0(x_0)$~; en vertu de la transitivité
des fibres $(\mathbf I,3.6.4)$, le morphisme projection
$f^{-1}(y)\to f_0^{-1}(y_0)$ est fidèlement plat (2.2.13), et comme
\[
(u\otimes1)_x=((u_0\otimes1)_{x_0})\otimes1_{k(y)},
\]
l'hypothèse que $(u\otimes1)_x$ est injectif entraîne qu'il en est de même de
$(u_0\otimes1)_{x_0}$ (2.2.7). Or, cela entraîne, par
$(\mathbf 0_{\mathrm{III}},10.2.4)$ appliqué aux anneaux locaux noethériens
$\mathcal O_{Y_0,y_0}$ et $\mathcal O_{X_0,x_0}$ et aux
$\mathcal O_{X_0,x_0}$-modules $(\mathcal F'_0)_{x_0}$ et
$(\mathcal F_0)_{x_0}$ (dont le second est plat sur $\mathcal O_{Y_0,y_0}$),
que l'on a $(\ker u_0)_{x_0}=0$ et que $\operatorname{Coker}u_0$ est
$f_0$-plat au point $x_0$. On déduit d'abord de là que
$\operatorname{Coker}u$ est $f$-plat au point $x$ (2.1.4)~; en vertu de
l'exactitude à droite du produit tensoriel, on a d'ailleurs
\[
\operatorname{Coker}u=(\operatorname{Coker}u_0)
 \otimes_{\mathcal O_{Y_0}}\mathcal O_Y~;
\]
appliquant alors $(\mathbf 0_{\mathrm I},6.1.2)$ à la suite (exacte par
hypothèse)
\[
0\longrightarrow(\mathcal F'_0)_{x_0}\longrightarrow
(\mathcal F_0)_{x_0}\longrightarrow(\operatorname{Coker}u_0)_{x_0}
\longrightarrow0,
\]
on en déduit que $u_x$ est un homomorphisme injectif.

Enfin, il résulte de (11.1.2) que l'ensemble $U_0$ des points $z_0\in X_0$ tels
que le morphisme $(u_0\otimes1)_{z_0}$ soit injectif est ouvert~; par platitude
on en déduit que pour tout $z\in X$ au-dessus de $z_0\in U_0$ le morphisme
$(u\otimes1)_z$ est injectif, donc l'ensemble de ces points contient l'image
réciproque $U$ de $U_0$ dans $X$ et est par suite un voisinage de $x$, ce qui
achève la démonstration.
\end{proof}

\begin{theorem}[11.3.8]
\label{IV.11.3.8-fr}
Soient $f:X\to Y$ un morphisme localement de présentation finie,
$\mathcal F$ un $\mathcal O_X$-Module quasi-cohérent de présentation finie,
$(g_i)_{1\leq i\leq n}$ une suite de sections de $\mathcal O_X$ au-dessus de
$X$~; posons
\[
\mathcal F_i=\mathcal F\big/\sum_{j=1}^{i}g_j\mathcal F
\quad\text{pour }0\leq i\leq n
\quad(\text{avec }\mathcal F_0=\mathcal F).
\]
Soient $x$ un point de $X$, $y=f(x)$, et posons
$X_y=f^{-1}(y)=X\times_Y\operatorname{Spec}(k(y))$,
$\mathcal F_y=\mathcal F\otimes_{\mathcal O_Y}k(y)$, qui est un
$\mathcal O_{X_y}$-Module. On suppose que les $(g_i)_x$ appartiennent à l'idéal
maximal $\mathfrak m_x$. Les conditions suivantes sont équivalentes~:
\begin{enumerate}[label=\emph{\alph*)}]
\item La suite $((g_i)_x)$ est $\mathcal F_x$-régulière et les $\mathcal F_i$
$(0\leq i\leq n)$ sont $f$-plats au point $x$.
\item La suite $((g_i)_x)$ est $\mathcal F_x$-régulière et $\mathcal F_n$ est
$f$-plat au point $x$.
\item[\emph{b$'$)}] Il existe un voisinage $U$ de $x$ tel que la suite
$(g_i|U)$ soit $(\mathcal F|U)$-quasi-régulière, et $\mathcal F_n$ est
$f$-plat au point $x$.

\oldpage[IV-3]{136}

\item $\mathcal F$ est $f$-plat au point $x$, et la suite
$((g_i\otimes1)_x)$ d'éléments de $\mathcal O_{X_y,x}$ images des $(g_i)_x$ est
$(\mathcal F_y)_x$-régulière.
\item $\mathcal F$ est $f$-plat au point $x$, et pour tout morphisme
$Y'\to Y$ et tout point $x'\in X'=X\times_YY'$ au-dessus de $x$, si l'on pose
$\mathcal F'=\mathcal F\otimes_YY'$, la suite $(g_i\otimes1)_{x'}$
d'éléments de $\mathcal O_{X',x'}$ est $\mathcal F'_{x'}$-régulière.
\end{enumerate}
En outre l'ensemble des $x\in X$ vérifiant ces conditions est ouvert dans
l'ensemble $Z$ des $x$ tels que $g_i(x)=0$ pour tout $i$.
\end{theorem}

\begin{proof}
Le fait que a) entraîne d) résulte aussitôt de $(\mathbf 0,15.1.15)$, et c) est
un cas particulier de d)~; par ailleurs a) implique b) trivialement. Montrons
que b) ou c) entraîne a). Les $\mathcal F_i$ sont de présentation finie~; le
fait que c) implique a) résulte alors aussitôt de (11.3.7) par récurrence sur
$n$, et cela montre aussi que l'ensemble des $x\in X$ vérifiant c) est ouvert
dans $Z$. Pour montrer que b) entraîne a), on est aussitôt ramené, par
récurrence sur $n$, au cas $n=1$~; nous écrirons $g$ au lieu de $g_1$. La
question étant locale sur $X$ et $Y$, on peut supposer
$Y=\operatorname{Spec}(A)$, $X=\operatorname{Spec}(B)$ affines, $B$ étant une
$A$-algèbre de présentation finie, $\mathcal F=\widetilde M$, où $M$ est un
$B$-module de présentation finie. On peut alors (8.9.1) écrire
$B=B_0\otimes_{A_0}A$, $M=M_0\otimes_{A_0}A$, où $A_0$ est un sous-anneau
noethérien de $A$, $B_0$ une $A_0$-algèbre de type fini et $M_0$ un
$B_0$-module de type fini. En outre $A$ est limite inductive filtrante de ses
sous-anneaux $A_\lambda$ qui sont des $A_0$-algèbres de type fini (donc
noethériens), et si l'on pose
$B_\lambda=B_0\otimes_{A_0}A_\lambda$,
$M_\lambda=M_0\otimes_{A_0}A_\lambda$, $B_\lambda$ est une
$A_\lambda$-algèbre de type fini, $M_\lambda$ un $B_\lambda$-module de type
fini et l'on a $B=\varinjlim B_\lambda$, $M=\varinjlim M_\lambda$. Il existe
donc un indice $\lambda$ et un élément $g_\lambda\in B_\lambda$ tels que
$g=g_\lambda\otimes1$~; revenant aux notations géométriques et posant
$Y_\lambda=\operatorname{Spec}(A_\lambda)$,
$X_\lambda=\operatorname{Spec}(B_\lambda)$ et
$\mathcal F_\lambda=\widetilde{M_\lambda}$, il nous suffira de prouver qu'il y
a un $\mu\geq\lambda$ tel qu'au point $x_\mu\in X_\mu$ projection de $x$,
$(g_\mu)_{x_\mu}$ soit $(\mathcal F_\mu)_{x_\mu}$-régulier et
$\mathcal F_\mu/g_\mu\mathcal F_\mu$ plat sur $Y_\mu$ au point $x_\mu$. On en
déduira en effet, par $(\mathbf 0,15.1.16)$ que $\mathcal F_\mu$ est plat sur
$Y_\mu$ au point $x_\mu$, donc $\mathcal F$ plat sur $Y$ au point $x$.

Le fait que b) entraîne a) résultera donc de la proposition suivante~:

\begin{proposition}[11.3.9]
\label{IV.11.3.9-fr}
Les notations et hypothèses générales étant celles de (8.5.1) et (8.8.1),
supposons que $f_\alpha:X_\alpha\to Y_\alpha$ soit un morphisme localement de
présentation finie. Soient $\mathcal F_\alpha$, $\mathcal G_\alpha$ deux
$\mathcal O_{X_\alpha}$-Modules quasi-cohérents de présentation finie,
$h_\alpha:\mathcal F_\alpha\to\mathcal G_\alpha$ un
$\mathcal O_{X_\alpha}$-homomorphisme, $\mathcal H_\alpha$ son conoyau. Soient
enfin $x$ un point de $X$, $x_\lambda$ sa projection dans $X_\lambda$ pour
$\lambda\geq\alpha$. Pour que $h_x$ soit injectif et que
$\mathcal H=\operatorname{Coker}h$ soit $f$-plat au point $x$, il faut et il
suffit qu'il existe $\lambda\geq\alpha$ tel que $(h_\lambda)_{x_\lambda}$ soit
injectif et que $\mathcal H_\lambda=\operatorname{Coker}h_\lambda$ soit
$f_\lambda$-plat au point $x_\lambda$. En outre, l'ensemble des $x\in X$ ayant
les propriétés précédentes est ouvert dans $X$.
\end{proposition}

Rappelons qu'en vertu de l'exactitude à droite du produit tensoriel, on a
$\mathcal H_\mu=v_{\lambda\mu}^*(\mathcal H_\lambda)$ pour $\lambda\leq\mu$ et
$\mathcal H=v_\lambda^*(\mathcal H_\lambda)$, ce qui justifie les notations.

La suffisance de la condition provient de ce que, si la suite
\[
0\longrightarrow(\mathcal F_\lambda)_{x_\lambda}
\longrightarrow(\mathcal G_\lambda)_{x_\lambda}
\longrightarrow(\mathcal H_\lambda)_{x_\lambda}\longrightarrow0
\]
est exacte et $(\mathcal H_\lambda)_{x_\lambda}$ un
$\mathcal O_{f_\lambda(x_\lambda)}$-module plat, $\mathcal H_x$ est un
$\mathcal O_{f(x)}$-module plat par changement de base (2.1.4) et la suite
$0\to\mathcal F_x\to\mathcal G_x\to\mathcal H_x\to0$ est exacte (2.1.8).
Pour prouver que

\oldpage[IV-3]{137}

la condition est nécessaire, notons que $\mathcal H$ est de présentation
finie~; la question étant locale sur $X$, on peut supposer $X$ et $Y$ affines,
et, en vertu de (11.3.1), supposer que $\mathcal H$ est $f$-plat. Notons
maintenant le

\begin{lemma}[11.3.9.1]
\label{IV.11.3.9.1-fr}
Soient $f:X\to Y$ un morphisme localement de présentation finie,
$\mathcal G$, $\mathcal H$ deux $\mathcal O_X$-Modules quasi-cohérents de
présentation finie tels que $\mathcal H$ soit $f$-plat,
$p:\mathcal G\to\mathcal H$ un homomorphisme surjectif. Alors $\ker(p)$ est un
$\mathcal O_X$-Module de présentation finie.
\end{lemma}

On peut en effet supposer $X$, $Y$ affines, et qu'il existe un morphisme
$Y\to Y_0$ où $Y_0$ est noethérien, un morphisme $f_0:X_0\to Y_0$ de type fini
tels que $X=X_0\times_{Y_0}Y$, $f=(f_0)_{(Y)}$, deux
$\mathcal O_{X_0}$-Modules cohérents $\mathcal G_0$, $\mathcal H_0$ et un
homomorphisme $p_0:\mathcal G_0\to\mathcal H_0$ tels que $\mathcal G$,
$\mathcal H$ et $p$ se déduisent de $\mathcal G_0$, $\mathcal H_0$ et $p_0$ par
changement de base (8.9.1 et 8.5.2). En outre, on peut supposer $p_0$ surjectif
(8.5.7) et $\mathcal H_0$ $f_0$-plat (11.2.7). Alors, si
$\mathcal K_0=\ker(p_0)$, $\mathcal K_0$ est un
$\mathcal O_{X_0}$-Module cohérent $(\mathbf 0_{\mathrm I},5.3.4)$ et en vertu
de $(\mathbf 0_{\mathrm I},6.1.2)$, $\mathcal K=\ker(p)$ se déduit de
$\mathcal K_0$ par changement de base, donc est de présentation finie.

Ce lemme étant démontré, posons $\mathcal K=\ker(\mathcal G\to\mathcal H)$,
qui est donc de présentation finie. On a un homomorphisme canonique
$q:\mathcal F\to\mathcal K$, et par hypothèse
$q_x:\mathcal F_x\to\mathcal K_x$ est un isomorphisme~; par suite
$(\mathbf 0_{\mathrm I},5.2.7)$ il existe un voisinage $U$ de $x$ tel que
$q|U$ soit un isomorphisme, et en restreignant $X$, on peut supposer que la
suite
\[
0\longrightarrow\mathcal F\xrightarrow{h}\mathcal G
\longrightarrow\mathcal H\longrightarrow0
\]
est \emph{exacte}. Cela étant, il résulte de (11.2.6) que pour un $\lambda$
assez grand, $\mathcal H_\lambda$ est un $\mathcal O_{X_\lambda}$-Module plat~;
posons $\mathcal K_\lambda=\ker(\mathcal G_\lambda\to\mathcal H_\lambda)$ et
$\mathcal K_\mu=v_{\lambda\mu}^*(\mathcal K_\lambda)$ pour $\mu\geq\lambda$~;
il résulte de $(\mathbf 0_{\mathrm I},6.1.2)$ que l'on a
$\mathcal K_\mu=\ker(\mathcal G_\mu\to\mathcal H_\mu)$ et
$\mathcal K=v_\lambda^*(\mathcal K_\lambda)
=\ker(\mathcal G\to\mathcal H)=\mathcal F$. On a d'autre part pour tout
$\mu\geq\lambda$ un homomorphisme canonique
$q_\mu:\mathcal F_\mu\to\mathcal K_\mu$ avec
$q_\mu=v_{\lambda\mu}^*(q_\lambda)$ pour $\lambda\leq\mu$, et
$q=v_\mu^*(q_\mu)$. Comme $q$ est un isomorphisme, il résulte de (8.5.2.4) et
(11.3.9.1) qu'il existe $\mu\geq\lambda$ tel que $q_\mu$ soit un isomorphisme,
et par suite $u_\mu:\mathcal F_\mu\to\mathcal G_\mu$ un homomorphisme
injectif.

Revenons à la démonstration de (11.3.8).

Puisque l'ensemble des $x\in X$ vérifiant b) est ouvert dans $X$, il est clair
que b) entraîne b$'$. Montrons enfin que b$'$ entraîne c). En premier lieu,
$\mathcal F_n$ est $f$-plat dans un voisinage de $x$ (11.3.1) et on peut donc
se borner au cas où $U=X$ et où $\mathcal F_n$ est $f$-plat. Comme par
définition $\operatorname{gr}_{\mathcal J}^*(\mathcal F)$ est isomorphe à
$\mathcal F_n\otimes_{\mathcal O_Y}\mathcal O_Y[T_1,\ldots,T_n]$
$(\mathbf 0,15.2.2)$, il est $f$-plat, et on en conclut par (11.3.4, (i)) que
$\mathcal F$ lui-même est $f$-plat au voisinage de $x$. D'autre part, si
$(\mathcal J_y)_x$ est l'idéal de $\mathcal O_{X_y,x}$ engendré par les
$(g_i\otimes1)_x$, il résulte de (11.2.9.2) que, dans le diagramme
\[
\xymatrix{
(\operatorname{gr}_{\mathcal J_x}^0(\mathcal F_x)[T_1,\ldots,T_n])
 \otimes_{\mathcal O_y}k(y) \ar[r] \ar[d] &
\operatorname{gr}_{\mathcal J_x}^*(\mathcal F_x)
 \otimes_{\mathcal O_y}k(y) \ar[d] \\
(\operatorname{gr}_{(\mathcal J_y)_x}^0((\mathcal F_y)_x))
 [T_1,\ldots,T_n] \ar[r] &
\operatorname{gr}_{(\mathcal J_y)_x}^*((\mathcal F_y)_x)
}
\]

\oldpage[IV-3]{138}

les flèches verticales sont des isomorphismes~; comme la première ligne est un
isomorphisme, il en est de même de la seconde, donc la suite
$((g_i\otimes1)_x)$ est $(\mathcal F_y)_x$-quasi-régulière, et par suite aussi
$(\mathcal F_y)_x$-régulière, puisque $X_y$ est localement de type fini sur
$k(y)$. C.Q.F.D.
\end{proof}

\begin{theorem}[11.3.10]
\label{IV.11.3.10-fr}
(critère de platitude par fibres). --- Soient $S$ un préschéma,
$g:X\to S$, $h:Y\to S$ deux morphismes, $f:X\to Y$ un $S$-morphisme,
$\mathcal F$ un $\mathcal O_X$-Module quasi-cohérent, $x$ un point de $X$,
$y=f(x)$, $s=h(y)=g(x)$. On suppose vérifiée l'une des deux hypothèses
suivantes~:
\begin{enumerate}[label=\arabic{enumi}\textsuperscript{o}]
\item $S$, $X$ et $Y$ sont localement noethériens et $\mathcal F$ est cohérent.
\item $g$ et $h$ sont localement de présentation finie et $\mathcal F$ est de
présentation finie.
\end{enumerate}
Alors, avec les notations de (9.4.1), si $\mathcal F_x\neq0$, les conditions
suivantes sont équivalentes~:
\begin{enumerate}[label=\emph{\alph*)}]
\item $\mathcal F$ est $g$-plat au point $x$ et $\mathcal F_s$ est $f_s$-plat
au point $x$.
\item $h$ est plat au point $y$ et $\mathcal F$ est $f$-plat au point $x$.
\end{enumerate}
En outre, dans l'hypothèse 2\textsuperscript{o}, l'ensemble des $x\in X$
vérifiant la condition b) est ouvert dans $X$.
\end{theorem}

\begin{proof}
La dernière assertion résulte de (11.3.1) appliqué à $\mathcal O_Y$ et au
morphisme $h$ d'une part, à $\mathcal F$ et au morphisme $f$ (qui est
localement de présentation finie) de l'autre.

Comme $g=h\circ f$, il est clair que b) implique a) sans supposer
1\textsuperscript{o} ni 2\textsuperscript{o} (2.1.6 et 2.1.4). Pour prouver que
a) entraîne b), on peut se borner au cas où $S$, $X$ et $Y$ sont affines~;
sous l'hypothèse 2\textsuperscript{o}, en appliquant (11.2.7), on se ramène au
cas où de plus $S$ est noethérien, c'est-à-dire qu'on peut se borner à
considérer le cas où l'hypothèse 1\textsuperscript{o} est satisfaite. Alors
l'assertion équivaut au lemme suivant, qui améliore
$(\mathbf 0_{\mathrm{III}},10.2.5)$~:

\begin{lemma}[11.3.10.1]
\label{IV.11.3.10.1-fr}
Soient $\rho:A\to B$, $\sigma:B\to C$ des homomorphismes locaux d'anneaux
locaux noethériens, $k$ le corps résiduel de $A$, $M$ un $C$-module
$\neq0$ de type fini. Alors les conditions suivantes sont équivalentes~:
\begin{enumerate}[label=\emph{\alph*)}]
\item $M$ est un $A$-module plat et $M\otimes_Ak$ est un
$(B\otimes_Ak)$-module plat.
\item $B$ est un $A$-module plat et $M$ est un $B$-module plat.
\end{enumerate}
\end{lemma}

Nous établirons d'abord le lemme plus général suivant~:

\begin{lemma}[11.3.10.2]
\label{IV.11.3.10.2-fr}
Soient $A$ un anneau commutatif, $B$ une $A$-algèbre commutative,
$\mathfrak J$ un idéal de $A$, $M$ un $B$-module. On considère d'une part les
conditions suivantes~:
\begin{enumerate}[label=\emph{(\roman*)}]
\item $\mathfrak J$ est nilpotent.
\item $B$ est noethérien et $M$ est idéalement séparé pour la topologie
$\mathfrak JB$-préadiques $(\mathbf 0_{\mathrm{III}},10.2.1)$.
\item $\mathfrak JB$ est contenu dans le radical de $B$.
\end{enumerate}
On considère d'autre part les quatre propriétés~:
\begin{enumerate}[label=\emph{\alph*)}]
\item $M$ est un $B$-module plat.
\item $B$ est un $A$-module plat.
\item $M$ est un $A$-module plat et $M/\mathfrak JM$ un
$(B/\mathfrak JB)$-module plat.
\item $B/\mathfrak JB$ est un $(A/\mathfrak J)$-module plat et pour tout idéal
maximal $\mathfrak m\supset\mathfrak JB$ de $B$ on a $\mathfrak mM\neq M$.
\end{enumerate}
Alors~:
\begin{enumerate}[label=\arabic{enumi}\textsuperscript{o}]
\item Si l'une des conditions (i), (ii) est vérifiée, la conjonction de a) et
b) implique c), et c) implique a).

\oldpage[IV-3]{139}

\item Si la condition (i) ou la conjonction de (ii) et (iii) est vérifiée, la
conjonction de c) et d) implique la conjonction de a) et b).
\end{enumerate}
\end{lemma}

\noindent 1\textsuperscript{o} La première assertion est immédiate
$(\mathbf 0_{\mathrm I},6.2.1)$. Supposons donc c) vérifiée, et prouvons a).
Considérons les anneaux gradués $\operatorname{gr}_{\mathfrak J}^*(A)$,
$\operatorname{gr}_{\mathfrak J}^*(B)$ et le module gradué
$\operatorname{gr}_{\mathfrak J}^*(M)$ (à la fois sur
$\operatorname{gr}_{\mathfrak J}^*(A)$ et
$\operatorname{gr}_{\mathfrak J}^*(B)$) relatifs aux filtrations
$\mathfrak J$-préadique, ainsi que les applications canoniques surjectives
$(\mathbf 0_{\mathrm{III}},10.1.1.2)$
\begin{align*}
u&:\operatorname{gr}^0(B)\otimes_{\operatorname{gr}^0(A)}
   \operatorname{gr}^*(A)\longrightarrow\operatorname{gr}^*(B),\\
v&:\operatorname{gr}^0(M)\otimes_{\operatorname{gr}^0(A)}
   \operatorname{gr}^*(A)\longrightarrow\operatorname{gr}^*(M),\\
w&:\operatorname{gr}^0(M)\otimes_{\operatorname{gr}^0(B)}
   \operatorname{gr}^*(B)\longrightarrow\operatorname{gr}^*(M).
\end{align*}
Il est clair qu'on a un diagramme commutatif
\begin{equation}
\xymatrix{
\operatorname{gr}^0(M)\otimes_{\operatorname{gr}^0(A)}\operatorname{gr}^*(A)
 \ar[rr]^v \ar[dr]_{1\otimes u} && \operatorname{gr}^*(M) \\
& \operatorname{gr}^0(M)\otimes_{\operatorname{gr}^0(B)}
  \operatorname{gr}^*(B) \ar[ur]_w &
}
\label{IV.11.3.10.3-fr}
\tag{11.3.10.3}
\end{equation}
L'hypothèse c) entraîne que $v$ est bijective
$(\mathbf 0_{\mathrm{III}},10.2.1)$~; comme les deux autres applications du
diagramme sont surjectives, elles sont aussi bijectives. Mais comme en vertu de
l'hypothèse c), $M/\mathfrak JM$ est un $(B/\mathfrak JB)$-module plat, il
résulte de $(\mathbf 0_{\mathrm{III}},10.2.1)$ que $M$ est un $B$-module plat.

\noindent 2\textsuperscript{o} L'une ou l'autre des conditions (i), (iii)
implique que tout idéal maximal de $B$ contient $\mathfrak JB$. Il résulte donc
de 1\textsuperscript{o} et de la conjonction de c) et d) que $M$ est un
$B$-module \emph{fidèlement plat}, et par suite $\operatorname{gr}^0(M)$ un
$\operatorname{gr}^0(B)$-module \emph{fidèlement plat}
$(\mathbf 0_{\mathrm I},6.2.1)$. On a vu dans 1\textsuperscript{o} que
l'hypothèse c) entraîne que les trois applications du diagramme (11.3.10.3)
sont bijectives~; le fait que $\operatorname{gr}^0(M)$ soit un
$\operatorname{gr}^0(B)$-module fidèlement plat implique donc que $u$ est aussi
bijective $(\mathbf 0_{\mathrm I},6.4.1)$. D'autre part, les conditions (ii) et
(iii) impliquent que $B$ est un $A$-module idéalement séparé pour la filtration
$\mathfrak J$-préadique (Bourbaki, \emph{Alg. comm.}, chap.~III, \S~5,
n$^\circ$~4, prop.~2)~; on déduit donc encore de
$(\mathbf 0_{\mathrm{III}},10.2.1)$ que si la condition (i), ou la conjonction
de (ii) et (iii), est vérifiée, $B$ est un $A$-module plat.

\phantomsection\label{IV.11.3.10.4-fr}
\noindent\textbf{(11.3.10.4)} Le lemme (11.3.10.2) étant établi, on en déduit
(11.3.10.1) en prenant pour $\mathfrak J$ l'idéal maximal de $A$, et en
remarquant qu'alors les conditions (ii) et (iii) de (11.3.10.2) sont
satisfaites (Bourbaki, \emph{Alg. comm.}, chap.~III, \S~5, n$^\circ$~4,
prop.~2). Ceci achève donc aussi la démonstration de (11.3.7).
\end{proof}

\begin{corollary}[11.3.11]
\label{IV.11.3.11-fr}
Soient $g:X\to S$, $h:Y\to S$ deux morphismes localement de présentation
finie, $f:X\to Y$ un $S$-morphisme. Les conditions suivantes sont
équivalentes~:
\begin{enumerate}[label=\emph{\alph*)}]
\item $g$ est plat et $f_s:X_s\to Y_s$ est plat pour tout $s\in S$.
\item $h$ est plat en tous les points de $f(X)$ et $f$ est plat.
\end{enumerate}
\end{corollary}

Il suffit d'appliquer (11.3.10) pour $\mathcal F=\mathcal O_X$.

\begin{remark}[11.3.12]
\label{IV.11.3.12-fr}
Il serait intéressant de pouvoir donner de (11.3.2) et (11.3.10) des
démonstrations n'utilisant pas le passage à la limite inductive~; il suffirait
pour cela de démontrer le critère suivant~:

Soient $A$ un anneau, $B$ une $A$-algèbre de présentation finie, $M$ un
$B$-module de présentation finie, $\mathfrak J$ un

\oldpage[IV-3]{140}

idéal de $A$, $\mathfrak p$ un idéal premier de $B$ contenant $\mathfrak JB$.
Alors les deux conditions suivantes sont équivalentes~:
\begin{enumerate}[label=\emph{\alph*)}]
\item $M_\mathfrak p$ est un $A$-module plat.
\item $M_\mathfrak p/\mathfrak JM_\mathfrak p$ est un
$(A/\mathfrak J)$-module plat et
$\operatorname{Tor}_1^A(A/\mathfrak J,M_\mathfrak p)=0$.
\end{enumerate}
Lorsque $A$ est \emph{noethérien}, c'est une conséquence de
$(\mathbf 0_{\mathrm{III}},10.2.2)$ appliqué à la $A$-algèbre noethérienne
$B_\mathfrak p$~; peut-on en déduire l'énoncé général par un passage à la
limite inductive~?

Il convient de noter à ce propos qu'une telle généralisation n'est certainement
pas possible lorsqu'on remplace la condition b) ci-dessus par l'une des
conditions c), d) de $(\mathbf 0_{\mathrm{III}},10.2.1)$ où $M$ est remplacé
par $M_\mathfrak p$. Prenons par exemple pour $A$ un anneau local dont l'idéal
maximal $\mathfrak J$ est principal et tel que l'intersection
\[
\mathfrak R=\bigcap_{n\geq0}\mathfrak J^{n+1}
\]
ne soit pas réduite à $0$ (par exemple l'anneau des germes de fonctions
numériques indéfiniment différentiables au voisinage de $0$ dans $\mathbf R$).
Prenons $B=A$, $\mathfrak p=\mathfrak J$, et $M=A/\mathfrak N$, où
$\mathfrak N$ est un sous-module monogène $\neq0$ de $\mathfrak R$~; $M$ est
donc de présentation finie. Il est clair que $M$ n'est pas un $A$-module plat,
car étant de présentation finie, il serait libre (Bourbaki, \emph{Alg. comm.},
chap.~II, \S~3, n$^\circ$~2, cor.~2 de la prop.~5), ce qui est absurde puisque
$\mathfrak N\neq0$. Cependant,
\[
\mathfrak J^kM/\mathfrak J^{k+n}M\simeq A/\mathfrak J^n
\]
quels que soient $k$ et $n$ positifs, donc $M$ vérifie bien les conditions c)
et d) de $(\mathbf 0_{\mathrm{III}},10.2.1)$ puisque $A/\mathfrak J$ est un
corps.
\end{remark}
\begin{proposition}[11.3.13]
\label{IV.11.3.13-fr}
Soient $f:X\to Y$ un morphisme de préschémas, $x$ un point de $X$ tel que
$f$ soit plat au point $x$~; posons $y=f(x)$.
\begin{enumerate}[label=\emph{(\roman*)}]
\item Si $X$ est réduit (resp. normal) au point $x$, alors $Y$ est réduit
(resp. normal) au point $y$.
\item On suppose de plus que $f$ soit de présentation finie au point $x$.
Alors, si $Y$ est réduit (resp. normal) au point $y$ et si le morphisme $f$
est réduit (resp. normal) au point $x$ (6.8.1), $X$ est réduit (resp. normal)
au point $x$.
\end{enumerate}
\end{proposition}

\begin{proof}
La première assertion n'est mise que pour mémoire, ayant été démontrée dans
(2.1.13). Pour prouver (ii), on peut se borner au cas où $X=\operatorname{Spec}(B)$,
$Y=\operatorname{Spec}(A)$, où $A$ est un anneau local réduit (resp. intègre et
intégralement clos) et $B$ une $A$-algèbre de présentation finie. On peut alors
(8.9.1) écrire
\[
B=B_0\otimes_{A_0}A,
\]
où $A_0$ est une sous-$\mathbf Z$-algèbre de type fini de $A$ et $B_0$ une
$A_0$-algèbre de type fini. Soit $(C_\alpha)$ la famille filtrante croissante
des sous-$A_0$-algèbres de type fini de $A$, qui sont donc des
$\mathbf Z$-algèbres de type fini~; on a $A=\varinjlim C_\alpha$. Distinguons
maintenant les deux cas~:

\noindent 1\textsuperscript{o} Supposons $A$ réduit et $f$ réduit au point
$x$. Si $\mathfrak m$ est l'idéal maximal de $A$, soit
$\mathfrak p_\alpha$ l'idéal premier $\mathfrak m\cap C_\alpha$, et posons
$A'_\alpha=(C_\alpha)_{\mathfrak p_\alpha}$, de sorte que l'on a aussi
$A=\varinjlim A'_\alpha$ (5.13.3). Posons
$Y_\alpha=\operatorname{Spec}(A'_\alpha)$,
$X_\alpha=\operatorname{Spec}(B_0\otimes_{A_0}A'_\alpha)$ et soient
$f_\alpha:X_\alpha\to Y_\alpha$, $x_\alpha$ (resp. $y_\alpha$) la projection
de $x$ (resp. $y$) dans $X_\alpha$ (resp. $Y_\alpha$). Puisque $f$ est réduit
au point $x$, il existe $\alpha_0$ tel que $f_\alpha$ soit réduit au point
$x_\alpha$ pour $\alpha\geq\alpha_0$ ((6.7.8) et (11.2.6))~; comme
$Y_\alpha$ et $X_\alpha$ sont noethériens et que $A'_\alpha$ est réduit
(puisqu'il en est ainsi de $C_\alpha$), on déduit de (3.3.5) que $X_\alpha$
est réduit au point $x_\alpha$. Mais puisque
$B=\varinjlim(B_0\otimes_{A_0}A'_\alpha)$, on a aussi
$\mathcal O_{X,x}=\varinjlim\mathcal O_{X_\alpha,x_\alpha}$ (5.13.3) et par
suite $\mathcal O_{X,x}$ est réduit (5.13.2).

\noindent 2\textsuperscript{o} Supposons $A$ intégralement clos et $f$ normal
au point $x$. Comme $C_\alpha$ est universellement japonais (7.7.4), sa
clôture intégrale $C'_\alpha$ est une $C_\alpha$-algèbre finie, évidemment
contenue dans $A$. Soit $\mathfrak p'_\alpha$ l'idéal premier
$\mathfrak m\cap C'_\alpha$, et posons
$A''_\alpha=(C'_\alpha)_{\mathfrak p'_\alpha}$, de sorte que $A''_\alpha$
est un anneau local noethérien intègre et intégralement clos, et que l'on a
$A=\varinjlim A''_\alpha$ (5.13.3). Posons
$Y'_\alpha=\operatorname{Spec}(A''_\alpha)$,
$X'_\alpha=\operatorname{Spec}(B_0\otimes_{A_0}A''_\alpha)$ et soient
$f'_\alpha:X'_\alpha\to Y'_\alpha$, $x'_\alpha$ (resp. $y'_\alpha$) la
projection de $x$ (resp. $y$) dans $X'_\alpha$ (resp. $Y'_\alpha$). Puisque
$f$ est normal au point $x$, il existe $\alpha_0$ tel que, pour
$\alpha\geq\alpha_0$, $f'_\alpha$ soit normal au point $x'_\alpha$ ((6.7.8)
et (11.2.6))~;

\oldpage[IV-3]{141}

comme $X'_\alpha$ et $Y'_\alpha$ sont noethériens et que $A''_\alpha$ est
intégralement clos, on déduit de (6.5.4) que $X'_\alpha$ est normal au point
$x'_\alpha$. Mais les morphismes
$\operatorname{Spec}(A''_\beta)\to\operatorname{Spec}(A''_\alpha)$ pour
$\alpha\leq\beta$ sont dominants, donc, comme en vertu de (11.3.1) on peut
supposer que les $f'_\alpha$ sont plats pour $\alpha\geq\alpha_0$, il résulte
de (2.3.7) que toute composante irréductible de $X'_\beta$ domine une
composante irréductible de $X'_\alpha$. On conclut alors de (5.13.4), appliqué
aux préschémas
\[
\operatorname{Spec}(\mathcal O_{X_0,x_0}\otimes_{A_0}A''_\alpha),
\]
que $X$ est normal au point $x$.
\end{proof}

\begin{corollary}[11.3.14]
\label{IV.11.3.14-fr}
Soient $f:X\to Y$ un morphisme localement de présentation finie, $x$ un point
de $X$, $y=f(x)$. Si $Y$ est géométriquement unibranche au point $y$ (6.15.1)
et si le morphisme $f$ est normal au point $x$ (6.8.1), alors $X$ est
géométriquement unibranche au point $x$.
\end{corollary}

\begin{proof}
On peut évidemment se borner au cas où $Y=\operatorname{Spec}(A)$, $A$ étant
un anneau local intègre géométriquement unibranche, $y$ étant le point fermé de
$Y$. Soit $A'$ la clôture intégrale de $A$~; posons
$Y'=\operatorname{Spec}(A')$ et soit $y'$ le point fermé de $Y'$, de sorte que
le morphisme $g:Y'\to Y$ est radiciel au point $y$ (6.15.3), entier et
birationnel. Si l'on pose $X'=X\times_Y Y'$, et si $f':X'\to Y'$ et
$g':X'\to X$ sont les projections, $g'$ est entier et radiciel au point $x$
(6.15.3.1). Soit donc $x'$ l'unique point de $g'^{-1}(x)$, qui est au-dessus
de $y'$. Le morphisme $f'$ est de présentation finie et normal au point $x'$
(6.7.8), et par suite l'anneau local $\mathcal O_{X',x'}$ est intègre et
intégralement clos (11.3.13). D'autre part, on peut supposer $f$ plat (11.3.1),
donc $g'$ est un morphisme birationnel (6.15.4.1)~; par suite
$\operatorname{Spec}(\mathcal O_{X,x})$ est irréductible, et comme
$\mathcal O_{X,x}$ est réduit (11.3.13) il est intègre et géométriquement
unibranche en vertu de (6.15.5).
\end{proof}

\begin{proposition}[11.3.15]
\label{IV.11.3.15-fr}
Soient $A$ un anneau, $B$ une $A$-algèbre de présentation finie, $M$ un
$B$-module de présentation finie, qui est un $A$-module plat. Alors il existe
une suite finie $(f_i)_{1\leq i\leq n}$ d'éléments de $A$, tels que l'idéal
engendré par les $f_i$ soit égal à $A$, et que, pour $0\leq i<n$,
\[
M_{f_{i+1}}\Big/\Big(\sum_{j=1}^{i}f_jM_{f_{i+1}}\Big)
\]
soit un
$A_{f_{i+1}}\big/(\sum_{j=1}^{i}f_jA_{f_{i+1}})$-module libre.
\end{proposition}

On peut encore dire que les $D(f_i)$ forment un recouvrement ouvert de
$X=\operatorname{Spec}(A)$, et que si l'on pose $Z_1=D(f_1)$, puis, par
récurrence,
\[
Z_{i+1}=D(f_{i+1})\cap V(f_1A+\cdots+f_iA),
\]
les $Z_i$ forment une \emph{partition} de $X$ en ensembles localement fermés,
$A_{f_{i+1}}/(\sum_{j=1}^{i}f_jA_{f_{i+1}})$ étant l'anneau d'un sous-schéma
affine de $X$ ayant $Z_{i+1}$ pour espace sous-jacent.

\begin{proof}
Pour démontrer la proposition, on peut d'abord, en vertu de (11.2.7), supposer
qu'il existe un sous-anneau noethérien $A_0$ de $A$, une $A_0$-algèbre de type
fini $B_0$ et un $B_0$-module de type fini $M_0$, plat sur $A_0$ et tels que
$B=B_0\otimes_{A_0}A$, $M=M_0\otimes_{A_0}A$~; il est clair qu'il suffira de
prouver la proposition pour $A_0$, $B_0$ et $M_0$, car si les éléments
$f_i\in A_0$ vérifient dans ce cas les conditions de l'énoncé, ils les
vérifieront aussi pour $A$, $B$, $M$, puisque
\[
M_{f_{i+1}}\Big/\Big(\sum_{j=1}^{i}f_jM_{f_{i+1}}\Big)
=\left((M_0)_{f_{i+1}}\Big/\Big(\sum_{j=1}^{i}f_j(M_0)_{f_{i+1}}\Big)\right)
\otimes_{A_0}A
\]
(Bourbaki, \emph{Alg. comm.}, chap.~II, \S~2, n$^\circ$~7, prop.~18). On
peut donc se borner désormais au cas où $A$ est noethérien.

Notons maintenant que si $C$ est un anneau noethérien, $\mathfrak N$ son
nilradical et $P$ un $C$-module plat, alors il résulte de
$(\mathbf 0_{\mathrm{III}},10.1.2)$ que pour que $P$ soit un $C$-module libre,

\oldpage[IV-3]{142}

il faut et il suffit que $P\otimes_C(C/\mathfrak N)$ soit un
$(C/\mathfrak N)$-module libre. Notons d'autre part que si $C$ est un anneau
noethérien réduit, il existe $g\in C$ tel que $C_g$ soit un anneau intègre.
Utilisons maintenant le lemme (6.9.2)~: on peut définir par récurrence une suite
$(f_i)_{i\geq1}$ d'éléments de $A$ de la façon suivante~:
\begin{enumerate}
\item[$1^\circ$] $f_1$ est tel que $A_1=(A_{\mathrm{red}})_{f_1}$ soit
intègre et $M\otimes_A A_1$ un $A_1$-module libre~;
\item[$2^\circ$] si l'idéal $\mathfrak J_i$ engendré par $f_1,\ldots,f_i$
est $A$, $f_{i+1}=f_i$~;
\item[$3^\circ$] si au contraire $\mathfrak J_i\neq A$, $f_{i+1}$ est un
élément n'appartenant pas à $\mathfrak J_i$ tel que
$A_{i+1}=((A/\mathfrak J_i)_{\mathrm{red}})_{f_{i+1}}$ soit intègre et
$M\otimes_A A_{i+1}$ un $A_{i+1}$-module libre.
\end{enumerate}
Comme $A$ est noethérien, la suite croissante $(\mathfrak J_i)$ est
stationnaire, donc il existe $n$ tel que $f_1,\ldots,f_n$ engendrent l'idéal
$A$, et les $f_i$ répondent à la question puisque
\[
((A/\mathfrak J_i)_{\mathrm{red}})_{f_{i+1}}
=((A/\mathfrak J_i)_{f_{i+1}})_{\mathrm{red}}.
\]
\end{proof}

\begin{proposition}[11.3.16]
\label{IV.11.3.16-fr}
Soient $f:X\to Y$ un morphisme fidèlement plat de présentation finie,
$g:Y\to Z$ un morphisme tel que $g\circ f:X\to Z$ soit un morphisme de type
fini (resp. de présentation finie). Alors $g$ est un morphisme de type fini
(resp. de présentation finie).
\end{proposition}

\begin{proof}
Comme $f$ est surjectif et $g\circ f$ quasi-compact, $g$ est quasi-compact
(1.1.3). Montrons d'abord que si $g\circ f$ est de présentation finie, $g$ est
quasi-séparé. En effet, considérons un ouvert affine $W\subset Z$, et soit
$(V_i)$ un recouvrement affine fini de $g^{-1}(W)$~; il s'agit de voir que les
$V_i\cap V_j$ sont quasi-compacts (1.2.6 et 1.2.7). Pour chaque $i$, soit
$(U_{ih})$ un recouvrement ouvert affine fini de $f^{-1}(V_i)$~; comme $f$ est
de présentation finie, les $U_{ih}\cap U_{jk}$ sont tous quasi-compacts~; or,
puisque $f$ est surjectif, $V_i\cap V_j$ est réunion des images
$f(U_{ih}\cap U_{jk})$ pour $h$, $k$ variables, donc est quasi-compact puisque
$f$ est continue.

La question est donc locale sur $Y$ et on peut supposer que
$Y=\operatorname{Spec}(B)$ et $Z=\operatorname{Spec}(A)$ sont affines, $X$
étant réunion finie d'ouverts affines $X_i=\operatorname{Spec}(C_i)$, où les
$C_i$ sont des $B$-algèbres de type fini (resp. de présentation finie). Si
$X'$ est le préschéma somme des $X_i$, $p:X'\to X$ le morphisme qui coïncide
avec l'injection canonique dans chaque $X_i$, $p$ est un morphisme fidèlement
plat de présentation finie (1.6.5), donc $g\circ f\circ p$ est un morphisme de
type fini (resp. de présentation finie) et $f\circ p$ un morphisme fidèlement
plat de présentation finie. On peut par suite supposer aussi que
$X=\operatorname{Spec}(C)$ est affine, et on est donc ramené à prouver le
corollaire suivant.
\end{proof}

\begin{corollary}[11.3.17]
\label{IV.11.3.17-fr}
Soient $A$ un anneau, $B$ une $A$-algèbre, $C$ une $B$-algèbre de présentation
finie et qui soit un $B$-module fidèlement plat. Alors, si $C$ est une
$A$-algèbre de type fini (resp. de présentation finie), $B$ est une
$A$-algèbre de type fini (resp. de présentation finie).
\end{corollary}

\begin{proof}
Supposons d'abord que $g\circ f$ soit de type fini. Soit $(B_\lambda)$ la
famille filtrante croissante des sous-$A$-algèbres de type fini de $B$~; en
vertu de (8.8.2), il existe un indice $\alpha$ tel que
$C=C_\alpha\otimes_{B_\alpha}B$, où $C_\alpha$ est une
$B_\alpha$-algèbre de présentation finie~; en outre (8.10.5 et 11.2.6) on peut
supposer que $C_\alpha$ est un $B_\alpha$-module fidèlement plat. Pour
$\lambda\geq\alpha$, on a donc
$C=C_\lambda\otimes_{B_\lambda}B$, où $C_\lambda$ est un
$B_\lambda$-module fidèlement plat~; comme l'application $B_\lambda\to B$ est
injective, on en déduit qu'il en est de même de $C_\lambda\to C$. En outre,
comme $C=\varinjlim C_\lambda$, tout élément de $C$ appartient à un
$C_\lambda$, et par suite il existe $\lambda$ tel que l'application
$C_\lambda\to C$

\oldpage[IV-3]{143}

soit bijective, puisque $C$ est une $B$-algèbre de type fini. Mais alors
l'application $B_\lambda\to B$ est bijective par fidèle platitude, donc
$B=B_\lambda$ est une $A$-algèbre de type fini.

Supposons maintenant que $g\circ f$ soit de présentation finie, $C$ étant donc
une $A$-algèbre de présentation finie. D'après la première partie du
raisonnement, on a
\[
B=A[T_1,\ldots,T_n]/\widetilde{\mathfrak J}
\]
pour un idéal $\widetilde{\mathfrak J}$~; soit
$(\widetilde{\mathfrak J}_\lambda)$ la famille filtrante des idéaux de type
fini de $A[T_1,\ldots,T_n]$ contenus dans $\widetilde{\mathfrak J}$, de sorte
que $\widetilde{\mathfrak J}=\varinjlim\widetilde{\mathfrak J}_\lambda$ et
$B=\varinjlim B_\lambda$, avec
$B_\lambda=A[T_1,\ldots,T_n]/\widetilde{\mathfrak J}_\lambda$. Appliquant
comme ci-dessus (8.8.2), (8.10.5) et (11.2.6), on a
$C=C_\alpha\otimes_{B_\alpha}B$, où $C_\alpha$ est une
$B_\alpha$-algèbre de présentation finie et un $B_\alpha$-module fidèlement
plat~; on posera encore $C_\lambda=C_\alpha\otimes_{B_\alpha}B_\lambda$ pour
$\lambda\geq\alpha$, de sorte que
\[
C_\lambda=C_\alpha/\big(C_\alpha\otimes_{B_\alpha}
(\widetilde{\mathfrak J}_\lambda/\widetilde{\mathfrak J}_\alpha)\big)
\]
par platitude, et de même
\[
C=C_\alpha/\big(C_\alpha\otimes_{B_\alpha}
(\widetilde{\mathfrak J}/\widetilde{\mathfrak J}_\alpha)\big).
\]
Comme par hypothèse $C$ est une $A$-algèbre de présentation finie ainsi que
$C_\alpha$, $C_\alpha\otimes_{B_\alpha}
(\widetilde{\mathfrak J}/\widetilde{\mathfrak J}_\alpha)$ est un idéal de
type fini de $C_\alpha$ (1.4.4), et les
$C_\alpha\otimes_{B_\alpha}
(\widetilde{\mathfrak J}_\lambda/\widetilde{\mathfrak J}_\alpha)$
s'identifient à des idéaux de type fini de $C_\alpha$ dont
$C_\alpha\otimes_{B_\alpha}
(\widetilde{\mathfrak J}/\widetilde{\mathfrak J}_\alpha)$ est la réunion. Il
existe par suite un $\lambda\geq\alpha$ tel que
\[
C_\alpha\otimes_{B_\alpha}
(\widetilde{\mathfrak J}/\widetilde{\mathfrak J}_\alpha)
=C_\alpha\otimes_{B_\alpha}
(\widetilde{\mathfrak J}_\lambda/\widetilde{\mathfrak J}_\alpha),
\]
d'où $\widetilde{\mathfrak J}=\widetilde{\mathfrak J}_\lambda$ par fidèle
platitude, ce qui prouve que $B$ est une $A$-algèbre de présentation finie.
\end{proof}

\subsection{Descente de la platitude par des morphismes quelconques~: cas d'un préschéma de base artinien}
\label{subsection:IV.11.4-fr}

\begin{theorem}[11.4.1]
\label{IV.11.4.1-fr}
Soient $A$ un anneau, $\mathfrak J$ un idéal nilpotent de $A$,
$u_\lambda:A\to B_\lambda$ $(\lambda\in L)$ une famille d'homomorphismes
d'anneaux telle que l'intersection des noyaux des $u_\lambda$ soit réduite à
$0$. Soit $M$ un $A$-module tel que pour tout $\lambda\in L$,
$M\otimes_A B_\lambda$ soit un $B_\lambda$-module libre et que
$M\otimes_A(A/\mathfrak J)$ soit un $(A/\mathfrak J)$-module libre. Alors $M$
est un $A$-module libre. Si de plus l'ensemble d'indices $L$ est fini, on peut
remplacer partout «~libre~» par «~plat~» dans l'énoncé précédent.
\end{theorem}

\begin{proof}
Dans les deux cas il suffira de prouver que $M$ est un $A$-module \emph{plat}~:
en effet, lorsque $M\otimes_A(A/\mathfrak J)$ est un
$(A/\mathfrak J)$-module libre, il en résultera que $M$ est un $A$-module libre
en vertu de $(\mathbf 0_{\mathrm{III}},10.1.2)$.

Nous utiliserons le lemme suivant, qui généralise
$(\mathbf 0_{\mathrm{III}},10.2.1)$~:

\begin{lemma}[11.4.1.1]
\label{IV.11.4.1.1-fr}
Soit $A$ un anneau muni d'une filtration finie
$(\mathfrak J_n)_{0\leq n\leq N+1}$, avec $\mathfrak J_0=A$,
$\mathfrak J_{N+1}=0$. Soit $M$ un $A$-module muni de la filtration
$(\mathfrak J_nM)_{0\leq n\leq N+1}$, et désignons par
$\operatorname{gr}(A)$ et $\operatorname{gr}(M)$ l'anneau et le module gradués
correspondants. Supposons que $M\otimes_A(A/\mathfrak J_1)$ soit un
$(A/\mathfrak J_1)$-module plat et que l'homomorphisme canonique
\begin{equation}
\operatorname{gr}_0(M)\otimes_{\operatorname{gr}_0(A)}\operatorname{gr}(A)
\longrightarrow\operatorname{gr}(M)
\label{IV.11.4.1.2-fr}
\tag{11.4.1.2}
\end{equation}
soit injectif. Alors $M$ est un $A$-module plat.
\end{lemma}

L'homomorphisme canonique (11.4.1.2) se définit de la même façon que celui de
$(\mathbf 0_{\mathrm{III}},10.1.1)$ comme étant en degré $n$ l'homomorphisme
\[
(M/\mathfrak J_1M)\otimes(\mathfrak J_n/\mathfrak J_{n+1})
=(M\otimes\mathfrak J_n)/
(\operatorname{Im}(M\otimes\mathfrak J_{n+1})
+\operatorname{Im}(\mathfrak J_1M\otimes\mathfrak J_n))
\longrightarrow\mathfrak J_nM/\mathfrak J_{n+1}M
\]
provenant de l'homomorphisme canonique $M\otimes\mathfrak J_n\to\mathfrak J_nM$
par passage aux quotients. Le lemme se démontre par récurrence sur $N$,
puisqu'il n'y a rien à démontrer pour $N=0$.

\oldpage[IV-3]{144}

Les conditions sur $M$ entraînent, en vertu de l'hypothèse de récurrence, que
$M\otimes_A(A/\mathfrak J_N)$ est un $(A/\mathfrak J_N)$-module plat. Notons
maintenant que l'on a $\mathfrak J_N^2\subset\mathfrak J_{N+1}=0$ si $N\geq1$,
et
\[
(M/\mathfrak J_1M)\otimes(\mathfrak J_N/\mathfrak J_{N+1})
=(M/\mathfrak J_NM)\otimes(\mathfrak J_N/\mathfrak J_{N+1})
=(M/\mathfrak J_NM)\otimes\mathfrak J_N\,;
\]
donc l'homomorphisme canonique
\[
(M/\mathfrak J_NM)\otimes\mathfrak J_N
\longrightarrow\mathfrak J_NM=\mathfrak J_NM/\mathfrak J_N^2M
\]
est injectif. Appliquant $(\mathbf 0_{\mathrm{III}},10.2.1)$ à la filtration
$\mathfrak J_N$-préadique, on en conclut bien que $M$ est un $A$-module plat.

Pour appliquer ce lemme à (11.4.1), nous désignerons par $\mathfrak J_n$
l'idéal de $A$ intersection des images réciproques
$u_\lambda^{-1}(\mathfrak J^nB_\lambda)$~; il est immédiat que l'on a
$\mathfrak J_0=A$, $\mathfrak J_m\mathfrak J_n\subset\mathfrak J_{m+n}$ pour
$m\geq0$, $n\geq0$~; en outre, si $\mathfrak J^{N+1}=0$, on a aussi
$\mathfrak J_{N+1}=0$ puisque l'intersection des noyaux des $u_\lambda$ est
réduite à $0$.

Munissons $A$ de la filtration $(\mathfrak J_n)$, $M$ de la filtration
$(\mathfrak J_nM)$, et, pour chaque $\lambda$, munissons $B_\lambda$ et
$N_\lambda=M\otimes_A B_\lambda$ des filtrations $\mathfrak J$-préadiques~;
considérons pour chaque $\lambda$ le diagramme commutatif
\[
\xymatrix{
\operatorname{gr}^0(M)\otimes_{\operatorname{gr}^0(A)}\operatorname{gr}(A)
 \ar[r]^{\varphi} \ar[d]_{f_\lambda} & \operatorname{gr}(M) \ar[d]^{g_\lambda} \\
\operatorname{gr}_{\mathfrak J}^0(N_\lambda)
 \otimes_{\operatorname{gr}_{\mathfrak J}^0(B_\lambda)}
 \operatorname{gr}_{\mathfrak J}(B_\lambda)
 \ar[r]_{\varphi_\lambda} & \operatorname{gr}_{\mathfrak J}(N_\lambda)
}
\]
où les flèches horizontales sont les homomorphismes canoniques (11.4.1.2) et
les flèches verticales sont déduites des homomorphismes canoniques
$A\to B_\lambda$ et $M\to N_\lambda$. L'hypothèse que $N_\lambda$ est un
$B_\lambda$-module plat entraîne que $\varphi_\lambda$ est injective
$(\mathbf 0_{\mathrm{III}},10.2.1)$, donc
$\operatorname{Ker}(\varphi)\subset\operatorname{Ker}(f_\lambda)$. Posant
$A_0=A/\mathfrak J_1$, $M_0=M\otimes_A A_0$, notons que
\[
\operatorname{gr}_{\mathfrak J}^0(N_\lambda)
=N_\lambda/\mathfrak JN_\lambda
=(M/\mathfrak JM)\otimes_A B_\lambda
=(M/\mathfrak J_1M)\otimes_A(B_\lambda/\mathfrak JB_\lambda),
\]
car ce dernier produit tensoriel est égal à
\[
(M\otimes_A B_\lambda)\Big/
\big(\operatorname{Im}(M\otimes_A\mathfrak JB_\lambda)
 +\operatorname{Im}(\mathfrak J_1M\otimes_A B_\lambda)\big),
\]
et l'on a
$\operatorname{Im}(\mathfrak J_1M\otimes_A B_\lambda)
=\operatorname{Im}(M\otimes_A\mathfrak J_1B_\lambda)
\subset\operatorname{Im}(M\otimes_A\mathfrak JB_\lambda)$ puisque
$\mathfrak J_1B_\lambda\subset\mathfrak JB_\lambda$ par définition~; enfin,
la relation $\mathfrak J_1B_\lambda\subset\mathfrak JB_\lambda$ montre que l'on
a aussi
\[
(M/\mathfrak J_1M)\otimes_A(B_\lambda/\mathfrak JB_\lambda)
=(M/\mathfrak J_1M)\otimes_{A/\mathfrak J_1}
 (B_\lambda/\mathfrak JB_\lambda),
\]
si bien que finalement on a
\[
\operatorname{gr}_{\mathfrak J}^0(N_\lambda)
=M_0\otimes_{A_0}(B_\lambda/\mathfrak JB_\lambda)
\]
et par suite
\[
\operatorname{gr}_{\mathfrak J}^0(N_\lambda)
\otimes_{\operatorname{gr}_{\mathfrak J}^0(B_\lambda)}
\operatorname{gr}_{\mathfrak J}(B_\lambda)
=M_0\otimes_{A_0}\operatorname{gr}_{\mathfrak J}(B_\lambda).
\]
L'homomorphisme $f_\lambda$ peut donc s'écrire
\[
1\otimes\operatorname{gr}(u_\lambda):
M_0\otimes_{A_0}\operatorname{gr}(A)
\longrightarrow M_0\otimes_{A_0}\operatorname{gr}_{\mathfrak J}(B_\lambda),
\]

\oldpage[IV-3]{145}

et comme $M_0$ est par hypothèse un $A_0$-module \emph{plat}, le noyau de
$f_\lambda$ est égal à $M_0\otimes_{A_0}R_\lambda$, où $R_\lambda$ est le
noyau de
$\operatorname{gr}(u_\lambda):\operatorname{gr}(A)
\to\operatorname{gr}_{\mathfrak J}(B_\lambda)$. Tout revient donc à prouver
que
\[
\bigcap_{\lambda\in L}(M_0\otimes_{A_0}R_\lambda)=0.
\]
Or, par définition des $\mathfrak J_n$, l'intersection des noyaux des
homomorphismes
\[
\mathfrak J_nA/\mathfrak J_{n+1}A
\longrightarrow\mathfrak J^nB_\lambda/\mathfrak J^{n+1}B_\lambda,
\]
lorsque $\lambda$ parcourt $L$, est réduite à $0$, autrement dit
$\bigcap_{\lambda\in L}R_\lambda=0$. Cela étant, supposons d'abord que $M_0$
soit un $A_0$-module libre~; prenant une base de $M_0$, on voit aussitôt que
l'on a
\[
M_0\otimes_{A_0}\Big(\bigcap_{\lambda\in L}R_\lambda\Big)
=\bigcap_{\lambda\in L}(M_0\otimes_{A_0}R_\lambda),
\]
d'où la proposition dans ce cas. Lorsque $L$ est fini, la formule précédente
est encore vraie sous la seule hypothèse que $M_0$ est un $A_0$-module plat
$(\mathbf 0_{\mathrm I},6.1.3)$, ce qui termine la démonstration.
\end{proof}

\begin{remark}[11.4.2]
\label{IV.11.4.2-fr}
La conclusion de (11.4.1) peut être inexacte si $L$ est infini et si l'on
suppose seulement que $M\otimes_A(A/\mathfrak J)$ est un
$(A/\mathfrak J)$-module plat. Par exemple, soient $V$ un anneau de valuation
discrète, $K$ son corps des fractions, et soit $A=V[T]/(T^2)$ ($T$
indéterminée)~; prenons pour $\mathfrak J$ l'image de $(T)$ dans $A$, de sorte
que $A/\mathfrak J=V$, et prenons $M=K$, qui est un
$(A/\mathfrak J)$-module, donc égal à $M\otimes_A(A/\mathfrak J)$~; en outre
$M$ est un $(A/\mathfrak J)$-module plat, mais non un $A$-module plat, car
puisque $\mathfrak J$ est nilpotent, il résulterait de
$(\mathbf 0_{\mathrm{III}},10.1.2)$ que $M$ serait un $A$-module libre, ce qui
est absurde puisque $\mathfrak JK=0$. Considérons d'autre part l'idéal maximal
$\mathfrak m$ de l'anneau local noethérien $A$~; on a $\mathfrak mK=K$, donc
$M\otimes_A(A/\mathfrak m^n)=0$ quel que soit l'entier $n$~; les
$(A/\mathfrak m^n)$-modules $M\otimes_A(A/\mathfrak m^n)$ sont donc plats pour
tout $n$, et l'intersection des $\mathfrak m^n$ est réduite à $0$.
\end{remark}

\begin{corollary}[11.4.3]
\label{IV.11.4.3-fr}
Soient $A$ un anneau semi-local dont le radical $\mathfrak J$ est nilpotent
(par exemple un anneau artinien), $u_\lambda:A\to B_\lambda$ $(\lambda\in L)$
une famille d'homomorphismes d'anneaux telle que l'intersection des noyaux des
$u_\lambda$ soit réduite à $0$. Pour qu'un $A$-module $M$ soit plat, il faut et
il suffit que pour tout $\lambda\in L$, $M\otimes_A B_\lambda$ soit un
$B_\lambda$-module plat.
\end{corollary}

\begin{proof}
Comme $A/\mathfrak J$ est composé direct d'un nombre fini de corps (Bourbaki,
\emph{Alg. comm.}, chap.~II, \S~3, n$^\circ$~5, prop.~16) et que $\mathfrak J$
est nilpotent, $A$ est composé direct d'un nombre fini d'anneaux locaux $A_i$
dont le radical est nilpotent (\emph{loc. cit.}, \S~4, n$^\circ$~3, cor.~1 de
la prop.~15) et $M$ est par suite somme directe de $A_i$-modules $M_i$, chaque
$M_i$ étant annulé par les $A_j$ d'indice $j\neq i$~; pour que $M$ soit un
$A$-module plat, il faut et il suffit que chacun des $M_i$ soit un $A_i$-module
plat~; d'ailleurs l'intersection des noyaux des homomorphismes
\[
A_i\longrightarrow A\xrightarrow{u_\lambda}B_\lambda
\]
est réduite à $0$, et $M_i\otimes_{A_i}B_\lambda$ est facteur direct de
$M\otimes_A B_\lambda$. On peut donc se borner au cas où $A$ est en outre
local. Alors $A/\mathfrak J$ est un corps, donc
$M\otimes_A(A/\mathfrak J)$ est un $(A/\mathfrak J)$-module libre, et il
suffit de voir que pour tout $\lambda$, $M\otimes_A B_\lambda$ est un
$B_\lambda$-module libre, en vertu de (11.4.1). Mais si l'on pose
$\mathfrak J_\lambda=\mathfrak JB_\lambda$,
\[
(M\otimes_A B_\lambda)\otimes_{B_\lambda}(B_\lambda/\mathfrak J_\lambda)
=(M\otimes_A(A/\mathfrak J))\otimes_{A/\mathfrak J}
 (B_\lambda/\mathfrak J_\lambda)
\]
est un $(B_\lambda/\mathfrak J_\lambda)$-module libre, et
$\mathfrak J_\lambda$ est nilpotent. La conclusion résulte donc de l'hypothèse
que $M\otimes_A B_\lambda$ est un $B_\lambda$-module plat et de
$(\mathbf 0_{\mathrm{III}},10.1.2)$.
\end{proof}

\begin{corollary}[11.4.4]
\label{IV.11.4.4-fr}
Soient $A$ un anneau, $M$ un $A$-module~; on suppose qu'il existe un idéal
nilpotent $\mathfrak J$ de $A$ tel que $M\otimes_A(A/\mathfrak J)$ soit un
$(A/\mathfrak J)$-module libre. Alors l'ensemble des idéaux $\mathfrak K$ de
$A$ tels que $M\otimes_A(A/\mathfrak K)$ soit un $(A/\mathfrak K)$-module
libre admet un plus petit élément $\mathfrak K_0$ (qui est aussi le plus petit
des idéaux $\mathfrak K$ tels que $M\otimes_A(A/\mathfrak K)$ soit un
$(A/\mathfrak K)$-module plat).
\end{corollary}

\oldpage[IV-3]{146}

\begin{proof}
Pour qu'un homomorphisme $u:A\to A'$ soit tel que $M\otimes_A A'$ soit un
$A'$-module libre (resp. un $A'$-module plat), il faut et il suffit que $u$ se
factorise en $A\to A/\mathfrak K_0\to A'$ (ou encore que
$\mathfrak K_0A'=0$).

Le fait que l'intersection $\mathfrak K_0$ des idéaux $\mathfrak K$ pour
lesquels $M\otimes_A(A/\mathfrak K)$ est un $(A/\mathfrak K)$-module libre
soit le plus petit de ces idéaux résulte de (11.4.1) appliqué à l'anneau
$A/\mathfrak K_0$, à son idéal nilpotent $\mathfrak J/\mathfrak K_0$ et aux
homomorphismes $A/\mathfrak K_0\to A/\mathfrak K$, dont les noyaux ont $0$
pour intersection. Si $A'$ est une $(A/\mathfrak K_0)$-algèbre, on a
\[
M\otimes_A A'=
(M\otimes_A(A/\mathfrak K_0))\otimes_{A/\mathfrak K_0}A',
\]
donc $M\otimes_A A'$ est un $A'$-module libre. Réciproquement, si $A'$ est
une $A$-algèbre telle que $M\otimes_A A'$ soit un $A'$-module libre et si
$\mathfrak K$ est le noyau de l'homomorphisme $A\to A'$, il résulte de
(11.4.1) appliqué à l'anneau $A/\mathfrak K$, au
$(A/\mathfrak K)$-module $M\otimes_A(A/\mathfrak K)$, à l'idéal nilpotent
$(\mathfrak J+\mathfrak K)/\mathfrak K$ et à l'homomorphisme injectif
$A/\mathfrak K\to A'$, que $M\otimes_A(A/\mathfrak K)$ est un
$(A/\mathfrak K)$-module libre, donc que $\mathfrak K\supset\mathfrak K_0$.
Le fait que l'on puisse remplacer «~libre~» par «~plat~» dans ce qui précède
(en conservant naturellement l'hypothèse que $M\otimes_A(A/\mathfrak J)$ est
un $(A/\mathfrak J)$-module libre) résulte de
$(\mathbf 0_{\mathrm{III}},10.1.2)$, comme on l'a vu au début de la
démonstration de (11.4.1).
\end{proof}

\begin{proposition}[11.4.5]
\label{IV.11.4.5-fr}
Soient $Y$ un préschéma irréductible, $f:X\to Y$ un morphisme de présentation
finie, $\mathcal F$ un $\mathcal O_X$-Module de présentation finie. Il existe
alors un ouvert non vide $U$ dans $Y$ et un sous-schéma fermé $Z$ de $U$, de
présentation finie sur $U$, ayant la propriété suivante~: pour tout changement
de base $U'\to U$, si l'on pose $X'=X\times_YU'$, $f'=f_{(U')}$ et
$\mathcal F'=\mathcal F\otimes_{\mathcal O_Y}\mathcal O_{U'}$, alors, pour que
$\mathcal F'$ soit $f'$-plat, il faut et il suffit que le morphisme $U'\to U$
se factorise en $U'\to Z\to U$. Un tel schéma $Z$ a même espace sous-jacent que
$U$. Supposons de plus $Y$ affine, et soit $(W_i)$ un recouvrement fini de $X$
par des ouverts affines~; alors, on peut supposer $U$ choisi de sorte que, si
$U'=\operatorname{Spec}(A')$ est un schéma affine et $U'\to U$ un morphisme qui
se factorise en $U'\to Z\to U$, chacun des
$\Gamma(W_i\times_YU',\mathcal F')$ est un $A'$-module libre.
\end{proposition}

\begin{proof}
On peut évidemment se borner au cas où $Y=\operatorname{Spec}(A)$ est affine.
Utilisant (8.9.1), il y a un sous-anneau noethérien $A_1$ de $A$, un morphisme
de type fini $f_1:X_1\to Y_1=\operatorname{Spec}(A_1)$ et un
$\mathcal O_{X_1}$-Module cohérent $\mathcal F_1$ tels que $f=(f_1)_{(Y)}$ et
$\mathcal F=\mathcal F_1\otimes_{\mathcal O_{Y_1}}\mathcal O_Y$~; on peut en
outre supposer que les $W_i$ sont images réciproques d'ouverts affines de
$X_1$. Notons en outre que $Y_1$ est irréductible, le morphisme $Y\to Y_1$
étant dominant $(\mathbf I,1.2.7)$. Supposons la proposition démontrée pour
$Y_1$, $X_1$ et $\mathcal F_1$, et soient $U_1$ l'ouvert de $Y_1$ et $Z_1$ le
sous-schéma fermé de $U_1$ ayant les propriétés voulues, $U=U_1\times_{Y_1}Y$
et $Z=Z_1\times_{Y_1}Y$ leurs images réciproques. Alors, si $U'\to U$ est un
changement de base tel que
$\mathcal F'=\mathcal F\otimes_{\mathcal O_Y}\mathcal O_{U'}
=\mathcal F_1\otimes_{\mathcal O_{Y_1}}\mathcal O_{U'}$ soit $f'$-plat, le
morphisme $U'\to U_1$ se factorise en $U'\to Z_1\to U_1$~; mais comme
$U'\to U_1$ se factorise aussi en $U'\to U\to U_1$, la définition du produit
de préschémas montre que $U'\to U$ se factorise en $U'\to Z\to U$.

On peut donc se borner au cas où $A$ est \emph{noethérien}~; soit
$\mathfrak N$ son nilradical, qui est nilpotent, et posons
$A_0=A_{\mathrm{red}}=A/\mathfrak N$,
$Y_0=\operatorname{Spec}(A_0)=Y_{\mathrm{red}}$,
$X_0=X\times_YY_0$, $\mathcal F_0=\mathcal F\otimes_{\mathcal O_Y}
\mathcal O_{Y_0}$, $W_{i0}=W_i\times_YY_0$~; si
$M_i=\Gamma(W_i,\mathcal F)$, $M_{i0}=\Gamma(W_{i0},\mathcal F_0)$ est donc
égal à $M_i\otimes_A A_0$. Comme $A_0$ est \emph{intègre}, on peut, en vertu du
théorème de platitude générique (6.9.2), en remplaçant au besoin $Y$ par un
ouvert non vide de $Y$, supposer que les $M_{i0}$ soient des $A_0$-modules
libres. En vertu de (11.4.4), il y a donc pour chaque $i$ un idéal
$\mathfrak J_i\subset\mathfrak N$ tel

\oldpage[IV-3]{147}

que les $A$-algèbres $A'$ pour lesquelles les $M_i\otimes_A A'$ sont des
$A'$-modules libres (ou plats) sont exactement celles pour lesquelles
$\mathfrak J_iA'=0$. Soit $\mathfrak J=\sum_i\mathfrak J_i$, qui est un idéal
contenu dans $\mathfrak N$~; dire que $\mathcal F\otimes_A A'$ est $A'$-plat
équivaut à dire que les $M_i\otimes_A A'$ sont tous des $A'$-modules plats,
donc que $\mathfrak JA'=0$. Il en résulte que si l'on prend
$Z=V(\mathfrak J)$, on répond à la question, car pour qu'un morphisme
$U'\to U$ ait la propriété de l'énoncé, il faut et il suffit évidemment que
pour tout ouvert affine $V'\subset U'$, le morphisme $V'\to U$ ait la même
propriété.
\end{proof}

\begin{corollary}[11.4.6]
\label{IV.11.4.6-fr}
Soient $A$ un anneau tel que $\operatorname{Spec}(A)$ soit irréductible,
$\mathfrak p$ son unique idéal premier minimal, $B$ une $A$-algèbre de
présentation finie, $M$ un $B$-module de présentation finie. Supposons que
$M_\mathfrak p$ soit un $A_\mathfrak p$-module plat~; alors il existe
$t\in A-\mathfrak p$ tel que $M_t$ soit un $A_t$-module libre.
\end{corollary}

\begin{proof}
Appliquant (11.4.5) à $Y=\operatorname{Spec}(A)$, $X=\operatorname{Spec}(B)$,
$\mathcal F=\widetilde M$, on peut (en remplaçant au besoin $A$ par $A_h$, où
$h\in A-\mathfrak p$) supposer qu'il existe un idéal de type fini
$\mathfrak J$ dans $A$ tel que les $A$-algèbres $A'$ pour lesquelles
$M\otimes_A A'$ soit un $A'$-module libre (ou plat) sont exactement celles
telles que $\mathfrak JA'=0$. En particulier, l'hypothèse que
$M_\mathfrak p$ soit un $A_\mathfrak p$-module plat entraîne
$\mathfrak JA_\mathfrak p=0$, ou encore $\mathfrak J_\mathfrak p=0$. Mais
comme $\mathfrak J$ est de type fini, il existe $t\in A-\mathfrak p$ tel que
$\mathfrak J_t=0$, ou encore $\mathfrak JA_t=0$, et par suite $M_t$ est un
$A_t$-module libre.
\end{proof}

\begin{proposition}[11.4.7]
\label{IV.11.4.7-fr}
Soient $A$ un anneau noethérien, $\mathfrak J$ un idéal de $A$, $M$ un
$A$-module idéalement séparé $(\mathbf 0_{\mathrm{III}},10.2.1)$ pour la
topologie $\mathfrak J$-préadique. Soit $(\mathfrak p_i)_{1\leq i\leq r}$ une
famille finie d'idéaux premiers de $A$ contenant $\mathfrak J$~; pour tout
entier $n\geq0$, soit $\mathfrak p_i^{(n)}$ la puissance symbolique $n$-ième de
$\mathfrak p_i$ (noyau de l'homomorphisme
$A\to A_{\mathfrak p_i}/\mathfrak p_i^nA_{\mathfrak p_i}$)~; posons
$\mathfrak J_n=\bigcap_{i=1}^{r}\mathfrak p_i^{(n)}$, de sorte que
$\mathfrak J_n\supset\mathfrak J^n$, et supposons que la topologie définie par
la filtration $(\mathfrak J_n)$ soit identique à la topologie
$\mathfrak J$-préadique (autrement dit que, pour tout $n$, il existe $m$ tel
que $\mathfrak J^n\supset\mathfrak J_m$). Pour que $M$ soit un $A$-module
plat, il faut et il suffit que $M\otimes_A(A/\mathfrak J)$ soit un
$(A/\mathfrak J)$-module plat et que, pour tout $i$,
$M\otimes_AA_{\mathfrak p_i}=M_{\mathfrak p_i}$ soit un
$A_{\mathfrak p_i}$-module plat.
\end{proposition}

\begin{proof}
Comme $M$ est idéalement séparé, il suffit, en vertu de
$(\mathbf 0_{\mathrm{III}},10.2.1)$, de montrer que, pour tout $n\geq0$,
$M\otimes_A(A/\mathfrak J^n)$ est un $(A/\mathfrak J^n)$-module plat~;
puisque tout $\mathfrak J^n$ contient un $\mathfrak J_m$, il revient au même
de prouver que pour tout $n$, $M\otimes_A(A/\mathfrak J_n)$ est un
$(A/\mathfrak J_n)$-module plat. Or, comme
$\mathfrak J_1\supset\mathfrak J$, $M\otimes_A(A/\mathfrak J_1)$ est un
$(A/\mathfrak J_1)$-module plat~; dans l'anneau $A/\mathfrak J_n$, l'idéal
$\mathfrak J_1/\mathfrak J_n$ est nilpotent et enfin l'intersection des noyaux
des homomorphismes
\[
A/\mathfrak J_n\longrightarrow
A_{\mathfrak p_i}/\mathfrak p_i^nA_{\mathfrak p_i}
\]
est nulle dans $A/\mathfrak J_n$, par définition de $\mathfrak J_n$. Il suffit
donc, par (11.4.1), de vérifier que
\[
M\otimes_A(A_{\mathfrak p_i}/\mathfrak p_i^nA_{\mathfrak p_i})
\]
est un $(A_{\mathfrak p_i}/\mathfrak p_i^nA_{\mathfrak p_i})$-module plat, ce
qui résulte de l'hypothèse que $M\otimes_AA_{\mathfrak p_i}$ est un
$A_{\mathfrak p_i}$-module plat.
\end{proof}

\begin{remark}[11.4.8]
\label{IV.11.4.8-fr}
L'hypothèse faite dans (11.4.7) sur la topologie définie par les
$\mathfrak J_n$ est vérifiée si, pour tout $n$ assez grand,
$\operatorname{Ass}(A/\mathfrak J^n)$ est contenu dans l'ensemble des
$\mathfrak p_i$. En effet, $\mathfrak J^n$ est alors intersection d'idéaux
primaires pour les $\mathfrak p_i$, dont chacun contient une puissance
symbolique de $\mathfrak p_i$, d'où la conclusion. En particulier~:
\end{remark}

\begin{corollary}[11.4.9]
\label{IV.11.4.9-fr}
Soient $A$ un anneau noethérien, $\mathfrak J$ un idéal nilpotent de $A$, $M$
un $A$-module. Pour que $M$ soit un $A$-module plat (resp. libre), il faut et
il suffit que $M\otimes_A(A/\mathfrak J)$ soit un $(A/\mathfrak J)$-module
plat (resp. libre) et que pour tout idéal premier
$\mathfrak p\in\operatorname{Ass}(A)$, $M_\mathfrak p$ soit un
$A_\mathfrak p$-module plat.
\end{corollary}
% IV3_B22_APPEND_ANCHOR

\oldpage[IV-3]{148}

L'assertion relative au cas où $M\otimes_A(A/\mathfrak J)$ est libre se déduit
encore de l'assertion relative au cas où $M\otimes_A(A/\mathfrak J)$ est plat
par $(\mathbf 0_{\mathrm{III}},10.1.2)$.

\begin{corollary}[11.4.10]
\label{IV.11.4.10-fr}
Soient $A$ un anneau noethérien, $\mathfrak J$ un idéal de $A$, $M$ un
$A$-module. On suppose que $M$ est idéalement séparé pour la topologie
$\mathfrak J$-préadique et que $\operatorname{gr}_{\mathfrak J}(A)$ est un
$(A/\mathfrak J)$-module plat. Pour que $M$ soit un $A$-module plat, il faut et
il suffit que $M\otimes_A(A/\mathfrak J)$ soit un $(A/\mathfrak J)$-module plat
et que, pour tout $\mathfrak p\in\operatorname{Ass}(A/\mathfrak J)$,
$M_{\mathfrak p}$ soit un $A_{\mathfrak p}$-module plat.
\end{corollary}

\begin{proof}
Compte tenu de (11.4.8), tout revient à montrer que
$\operatorname{Ass}(A/\mathfrak J^n)$ est contenu dans
$\operatorname{Ass}(A/\mathfrak J)$ pour tout $n$. Or, si $a\in A$ n'appartient
à aucun des $\mathfrak p\in\operatorname{Ass}(A/\mathfrak J)$, l'homothétie de
rapport $a$ est injective dans $A/\mathfrak J$~; comme chacun des
$\mathfrak J^k/\mathfrak J^{k+1}$ est un $(A/\mathfrak J)$-module plat, $a$ est
aussi un élément $(\mathfrak J^k/\mathfrak J^{k+1})$-régulier, donc $a$ est
$(A/\mathfrak J^n)$-régulier pour tout $n$ (Bourbaki, \emph{Alg. comm.},
chap.~III, \S~2, n$^\circ$~8, cor.~1 du th.~1), et par suite n'appartient à
aucun idéal premier associé à $A/\mathfrak J^n$, d'où le corollaire (Bourbaki,
\emph{Alg. comm.}, chap.~II, \S~1, n$^\circ$~1, prop.~2).
\end{proof}

\begin{proposition}[11.4.11]
\label{IV.11.4.11-fr}
Soient $A$ un anneau local artinien, d'idéal maximal $\mathfrak m$,
$Y=\operatorname{Spec}(A)$, $y$ l'unique point de $Y$, $f:X\to Y$ un morphisme
localement de type fini, $\mathcal F$ un $\mathcal O_X$-Module cohérent. Soient
$(A'_\alpha)$ une famille d'anneaux locaux, et pour chaque $\alpha$,
$u_\alpha:A\to A'_\alpha$ un homomorphisme d'anneaux (nécessairement local).
Posons $Y'_\alpha=\operatorname{Spec}(A'_\alpha)$,
$X'_\alpha=X\times_Y Y'_\alpha$,
$\mathcal F'_\alpha=\mathcal F\otimes_A A'_\alpha$. Soit $x$ un point de $X$,
et supposons vérifiées les conditions suivantes~:
\begin{enumerate}[label=(\roman*)]
  \item L'intersection des noyaux des $u_\alpha$ est réduite à $0$.
  \item L'extension $k(x)$ du corps résiduel $k$ de $A$ est primaire (4.3.1)
  (condition automatiquement vérifiée si $k$ est séparablement clos).
  \item Pour tout $\alpha$, il existe un point $x'_\alpha\in X'_\alpha$ dont
  les projections respectives dans $X$ et $Y'_\alpha$ sont $x$ et le point
  fermé $y'_\alpha$ de $Y'_\alpha$, et tel que $\mathcal F'_\alpha$ soit
  $Y'_\alpha$-plat au point $x'_\alpha$.
\end{enumerate}
Dans ces conditions, $\mathcal F$ est $f$-plat au point $x$.
\end{proposition}

\begin{proof}
La question étant locale sur $X$, on peut évidemment se borner au cas où $f$
est un morphisme de type fini et supposer $\mathcal F_x\neq0$. Nous procéderons
en plusieurs étapes.

\medskip
\noindent\textbf{I)} \emph{Réduction au cas où $A'_\alpha$ est un anneau local
d'un $Y$-préschéma de type fini et où le corps résiduel $k'_\alpha$ de
$A'_\alpha$ est une extension finie de $k$.}

\medskip
Comme la réduction se fait séparément sur chaque $A'_\alpha$, on peut supprimer
dans cette partie l'indice $\alpha$. Soit $\mathfrak m'$ l'idéal maximal de
$A'$. Considérons $A'$ comme limite inductive de ses sous-$A$-algèbres de type
fini $B_\lambda$, et posons $\mathfrak n_\lambda=\mathfrak m'\cap B_\lambda$~;
$A'$ est aussi limite inductive de ses sous-anneaux locaux
$B'_\lambda=(B_\lambda)_{\mathfrak n_\lambda}$ (5.13.3), et on a évidemment
$\mathfrak m'\cap B'_\lambda=\mathfrak n'_\lambda$, idéal maximal de
$B'_\lambda$. Il existe donc (11.2.6) un indice $\lambda$ tel que, en posant
$Z'_\lambda=\operatorname{Spec}(B'_\lambda)$,
$\mathcal F\otimes_A B'_\lambda$ soit $Z'_\lambda$-plat au point $x'_\lambda$,
projection de $x'$, la projection de $x'_\lambda$ sur $Z'_\lambda$ étant le
point fermé $z'_\lambda$ de $Z'_\lambda$.

On peut donc supposer qu'il existe un schéma affine $Z'$ de type fini sur $Y$
et un point $z'$ de $Z'$ tel que $A'=\mathcal O_{Z',z'}$, et que, si l'on pose
$T'=X\times_Y Z'$, il existe un point $t'\in T'$ dont les projections dans
$Z'$ et $X$ sont $z'$ et $x$, et tel que $\mathcal F\otimes_Y Z'$ soit
$Z'$-plat au point $t'$. Soient $Z'_1$ (resp. $X_1$) un sous-préschéma fermé de
$Z'$ (resp. $X$) ayant pour espace sous-jacent l'adhérence de $\{z'\}$ (resp.
$\{x\}$), et posons $T'_1=X_1\times_Y Z'_1$. L'ensemble $U$ des points de
$T'$ où $\mathcal F\otimes_Y Z'$ est $Z'$-plat est ouvert dans $T'$ (11.1.1) et
contient $t'$, donc

\oldpage[IV-3]{149}

$V=U\cap T'_1$ est ouvert non vide dans $T'_1$. L'anneau $A$, étant artinien,
est un anneau de Jacobson, donc (10.4.6) $Z'_1$ et $T'_1$ sont des préschémas de
Jacobson~; il existe par suite dans $V$ un point $t'_0$ fermé dans $V$, et son
image $z'_0$ dans $Z'_1$ est un point fermé de $Z'_1$ (10.4.7). Soient $f_1$ la
restriction de $f$ à $X_1$ et $f'_1:T'_1\to Z'_1$ la projection canonique~;
$V\cap f_1'^{-1}(z'_0)$ est un ouvert non vide dans
$f_1^{-1}(y)\otimes_{k(y)}k(z'_0)$, et comme ce dernier préschéma est plat sur
$f_1^{-1}(y)$, un point maximal $t'_1$ de $V\cap f_1'^{-1}(z'_0)$ est
nécessairement au-dessus de l'unique point maximal $x$ de $X_1$ (2.3.4).
Enfin, $k(z'_0)$ est une extension finie de $k$ (\textbf I, 6.4.2) et
l'homomorphisme $A=\mathcal O_{Y,y}\to A'=\mathcal O_{Z',z'}$ se factorise en
$\mathcal O_{Y,y}\to\mathcal O_{Z',z'_0}\to\mathcal O_{Z',z'}$, donc son noyau
est contenu dans celui de $\mathcal O_{Y,y}\to\mathcal O_{Z',z'_0}$. Ceci
achève la réduction annoncée.

\medskip
\noindent\textbf{II)} \emph{Réduction au cas où les $A'_\alpha$ sont en nombre
fini et sont des $A$-algèbres finies.} --- Soit $\mathfrak m'_\alpha$ l'idéal
maximal de $A'_\alpha$~; comme $A'_\alpha$ est un anneau local noethérien,
l'intersection des $\mathfrak m'^n_\alpha$ est $0$ (\textbf I, 7.3.5)~;
l'intersection des $u_\alpha^{-1}(\mathfrak m'^n_\alpha)$, pour tous les indices
$n$ et $\alpha$, est donc égale à l'intersection des noyaux des $u_\alpha$, donc
est réduite à $0$ par l'hypothèse (i). Puisque $A$ est artinien, il y a un nombre
fini de ces idéaux dont l'intersection est déjà $0$~; notons-les
$u_i^{-1}(\mathfrak m_i'^{\,n_i})$ $(1\leq i\leq r)$. Comme les $A'_i$ sont
noethériens, les $\mathfrak m_i'^{\,h}/\mathfrak m_i'^{\,h+1}$ sont des
$(A'_i/\mathfrak m'_i)$-modules de longueur finie, et comme
$A'_i/\mathfrak m'_i$ est un $k$-espace vectoriel de rang fini, on voit que
$A''_i=A'_i/\mathfrak m_i'^{\,n_i}$ est une $A$-algèbre finie et un anneau local
artinien. La réduction annoncée résulte donc de (2.1.4).

\medskip
\noindent\textbf{III)} \emph{Fin de la démonstration.} --- Supposons désormais
que les $A'_i$ $(1\leq i\leq r)$ soient en nombre fini et soient des
$A$-algèbres finies. Pour tout $i$, le corps résiduel $k_i$ de $A'_i$ est une
extension finie de $k$~; utilisant l'hypothèse (ii), on en conclut que l'image
réciproque de $x$ dans $X'_i$ est réduite au seul point $x'_i$ (4.3.2). Soient
$Y'$ le préschéma somme des $Y'_i$, $X'=X\times_Y Y'$, somme des $X'_i$,
$\mathcal F'=\mathcal F\otimes_Y Y'$. L'hypothèse entraîne que
$\mathcal F'$ est $Y'$-plat aux points de l'image réciproque de $x$ par la
projection $p:X'\to X$. Comme $p$ est de type fini, il existe par suite un
ouvert $U'\supset p^{-1}(x)$ tel que $\mathcal F'$ soit $Y'$-plat aux points de
$U'$ (11.1.1). En outre, le morphisme $Y'\to Y$ est fini puisque les $A'_i$
sont des $A$-algèbres finies~; donc $p$ est un morphisme fini (\textbf{II},
6.1.5), par suite fermé (\textbf{II}, 6.1.10), et il existe donc un voisinage
affine $U$ de $x$ dans $X$ tel que $p^{-1}(U)\subset U'$. Soient $B$ et $A'$
les anneaux des schémas $U$ et $Y'$ ($A'$ étant le composé direct des $A'_i$)~;
remplaçant $X$ par $U$, on a donc $\mathcal F=\widetilde M$, où $M$ est un
$B$-module et, par hypothèse, $M'=M\otimes_A A'$ est un $A'$-module plat
(2.1.2)~; comme l'homomorphisme $A\to A'$ est injectif par construction, on
peut appliquer (11.4.3), qui prouve que $M$ est un $A$-module plat, C.Q.F.D.
\end{proof}

\begin{proposition}[11.4.12]
\label{IV.11.4.12-fr}
Soient $A$ un anneau local artinien de corps résiduel $k$,
$Y=\operatorname{Spec}(A)$, $f:X\to Y$ un morphisme localement de type fini,
$\mathcal F$ un $\mathcal O_X$-Module cohérent. Soient $(A'_\alpha)$ une famille
d'anneaux locaux, et pour chaque $\alpha$, $u_\alpha:A\to A'_\alpha$ un
homomorphisme d'anneaux. Posons $Y'_\alpha=\operatorname{Spec}(A'_\alpha)$,
$X'_\alpha=X\times_Y Y'_\alpha$, $f'_\alpha=f_{(Y'_\alpha)}$,
$\mathcal F'_\alpha=\mathcal F\otimes_A A'_\alpha$. Soit $x$ un point de $X$,
et supposons vérifiées les conditions suivantes~:
\begin{enumerate}[label=(\roman*)]
  \item L'intersection des noyaux des $u_\alpha$ est réduite à $0$.
  \item Pour tout $\alpha$, $\mathcal F'_\alpha$ est $f'_\alpha$-plat en tous
  les points $x'_\alpha\in X'_\alpha$ dont les projections respectives dans
  $X$ et $Y'_\alpha$ sont $x$ et le point fermé $y'_\alpha$ de $Y'_\alpha$.
\end{enumerate}
Alors $\mathcal F$ est $f$-plat au point $x$.
\end{proposition}

\begin{proof}
% IV3_B22_P149_ANCHOR

\oldpage[IV-3]{150}

Par hypothèse, l'intersection des noyaux des $u_\alpha$ est réduite à $0$~;
comme $A$ est artinien, il y a déjà un nombre fini de ces noyaux dont
l'intersection est $0$, donc on peut se borner au cas où la famille
$(A'_\alpha)$ est finie. Soit $k'$ une clôture algébrique de $k$~; on sait
$(\mathbf 0_{\mathrm{III}},10.3.1)$ qu'il existe un homomorphisme local
$A\to B$ faisant de $B$ un $A$-module plat, tel que $B$ soit un anneau local
artinien, entier sur $A$ et que $B\otimes_A k$ soit isomorphe à $k'$.

Par platitude, les noyaux des homomorphismes
$B\to B'_\alpha=B\otimes_A A'_\alpha$ se déduisent de ceux des $u_\alpha$ par
tensorisation avec $B$, et comme ils sont en nombre fini, leur intersection est
réduite à $0$ $(\mathbf 0_{\mathrm I},6.1.3)$. Considérons les anneaux
$B'_{\alpha\beta}$, localisés de $B'_\alpha$ en ses idéaux maximaux~; on sait
(Bourbaki, \emph{Alg. comm.}, chap.~II, \S~3, n$^\circ$~3, cor.~2 du th.~1)
que l'intersection des noyaux des homomorphismes
$B'_\alpha\to B'_{\alpha\beta}$ (pour un $\alpha$ donné) est réduite à $0$~;
on en conclut que l'intersection des noyaux des homomorphismes composés
$v_{\alpha\beta}:B\to B'_\alpha\to B'_{\alpha\beta}$ ($\alpha$ et $\beta$
variables) est réduite à $0$. D'autre part, comme $B$ est entier sur $A$,
$B'_\alpha$ est entier sur $A'_\alpha$, donc ses idéaux maximaux sont au-dessus
de l'idéal maximal de $A'_\alpha$. Si l'on pose
$Z'_{\alpha\beta}=\operatorname{Spec}(B'_{\alpha\beta})$,
$T'_{\alpha\beta}=X'_\alpha\times_{Y'_\alpha}Z'_{\alpha\beta}$,
$f'_{\alpha\beta}=f_{(Z'_{\alpha\beta})}$,
$\mathcal G=\mathcal F\otimes_A B$ et
$\mathcal G'_{\alpha\beta}=\mathcal G\otimes_B B'_{\alpha\beta}$, on voit donc
que les hypothèses (i) et (ii) sont encore vérifiées lorsqu'on remplace
respectivement $A$, $A'_\alpha$, $u_\alpha$, $\mathcal F$ et $x$ par $B$,
$B'_{\alpha\beta}$, $v_{\alpha\beta}$, $\mathcal G$ et un point $t$ de
$T=X\otimes_A B$ au-dessus de $x$. Comme le corps résiduel de $B$ est
séparablement clos, on déduit de (11.4.11) que $\mathcal G$ est plat sur $B$ au
point $t$. Mais puisque $B$ est un $A$-module fidèlement plat, on conclut de
(2.1.4) que $\mathcal F$ est plat sur $A$ au point $x$, ce qui prouve (11.4.12).
\end{proof}

\begin{corollary}[11.4.13]
\label{IV.11.4.13-fr}
Soient $A$ un anneau artinien local, $Y=\operatorname{Spec}(A)$,
$f:X\to Y$ un morphisme localement de type fini, $\mathcal F$ un
$\mathcal O_X$-Module cohérent. Soit $(Y'_\alpha)$ une famille de
$Y$-préschémas, et pour tout $\alpha$, posons
$X'_\alpha=X\times_Y Y'_\alpha$, $f'_\alpha=f_{(Y'_\alpha)}$,
$\mathcal F'_\alpha=\mathcal F\otimes_Y Y'_\alpha$. Soit $x$ un point de $X$ et
supposons vérifiées les hypothèses suivantes~:
\begin{enumerate}[label=(\roman*)]
  \item L'intersection des noyaux des homomorphismes
  $A\to\Gamma(Y'_\alpha,\mathcal O_{Y'_\alpha})$ correspondant aux morphismes
  structuraux est réduite à $0$.
  \item Pour tout $\alpha$, $\mathcal F'_\alpha$ est $f'_\alpha$-plat en tous
  les points $x'_\alpha\in X'_\alpha$ dont la projection dans $X$ est $x$.
\end{enumerate}
Alors $\mathcal F$ est $f$-plat au point $x$.
\end{corollary}

\begin{proof}
En effet, pour tout $y'_\alpha\in Y'_\alpha$, considérons le schéma local
$Y''_{\alpha,y'_\alpha}=\operatorname{Spec}(\mathcal O_{Y'_\alpha,y'_\alpha})$~;
en vertu de (2.1.4),
$\mathcal F'_\alpha\otimes_{Y'_\alpha}Y''_{\alpha,y'_\alpha}$ est
$Y''_{\alpha,y'_\alpha}$-plat aux points dont les projections sur $X$ et
$Y''_{\alpha,y'_\alpha}$ sont $x$ et le point fermé de
$Y''_{\alpha,y'_\alpha}$. En outre le noyau de l'homomorphisme
$A\to\Gamma(Y'_\alpha,\mathcal O_{Y'_\alpha})$ est l'intersection des noyaux
des homomorphismes
$A\to\Gamma(Y''_{\alpha,y'_\alpha},\mathcal O_{Y''_{\alpha,y'_\alpha}})$, car
on se ramène aussitôt au cas où $Y'_\alpha$ est affine, et il suffit alors
d'appliquer Bourbaki, \emph{Alg. comm.}, chap.~II, \S~3, n$^\circ$~3, cor.~2 du
th.~1. En remplaçant la famille $(Y'_\alpha)$ par la famille des
$Y''_{\alpha,y'_\alpha}$, on est donc ramené à (11.4.12).
\end{proof}

\subsection{Descente de la platitude par des morphismes quelconques : cas général}
\label{subsection:IV.11.5-fr}

\begin{theorem}[11.5.1]
\label{IV.11.5.1-fr}
Soient $Y$ un préschéma localement noethérien, $f:X\to Y$ un morphisme
localement de type fini, $\mathcal F$ un $\mathcal O_X$-Module cohérent, $x$ un
point de $X$, $y=f(x)$. Soit $(Y'_\alpha)$ une famille de $Y$-préschémas locaux
$Y'_\alpha=\operatorname{Spec}(A'_\alpha)$ telle que les images des points
fermés $y'_\alpha$ des $Y'_\alpha$ soient toutes égales à $y$. Pour tout
$\alpha$, soient $\mathfrak m'_\alpha$ l'idéal maximal de $A'_\alpha$, et
$u_\alpha:\mathcal O_y\to A'_\alpha$ l'homomorphisme canonique
(\textbf I, 2.4.4)~; on suppose que les intersections finies des idéaux
$u_\alpha^{-1}(\mathfrak m_\alpha'^{\,n_\alpha})$ forment un système fondamental
de voisinages de $0$ dans $\mathcal O_y$. Posons
$X'_\alpha=X\times_Y Y'_\alpha$, $f'_\alpha=f_{(Y'_\alpha)}$,
$\mathcal F'_\alpha=\mathcal F\otimes_Y Y'_\alpha$, et supposons vérifiée l'une
des hypothèses suivantes~:
\begin{enumerate}[label=(\roman*)]
  \item Pour tout $\alpha$, $\mathcal F'_\alpha$ est $f'_\alpha$-plat en tous
  les points dont la projection dans $X$ est égale à $x$ et la projection dans
  $Y'_\alpha$ égale à $y'_\alpha$.
  \item Pour tout $\alpha$, il existe un $x'_\alpha\in X'_\alpha$ dont la
  projection sur $X$ est $x$ et la projection dans $Y'_\alpha$ égale à
  $y'_\alpha$, tel que $\mathcal F'_\alpha$ soit $f'_\alpha$-plat au point
  $x'_\alpha$, et $k(x)$ est une extension primaire de $k(y)$.
\end{enumerate}
Dans ces conditions, $\mathcal F$ est $f$-plat au point $x$.
\end{theorem}

\begin{proof}
Soit $\mathfrak m_y$ l'idéal maximal de $\mathcal O_y$~; comme $\mathcal O_y$
et $\mathcal O_x$ sont noethériens et $\mathcal O_y\to\mathcal O_x$ un
homomorphisme local, il suffit, en vertu de
$(\mathbf 0_{\mathrm{III}},10.2.2)$, de prouver que pour tout $n>0$,
$\mathcal F_x/\mathfrak m_y^n\mathcal F_x$ est un
$(\mathcal O_y/\mathfrak m_y^n)$-module plat. Désignons par
$(\mathfrak J_\lambda)$ la famille des intersections finies des
$u_\alpha^{-1}(\mathfrak m_\alpha'^{\,n_\alpha})$~; par hypothèse, il existe
$\lambda$ tel que $\mathfrak J_\lambda\subset\mathfrak m_y^n$, et comme
\[
\mathcal F_x\otimes_{\mathcal O_y}(\mathcal O_y/\mathfrak m_y^n)
=\bigl(\mathcal F_x\otimes_{\mathcal O_y}(\mathcal O_y/\mathfrak J_\lambda)\bigr)
 \otimes_{\mathcal O_y/\mathfrak J_\lambda}(\mathcal O_y/\mathfrak m_y^n),
\]
il suffira de prouver que $\mathcal F_x/\mathfrak J_\lambda\mathcal F_x$ est un
$(\mathcal O_y/\mathfrak J_\lambda)$-module plat. Or, pour chaque $\alpha$ tel
que $\mathfrak J_\lambda\subset u_\alpha^{-1}(\mathfrak m_\alpha'^{\,n_\alpha})$,
on a, par passage aux quotients, un homomorphisme
$u'_\alpha:\mathcal O_y/\mathfrak J_\lambda\to
A'_\alpha/\mathfrak m_\alpha'^{\,n_\alpha}$, et l'intersection des noyaux des
$u'_\alpha$ est réduite à $0$. Compte tenu de (\textbf I, 3.6.1), on voit qu'on
est ramené à (11.4.11) dans le cas (ii), et à (11.4.12) dans le cas (i).
\end{proof}

\begin{corollary}[11.5.2]
\label{IV.11.5.2-fr}
Soient $Y$ un préschéma localement noethérien, $f:X\to Y$ un morphisme
localement de type fini, $\mathcal F$ un $\mathcal O_X$-Module cohérent, $x$ un
point de $X$, $y=f(x)$~; posons $A=\mathcal O_y$. Soit $u:A\to A'$ un
homomorphisme de $A$ dans un anneau de Zariski $A'$, tel que l'image réciproque
$u^{-1}(\mathfrak r')$ du radical $\mathfrak r'$ de $A'$ soit l'idéal maximal
$\mathfrak m$ de $A$~; supposons en outre que l'homomorphisme
$\widehat u:\widehat A\to\widehat{A'}$ soit injectif. Posons
$Y'=\operatorname{Spec}(A')$, $X'=X\times_Y Y'$, $f'=f_{(Y')}$,
$\mathcal F'=\mathcal F\otimes_Y Y'$. Pour que $\mathcal F$ soit $f$-plat au
point $x$, il faut et il suffit que $\mathcal F'$ soit $f'$-plat en tout point
dont la projection dans $X$ est égale à $x$ et dont la projection dans $Y'$ est
égale à un point fermé de $Y'$.

Si en outre $A'$ est une $A$-algèbre finie, on peut dans ce qui précède
remplacer l'hypothèse que $\widehat u$ est injectif par l'hypothèse que $u$ est
injectif.
\end{corollary}

\begin{proof}
\oldpage[IV-3]{151}
% IV3_B22_P151_ANCHOR

Comme $A$ (resp. $A'$) s'identifie à un sous-anneau de $\widehat A$ (resp.
$\widehat{A'}$) (Bourbaki, \emph{Alg. comm.}, chap.~III, \S~3, n$^\circ$~3,
prop.~6), on voit d'abord que $u$ lui-même est injectif et que $\widehat u$ est
son prolongement par continuité à $\widehat A$.

Soit $(\mathfrak m'_\alpha)$ la famille des idéaux maximaux de $A'$~; comme on a
\[
\mathfrak m_\alpha'^{\,n}\widehat{A'}
=\widehat{\mathfrak m}_\alpha'^{\,n},
\qquad\text{et}\qquad
\widehat{\mathfrak m}_\alpha'^{\,n}\cap A'
=\mathfrak m_\alpha'^{\,n},
\]
et que les $\mathfrak m_\alpha'^{\,n}$ sont ouverts dans $A'$, on a
$u^{-1}(\mathfrak m_\alpha'^{\,n})
=A\cap\widehat u^{-1}(\widehat{\mathfrak m}_\alpha'^{\,n})$, et il suffira de
montrer que dans $\widehat A$, les intersections finies des
$\widehat u^{-1}(\widehat{\mathfrak m}_\alpha'^{\,n_\alpha})$ forment un système
fondamental de voisinages de $0$, ce qui permettra d'appliquer (11.5.1) aux
homomorphismes composés $v_\alpha\circ u:A\to A'_{\mathfrak m'_\alpha}$,
où $v_\alpha:A'\to A'_{\mathfrak m'_\alpha}$ est l'homomorphisme canonique.
Comme $\widehat A$ est complet, il suffira de montrer que l'intersection des
$\widehat u^{-1}(\widehat{\mathfrak m}_\alpha'^{\,n_\alpha})$ est réduite à
$0$ (Bourbaki, \emph{Alg. comm.}, chap.~III, \S~2, n$^\circ$~7, prop.~8, où on
peut dans la démonstration remplacer la suite décroissante par un ensemble
filtrant quelconque). Or, pour tout $\alpha$ fixé, l'intersection des idéaux
$\widehat{\mathfrak m}_\alpha'^{\,n}\widehat{A'}_{\widehat{\mathfrak m}'_\alpha}$
pour $n>0$ est réduite à $0$ dans l'anneau local noethérien
$\widehat{A'}_{\widehat{\mathfrak m}'_\alpha}$. D'autre part les
$\widehat{\mathfrak m}'_\alpha$ sont les idéaux maximaux de $\widehat{A'}$, donc
l'homomorphisme canonique
$\widehat{A'}\to\prod_\alpha\widehat{A'}_{\widehat{\mathfrak m}'_\alpha}$ est
injectif (Bourbaki, \emph{Alg. comm.}, chap.~II, \S~3, n$^\circ$~3, cor.~2 du
th.~1), et comme par hypothèse $\widehat u:\widehat A\to\widehat{A'}$ est aussi
injectif, cela achève la démonstration dans le cas général. La dernière
assertion résulte de ce que $\widehat A$ est un $A$-module fidèlement plat
($\mathbf 0_{\mathrm I},7.3.5$) et $\widehat{A'}=A'\otimes_A\widehat A$ puisque
$A'$ est par hypothèse un $A$-module de type fini (Bourbaki, \emph{Alg. comm.},
chap.~III, \S~3, n$^\circ$~4, th.~3 et chap.~IV, \S~2, n$^\circ$~5, cor.~3 de la
prop.~9).
\end{proof}

\oldpage[IV-3]{152}

\begin{proposition}[11.5.3]
\label{IV.11.5.3-fr}
Soient $f:X\to Y$ un morphisme localement de présentation finie,
$\mathcal F$ un $\mathcal O_X$-Module quasi-cohérent de présentation finie, $x$
un point de $X$, $y=f(x)$, $g=(\psi,\theta):Y'\to Y$ un morphisme propre de
présentation finie. Supposons que~:
\begin{enumerate}[label=(\roman*)]
  \item l'homomorphisme $\theta_y:\mathcal O_y\to
  (g_*\mathcal O_{Y'})_y$ soit injectif~;
  \item pour tout $x'\in X'=X\times_Y Y'$ dont la projection dans $X$ est $x$,
  $\mathcal F'=\mathcal F\otimes_Y Y'$ soit $Y'$-plat au point $x'$.
\end{enumerate}
Alors $\mathcal F$ est $f$-plat au point $x$.
\end{proposition}

\begin{proof}
La question étant locale sur $X$, on peut supposer $f$ de présentation finie.
Soit $p:X'\to X$ la projection canonique. Comme
$f'=f_{(Y')}:X'\to Y'$ est de présentation finie (1.6.2) et que $\mathcal F'$
est un $\mathcal O_{X'}$-Module de présentation finie
$(\mathbf 0_{\mathrm I},5.2.5)$, il résulte de (11.3.1) que l'ensemble $U'$ des
points de $X'$ où $\mathcal F'$ est $f'$-plat est ouvert. Comme $U'$ contient
$p^{-1}(x)$ par hypothèse, et que $p$ est propre, donc fermé, $U'$ contient un
ensemble de la forme $p^{-1}(U)$, où $U$ est un voisinage de $x$. Remplaçant $X$
par $U$, on peut donc supposer déjà que $\mathcal F'$ est $f'$-plat.

D'autre part, en tenant compte de (\textbf I, 3.6.5), (\textbf{II}, 5.4.2) et
(2.1.4), on peut remplacer $Y$ par $\operatorname{Spec}(\mathcal O_y)$,
autrement dit supposer que $Y=\operatorname{Spec}(A)$, où $A$ est un anneau
local. Sous ces conditions, on va prouver que $\mathcal F$ est $f$-plat. En
vertu de (5.13.3), $A$ est limite inductive de sous-anneaux locaux noethériens
$A_\lambda$ tels que l'injection canonique $A_\lambda\to A$ soit un
homomorphisme local. En vertu de (8.9.1), on peut supposer que
$X=X_\lambda\otimes_{A_\lambda}A$, $f=f_\lambda\otimes1_A$,
$\mathcal F=\mathcal F_\lambda\otimes_{A_\lambda}A$, pour un $\lambda$
convenable, $f_\lambda:X_\lambda\to Y_\lambda=\operatorname{Spec}(A_\lambda)$
étant un morphisme de type fini et $\mathcal F_\lambda$ un
$\mathcal O_{X_\lambda}$-Module cohérent. De même, on peut supposer que
$Y'=Y'_\lambda\otimes_{A_\lambda}A$, $g=g_\lambda\otimes1_A$, où
$g_\lambda:Y'_\lambda\to Y_\lambda$ est un morphisme de présentation finie~;
en outre, (8.10.5, (xii)) permet de supposer $g_\lambda$ propre. Comme par
hypothèse l'homomorphisme $A\to\Gamma(Y',\mathcal O_{Y'})$ est injectif et que
$A_\lambda$ est un sous-anneau de $A$, l'homomorphisme
$A_\lambda\to\Gamma(Y',\mathcal O_{Y'})$ est aussi injectif. Enfin, en vertu de
(11.2.6), on peut supposer $\lambda$ pris tel que
$\mathcal F'_\lambda=\mathcal F_\lambda\otimes_{Y_\lambda}Y'_\lambda$ soit
$f'_\lambda$-plat, puisque $\mathcal F'=\mathcal F'_\lambda\otimes_{Y'_\lambda}Y'$.
Ces remarques prouvent que l'on peut désormais supposer l'anneau $A$
noethérien, les autres hypothèses de (11.5.3) étant vérifiées.

Posons alors $B=\widehat A$, $Z=\operatorname{Spec}(B)$~; comme $B$ est un
$A$-module fidèlement plat $(\mathbf 0_{\mathrm I},7.3.5)$, il revient au même
de dire que $\mathcal F$ est $f$-plat ou que $\mathcal F\otimes_Y Z$ est
$Z$-plat (2.1.4). De même, si l'on pose $Z'=Y'\times_Y Z$, le morphisme
$Z'\to Y'$ est fidèlement plat (2.2.13), donc il revient au même de dire que
$\mathcal F'$ est $f'$-plat ou que $\mathcal F'\otimes_{Y'}Z'$ est $Z'$-plat~;
enfin $h=g_{(Z)}:Z'\to Z$ est propre (\textbf{II}, 5.4.2) et de type fini
(1.5.2), $\widehat A$ est noethérien, et si $z$ est son point fermé,
l'homomorphisme $\mathcal O_z\to(h_*\mathcal O_{Z'})_z$ est injectif, car il
résulte de (2.3.1) que $h_*\mathcal O_{Z'}=g_*\mathcal O_{Y'}\otimes_Y Z$,
et notre assertion résulte de l'hypothèse (i) et de la définition des modules
plats $(\mathbf 0_{\mathrm I},6.1.1)$.

\oldpage[IV-3]{153}

On peut donc désormais supposer l'anneau local noethérien $A$ complet~; la
démonstration sera achevée si l'on prouve que l'intersection des noyaux des
homomorphismes $A=\mathcal O_y\to\mathcal O_{y'}$ (où $y'$ parcourt
$g^{-1}(y)$) est réduite à $0$. En effet, les $\mathcal O_{y'}$ sont des anneaux
locaux noethériens, donc pour chaque $y'$ l'intersection des
$\mathfrak m_{y'}^n$ $(n\geq0)$ est réduite à $0$~; si $\mathfrak a_{n,y'}$ est
l'image réciproque dans $A$ de $\mathfrak m_{y'}^n$, les intersections finies
des $\mathfrak a_{n,y'}$ sont des voisinages de $0$ dans $A$ et l'intersection de
tous les $\mathfrak a_{n,y'}$ est réduite à $0$~; elles forment donc un système
fondamental de voisinages de $0$ dans $A$ (Bourbaki, \emph{Alg. comm.}, chap.~III,
\S~2, n$^\circ$~7, prop.~8, où l'on peut remplacer la suite décroissante par un
ensemble filtrant quelconque), et l'on pourra par suite appliquer (11.5.1).
Or, soit $s\in A$ un élément appartenant au noyau de chacun des homomorphismes
$A\to\mathcal O_{y'}$~; l'image $s'$ de $s$ dans
$\Gamma(Y',\mathcal O_{Y'})$ est donc une section de $\mathcal O_{Y'}$ au-dessus
de $Y'$ telle que $s'_{y'}=0$ pour tout $y'\in g^{-1}(y)$~; il y a par suite un
voisinage $V'$ de $g^{-1}(y)$ dans $Y'$ tel que $s'|V'=0$. Mais comme $g$ est
fermé, $V'$ contient un ensemble de la forme $g^{-1}(V)$, où $V$ est un
voisinage ouvert de $y$ dans $Y$~; or $Y$ est un schéma local, donc le seul
voisinage du point fermé $y$ de $Y$ est $Y$ tout entier, autrement dit $V'=Y'$,
$s'=0$, et comme $A\to\Gamma(Y',\mathcal O_{Y'})$ est injectif par hypothèse,
$s=0$. C.Q.F.D.
\end{proof}

Le cas particulier suivant de (11.5.3) nous sera utile au chap.~V.

\begin{corollary}[11.5.4]
\label{IV.11.5.4-fr}
Soit $f=(\psi,\theta):X\to Y$ un morphisme propre de présentation finie, et
soit $p:X\times_YX\to X$ la première projection. On suppose que
$\theta:\mathcal O_Y\to f_*(\mathcal O_X)$ est injectif. Alors, pour que $f$
soit plat, il faut et il suffit que $p$ le soit.
\end{corollary}

\begin{proposition}[11.5.5]
\label{IV.11.5.5-fr}
Soient $A$ un anneau, $Y=\operatorname{Spec}(A)$, $f:X\to Y$ un morphisme
localement de présentation finie, $\mathcal F$ un $\mathcal O_X$-Module
quasi-cohérent de présentation finie, $x$ un point de $X$. Soit
$A\to A'$ un homomorphisme injectif faisant de $A'$ une algèbre entière sur $A$.
Posons $Y'=\operatorname{Spec}(A')$, $X'=X\times_Y Y'$, $f'=f\times1$,
$\mathcal F'=\mathcal F\otimes_Y Y'$. Alors, si $\mathcal F'$ est $f'$-plat en
tout point de $X'$ dont la projection dans $X$ est égale à $x$, $\mathcal F$ est
$f$-plat au point $x$.
\end{proposition}

\begin{proof}
Supposons d'abord que $A'$ soit une $A$-algèbre finie et de présentation finie~;
alors le morphisme $Y'\to Y$ est propre (\textbf{II}, 6.1.11) et de présentation
finie, donc les hypothèses de (11.5.3) sont vérifiées, d'où la conclusion. Dans
le cas général, la proposition résultera de ce cas particulier, du fait que $A'$
est limite inductive de ses sous-$A$-algèbres finies $A'_\lambda$ et des deux
lemmes suivants~:
\end{proof}

\begin{lemma}[11.5.5.1]
\label{IV.11.5.5.1-fr}
Toute $A$-algèbre finie $A'$ est limite inductive de $A$-algèbres $A'_\lambda$
qui sont finies et de présentation finie.
\end{lemma}

\begin{proof}
On raisonne comme dans (1.9.3.1). On a en effet $A'=B/\mathfrak J$, où $B$ est
une $A$-algèbre finie qui est un $A$-module libre, et $\mathfrak J$ un idéal de
$B$ (1.4.7.1). Or, $\mathfrak J$ est limite inductive des idéaux
$\mathfrak J_\lambda\subset\mathfrak J$ de $B$ qui sont de type fini (et a
fortiori des $A$-modules de type fini), donc, par l'exactitude du foncteur
$\varinjlim$, $A'$ est limite inductive des $A$-algèbres
$B/\mathfrak J_\lambda$~; or $B/\mathfrak J_\lambda$ est par définition un
$A$-module de présentation finie, donc aussi (1.4.7) une $A$-algèbre de
présentation finie.
\end{proof}
% IV3_B22_P153_END

\oldpage[IV-3]{154}

Pour appliquer ce lemme à la situation de (11.5.5) on notera en outre que si
l'homomorphisme $A\to A'$ est injectif, il en est ainsi \emph{a fortiori} de
$A\to A'_{\lambda}$, pour tout $\lambda$.

\begin{lemma}[11.5.5.2]
\label{IV.11.5.5.2-fr}
Soient $A$ un anneau, $A'$ une $A$-algèbre, $(A'_{\lambda})$ un système inductif
de $A$-algèbres tel que $A'=\varinjlim A'_{\lambda}$~; on pose
$Y=\operatorname{Spec}(A)$, $Y'=\operatorname{Spec}(A')$,
$Y'_{\lambda}=\operatorname{Spec}(A'_{\lambda})$. Soient $f:X\to Y$ un
morphisme de présentation finie, $\mathcal F$ un $\mathcal O_X$-Module
quasi-cohérent de présentation finie~; on pose $X'=X\times_Y Y'$,
$f'=f_{(Y')}$, $\mathcal F'=\mathcal F\otimes_Y Y'$,
$X'_{\lambda}=X\times_Y Y'_{\lambda}$, $f'_{\lambda}=f_{(Y'_{\lambda})}$,
$\mathcal F'_{\lambda}=\mathcal F\otimes_Y Y'_{\lambda}$. Soit $x$ un point de
$X$, tel que $\mathcal F'$ soit $f'$-plat en tous les points $x'\in X'$
au-dessus de $x$~; alors il existe $\lambda$ tel que $\mathcal F'_{\lambda}$
soit $f'_{\lambda}$-plat en tous les points de $X'_{\lambda}$ au-dessus de $x$.
\end{lemma}

\begin{proof}
Soit $U'$ l'ensemble des $x'\in X'$ tels que $\mathcal F'$ soit $f'$-plat au
point $x'$~; on sait (11.3.1) que $U'$ est ouvert dans $X'$ puisque $f'$ est de
présentation finie (1.6.2)~; de même l'ensemble $U'_{\lambda}$ des points de
$X'_{\lambda}$ où $\mathcal F'_{\lambda}$ est $f'_{\lambda}$-plat est ouvert
dans $X'_{\lambda}$, et on sait en outre (11.2.6) que $U'$ est réunion des
$v_{\lambda}^{-1}(U'_{\lambda})$, où
$v_{\lambda}:X'\to X'_{\lambda}$ est la projection canonique. Considérons le
schéma $T=\operatorname{Spec}(k(x))$~; posons $T'=T\times_Y Y'$,
$T'_{\lambda}=T\times_Y Y'_{\lambda}$, et soient
$w_{\lambda}:T'\to T'_{\lambda}$, $g':T'\to X'$,
$g'_{\lambda}:T'_{\lambda}\to X'_{\lambda}$ les projections canoniques. Posons
$V'_{\lambda}=g_{\lambda}'^{-1}(U'_{\lambda})$,
$V'=g'^{-1}(U')=\bigcup_{\lambda}w_{\lambda}^{-1}(V'_{\lambda})$. Par
hypothèse on a (compte tenu de (\textbf I, 3.6.1)) $V'=T'$~; comme $T'$ est
quasi-compact, il existe $\lambda$ tel que
$w_{\lambda}^{-1}(V'_{\lambda})=T'$. On déduit alors de (8.3.3) appliqué aux
parties fermées quasi-compactes $T'_{\lambda}-V'_{\lambda}$ de
$T'_{\lambda}$, qu'il existe $\mu\geq\lambda$ tel que
$T'_{\mu}=V'_{\mu}$~; cela signifie que $\mathcal F'_{\mu}$ est
$f'_{\mu}$-plat en tous les points de $X'_{\mu}$ dont la projection dans $X$
est $x$. C.Q.F.D.
\end{proof}

\subsection{Descente de la platitude par des morphismes quelconques : cas d'un
préschéma de base unibranche}
\label{subsection:IV.11.6-fr}

\begin{theorem}[11.6.1]
\label{IV.11.6.1-fr}
Soient $A$ un anneau local intègre géométriquement unibranche
$(\mathbf 0,23.2.1)$, $Y=\operatorname{Spec}(A)$, $f:X\to Y$ un morphisme
localement de présentation finie, $\mathcal F$ un $\mathcal O_X$-Module
quasi-cohérent de présentation finie. Soient $A'$ un anneau local,
$u:A\to A'$ un homomorphisme local injectif~; on pose
$Y'=\operatorname{Spec}(A')$, $X'=X\times_Y Y'$, $f'=f_{(Y')}$,
$\mathcal F'=\mathcal F\otimes_Y Y'$. Soient $x$ un point de $X$ dont la
projection $f(x)=y$ est le point fermé de $Y$, $x'$ un point de $X'$ dont les
projections dans $X$ et $Y'$ sont respectivement $x$ et le point fermé $y'$ de
$Y'$. Alors, si $\mathcal F'$ est $f'$-plat au point $x'$, $\mathcal F$ est
$f$-plat au point $x$.
\end{theorem}

\begin{proof}
On peut se borner au cas où $f$ est de présentation finie, la question étant
locale sur $X$. Nous procéderons en plusieurs étapes.

\medskip
\noindent\textbf{I)} \emph{Réduction au cas où $A$ et $A'$ sont des anneaux
locaux intégralement clos.}

Comme $u$ est injectif et $A$ intègre, il existe un idéal premier
$\mathfrak p'$ de $A'$ tel que $u^{-1}(\mathfrak p')=0$
$(\mathbf 0_{\mathrm I},1.5.8)$~; l'homomorphisme composé
$A\xrightarrow{u}A'\to A''=A'/\mathfrak p'$ est donc injectif et local, et si
$Y''=\operatorname{Spec}(A'')$, $X''=X'\otimes_{A'}A''=X\times_Y Y''$,
$\mathcal F''=\mathcal F'\otimes_{A'}A''$, $\mathcal F''$ est $Y''$-plat aux
points de $X''$ au-dessus de $x'$ (2.1.4)~; remplaçant au besoin $A'$ par $A''$
et tenant compte de (\textbf I, 3.4.7), on peut donc supposer d'abord que $A'$
est intègre. Si $K'$ est le corps des fractions de $A'$, il y a alors un anneau
de valuation $B'$ dans $K'$ qui domine $A'$~; l'homomorphisme composé
$A\xrightarrow{u}A'\to B'$ étant injectif et local, le même raisonnement que
% IV3_B23_P154_END

\oldpage[IV-3]{155}

précédemment permet de remplacer $A'$ par $B'$~; on peut donc supposer l'anneau
local $A'$ intégralement clos, $A$ étant un sous-anneau local de $A'$ dominé par
$A'$. Soit $A_1$ la clôture intégrale de $A$~; il est clair que
$A\subset A_1\subset A'$, et par hypothèse $A_1$ est un anneau local~; si
$\mathfrak m$, $\mathfrak m_1$, $\mathfrak m'$ sont les idéaux maximaux de
$A$, $A_1$, $A'$, on a $\mathfrak m\subset\mathfrak m_1\subset\mathfrak m'$~;
en effet, $\mathfrak m_1$ est le \emph{seul} idéal premier de $A_1$ au-dessus de
$\mathfrak m$, puisque $A_1$ est un anneau local (Bourbaki,
\emph{Alg. comm.}, chap.~V, \S~2, n$^\circ$~1, prop.~1)~; comme
$\mathfrak m'\cap A=\mathfrak m$, on a
$\mathfrak m'\cap A_1\cap A=\mathfrak m$, donc
$\mathfrak m'\cap A_1=\mathfrak m_1$. Posons
$Y_1=\operatorname{Spec}(A_1)$, $X_1=X\times_Y Y_1$, $f_1=f_{(Y_1)}$,
$\mathcal F_1=\mathcal F\otimes_Y Y_1$, et soit $x_1$ la projection de $x'$ dans
$X_1$~; désignons d'autre part par $y_1$ l'unique point fermé de $Y_1$, de sorte
que $y_1=f_1(x_1)$. Par hypothèse le morphisme
$\operatorname{Spec}(k(y_1))\to\operatorname{Spec}(k(y))$ est radiciel, d'où
l'on conclut, par la transitivité des fibres (\textbf I, 3.6.4) et
(\textbf I, 3.5.7), que le morphisme $f_1^{-1}(y_1)\to f^{-1}(y)$ est radiciel,
et en particulier que $x_1$ est le \emph{seul} point de $X_1$ dont les
projections dans $X$ et $Y_1$ soient $x$ et $y_1$ respectivement~; en outre, on
a vu que $y_1$ est le seul point de $Y_1$ dont la projection dans $Y$ soit $y$,
donc $x_1$ est le seul point de $X_1$ dont la projection dans $X$ soit $x$. Si
l'on prouve que $\mathcal F_1$ est $f_1$-plat au point $x_1$, on pourra donc
appliquer (11.5.5), d'où résultera la conclusion. On est donc ramené au cas où
$A$ lui aussi est \emph{intégralement clos}.

\medskip
\noindent\textbf{II)} \emph{Réduction au cas où $A$ et $A'$ sont des anneaux
locaux de $\mathbf Z$-algèbres de type fini intégralement clos.}

On peut considérer $A'$ comme limite inductive filtrante de ses
sous-$\mathbf Z$-algèbres (intègres) de type fini $B'_{\lambda}$~; en outre,
comme $A'$ est intégralement clos et que la clôture intégrale d'une
$\mathbf Z$-algèbre de type fini est aussi une $\mathbf Z$-algèbre de type fini
(7.8.3), on voit que $A'$ est limite inductive de ses sous-$\mathbf Z$-algèbres
de type fini $B'_{\lambda}$ \emph{intégralement closes}~; si
$\mathfrak m'_{\lambda}=\mathfrak m'\cap B'_{\lambda}$, $A'$ est aussi limite
inductive des sous-anneaux locaux
$(B'_{\lambda})_{\mathfrak m'_{\lambda}}=A'_{\lambda}$ dominés par $A'$
(5.13.3). Pour tout $\lambda$, $B'_{\lambda}\cap A$ est aussi limite inductive
de ses sous-$\mathbf Z$-algèbres de type fini $B_{\alpha\lambda}$, donc
$A=\varinjlim_{\alpha,\lambda}B_{\alpha\lambda}$, et comme plus haut on peut
remplacer $B_{\alpha\lambda}$ dans cette formule par sa clôture intégrale
(contenue par hypothèse dans $B'_{\lambda}\cap A$), puis par l'anneau local
$A_{\alpha\lambda}=(B_{\alpha\lambda})_{\mathfrak m_{\alpha\lambda}}$, où
$\mathfrak m_{\alpha\lambda}=\mathfrak m\cap B_{\alpha\lambda}
=\mathfrak m'_{\lambda}\cap B_{\alpha\lambda}$, de sorte que
$A_{\alpha\lambda}$ est dominé par $A'_{\lambda}$ et par $A$. Posons
$Y_{\alpha\lambda}=\operatorname{Spec}(A_{\alpha\lambda})$~; il résulte de
(8.9.1) qu'il existe un couple $(\alpha,\lambda)$ assez grand, un morphisme
$f_{\alpha\lambda}:X_{\alpha\lambda}\to Y_{\alpha\lambda}$ de type fini et un
$\mathcal O_{X_{\alpha\lambda}}$-Module cohérent $\mathcal F_{\alpha\lambda}$
tels que
\[
X=X_{\alpha\lambda}\times_{Y_{\alpha\lambda}}Y,
\qquad f=f_{\alpha\lambda}\times1_Y,
\qquad \mathcal F=\mathcal F_{\alpha\lambda}\otimes_{Y_{\alpha\lambda}}Y~;
\]
si $x_{\alpha\lambda}$ est la projection de $x$ dans $X_{\alpha\lambda}$, il
suffira de montrer que $\mathcal F_{\alpha\lambda}$ est
$f_{\alpha\lambda}$-plat au point $x_{\alpha\lambda}$. Comme $A'$ est limite
inductive des $A'_{\mu}$ pour $\mu\geq\lambda$, $X'$ est limite projective des
$X_{\alpha\lambda}\times_{Y_{\alpha\lambda}}Y'_{\mu}=X'_{\mu}$, et on a aussi
$\mathcal F'=\mathcal F'_{\alpha\mu}\otimes_{Y'_{\mu}}Y'$, où
$\mathcal F'_{\alpha\mu}
=\mathcal F_{\alpha\lambda}\otimes_{Y_{\alpha\lambda}}Y'_{\mu}$.
Appliquant (11.2.6), on voit qu'on peut prendre $\mu$ assez grand pour que
$\mathcal F'_{\alpha\mu}$ soit $Y'_{\mu}$-plat au point $x'_{\mu}$, projection
de $x'$ dans $X'_{\mu}$, et en outre, par construction des $A'_{\mu}$, la
projection de $x'_{\mu}$ dans $Y'_{\mu}$ est le point fermé de $Y'_{\mu}$.

\medskip
\noindent\textbf{III)} \emph{Réduction au cas où le corps résiduel $k'$ de
$A'$ est une extension finie et radicielle du corps résiduel $k$ de $A$.}

On peut en premier lieu répéter le raisonnement de la partie I) de la
démonstration de (11.4.11), tenant compte du fait que $\mathbf Z$ est un anneau
de
% IV3_B23_P155_END

\oldpage[IV-3]{156}

Jacobson~; on se réduit ainsi au cas où $k'$ est une extension \emph{finie} de
$k$, ce que l'on va supposer dans ce qui suit. Soient $k''$ la plus grande
extension séparable de $k$ contenue dans $k'$, $k_1$ une extension galoisienne
finie de $k$ contenant $k''$, de sorte que $k''\otimes_k k_1$ est composée
directe de corps isomorphes à $k_1$~; comme $k'$ est une extension radicielle de
$k''$, $k'\otimes_k k_1$ est donc composée directe d'extensions radicielles de
$k_1$. Il existe un anneau local $A_1$ qui est une $A$-algèbre et un $A$-module
\emph{libre de type fini}, tel que $A_1/\mathfrak m A_1$ soit $k$-isomorphe à
$k_1$ $(\mathbf 0_{\mathrm{III}},10.3.1.2)$~; de façon précise, on peut supposer
que $k_1=k[T]/(r)$, où $r$ est un polynôme unitaire irréductible et séparable de
$k[T]$ de degré $n$~; si $R$ est un polynôme unitaire de $A[T]$ dont l'image
canonique est $r$ (et qui est donc de degré $n$), on peut prendre
$A_1=A[T]/(R)$. Or, si $K$ est le corps des fractions de $A$, il est clair que
$R$ est un polynôme irréductible et séparable de $K[T]$~; on en déduit donc
d'abord que $A_1$ est un anneau intègre dont le corps des fractions
$K_1=K[T]/(R)$ est une extension séparable de $K$. En outre, si $t$ est l'image
canonique de $T$ dans $A_1$, les $t^j$ $(0\leq j<n)$ forment une base du
$A$-module $A_1$, et leurs images dans $k_1$ une base sur $k$~; on déduit de là
que
\[
d=\det\bigl(\operatorname{Tr}_{A_1/A}(t^{i+j})\bigr)
\]
est un élément de $A$ dont la classe dans $k$ est $\ne0$, et qui est par suite
inversible. Le même raisonnement que dans (6.12.4.1, I)) prouve alors que le
morphisme $\operatorname{Spec}(A_1)\to\operatorname{Spec}(A)$ est plat et a ses
fibres régulières~; on conclut par suite de (6.5.4, (ii)) que $A_1$ est
\emph{intégralement clos}. Considérons alors l'anneau
$A'_1=A'\otimes_A A_1$~; c'est un $A'$-module libre de type fini, donc un anneau
semi-local (Bourbaki, \emph{Alg. comm.}, chap.~IV, \S~2, n$^\circ$~5, cor.~3 de
la prop.~9)~; en outre, les idéaux maximaux de cette $A'$-algèbre finie sont
tous au-dessus de l'idéal maximal $\mathfrak m'$ de $A'$, et \emph{a fortiori}
contiennent $\mathfrak m A'_1$. Mais
\[
A'_1/\mathfrak m A'_1=(A'/\mathfrak m A')\otimes_k k_1,
\]
et comme $k_1$ est une extension séparable et finie de $k$, le radical de
$A'_1/\mathfrak m A'_1$ est égal à
$(\mathfrak m'/\mathfrak m A')\otimes_k k_1$ (Bourbaki, \emph{Alg.},
chap.~VIII, \S~7, n$^\circ$~2, cor.~2 de la prop.~3)~; si $\mathfrak n_i$
$(1\leq i\leq r)$ sont les idéaux maximaux de $A'_1$, les corps
$A'_1/\mathfrak n_i$ sont donc les corps composant l'algèbre
$k'\otimes_k k_1$, autrement dit ce sont des extensions finies
\emph{radicielles} de $k_1$. D'ailleurs, comme $A\to A'$ est un homomorphisme
injectif, il en est de même de $A_1\to A'_1$, $A_1$ étant un $A$-module plat~;
l'homomorphisme canonique
\[
A'_1\longrightarrow\prod_{i=1}^{r}(A'_1)_{\mathfrak n_i}
\]
étant lui aussi injectif (Bourbaki, \emph{Alg. comm.}, chap.~II, \S~3,
n$^\circ$~3, th.~1), il en est de même du composé de $A_1$ dans ce produit.
Mais $A_1$ est intègre, et les noyaux des homomorphismes
$A_1\to(A'_1)_{\mathfrak n_i}$ sont en nombre fini~; comme leur intersection
est nulle, c'est que l'un d'eux déjà est nul. En d'autres termes, il y a un
$B_1=(A'_1)_{\mathfrak n_i}$ tel que l'homomorphisme $A_1\to B_1$ soit
\emph{injectif} et \emph{local}. Posons $Y'_1=\operatorname{Spec}(B_1)$,
$X'_1=X'\times_{Y'}Y'_1$~; $\mathcal F'_1=\mathcal F'\otimes_{Y'}Y'_1$ est
$Y'_1$-plat en tous les points de $X'_1$ au-dessus de $x'$~; d'ailleurs l'idéal
maximal de $B_1$ est le seul au-dessus de $\mathfrak m'$, par suite tous ces
points ont pour projection dans $Y'_1$ le point fermé $y'_1$. Soit $x'_1$ un de
ces points. Posons d'autre part $Y_1=\operatorname{Spec}(A_1)$,
$X_1=X\times_Y Y_1$, $\mathcal F_1=\mathcal F\otimes_Y Y_1$~; si $x_1$ est la
projection de $x'_1$ dans $X_1$, la projection de $x_1$ dans $X$ est $x$ et sa
projection dans $Y_1$ est le point fermé $y_1$. Si l'on prouve que
$(\mathcal F_1)_{x_1}$ est un $\mathcal O_{y_1}$-module plat, il en résultera
que $\mathcal F_x$ est un $\mathcal O_y$-module plat~; en effet
$\mathcal O_{y_1}$ est un $\mathcal O_y$-module plat, donc
$(\mathcal F_1)_{x_1}$ est un $\mathcal O_y$-module plat
$(\mathbf 0_{\mathrm I},6.2.1)$.
% IV3_B23_P156_END

\oldpage[IV-3]{157}

Mais $(\mathcal F_1)_{x_1}=\mathcal F_x\otimes_{\mathcal O_x}\mathcal O_{x_1}$,
et $\mathcal O_{x_1}$ est un $\mathcal O_x$-module \emph{fidèlement plat}~;
donc (2.2.11, (iii)) $\mathcal F_x$ est un $\mathcal O_y$-module plat. Comme
$X'_1=X_1\times_{Y_1}Y'_1$, $\mathcal F'_1=\mathcal F_1\otimes_{Y_1}Y'_1$,
on est bien ramené à la situation de l'énoncé (11.6.1), avec $A$ et $A'$
remplacés par $A_1$ et $B_1$ respectivement.

\medskip
\noindent\textbf{IV)} \emph{Fin de la démonstration.} --- On est finalement
réduit à prouver (11.6.1) sous les hypothèses supplémentaires suivantes~:
\begin{enumerate}[label=(\roman*)]
  \item $A$ et $A'$ sont des anneaux locaux de $\mathbf Z$-algèbres de type fini
  (donc des anneaux \emph{excellents} (7.8.3))~;
  \item $A$ est intégralement clos~;
  \item le corps résiduel $k'$ de $A'$ est une extension finie et radicielle du
  corps résiduel $k$ de $A$.
\end{enumerate}
On sait alors ((7.8.3) et (2.3.8)) que sous les conditions (i) et (ii), si
$\mathfrak m$ et $\mathfrak m'$ sont les idéaux maximaux de $A$ et $A'$
respectivement, la topologie $\mathfrak m$-adique sur $A$ est induite par la
topologie $\mathfrak m'$-adique de $A'$. Le complété
$\widehat u:\widehat A\to\widehat{A'}$ est donc un homomorphisme injectif.
D'autre part, comme le morphisme
$\operatorname{Spec}(k')\to\operatorname{Spec}(k)$ est radiciel, il en est de
même du morphisme $f'^{-1}(y')\to f^{-1}(y)$ (\textbf I, 3.5.7), et il n'y a
donc qu'un seul point $x'$ de $X'$ dont les projections dans $X$ et $Y'$ soient
$x$ et $y'$ respectivement. On peut donc appliquer le résultat de (11.5.2).
C.Q.F.D.
\end{proof}

\begin{corollary}[11.6.2]
\label{IV.11.6.2-fr}
Soient $A$ un anneau local intègre unibranche, $Y=\operatorname{Spec}(A)$,
$f:X\to Y$ un morphisme localement de présentation finie, $\mathcal F$ un
$\mathcal O_X$-Module quasi-cohérent de présentation finie. Soient $A'$ un
anneau local, $A\to A'$ un homomorphisme local injectif~; on pose
$Y'=\operatorname{Spec}(A')$, $X'=X\times_Y Y'$, $f'=f_{(Y')}$,
$\mathcal F'=\mathcal F\otimes_Y Y'$. Soit $x$ un point de $X$ dont la
projection $f(x)=y$ est le point fermé de $Y$~; on suppose que $\mathcal F'$ est
$f'$-plat en tous les points $x'$ de $X'$ dont les projections dans $X$ et $Y'$
sont respectivement $x$ et le point fermé $y'$ de $Y'$. Alors $\mathcal F$ est
$f$-plat au point $x$.
\end{corollary}

\begin{proof}
On peut en effet reprendre la partie I) de la démonstration de (11.6.1), qui
prouve (avec les mêmes notations) que si $\mathcal F_1$ est $f_1$-plat en tous
les points $x_1$ de $X_1$ dont les projections respectives dans $X$ et $Y_1$
sont $x$ et $y_1$, alors $\mathcal F$ est $f$-plat au point $x$~; on est ainsi
ramené au cas où $A$ est intégralement clos, donc géométriquement unibranche, et
la conclusion résulte alors de (11.6.1).
\end{proof}

\subsection{Contre-exemples}
\label{subsection:IV.11.7-fr}

\begin{env}[11.7.1]
\label{IV.11.7.1-fr}
Considérons d'abord le cas où $A$ est un anneau local \emph{artinien}, et où les
hypothèses de (11.4.11) sont remplies \emph{sauf} la condition (ii) concernant
le corps résiduel $k$ de $A$. Nous allons voir que la conclusion de (11.4.11)
peut alors être en défaut. Soit $k$ un corps admettant une extension
\emph{galoisienne} $k'$ de degré $[k':k]>1$, et désignons par $\Gamma$ le groupe
de Galois de $k'$. Soit $A$ une $k$-algèbre ayant une base de 3 éléments
$1,a,b$ avec la table de multiplication $a^2=ab=ba=b^2=0$, de sorte que $A$ est
un anneau local artinien dont l'idéal maximal $\mathfrak m=ka+kb$ est de carré
nul. Soit $A'=A\otimes_k k'$, qui est une $k'$-algèbre de base $1,a,b$, anneau
local artinien d'idéal maximal
$\mathfrak m'=k'a+k'b=\mathfrak m A'$, de carré nul~; $A$ s'identifie
canoniquement à un sous-anneau de $A'$. Soit $\mathfrak J$ le sous-$k'$-espace
vectoriel de $\mathfrak m'$ engendré par $a+\gamma b$, où $\gamma\in k'$
n'appartient pas à $k$~; il est clair que $\mathfrak J$ est un idéal de $A'$.
Posons $B=A'/\mathfrak J$~; c'est un anneau artinien qui est un $A$-module
\emph{non plat}~; sinon (Bourbaki, \emph{Alg. comm.}, chap.~II, \S~3,
n$^\circ$~2, cor.~2 de la prop.~5), $B$ serait un $A$-module \emph{libre}~;
comme $A'$ est lui aussi un $A$-module libre, et que l'homomorphisme canonique
$A'/\mathfrak m A'\to B/\mathfrak m B$ est bijectif, l'homomorphisme canonique
$A'\to B=A'/\mathfrak J$ serait lui
% IV3_B23_P157_END

\oldpage[IV-3]{158}

aussi bijectif (\emph{loc. cit.}, n$^\circ$~2, cor. de la prop.~6) ce qui est
absurde. Autrement dit, si l'on pose $Y=\operatorname{Spec}(A)$,
$X=\operatorname{Spec}(B)$, $\mathcal O_X$ n'est pas $Y$-plat en l'unique point
$x$ de $X$. Mais soit $A_1=B$, et $Y_1=\operatorname{Spec}(A_1)$, et considérons
l'homomorphisme canonique $A\to A_1$, qui est local et injectif puisque
$\mathfrak J\cap A=0$~; si
$B_1=B\otimes_A A_1=B\otimes_A B$, nous allons voir qu'il y a un point de
$X_1=X\times_Y Y_1=\operatorname{Spec}(B_1)$ où $\mathcal O_{X_1}$ est
$Y_1$-plat. Pour cela, remarquons que l'on a
\[
B\otimes_A B
=\frac{A'\otimes_A A'}
{\operatorname{Im}(\mathfrak J\otimes_A A')
+\operatorname{Im}(A'\otimes_A\mathfrak J)}.
\]
Or la structure de $C'=A'\otimes_A A'$ s'obtient aisément~; on considère la
$A$-algèbre produit $C=\prod_{\sigma\in\Gamma}A'_{\sigma}$, où tous les
$A'_{\sigma}$ sont égaux à $A'$, et l'application canonique
$\varphi:A'\otimes_A A'\to C$ telle que
$\varphi(x\otimes y)=(\sigma(x)y)_{\sigma\in\Gamma}$ (le groupe $\Gamma$
opérant canoniquement sur $A'$)~; en passant aux quotients, on en déduit un
homomorphisme $C'/\mathfrak m C'\to C/\mathfrak m C$ qui n'est autre que
l'homomorphisme canonique
$k'\otimes_k k'\to\prod_{\sigma}k'_{\sigma}$ (avec $k'_{\sigma}=k'$ pour tout
$\sigma\in\Gamma$)~; on sait que ce dernier est bijectif (Bourbaki,
\emph{Alg.}, chap.~VIII, \S~8, prop.~4), donc il en est de même de $\varphi$,
puisque $C'$ et $C$ sont des $A$-modules libres (Bourbaki,
\emph{Alg. comm.}, chap.~II, \S~3, n$^\circ$~2, cor. de la prop.~6). De ce qui
précède, il résulte que $B\otimes_A B$ est une $A$-algèbre semi-locale, composée
directe des $A$-algèbres locales $A'/(\sigma(\mathfrak J)+\mathfrak J)$. Celle
de ces algèbres correspondant à l'identité de $\Gamma$ est isomorphe à
$A'/\mathfrak J$, donc est un $A_1$-module plat, mais comme $\gamma\notin k$, il
y a au moins un $\sigma\in\Gamma$ tel que
$\sigma(\mathfrak J)\ne\mathfrak J$~; alors
$\sigma(\mathfrak J)+\mathfrak J=\mathfrak m'$, et $A'/\mathfrak m'$ n'est pas
un $(A'/\mathfrak J)$-module plat, puisque c'est le quotient de
$A'/\mathfrak J$ par un idéal non nul.
\end{env}

\begin{env}[11.7.2]
\label{IV.11.7.2-fr}
Nous allons maintenant montrer que le résultat de (11.5.4) perd sa validité
lorsqu'on ne suppose plus que $f$ soit un morphisme \emph{propre} (et
\emph{a fortiori} (11.5.3) cesse d'être exact lorsqu'on ne suppose plus $g$
propre). Soient $k$ un corps, $A_0$ l'anneau de polynômes $k[S,T]$, $A$
l'anneau quotient $A_0/A_0ST$~; $Y=\operatorname{Spec}(A)$ est donc la courbe
réductible formée des deux «~axes de coordonnées~» dans le plan affine
$V_k^2=\operatorname{Spec}(A_0)$. L'anneau $A$ admet deux idéaux premiers
minimaux $\mathfrak p_1=A_0S/A_0ST$, $\mathfrak p_2=A_0T/A_0ST$, et comme $A$
est réduit, il se plonge canoniquement dans $B=B_1\oplus B_2$, où
$B_1=A/\mathfrak p_1$, $B_2=A/\mathfrak p_2$~; d'ailleurs $B_1$ s'identifie à
$k[T]$ et $B_2$ à $k[S]$, donc ce sont des anneaux intègres intégralement clos
et par suite $Z=\operatorname{Spec}(B)$ n'est autre que le normalisé du
préschéma $Y$ par rapport à $R(Y)$ (\textbf{II}, 6.3.8), somme des deux schémas
$Z_1=\operatorname{Spec}(B_1)$, $Z_2=\operatorname{Spec}(B_2)$. Désignons par
$y$ le «~point double~» de $Y$, correspondant à l'idéal maximal
$\mathfrak p_1+\mathfrak p_2=\mathfrak m$ de $A$, par $z_1$ et $z_2$ les
points de $Z$ qui se projettent en $y$, correspondant respectivement aux idéaux
maximaux $\mathfrak m_1=(T)$ et $\mathfrak m_2=(S)$ de $B_1$ et $B_2$. Nous
désignerons par $X$ le sous-préschéma de $Z$ induit sur le complémentaire de
$z_2$ dans $Z$~; on a donc
$X=\operatorname{Spec}(B_1\oplus(B_2)_S)$~; il est immédiat que
l'homomorphisme $A\to A'=B_1\oplus(B_2)_S$ est injectif, mais le morphisme
correspondant $f:X\to Y$ n'est pas fermé (car $f(Z_2-\{z_2\})$ n'est pas fermé
dans $Y$, bien que $Z_2-\{z_2\}$ soit fermé dans $X$)~; \emph{a fortiori} il
n'est pas propre. Nous allons voir maintenant que $f$ n'est pas plat au point
$z_1$~; il suffira de montrer $(\mathbf 0_{\mathrm I},6.6.2)$ que
$\mathcal O_{z_1}$ n'est pas un $\mathcal O_y$-module fidèlement plat, et pour
cela il suffit de voir que l'homomorphisme canonique
$\mathcal O_y\to\mathcal O_{z_1}$ n'est pas injectif~; mais cela est immédiat
car $\mathcal O_{z_1}$ est un anneau intègre, tandis que $\mathcal O_y$ admet
des diviseurs de zéro. Cependant, la première projection
$p:X\times_Y X\to X$ est un isomorphisme~: en effet, on a
\[
B_1\otimes_A B_1=(A/\mathfrak p_1)\otimes_A(A/\mathfrak p_1)=A/\mathfrak p_1,
\]
$(B_2)_S\otimes_A(B_2)_S=(B_2\otimes_A B_2)_S=(B_2)_S$ pour la même raison,
et enfin $B_1\otimes_A(B_2)_S=0$, car l'image canonique de $S$ dans $B_1$ est
nulle.
\end{env}

\begin{env}[11.7.3]
\label{IV.11.7.3-fr}
L'exemple précédent peut se généraliser~: on considère sur un corps $k$ une
courbe algébrique réduite $Y$ admettant un seul «~point double ordinaire~» $y$
(notion qui sera définie plus tard de façon générale), et sa normalisée $Z$, de
sorte que le morphisme $g:Z\to Y$ est fini, que la restriction de $g$ à
$Z-g^{-1}(y)$ est un isomorphisme sur $Y-\{y\}$, et que $g^{-1}(y)$ se réduit à
deux points «~simples~» $z_1,z_2$~; en outre le préschéma $g^{-1}(y)$ est somme
des deux préschémas $\operatorname{Spec}(k(z_1))$,
$\operatorname{Spec}(k(z_2))$, canoniquement isomorphes à
$\operatorname{Spec}(k(y))$. Soit $X$ le sous-préschéma de $Z$ induit sur
l'ouvert $Z-\{z_2\}$~; le morphisme $f:X\to Y$, restriction de $g$ à $X$, n'est
pas propre, sans quoi (\textbf{II}, 5.4.3) il en serait de même de l'injection
canonique $j:X\to Z$, qui n'est pas fermée. Le morphisme $f$ est radiciel, car
pour tout $y'\in Y$, la fibre $f^{-1}(y')$ ne comprend qu'un seul point $x'$,
$k(y')\to k(x')$ étant un isomorphisme~; on en conclut d'abord que le morphisme
diagonal $\Delta_f:X\to X\times_Y X$ est bijectif (1.7.7.1) et d'autre part,
comme $f$ est non ramifié (17.4.2, d$'$), $\Delta_f$ est une immersion ouverte
(17.4.2, b$'$)~; par suite $\Delta_f$ est un isomorphisme, et la première
projection $p:X\times_Y X\to X$ l'isomorphisme réciproque. Cependant $f$ n'est
pas plat au point $z_1$~; sinon $\mathcal O_{z_1}$ serait un
$\mathcal O_y$-module fidèlement plat $(\mathbf 0_{\mathrm I},6.6.2)$, et comme
$\mathcal O_y$ contient deux idéaux premiers minimaux $\mathfrak p_1$,
$\mathfrak p_2$ distincts (correspondant aux deux «~branches~» de $Y$ au point
$y$) il existerait dans $\mathcal O_{z_1}$ deux idéaux premiers dont les images
réciproques par $u:\mathcal O_y\to\mathcal O_{z_1}$ seraient
$\mathfrak p_1$ et $\mathfrak p_2$ $(\mathbf 0_{\mathrm I},6.5.1)$~; mais cela
est absurde, car $\mathcal O_{z_1}$ n'a que deux idéaux premiers distincts, $0$
et l'idéal maximal $\mathfrak m_1$, et $u^{-1}(\mathfrak m_1)$ est l'idéal
maximal $\mathfrak m$ de $\mathcal O_y$.
\end{env}

\begin{env}[11.7.4]
\label{IV.11.7.4-fr}
On notera que dans l'exemple précédent l'homomorphisme $u$ est \emph{injectif}
lorsque $Y$ est \emph{irréductible} (on peut par exemple prendre
$Y=\operatorname{Spec}(k[S,T]/(S(S^2+T^2)-(S^2-T^2)))$, «~cubique à point
double~»)~; on peut en effet alors (en remplaçant $Y$ par un voisinage affine de
$y$) supposer que $Y=\operatorname{Spec}(A)$, où $A$ est intègre, d'où
$Z=\operatorname{Spec}(B)$, où $B$ est la clôture intégrale de $A$~; comme $B$,
$\mathcal O_y$ et $\mathcal O_{z_1}$ s'identifient alors à des sous-anneaux du
corps
% IV3_B23_P158_END

\oldpage[IV-3]{159}

des fractions de $A$, $u$ est évidemment injectif. On notera par contre que
l'homomorphisme $\widehat u:\widehat{\mathcal O}_y\to
\widehat{\mathcal O}_{z_1}$ n'est pas injectif, car
$\widehat{\mathcal O}_{z_1}$ est un anneau intègre local ($z_1$ étant un point
simple), tandis que $\widehat{\mathcal O}_y$ a deux idéaux premiers minimaux
distincts (correspondant aux deux «~branches~» de $Y$) et admet donc des
diviseurs de zéro. Cela donne un exemple montrant que dans l'énoncé (11.5.2), on
ne peut pas remplacer l'hypothèse que $\widehat u$ est injectif par l'hypothèse
que $u$ lui-même est injectif, même lorsque $A'$ est un anneau local. Il suffit
en effet de prendre (avec les notations précédentes) $A=\mathcal O_y$,
$A'=\mathcal O_{z_1}$, $Y=\operatorname{Spec}(A)$,
$X=\operatorname{Spec}(A')$~; le raisonnement de (11.7.3) prouve encore que la
première projection $p:X\times_Y X\to X$ est un isomorphisme, bien que $f$ ne
soit pas plat au point $z_1$.
\end{env}

\begin{env}[11.7.5]
\label{IV.11.7.5-fr}
Les exemples de (11.7.2) et (11.7.3) expliquent la restriction aux anneaux
locaux \emph{unibranches} dans (11.6.1) et (11.6.2). Nous allons voir maintenant
que dans (11.6.1), on ne peut affaiblir l'hypothèse sur $A$ en supposant
seulement $A$ unibranche. Considérons en effet l'anneau local intègre complet
$A=\mathbf R[[U,V]]/(U^2+V^2)$ qui est unibranche mais non géométriquement
unibranche (6.5.11). On sait (\emph{loc. cit.}) que si $u,v$ sont les images de
$U$ et $V$ dans $A$, la clôture intégrale de $A$ est $\overline A=A[t]$ avec
$t=v/u$, tel que $t^2=-1$, si bien que $\overline A$ est isomorphe à
$\mathbf C[[U]]$. Posons $Y=\operatorname{Spec}(A)$,
$X=\operatorname{Spec}(\overline A)$ (normalisé de $Y$ (\textbf{II}, 6.3.8)) et
soient $y$ et $x$ les points fermés de $Y$ et $X$ respectivement~; nous allons
montrer que pour une $A$-algèbre locale convenable $A'$, si l'on pose
$Y'=\operatorname{Spec}(A')$, $X'=X\times_Y Y'$, et si $y'$ désigne le point
fermé de $Y'$, $\mathcal O_{X'}$ est $Y'$-plat en un point de $X'$ dont les
projections dans $X$ et $Y'$ sont respectivement $x$ et $y'$, mais n'est pas
$Y'$-plat en tous les points ayant ces projections~; il en résultera évidemment
(2.1.4) que $\mathcal O_X$ n'est pas $Y$-plat au point $x$ (ce qui est d'ailleurs
trivial \emph{a priori}, $\overline A$ n'étant pas un $A$-module libre).

Soit $B=A\otimes_{\mathbf R}\mathbf C$, isomorphe à
$\mathbf C[[U,V]]/(U+iV)(U-iV)$~; $B$ a deux idéaux premiers minimaux
$\mathfrak p'$, $\mathfrak p''$ engendrés respectivement par $u+iv$ et $u-iv$,
et $\mathfrak n=\mathfrak p'+\mathfrak p''$ est l'idéal maximal de l'anneau
local complet $B$. Soit $\overline B=\overline A\otimes_{\mathbf R}\mathbf C$~;
$\overline B$ est composée directe de deux algèbres isomorphes à
$\mathbf C[[U]]$, engendrées par les idempotents
$e'=(1\otimes1+t\otimes i)/2$ et $e''=(1\otimes1-t\otimes i)/2$~; comme
l'homomorphisme $A\to\overline A$ est injectif, il en est de même de
$B\to\overline B$, et les images de $u+iv$ et de $u-iv$ par cette injection sont
respectivement $ue'$ et $ue''$~; on en conclut aussitôt que $\overline B$
s'identifie canoniquement à $(B/\mathfrak p')\oplus(B/\mathfrak p'')$. Cela
étant, prenons pour $A'$ l'anneau local $B/\mathfrak p'$~; alors
$\overline A\otimes_A A'$ s'identifie à $\overline B\otimes_B A'$. Mais on a
$(B/\mathfrak p')\otimes_B(B/\mathfrak p')=B/\mathfrak p'$ et
$(B/\mathfrak p')\otimes_B(B/\mathfrak p'')=B/\mathfrak n$, donc
$\overline A\otimes_A A'$ est isomorphe à $A'\oplus(B/\mathfrak n)$. Cela
établit notre assertion, car $B/\mathfrak n=A'/(\mathfrak n/\mathfrak p')$
n'est pas un $A'$-module plat (sans quoi ce serait un $A'$-module libre
(Bourbaki, \emph{Alg. comm.}, chap.~II, \S~3, n$^\circ$~2, cor.~2 de la
prop.~5), ce qui est absurde).
\end{env}

\subsection{Un critère valuatif de platitude}
\label{subsection:IV.11.8-fr}

\begin{theorem}[11.8.1]
\label{IV.11.8.1-fr}
Soient $f:X\to Y$ un morphisme localement de présentation finie, $\mathcal F$ un
$\mathcal O_X$-Module quasi-cohérent de présentation finie, $x$ un point de
$X$, $y=f(x)$. On suppose l'anneau local $\mathcal O_y$ intègre (resp. réduit et
noethérien). Pour que $\mathcal F$ soit $f$-plat au point $x$, il faut et il
suffit que, pour tout anneau de valuation (resp. tout anneau de valuation
discrète) $A'$ et tout homomorphisme local $\mathcal O_y\to A'$, la condition
suivante soit satisfaite~: en posant $Y'=\operatorname{Spec}(A')$,
$X'=X\times_Y Y'$, $f'=f_{(Y')}$, le $\mathcal O_{X'}$-Module
$\mathcal F'=\mathcal F\otimes_Y Y'$ est $f'$-plat en tous les points $x'$ de
$X'$ dont les projections respectives dans $X$ et $Y'$ sont $x$ et le point
fermé $y'$ de $Y'$.
\end{theorem}

\begin{proof}
La condition étant évidemment nécessaire (2.1.4), reste à prouver qu'elle est
suffisante. On peut évidemment ((\textbf I, 3.6.5) et (\textbf I, 2.4.4)) se
borner au cas où $Y=\operatorname{Spec}(A)$ est le spectre d'un anneau local $A$
et $y$ le point fermé de $Y$.

\medskip
\noindent\textbf{(i)} \emph{Cas où $A$ est intègre.} --- Soient $K$ le corps
des fractions de $A$, $A_1$ la clôture intégrale de $A$~; si l'on pose
$Y_1=\operatorname{Spec}(A_1)$, $X_1=X\times_Y Y_1$, $f_1=f_{(Y_1)}$, il suffit,
en vertu de (11.5.5), de montrer que
$\mathcal F_1=\mathcal F\otimes_Y Y_1$ est $f_1$-plat en \emph{tous} les points
$x_1$ de $X_1$ dont $x$ est la projection dans $X$. Or, si $f_1(x_1)=y_1$, on a
$\mathfrak i_{y_1}=\mathfrak m_1$, où $\mathfrak m_1$ est un idéal premier de
$A_1$ dont $\mathfrak m=\mathfrak i_y$ est la trace sur $A$ ($\mathfrak m_1$
est d'ailleurs nécessairement un idéal maximal). Soit alors $A'$ un anneau de
valuation pour $K$ qui domine $\mathcal O_{y_1}=(A_1)_{\mathfrak m_1}$~;
l'homomorphisme $A\to\mathcal O_{y_1}$ étant local, il en est de même de
$A\to A'$. Il existe alors au moins un
% IV3_B23_P159_END

\oldpage[IV-3]{160}

point $x'\in X'$ dont les projections dans $X_1$ et $Y'$ sont respectivement
$x_1$ et $y'$ (\textbf I, 3.4.7)~; comme par hypothèse $\mathcal F'$ est
$f'$-plat au point $x'$, et que $\mathcal O_{y_1}$ est intégralement clos, on
peut appliquer (11.6.1), et on en déduit que $\mathcal F_1$ est $f_1$-plat au
point $x_1$, d'où le théorème dans ce cas.

\medskip
\noindent\textbf{(ii)} \emph{Cas où $A$ est réduit et noethérien.} --- Soient
$\mathfrak p_i$ $(1\leq i\leq m)$ les idéaux minimaux de $A$~; comme $A$ est
réduit, il s'identifie canoniquement à un sous-anneau du produit des
$A_i=A/\mathfrak p_i$, qui sont des anneaux locaux noethériens~; posant
$Y_i=\operatorname{Spec}(A_i)$, $X_i=X\times_Y Y_i$, $f_i=f_{(Y_i)}$, il
résulte de (11.5.2) qu'il suffit de démontrer que pour chaque $i$,
$\mathcal F_i=\mathcal F\otimes_Y Y_i$ est $f_i$-plat en tout point $x_i$ de
$X_i$ dont les projections dans $X$ et $Y_i$ sont $x$ et le point fermé $y_i$
de $Y_i$ respectivement. Or, comme $A_i$ est un anneau local noethérien intègre,
il existe dans son corps des fractions $K_i$ un sous-anneau $A''_i$ contenant
$A_i$, qui soit une $A_i$-algèbre \emph{finie} (donc un anneau semi-local
noethérien) et dont les anneaux locaux soient \emph{géométriquement
unibranches} ((6.15.5) et $(\mathbf 0,23.2.5)$). Comme les idéaux maximaux de
$A''_i$ sont alors nécessairement au-dessus de l'idéal maximal de $A_i$, on
déduit encore de (\textbf I, 3.4.7) et de (11.5.2) qu'il suffit, en posant
$Y''_i=\operatorname{Spec}(A''_i)$, $X''_i=X\times_Y Y''_i$,
$f''_i=f_{(Y''_i)}$, de prouver que
$\mathcal F''_i=\mathcal F\otimes_Y Y''_i$ est $f''_i$-plat en tout point
$x''_i$ de $X''_i$ dont les projections dans $X$ et $Y''_i$ sont $x$ et le
point fermé $y''_i$ de $Y''_i$ respectivement. Or, soit $A'$ un anneau de
valuation discrète pour $K_i$ dominant $A''_i$, et soit $x'$ un point de $X'$
dont les projections dans $X''_i$ et dans $Y'$ sont $x''_i$ et $y'$
respectivement (\textbf I, 3.4.7)~; comme $\mathcal O_{y''_i}$ est
géométriquement unibranche, on peut encore appliquer (11.6.1) et on en déduit
bien que $\mathcal F''_i$ est $f''_i$-plat au point $x''_i$.
\end{proof}

\begin{remarks}[11.8.2]
\label{IV.11.8.2-fr}
\begin{enumerate}[label=(\roman*)]
  \item Dans l'énoncé de (11.8.1), on peut se borner à supposer que la condition
  sur $\mathcal F'$ est vérifiée pour les anneaux de valuation $A'$ \emph{complets}
  et dont le corps résiduel est \emph{algébriquement clos}. On sait en effet que
  tout anneau de valuation $A'$ est dominé par un tel anneau $A''$
  (\textbf{II}, 7.1.2), et que si $A'$ est un anneau de valuation discrète, on
  peut supposer qu'il en est de même de $A''$
  $(\mathbf 0_{\mathrm{III}},10.3.1)$.
  \item La démonstration de (11.8.1) se simplifie lorsqu'on suppose non seulement
  que $A$ est intègre et noethérien, mais que son complété $\widehat A$ est lui
  aussi \emph{intègre}. Remplaçant $X$ par $X\otimes_A\widehat A$ et raisonnant
  comme dans la démonstration de (11.5.3), on peut en effet se ramener alors à
  prouver (11.8.1) lorsque $A=\mathcal O_y$ est intègre, noethérien et
  \emph{complet}. Or, on sait (\textbf{II}, 7.1.7) qu'un tel anneau $A$ est
  dominé par un anneau de valuation discrète complet~; la conclusion résulte
  donc directement de (11.5.2).
\end{enumerate}
\end{remarks}

\subsection{Familles séparantes et universellement séparantes d'homomorphismes
de faisceaux de modules}
\label{subsection:IV.11.9-fr}

\begin{env}[11.9.1]
\label{IV.11.9.1-fr}
Soient $X$ un préschéma, $(f_{\lambda})_{\lambda\in L}$ une famille de morphismes
$f_{\lambda}:Z_{\lambda}\to X$, $\mathcal F$ un $\mathcal O_X$-Module
quasi-cohérent~; pour tout $\lambda\in L$, supposons donnés un
$\mathcal O_{Z_{\lambda}}$-Module quasi-cohérent $\mathcal G_{\lambda}$ et un
homomorphisme
\begin{equation}
u_{\lambda}:\mathcal F\longrightarrow(f_{\lambda})_*(\mathcal G_{\lambda}).
\label{IV.11.9.1.1-fr}
\tag{11.9.1.1}
\end{equation}
\end{env}
% IV3_B23_P160_END

\oldpage[IV-3]{161}

On dit que la famille $(u_\lambda)$ (ou la famille correspondante des
$u^\#_\lambda:(f_\lambda)^*(\mathcal F)\to\mathcal G_\lambda$) est
\emph{séparante} si l'intersection des noyaux des $u_\lambda$ est
\emph{nulle}. En d'autres termes, cela signifie que pour \emph{tout} ouvert
$U$ de $X$, et toute section $t$ de $\mathcal F$ au-dessus de $U$, telle que,
pour \emph{tout} $\lambda$, la section $u_\lambda(t)$ (qui, par définition est
une section de $\mathcal G_\lambda$ au-dessus de $f_\lambda^{-1}(U)$) soit
nulle, alors $t$ est elle-même nulle.

\begin{env}[11.9.2]
\label{IV.11.9.2-fr}
Avec les notations de (11.9.1), soit $M$ un second ensemble d'indices~; pour
tout $\lambda\in L$, soit $(g_{\lambda\mu})_{\mu\in M}$ une famille de
morphismes $g_{\lambda\mu}:Z'_{\lambda\mu}\to Z_\lambda$~; pour tout couple
$(\lambda,\mu)$, supposons donnés un $\mathcal O_{Z'_{\lambda\mu}}$-Module
quasi-cohérent $\mathcal H_{\lambda\mu}$ et un homomorphisme
$v_{\lambda\mu}:\mathcal G_\lambda\to
(g_{\lambda\mu})_*(\mathcal H_{\lambda\mu})$~; posons
$h_{\lambda\mu}=f_\lambda\circ g_{\lambda\mu}:Z'_{\lambda\mu}\to X$ et
considérons l'homomorphisme composé
\[
w_{\lambda\mu}:\mathcal F\xrightarrow{u_\lambda}
(f_\lambda)_*(\mathcal G_\lambda)
\xrightarrow{(f_\lambda)_*(v_{\lambda\mu})}
(h_{\lambda\mu})_*(\mathcal H_{\lambda\mu}).
\]
Supposons que, pour tout $\lambda\in L$, la famille
$(v_{\lambda\mu})_{\mu\in M}$ soit séparante~: alors, il en est de même de la
famille des $(f_\lambda)_*(v_{\lambda\mu})$ ($\mu\in M$), comme on le voit
aussitôt. On en conclut que, pour que la famille $(u_\lambda)$ soit séparante,
il faut et il suffit que la famille
$(w_{\lambda\mu})_{(\lambda,\mu)\in L\times M}$ le soit.
\end{env}

\begin{env}[11.9.3]
\label{IV.11.9.3-fr}
Pour vérifier que la famille $(u_\lambda)$ est séparante (avec les notations de
(11.9.1)), on peut évidemment se ramener tout d'abord au cas où $X$ est
\emph{affine}, la propriété étant locale sur $X$. On peut en outre supposer que
$Z_\lambda=X$ \emph{pour tout} $\lambda\in L$. En effet, soit $M$ un ensemble
d'indices, somme d'une famille $(M_\lambda)_{\lambda\in L}$, et pour tout
$\lambda\in L$, soit $(Y_{\lambda\mu})_{\mu\in M_\lambda}$ un recouvrement
ouvert affine de $Z_\lambda$~; soit
$j_{\lambda\mu}:Y_{\lambda\mu}\to Z_\lambda$ l'injection canonique et posons
$\mathcal H_{\lambda\mu}=j_{\lambda\mu}^*(\mathcal G_\lambda)
=\mathcal G_\lambda|Y_{\lambda\mu}$. Si l'on considère l'homomorphisme
canonique
$v_{\lambda\mu}=\rho_{\mathcal G_\lambda}:\mathcal G_\lambda\to
(j_{\lambda\mu})_*(j_{\lambda\mu}^*(\mathcal G_\lambda))
=(j_{\lambda\mu})_*(\mathcal H_{\lambda\mu})$ relatif à $j_{\lambda\mu}$
$(\mathbf 0_{\mathrm I},4.4.3.2)$, il est immédiat que pour chaque
$\lambda\in L$, la famille $(v_{\lambda\mu})_{\mu\in M_\lambda}$ est
séparante. En vertu de (11.9.2), on est donc ramené à prouver que la famille des
homomorphismes composés $((f_\lambda)_*(v_{\lambda\mu}))\circ u_\lambda$ est
séparante, autrement dit on est ramené au cas où les $Z_\lambda$ sont affines.
Mais alors les $(f_\lambda)_*(\mathcal G_\lambda)$ sont des
$\mathcal O_X$-Modules quasi-cohérents (\textbf I, 1.6.2) et en vertu de la
définition, on peut remplacer les $Z_\lambda$ par $X$ et les
$\mathcal G_\lambda$ par les $(f_\lambda)_*(\mathcal G_\lambda)$, d'où notre
assertion.

On notera en outre que si $L$ est \emph{fini} et les $f_\lambda$
\emph{quasi-compacts}, on peut, dans la réduction précédente, supposer que les
$M_\lambda$ sont aussi \emph{finis}, donc on est dans ce cas ramené à vérifier
qu'une famille \emph{finie} d'homomorphismes de $\mathcal F$ dans des
$\mathcal O_X$-Modules quasi-cohérents est séparante.
\end{env}

\begin{env}[11.9.4]
\label{IV.11.9.4-fr}
Considérons donc le cas où $Z_\lambda=X$ pour tout $\lambda$, et où
$X=\operatorname{Spec}(A)$ est affine~; alors on a
$\mathcal F=\widetilde M$ et $\mathcal G_\lambda=\widetilde N_\lambda$, où
$M$ et $N_\lambda$ sont des $A$-modules, et $u_\lambda=\widetilde\varphi_\lambda$,
où les $\varphi_\lambda:M\to N_\lambda$ sont des $A$-homomorphismes. Dire que
la famille $(u_\lambda)$ est séparante signifie alors que, pour tout $s\in A$,
l'intersection des noyaux des
$(\varphi_\lambda)_s:M_s\to(N_\lambda)_s$ est réduite à 0. On dit alors aussi
que la famille $(\varphi_\lambda)$ est \emph{séparante}. On notera que si $L$
est \emph{fini}, il revient au même de dire que l'intersection des noyaux des
$\varphi_\lambda$ est 0, car on a alors
\[
\bigcap_{\lambda\in L}\operatorname{Ker}((\varphi_\lambda)_s)
=\left(\bigcap_{\lambda\in L}\operatorname{Ker}(\varphi_\lambda)\right)_s
\quad(\mathbf 0_{\mathrm I},1.3.2).
\]
Mais cette relation n'est plus exacte en général lorsque $L$ est \emph{infini},
et le fait que l'intersection des noyaux des $\varphi_\lambda$ soit 0
\emph{n'entraîne donc pas},
\oldpage[IV-3]{162}
en général, que la famille $(\varphi_\lambda)$ soit séparante. Par exemple,
supposons que $A$ soit un anneau de valuation discrète, d'idéal maximal
$\mathfrak m$, et considérons la famille des homomorphismes
$\varphi_k:A\to A/\mathfrak m^k$, dont l'intersection des noyaux
$\mathfrak m^k$ est réduite à 0~; cette famille n'est toutefois pas séparante,
car les fibres de tous les $\mathfrak m^k$ au point générique $x$ de
$X=\operatorname{Spec}(A)$ (qui est ouvert dans $X$) sont égales à
$\mathcal O_x=k(x)$, corps des fractions de $A$, et leur intersection n'est
donc pas réduite à 0.
\end{env}

\begin{env}[11.9.5]
\label{IV.11.9.5-fr}
Nous allons principalement nous occuper dans ce qui suit du problème du
\emph{changement de base} pour les familles séparantes. Les notations étant
celles de (11.9.1), considérons un morphisme $g:X'\to X$ et posons
$\mathcal F'=\mathcal F\otimes_{\mathcal O_X}\mathcal O_{X'}
=g^*(\mathcal F)$, et, pour tout $\lambda$,
$Z'_\lambda=Z_\lambda\times_X X'$, $f'_\lambda=f_\lambda\times 1$,
$\mathcal G'_\lambda=\mathcal G_\lambda\otimes_{\mathcal O_{Z_\lambda}}
\mathcal O_{Z'_\lambda}=(g'_\lambda)^*(\mathcal G_\lambda)$, où
$g'_\lambda:Z'_\lambda\to Z_\lambda$ est la projection canonique. Pour tout
$\lambda$, on désigne alors par
$u'_\lambda:\mathcal F'\to(f'_\lambda)_*(\mathcal G'_\lambda)$
l'homomorphisme obtenu comme suit~: soit
\[
\psi:(f_\lambda)_*(\mathcal G_\lambda)\longrightarrow
(f_\lambda)_*((g'_\lambda)_*((g'_\lambda)^*(\mathcal G_\lambda)))
=g_*((f'_\lambda)_*(\mathcal G'_\lambda))
\]
l'homomorphisme $(f_\lambda)_*(\rho_{\mathcal G_\lambda})$, où
$\rho_{\mathcal G_\lambda}$ est l'homomorphisme canonique
$(\mathbf 0_{\mathrm I},4.4.3.2)$ correspondant à $g'_\lambda$. Alors
$u'_\lambda$ est défini comme le composé
\[
g^*(\mathcal F)\xrightarrow{g^*(u_\lambda)}
g^*((f_\lambda)_*(\mathcal G_\lambda))
\xrightarrow{g^*(\psi)}g^*g_*((f'_\lambda)_*(\mathcal G'_\lambda))
\xrightarrow{\sigma}(f'_\lambda)_*(\mathcal G'_\lambda)
\]
où $\sigma=\sigma_{(f'_\lambda)_*(\mathcal G'_\lambda)}$ est
l'homomorphisme canonique $(\mathbf 0_{\mathrm I},4.4.3.3)$ correspondant à
$g$. On dira pour abréger que $u'_\lambda$ est \emph{déduit de $u_\lambda$
par le changement de base} $g$. Lorsque $f_\lambda:Z_\lambda\to X$ est un
morphisme \emph{affine}, on a
$(f'_\lambda)_*(\mathcal G'_\lambda)=g^*((f_\lambda)_*(\mathcal G_\lambda))$,
et dans ce cas on a donc simplement $u'_\lambda=g^*(u_\lambda)$.

On peut aussi interpréter $u'_\lambda$ de la façon suivante~: il suffit de
connaître la valeur de $u'_\lambda(t')$, lorsque $t'$ est une section de
$\mathcal F'$ au-dessus d'un ouvert $U'$ de $X'$, du type particulier
suivant~: $t'$ est la restriction à $U'$ de l'image canonique par
$\Gamma(U,\mathcal F)\to\Gamma(g^{-1}(U),\mathcal F')$ d'une section $t$ de
$\mathcal F$ au-dessus d'un ouvert $U$ de $X$ contenant $g(U')$ (ces sections
engendrant en effet le $\mathcal O_{X'}$-Module $\mathcal F'$
$(\mathbf 0_{\mathrm I},3.7.1)$). Considérons la section $u_\lambda(t)$ de
$\mathcal G_\lambda$ au-dessus de $f_\lambda^{-1}(U)$, et son image canonique
$t''$ par
$\Gamma(f_\lambda^{-1}(U),\mathcal G_\lambda)\to
\Gamma((g'_\lambda)^{-1}(f_\lambda^{-1}(U)),\mathcal G'_\lambda)$~; alors
$u'_\lambda(t')$ est la restriction de $t''$ à $(f'_\lambda)^{-1}(U')$.
Considérons en particulier le cas où $Z_\lambda$ est un sous-préschéma induit
sur un ouvert de $X$, $\mathcal G_\lambda=j_\lambda^*(\mathcal F)$, où
$j_\lambda:Z_\lambda\to X$ est l'injection canonique, et où $u_\lambda$ est
l'homomorphisme canonique
$\rho_{\mathcal F}:\mathcal F\to(j_\lambda)_*(j_\lambda^*(\mathcal F))
=(j_\lambda)_*(\mathcal G_\lambda)$ $(\mathbf 0_{\mathrm I},4.4.3.2)$
correspondant à $j_\lambda$. Alors $Z'_\lambda$ est induit sur un ouvert de
$X'$, et l'interprétation précédente montre que $u'_\lambda$ n'est autre que
l'homomorphisme canonique
$\rho_{\mathcal F'}:\mathcal F'\to
(j'_\lambda)_*((j'_\lambda)^*(\mathcal F'))
=(j'_\lambda)_*(\mathcal G'_\lambda)$ correspondant à l'injection canonique
$j'_\lambda:Z'_\lambda\to X'$.
\end{env}

\begin{env}[11.9.6]
\label{IV.11.9.6-fr}
Sous les conditions de (11.9.5), supposons que $X$ et $X'$ soient affines, et
que l'on veuille prouver que pour toute section $t'$ de $\mathcal F'$ au-dessus
de $X'$, dont les images par tous les $u'_\lambda$ sont nulles, alors $t'$ est
elle-même nulle. Alors, on peut aussi se borner au cas où $Z_\lambda=X$ pour
tout $\lambda\in L$. En effet, avec les notations de (11.9.3) et de (11.9.5),
si l'on pose $Y'_{\lambda\mu}=(g'_\lambda)^{-1}(Y_{\lambda\mu})$,
l'homomorphisme $v'_{\lambda\mu}$ déduit de $v_{\lambda\mu}$ par le changement
de base $g$ n'est autre que l'homomorphisme canonique
\oldpage[IV-3]{163}
$\rho_{\mathcal G'_\lambda}:\mathcal G'_\lambda\to
(j'_{\lambda\mu})_*((j'_{\lambda\mu})^*(\mathcal G'_\lambda))$
correspondant à l'injection canonique
$j'_{\lambda\mu}:Y'_{\lambda\mu}\to Z'_\lambda$, comme on l'a vu dans
(11.9.5). L'assertion résulte alors des raisonnements de (11.9.2) et (11.9.3),
$Y_{\lambda\mu}$ et $v_{\lambda\mu}$ étant remplacés par
$Y'_{\lambda\mu}$ et $v'_{\lambda\mu}$.

On a une réduction semblable lorsque l'on veut prouver que la famille
$(u'_\lambda)$ est \emph{séparante} ($X$ et $X'$ étant affines)~: cela résulte
encore de (11.9.2) et (11.9.3).
\end{env}

\begin{proposition}[11.9.7]
\label{IV.11.9.7-fr}
Avec les notations de (11.9.1) et (11.9.5), supposons
$X=\operatorname{Spec}(A)$ et $X'=\operatorname{Spec}(A')$ affines, et
supposons en outre que $A'$ soit un $A$-module projectif. Alors, si la famille
$(u_\lambda)$ est séparante, toute section $t'$ de $\mathcal F'$ au-dessus de
$X'$ dont les images par tous les $u'_\lambda$ sont nulles, est elle-même
nulle.
\end{proposition}

\begin{proof}
On a vu (11.9.6) qu'on peut se borner au cas où tous les $Z_\lambda$ sont égaux
à $X$. La proposition est alors conséquence du lemme suivant~:
\end{proof}

\begin{lemma}[11.9.7.1]
\label{IV.11.9.7.1-fr}
Soient $A$ un anneau, $(M_\lambda)_{\lambda\in L}$ une famille de $A$-modules,
$M$ un $A$-module et pour chaque $\lambda$,
$u_\lambda:M\to M_\lambda$ un homomorphisme. Supposons que l'intersection des
noyaux des $u_\lambda$ soit réduite à 0. Alors, pour tout $A$-module projectif
$N$, l'intersection des noyaux des homomorphismes
$u_\lambda\otimes 1:M\otimes_A N\to M_\lambda\otimes_A N$ est réduite à 0.
\end{lemma}

\begin{proof}
En effet, $N$ est facteur direct d'un $A$-module libre $L$, et il suffit
évidemment de prouver que l'intersection des noyaux des homomorphismes
$v_\lambda:M\otimes_A L\to M_\lambda\otimes_A L$ est réduite à 0, puisque
$u_\lambda\otimes 1:M\otimes_A N\to M_\lambda\otimes_A N$ est la restriction
de $v_\lambda$. Mais l'assertion résulte alors trivialement de l'hypothèse.
\end{proof}

\begin{remark}[11.9.8]
\label{IV.11.9.8-fr}
Nous ignorons si, sous les hypothèses de la proposition (11.9.7), la famille
$(u'_\lambda)$ est \emph{séparante}~: il faudrait en effet (11.9.4) prouver
qu'une section $t'$ de $\mathcal F'$ \emph{au-dessus d'un ouvert}
$D(h')\subset X'$ (où $h'\in A'$) telle que les $u'_\lambda(t')$ soient toutes
nulles, est elle-même nulle. Or, on ne peut appliquer la proposition (11.9.4)
à $D(h')=\operatorname{Spec}(A'_{h'})$, car du fait que $A'$ soit un
$A$-module projectif (même libre), il ne résulte pas que $A'_{h'}$ soit un
$A$-module projectif. Par exemple, on peut prendre pour $A$ un anneau de
valuation discrète, pour $A'$ un anneau de valuation discrète qui soit un
$A$-module libre de rang 2, et pour $A'_{h'}$ le corps des fractions de $A'$.
\end{remark}

On a toutefois le résultat suivant~:

\begin{corollary}[11.9.9]
\label{IV.11.9.9-fr}
Soient $X$ un préschéma artinien, $g:X'\to X$ un morphisme plat (on notera que
ces deux conditions sont satisfaites si $X$ est le spectre d'un \emph{corps}
et $g$ un morphisme \emph{quelconque}). Alors, avec les notations de (11.9.1)
et (11.9.5), si la famille $(u_\lambda)$ est séparante, il en est de même de la
famille $(u'_\lambda)$.
\end{corollary}

\begin{proof}
On peut évidemment se borner au cas où $X=\operatorname{Spec}(A)$ est le
spectre d'un anneau local artinien $A$ (\textbf I, 6.2.2)~; on note ensuite que
pour tout ouvert affine $U'=\operatorname{Spec}(A')$ de $X'$, $A'$ est un
$A$-module plat, donc \emph{projectif} $(\mathbf 0_{\mathrm{III}},10.1.3)$.
Il suffit donc d'appliquer (11.9.7) à tout ouvert affine de $X'$ pour obtenir
le corollaire.
\end{proof}

\begin{theorem}[11.9.10]
\label{IV.11.9.10-fr}
Soient $X$ un préschéma, $(u_\lambda)$ une famille d'homomorphismes
(11.9.1.1), $g:X'\to X$ un morphisme, $(u'_\lambda)$ la famille
d'homomorphismes déduite de $(u_\lambda)$ par le changement de base $g$
(11.9.5).
\begin{enumerate}[label=(\roman*)]
  \item Si $g$ est un morphisme fidèlement plat et si la famille
  $(u'_\lambda)$ est séparante, alors la famille $(u_\lambda)$ est séparante.
  \oldpage[IV-3]{164}
  \item Supposons que $g$ soit un morphisme plat, et en outre que l'une des
  deux conditions suivantes soit vérifiée~:
  \begin{enumerate}[label=\emph{\alph*)}]
    \item $L$ est fini et les $f_\lambda$ sont quasi-compacts.
    \item Le morphisme $g$ est localement de présentation finie.
  \end{enumerate}
  Alors, si la famille $(u_\lambda)$ est séparante, il en est de même de la
  famille $(u'_\lambda)$.
\end{enumerate}
\end{theorem}

\begin{proof}
\begin{enumerate}[label=(\roman*)]
  \item En vertu de (2.2.8), il suffit de montrer que si une section $t$ de
  $\mathcal F$ au-dessus d'un ouvert $U$ de $X$ appartient au noyau de chacun
  des $u_\lambda$, son image $t'$ par l'homomorphisme canonique
  \[
  \Gamma(\rho):\Gamma(U,\mathcal F)\longrightarrow
  \Gamma(g^{-1}(U),g^*(\mathcal F))
  =\Gamma(U,g_*(g^*(\mathcal F)))
  \]
  est nulle. Or les images de $t'$ par les $g^*(u_\lambda)$ sont les images
  des $u_\lambda(t)$ par l'homomorphisme
  \[
  \Gamma(U,(f_\lambda)_*(\mathcal G_\lambda))\longrightarrow
  \Gamma(g^{-1}(U),g^*((f_\lambda)_*(\mathcal G_\lambda))),
  \]
  donc sont nulles, et \emph{a fortiori} on a $u'_\lambda(t')=0$ pour tout
  $\lambda$, donc $t'=0$ par hypothèse, ce qui prouve (i).

  \item La question étant locale sur $X$ et $X'$, on peut se borner au cas où
  $X=\operatorname{Spec}(A)$ et $X'=\operatorname{Spec}(A')$ sont affines,
  $A'$ étant un $A$-module plat, et à prouver que, pour toute section $z'$ de
  $\mathcal F'$ au-dessus de $X'$ dont les images par tous les $u'_\lambda$
  sont nulles, alors $z'$ est elle-même nulle. On a vu en outre (11.9.6) que
  l'on peut alors supposer que $Z_\lambda=X$ pour tout $\lambda$.

  Distinguons maintenant les deux cas.

  \emph{a)} Si $L$ est fini et les $f_\lambda$ quasi-compacts, on a vu (11.9.3)
  qu'on peut encore se ramener au cas où $Z_\lambda=X$ pour tout
  $\lambda\in L$, et où en outre $L$ est fini. Il revient au même alors de dire
  que l'intersection des noyaux des
  $u_\lambda:\mathcal F\to\mathcal G_\lambda$ est nulle, ou que
  l'homomorphisme
  $u=(u_\lambda):\mathcal F\to\mathcal G=\bigoplus_\lambda\mathcal G_\lambda$
  est injectif. Comme $u'=g^*(u)$ est injectif puisque $g$ est plat, la
  proposition est démontrée dans ce cas.

  \emph{b)} Soient $\mathcal F=\widetilde M$,
  $\mathcal G_\lambda=\widetilde N_\lambda$, et posons
  $M'=M\otimes_A A'$, $N'_\lambda=N_\lambda\otimes_A A'$, de sorte que
  $\mathcal F'=\widetilde{M'}$, $\mathcal G'_\lambda=\widetilde{N'_\lambda}$~;
  par abus de langage, nous noterons encore $u_\lambda$ l'homomorphisme
  $M\to N_\lambda$, et $u'_\lambda$ l'homomorphisme
  $u_\lambda\otimes 1:M'\to N'_\lambda$. Donnons-nous un élément $z'\in M'$
  tel que $u'_\lambda(z')=0$ pour tout $\lambda$~; il s'agit de prouver que
  l'on a $z'=0$. Or, l'hypothèse que $g$ est plat et de présentation finie
  entraîne, d'après (11.3.15), qu'il existe une suite finie
  $(s_i)_{1\leq i\leq n}$ d'éléments de $A$, telle que, si l'on pose
  $\mathfrak J_k=s_1A+\cdots+s_kA$, et
  $A_i=A_{s_i}/\mathfrak J_{i-1}A_{s_i}$, l'anneau
  $A'_i=A'\otimes_A A_i$ soit un $A_i$-module \emph{libre} pour
  $1\leq i\leq n$, et $\mathfrak J_n=A$. La proposition sera établie si nous
  prouvons pour $1\leq i\leq n$ l'assertion suivante~:
  \begin{enumerate}[label=$(*_i)$]
    \item \emph{Il existe un entier $m_i>0$ tel que $s_j^{m_i}z'=0$ pour
    $j\leq i$.}
  \end{enumerate}
  En effet, posant alors $k=m_n$ et notant que les $s_i^k$ ($1\leq i\leq n$)
  engendrent aussi l'idéal unité de $A$, l'assertion $(*_n)$ montrera que $z'$,
  combinaison linéaire des $s_i^kz'$, est nul.

  Prouvons $(*_i)$ par récurrence sur $i$, l'assertion étant vide pour $i=0$.
  Supposons donc $i>0$ et soit $m$ un multiple commun des $m_j$ pour
  $j\leq i-1$. Remarquons que (pour $1\leq h\leq n$) si
  $\mathfrak J'_h$ est l'idéal engendré par les $s_j^m$ ($0\leq j\leq h$),
  $\mathfrak J_h/\mathfrak J'_h$ est nilpotent~; remplacer les $s_j$ par
  $s_j^m$ pour $1\leq j\leq n$ revient donc à remplacer, pour
  $1\leq h\leq n$, $A_h$ par
  $A_{s_h}/\mathfrak J'_{h-1}A_{s_h}=B_h$, de sorte que
  $A_h=B_h/(\mathfrak J_{h-1}/\mathfrak J'_{h-1})B_h$~; comme $A'$ est un
  $A$-module plat, il résulte de $(\mathbf 0_{\mathrm{III}},10.1.2)$ que
  $B'_h=A'\otimes_A B_h$ est encore un $B_h$-module libre. On peut donc
  \oldpage[IV-3]{165}
  remplacer tous les $s_j$ ($1\leq j\leq n$) par $s_j^m$ sans changer les
  propriétés des $\mathfrak J_h$ et des $A'_h$, et supposer par la suite que
  $m=1$. Alors $z'$, étant annulé par $\mathfrak J_{i-1}$, s'identifie à un
  élément de
  $\operatorname{Hom}_{A'}(A'/\mathfrak J_{i-1}A',M')$, et comme
  $\mathfrak J_{i-1}A'$ est un idéal de type fini de $A'$, ce module
  d'homomorphismes s'identifie lui-même à
  $\operatorname{Hom}_A(A/\mathfrak J_{i-1},M)\otimes_A A'$
  $(\mathbf 0_{\mathrm I},6.2.2)$. Soit
  \[
  v_\lambda=\operatorname{Hom}(1,u_\lambda):
  \operatorname{Hom}_A(A/\mathfrak J_{i-1},M)\longrightarrow
  \operatorname{Hom}_A(A/\mathfrak J_{i-1},N_\lambda)
  \]
  l'homomorphisme déduit de $u_\lambda$~; la famille $(v_\lambda)$ est elle
  aussi \emph{séparante}. En effet, pour tout $t\in A$, on a
  \[
  (\operatorname{Hom}_A(A/\mathfrak J_{i-1},M))_t
  =\operatorname{Hom}_{A_t}(A_t/(\mathfrak J_{i-1})_t,M_t)
  \]
  et de même en remplaçant $M$ par $N_\lambda$, puisque l'idéal
  $\mathfrak J_{i-1}$ est de type fini $(\mathbf 0_{\mathrm I},1.3.5)$~;
  comme par hypothèse l'intersection des noyaux des $(u_\lambda)_t$ est nulle,
  il en est de même de l'intersection des noyaux des
  $(v_\lambda)_t=\operatorname{Hom}(1,(u_\lambda)_t)$, d'où notre assertion
  (11.9.4). Remplaçant $A$ par $A/\mathfrak J_{i-1}$, $M$ par
  $\operatorname{Hom}_A(A/\mathfrak J_{i-1},M)$, $N_\lambda$ par
  $\operatorname{Hom}_A(A/\mathfrak J_{i-1},N_\lambda)$ (qui sont des
  $(A/\mathfrak J_{i-1})$-modules), $u_\lambda$ par $v_\lambda$ et enfin $A'$
  par $A'/\mathfrak J_{i-1}A'$, on voit qu'on peut se ramener au cas où, dans
  la situation initiale, l'élément $s=s_i\in A$ est tel que $A'_s$ soit un
  $A_s$-module libre. Or la famille des
  $(u_\lambda)_s:M_s\to(N_\lambda)_s$ est séparante par hypothèse~; comme on a
  $(u'_\lambda)_s(z'/1)=0$, il résulte de (11.9.7) que l'on a $z'/1=0$ dans
  $M'_s$~; mais cela signifie qu'il existe un entier $r$ tel que $s^rz'=0$
  dans $M'$, ce qui achève de prouver $(*_i)$ par récurrence.
\end{enumerate}
\end{proof}

\begin{remark}[11.9.11]
\label{IV.11.9.11-fr}
Bornons-nous pour simplifier au cas où $Z_\lambda=X$ pour tout $\lambda$. Il
faut noter alors que si la famille des homomorphismes
$u_\lambda:\mathcal F\to\mathcal G_\lambda$ est séparante, \emph{il ne
s'ensuit pas nécessairement} que, pour tout $x\in X$, l'intersection des
noyaux des homomorphismes
$(u_\lambda)_x:\mathcal F_x\to(\mathcal G_\lambda)_x$ soit réduite à 0. Par
exemple, soient $X$ un préschéma de Jacobson localement noethérien, de
dimension $\geq1$, et $\mathcal F$ un $\mathcal O_X$-Module cohérent~; pour
tout point fermé $x$ de $X$, et tout entier $n\geq0$,
$\mathfrak m_x^n\mathcal F$ est un $\mathcal O_X$-Module cohérent de support
contenu dans $\{x\}$. La famille des homomorphismes canoniques
$\mathcal F\to\mathcal F/\mathfrak m_x^n\mathcal F$ (où $x$ parcourt
l'ensemble $X_0$ des points fermés de $X$ et $n$ l'ensemble des entiers
$\geq0$) est \emph{séparante}~: en effet, si $t$ est une section de
$\mathcal F$ au-dessus d'un ouvert $U$ de $X$ dont les images dans les
$\Gamma(U,\mathcal F/\mathfrak m_x^n\mathcal F)$ sont toutes nulles, il en
résulte aussitôt que pour tout point fermé $x\in U$, on a $t_x=0$~; comme
l'ensemble des points fermés contenus dans $U$ est très dense dans $U$, cela
entraîne bien $t=0$ (10.2.1). Cependant, si l'on prend
$\mathcal F=\mathcal O_X$, et si $y\in X$ est un point \emph{non fermé} de
$X$, on a $(\mathcal O_X/\mathfrak m_x^n\mathcal O_X)_y=0$ pour tout point
fermé $x$ de $X$, mais $\mathcal O_{X,y}\neq0$, et l'intersection des noyaux
des homomorphismes
$\mathcal O_{X,y}\to(\mathcal O_X/\mathfrak m_x^n\mathcal O_X)_y$ est égale
à $\mathcal O_{X,y}$.
\end{remark}

\begin{lemma}[11.9.12]
\label{IV.11.9.12-fr}
Les notations étant celles de (11.9.1) et (11.9.5), supposons la famille
$(u_\lambda)$ séparante~; supposons en outre que $X$ soit un $S$-préschéma, où
$S=\operatorname{Spec}(A)$ est un schéma affine, et que
$X'=X\times_S S'$, où $S'=\operatorname{Spec}(A')$, $A'$ étant une
$A$-algèbre~; supposons enfin que les $\mathcal G_\lambda$ soient $S$-plats.
Soit $(A'_\alpha)_{\alpha\in I}$ la famille filtrante des sous-$A$-algèbres de
type fini de $A'$~; pour tout $\alpha\in I$, posons
$S''_\alpha=\operatorname{Spec}(A'_\alpha)$,
$X''_\alpha=X\times_S S''_\alpha$,
$Z''_{\alpha\lambda}=Z_\lambda\times_X X''_\alpha$, et soient
$f''_{\alpha\lambda}:Z''_{\alpha\lambda}\to X''_\alpha$,
$\mathcal F''_\alpha$, $\mathcal G''_{\alpha\lambda}$ et
$u''_{\alpha\lambda}$ les morphismes, Modules et homomorphismes de Modules
déduits de $f_\lambda$, $\mathcal F$, $\mathcal G_\lambda$ et $u_\lambda$ au
moyen du changement de base $X''_\alpha\to X$. Alors, si, pour tout
$\alpha\in I$, la famille $(u''_{\alpha\lambda})_{\lambda\in L}$ est
séparante, il en est de même de $(u'_\lambda)_{\lambda\in L}$.
\end{lemma}

\begin{proof}
Il s'agit de prouver que si une section $t'$ de $\mathcal F'$ au-dessus d'un
ouvert affine $U'$ de $X'$ est telle que $u'_\lambda(t')=0$ pour tout
$\lambda\in L$, on a $t'=0$. Si $h_\alpha:X'\to X''_\alpha$ est la projection
\oldpage[IV-3]{166}
canonique, il résulte de (8.2.11) qu'il existe un indice $\alpha$ et un ouvert
quasi-compact $U''_\alpha\subset X''_\alpha$ tels que
$U'=h_\alpha^{-1}(U''_\alpha)$~; en outre, en vertu de (8.5.2, (i)), on peut
supposer qu'il y a une section $t''_\alpha$ de $\mathcal F''_\alpha$ au-dessus
de $U''_\alpha$ telle que $t'$ soit l'image canonique de $t''_\alpha$. Quitte à
remplacer $S$, $X$, $Z_\lambda$, $f_\lambda$, $\mathcal F$,
$\mathcal G_\lambda$ et $u_\lambda$ par $S''_\alpha$, $U''_\alpha$,
$(f''_{\alpha\lambda})^{-1}(U''_\alpha)$ et les restrictions correspondantes
de $\mathcal F''_\alpha$, $\mathcal G''_{\alpha\lambda}$ et
$u''_{\alpha\lambda}$, on peut donc supposer que $U'=X'$, que $t'$ est
l'image canonique d'une section $t$ de $\mathcal F$ au-dessus de $X$ et que
l'homomorphisme $A\to A'$ est injectif, ou encore, si $p:S'\to S$ est le
morphisme correspondant, que l'homomorphisme
$\mathcal O_S\to p_*(\mathcal O_{S'})$ intervenant dans la définition de $p$
est injectif. Il en résulte aussitôt en vertu de la platitude de
$\mathcal G_\lambda$ sur $S$ (et en se ramenant par exemple au cas où
$Z_\lambda$ est affine sur $S$ (\textbf I, 1.6.3 et 1.6.5)) que
l'homomorphisme canonique
$\rho:\mathcal G_\lambda\to(g'_\lambda)_*(\mathcal G'_\lambda)$ est lui aussi
injectif. Mais l'homomorphisme composé
\[
\Gamma(X,\mathcal F)\xrightarrow{\Gamma(u_\lambda)}
\Gamma(Z_\lambda,\mathcal G_\lambda)\xrightarrow{\Gamma(\rho)}
\Gamma(Z'_\lambda,\mathcal G'_\lambda)
\]
est égal par définition à l'homomorphisme composé
\[
\Gamma(X,\mathcal F)\xrightarrow{\Gamma(\rho)}
\Gamma(X',\mathcal F')\xrightarrow{\Gamma(u'_\lambda)}
\Gamma(Z'_\lambda,\mathcal G'_\lambda)
\]
donc l'image de $t$ par ces homomorphismes composés est
$u'_\lambda(t')=0$~; en vertu de l'injectivité de l'homomorphisme
$\Gamma(Z_\lambda,\mathcal G_\lambda)\xrightarrow{\Gamma(\rho)}
\Gamma(Z'_\lambda,\mathcal G'_\lambda)$ on en conclut que $u_\lambda(t)=0$
pour tout $\lambda\in L$, d'où $t=0$ par hypothèse, et finalement $t'=0$.
\end{proof}

\begin{proposition}[11.9.13]
\label{IV.11.9.13-fr}
Les notations étant celles de (11.9.1) et (11.9.5), supposons que $X$ soit un
préschéma sur un corps $k$, et que, en posant $S=\operatorname{Spec}(k)$, on
ait $X'=X\times_S S'$, où $S'$ est un $k$-préschéma quelconque. Alors, si la
famille $(u_\lambda)_{\lambda\in L}$ est séparante, il en est de même de
$(u'_\lambda)$.
\end{proposition}

\begin{proof}
On peut se borner au cas où $S'=\operatorname{Spec}(A')$ est affine. Si $A'$
est une $k$-algèbre de type fini, le morphisme $g:X'\to X$ est plat et de
présentation finie, et on est donc dans les conditions d'application de
(11.9.10, (ii), b)). Dans le cas général, on considère $A'$ comme limite
inductive de ses $k$-sous-algèbres $A'_\alpha$ de type fini, et on applique à
chaque $A'_\alpha$ le résultat de (11.9.10, (ii), b))~; on conclut alors à
l'aide du lemme (11.9.12), puisque les $\mathcal G_\lambda$ sont $S$-plats.
\end{proof}

\begin{env}[11.9.14]
\label{IV.11.9.14-fr}
Gardons toujours les notations de (11.9.1) et (11.9.5) et supposons que $X$
soit un $S$-préschéma. Si pour tout changement de base
$g:X\times_S S'\to X$, où $S'\to S$ est un morphisme quelconque, la famille
$(u'_\lambda)$ correspondante est séparante, nous dirons que la famille
$(u_\lambda)$ est \emph{universellement séparante relativement à} $S$.
Lorsque la famille $(u_\lambda)$ est réduite à un seul élément $u$, nous
dirons encore que $u$ est \emph{universellement injectif}, relativement à
$S$. Il est clair alors que pour tout morphisme $h:S'\to S$, la famille
$(u'_\lambda)$ correspondante est universellement séparante relativement à
$S'$~; inversement, si $h$ est \emph{fidèlement plat} et si $(u'_\lambda)$ est
universellement séparante relativement à $S'$, alors $(u_\lambda)$ est
universellement séparante relativement à $S$, comme il résulte aussitôt de
(11.9.10, (i)) et du fait que pour tout morphisme $S''\to S$, le morphisme
correspondant $S''\times_S S'\to S''$ est fidèlement plat.
\end{env}

\oldpage[IV-3]{167}

\begin{proposition}[11.9.15]
\label{IV.11.9.15-fr}
Les notations étant celles de (11.9.1), supposons que $X$ soit un
$S$-préschéma, les $\mathcal G_\lambda$ étant $S$-plats. Soit $S_0$ un
sous-préschéma fermé de $S$ défini par un Idéal quasi-cohérent nilpotent
$\mathcal J$ de $\mathcal O_S$, tel que les
$(\mathcal O_S/\mathcal J)$-Modules $\mathcal J^k/\mathcal J^{k+1}$ soient
tous localement libres. Soit $(u_{\lambda0})$ la famille d'homomorphismes
obtenue à partir du changement de base
$X\leftarrow X_0=X\times_S S_0$. Alors, si la famille $(u_{\lambda0})$ est
séparante (resp. universellement séparante relativement à $S_0$), la famille
$(u_\lambda)$ est séparante (resp. universellement séparante relativement à
$S$).
\end{proposition}

\begin{proof}
Notons que si $S'\to S$ est un changement de base quelconque et
$S'_0=S'\times_S S_0$, $S'_0$ est un sous-préschéma fermé de $S'$ défini par
un idéal quasi-cohérent nilpotent $\mathcal J'$ de $\mathcal O_{S'}$, tel que
$\mathcal J'^k/\mathcal J'^{k+1}$ soit un
$(\mathcal O_{S'}/\mathcal J')$-Module localement libre pour tout $k$
(2.1.8, (i))~; comme en outre les
$\mathcal G'_\lambda=\mathcal G_\lambda\otimes_S S'$ sont $S'$-plats, on voit
que l'assertion relative aux familles universellement séparantes est
conséquence de l'assertion relative aux familles séparantes. Pour prouver
cette dernière, on peut (11.9.3) se ramener au cas où
$S=\operatorname{Spec}(A)$, $X=\operatorname{Spec}(B)$ sont affines,
$Z_\lambda=X$ pour tout $\lambda$, $\mathcal F=\widetilde M$,
$\mathcal G_\lambda=\widetilde N_\lambda$, où $M$ et les $N_\lambda$ sont des
$B$-modules, les $N_\lambda$ étant des $A$-modules plats. En outre, la question
étant locale sur $X$ et $S$, il suffit de voir que si $t\in M$ est tel que
$u_\lambda(t)=0$ pour tout $\lambda$, alors $t=0$. On a
$\mathcal J=\widetilde{\mathfrak J}$, où $\mathfrak J$ est un idéal nilpotent
de $A$, tel que les $\mathfrak J^k/\mathfrak J^{k+1}$ soient des
$(A/\mathfrak J)$-modules libres, et par hypothèse les
$u_{\lambda0}:M/\mathfrak JM\to N_\lambda/\mathfrak JN_\lambda$ forment une
famille séparante. Supposons que $\mathfrak J^{n+1}=0$ ($n$ entier $\geq0$) et
raisonnons par récurrence sur $n$, l'assertion étant triviale pour $n=0$. Si
$\bar t$ est la classe de $t$ dans $M/\mathfrak J^nM$, la classe de
$u_\lambda(t)$ dans $N_\lambda/\mathfrak J^nN_\lambda$ est nulle pour tout
$\lambda$, donc, par l'hypothèse de récurrence, $\bar t=0$, autrement dit on a
$t\in\mathfrak J^nM$. Or
$\mathfrak J^n=\mathfrak J^n/\mathfrak J^{n+1}$ est un
$(A/\mathfrak J)$-module libre~; si $(e_\alpha)$ est une base de ce module, on
peut donc écrire $t=\sum_\alpha e_\alpha t^0_\alpha$, avec
$t^0_\alpha\in M/\mathfrak JM$, nul sauf pour un nombre fini d'indices.
D'autre part, puisque $N_\lambda$ est un $A$-module plat,
$\mathfrak J^nN_\lambda$ s'identifie à
$\mathfrak J^n\otimes_{(A/\mathfrak J)}(N_\lambda/\mathfrak JN_\lambda)$ et
l'on peut par suite écrire
\[
u_\lambda(t)=\sum_\alpha e_\alpha u_{\lambda0}(t^0_\alpha)
=\sum_\alpha e_\alpha\otimes u_{\lambda0}(t^0_\alpha).
\]
Comme par hypothèse $u_\lambda(t)=0$, on en déduit que
$u_{\lambda0}(t^0_\alpha)=0$ pour tout $\alpha$ et tout $\lambda$~; d'où
$t^0_\alpha=0$ pour tout $\alpha$ puisque la famille $(u_{\lambda0})$ est
séparante, et par suite $t=0$. C.Q.F.D.
\end{proof}

\begin{theorem}[11.9.16]
\label{IV.11.9.16-fr}
Les notations étant celles de (11.9.1), supposons que $X$ soit un
$S$-préschéma localement noethérien, $\mathcal F$ un $\mathcal O_X$-Module
cohérent et que les $\mathcal G_\lambda$ soient $S$-plats. Pour tout $s\in S$,
soit $((u_\lambda)_s)_{\lambda\in L}$ la famille obtenue à partir de
$(u_\lambda)$ par le changement de base
$X_s=X\times_S\operatorname{Spec}(k(s))\to X$. Alors, pour que la famille
$(u_\lambda)$ soit universellement séparante relativement à $S$, il faut et il
suffit que pour tout $s\in S$, la famille $((u_\lambda)_s)$ soit séparante.
\end{theorem}

\begin{proof}
La nécessité de la condition découle trivialement des définitions.
Inversement, supposons la condition de l'énoncé vérifiée, et prouvons d'abord
que la famille $(u_\lambda)$ est séparante. On peut (11.9.3) se réduire au cas
où $S=\operatorname{Spec}(A)$, $X=\operatorname{Spec}(B)$ sont affines,
$Z_\lambda=X$ pour tout $\lambda$, $\mathcal F=\widetilde M$,
$\mathcal G_\lambda=\widetilde N_\lambda$, où $B$ est noethérien, $M$ est un
$B$-module de type fini et les $N_\lambda$ des $A$-modules plats et se borner à
prouver que, si $t\in M$ est tel que $u_\lambda(t)=0$ pour tout $\lambda$,
alors $t=0$. Pour montrer que $t=0$, il suffit de prouver que pour tout idéal
maximal $\mathfrak p$ de $B$, l'image $t_\mathfrak p$ de $t$ dans
$M_\mathfrak p$ est nulle (Bourbaki, \emph{Alg. comm.}, chap. II, \S~3,
n\textsuperscript{o}~3, cor. 1 du th. 1). On peut donc se borner à montrer que
l'intersection
\oldpage[IV-3]{168}
des noyaux des
$(u_\lambda)_\mathfrak p:M_\mathfrak p\to(N_\lambda)_\mathfrak p$ déduits
des $(u_\lambda)$ par le changement de base
$\operatorname{Spec}(B_\mathfrak p)\to\operatorname{Spec}(B)$ est réduite à
0. Autrement dit, on est ramené au cas où $B$ est un anneau local noethérien,
et en considérant l'idéal premier de $A$ image réciproque de l'idéal maximal de
$B$, on peut aussi supposer que $A$ est un anneau local, d'idéal maximal
$\mathfrak m$. Alors, comme $\mathfrak mB$ est contenu dans l'idéal maximal de
$B$, et que $M$ est un $B$-module de type fini, l'intersection des
$\mathfrak m^nM$ est réduite à 0 $(\mathbf 0_{\mathrm I},7.3.5)$, donc il
suffit de prouver que pour tout $n$, l'image de $t$ dans
$M/\mathfrak m^{n+1}M$ est nulle. Il suffit donc de prouver que la famille
déduite de $(u_\lambda)$ par le changement de base
$\operatorname{Spec}(B/\mathfrak m^{n+1}B)\to\operatorname{Spec}(B)$ est
séparante, ce qui signifie encore qu'on peut se borner au cas où $A$ est un
anneau local dont l'idéal maximal $\mathfrak m$ est nilpotent. Mais alors les
$\mathfrak m^k/\mathfrak m^{k+1}$ sont des $(A/\mathfrak m)$-modules libres,
et en vertu de l'hypothèse sur les $(u_\lambda)_s$, on est précisément dans les
conditions d'application de (11.9.15), d'où la conclusion annoncée.

Soient maintenant $h:S'\to S$ un morphisme changement de base, et
$(u'_\lambda)$ la famille obtenue à partir de $(u_\lambda)$ par le changement
de base $h':X\times_S S'\to X$~; prouvons que $(u'_\lambda)$ est aussi
séparante. Supposons d'abord que $h$ soit \emph{localement de type fini}~; il
en est alors de même de $h'$, donc $X'=X\times_S S'$ est localement
noethérien~; de plus, si $s'\in S'$ est au-dessus du point $s\in S$, il résulte
de (11.9.13) appliqué à $X_s$ et à
$X_{s'}=X_s\otimes_{k(s)}k(s')$ que, pour tout $s'\in S'$, la famille
$((u'_\lambda)_{s'})$ est séparante~; on peut par suite conclure de la première
partie de la démonstration que dans ce cas $(u'_\lambda)$ est séparante.

Enfin, si $h$ est quelconque, on peut évidemment se limiter au cas où
$S=\operatorname{Spec}(A)$ et $S'=\operatorname{Spec}(A')$ sont affines, et
considérer $A'$ comme limite inductive de ses sous-$A$-algèbres de type fini.
Comme les $\mathcal G_\lambda$ sont $S$-plats, il suffit d'appliquer ce qui
précède et le lemme (11.9.5) pour terminer la démonstration.
\end{proof}

\begin{proposition}[11.9.17]
\label{IV.11.9.17-fr}
Soient $f:X\to S$ un morphisme localement de présentation finie, $\mathcal F$
un $\mathcal O_X$-Module quasi-cohérent de présentation finie et $f$-plat,
$U$ un ouvert de $X$, $j:U\to X$ l'injection canonique,
$u:\mathcal F\to j_*(j^*(\mathcal F))$ l'homomorphisme canonique
$(\mathbf 0_{\mathrm I},4.4.3.2)$. Pour tout $s\in S$, soient $X_s$ la fibre
$X\times_S\operatorname{Spec}(k(s))$, $U_s$ l'ouvert $U\cap X_s$ de $X_s$,
$\mathcal F_s=\mathcal F\otimes_{\mathcal O_S}k(s)$,
$j_s:U_s\to X_s$ l'injection canonique,
$u_s:\mathcal F_s\to(j_s)_*((j_s)^*(\mathcal F_s))$ l'homomorphisme
canonique. Pour que $u$ soit universellement injectif relativement à $S$, il
faut et il suffit que $u_s$ soit injectif pour tout $s\in S$.
\end{proposition}

\begin{proof}
Il n'y a encore à prouver que la suffisance de la condition. Lorsque $X$ est
localement noethérien, la proposition est un corollaire immédiat de (11.9.16).
Nous allons nous ramener à ce cas en deux étapes, en nous bornant, comme on
peut évidemment le faire, au cas où $S=\operatorname{Spec}(A)$ et $X$ sont
affines.

\emph{A) Cas où $U$ est quasi-compact.} --- Nous utiliserons le lemme suivant~:

\begin{lemma}[11.9.17.1]
\label{IV.11.9.17.1-fr}
Sous les hypothèses générales de (11.9.17), et en supposant en outre $S$ et
$X$ affines et $U$ quasi-compact, l'ensemble $E$ des $s\in S$ tels que $u_s$
soit injectif est constructible.
\end{lemma}

\begin{proof}
En effet, les fibres $X_s$ sont des préschémas localement noethériens, donc
$E$ peut aussi, en vertu de (5.10.2), être défini comme l'ensemble des
$s\in S$ tels que $\operatorname{Ass}(\mathcal F_s)\subset U_s$. Notons en
outre que $U$, étant quasi-compact dans un schéma affine, est constructible.
Alors, la vérification de la condition (9.2.1, (i)) découle aussitôt de
(4.2.7)~; d'autre
\oldpage[IV-3]{169}
part, la vérification de (9.2.1, (ii)) se fait aisément en utilisant l'étude des
cycles premiers associés au voisinage de la fibre générique (9.8.3), ainsi que
(9.5.2) et (9.5.3).
\end{proof}

Ce lemme étant établi, on peut, en vertu de (8.9.1) et (8.2.11), supposer qu'il
existe un sous-anneau noethérien $A_0$ de $A$, un préschéma de type fini $X_0$
sur $S_0=\operatorname{Spec}(A_0)$, un ouvert $U_0$ de $X_0$ et un
$\mathcal O_{X_0}$-Module cohérent $\mathcal F_0$ tels que
$X=X_0\times_{S_0}S$ et que, si $p_0:X\to X_0$ est la projection canonique,
on ait $U=p_0^{-1}(U_0)$ et $\mathcal F=p_0^*(\mathcal F_0)$. Soit
$(A_\alpha)$ la famille filtrante des sous-anneaux de $A$ qui sont des
$A_0$-algèbres de type fini, et posons
$S_\alpha=\operatorname{Spec}(A_\alpha)$,
$X_\alpha=X_0\times_{S_0}S_\alpha$,
$U_\alpha=U_0\times_{S_0}S_\alpha$,
$\mathcal F_\alpha=\mathcal F_0\otimes_{S_0}S_\alpha$~; soit $u_\alpha$
l'homomorphisme canonique relatif à $\mathcal F_\alpha$ et $U_\alpha$, défini
comme dans (11.9.17). Pour tout $s\in S$, l'hypothèse que $u_s$ soit injectif
entraîne que $(u_\alpha)_{s_\alpha}$ l'est aussi (11.9.10, (i)), où
$s_\alpha=q_\alpha(s)$ est l'image de $s$ par le morphisme
$q_\alpha:S\to S_\alpha$. Si $E_\alpha$ est l'ensemble des
$s_\alpha\in S_\alpha$ tels que $(u_\alpha)_{s_\alpha}$ soit injectif, on a
donc $S=q_\alpha^{-1}(E_\alpha)$, et les $E_\alpha$ forment un système
projectif d'ensembles. Mais le lemme (11.9.17.1) appliqué à $u_\alpha$ montre
que $E_\alpha$ est constructible dans $S_\alpha$~; on déduit donc de (8.3.4)
qu'il existe un indice $\alpha$ tel que $E_\alpha=S_\alpha$. Mais alors, comme
$X_\alpha$ est noethérien, $u_\alpha$ est universellement injectif en vertu de
(11.9.16), donc il en est de même de $u$.

\emph{B) Cas général.} --- L'ouvert $U$ est réunion filtrante croissante
d'ouverts quasi-compacts $U_\lambda$~; si $j_\lambda:U_\lambda\to X$ est
l'injection canonique et
$u_\lambda:\mathcal F\to(j_\lambda)_*(j_\lambda^*(\mathcal F))$
l'homomorphisme correspondant, il résulte du lemme (11.9.17.1) que l'ensemble
$E_\lambda$ des $s\in S$ tels que $(u_\lambda)_s$ soit injectif est
constructible. D'autre part, pour un $s\in S$, dire que $u_s$ (resp.
$(u_\lambda)_s$) est injectif signifie que
$\operatorname{Ass}(\mathcal F_s)\subset U\cap X_s$ (resp.
$\operatorname{Ass}(\mathcal F_s)\subset U_\lambda\cap X_s$) (5.10.2).
Comme $\operatorname{Ass}(\mathcal F_s)$ est fini (3.1.6), dire que $u_s$ est
injectif signifie donc qu'il existe $\lambda$ tel que $s\in E_\lambda$~;
l'hypothèse signifie par suite que $S=\bigcup_\lambda E_\lambda$. En vertu de
(1.9.9), il existe un indice $\lambda$ tel que $S=E_\lambda$, d'où l'on
conclut par la première partie du raisonnement que $u_\lambda$ est
universellement injectif. Il résulte alors de la factorisation de $u_\lambda$~:
\[
\mathcal F\xrightarrow{u}j_*(\mathcal F|U)\longrightarrow
(j_\lambda)_*(\mathcal F|U_\lambda)
\]
que $u$ est aussi universellement injectif.
\end{proof}

\begin{remark}[11.9.18]
\label{IV.11.9.18-fr}
Soient $X$ un préschéma, $\mathcal F$, $\mathcal G$ deux
$\mathcal O_X$-Modules quasi-cohérents de présentation finie, $\mathcal G$
étant en outre supposé \emph{localement libre},
$u:\mathcal F\to\mathcal G$ un homomorphisme. Les conditions suivantes sont
équivalentes~:
\begin{enumerate}[label=\emph{\alph*)}]
  \item pour tout morphisme $g:X'\to X$, l'homomorphisme
  $g^*(u):g^*(\mathcal F)\to g^*(\mathcal G)$ est injectif~;
  \item pour tout $x\in X$, l'homomorphisme
  $u\otimes 1_{k(x)}:\mathcal F\otimes_{\mathcal O_X}k(x)
  \to\mathcal G\otimes_{\mathcal O_X}k(x)$ est injectif~;
  \item pour tout $x\in X$, il existe un voisinage ouvert $U$ de $x$ tel que
  $u|U:\mathcal F|U\to\mathcal G|U$ soit un isomorphisme de $\mathcal F|U$
  sur un facteur direct de $\mathcal G|U$.
\end{enumerate}

En effet, il est clair que c) entraîne a) et que a) entraîne b). Le fait que
b) entraîne c) résulte de $(\mathbf 0,19.1.12)$,
$(\mathbf 0_{\mathrm I},5.2.5)$ et $(\mathbf 0_{\mathrm I},5.5.5)$.

Lorsque les conditions équivalentes précédentes sont vérifiées, on dit que
$u$ est \emph{universellement injectif}.
\end{remark}
\oldpage[IV-3]{170}
\subsection{Familles schématiquement dominantes de morphismes et familles
schématiquement denses de sous-préschémas}
\label{subsection:IV.11.10-fr}

\begin{proposition}[11.10.1]
\label{IV.11.10.1-fr}
Soient $X$ un préschéma, $(Z_\lambda)_{\lambda\in L}$ une famille de
préschémas, et pour tout $\lambda\in L$, soit
$f_\lambda=(\psi_\lambda,\theta_\lambda):Z_\lambda\to X$ un morphisme. Les
conditions suivantes sont équivalentes~:
\begin{enumerate}[label=\emph{\alph*)}]
  \item La famille des homomorphismes
  $\theta_\lambda:\mathcal O_X\to(f_\lambda)_*(\mathcal O_{Z_\lambda})$ est
  séparante (autrement dit (11.9.1), l'intersection des noyaux des
  $\theta_\lambda$ est nulle).
  \item Pour tout ouvert $U$ de $X$, toute section $t$ de $\mathcal O_X$
  au-dessus de $U$ dont toutes les images par les homomorphismes canoniques
  \[
  (\theta_\lambda)_U:\Gamma(U,\mathcal O_X)
  \longrightarrow\Gamma(f_\lambda^{-1}(U),\mathcal O_{Z_\lambda})
  \tag{11.10.1.1}
  \label{IV.11.10.1.1-fr}
  \]
  sont nulles, est elle-même nulle.
  \item Pour tout ouvert $U$ de $X$, et tout sous-préschéma fermé $Y$ de $U$
  tel que pour tout $\lambda\in L$, il existe une factorisation
  \[
  f_\lambda^{-1}(U)\xrightarrow{g_\lambda}Y\xrightarrow{j}U
  \tag{11.10.1.2}
  \label{IV.11.10.1.2-fr}
  \]
  de la restriction $f_\lambda^{-1}(U)\to U$ de $f_\lambda$ (où $j$ est
  l'injection canonique), on a $Y=U$.
\end{enumerate}

Si de plus $X$ est un $S$-préschéma, ces conditions sont aussi équivalentes à
la suivante~:
\begin{enumerate}[label=\emph{\alph*)}]
\setcounter{enumi}{3}
  \item Pour tout morphisme séparé $g:X'\to S$ et tout couple de
  $S$-morphismes $u$, $v$ d'un ouvert $U$ de $X$ dans $X'$, tel que pour tout
  $\lambda$, les composés de $u$ et $v$ avec le morphisme
  $f_\lambda^{-1}(U)\to U$, restriction de $f_\lambda$, soient égaux, on a
  $u=v$.
\end{enumerate}
\end{proposition}

\begin{proof}
L'équivalence de a) et b) résulte des définitions. Pour voir que b) entraîne
c), il suffit de considérer l'Idéal quasi-cohérent $\mathcal J$ de
$\mathcal O_U$ définissant $Y$, et de noter que, pour tout ouvert $V\subset U$,
l'hypothèse entraîne que le morphisme $(\theta_\lambda)_V$ se factorise en
\[
\Gamma(V,\mathcal O_U)\longrightarrow
\Gamma(V,\mathcal O_U/\mathcal J)\longrightarrow
\Gamma(f_\lambda^{-1}(V),\mathcal O_{Z_\lambda}).
\]
On en conclut que toute section $t$ de $\mathcal J$ au-dessus de $V$ a pour
image 0 dans tous les $\Gamma(f_\lambda^{-1}(V),\mathcal O_{Z_\lambda})$,
donc, en vertu de b), $\mathcal J=0$ et $Y=U$. Inversement, si c) est vérifiée,
il suffit, pour prouver b), d'appliquer c) au sous-préschéma fermé $Y$ de $U$
défini par l'Idéal $t\mathcal O_U$~: l'hypothèse que les images de $t$ par les
$(\theta_\lambda)_U$ sont toutes nulles entraîne que l'on a des factorisations
(11.10.1.2) pour tout $\lambda$ ($\mathbf I$, 4.1.9). Pour prouver que c)
entraîne d), il suffit d'appliquer c) au sous-préschéma fermé $Y$ de $U$, image
réciproque de la diagonale de $X'\times_SX'$ par $(u,v)_S$ et d'utiliser
($\mathbf I$, 4.4.1). Inversement, on déduit b) de d) en considérant le
$S$-schéma $X'=S[T]=S\otimes_{\mathbf Z}\mathbf Z[T]$ ($T$ indéterminée) et en
se rappelant que les sections de $\mathcal O_U$ au-dessus de $U$ correspondent
biunivoquement aux $S$-morphismes $U\to S[T]$ ($\mathbf I$, 3.3.15)~; dire que
deux sections de $\mathcal O_U$ au-dessus de $U$ ont mêmes images par tous les
$(\theta_\lambda)_U$ équivaut à dire que les composés des deux morphismes
correspondants avec tous les morphismes $f_\lambda^{-1}(U)\to U$ sont égaux.
\end{proof}

\oldpage[IV-3]{171}
\begin{definition}[11.10.2]
\label{IV.11.10.2-fr}
Lorsque les conditions équivalentes de (11.10.1) sont vérifiées, on dit que la
famille $(f_\lambda)$ est \emph{schématiquement dominante}. Lorsque les
$f_\lambda$ sont les immersions canoniques dans $X$ d'une famille
$(Z_\lambda)$ de sous-préschémas de $X$, on dit aussi que la famille
$(Z_\lambda)$ est \emph{schématiquement dense}.
\end{definition}

\begin{remarks}[11.10.3]
\label{IV.11.10.3-fr}
\begin{enumerate}[label=(\roman*)]
  \item La notion de famille schématiquement dominante est \emph{locale} sur
  $X$, comme il résulte par exemple de la forme b) dans (11.10.1)~: si
  $(W_\alpha)$ est un recouvrement ouvert de $X$, la famille $(f_\lambda)$ est
  schématiquement dominante si et seulement s'il en est ainsi de chacune des
  familles formées des morphismes
  $f_\lambda^{-1}(W_\alpha)\to W_\alpha$, restrictions des $f_\lambda$.
  \item Si $Z$ est le préschéma somme des $Z_\lambda$, $f:Z\to X$ le morphisme
  coïncidant avec $f_\lambda$ dans chacun des $Z_\lambda$, il revient au même
  de dire que la famille $(f_\lambda)$ est schématiquement dominante, ou que
  la famille réduite au seul élément $f$ l'est.
  \item Soit $M$ un second ensemble d'indices et, pour tout $\lambda\in L$,
  soit $(g_{\lambda\mu})_{\mu\in M}$ une famille de morphismes
  $g_{\lambda\mu}:Z'_{\lambda\mu}\to Z_\lambda$~; si, pour chaque
  $\lambda\in L$, la famille $(g_{\lambda\mu})_{\mu\in M}$ est
  schématiquement dominante, alors, il revient au même de dire que la famille
  $(f_\lambda)_{\lambda\in L}$ est schématiquement dominante, ou que la
  famille
  $(f_\lambda\circ g_{\lambda\mu})_{(\lambda,\mu)\in L\times M}$ l'est
  (11.9.2).
  \item Soit $f:Z\to X$ un morphisme tel que $f_*(\mathcal O_Z)$ soit un
  $\mathcal O_X$-Module quasi-cohérent (par exemple un morphisme quasi-compact
  et quasi-séparé (1.7.4)). Alors, dire que $f$ est schématiquement dominant
  signifie que \emph{l'image fermée de $Z$ par $f$} ($\mathbf I$, 9.5.3) est
  \emph{identique à $X$}.
\end{enumerate}
\end{remarks}

\begin{proposition}[11.10.4]
\label{IV.11.10.4-fr}
Si la famille de morphismes $f_\lambda:Z_\lambda\to X$ est schématiquement
dominante, la réunion des $f_\lambda(Z_\lambda)$ est dense dans $X$.
Inversement, si cette réunion est dense dans $X$ et si de plus $X$ est réduit,
la famille $(f_\lambda)$ est schématiquement dominante.
\end{proposition}

\begin{proof}
La première assertion résulte aussitôt de (11.10.1, b)). D'autre part, si $X$
est réduit, et si la réunion des $f_\lambda(Z_\lambda)$ est dense dans $X$,
alors, si l'on a des factorisations (11.10.1.2) pour tout $\lambda\in L$, on a
aussi des factorisations
\[
(f_\lambda^{-1}(U))_{\mathrm{red}}
\xrightarrow{(g_\lambda)_{\mathrm{red}}}Y_{\mathrm{red}}
\xrightarrow{j_{\mathrm{red}}}U,
\]
et l'hypothèse entraîne que l'espace sous-jacent à $Y$ est identique à $U$,
donc $Y=Y_{\mathrm{red}}=U$ puisque $U$ est réduit.
\end{proof}

Les résultats du n$^\circ$~11.9 sur les familles séparantes se traduisent en
résultats sur les familles schématiquement dominantes~:

\begin{theorem}[11.10.5]
\label{IV.11.10.5-fr}
Soient $(f_\lambda)$ une famille de morphismes
$f_\lambda:Z_\lambda\to X$, $g:X'\to X$ un morphisme, et posons
$Z'_\lambda=Z_\lambda\times_XX'$,
$f'_\lambda=(f_\lambda)_{(X')}:Z'_\lambda\to X'$.
\begin{enumerate}[label=(\roman*)]
  \item Si $g$ est fidèlement plat et si la famille $(f'_\lambda)$ est
  schématiquement dominante, alors la famille $(f_\lambda)$ est
  schématiquement dominante.
  \item Supposons que $g$ soit un morphisme plat et en outre que l'une des
  conditions suivantes soit vérifiée~:
  \begin{enumerate}[label=\emph{\alph*)}]
    \item $L$ est fini et les $f_\lambda$ sont quasi-compacts.
    \item Le morphisme $g$ est localement de présentation finie.
  \end{enumerate}
  Alors, si la famille $(f_\lambda)$ est schématiquement dominante, il en est
  de même de la famille $(f'_\lambda)$.
\end{enumerate}
\end{theorem}

\oldpage[IV-3]{172}
\begin{proposition}[11.10.6]
\label{IV.11.10.6-fr}
Les notations étant celles de (11.10.5), supposons que $X$ soit un préschéma
sur un corps $k$, et qu'en posant $S=\operatorname{Spec}(k)$, on ait
$X'=X\times_SS'$, où $S'$ est un $k$-préschéma quelconque. Alors, si la famille
$(f_\lambda)$ est schématiquement dominante, il en est de même de
$(f'_\lambda)$.
\end{proposition}

\begin{corollary}[11.10.7]
\label{IV.11.10.7-fr}
Soient $X$ un préschéma sur un corps $k$, $(Z_\lambda)_{\lambda\in L}$ une
famille de $k$-préschémas géométriquement réduits sur $k$, et pour chaque
$\lambda\in L$, soit $f_\lambda:Z_\lambda\to X$ un $k$-morphisme. Désignons
par $Y$ le sous-préschéma réduit de $X$ ayant pour espace sous-jacent
l'adhérence de $\bigcup_{\lambda\in L}f_\lambda(Z_\lambda)$. Soient $k'$ une
extension de $k$, $X'$, $Z'_\lambda$,
$f'_\lambda:Z'_\lambda\to X'$ les préschémas et morphismes déduits de $X$,
$Z_\lambda$ et $f_\lambda$ par extension de la base à $k'$. Alors, si $Y'$ est
le sous-préschéma réduit de $X'$ ayant pour espace sous-jacent l'adhérence de
$\bigcup_{\lambda\in L}f'_\lambda(Z'_\lambda)$, on a
$Y'=Y\otimes_k k'$. En particulier, $Y$ est géométriquement réduit sur $k$.
\end{corollary}

\begin{proof}
Comme les $Z_\lambda$ sont réduits, les morphismes $f_\lambda$ se factorisent
en
\[
Z_\lambda\xrightarrow{g_\lambda}Y\xrightarrow{j}X,
\]
où $j$ est l'injection canonique ($\mathbf I$, 5.2.2). Il résulte alors de
(11.10.4) que $(g_\lambda)$ est une famille schématiquement dominante. Posons
$Y'_1=Y\otimes_k k'$, sous-préschéma fermé de $X'$, et soit $g'_\lambda$ le
morphisme déduit de $g_\lambda$ par extension de la base à $k'$, de sorte que
$f'_\lambda$ se factorise en
\[
Z'_\lambda\xrightarrow{g'_\lambda}Y'_1\xrightarrow{j'}X',
\]
où $j'$ est l'injection canonique. Il résulte de (11.10.6) que la famille
$(g'_\lambda)$ est schématiquement dominante. Mais par hypothèse les
$Z'_\lambda$ sont \emph{réduits}, donc ($\mathbf I$, 5.2.2) les
$g'_\lambda$ se factorisent en $Z'_\lambda\to Y'\to Y'_1$~; on conclut donc
de (11.10.2) que $Y'=Y'_1$ et $h'_\lambda=g'_\lambda$, ce qui établit le
corollaire.
\end{proof}

\begin{definition}[11.10.8]
\label{IV.11.10.8-fr}
Les notations étant celles de (11.10.5), supposons que $X$ soit un
$S$-préschéma. On dit que $(f_\lambda)$ est \emph{universellement
schématiquement dominante} (relativement à $S$) si, pour tout changement de
base $S'\to S$, la famille $(f'_\lambda)$ correspondant au changement de base
$X'=X\times_SS'\to X$, est schématiquement dominante.
\end{definition}

Lorsque $S$ est le spectre d'un corps, la prop. (11.10.6) signifie donc qu'une
famille schématiquement dominante est universellement schématiquement
dominante (relativement à $S$).

Lorsque les $f_\lambda$ sont des immersions canoniques $Z_\lambda\to X$, on
dira aussi que la famille des sous-préschémas $Z_\lambda$ est
\emph{universellement schématiquement dense dans $X$} (relativement à $S$) au
lieu de dire que la famille $(f_\lambda)$ est universellement schématiquement
dominante (relativement à $S$).

\begin{theorem}[11.10.9]
\label{IV.11.10.9-fr}
Les notations étant celles de (11.10.5), supposons que $X$ soit un
$S$-préschéma localement noethérien et que les $Z_\lambda$ soient tous
$S$-plats. Pour tout $s\in S$, soient
$X_s=X\times_S\operatorname{Spec}(k(s))$,
$(Z_\lambda)_s=Z_\lambda\times_S\operatorname{Spec}(k(s))$,
$(f_\lambda)_s=f_\lambda\times1:(Z_\lambda)_s\to X_s$. Pour que la famille
$(f_\lambda)$ soit universellement schématiquement dominante relativement à
$S$, il faut et il suffit que, pour tout $s\in S$, la famille
$((f_\lambda)_s)$ soit schématiquement dominante.
\end{theorem}

\begin{proposition}[11.10.10]
\label{IV.11.10.10-fr}
Soient $f:X\to S$ un morphisme plat localement de présentation finie, $U$ un
ouvert de $X$. Pour que $U$ soit universellement schématiquement dense dans
$X$ relativement à $S$, il faut et il suffit que, pour tout $s\in S$,
$U_s=U\cap X_s$ soit schématiquement dense dans $X_s$ (ou, ce qui revient au
même, que l'on ait $\operatorname{Ass}(\mathcal O_{X_s})\subset U_s$).
\end{proposition}
% IV3_B25_P172_END

\oldpage[IV-3]{173}
\section{Étude des fibres des morphismes plats de présentation finie}
\label{section:IV.12-fr}

\setcounter{subsection}{-1}
\subsection{Introduction}
\label{subsection:IV.12.0-fr}

\emph{Nous utiliserons dans tout ce paragraphe les notations générales de}
(9.4.1).

\begin{env}[12.0.1]
\label{IV.12.0.1-fr}
Étant donné un morphisme localement de présentation finie $f:X\to Y$, nous
avons vu (9.9) que pour certaines propriétés locales $\mathbf P$ de préschémas
sur des corps ou de Modules sur ces préschémas, l'ensemble $E$ des $x\in X$
tels que la propriété $\mathbf P$ ait lieu, pour la fibre $X_{f(x)}$, au point
$x$ de cette fibre, est localement constructible dans $X$. Nous nous proposons
de montrer que, pour la plupart de ces propriétés, si l'on suppose de plus que
le morphisme $f$ est plat, alors l'ensemble $E$ est même ouvert dans $X$. De
même ((9.2) à (9.8)) nous avons montré que si $f$ est de présentation finie,
et si cette fois $\mathbf P$ désigne certaines propriétés globales de
préschémas sur des corps ou de Modules sur ces préschémas, l'ensemble $F$ des
$y\in Y$ tels que la propriété $\mathbf P$ ait lieu pour la fibre $X_y$ est
localement constructible dans $Y$. Nous montrerons que si l'on suppose de plus
que le morphisme $f$ est propre et plat, $F$ est même ouvert dans $Y$.
\end{env}

\begin{env}[12.0.2]
\label{IV.12.0.2-fr}
La méthode générale de démonstration des propriétés en question comporte trois
étapes. On se ramène d'abord au cas où $Y$ est affine et $X$ de présentation
finie~; puis~:
\begin{enumerate}[label=\emph{\Alph*)}]
  \item À l'aide de (8.9.1) (et éventuellement d'autres résultats du \S~8) et
  de (11.2.6), on se ramène au cas où $X$ et $Y$ sont noethériens.
  \item On applique les résultats du \S~9 rappelés en (12.0.1) prouvant que
  $E$ (resp. $F$) est constructible.
  \item Pour voir que $E$ est ouvert, il suffit, en vertu de
  ($\mathbf 0_{\mathrm{III}}$, 9.2.5) de montrer que si $x\in E$, alors toute
  générisation $x'$ de $x$ appartient aussi à $E$. Utilisant ($\mathbf{II}$,
  7.1.7), on voit, puisque $Y$ est noethérien, qu'il existe un spectre d'anneau
  de valuation discrète $Y_1$ et un morphisme $h:Y_1\to X$ tel que, si $y_1$
  (resp. $y'_1$) est le point fermé (resp. le point générique) de $Y_1$, on ait
  $h(y_1)=x$, $h(y'_1)=x'$. On fait alors le changement de base
  $g=f\circ h:Y_1\to Y$~; compte tenu des résultats des \S\S~4 et 6 sur les
  préschémas localement noethériens sur des corps et les changements du corps
  de base, on est ramené à prouver l'assertion en question pour un point $x_1$
  de $X_1=X\times_YY_1$ au-dessus de $y_1$ et pour une générisation $x'_1$ de
  $x_1$ au-dessus de $y'_1$. Comme $Y_1=\operatorname{Spec}(A)$, où $A$ est
  un anneau de valuation discrète, d'uniformisante $t$, on doit en définitive,
  pour un $A$-module $M$, prouver une propriété de $M$, sachant que $M/tM$ a
  la même propriété et que $t$ est $M$-régulier (ce qui résulte de l'hypothèse
  de platitude)~; on utilise pour cela les résultats de (3.4) et de (5.12). On
  procède de même pour l'ensemble $F\subset Y$.
\end{enumerate}
\end{env}

\oldpage[IV-3]{174}
\subsection{Propriétés locales des fibres d'un morphisme plat localement de
présentation finie}
\label{subsection:IV.12.1-fr}

\begin{theorem}[12.1.1]
\label{IV.12.1.1-fr}
Soient $f:X\to Y$ un morphisme localement de présentation finie, $\mathcal F$
un $\mathcal O_X$-Module $f$-plat et de présentation finie, $\Phi$ une partie
finie de $\mathbf Z\cup\{\pm\infty\}$, $k$ un entier. Les parties suivantes de
$X$ sont ouvertes~:
\begin{enumerate}[label=(\roman*)]
  \item L'ensemble des $x\in X$ tels que les dimensions des cycles premiers
  associés à $\mathcal F_{f(x)}$ et contenant $x$ soient des éléments de
  $\Phi$.
  \item L'ensemble des $x\in X$ tels que les cycles premiers associés à
  $\mathcal F_{f(x)}$ contenant $x$ aient tous la dimension $k$.
  \item L'ensemble des $x\in X$ tels que $x$ n'appartienne à aucun cycle
  premier immergé associé à $\mathcal F_{f(x)}$.
  \item L'ensemble des $x\in X$ tels que $\mathcal F_{f(x)}$ soit
  équidimensionnel au point $x$ et possède la propriété $(S_k)$ au point $x$
  (5.7.2).
  \item L'ensemble des $x\in X$ tels que
  $\operatorname{coprof}((\mathcal F_{f(x)})_x)\leq k$
  ($\mathbf 0$, 16.4.9).
  \item L'ensemble des $x\in X$ tels que $\mathcal F_{f(x)}$ soit un
  $\mathcal O_{X_{f(x)}}$-Module de Cohen-Macaulay au point $x$ (5.7.1).
  \item L'ensemble des $x\in X$ tels que $\mathcal F_{f(x)}$ soit
  géométriquement réduit au point $x$ (4.6.22).
  \item L'ensemble des $x\in X$ tels que $\mathcal F_{f(x)}$ soit
  géométriquement ponctuellement intègre au point $x$ (4.6.22).
\end{enumerate}
\end{theorem}

\begin{proof}
Les questions étant locales sur $X$, on se ramène d'abord au cas où
$X=\operatorname{Spec}(B)$ et $Y=\operatorname{Spec}(A)$ sont affines, $f$ un
morphisme de présentation finie. Il y a alors un sous-anneau noethérien $A_0$
de $A$, un $A_0$-préschéma de type fini $X_0$, un
$\mathcal O_{X_0}$-Module cohérent $\mathcal F_0$ tel que
$\mathcal F_0\otimes_{A_0}A$ soit isomorphe à $\mathcal F$ (8.9.1)~; en
outre, en vertu de (11.2.6), on peut supposer que $\mathcal F_0$ est
$Y_0$-plat (avec $Y_0=\operatorname{Spec}(A_0)$). Si $h:X\to X_0$ est la
projection canonique, l'ensemble $E$ des points $x\in X$ où l'une des
propriétés (i) à (viii) est vérifiée est égal à $h^{-1}(E_0)$, où $E_0$ est
l'ensemble des $x_0\in X_0$ où la propriété correspondante pour $\mathcal F_0$
est vérifiée~: cela résulte, pour les propriétés (i) à (iii), de (4.2.7)~; pour
les propriétés (iv) à (vi), de (6.7.1)~; pour les propriétés (vii) et (viii),
de (4.7.11).

On est ainsi ramené au cas où $A$ et $B$ sont noethériens, $f$ de type fini,
et $\mathcal F=\widetilde M$, où $M$ est un $B$-module de type fini. En vertu
de (9.9.2), on sait que l'ensemble $E$ est constructible pour toutes les
propriétés envisagées, et il reste dans chaque cas l'étape C) de (12.0.2), où
l'on doit prouver que $E$ est stable par générisation.

\begin{env}[12.1.1.1]
\label{IV.12.1.1.1-fr}
Commençons par le cas (iii) qui est le plus simple. Soit $x'\ne x$ une
générisation de $x$ dans $X$. Posons $y=f(x)$, $y'=f(x')$~; il existe un
spectre d'anneau de valuation discrète $Y_1$ et un morphisme $u:Y_1\to X$ tels
que, si $y_1$ et $y'_1$ sont le point fermé et le point générique de $Y_1$, on
ait $u(y_1)=x$ et $u(y'_1)=x'$ ($\mathbf{II}$, 7.1.7)~; posons
$g=f\circ u:Y_1\to Y$~; si $X_1=X\times_YY_1$, il existe une $Y_1$-section
$u_1:Y_1\to X_1$ telle que $u=g_1\circ u_1$, où $g_1:X_1\to X$ est la
projection canonique. Si $x_1=u_1(y_1)$, $x'_1=u_1(y'_1)$, on a donc
$g_1(x_1)=x$ et $g_1(x'_1)=x'$, et $x'_1$ est une générisation de $x_1$.
Appliquant de nouveau (4.2.7), on voit
\oldpage[IV-3]{175}
qu'on est ramené au cas où $Y=\operatorname{Spec}(A)$ est le spectre d'un
anneau de valuation discrète, $y$ étant le point fermé et $y'$ le point
générique de $Y$. Si $t$ est une uniformisante de $A$, l'hypothèse que
$\mathcal F$ est $f$-plat entraîne que $t$ est $\mathcal F$-régulier
($\mathbf 0_{\mathrm I}$, 6.3.4). Par hypothèse aucun des cycles premiers
associés immergés de $\mathcal F/t\mathcal F$ ne contient $x$. Alors, il
résulte de (3.4.4) que $x$ n'appartient à aucun cycle premier immergé associé à
$\mathcal F$, et il en est donc de même de toute générisation de $x$. En
particulier, $x'$ n'appartient à aucun cycle premier immergé associé à
$\mathcal F$, ni \emph{a fortiori} à aucun de ceux associés à
$\mathcal F_{y'}$ ($X_{y'}$ étant un sous-préschéma induit sur un ouvert de
$X$).
\end{env}

\begin{env}[12.1.1.2]
\label{IV.12.1.1.2-fr}
Envisageons ensuite les cas (iv) et (v) ((vi) n'est qu'un cas particulier de
(v)). On procède comme ci-dessus (en utilisant (6.7.1) au lieu de (4.2.7)) et
l'on est ramené au cas où $Y$ est le spectre d'un anneau de valuation discrète
$A$.

L'anneau $C=\mathcal O_{X,x}$ est alors localisé d'une $A$-algèbre de type
fini, donc est caténaire (5.6.4), et par hypothèse le $C$-module
$M_x/tM_x$ vérifie $(S_k)$ et est équidimensionnel (resp. on a
$\operatorname{coprof}(M_x/tM_x)\leq k$)~; on déduit donc de (5.12.2) (resp.
de ($\mathbf 0$, 16.4.10)) que $M_x$ vérifie $(S_k)$ et est équidimensionnel
(resp. que $\operatorname{coprof}(M_x)\leq k$) puisque $t$ est $M_x$-régulier
et appartient à l'idéal maximal de $C$. D'où les conclusions puisque $x'$ est
une générisation de $x$ et que $(\mathcal F_{y'})_{x'}=\mathcal F_{x'}$.
\end{env}

\begin{env}[12.1.1.3]
\label{IV.12.1.1.3-fr}
Passons aux cas (vii) et (viii). On peut évidemment remplacer $Y$ par le
sous-schéma intègre ayant $\overline{\{y'\}}$ pour espace sous-jacent, et $X$
par $f^{-1}(\overline{\{y'\}})$, sans changer les fibres en $y$ et $y'$, donc
on peut supposer $A$ intègre, de corps des fractions $K=k(y')$. On sait
((4.5.11) et (4.7.8)) qu'il existe une extension finie $K'$ de $K$ telle que,
pour le $(\mathcal O_X\otimes_AK')$-Module
$\mathcal F'=\mathcal F\otimes_AK'$, les cycles premiers associés soient
géométriquement irréductibles et que les longueurs géométriques de
$\mathcal F'$ aux points maximaux $z'_i$ de son support soient respectivement
égales aux longueurs de $\mathcal F'_{z'_i}$ en ces points~; il revient donc
au même de dire que $\mathcal F$ est géométriquement réduit (resp.
géométriquement ponctuellement intègre) en un point $x'$ de
$X_{y'}=X\otimes_AK$, ou de dire que $\mathcal F'$ est réduit (resp. intègre)
aux points de $X\otimes_AK'$ au-dessus de $x'$, compte tenu de (4.2.7). Soit
$A'$ une $A$-algèbre finie dont $K'$ est le corps des fractions~; il résulte
de (4.7.11) qu'en tout point de $X\otimes_AA'$ au-dessus de $x$,
$\mathcal F\otimes_AA'$ a la propriété (vii) (resp. (viii)). Remplaçant $A$
par $A'$ et $X$ par $X\otimes_AA'$, on voit qu'on peut se borner à prouver que
$\mathcal F$ est réduit (resp. intègre) en toute générisation $x'$ de $x$. On
procède alors comme dans (12.1.1.1), en utilisant cette fois (4.7.11), et l'on
est encore ramené au cas où $A$ est un anneau de valuation discrète, $y$ étant
le point fermé, $y'$ le point générique de $Y$, et l'uniformisante $t$ de $A$
étant un élément $M$-régulier. Comme par hypothèse
$\mathcal F/t\mathcal F$ est réduit (resp. intègre) au point $x$, il résulte de
(3.4.6) (resp. (3.4.5)) que $\mathcal F$ est réduit (resp. intègre) au point
$x$, donc aussi dans un voisinage de $x$, et en particulier au point $x'$, ce
qui achève la démonstration dans les cas (vii) et (viii).
\end{env}

\begin{env}[12.1.1.4]
\label{IV.12.1.1.4-fr}
Reste à examiner les cas (i) et (ii). Remplaçant $Y$ par le sous-schéma intègre
ayant $\overline{\{y'\}}$ pour espace sous-jacent, on peut se borner au cas où
$Y$ est irréductible et $y'$ son point générique. Soit $z'$ un point générique
d'un cycle premier associé à $\mathcal F_{y'}$
\oldpage[IV-3]{176}
contenant $x'$, et soit $Z'$ l'adhérence de $z'$ dans $X$, de sorte que l'on a
$x\in Z'$. Pour traiter le cas (i), il va suffire de prouver la

\begin{proposition}[12.1.1.5]
\label{IV.12.1.1.5-fr}
Soient $Y$ un préschéma localement noethérien, $f:X\to Y$ un morphisme
localement de type fini, $\mathcal F$ un $\mathcal O_X$-Module cohérent et
$f$-plat~; pour tout $y\in Y$, posons
$\mathcal F_y=\mathcal F\otimes_{\mathcal O_Y}k(y)$. Soient $z'$ un point de
$X$, $y'=f(z')$, et supposons que $z'\in\operatorname{Ass}(\mathcal F_{y'})$.
Soient $Z'$ (resp. $Y'$) le sous-préschéma réduit de $X$ (resp. $Y$) ayant
pour espace sous-jacent l'adhérence de $\{z'\}$ dans $X$ (resp. de $\{y'\}$
dans $Y$). Alors, pour tout $y\in f(Z')$, les dimensions de toutes les
composantes irréductibles de $Z'_y$ sont égales à $\dim(Z'_{y'})$ (dimension
de l'adhérence de $z'$ dans $X_{y'}$) (ce que nous exprimerons plus tard
(13.2.2) en disant que la restriction $Z'\to Y'$ de $f$ est un morphisme
équidimensionnel)~; en outre, en tout point maximal $z$ de $Z'_y$, on a
$z\in\operatorname{Ass}(\mathcal F_y)$.
\end{proposition}

Pour cela, nous allons nous ramener au cas où $Y$ est spectre d'un anneau de
valuation discrète. Prenons un spectre d'anneau de valuation discrète $Y_1$ et
un morphisme $h:Y_1\to X$ tel que $h(y_1)=z$, $h(y'_1)=z'$, $y_1$ et $y'_1$
étant le point fermé et le point générique de $Y_1$ respectivement
($\mathbf{II}$, 7.1.7). Soient $g=f\circ h:Y_1\to Y$,
$X_1=X\times_YY_1$, et soient $f_1:X_1\to Y_1$, $g_1:X_1\to X$ les projections
canoniques~; il y a une $Y_1$-section $h_1:Y_1\to X_1$ telle que
$h=g_1\circ h_1$ ($\mathbf I$, 3.3.14). Si
$Z'_1=Z'_{y'}\otimes_{k(y')}k(y'_1)$, on sait (4.2.6) que les composantes
irréductibles de $Z'_1$ dominent $Z'_{y'}$~; soit $z'_1$ le point générique de
l'une de ces composantes qui contient $h_1(y'_1)$, de sorte que
$f_1(z'_1)=y'_1$ et $g_1(z'_1)=z'$. Comme $h_1(y_1)$ est spécialisation de
$h_1(y'_1)$, il est \emph{a fortiori} spécialisation de $z'_1$~; soit $z_1$ un
point générique d'une des composantes irréductibles de
$\overline{\{z'_1\}}\cap(X_1)_{y_1}$ contenant $h_1(y_1)$. Alors $g_1(z_1)$
est spécialisation de $g_1(z'_1)=z'$ dans $X$, et
$z=h(y_1)=g_1(h_1(y_1))$ est spécialisation de $g_1(z_1)$~; mais comme
$f(g_1(z_1))=y$ et que $z$ est point générique de $Z'_y$, on a nécessairement
$z=g_1(z_1)$.

Supposons maintenant que l'on ait prouvé que, si l'on pose
$\mathcal F_1=\mathcal F\otimes_{\mathcal O_Y}\mathcal O_{Y_1}$, on a
$z_1\in\operatorname{Ass}((\mathcal F_1)_{y_1})$~; il résultera de (4.2.7) que
$z\in\operatorname{Ass}(\mathcal F_y)$~; en outre, les dimensions de
$\overline{\{z_1\}}\cap(X_1)_{y_1}$ et de $\overline{\{z\}}\cap X_y$ sont
égales, et de même les dimensions de
$\overline{\{z'_1\}}\cap(X_1)_{y'_1}$ et de
$\overline{\{z'\}}\cap X_{y'}$ (4.2.7)~; on s'est donc bien ramené, comme
annoncé, au cas où $Y$ est spectre d'un anneau de valuation discrète $A$, $y$
et $y'$ étant son point fermé et son point générique respectivement.

Cela étant, comme $z'\in\operatorname{Ass}(\mathcal F_{y'})$, on a
\emph{a fortiori} $z'\in\operatorname{Ass}(\mathcal F)$. Si $t$ est une
uniformisante de $A$, $t$ est $\mathcal F$-régulier par platitude, donc (3.4.3)
on a $z\in\operatorname{Ass}(\mathcal F/t\mathcal F)
=\operatorname{Ass}(\mathcal F_y)$. Quant à l'assertion relative aux
dimensions, elle résulte de (7.1.13), appliqué à un sous-préschéma fermé de $X$
ayant $Z'$ pour espace sous-jacent.

Abordons enfin le cas (ii) de (12.1.1). Avec les mêmes notations, il faut
prouver que si tous les cycles premiers associés à $\mathcal F_y$ et contenant
$x$ ont la dimension $k$, il en est de même de tous les cycles premiers
associés à $\mathcal F_{y'}$ et contenant $x'$. Or on vient de voir que tout
cycle premier associé à $\mathcal F_{y'}$ et contenant $x'$ a même dimension
que l'un des cycles premiers associés à $\mathcal F_y$ et contenant $x$~; cela
démontre donc (ii).
\end{env}
\end{proof}

\begin{remarks}[12.1.2]
\label{IV.12.1.2-fr}
\begin{enumerate}[label=(\roman*)]
  \item Sous les conditions de (12.1.1.5) avec $\mathcal F=\mathcal O_X$ (de
  sorte que $f$ est plat), on ne peut en général affirmer que la restriction
  $Z'\to Y'$ de $f$ est un morphisme plat ni même un morphisme ouvert. C'est ce
  que montre l'exemple (6.5.5,
  \oldpage[IV-3]{177}
  (ii)) où l'on prend pour $Z'$ une des deux composantes irréductibles de
  $X=\operatorname{Spec}(B)$. Il est immédiat que ce morphisme restriction
  n'est pas ouvert aux points de $Z'$ au-dessus du «~point double~» de $Y'=Y$.
  \item Dans les hypothèses de (12.1.1.5), avec
  $\mathcal F=\mathcal O_X$, on ne peut affaiblir la condition «~$f$ est
  plat~» en «~$f$ est universellement ouvert~» (cf. (2.4.6)), comme nous le
  verrons plus loin sur un exemple (14.4.10, (i)).
\end{enumerate}
\end{remarks}

\begin{corollary}[12.1.3]
\label{IV.12.1.3-fr}
Sous les hypothèses de (12.1.1), la fonction
$x\mapsto\operatorname{coprof}((\mathcal F_{f(x)})_x)$ est semi-continue
supérieurement dans $X$ et la fonction
$x\mapsto\operatorname{prof}_x^*(\mathcal F_{f(x)})$ (10.8.1) est
semi-continue inférieurement dans $X$.
\end{corollary}

\begin{proof}
La première assertion n'est autre qu'une formulation équivalente de (12.1.1,
(v)). Pour prouver la seconde, on peut d'abord, compte tenu de (10.8.7) et
(10.8.8), se ramener comme dans (12.1.1) au cas où $Y$ est le spectre d'un
anneau de valuation discrète $A$ de point fermé $y$ et de point générique $y'$,
$X=\operatorname{Spec}(B)$ étant affine, $\mathcal F=\widetilde M$, où $M$ est
un $B$-module de type fini. Puisque $\mathcal F$ est $f$-plat, toute composante
irréductible de $\operatorname{Supp}(\mathcal F)$ qui contient $x\in X_y$
domine $Y$ (2.3.4), autrement dit son point générique $z'$ appartient à
$X_{y'}$~; les composantes irréductibles $Z$ de
$\operatorname{Supp}(\mathcal F)$ qui contiennent une générisation $x'$ de
$x$ appartenant à $X_{y'}$ sont donc exactement celles qui contiennent $x$ et
en outre, d'après (12.1.1.5), on a $\dim(Z_{y'})=\dim(Z_y)$, d'où, si
$T=\operatorname{Supp}(\mathcal F)$,
$\dim_x(T_y)=\dim_{x'}(T_{y'})$. De plus, pour une uniformisante $t$ de $A$, on
a vu dans (12.1.1, (v)) que l'on a
$\operatorname{coprof}(M_x)=\operatorname{coprof}(M_x/tM_x)$, et puisque
d'autre part $\operatorname{coprof}(M_{x'})\leq\operatorname{coprof}(M_x)$
(6.11.5), on voit que l'on a
$\operatorname{coprof}((\mathcal F_{y'})_{x'})\leq
\operatorname{coprof}((\mathcal F_y)_x)$, la relation
$\operatorname{prof}_x^*(\mathcal F_y)\leq
\operatorname{prof}_{x'}^*(\mathcal F_{y'})$ découle alors de (10.8.7).
\end{proof}

\begin{remark}[12.1.4]
\label{IV.12.1.4-fr}
On déduit aussi de (12.1.1, (i)) que la fonction
$x\mapsto\dim_x(X_{f(x)})$ est semi-continue supérieurement dans $X$, mais
nous verrons plus loin (13.1.3) que cette propriété est vraie même sans
supposer $f$ plat.
\end{remark}

\begin{corollary}[12.1.5]
\label{IV.12.1.5-fr}
Soient $Y$ un préschéma localement noethérien, $f:X\to Y$ un morphisme
localement de type fini, $\mathcal F$ un $\mathcal O_X$-Module cohérent et
$f$-plat, $\mathcal G$ un $\mathcal O_Y$-Module cohérent, $x$ un point de $X$,
$y=f(x)$. Supposons que $\mathcal G$ ait la propriété $(S_k)$ au point $y$, et
que $\mathcal F_y$ ait la propriété $(S_k)$ au point $x$ et soit
équidimensionnel au point $x$. Alors~:
\begin{enumerate}[label=(\roman*)]
  \item $\mathcal F\otimes_{\mathcal O_Y}\mathcal G$ possède la propriété
  $(S_k)$ au point $x$.
  \item Si de plus $Y$ est localement immersible dans un schéma régulier, ou
  est un préschéma excellent (7.8.5), il existe un voisinage de $x$ dans $X$
  tel que $\mathcal F\otimes_{\mathcal O_Y}\mathcal G$ ait la propriété
  $(S_k)$ dans ce voisinage.
\end{enumerate}
\end{corollary}

\begin{proof}
En effet, par (12.1.1, (iv)), il existe un voisinage ouvert $V$ de $x$ tel que
pour tout $x'\in V$, $\mathcal F_{f(x')}$ vérifie $(S_k)$ au point $x'$.
D'autre part, $\mathcal G$ vérifie $(S_k)$ en tous les points de
$Y'=\operatorname{Spec}(\mathcal O_y)$ (5.7.2). Remplaçant $Y$ par $Y'$ et $X$
par $V\times_YY'$, on est ramené (compte tenu de ($\mathbf I$, 3.6.5)) au cas
où $\mathcal G$ vérifie $(S_k)$ dans $Y$ tout entier et où, pour tout
$y\in f(X)$, $\mathcal F_y$ possède la propriété $(S_k)$ en tous les points de
$f^{-1}(y)$~; comme $\mathcal F$ est $f$-plat, on sait alors (6.4.1) que
$\mathcal F\otimes_{\mathcal O_Y}\mathcal G$ possède la propriété $(S_k)$ dans
$X$, ce qui prouve (i). Pour démontrer (ii), il suffit d'observer que, sous les
hypothèses faites, $\mathcal G$ possède la
\oldpage[IV-3]{178}
propriété $(S_k)$ dans un voisinage $U$ de $y$ dans $Y$ ((6.11.2) et (7.8.3,
(iv)))~; on conclut alors de la même manière en remplaçant cette fois $Y$ par
$U$ et $X$ par $V\times_YU$.
\end{proof}

\begin{theorem}[12.1.6]
\label{IV.12.1.6-fr}
Soient $f:X\to Y$ un morphisme plat et localement de présentation finie, $k$
un entier. Les parties suivantes de $X$ sont ouvertes~:
\begin{enumerate}[label=(\roman*)]
  \item L'ensemble des points $x\in X$ tels que $X_{f(x)}$ vérifie la propriété
  $(S_k)$ au point $x$.
  \item L'ensemble des $x\in X$ tels que $X_{f(x)}$ vérifie la propriété
  $(R_k)$ géométrique au point $x$, soit équidimensionnel au point $x$ et tels
  en outre que $x$ n'appartienne à aucun cycle premier immergé associé à
  $X_{f(x)}$.
  \item L'ensemble des $x\in X$ tels que $X_{f(x)}$ soit géométriquement
  régulier (i.e. lisse) au point $x$.
  \item L'ensemble des $x\in X$ tels que $X_{f(x)}$ soit géométriquement
  normal au point $x$.
\end{enumerate}
\end{theorem}

\begin{proof}
Les étapes A) et B) de (12.0.2) se font ici comme dans (12.1.1)~; pour l'étape
A), on utilise (6.7.8), ainsi que (4.2.7) pour (ii)~; pour l'étape B), on
utilise (9.9.2) et (9.9.4). Reste à examiner l'étape C) dans chaque cas.
\begin{enumerate}[label=(\roman*)]
  \item[(i)] Comme dans (12.1.1.2), on se ramène (en utilisant (6.7.8)) au cas
  où $Y$ est le spectre d'un anneau de valuation discrète $A$,
  $X=\operatorname{Spec}(B)$, où $B$ est une $A$-algèbre de type fini, $x$,
  $x'$ deux points de $X$ tels que $f(x)=y$ soit le point fermé et
  $y'=f(x')$ le point générique de $Y$, $x'$ étant en outre une générisation
  de $x$. Comme il s'agit de prouver que $X$ vérifie $(S_k)$ au point $x'$,
  on peut remplacer $X$ par $\operatorname{Spec}(\mathcal O_{X,x})$,
  c'est-à-dire supposer l'anneau $B$ local, l'homomorphisme $A\to B$ étant
  local et $B$ étant une $A$-algèbre essentiellement de type fini (1.3.8).
  Alors, par hypothèse, si $t$ est une uniformisante de $A$, $t$ est un
  élément $B$-régulier par platitude, et $B/tB$ vérifie $(S_k)$. Mais comme
  $A$ est un anneau universellement caténaire (5.6.4), $B$ est caténaire,
  donc, par (5.12.4), il en résulte que $B$ vérifie $(S_k)$, ce qui achève la
  démonstration dans le cas (i).
  \item[(iii)] Ici l'étape C) est inutile~; comme $f$ est plat, on sait en
  effet (6.8.7) que lorsque $Y$ est localement noethérien et $f$ localement de
  type fini, l'ensemble des $x\in X$ tels que $X_{f(x)}$ soit géométriquement
  régulier au point $x$ est ouvert dans $X$.
  \item[(ii)] En raisonnant comme dans (12.1.1.3), pour prouver que la propriété
  considérée est stable par générisation, on peut d'abord, en considérant une
  extension finie quelconque $K'$ de $k(y')$, et tenant compte de la définition
  (6.7.6) de la propriété $(R_k)$ géométrique, ainsi que de l'invariance par
  changement du corps de base des deux autres propriétés figurant dans (ii),
  remplacer dans (ii) la propriété $(R_k)$ géométrique par la propriété
  $(R_k)$. Procédant alors comme dans (i), on se ramène au cas où $Y$ est
  spectre d'un anneau de valuation discrète $A$, et puisque $A$ est caténaire,
  il suffit d'appliquer (5.12.5) pour conclure.
  \item[(iv)] L'ensemble des points $x\in X$ tels que $X_{f(x)}$ soit
  géométriquement normal au point $x$ est contenu dans celui des $x\in X$ tels
  que $X_{f(x)}$ soit géométriquement ponctuellement intègre et vérifie $(S_2)$
  au point $x$, et ce dernier est ouvert dans $X$ en vertu de (12.1.1,
  (viii)). On peut donc déjà supposer que, pour tout $x\in X$, $X_{f(x)}$ est
  géométriquement ponctuellement intègre et vérifie $(S_2)$ au point $x$, et
  \emph{a fortiori} il est équidimensionnel. D'autre part, en vertu du critère
  de Serre (5.8.6), dire que $X_{f(x)}$ est géométriquement normal au point $x$
  signifie que $X_{f(x)}$ vérifie $(S_2)$ et la propriété $(R_1)$ géométrique
  en $x$~;
  \oldpage[IV-3]{179}
  mais en vertu de (ii) et des remarques précédentes, cet ensemble est
  intersection de deux ouverts de $X$, donc est ouvert dans $X$. C.Q.F.D.
\end{enumerate}
\end{proof}

\begin{corollary}[12.1.7]
\label{IV.12.1.7-fr}
Soit $f:X\to Y$ un morphisme localement de présentation finie. Alors
l'ensemble des points $x\in X$ où $f$ possède l'une des propriétés suivantes
(6.8.1)~:
\begin{enumerate}[label=(\roman*)]
  \item vérifier la propriété $(S_k)$~;
  \item être de coprofondeur $\leq n$~;
  \item être de Cohen-Macaulay~;
  \item être régulier (ou lisse, ce qui revient au même)~;
  \item être normal~;
  \item être réduit~;
\end{enumerate}
est ouvert dans $X$.
\end{corollary}

\begin{proof}
En effet, il résulte de (11.3.1) que l'ensemble des $x\in X$ où $f$ est plat
est ouvert. On peut donc se borner au cas où $f$ est plat, et alors le
corollaire résulte de (12.1.1, (iv), (vi) et (vii)) et de (12.1.6, (i), (ii)
et (iv)).
\end{proof}

\begin{remarks}[12.1.8]
\label{IV.12.1.8-fr}
\begin{enumerate}[label=(\roman*)]
  \item Dans (12.1.6, (ii)), on ne peut supprimer l'hypothèse que $x$
  n'appartienne à aucun cycle premier associé immergé de $X_{f(x)}$. C'est ce
  que montre l'exemple (5.12.3), où l'on prend $X=\operatorname{Spec}(A)$, $Y$
  étant le spectre de l'anneau local $A_0$ de $k[T]$ correspondant à l'idéal
  premier $(T)$ et le morphisme $X\to Y$ correspondant à l'homomorphisme
  injectif $A_0\to A$, déduit par localisation et passage au quotient de
  l'injection $k[T]\to k[T,U,V]$~; ce morphisme est plat puisque $A$ est un
  $A_0$-module sans torsion ($\mathbf 0_{\mathrm I}$, 6.3.4). Ici la fibre
  $X_y$ au point fermé $y$ de $Y$, égale à $\operatorname{Spec}(A/tA)$, est
  irréductible, de dimension 1 et vérifie la condition $(R_0)$ géométrique
  puisque l'anneau local en son point générique est un corps. Par contre, au
  point générique $y'$ de $Y$, la fibre $X_{y'}$ a deux composantes
  irréductibles de dimensions respectives 0 et 1, et celle de dimension 0
  n'est pas réduite, donc $X_{y'}$ ne vérifie pas la condition $(R_0)$.
  \item Dans (12.1.6, (ii)), on ne peut pas non plus supprimer l'hypothèse que
  $X_{f(x)}$ est équidimensionnel au point $x$. On le voit ici sur l'exemple
  (5.12.6) avec $X=\operatorname{Spec}(A)$, $Y=\operatorname{Spec}(A_0)$, où
  $A_0$ est défini comme dans (i), le morphisme $X\to Y$ provenant encore de
  l'injection $k[T]\to k[T,U,V,W]$ par localisation et passage au quotient, et
  étant plat puisque $A$ est un $A_0$-module sans torsion
  ($\mathbf 0_{\mathrm I}$, 6.3.4). La fibre $X_y$ au point fermé $y$ de $Y$
  est réduite (donc vérifie $(S_1)$) et vérifie $(R_1)$, mais a deux composantes
  irréductibles de dimensions 2 et 1, tandis que la fibre $X_{y'}$ au point
  générique $y'$ de $Y$ ne vérifie pas $(R_1)$.
\end{enumerate}
\end{remarks}
\subsection{Propriétés locales et globales des fibres d'un morphisme propre,
plat et de présentation finie}
\label{subsection:IV.12.2-fr}

\begin{theorem}[12.2.1]
\label{IV.12.2.1-fr}
Soient $f:X\to Y$ un morphisme \emph{propre} de présentation finie,
$\mathcal F$ un $\mathcal O_X$-Module $f$-plat et de présentation finie,
$\Phi$ une partie finie de $\mathbf Z\cup\{\pm\infty\}$, $k$ un entier. Les
parties suivantes de $Y$ sont ouvertes~:
\begin{enumerate}[label=(\roman*)]
  \item L'ensemble des $y\in Y$ tels que l'ensemble des dimensions des cycles
  premiers associés à $\mathcal F_y$ soit contenu dans $\Phi$.
  \item L'ensemble des $y\in Y$ tels que l'ensemble des dimensions des
  composantes irréductibles de $\operatorname{Supp}(\mathcal F_y)$ contienne
  $\Phi$.
  \item L'ensemble des $y\in Y$ tels que tous les cycles premiers associés à
  $\mathcal F_y$ aient une même dimension égale à $k$.
  \item L'ensemble des $y\in Y$ tels que $\mathcal F_y$ n'ait pas de cycle
  premier associé immergé et que l'ensemble des dimensions des composantes
  irréductibles de $\operatorname{Supp}(\mathcal F_y)$ soit égal à $\Phi$.
  \oldpage[IV-3]{180}
  \item L'ensemble des $y\in Y$ tels que $\mathcal F_y$ ait la propriété
  $(S_k)$ et soit équidimensionnel en tout point de $X_y$.
  \item L'ensemble des $y\in Y$ tels que
  $\operatorname{coprof}(\mathcal F_y)\leq k$.
  \item L'ensemble des $y\in Y$ tels que $\mathcal F_y$ soit un
  $\mathcal O_{X_y}$-Module de Cohen-Macaulay.
  \item L'ensemble des $y\in Y$ tels que $\mathcal F_y$ soit géométriquement
  réduit.
  \item L'ensemble des $y\in Y$ tels que $\mathcal F_y$ soit géométriquement
  ponctuellement intègre en chaque point de $X_y$.
  \item L'ensemble des $y\in Y$ tels que $\mathcal F_y$ soit géométriquement
  intègre.
  \item L'ensemble des $y\in Y$ tels que $\mathcal F_y$ n'ait pas de cycle
  premier associé immergé et que la somme $l(y)$ des multiplicités totales
  (4.7.12) de $\mathcal F_y$ aux points maximaux de
  $\operatorname{Supp}(\mathcal F_y)$ soit $\leq k$.
\end{enumerate}
\end{theorem}

\begin{proof}
À l'exception de (ii), (iii), (iv), (x) et (xi), les propriétés $P(y)$
considérées sont de la forme suivante~: «~pour tout $x\in X_y$, $\mathcal F_y$
a la propriété $Q(x)$~», où $Q(x)$ est l'une des propriétés (i) à (viii) de
(12.1.1). Il résulte de (12.1.1) que l'ensemble $U$ des $x\in X$ tels que
$Q(x)$ soit vraie est ouvert, et l'ensemble $V$ des $y\in Y$ tels que $P(y)$
soit vraie n'est autre que l'ensemble $Y-f(X-U)$~; dans tous ces cas, le
théorème est donc déjà vrai en supposant seulement le morphisme $f$ \emph{fermé}.
Pour (iii), on applique (i) avec $\Phi$ réduit à un seul élément. Il reste à
examiner les cas (ii), (iv) et (xi) séparément, ((x) se déduisant aussitôt de
(xi) et de (4.7.14)) en appliquant toujours la méthode décrite dans (12.0.2).
\end{proof}

\begin{env}[12.2.1.1]
\label{IV.12.2.1.1-fr}
Cas (ii)~: Les étapes A) et B) de la démonstration procèdent comme dans le
début de (12.1.1)~; pour l'étape A), on utilise (8.9.1), (11.2.6) ainsi que
(4.2.7)~; pour l'étape B), on utilise (9.5.5) appliqué à
$\operatorname{Supp}(\mathcal F)$. Il reste à démontrer que si (en supposant
$Y$ noethérien) l'ensemble des dimensions des composantes irréductibles de
$\operatorname{Supp}(\mathcal F_y)$ contient $\Phi$, alors il en est de même
de l'ensemble des dimensions des composantes irréductibles de
$\operatorname{Supp}(\mathcal F_{y'})$, pour toute \emph{générisation} $y'$ de
$y$ dans $Y$. Soit $Y_1$ un spectre d'anneau de valuation discrète tel que, si
$y_1$ et $y'_1$ sont le point fermé et le point générique de $Y_1$, il y ait un
morphisme $h:Y_1\to Y$ tel que $h(y_1)=y$, $h(y'_1)=y'$ (\textbf{II}, 7.1.7).
Appliquant (4.2.7), on voit qu'on peut remplacer $Y$ par $Y_1$ et $X$ par
$X_1=X\times_YY_1$, autrement dit, se borner au cas où $Y$ est spectre
d'anneau de valuation discrète, $y$ son point fermé et $y'$ son point générique.

Utilisant (\textbf{Err}\textsubscript{\textbf{III}}, 30), on peut se borner au
cas où $\operatorname{Supp}(\mathcal F)=X$, de sorte que $f$ est quasi-plat
(2.3.3)~; les composantes irréductibles $Z_i$ de $X$ dominent alors $Y$
(2.3.4), autrement dit leurs points génériques appartiennent à $X_{y'}$, et
les $Z_i\cap X_{y'}$ sont les composantes irréductibles de $X_{y'}$. Mais toute
composante irréductible de $X_y$ est contenue dans un des $Z_i$, donc est une
composante irréductible de $(Z_i)_y$~; or, on sait (7.1.13) que les dimensions
de toutes les composantes irréductibles de $(Z_i)_y$ sont égales à
$\dim((Z_i)_{y'})$, ce qui achève de prouver (ii).
\end{env}

\begin{env}[12.2.1.2]
\label{IV.12.2.1.2-fr}
Cas (iv)~: Le même raisonnement qu'au début montre, en utilisant (12.1.1,
(iii)) que l'on peut déjà supposer (en remplaçant $Y$ par un ouvert de $Y$) que
pour tout $y\in Y$, $\mathcal F_y$ n'a pas de cycle premier associé immergé.
L'assertion du cas (iv) est alors une conséquence immédiate des assertions des
cas (i) et (ii).
\end{env}

\begin{env}[12.2.1.3]
\label{IV.12.2.1.3-fr}
\oldpage[IV-3]{181}
Cas (xi)~: Pour l'étape A) du raisonnement, on utilise (8.9.1), (11.2.6),
(8.10.5, (xii)) (pour conserver l'hypothèse que $f$ est propre), ainsi que
(4.2.7), (4.5.6) et (4.7.9). Pour l'étape B), on voit, comme dans le cas (iv),
qu'on peut supposer que pour tout $y\in Y$, $\mathcal F_y$ n'a pas de cycle
premier associé immergé, et l'on applique (9.8.8) qui prouve que la fonction
$z\mapsto l(z)$ est constructible. Il reste donc à voir (en supposant $Y$
noethérien et $f$ propre) que pour toute générisation $y'$ d'un point $y\in Y$,
on a $l(y')\leq l(y)$. En raisonnant comme dans (12.1.1.3), on voit (à l'aide
de (4.7.8) et (4.5.11)) qu'on peut supposer que les composantes irréductibles
de $\operatorname{Supp}(\mathcal F_{y'})$ sont géométriquement irréductibles
et que $l(y')$ est la \emph{somme des longueurs} des
$(\mathcal F_{y'})_{z'_j}$ aux points maximaux $z'_j$ de
$\operatorname{Supp}(\mathcal F_{y'})$. Appliquant de nouveau (4.2.7), (4.5.6)
et (4.7.9), on se ramène, comme dans (12.2.1.1), au cas où $Y$ est un spectre
d'anneau de valuation discrète, $y$ son point fermé et $y'$ son point générique.
Le fait que $l(y')\leq l(y)$ sera alors conséquence du

\begin{lemma}[12.2.1.4]
\label{IV.12.2.1.4-fr}
Soient $Y$ un spectre d'anneau de valuation discrète, $y$ son point fermé,
$y'$ son point générique, $f:X\to Y$ un morphisme propre, $\mathcal F$ un
$\mathcal O_X$-Module cohérent et $f$-plat, $z_i$ (resp. $z'_j$) les points
maximaux de $\operatorname{Supp}(\mathcal F_y)$ (resp.
$\operatorname{Supp}(\mathcal F_{y'})$). Supposons que $\mathcal F_y$ n'ait
pas de cycle premier associé immergé. Alors on a
\[
  \sum_j\operatorname{long}((\mathcal F_{y'})_{z'_j})
  \leq\sum_i\operatorname{long}((\mathcal F_y)_{z_i}).
\]
\end{lemma}

On a $Y=\operatorname{Spec}(A)$, où $A$ est un anneau de valuation discrète,
dont nous désignerons par $t$ une uniformisante, de sorte que
$\mathcal F_y=\mathcal F/t\mathcal F$. Comme $\mathcal F$ est $f$-plat, les
$z'_j$ sont aussi les points maximaux de $\operatorname{Supp}(\mathcal F)$
(2.3.4)~; pour tout $z_i$, désignons par $z'_{ik}$ ceux des $z'_j$ qui sont des
générisations de $z_i$~; il résulte de (3.4.1.1) que l'on a
\[
  \operatorname{long}((\mathcal F_y)_{z_i})
  \geq\sum_k\operatorname{long}((\mathcal F_{y'})_{z'_{ik}})
\]
d'où en sommant
\[
  \sum_i\operatorname{long}((\mathcal F_y)_{z_i})
  \geq\sum_i\sum_k\operatorname{long}((\mathcal F_{y'})_{z'_{ik}}).
\]

Le lemme sera donc prouvé si nous établissons que pour tout $z'_j$, il y a au
moins un indice $i$ tel que $z'_j$ soit l'un des $z'_{ik}$. Or, comme $f$ est
propre (donc fermé) et $f(z'_j)=y'$, il existe $x\in X$ qui est une
spécialisation de $z'_j$ et est tel que $f(x)=y$, autrement dit
$\overline{\{z'_j\}}\cap X_y$ n'est pas vide. Comme $t$ est
$\mathcal F$-régulier par platitude, on déduit de (3.4.3) qu'il y a au moins un
point de $\operatorname{Ass}(\mathcal F_y)$ dont $z'_j$ est une
générisation~; mais comme les points de $\operatorname{Ass}(\mathcal F_y)$
sont par hypothèse les $z_i$, cela termine la démonstration.
\end{env}

\begin{corollary}[12.2.2]
\label{IV.12.2.2-fr}
Sous les hypothèses de (12.2.1)~:
\begin{enumerate}[label=(\roman*)]
  \item La fonction $y\mapsto\dim(\operatorname{Supp}(\mathcal F_y))$ est
  continue (donc localement constante) dans $Y$.
  \item La fonction $y\mapsto\operatorname{coprof}(\mathcal F_y)$ (5.7.1) est
  semi-continue supérieurement dans $Y$.
  \item La fonction $y\mapsto l(y)$ (12.2.1, (xi)) est semi-continue
  supérieurement dans $Y$ lorsque les $\mathcal F_y$ n'ont pas de cycle
  premier associé immergé.
  \oldpage[IV-3]{182}
  \item La fonction $y\mapsto\operatorname{prof}^{*}(\mathcal F_y)$ (10.8.1)
  est semi-continue inférieurement dans $Y$.
\end{enumerate}
\end{corollary}

\begin{proof}
Il résulte de (12.2.1, (i)) appliqué à $\Phi$ égal au plus petit intervalle de
$\mathbf N$ contenant les dimensions des cycles premiers associés à
$\mathcal F_y$, que $z\mapsto\dim(\operatorname{Supp}(\mathcal F_z))$ est
semi-continue supérieurement au point $y$~; il résulte d'autre part de
(12.2.1, (ii)) appliqué à $\Phi$ égal à l'ensemble des dimensions des
composantes irréductibles de $\operatorname{Supp}(\mathcal F_y)$, que cette
même fonction est semi-continue inférieurement au point $y$, d'où la
conclusion. Les assertions (ii) et (iii) ne sont que d'autres formulations de
(12.2.1, (vi) et (xi)). Enfin, d'après (12.1.3), l'ensemble des $x\in X$ tels
que $\operatorname{prof}^{*}_{f(x)}(\mathcal F_{f(x)})\leq k$ est ouvert, et la
conclusion de (iv) en résulte par le même raisonnement qu'au début de (12.2.1).
\end{proof}

\begin{remarks}[12.2.3]
\label{IV.12.2.3-fr}
\begin{enumerate}[label=(\roman*)]
  \item On observera que pour (12.2.1, (ii)), on peut se dispenser de
  l'hypothèse que $f$ est \emph{propre}.

  \item Dans (12.2.1, (ii)), on ne peut remplacer la condition que l'ensemble
  $E(y)$ des dimensions des composantes irréductibles de
  $\operatorname{Supp}(\mathcal F_y)$ contienne $\Phi$, par la condition que
  $E(y)$ soit \emph{égal à} $\Phi$. En effet, soit $k$ un corps, et considérons
  l'espace projectif à 3 dimensions $P=\mathbf P^3_k=\operatorname{Proj}(C)$,
  où $C=k[t_0,t_1,t_2,t_3]$ ($t_i$ indéterminées)~; dans $P$, soit $X_0$ le
  sous-préschéma fermé «~réunion de la droite $X_1$ définie par $t_1=t_3=0$ et
  du plan $X_2$ défini par $t_2-t_3=0$~», ce qui décrit en termes géométriques
  le fait que $X_0$ est égal à $\operatorname{Proj}(C/\mathfrak a)$, où l'idéal
  gradué $\mathfrak a$ est égal à $\mathfrak b\mathfrak c$, avec
  $\mathfrak b=Ct_1+Ct_3$, $\mathfrak c=C(t_2-t_3)$~; considérons le morphisme
  $f_0:X_0\to X_1$ qui, en termes géométriques, est la «~projection de $X_0$
  sur $X_1$ à partir de la droite à l'infini $D$ définie par $t_0=t_2=0$~»~;
  sous forme algébrique, $f_0$ correspond à l'homomorphisme d'algèbres graduées
  \[
    k[t_0,t_2]\longrightarrow k[t_0,t_1,t_2,t_3]/\mathfrak a
  \]
  obtenu par passage au quotient à partir de l'injection canonique
  $k[t_0,t_2]\to k[t_0,t_1,t_2,t_3]$. Il est clair que $f_0$ est un morphisme
  projectif, donc \emph{propre}~; il en est donc de même de sa restriction
  $f:X\to Y$, où $Y=D_+(t_0)$ et $X=f_0^{-1}(Y)$~; pour voir que $f$ est
  \emph{plat}, il suffit de remarquer que $k[t_2]$ est un anneau principal et
  l'anneau
  \[
    C'=k[t_1,t_2,t_3]/(t_2-t_3)((t_1)+(t_3))
  \]
  un $k[t_2]$-module sans torsion ($\mathbf 0_{\mathrm I}$, 6.3.4). Pour tout
  $y\in Y$ distinct du point $y_0$ défini par $t_2=0$, $f^{-1}(y)$ a alors
  deux composantes irréductibles, de dimensions 0 et 1, tandis que
  $f^{-1}(y_0)$ n'a qu'une seule composante irréductible de dimension 1.

  \item Dans (12.2.1, (i)), on ne peut pas non plus remplacer la condition que
  l'ensemble $E'(y)\supset E(y)$ des dimensions des cycles premiers associés à
  $\mathcal F_y$ soit contenu dans $\Phi$ par la condition qu'il soit
  \emph{égal à} $\Phi$. Soient en effet $k$ un corps, $A$ l'anneau $k[t]$ de
  polynômes, $B$ le sous-anneau $k[t^4,t^3u,tu^3,u^4]$ de l'anneau de polynômes
  $k[t,u]$ ($t,u$ indéterminées)~; $B$ n'est pas intégralement clos, l'élément
  $t^2u^2$ appartenant à sa clôture intégrale mais n'étant pas dans $B$~; si
  $\mathfrak m$ est l'idéal maximal de $B$, engendré par $t^4,t^3u,tu^3$ et
  $u^4$, $\mathfrak m$ est idéal premier associé immergé de $A/At^4$. Posons
  $Y=\operatorname{Spec}(A)$, $X=\operatorname{Spec}(B)$, et soit $f:X\to Y$
  le morphisme correspondant à l'homomorphisme $A\to B$ qui transforme $t$ en
  $t^4$. Comme cet homomorphisme fait de $B$ un $A$-module sans torsion, $f$
  est plat ($\mathbf 0_{\mathrm I}$, 6.3.4). Si $y$ est le point de $Y$
  correspondant à l'idéal maximal $tA$, le préschéma $f^{-1}(y)$ est
  irréductible et de dimension 1, mais admet un cycle premier associé immergé
  de dimension 0~; au contraire, si $y'$ est le point générique de $Y$,
  $f^{-1}(y')$ n'a pas de cycle premier associé immergé, étant le spectre d'une
  $k(y')$-algèbre intègre de dimension 1. Dans cet exemple, $f$ est un morphisme
  affine non propre~; on peut en déduire aussitôt un exemple analogue où $f$
  est propre et plat en considérant $X$ comme un ouvert partout dense d'un
  schéma projectif sur $Y$ (\textbf{II}, 5.3.4 et 5.3.2), ou en procédant
  directement comme dans la Remarque (ii).

  Les Remarques (ii) et (iii) montrent la nécessité d'inclure dans (12.2.1,
  (iv)) la condition que $\mathcal F_y$ n'ait pas de cycle premier associé
  immergé.

  \item Dans (12.2.1, (xi)), on ne peut supprimer l'hypothèse que $f$ soit
  propre. En effet, soient $k$ un corps, $Y$ la «~droite affine~», spectre de
  l'anneau de polynômes $k[t]$ ($t$ indéterminée), $X$ le schéma \emph{somme}
  de $Y$ et de l'ouvert complémentaire du point fermé $y$ de $Y$ défini par
  $t=0$. Il est clair que le morphisme $f:X\to Y$ qui est égal aux injections
  canoniques dans chacune des composantes de $X$ est plat et que si l'on prend
  $\mathcal F=\mathcal O_X$, $\mathcal F_z$ n'a de cycle premier associé
  immergé pour aucun $z\in Y$. Toutefois, on voit aussitôt que l'on a $l(y)=1$
  et $l(z)=2$ pour tout $z\neq y$.

  \item Dans (12.2.1, (xi)), on ne peut pas non plus supprimer l'hypothèse que
  $\mathcal F_y$ soit sans cycle premier associé immergé. C'est ce que montre
  l'exemple de la Remarque (ii)~: on voit aussitôt en effet que l'on a, avec
  $\mathcal F=\mathcal O_X$, $l(y_0)=1$ et $l(y)=2$ pour $y\neq y_0$ dans $Y$.

  \item Nous ignorons si dans (12.2.1, (xi)) on peut remplacer l'hypothèse que
  $\mathcal F$ est $f$-plat par celle que $\operatorname{Supp}(\mathcal F)=X$
  et que $f$ est universellement ouvert. En reprenant la démonstration de
  (12.2.1.4), on voit (en utilisant aussi (14.3.6)) qu'il faudrait résoudre la
  question suivante~: Soient $A$ un anneau local noethérien, $M$ un $A$-module
  \oldpage[IV-3]{183}
  de type fini, $t$ un élément de l'idéal maximal $\mathfrak m$ de $A$, qui
  n'est contenu dans aucun idéal premier minimal de $A$~; s'il existe un de ces
  idéaux premiers minimaux $\mathfrak p$ tels que $\dim(A/\mathfrak p)=1$,
  est-il vrai que $\mathfrak m$ est un idéal premier associé à $M/tM$ ($t$
  n'étant pas supposé $M$-régulier)~?

  On notera toutefois qu'on ne peut, dans (12.2.1, (xi)), supprimer purement et
  simplement l'hypothèse que $\mathcal F$ est $f$-plat. On prendra en effet ici
  $P$ comme dans la Remarque (ii), et dans $P$ le sous-préschéma fermé $X_0$
  réunion des trois droites $X_1,X_2,X_3$ définies respectivement par
  $t_1=t_3=0$, $t_1-t_0=t_3=0$, $t_2=t_3=0$~; on définit la projection $f_0$ de
  $X_0$ sur $X_1$ comme précédemment et sa restriction $f:X\to Y$, où
  $Y=D_+(t_0)$ est la droite affine et $X=f_0^{-1}(Y)$~; $f$ est propre mais
  non plat, et si l'on prend $\mathcal F=\mathcal O_X$, $\mathcal F_y$ n'a de
  cycles premiers associés immergés pour aucun $y\in Y$. Mais on a $l(y_0)=1$
  et $l(y)=2$ pour $y\neq y_0$ dans $Y$.
\end{enumerate}
\end{remarks}

\begin{theorem}[12.2.4]
\label{IV.12.2.4-fr}
Soient $f:X\to Y$ un morphisme propre, plat et de présentation finie, $k$ un
entier $\geq 1$. Les parties suivantes de $Y$ sont ouvertes~:
\begin{enumerate}[label=(\roman*)]
  \item L'ensemble des $y\in Y$ tels que $X_y$ possède la propriété $(S_k)$.
  \item L'ensemble des $y\in Y$ tels que $X_y$ vérifie la propriété $(R_k)$
  géométrique, soit équidimensionnel en chaque point et n'ait pas de cycle
  premier associé immergé.
  \item L'ensemble des $y\in Y$ tels que $X_y$ soit géométriquement régulier
  (i.e. lisse sur $k(y)$).
  \item L'ensemble des $y\in Y$ tels que $X_y$ soit géométriquement normal.
  \item L'ensemble des $y\in Y$ tels que $X_y$ soit géométriquement réduit.
  \item L'ensemble des $y\in Y$ tels que $X_y$ soit géométriquement réduit et
  que le nombre géométrique de composantes connexes de $X_y$ soit égal à $k$.
  \item L'ensemble des $y\in Y$ tels que $X_y$ soit géométriquement
  ponctuellement intègre (4.6.9).
  \item L'ensemble des $y\in Y$ tels que $X_y$ soit géométriquement intègre.
  \item L'ensemble des $y\in Y$ tels que $X_y$ n'ait pas de cycle premier
  associé immergé et que la multiplicité totale (4.7.4) de $X_y$ sur $k(y)$
  soit $\leq k$.
\end{enumerate}
\end{theorem}

\begin{proof}
Sauf pour (vi), ces assertions sont des cas particuliers d'assertions de
(12.2.1), ou se déduisent d'assertions de (12.1.6) comme au début de la
démonstration de (12.2.1). Pour (vi), on se ramène comme toujours (compte tenu
de l'invariance du nombre géométrique de composantes connexes par changement
de base (4.5.6)) au cas où $Y$ est noethérien. L'ensemble $U$ des $y\in Y$ tels
que $X_y$ soit géométriquement réduit est ouvert dans $Y$ en vertu de (v)~; il
résulte alors de (\textbf{III}, 7.8.7 et 7.8.6) que pour tout $y\in U$, il y a
un voisinage $V\subset U$ de $y$ et un entier $m$ tels que, pour \emph{tout}
$z\in V$, $\Gamma(X_z,\mathcal O_{X_z})$ soit isomorphe à $(k(z))^m$~; mais en
vertu de (\textbf{III}, 4.3.4), $m$ est alors le nombre géométrique de
composantes connexes de $X_z$, d'où la conclusion. On donnera une autre
démonstration de (vi) dans (15.5.9).
\end{proof}

\subsection{Propriétés cohomologiques locales des fibres d'un morphisme plat et
localement de présentation finie}
\label{subsection:IV.12.3-fr}

\begin{lemma}[12.3.1]
\label{IV.12.3.1-fr}
Soient $A$ un anneau, $B$ une $A$-algèbre de présentation finie qui soit un
$A$-module plat, $M$ un $B$-module de présentation finie, qui soit un
$A$-module plat. Alors les conditions suivantes sont équivalentes~:
\begin{enumerate}[label=\emph{\alph*)}]
  \item $M$ est un $B$-module projectif.
  \item $M$ est un $B$-module plat.
  \item Pour tout $y\in\operatorname{Spec}(A)$, $M\otimes_Ak(y)$ est un
  $(B\otimes_Ak(y))$-module projectif.
\end{enumerate}
\end{lemma}

\begin{proof}
\oldpage[IV-3]{184}
L'équivalence de $a)$ et $b)$ résulte de Bourbaki, \emph{Alg. comm.}, chap. II,
\S~5, n\textsuperscript{o} 2, cor. 2 du th. 1. Comme $a)$ implique $c)$
trivialement, il reste à prouver que $c)$ implique $b)$, ce qui résulte du
critère de platitude par fibres (11.3.10), appliqué avec $g=h$, $f=1_X$.
\end{proof}

\begin{proposition}[12.3.2]
\label{IV.12.3.2-fr}
Soient $A$ un anneau, $B$ une $A$-algèbre de présentation finie qui soit un
$A$-module plat, $M$ un $B$-module de présentation finie qui soit un
$A$-module plat. Alors~:
\begin{enumerate}[label=(\roman*)]
  \item Il existe une résolution gauche de $M$ par des $B$-modules libres de
  type fini.
  \item On a
  \begin{equation}
    \tag{12.3.2.1}
    \label{IV.12.3.2.1-fr}
    \operatorname{dim.proj}_B(M)
    =\sup_{y\in\operatorname{Spec}(A)}
    \operatorname{dim.proj}_{B\otimes_Ak(y)}(M\otimes_Ak(y))
    =\operatorname{Tor.dim}_B(M)
  \end{equation}
  où $\operatorname{Tor.dim}_B(M)$ est le \emph{plus petit entier} $i$ tel que
  $\operatorname{Tor}^B_j(M,N)=0$ pour tout $j\geq i$ et tout $B$-module $N$
  (et $+\infty$ si un tel entier $i$ n'existe pas).
\end{enumerate}
\end{proposition}

\begin{proof}
\begin{enumerate}[label=(\roman*)]
  \item En vertu de (8.9.1), il existe un sous-anneau noethérien $A_0$ de $A$,
  une $A_0$-algèbre de type fini $B_0$ et un $B_0$-module de type fini $M_0$
  tels que $B$ soit isomorphe à $B_0\otimes_{A_0}A$ et $M$ à
  $M_0\otimes_{A_0}A$. De plus (11.2.7) on peut supposer que $B_0$ et $M_0$
  sont des $A_0$-modules plats. Il y a alors une résolution gauche
  $(L_0)_\bullet$ de $M_0$ formée de $B_0$-modules libres de type fini, et
  comme $M_0$ et les $L_{0i}$ sont des $A_0$-modules plats,
  $(L_0)_\bullet\otimes_{A_0}A$ est une résolution gauche de $M$ formée de
  $B$-modules libres de type fini (2.1.10).

  \item En vertu de ($\mathbf 0$, 17.2.2, (ii)), on a
  $\operatorname{Tor.dim}_B(M)\leq\operatorname{dim.proj}_B(M)$, et la
  définition de la dimension projective prouve aussitôt que, pour tout
  $y\in\operatorname{Spec}(A)$, on a
  \[
    \operatorname{dim.proj}_{B\otimes_Ak(y)}(M\otimes_Ak(y))
    \leq\operatorname{dim.proj}_B(M).
  \]
  Pour prouver les inégalités inverses, considérons une résolution gauche
  $(L_i)$ de $M$ par des $B$-modules libres de type fini, et supposons que
  $\operatorname{Tor.dim}_B(M)=n$ (resp.
  $\operatorname{dim.proj}_{B\otimes_Ak(y)}(M\otimes_Ak(y))\leq n$ pour tout
  $y\in\operatorname{Spec}(A)$). Alors
  $R=\operatorname{Im}(L_n\to L_{n-1})=\operatorname{Ker}(L_{n-1}\to L_{n-2})$
  est un $B$-module de type fini et aussi un $A$-module plat, en vertu de
  l'hypothèse sur $M$ et $B$ et de (2.1.10). En outre, on a
  $\operatorname{Tor}^B_{n+1}(M,N)=\operatorname{Tor}^B_1(R,N)$ pour tout
  $B$-module $N$ (M, V, 7). L'hypothèse $\operatorname{Tor.dim}_B(M)=n$
  entraîne donc $\operatorname{Tor}^B_1(R,N)=0$ pour tout $B$-module $N$,
  c'est-à-dire que $R$ est un $B$-module plat, donc projectif en vertu de
  (12.3.1)~; cela établit que $\operatorname{dim.proj}_B(M)\leq n$.
  L'hypothèse
  $\operatorname{dim.proj}_{B\otimes_Ak(y)}(M\otimes_Ak(y))\leq n$ pour tout
  $y\in\operatorname{Spec}(A)$ entraîne d'autre part, par tensorisation avec
  $k(y)$, que dans chacune des suites (exactes en vertu de la platitude sur $A$
  de $M$, des $L_i$ et de $R$ (2.1.10))
  \[
    0\to R\otimes_Ak(y)\to L_{n-1}\otimes_Ak(y)\to\cdots
    \to L_0\otimes_Ak(y)\to M\otimes_Ak(y)\to0
  \]
  $R\otimes_Ak(y)$ est un $(B\otimes_Ak(y))$-module projectif (pour
  $y\in\operatorname{Spec}(A)$). On conclut alors encore de (12.3.1) que $R$
  est un $B$-module projectif, donc $\operatorname{dim.proj}_B(M)\leq n$.
\end{enumerate}
\end{proof}

\begin{proposition}[12.3.3]
\label{IV.12.3.3-fr}
Soient $f:X\to S$ un morphisme localement de présentation finie,
$\mathcal L_\bullet$ un complexe formé de $\mathcal O_X$-Modules
quasi-cohérents de présentation finie~; pour tout $s\in S$, soit
$(\mathcal L_\bullet)_s$ le complexe
$\mathcal L_\bullet\otimes_{\mathcal O_S}k(s)$ de $\mathcal O_{X_s}$-Modules de
type fini. Supposons que $\mathcal L_{n-1}$ soit $f$-plat. Alors l'ensemble $U$
des $x\in X$ tels que $(\mathcal H_n((\mathcal L_\bullet)_{f(x)}))_x=0$ est
ouvert dans $X$. Si de plus $\mathcal L_n$ est $f$-plat, alors on a
$\mathcal H_n(\mathcal L_\bullet)|U=0$.
\end{proposition}

\begin{proof}
On peut évidemment se limiter au cas où $n=1$ et où le complexe
$\mathcal L_\bullet$ se réduit à
$0\to\mathcal L_2\xrightarrow{u}\mathcal L_1\xrightarrow{v}\mathcal L_0\to0$.
On peut d'abord se ramener au cas où $S=\operatorname{Spec}(A)$ et $X$
\oldpage[IV-3]{185}
sont affines, puis au cas où $S$ est \emph{noethérien}~; en effet, par (8.9.1)
et (8.5.2), on sait qu'il existe un préschéma noethérien
$S'=\operatorname{Spec}(A')$, où $A'$ est un sous-anneau de $A$, un morphisme
de type fini $f':X'\to S'$ et trois $\mathcal O_{X'}$-Modules cohérents
$\mathcal L'_i$ ($0\leq i\leq2$) tels que $X=X'\times_{S'}S$,
$f=f'_{(S)}$, $\mathcal L_i=\mathcal L'_i\otimes_{S'}S$, ainsi que deux
homomorphismes $u'$, $v'$ tels que $u=u'\otimes1$, $v=v'\otimes1$ et
$v'\circ u'=0$. On peut en outre supposer que $\mathcal L'_0$ est $f'$-plat
(11.2.7), et que $\mathcal L'_1$ est $f'$-plat lorsque $\mathcal L_1$ est
supposé $f$-plat. Par fidèle platitude, l'hypothèse
$(\mathcal H_n((\mathcal L_\bullet)_{f(x)}))_x=0$ équivaut à
$(\mathcal H_n((\mathcal L'_\bullet)_{f'(x')}))_{x'}=0$, si $x'$ est la
projection de $x$ dans $X'$~; si $U'$ est l'ensemble des $x'\in X'$ tels que
$(\mathcal H_n((\mathcal L'_\bullet)_{f'(x')}))_{x'}=0$, $U$ est donc l'image
réciproque de $U'$ par la projection $X\to X'$, ce qui ramène, pour la première
assertion, au cas où $S$ est noethérien. Pour la seconde assertion, on remarque
en outre que $A$ est limite inductive de ses sous-anneaux $A_\lambda$ qui sont
des $A'$-algèbres de type fini~; posons
$X_\lambda=X'\times_{S'}S_\lambda$, où
$S_\lambda=\operatorname{Spec}(A_\lambda)$,
$\mathcal L_i^{(\lambda)}=\mathcal L'_i\otimes_{S'}S_\lambda$, et soient
$u_\lambda=u'\otimes1:\mathcal L_2^{(\lambda)}\to\mathcal L_1^{(\lambda)}$,
$v_\lambda=v'\otimes1:\mathcal L_1^{(\lambda)}\to\mathcal L_0^{(\lambda)}$,
de sorte que l'on a
$\mathcal H_1(\mathcal L_\bullet^{(\lambda)})
=\operatorname{Ker}(v_\lambda)/\operatorname{Im}(u_\lambda)$. Or, comme le
foncteur $\varinjlim$ est exact dans la catégorie des groupes commutatifs, on a
$\operatorname{Ker}(v)=\varinjlim\operatorname{Ker}(v_\lambda)$,
$\operatorname{Im}(u)=\varinjlim\operatorname{Im}(u_\lambda)$, et
$\mathcal H_1(\mathcal L_\bullet)
=\varinjlim\mathcal H_1(\mathcal L_\bullet^{(\lambda)})$. Si l'on a supposé
que $\mathcal L_1$ est $f$-plat et (en se ramenant au cas où $X=U$) que l'on
ait prouvé $\mathcal H_1(\mathcal L_\bullet^{(\lambda)})=0$ pour tout
$\lambda$, on en déduira bien $\mathcal H_1(\mathcal L_\bullet)=0$.

\medskip
I) Supposons désormais $S$ noethérien. On sait (sans hypothèse de platitude sur
$\mathcal L_0$) que l'ensemble $U$ est \emph{constructible} dans $X$ (9.9.6).
Utilisant maintenant ($\mathbf 0_{\mathrm{III}}$, 9.2.5), il reste à montrer
que pour toute \emph{générisation} $x'$ de $x\in U$ dans $X$, on a aussi
$x'\in U$. La méthode exposée dans (12.0.2) s'applique sans changement (compte
tenu du fait que pour un préschéma $Z$ sur un corps $k$ et une extension $k'$
de $k$, la projection $Z\otimes_k k'\to Z$ est fidèlement plate). On peut donc
supposer que $S$ est le spectre d'un anneau de valuation discrète $A$, dont on
désigne par $t$ une uniformisante, $x$ étant au-dessus du point fermé de $S$ et
$x'$ au-dessus du point générique de $S$. L'hypothèse que $\mathcal L_0$ est
$f$-plat entraîne que $t$ est $\mathcal L_0$-régulier
($\mathbf 0_{\mathrm I}$, 6.3.4). On est alors ramené à prouver le lemme suivant
(où $B$, $M$, $N$, $P$ seront remplacés par $\mathcal O_x$,
$(\mathcal L_2)_x$, $(\mathcal L_1)_x$ et $(\mathcal L_0)_x$ respectivement)~:

\begin{lemma}[12.3.3.1]
\label{IV.12.3.3.1-fr}
Soient $B$ un anneau local noethérien, $M,N,P$ trois $B$-modules, $N$ étant
supposé de type fini, $M\xrightarrow{u}N\xrightarrow{v}P$ deux homomorphismes
tels que $v\circ u=0$. Soit $t$ un élément de l'idéal maximal de $B$ tel que
$t$ soit $P$-régulier et que la suite
\begin{equation}
  \tag{12.3.3.2}
  \label{IV.12.3.3.2-fr}
  M/tM\xrightarrow{u\otimes1_{B/tB}}N/tN
  \xrightarrow{v\otimes1_{B/tB}}P/tP
\end{equation}
soit exacte. Alors la suite $M\xrightarrow{u}N\xrightarrow{v}P$ est exacte.
\end{lemma}

Notons d'abord que si l'on remplace $M$ par son image $Q$ dans $N$ et $u$ par
l'injection $j:Q\to N$, l'image de $j\otimes1$ est la même que celle de
$u\otimes1$, donc on peut, sans changer ni l'hypothèse ni la conclusion,
supposer $u$ injectif. D'autre part, si $R$ est l'image de $v$ et
$\rho:N\to R$ la surjection canonique, $t$ est évidemment $R$-régulier, on a
$\rho\circ u=0$, et le noyau de $\rho\otimes1$ est contenu dans celui de
$v\otimes1$~; il lui est par suite égal si la suite (12.3.3.2) est exacte~;
comme d'autre part on a évidemment $\operatorname{Ker}(\rho)
=\operatorname{Ker}(v)$, on voit qu'on peut, pour prouver le lemme, supposer en
outre $v$ surjectif. Le lemme sera alors conséquence du suivant~:

\begin{lemma}[12.3.3.3]
\label{IV.12.3.3.3-fr}
Soient $B$ un anneau, $M,N,P$ trois $B$-modules,
$u:M\to N$, $v:N\to P$ deux homomorphismes tels que $u$ soit injectif, $v$
surjectif et $v\circ u=0$. Soit $t$ un élément de $B$ tel que $t$ soit
$P$-régulier. Alors on a
\begin{equation}
  \tag{12.3.3.4}
  \label{IV.12.3.3.4-fr}
  \operatorname{Ker}(v\otimes1)/\operatorname{Im}(u\otimes1)
  =(\operatorname{Ker}(v)/\operatorname{Im}(u))\otimes_B(B/tB)
\end{equation}
à un isomorphisme canonique près.
\end{lemma}

En effet, l'hypothèse que la suite (12.3.3.2) est exacte entraînera alors
$(\operatorname{Ker}v/\operatorname{Im}u)\otimes_B(B/tB)=0$, et comme $t$
appartient à l'idéal maximal de $B$ et que $\operatorname{Ker}(v)$ est un
$B$-module de type fini (puisque $B$ est supposé noethérien dans (12.3.3.1)),
le lemme de Nakayama prouvera que $\operatorname{Ker}(v)=\operatorname{Im}(u)$.

Reste donc à démontrer (12.3.3.3). Posons $u'=u\otimes1$, $v'=v\otimes1$,
$Z=\operatorname{Ker}(v)$, $H=Z/\operatorname{Im}(u)$, de sorte qu'on a les
suites exactes
\begin{gather*}
  0\to M\to Z\to H\to0,\\
  0\to Z\to N\xrightarrow{v}P\to0,
\end{gather*}
\oldpage[IV-3]{186}
d'où, en tensorisant par $B/tB$ et utilisant le lemme (3.4.1.4) et le fait que
$t$ est $P$-régulier, les suites exactes
\begin{gather*}
  M/tM\xrightarrow{w'}Z/tZ\to H/tH\to0,\\
  0\to Z/tZ\to N/tN\xrightarrow{v'}P/tP\to0,
\end{gather*}
d'où $\operatorname{Ker}(v')=Z/tZ$, et comme
$\operatorname{Im}(w')=\operatorname{Im}(u')$, on obtient (12.3.3.4).

\medskip
II) Supposons maintenant en outre que $\mathcal L_1$ soit $f$-plat~;
remplaçant en outre éventuellement $X$ par $U$, on peut supposer que $X=U$.
Il s'agit de voir que pour tout $x\in X$, on a
$(\mathcal H_1(\mathcal L_\bullet))_x=0$. Soit $s=f(x)$, et soit
$\mathfrak n$ l'idéal $\mathfrak m_s\mathcal O_{X,x}$ qui est contenu dans
l'idéal maximal de $\mathcal O_{X,x}$~; comme
$(\mathcal H_1(\mathcal L_\bullet))_x$ est un $\mathcal O_{X,x}$-module de type
fini ($X$ étant localement noethérien), il est séparé pour la topologie
$\mathfrak n$-adique ($\mathbf 0_{\mathrm I}$, 7.3.5), donc il suffit de
montrer que son \emph{séparé complété} pour cette topologie est 0. Or, en vertu
de (\textbf{III}, 7.4.7.2), ce séparé complété est égal à
$\varprojlim_k\mathcal H_1(\mathcal L_\bullet\otimes_{\mathcal O_S}
(\mathcal O_{S,s}/\mathfrak m_s^{k+1}))$, et il suffira donc de prouver que
chacun des termes de ce système projectif est nul. Mais cela est vrai par
hypothèse pour $k=0$~; raisonnons donc par récurrence sur $k$. La conclusion
résultera du lemme plus général suivant (qu'on appliquera pour
$A=\mathcal O_{S,s}/\mathfrak m_s^{k+1}$ et $\mathfrak J$ égal à l'idéal
maximal $\mathfrak m_s/\mathfrak m_s^{k+1}$ de $A$)~:

\begin{lemma}[12.3.3.5]
\label{IV.12.3.3.5-fr}
Soient $A$ un anneau, $\mathfrak J$ un idéal nilpotent de $A$,
$L_\bullet:P\xrightarrow{u}Q\xrightarrow{v}R$ un complexe de $A$-modules. On
pose $A_0=A/\mathfrak J$,
$L_\bullet^{(0)}=L_\bullet\otimes_AA_0:
P_0\xrightarrow{u_0}Q_0\xrightarrow{v_0}R_0$, et l'on suppose que
$\operatorname{gr}_{\mathfrak J}^{\bullet}(A)$ est un $A_0$-module plat, et
que $Q$ et $R$ sont des $A$-modules plats. Alors la relation
$\mathrm H_1(L_\bullet^{(0)})=0$ entraîne $\mathrm H_1(L_\bullet)=0$.
\end{lemma}

Posons
$L_\bullet^{(k)}=L_\bullet\otimes_A(A/\mathfrak J^{k+1}):
P_k\xrightarrow{u_k}Q_k\xrightarrow{v_k}R_k$~; comme il existe un entier
$k\geq1$ tel que $L_\bullet^{(k)}=L_\bullet$, il suffira de prouver, par
récurrence sur $k$, que la suite
$P_k\xrightarrow{u_k}Q_k\xrightarrow{v_k}R_k$ est
\oldpage[IV-3]{187}
exacte (cette assertion découlant de l'hypothèse pour $k=0$). Soit donc
$x\in Q_k$ un élément tel que $v_k(x)=0$. Comme par l'hypothèse de récurrence
la suite
$P_{k-1}\xrightarrow{u_{k-1}}Q_{k-1}\xrightarrow{v_{k-1}}R_{k-1}$ est exacte,
l'image canonique de $x$ dans
$Q_{k-1}=Q_k/(\mathfrak J^kQ/\mathfrak J^{k+1}Q)$ appartient à
$\operatorname{Im}(u_{k-1})$, donc il existe $z\in P_k$ tel que
$x=u_k(z)+x'$ avec $x'\in\mathfrak J^kQ/\mathfrak J^{k+1}Q$. Comme on a
$v_k\circ u_k=0$, la relation $v_k(x)=0$ entraîne $v_k(x')=0$, et par ailleurs
on a évidemment $v_k(x')\in\mathfrak J^kR/\mathfrak J^{k+1}R$~; tout revient
donc à prouver que la suite
\[
  \mathfrak J^kP/\mathfrak J^{k+1}P\xrightarrow{u'}
  \mathfrak J^kQ/\mathfrak J^{k+1}Q\xrightarrow{v'}
  \mathfrak J^kR/\mathfrak J^{k+1}R
\]
(où $u'$ et $v'$ proviennent de $u$ et $v$ par restriction et passage aux
quotients) est exacte. Or, par hypothèse,
$\mathfrak J^k/\mathfrak J^{k+1}$ est un $(A/\mathfrak J)$-module plat, donc la
suite
\[
  (P/\mathfrak JP)\otimes_{A/\mathfrak J}
  (\mathfrak J^k/\mathfrak J^{k+1})\xrightarrow{u''}
  (Q/\mathfrak JQ)\otimes_{A/\mathfrak J}
  (\mathfrak J^k/\mathfrak J^{k+1})\xrightarrow{v''}
  (R/\mathfrak JR)\otimes_{A/\mathfrak J}
  (\mathfrak J^k/\mathfrak J^{k+1})
\]
est exacte. Mais
$(Q/\mathfrak JQ)\otimes_{A/\mathfrak J}
(\mathfrak J^k/\mathfrak J^{k+1})
=Q\otimes_A(\mathfrak J^k/\mathfrak J^{k+1})$ s'identifie à
$\mathfrak J^kQ/\mathfrak J^{k+1}Q$ en vertu de la platitude de $Q$ sur $A$,
et de même
$(R/\mathfrak JR)\otimes_{A/\mathfrak J}
(\mathfrak J^k/\mathfrak J^{k+1})$ s'identifie à
$\mathfrak J^kR/\mathfrak J^{k+1}R$. Enfin, l'image de $u'$ s'identifie à
celle de $u''$, et cela termine la démonstration.
\end{proof}

\begin{corollary}[12.3.4]
\label{IV.12.3.4-fr}
Soient $f:X\to S$ un morphisme plat localement de présentation finie,
$\mathcal F$, $\mathcal G$ deux $\mathcal O_X$-Modules de présentation finie et
$f$-plats. Soient $n$ un entier $\geq0$, $U$ l'ensemble des $x\in X$ tels que
\[
  \bigl(\mathscr{E}\!xt^n_{\mathcal O_{X_{f(x)}}}
  (\mathcal F_{f(x)},\mathcal G_{f(x)})\bigr)_x=0
\]
(resp.
$\bigl(\mathscr{T}\!or_n^{\mathcal O_{X_{f(x)}}}
(\mathcal F_{f(x)},\mathcal G_{f(x)})\bigr)_x=0$). Alors $U$ est ouvert, et
l'on a
\[
  \mathscr{E}\!xt^n_{\mathcal O_X}(\mathcal F,\mathcal G)|U=0
\]
(resp.
$\mathscr{T}\!or_n^{\mathcal O_X}(\mathcal F,\mathcal G)|U=0$).
\end{corollary}

\begin{proof}
On peut évidemment se borner au cas où $S$ et $X$ sont affines. Alors (12.3.2),
il existe une résolution gauche $\mathcal L_\bullet=(\mathcal L_i)$ de
$\mathcal F$ formée de $\mathcal O_X$-Modules \emph{libres de type fini}. Si
l'on pose $\mathcal M^\bullet=\mathscr{H}\!om(\mathcal L_\bullet,\mathcal G)$
(resp. $\mathcal N_\bullet=\mathcal L_\bullet\otimes_{\mathcal O_X}\mathcal G$),
on en déduit que chacun des $\mathcal M^i$ (resp. $\mathcal N_i$) est isomorphe
à un $\mathcal O_X$-Module de la forme $\mathcal G^m$, donc les $\mathcal M^i$
et $\mathcal N_i$ sont des $\mathcal O_X$-Modules de présentation finie et
$f$-plats. En outre, pour tout changement de base $S'\to S$, si l'on pose
$X'=X\times_SS'$, $\mathcal F'=\mathcal F\otimes_SS'$,
$\mathcal G'=\mathcal G\otimes_SS'$, $\mathcal L'_\bullet
=\mathcal L_\bullet\otimes_SS'$, $\mathcal L'_\bullet$ est encore une
résolution gauche de $\mathcal F'$ par des $\mathcal O_{X'}$-Modules libres de
type fini (2.1.10), et $\mathcal M'^\bullet=\mathcal M^\bullet\otimes_SS'$ est
égal à $\mathscr{H}\!om(\mathcal L'_\bullet,\mathcal G')$ (resp.
$\mathcal N'_\bullet=\mathcal N_\bullet\otimes_SS'$ est égal à
$\mathcal L'_\bullet\otimes_{\mathcal O_{X'}}\mathcal G'$), d'après ce qui
précède. En particulier, on a, pour tout $s\in S$,
\[
  \mathscr{E}\!xt^n_{\mathcal O_{X_s}}(\mathcal F_s,\mathcal G_s)
  =\mathcal H^n((\mathcal M^\bullet)_s)
\]
(resp.
$\mathscr{T}\!or_n^{\mathcal O_{X_s}}(\mathcal F_s,\mathcal G_s)
=\mathcal H_n((\mathcal N_\bullet)_s)$). Appliquant (12.3.3) aux complexes de
Modules $f$-plats $\mathcal M^\bullet$ et $\mathcal N_\bullet$, on en déduit
aussitôt le corollaire.
\end{proof}

% IV3_B26_P187_END_BEFORE_SECTION_13

\section{Morphismes équidimensionnels}
\label{section:IV.13-fr}

Ce paragraphe est consacré à l'étude de la variation de la \emph{dimension}
des fibres d'un morphisme localement de type fini $f:X\to Y$ (qui est déjà
intervenue à propos de la «~formule des dimensions~» dans 5.5 et 5.6). On
prouve d'abord (13.1.3) que la fonction
$x\mapsto\dim_x(f^{-1}(f(x)))$ est toujours \emph{semi-continue
supérieurement} dans $X$ (théorème de semi-continuité de Chevalley). On étudie
ensuite plus particulièrement les morphismes, dits «~équidimensionnels~», pour
lesquels cette fonction est \emph{localement constante}.

\oldpage[IV-3]{188}
Malheureusement la notion de morphisme équidimensionnel n'est pas stable par
changement de base~; c'est pourquoi dans de nombreuses questions il est plus
commode de travailler avec la notion de morphisme universellement ouvert, dont
l'étude fait l'objet des \S\S~14 et 15.

\subsection{Le théorème de semi-continuité de Chevalley}
\label{subsection:IV.13.1-fr}

\begin{lemma}[13.1.1]
\label{IV.13.1.1-fr}
Soient $Y$ un préschéma irréductible localement noethérien, $X$ un préschéma
irréductible, $f:X\to Y$ un morphisme dominant de type fini. Soit $\xi$
(resp. $\eta$) le point générique de $X$ (resp. $Y$) et soit
$e=\dim(f^{-1}(\eta))$. Alors, pour tout $x\in X$, toutes les composantes
irréductibles de $f^{-1}(f(x))$ sont de dimension $\geq e$.
\end{lemma}

\begin{proof}
La proposition est immédiate lorsque, pour tout $y\in f(X)$,
$\mathcal O_y$ est un \emph{anneau universellement caténaire}~: en effet, si
$x$ est point générique d'une composante irréductible $Z$ de $f^{-1}(y)$, il
résulte de (5.6.5), joint à (5.2.1), que l'on a
$e+\dim(\mathcal O_y)=\dim(Z)+\dim(\mathcal O_x)$~; mais en vertu de
($\mathbf 0$, 16.3.9), on a $\dim(\mathcal O_x)\leq\dim(\mathcal O_y)$,
d'où la conclusion dans ce cas.

Nous allons ramener le cas général à ce cas particulier. La question est
évidemment locale sur $Y$, et, compte tenu de (4.1.1.3), elle est aussi locale
sur $X$~; on peut donc se borner au cas où $Y=\operatorname{Spec}(A)$ et
$X=\operatorname{Spec}(B)$ sont affines et irréductibles, $B$ étant une
$A$-algèbre de type fini. En outre (1.5.4), on peut supposer $X$ et $Y$
réduits, donc $A$ et $B$ \emph{intègres} et, comme $f$ est dominant, $A$ est
alors un \emph{sous-anneau} de $B$. Considérons $A$ comme limite inductive de
ses sous-$\mathbf Z$-algèbres de type fini~; il résulte alors de (8.9.1) qu'il
existe une telle sous-algèbre $A_0$ et une $A_0$-algèbre de type fini $B_0$
telles que $B=B_0\otimes_{A_0}A$. Posons $Y_0=\operatorname{Spec}(A_0)$,
$X_0=\operatorname{Spec}(B_0)$, et soit $f_0:X_0\to Y_0$ le morphisme
correspondant à l'homomorphisme $A_0\to B_0$, de sorte que
$f=f_0\times 1_Y$. Il n'est pas évident \emph{a priori} que le préschéma $X_0$
soit intègre, mais nous allons voir qu'on peut se ramener à ce cas. Soit
$\eta_0$ le point générique de $Y_0$, de sorte que si $g$ est le morphisme
$Y\to Y_0$, on a $g(\eta)=\eta_0$~; par transitivité des fibres
(\textbf{I}, 3.6.4), on a
$f^{-1}(\eta)=f_0^{-1}(\eta_0)\otimes_{k(\eta_0)}k(\eta)$, et comme
$f^{-1}(\eta)$ est irréductible par hypothèse, il en est de même de
$f_0^{-1}(\eta_0)$ (4.4.1). Notre assertion résultera alors du lemme suivant~:

\begin{lemma}[13.1.2]
\label{IV.13.1.2-fr}
Soient $Y_0$, $Y$ deux préschémas intègres de points génériques $\eta_0$,
$\eta$, $g:Y\to Y_0$ un morphisme dominant, $f_0:X_0\to Y_0$ un morphisme
dominant tel que $f_0^{-1}(\eta_0)$ soit irréductible. Soit $X'_0$ l'unique
composante irréductible de $X_0$ rencontrant $f_0^{-1}(\eta_0)$
($\mathbf 0_{\mathrm I}$, 2.1.8), et désignons encore par $X'_0$ le
sous-préschéma fermé réduit de $X_0$ ayant $X'_0$ pour espace sous-jacent.
Supposons que le préschéma $X=X_0\times_{Y_0}Y$ soit intègre~; alors $X$ est
isomorphe à $X'_0\times_{Y_0}Y$.
\end{lemma}

En effet, si $j_0:X'_0\to X_0$ est l'injection canonique, qui est une immersion
fermée, $j=j_0\times 1_Y:X'_0\times_{Y_0}Y\to X_0\times_{Y_0}Y=X$ est une
immersion fermée. D'autre part, $X'_0$ contient la fibre
$f_0^{-1}(\eta_0)$, donc $X'_0\times_{Y_0}Y$ contient $f^{-1}(\eta)$~;
notons d'autre part que $f^{-1}(\eta)$ n'est pas vide (\textbf{I}, 3.4.7),
donc contient le point générique de $X$~; par suite l'image de $j$ est
nécessairement $X$ tout entier. Mais comme $X$ est intègre, le seul
sous-préschéma fermé de $X$ ayant $X$ pour espace sous-jacent est $X$
lui-même, donc $j$ est un isomorphisme.

Ce lemme étant établi, on peut donc supposer que $X_0$ est intègre~; pour tout
$y_0\in Y_0$, $\mathcal O_{y_0}$ est une $\mathbf Z$-algèbre essentiellement
de type fini, donc un anneau universellement

\oldpage[IV-3]{189}
caténaire (5.6.4)~; par suite, pour tout $y_0\in f_0(X_0)$, toutes les
composantes irréductibles de $f_0^{-1}(y)$ ont une dimension $\geq e$, car on
sait que $\dim(f_0^{-1}(\eta_0))=e$ par transitivité des fibres (4.1.4). Pour
tout $y\in f(X)$, on a alors $y_0=g(y)\in f_0(X_0)$ et par transitivité des
fibres et (4.2.7), on achève la démonstration.
\end{proof}

\begin{theorem}[13.1.3]
\label{IV.13.1.3-fr}
\textnormal{(Chevalley).} --- Soit $f:X\to Y$ un morphisme localement de type
fini. Pour tout entier $n$, l'ensemble $F_n(X)$ des $x\in X$ tels que
$\dim_x(f^{-1}(f(x)))\geq n$ est fermé~; autrement dit, la fonction
$x\mapsto\dim_x(f^{-1}(f(x)))$ est semi-continue supérieurement dans $X$.
\end{theorem}

\begin{proof}
I) Supposons d'abord que $f$ soit localement de présentation finie.

La question est évidemment locale sur $X$ et sur $Y$, et on peut donc supposer
$Y=\operatorname{Spec}(A)$, $X=\operatorname{Spec}(B)$ affines, $B$ étant une
$A$-algèbre de présentation finie. On sait alors (8.9.1) qu'il y a un
sous-anneau noethérien $A_0$ de $A$, et une $A_0$-algèbre de type fini $B_0$
tels que si l'on pose $Y_0=\operatorname{Spec}(A_0)$,
$X_0=\operatorname{Spec}(B_0)$, on ait

\[
  X=X_0\times_{Y_0}Y,\qquad f=f_0\times 1_Y,
\]

où $f_0:X_0\to Y_0$ correspond à l'homomorphisme $A_0\to B_0$. Soient
$g:Y\to Y_0$ le morphisme correspondant à l'injection canonique $A_0\to A$,
$y$ un point de $Y$, $y_0=g(y)$~; on sait que
$f^{-1}(y)=f_0^{-1}(y_0)\otimes_{k(y_0)}k(y)$, et il résulte de (4.2.7) que
si $p$ est la projection canonique $f^{-1}(y)\to f_0^{-1}(y_0)$, les
composantes irréductibles de $f^{-1}(y)$ sont les composantes irréductibles des
ensembles $p^{-1}(Z_0)$, où $Z_0$ parcourt l'ensemble des composantes
irréductibles de $f_0^{-1}(y_0)$, et chacune des composantes irréductibles de
$p^{-1}(Z_0)$ domine $Z_0$ et a une dimension égale à $\dim(Z_0)$. Tenant
compte de ($\mathbf 0$, 14.1.5), on voit donc que si $g':X\to X_0$ est la
projection canonique, $x$ un point de $X$ et $x_0=g'(x)$, on a
$\dim_x(f^{-1}(f(x)))=\dim_{x_0}(f_0^{-1}(f_0(x_0)))$, d'où
$F_n(X)=g'^{-1}(F_n(X_0))$, et on est par suite ramené à démontrer le théorème
lorsque $Y$ est \emph{noethérien}, ce que nous supposerons désormais.

On peut évidemment supposer $X$ et $Y$ réduits (1.5.4). En considérant
l'ensemble des parties fermées $Y'$ de $Y$ telles que le théorème soit vrai
pour le sous-préschéma fermé de $Y$ ayant $Y'$ pour espace sous-jacent et pour
$X'=f^{-1}(Y')$, on peut raisonner par récurrence noethérienne
($\mathbf 0_{\mathrm I}$, 2.2.2), et supposer que le théorème est vrai pour
toute partie fermée $Y'\neq Y$ de $Y$. Si $X_i$ ($1\leq i\leq m$) sont les
sous-préschémas fermés réduits de $X$ ayant pour espaces sous-jacents les
composantes irréductibles de $X$, on a $F_n(X)=\bigcup_iF_n(X_i)$ en vertu de
($\mathbf 0$, 14.1.5), et l'on peut donc se borner à démontrer le théorème pour
chacun des $X_i$, autrement dit, on peut supposer $X$ irréductible. Si $Z$ est
le sous-préschéma fermé de $Y$ ayant $\overline{f(X)}$ pour espace sous-jacent,
$f$ se factorise en $X\xrightarrow{g}Z\xrightarrow{j}Y$, où $j$ est
l'injection canonique (\textbf{I}, 5.2.2), et $g$ est de type fini (1.5.4),
donc il suffit de démontrer le théorème pour $Z$ et $g$~; en vertu de
l'hypothèse de récurrence, on est donc ramené à considérer seulement le cas où
$Z=Y$, autrement dit où $Y$ est irréductible et $f$ dominant. Soit alors
$\eta$ le point générique de $Y$ et posons $e=\dim(f^{-1}(\eta))$~; il résulte
de (13.1.1) que pour $n\leq e$ on a $F_n(X)=X$, et par suite on peut se borner
au cas où $n>e$. Mais alors (9.5.6), il y a un voisinage ouvert $U$ de $\eta$
dans $Y$ tel que $F_n(X)\subset f^{-1}(Y-U)$~; comme $Y-U\neq Y$, l'hypothèse
de récurrence entraîne que $F_n(X)$ est fermé.

II) Passons maintenant au cas général, en supposant toujours
$S=\operatorname{Spec}(A)$ et $X=\operatorname{Spec}(B)$ affines, $B$ étant
une $A$-algèbre de type fini, donc de la forme

\oldpage[IV-3]{190}
$B=A[T_1,\ldots,T_n]/\mathfrak J$. Soit $(\mathfrak J_\lambda)$ la famille des
idéaux de type fini de $A[T_1,\ldots,T_n]$ contenus dans $\mathfrak J$, de
sorte que $\mathfrak J$ est la réunion filtrante des $\mathfrak J_\lambda$~;
si $X$ et les
$X_\lambda=\operatorname{Spec}(A[T_1,\ldots,T_n]/\mathfrak J_\lambda)$ sont
considérés comme des sous-préschémas fermés de
$Z=\operatorname{Spec}(A[T_1,\ldots,T_n])$, on a donc, pour les espaces
sous-jacents, $X=\bigcap_\lambda X_\lambda$. Si $f_\lambda:X_\lambda\to S$ est
le morphisme structural, on en déduit que
$f^{-1}(s)=\bigcap_\lambda f_\lambda^{-1}(s)$ pour tout $s\in S$, et comme les
ensembles $f_\lambda^{-1}(s)$ sont fermés dans l'espace noethérien $Z_s$, il
existe un $\lambda$ (dépendant de $s$) tel que
$f^{-1}(s)=f_\lambda^{-1}(s)$. Si alors, pour tout $x\in X$, on pose
$d(x)=\dim_x(f^{-1}(f(x)))$, $d_\lambda(x)=\dim_x(f_\lambda^{-1}(f_\lambda(x)))$,
ce qui précède prouve que l'on a $d(x)=\inf_\lambda d_\lambda(x)$~; les
fonctions $d_\lambda$ étant semi-continues supérieurement d'après la première
partie de la démonstration, il en est de même de $d$. C.Q.F.D.
\end{proof}

\begin{corollary}[13.1.4]
\label{IV.13.1.4-fr}
Sous les hypothèses de (13.1.3), l'ensemble des $x\in X$ tels que $x$ soit
isolé dans $f^{-1}(f(x))$ est ouvert dans $X$.
\end{corollary}

\begin{proof}
En effet, c'est le complémentaire de $F_1(X)$ ($\mathbf 0$, 14.1.10).
\end{proof}

On notera qu'on retrouve ainsi, sous des hypothèses plus générales, la
conséquence (\textbf{III}, 4.4.10) du «~Main theorem~» de Zariski.

\begin{corollary}[13.1.5]
\label{IV.13.1.5-fr}
Sous les hypothèses de (13.1.3), supposons en outre que $f$ soit un morphisme
fermé. Alors, pour tout entier $n$, l'ensemble des $y\in Y$ tels que
$\dim(f^{-1}(y))\geq n$ est fermé, autrement dit, l'application
$y\mapsto\dim(f^{-1}(y))$ est semi-continue supérieurement~; en particulier,
si $y$ est spécialisation de $y'$, on a
$\dim(f^{-1}(y))\geq\dim(f^{-1}(y'))$.
\end{corollary}

\begin{proof}
En effet, dire que $\dim(f^{-1}(y))\geq n$ signifie que $y\in f(F_n(X))$
($\mathbf 0$, 14.1.6).
\end{proof}

\begin{corollary}[13.1.6]
\label{IV.13.1.6-fr}
Soient $X$, $Y$ deux préschémas irréductibles, $f:X\to Y$ un morphisme
dominant localement de type fini. Soit $\xi$ (resp. $\eta$) le point générique
de $X$ (resp. $Y$) et soit $e=\dim(f^{-1}(\eta))$. Alors, pour tout $x\in X$,
on a $\dim_x(f^{-1}(f(x)))\geq e$.
\end{corollary}

\begin{proof}
En effet, l'ensemble des $x$ tels que $\dim_x(f^{-1}(f(x)))<e$ est ouvert en
vertu de (13.1.3), et comme il ne peut contenir $\xi$, il est vide.
\end{proof}

Notons enfin le résultat plus facile suivant~:

\begin{proposition}[13.1.7]
\label{IV.13.1.7-fr}
Soient $Y$ un préschéma quasi-compact, $f:X\to Y$ un morphisme de type fini.
Il existe un entier $n$ tel que, pour tout $y\in Y$, on ait
$\dim(f^{-1}(y))\leq n$.
\end{proposition}

\begin{proof}
Comme il y a un recouvrement ouvert affine fini $(V_i)$ de $Y$ tel que chaque
$f^{-1}(V_i)$ soit réunion finie d'ouverts affines, on est aussitôt ramené au
cas où $Y=\operatorname{Spec}(A)$ et $X=\operatorname{Spec}(B)$ sont affines,
$B$ étant une $A$-algèbre de type fini. Si $B$ admet un système de $n$
générateurs, alors, pour tout $y\in Y$, $B\otimes_Ak(y)$ est une
$k(y)$-algèbre admettant $n$ générateurs, donc $\dim(f^{-1}(y))\leq n$ en
vertu de (4.1.1).
\end{proof}

\subsection{Morphismes équidimensionnels : cas des morphismes dominants de préschémas irréductibles}
\label{subsection:IV.13.2-fr}

\begin{env}[13.2.1]
\label{IV.13.2.1-fr}
Soient $Y$ un préschéma irréductible, $X$ un préschéma irréductible,
$f:X\to Y$ un morphisme dominant localement de type fini~; soit $\eta$ le
point générique de $Y$. On sait (13.1.6) que pour tout $x\in X$, on a
$\dim_x(f^{-1}(f(x)))\geq\dim(f^{-1}(\eta))$.
\end{env}

\oldpage[IV-3]{191}
\begin{definition}[13.2.2]
\label{IV.13.2.2-fr}
Sous les hypothèses de (13.2.1), on dit que $f$ est équidimensionnel au point
$x$ (ou que $X$ est équidimensionnel sur $Y$ au point $x$) si
\[
  \dim_x(f^{-1}(f(x)))=\dim(f^{-1}(\eta)).
\]
On dit que $f$ est équidimensionnel (ou que $X$ est équidimensionnel sur $Y$)
si $f$ est équidimensionnel en tout point $x\in X$.
\end{definition}

Il résulte du théorème de Chevalley (13.1.3) que l'ensemble des $x\in X$ où
$f$ est équidimensionnel est \emph{ouvert non vide}. En outre, si $f$ est
équidimensionnel au point $x$, toutes les composantes irréductibles de
$f^{-1}(f(x))$ qui contiennent $x$ ont \emph{même dimension}, car chacune
d'elles a une dimension qui est $\geq\dim(f^{-1}(\eta))$ (13.1.6) et
$\leq\dim_x(f^{-1}(f(x)))$ en vertu de ($\mathbf 0$, 14.1.5).

\begin{proposition}[13.2.3]
\label{IV.13.2.3-fr}
Soient $Y$ un préschéma irréductible localement noethérien, $X$ un préschéma
irréductible, $f:X\to Y$ un morphisme dominant localement de type fini~;
soient $\eta$ le point générique de $Y$, $x$ un point de $X$, $y=f(x)$, et
supposons que l'on ait
\[
\tag{13.2.3.1}
\label{IV.13.2.3.1-fr}
  \dim(\mathcal O_x)=\dim(\mathcal O_y)
  +\dim(\mathcal O_x\otimes_{\mathcal O_y}k(y)).
\]
Alors $f$ est équidimensionnel au point $x$. La réciproque est vraie si les
deux membres de l'inégalité (5.6.5.2) sont égaux, en particulier si
$\mathcal O_y$ est universellement caténaire.
\end{proposition}

\begin{proof}
Cela résulte aussitôt de (5.6.5.2) et de l'inégalité $\delta(x)\geq0$
(13.1.1).
\end{proof}

\begin{env}[13.2.4]
\label{IV.13.2.4-fr}
Soient maintenant, de façon générale, $Y$ un préschéma localement noethérien,
$f:X\to Y$ un morphisme de type fini, $x$ un point de $X$, $X'$ une partie
fermée irréductible de $X$ contenant $x$, $y=f(x)$,
$Y'=\overline{f(X')}$, $\eta'$ le point générique de $Y'$. Désignons encore
par $X'$ et $Y'$ les sous-préschémas fermés réduits de $X$ et $Y$
respectivement ayant $X'$ et $Y'$ pour espaces sous-jacents~; la restriction
$X'\to Y$ de $f$ se factorise donc en
$X'\xrightarrow{f'}Y'\xrightarrow{j}Y$, où $j$ est l'injection canonique
(\textbf{I}, 5.2.2), et $f'$ est de type fini (1.5.4). Posons alors pour
abréger
\[
\tag{13.2.4.1}
\label{IV.13.2.4.1-fr}
  A=\mathcal O_{Y,y},\quad B=\mathcal O_{X,x},\quad
  A'=\mathcal O_{Y',y},\quad B'=\mathcal O_{X',x}.
\]
La formule (5.6.5.2) appliquée au morphisme dominant $f'$ et aux préschémas
irréductibles $X'$, $Y'$, donne
\[
\tag{13.2.4.2}
\label{IV.13.2.4.2-fr}
  \dim(B')\leq\dim(A')+\dim(B'\otimes_{A'}k(y))
  -(\dim_x(f'^{-1}(y))-\dim(f'^{-1}(\eta'))).
\]
D'autre part, l'anneau local $A'$ (resp. $B'$, $B'\otimes_{A'}k(y)$) est un
anneau quotient de $A$ (resp. $B$, $B\otimes_Ak(y)$), donc
($\mathbf 0$, 16.1.2.1) on a
\[
\tag{13.2.4.3}
\label{IV.13.2.4.3-fr}
  \dim(A')\leq\dim(A),\quad \dim(B')\leq\dim(B),\quad
  \dim(B'\otimes_{A'}k(y))\leq\dim(B\otimes_Ak(y)).
\]
On déduit donc en premier lieu de (13.2.4.3) et ($\mathbf 0$, 16.3.9)
\[
\tag{13.2.4.4}
\label{IV.13.2.4.4-fr}
  \dim(B')\leq\dim(A')+\dim(B'\otimes_{A'}k(y))
  \leq\dim(A)+\dim(B\otimes_Ak(y)).
\]
En outre, en vertu de ($\mathbf 0$, 16.3.9), on a aussi les inégalités
\[
\tag{13.2.4.5}
\label{IV.13.2.4.5-fr}
  \dim(B')\leq\dim(B)\leq\dim(A)+\dim(B\otimes_Ak(y)).
\]
\end{env}

\oldpage[IV-3]{192}
La comparaison de ces inégalités montre donc que~:

\begin{lemma}[13.2.5]
\label{IV.13.2.5-fr}
Avec les notations de (13.2.4), les conditions suivantes sont équivalentes~:
\begin{enumerate}[label=\emph{\alph*)}]
  \item $\dim(B')=\dim(A)+\dim(B\otimes_Ak(y))$.
  \item On a simultanément les relations suivantes~:
  \begin{enumerate}[label=(\roman*)]
    \item $\dim(A')=\dim(A)$.
    \item $\dim(B'\otimes_{A'}k(y))=\dim(B\otimes_Ak(y))$, autrement dit
    \[
      \dim_x(X'\cap f^{-1}(y))=\dim_x(f^{-1}(y)).
    \]
    \item $\dim_x(f'^{-1}(y))=\dim(f'^{-1}(\eta'))$, autrement dit
    $f':X'\to Y'$ est équidimensionnel (13.2.2) au point $x$.
    \item On a l'égalité
    \[
      \dim(B')=\dim(A')+\dim(B'\otimes_{A'}k(y))
      -(\dim_x(f'^{-1}(y))-\dim(f'^{-1}(\eta')))
    \]
  \end{enumerate}
  (relation qui est toujours vérifiée lorsque $A=\mathcal O_{Y,y}$ est un
  anneau universellement caténaire, en vertu de (5.6.5)).
  \item On a simultanément les relations suivantes~:
  \begin{enumerate}[label=(\roman*)]
    \item $\dim(B')=\dim(B)$.
    \item $\dim(B)=\dim(A)+\dim(B\otimes_Ak(y))$.
  \end{enumerate}
\end{enumerate}
\end{lemma}

\begin{env}[13.2.6]
\label{IV.13.2.6-fr}
Rappelons maintenant que les composantes irréductibles $X_i$ de $X$ contenant
$x$ sont en nombre fini et que l'on a ((5.1.2.1) et ($\mathbf 0$, 14.2.1.1))
\[
\tag{13.2.6.1}
\label{IV.13.2.6.1-fr}
  \dim(\mathcal O_{X,x})=\sup_i\dim(\mathcal O_{X_i,x}).
\]
L'équivalence des conditions $b)$ et $c)$ dans (13.2.5) implique par suite,
compte tenu de ($\mathbf 0$, 16.3.9) et ($\mathbf 0$, 14.2.1)~:
\end{env}

\begin{proposition}[13.2.7]
\label{IV.13.2.7-fr}
Soient $Y$ un préschéma localement noethérien, $f:X\to Y$ un morphisme de type
fini, $x$ un point de $X$, $f(x)=y$~; on a
\[
\tag{13.2.7.1}
\label{IV.13.2.7.1-fr}
  \dim(\mathcal O_{X,x})\leq\dim(\mathcal O_{Y,y})
  +\dim(\mathcal O_{X,x}\otimes_{\mathcal O_{Y,y}}k(y)).
\]
Pour que les deux membres de (13.2.7.1) soient égaux, il faut et il suffit
qu'il existe une partie fermée irréductible $X'$ de $X$ contenant $x$ et
vérifiant simultanément les conditions suivantes~:
\begin{enumerate}[label=(\roman*)]
  \item Si $Y'=\overline{f(X')}$, on a
  $\dim(\mathcal O_{Y',y})=\dim(\mathcal O_{Y,y})$.
  \item $\dim_x(X'\cap f^{-1}(y))=\dim_x(f^{-1}(y))$ (autrement dit, $X'$
  contient une des composantes irréductibles de $f^{-1}(y)$ contenant $x$, de
  dimension maxima parmi toutes ces composantes). Il revient au même de dire
  que
  $\dim(\mathcal O_{X',x}\otimes_{\mathcal O_{Y',y}}k(y))
  =\dim(\mathcal O_{X,x}\otimes_{\mathcal O_{Y,y}}k(y))$.
  \item $X'$ est équidimensionnel sur $Y'$ au point $x$, autrement dit, on a
  \[
    \dim_x(X'\cap f^{-1}(y))=\dim(X'\cap f^{-1}(\eta')),
  \]
  où $\eta'$ est le point générique de $Y'$ (et par suite, toutes les
  composantes irréductibles de $X'\cap f^{-1}(y)$ contenant $x$ ont même
  dimension (13.2.2)).

  \oldpage[IV-3]{193}
  \item On a l'égalité
  \[
    \dim(\mathcal O_{X',x})=\dim(\mathcal O_{Y',y})
    +\dim(\mathcal O_{X',x}\otimes_{\mathcal O_{Y',y}}k(y))
  \]
  (condition toujours impliquée par (iii) lorsque $\mathcal O_{Y,y}$ est un
  anneau universellement caténaire).
\end{enumerate}
De plus, $X'$ est alors une composante irréductible de $X$ et $Y'$ une
composante irréductible de $Y$.
\end{proposition}

Dans l'énoncé, les anneaux locaux $\mathcal O_{X',x}$ et
$\mathcal O_{Y',y}$ sont relatifs aux sous-préschémas fermés réduits de $X$,
$Y$ respectivement ayant $X'$ et $Y'$ pour espaces sous-jacents respectifs.

En outre, cela prouve que

\begin{corollary}[13.2.8]
\label{IV.13.2.8-fr}
Si les deux membres de (13.2.7.1) sont égaux, les parties fermées irréductibles
$X'$ de $X$ contenant $x$ et vérifiant les conditions (i) à (iv) de (13.2.7)
sont exactement celles pour lesquelles on a
$\dim(\mathcal O_{X',x})=\dim(\mathcal O_{X,x})$ (ce qui implique
nécessairement que $X'$ est une composante irréductible de $X$).
\end{corollary}

La proposition (13.2.7) entraîne~:

\begin{corollary}[13.2.9]
\label{IV.13.2.9-fr}
Soient $Y$ un préschéma localement noethérien, $f:X\to Y$ un morphisme de type
fini, $x$ un point de $X$, $f(x)=y$. Supposons que $\mathcal O_{Y,y}$ soit un
anneau universellement caténaire. Alors les conditions suivantes sont
équivalentes~:
\begin{enumerate}[label=\emph{\alph*)}]
  \item Les deux membres de (13.2.7.1) sont égaux.
  \item Il existe une composante irréductible $Z$ de $f^{-1}(y)$ contenant
  $x$, de dimension $\dim_x(f^{-1}(y))$, telle que, pour tout $x'$ dans un
  voisinage de $x$ dans $Z$, on ait
  \[
  \tag{13.2.9.1}
  \label{IV.13.2.9.1-fr}
    \dim(\mathcal O_{X,x'})=\dim(\mathcal O_{Y,y})
    +\dim(\mathcal O_{X,x'}\otimes_{\mathcal O_{Y,y}}k(y)).
  \]
  \item Il existe une composante irréductible $Z$ de $f^{-1}(y)$ contenant
  $x$, de dimension $\dim_x(f^{-1}(y))$, telle que pour le point générique $z$
  de $Z$, on ait
  \[
  \tag{13.2.9.2}
  \label{IV.13.2.9.2-fr}
    \dim(\mathcal O_{X,z})=\dim(\mathcal O_{Y,y}).
  \]
\end{enumerate}
\end{corollary}

Montrons que $a)$ entraîne $b)$. Soit $\dim_x(f^{-1}(y))=n$~; en vertu de
(13.2.7), il existe une composante irréductible $X'$ de $X$ vérifiant les
conditions (i) à (iv) de (13.2.7)~; soit $Z$ une composante irréductible de
dimension $n$ de $f^{-1}(y)$, contenant $x$ et contenue dans $X'$. Comme
$f^{-1}(y)$ est localement noethérien, il existe un voisinage ouvert $U$ de
$x$ dans $Z$ tel que $U$ ne rencontre aucune composante irréductible de
$f^{-1}(y)$ autre que celles qui contiennent $x$, donc (4.1.1.3)
$\dim_{x'}(f^{-1}(y))=n$ pour tout $x'\in U$~; il est clair alors que les
conditions (i) à (iii) de (13.2.7) sont satisfaites quand on y remplace $x$
par un point quelconque $x'\in U$, et il en est de même de la condition (iv)
puisque $\mathcal O_{Y,y}$ est universellement caténaire~; d'où la conclusion
par (13.2.7). La condition $b)$ entraîne trivialement $c)$ en vertu de
(5.1.2). Enfin, si $c)$ est vérifiée et si $X''$ est une composante
irréductible de $X$ contenant $Z$ et telle que les conditions (i) à (iv) de
(13.2.7) soient vérifiées quand on y remplace $X'$ par $X''$ et $x$ par $z$,
il est clair que ces conditions sont aussi vérifiées pour $X''$ et $x$ puisque
$\mathcal O_{Y,y}$ est universellement caténaire, donc $c)$ implique $a)$.

\begin{proposition}[13.2.10]
\label{IV.13.2.10-fr}
Soient $Y$ un préschéma irréductible localement noethérien, $\eta$ son point
générique, $f:X\to Y$ un morphisme de type fini, $y$ un point de $f(X)$.
Soient $Z_i$ les

\oldpage[IV-3]{194}
composantes irréductibles de $f^{-1}(y)$, $z_i$ le point générique de $Z_i$,
et considérons les conditions suivantes~:
\begin{enumerate}[label=\emph{\alph*)}]
  \item Pour tout $x\in f^{-1}(y)$, on a la relation
  \[
  \tag{13.2.10.1}
  \label{IV.13.2.10.1-fr}
    \dim(\mathcal O_{X,x})=\dim(\mathcal O_{Y,y})
    +\dim(\mathcal O_{X,x}\otimes_{\mathcal O_{Y,y}}k(y)).
  \]
  \item Pour tout $i$, on a
  \[
  \tag{13.2.10.2}
  \label{IV.13.2.10.2-fr}
    \dim(\mathcal O_{X,z_i})=\dim(\mathcal O_{Y,y}).
  \]
  \item Pour tout $i$, il existe une composante irréductible $X_i$ de $X$
  contenant $Z_i$ et telle que
  $\dim(X_i\cap f^{-1}(\eta))=\dim(Z_i)$ (autrement dit, telle que le
  \emph{sous-préschéma fermé réduit} $X_i$ de $X$ soit équidimensionnel sur $Y$
  au point $z_i$).
\end{enumerate}
Alors $a)$ entraîne $b)$ et $b)$ entraîne $c)$~; en outre, si
$\mathcal O_{Y,y}$ est universellement caténaire, les trois conditions $a)$,
$b)$, $c)$ sont équivalentes.
\end{proposition}

L'anneau $\mathcal O_{X,z_i}\otimes_{\mathcal O_{Y,y}}k(y)$ étant de dimension
$0$ (5.1.2), $a)$ entraîne évidemment $b)$~; $b)$ entraîne $c)$ en vertu de
(13.2.7) appliqué au point $z_i$. Inversement, supposons que
$\mathcal O_{Y,y}$ soit universellement caténaire~; comme la condition $c)$
implique que $f(X_i)$ est dense dans $Y$, il résulte de $c)$ que les conditions
(i) à (iv) de (13.2.7) sont vérifiées en remplaçant $X'$ par $X_i$ et $x$ par
$z_i$, donc $c)$ implique $b)$~; enfin $b)$ implique $a)$ en vertu de
(13.2.9).

\begin{corollary}[13.2.11]
\label{IV.13.2.11-fr}
Les notations étant celles de (13.2.10), supposons que $\mathcal O_{Y,y}$ soit
universellement caténaire. Pour tout $y'\in Y$, soit $E(y')$ l'ensemble des
dimensions des composantes irréductibles de $f^{-1}(y')$ et posons
$d(y')=\dim(f^{-1}(y'))=\sup(E(y'))$. Alors, si les conditions équivalentes
$a)$, $b)$, $c)$ de (13.2.10) sont vérifiées, on a
$E(y)\subset E(\eta)$, d'où $d(y)\leq d(\eta)$.
\end{corollary}

En effet, avec les notations de (13.2.10), $X_i\cap f^{-1}(\eta)$ n'est pas
vide et est par suite une composante irréductible de $f^{-1}(\eta)$
($\mathbf 0_{\mathrm I}$, 2.1.8).

\begin{remarks}[13.2.12]
\label{IV.13.2.12-fr}
\begin{enumerate}[label=(\roman*)]
  \item Rappelons (6.1.2) que la relation (13.2.10.1) est toujours vérifiée
  lorsque le morphisme $f$ est \emph{plat} au point $x$.
  \item Avec les notations de (13.2.10), supposons que $\mathcal O_{Y,y}$ soit
  universellement caténaire et en outre que le morphisme $f$ soit \emph{propre}~;
  alors, si les conditions équivalentes $a)$, $b)$, $c)$ de (13.2.10) sont
  satisfaites, on a même $d(y)=d(\eta)$, car il résulte de (13.1.5) que l'on a
  $d(y)\geq d(\eta)$.
  \item Le morphisme de (12.2.3, (ii)) est \emph{propre} et \emph{plat} et tous
  les anneaux locaux de $Y$ sont universellement caténaires~; en outre, les deux
  composantes irréductibles $X_1$, $X_2$ de $X$ sont \emph{équidimensionnelles}
  sur $Y$ en tout point~; mais $E(y)$ a \emph{deux} éléments pour tout
  $y\neq y_0$, alors que $E(y_0)$ est réduit à un \emph{seul} élément, donc
  $E(y)$ n'est \emph{pas constant} dans $Y$.
\end{enumerate}
\end{remarks}

\subsection{Morphismes équidimensionnels : cas général}
\label{subsection:IV.13.3-fr}

\begin{proposition}[13.3.1]
\label{IV.13.3.1-fr}
Soient $Y$ un préschéma, $f:X\to Y$ un morphisme localement de type fini,
$x$ un point de $X$, $y=f(x)$. Désignons par $Y_\alpha$ les composantes
irréductibles de $Y$ contenant $y$. Alors les conditions suivantes sont
équivalentes~:
\begin{enumerate}[label=\emph{\alph*)}]
\oldpage[IV-3]{195}
  \item Il existe un entier $e$ et un voisinage ouvert $U$ de $x$ tels que
  l'image par $f$ de toute composante irréductible de $U$ soit dense dans un
  $Y_\alpha$, et que, pour tout $x'\in U$, l'espace
  $U\cap f^{-1}(f(x'))$ soit équidimensionnel et de dimension $e$.
  \item[\emph{a$'$)}] Il existe un entier $e$ et un voisinage ouvert $U$ de
  $x$ tels que l'image par $f$ de toute composante irréductible de $U$ soit
  dense dans un $Y_\alpha$ et que, si on désigne par $y_\alpha$ le point
  générique de $Y_\alpha$, toute composante irréductible des espaces
  $U\cap f^{-1}(y)$, $U\cap f^{-1}(y_\alpha)$ soit de dimension $e$.
  \item[\emph{a$''$)}] Il existe un entier $e$ et un voisinage ouvert $U$ de
  $x$ tels que, pour chacune des composantes irréductibles $U_\lambda$ de
  $U$, $f(U_\lambda)$ soit dense dans un $Y_\alpha$ et que, pour tout
  $x'\in U_\lambda$, les composantes irréductibles de
  $U_\lambda\cap f^{-1}(f(x'))$ soient toutes de dimension $e$.
  \item Il existe un entier $e$, un voisinage ouvert $U$ de $x$ et un
  $Y$-morphisme quasi-fini
  \[
    g:U\longrightarrow Y\otimes_{\mathbf Z}\mathbf Z[T_1,\ldots,T_e]
  \]
  (préschéma que nous noterons aussi $Y[T_1,\ldots,T_e]$ pour abréger) tels
  que l'image par $g$ de toute composante irréductible de $U$ soit dense dans
  une composante irréductible de $Y[T_1,\ldots,T_e]$.
\end{enumerate}
\end{proposition}

\begin{proof}
Il est immédiat que $a''$) entraîne $a$), car pour tout $x\in U$, les
composantes irréductibles de $U\cap f^{-1}(f(x))$ sont chacune composante
irréductible d'un des $U_\lambda\cap f^{-1}(f(x))$, d'où la conclusion par
($\mathbf 0$, 14.1.4). La condition $a$) entraîne trivialement $a'$).
Montrons ensuite que $a'$) entraîne $a''$)~; on peut se borner au cas où $f$
est de type fini et $X$ et $Y$ réduits (1.5.4). Soit $U_\lambda$ une
composante irréductible de $U$, et supposons que $f(U_\lambda)$ soit dense
dans $Y_\alpha$~; alors la restriction de $f$ à $U_\lambda$ se factorise en
$U_\lambda\xrightarrow{f_\alpha}Y_\alpha\to Y$, où $f_\alpha$ est de type
fini et dominant (\textbf{I}, 5.2.2). Soit $x_\lambda$ le point générique de
$U_\lambda$~; en vertu de ($\mathbf 0_{\mathrm I}$, 2.1.8),
$U_\lambda\cap f_\alpha^{-1}(y_\alpha)$ est l'unique composante irréductible
de $U\cap f^{-1}(y_\alpha)$ contenant $x_\lambda$ et est par hypothèse de
dimension $e$, égale aux dimensions de toutes les composantes irréductibles
de $U_\lambda\cap f^{-1}(y)$, en vertu de l'hypothèse et de (13.1.1). Mais en
vertu du théorème de Chevalley (13.1.3) et de (13.1.1), l'ensemble $V_\lambda$
des $x'\in U_\lambda$ tels que
$\dim_{x'}(f_\lambda^{-1}(f_\lambda(x')))=e$ est ouvert et contient $x$, et il
suffit de prendre la réunion des $V_\lambda$ pour obtenir un ensemble ouvert
satisfaisant aux conditions de $a''$).

Prouvons maintenant que $a$) entraîne $b$)~; on peut se limiter au cas où
$Y=\operatorname{Spec}(A)$ et $X=\operatorname{Spec}(B)$ sont affines et où
$U=X$. Prouvons d'abord le lemme suivant~:

\begin{lemma}[13.3.1.1]
\label{IV.13.3.1.1-fr}
Soient $Y=\operatorname{Spec}(A)$, $X=\operatorname{Spec}(B)$ deux schémas
affines, $f:X\to Y$ un morphisme de type fini, $y$ un point de $f(X)$, et
posons $e=\dim(f^{-1}(y))$. Alors il existe un $Y$-morphisme
$g:X\to Y[T_1,\ldots,T_e]$ tel que (si l'on pose
$X_y=X\times_Y\operatorname{Spec}(k(y))$), le morphisme
$g_y:X_y\to\operatorname{Spec}(k(y)[T_1,\ldots,T_e])$ soit fini. En outre,
pour un tel morphisme $g$, $g_y$ est nécessairement surjectif et il existe un
voisinage ouvert $U$ de $f^{-1}(y)$ dans $X$ tel que $g|U$ soit quasi-fini.
\end{lemma}

Soit $\mathfrak p=\mathfrak j_y$~; l'anneau $B\otimes_Ak(\mathfrak p)$ est une
$k(\mathfrak p)$-algèbre de type fini, donc le lemme de normalisation
(Bourbaki, \emph{Alg. comm.}, chap.~V, \S~3, n\textsuperscript{o}~1, th.~1)
prouve qu'il y a dans $B\otimes_Ak(\mathfrak p)$ une suite finie
$(t_i)_{1\leq i\leq r}$ d'éléments algébriquement indépendants sur
$k(\mathfrak p)$ et tels que si l'on pose
$C'=k(\mathfrak p)[t_1,\ldots,t_r]$, $B\otimes_Ak(\mathfrak p)$ soit une
$C'$-algèbre \emph{finie}~; on a donc
$\dim(B\otimes_Ak(\mathfrak p))=\dim(C')$ ($\mathbf 0$, 16.1.5) et comme
$\dim(C')=r$ (5.2.1), on a $r=e$. Comme
$B\otimes_Ak(\mathfrak p)=(B/\mathfrak pB)\otimes_{A/\mathfrak p}k(\mathfrak p)$,
on peut, en multipliant les $t_i$ par un élément $\neq0$

\oldpage[IV-3]{196}
convenable de $A/\mathfrak p$, supposer que chaque $t_i$ est l'image canonique
dans $B\otimes_Ak(\mathfrak p)$ d'un élément $s_i\in B$. Soit alors
$u:A[T_1,\ldots,T_e]\to B$ l'homomorphisme tel que $u(T_i)=s_i$ pour tout
$i$, et soit $g:X\to Y[T_1,\ldots,T_e]$ le morphisme correspondant. Il est
clair qu'en raison du choix des $s_i$, $g_y$ est un morphisme fini. Pour tout
morphisme $g:X\to Y[T_1,\ldots,T_e]$, tel que $g_y$ soit fini, il résulte de
(5.4.2) et (4.1.2.1) que $g_y$ est nécessairement surjectif. D'autre part, en
vertu de (13.1.4), l'ensemble $U$ des $x\in X$ qui sont isolés dans leur fibre
$g^{-1}(g(x))$ est ouvert dans $X$ et contient $f^{-1}(y)$, et par définition
la restriction $g|U$ est un morphisme quasi-fini.

Ce lemme étant établi, pour prouver que $a$) implique $b$), il reste à voir
que si $x_\lambda$ est un point maximal de $U$, $g(x_\lambda)=z_\lambda$ est
un point maximal de $Z=Y[T_1,\ldots,T_e]$. En vertu de l'hypothèse $a$), on
peut (en restreignant au besoin $U$), supposer que $f(x_\lambda)$ est l'un des
points génériques $y_\alpha$ des composantes irréductibles de $Y$ contenant
$y$~; si $p:Z\to Y$ est le morphisme structural, on a donc
$p(z_\lambda)=y_\alpha$, et l'on déduit par suite de $g$ un
$k(y_\alpha)$-morphisme quasi-fini
\[
  h:U\cap f^{-1}(y_\alpha)\longrightarrow p^{-1}(y_\alpha).
\]
Or $p^{-1}(y_\alpha)=\operatorname{Spec}(k(y_\alpha)[T_1,\ldots,T_e])$ est
intègre et de dimension $e$~; si $z_\lambda$ n'était pas point maximal de
$Z$, il ne serait pas point maximal de $p^{-1}(y_\alpha)$
($\mathbf 0_{\mathrm I}$, 2.1.8) et son adhérence $Z'_\lambda$ dans
$p^{-1}(y_\alpha)$ serait donc de dimension $<e$ (4.1.2.1). Mais comme $h$
est quasi-fini, il résulte de (4.1.2) et de l'hypothèse $a$) que l'on a
$\dim(Z'_\lambda)\geq e$ (la restriction de $h$ à
$\overline{\{x_\lambda\}}$ se factorisant en
\[
  \overline{\{x_\lambda\}}\longrightarrow Z'_\lambda
  \longrightarrow p^{-1}(y_\alpha)
\]
en vertu de (\textbf{I}, 5.2.2))~; on aboutit ainsi à une contradiction, ce
qui montre que $a$) entraîne $b$).

Prouvons finalement que $b$) implique $a$). Notons que le morphisme structural
$p:Z=Y[T_1,\ldots,T_e]\to Y$ est fidèlement plat~; donc (2.3.4) les points
maximaux de $Z$ ont pour images par $p$ les points maximaux de $Y$~; ceci
prouve déjà que les $f(x_\lambda)$ sont points génériques des $Y_\alpha$. En
outre, si $f(x_\lambda)=y_\alpha$, le $k(y_\alpha)$-morphisme
$h:U\cap f^{-1}(y_\alpha)\to p^{-1}(y_\alpha)$ déduit de $g$ est dominant et
quasi-fini par hypothèse~; on a donc
$\dim(U_\lambda\cap f^{-1}(y_\alpha))=e$ en vertu de (4.1.2)
($U_\lambda$ étant la composante irréductible de $U$ de point générique
$x_\lambda$). De même, pour tout $x'\in U_\lambda$, le morphisme
$U_\lambda\cap f^{-1}(f(x'))\to p^{-1}(f(x'))$ déduit de $g$ est un
$k(f(x'))$-morphisme quasi-fini, donc (4.1.2) on a
$\dim(f^{-1}(f(x'))\cap U_\lambda)\leq e$~; mais par ailleurs on sait (13.1.6)
que toutes les composantes irréductibles de
$U_\lambda\cap f^{-1}(f(x'))$ sont de dimension $\geq e$~; on voit ainsi que
ces composantes sont exactement de dimension $e$, donc $b$) entraîne $a$).
C.Q.F.D.
\end{proof}

\begin{definition}[13.3.2]
\label{IV.13.3.2-fr}
Soient $Y$ un préschéma, $f:X\to Y$ un morphisme localement de type fini, $x$
un point de $X$. On dit que $f$ est équidimensionnel au point $x$ (ou que $X$
est équidimensionnel sur $Y$ au point $x$) si les conditions équivalentes de
(13.3.1) sont vérifiées. On dit que $f$ est équidimensionnel (ou que $X$ est
équidimensionnel sur $Y$) si $f$ est équidimensionnel en tout point $x\in X$.
\end{definition}

Il est clair que, pour que $f$ soit équidimensionnel en un point $x$, il faut
et il suffit que $f_{\mathrm{red}}$ le soit. En outre, les conditions de
(13.3.1) montrent que l'ensemble des points où $f$ est équidimensionnel est
\emph{ouvert} dans $X$.

On notera que lorsque $X$ et $Y$ sont \emph{irréductibles}, dire que $f$ est
équidimen-
\oldpage[IV-3]{197}
sionnel au point $x$ signifie, en vertu de (13.3.1, $a''$)), que $f$ est
dominant et que
$\dim_x(f^{-1}(f(x)))=\dim(f^{-1}(\eta))$ (où $\eta$ est le point générique de
$Y$)~; la définition (13.3.2) coïncide donc dans ce cas avec la définition
(13.2.2).

\begin{proposition}[13.3.3]
\label{IV.13.3.3-fr}
Soient $Y$ un préschéma localement noethérien, $f:X\to Y$ un morphisme
localement de type fini, $x$ un point de $X$, $X_j$ les composantes
irréductibles de $X$ (en nombre fini) contenant $x$. Pour que $f$ soit
équidimensionnel au point $x$, il faut et il suffit que, pour tout $j$,
$Y_j=\overline{f(X_j)}$ soit une composante irréductible de $Y$ et, en
désignant par $X_j$ et $Y_j$ les sous-préschémas fermés réduits de $X$ et $Y$
d'espaces sous-jacents $X_j$ et $Y_j$, par $f_j:X_j\to Y_j$ le morphisme
déduit de $f$ (\textbf{I}, 5.2.2), par $y_j$ le point générique de $Y_j$, que
$f_j$ soit équidimensionnel au point $x$ et que tous les nombres
$\dim(f_j^{-1}(y_j))$ soient égaux.
\end{proposition}

Cela résulte aussitôt de (13.3.1, $a'$)).

\begin{corollary}[13.3.4]
\label{IV.13.3.4-fr}
Avec les notations de (13.3.3), posons $y=f(x)$. Si $\mathcal O_{Y,y}$ est un
anneau universellement caténaire et si $f$ est équidimensionnel au point $x$,
on a
\[
\tag{13.3.4.1}
\label{IV.13.3.4.1-fr}
  \dim(\mathcal O_{X_j,x})=\dim(\mathcal O_{Y_j,y})+e
  -\deg.\operatorname{tr}_{k(y)}k(x)
\]
où $e$ est la valeur commune des nombres $\dim(f_j^{-1}(y_j))$.
\end{corollary}

Comme chacun des $\mathcal O_{Y_j,y}$, quotient de $\mathcal O_{Y,y}$, est un
anneau universellement caténaire (5.6.1), l'égalité est conséquence de (5.6.5).

\begin{corollary}[13.3.5]
\label{IV.13.3.5-fr}
Avec les notations de (13.3.4), supposons que $f$ soit équidimensionnel au
point $x$ et que $\mathcal O_{Y,y}$ soit un anneau universellement caténaire.
Si l'anneau $\mathcal O_{Y,y}$ est équidimensionnel, il en est de même de
$\mathcal O_{X,x}$~; la réciproque est vraie si l'image par $f$ de la réunion
des $X_j$ est dense dans un voisinage de $y$.
\end{corollary}

En effet, dire que $\mathcal O_{Y,y}$ est équidimensionnel signifie que les
nombres $\dim(\mathcal O_{Y'_i,y})$ sont égaux pour toutes les composantes
irréductibles $Y'_i$ de $Y$ contenant $y$, comme il résulte de (5.1.1.5)
appliqué au schéma local $\operatorname{Spec}(\mathcal O_{Y,y})$~; comme on a
la relation (13.3.4.1), il en résulte par le même raisonnement que
$\mathcal O_{X,x}$ est alors équidimensionnel. Inversement, si
$\mathcal O_{X,x}$ est équidimensionnel, tous les nombres
$\mathcal O_{Y_j,y}$ sont égaux par (13.3.4.1)~; on en déduit que
$\mathcal O_{Y,y}$ est équidimensionnel si les $Y_j$ sont \emph{toutes} les
composantes irréductibles de $Y$ contenant $y$, ce qui résulte de l'hypothèse
supplémentaire.

\begin{proposition}[13.3.6]
\label{IV.13.3.6-fr}
Soient $Y$ un préschéma localement noethérien, $f:X\to Y$ un morphisme
localement de type fini, $x$ un point de $X$, $y=f(x)$. Supposons que l'anneau
$\mathcal O_y$ soit équidimensionnel. Alors, si $\mathcal O_x$ est
équidimensionnel et si l'on a l'égalité
\[
\tag{13.3.6.1}
\label{IV.13.3.6.1-fr}
  \dim(\mathcal O_x)=\dim(\mathcal O_y)
  +\dim(\mathcal O_x\otimes_{\mathcal O_y}k(y))
\]
(cf. (13.2.7.1)) $f$ est équidimensionnel au point $x$, et la réciproque est
vraie si $\mathcal O_y$ est un anneau universellement caténaire.
\end{proposition}

Gardons les notations de (13.3.3)~; il résulte de (13.2.8) et de l'hypothèse
que $\mathcal O_x$ est équidimensionnel, que chacun des $Y_j$ est une
composante irréductible de $Y$ et que chacun des $f_j$ est équidimensionnel au
point $x$~; en outre (13.2.8), on a (compte tenu de (5.6.5))
\[
  \dim(\mathcal O_{X_j,x})=\dim(\mathcal O_{Y_j,y})
  +\dim_x(X_j\cap f^{-1}(y))-\deg.\operatorname{tr}_{k(y)}k(x).
\]

\oldpage[IV-3]{198}
Or, comme $\mathcal O_y$ est supposé équidimensionnel, cette égalité s'écrit
\[
\tag{13.3.6.2}
\label{IV.13.3.6.2-fr}
  \dim_x(X_j\cap f^{-1}(y))=\dim(\mathcal O_{X,x})
  -\dim(\mathcal O_{Y,y})+\deg.\operatorname{tr}_{k(y)}k(x).
\]
Le premier membre de (13.3.6.2) est donc indépendant de $j$~; mais comme $f_j$
est équidimensionnel au point $x$, on a
$\dim_x(X_j\cap f^{-1}(y))=\dim(f_j^{-1}(y_j))$, donc le critère (13.3.3)
montre que $f$ est équidimensionnel au point $x$.

Inversement, supposons que $\mathcal O_{Y,y}$ soit un anneau universellement
caténaire et que $f$ soit équidimensionnel au point $x$~; alors (13.3.5)
$\mathcal O_{X,x}$ est équidimensionnel, et il résulte alors de (13.3.4) que
l'on a la relation
\[
  \dim(\mathcal O_{X,x})=\dim(\mathcal O_{Y,y})+e
  -\deg.\operatorname{tr}_{k(y)}k(x)
\]
où $e$ est la valeur commune des nombres
$\dim(f_j^{-1}(y_j))=\dim_x(X_j\cap f^{-1}(y))$~; mais par définition $e$ est
aussi égal à $\dim_x(f^{-1}(y))$, d'où la relation (13.3.6.1), compte tenu de
(5.6.5.2).

\begin{proposition}[13.3.7]
\label{IV.13.3.7-fr}
Soient $Y$ un préschéma, $f:X\to Y$ un morphisme localement de type fini et
équidimensionnel, $Z$ une partie fermée de $X$. Alors la fonction
\[
\tag{13.3.7.1}
\label{IV.13.3.7.1-fr}
  x\longmapsto\operatorname{codim}_x
  (Z\cap f^{-1}(f(x)),f^{-1}(f(x)))
\]
est semi-continue inférieurement dans $X$.
\end{proposition}

Comme toutes les composantes irréductibles de $f^{-1}(f(x))$ contenant $x$
ont par hypothèse la même dimension (13.3.1), on a, en vertu de (5.2.1),
\[
  \operatorname{codim}_x(Z\cap f^{-1}(f(x)),f^{-1}(f(x)))
  =\dim_x(f^{-1}(f(x)))-\dim_x(Z\cap f^{-1}(f(x))).
\]
Mais par hypothèse le premier terme du second membre est fonction
\emph{continue} de $x$ (13.3.1) et le second est fonction semi-continue
supérieurement de $x$ en vertu de (13.3.3)~; d'où la conclusion.

\begin{proposition}[13.3.8]
\label{IV.13.3.8-fr}
Soient $Y$ un préschéma, $f:X\to Y$ un morphisme localement de type fini, $x$
un point de $X$. Soient $Y'$ un préschéma, $g:Y'\to Y$ un morphisme plat,
$X'=X\times_Y Y'$, $f'=f_{(Y')}:X'\to Y'$~; si $f$ est équidimensionnel au
point $x$, alors $f'$ est équidimensionnel en tout point $x'\in X'$ au-dessus
de $x$.
\end{proposition}

La question étant locale sur $X$ et $Y$, on peut se borner au cas où toute
composante irréductible de $X$ (resp. $Y$) contient $x$ (resp. $y$)~; comme
l'image par $f$ de toute composante irréductible de $X$ est alors dense dans
une composante irréductible de $Y$ (13.3.1), on sait (2.3.5) que l'image par
$f'$ de toute composante irréductible de $X'$ est dense dans une composante
irréductible de $Y'$. D'autre part, par transitivité des fibres
(\textbf{I}, 3.6.4), il résulte de (4.2.8) et de (13.3.1) que, pour tout
$z'\in X'$, l'ensemble des dimensions des composantes irréductibles de
${f'}^{-1}(f'(z'))$ est le même que l'ensemble des dimensions des composantes
irréductibles de $f^{-1}(f(z))$, où $z$ est la projection de $z'$ dans $X$~;
d'où la conclusion en vertu de (13.3.1, $a$)).

\begin{remark}[13.3.9]
\label{IV.13.3.9-fr}
La propriété d'équidimensionnalité d'un morphisme $f:X\to Y$ n'est \emph{pas}
stable par changement de base quelconque $g:Y'\to Y$, même lorsque $g$ est
l'injection canonique d'une composante irréductible $Y'$ de $Y$ dans $Y$
($Y'$ étant

\oldpage[IV-3]{199}
considérée comme sous-préschéma fermé réduit de $Y$). Par exemple, soient $k$
un corps, $A_0=k[S,T]$ l'anneau de polynômes à deux indéterminées,
$\mathfrak a=\mathfrak p_1\mathfrak p_2$, où $\mathfrak p_1$ et
$\mathfrak p_2$ sont les idéaux premiers $A_0S$ et $A_0T$ de $A_0$~; soient
$A=A_0/\mathfrak a$ et $Y=\operatorname{Spec}(A)$, qui a deux composantes
irréductibles $Y_1=\operatorname{Spec}(A/\mathfrak p_1)$,
$Y_2=\operatorname{Spec}(A/\mathfrak p_2)$~; prenons $X=Y_1$,
$f:X\to Y$ étant l'injection canonique, qui est évidemment un morphisme
équidimensionnel. Prenons d'autre part $Y'=Y_2$, $g:Y'\to Y$ étant l'injection
canonique~; alors on a
$X'=X\times_Y Y'=\operatorname{Spec}((A/\mathfrak p_1)\otimes_A
(A/\mathfrak p_2))=\operatorname{Spec}(A/(\mathfrak p_1+\mathfrak p_2))$,
et $\mathfrak p_1+\mathfrak p_2$ est un idéal maximal de $A$, donc $X'$ est
réduit à un point, et le morphisme $f':X'\to Y'$ n'est \emph{pas dominant},
l'image par $f'$ de l'unique point de $X'$ étant un point fermé de $Y'$~;
donc $f'$ n'est pas équidimensionnel.

On peut aussi donner un contre-exemple où $X$ et $Y$ sont \emph{intègres}, $f$
\emph{fini} et \emph{birationnel} (et \emph{a fortiori} équidimensionnel par
(13.3.1, $b$)), $g:Y'\to Y$ \emph{fini et dominant}. Soient $A$ et
$\overline A$ les anneaux locaux définis dans (\textbf{II}, 7.5), et prenons
$Y=\operatorname{Spec}(A)$, $X=\operatorname{Spec}(\overline A)$~; d'autre
part, avec les notations de (\textbf{II}, 7.5), prenons
$Y'=\operatorname{Spec}(\overline B)$~; alors
$X'=\operatorname{Spec}(\overline A\otimes_A\overline B)
=\operatorname{Spec}(\overline B\otimes_B\overline B)$~; mais on vérifie
aussitôt que $\overline B\otimes_B\overline B$ est composé direct des anneaux
$B/\mathfrak p'$, $B/\mathfrak p''$ et de deux anneaux isomorphes à
$B/\mathfrak n$, dont les spectres sont donc réduits à un point~; comme les
projections de ces points sont des points fermés de $Y'$, ici encore on voit
que $f'$ ne transforme pas une composante irréductible de $X'$ en une partie
partout dense d'une composante irréductible de $Y'$, donc $f'$ n'est pas
équidimensionnel.

Dans cet exemple, l'anneau $A$ n'est pas géométriquement unibranche~; nous
allons voir au paragraphe suivant (14.4.6) que de tels phénomènes ne peuvent se
produire lorsque les points de $Y$ sont \emph{géométriquement unibranches}. Le
manque de stabilité de la notion de morphisme équidimensionnel restreint
grandement son intérêt, au profit de la notion de morphisme universellement
ouvert, qui va être étudiée en détail au paragraphe suivant.
\end{remark}

% IV3_B28_P199_END_BEFORE_SECTION_14

\section{Morphismes universellement ouverts}
\label{section:IV.14-fr}

Les \S\S~14 et 15 sont consacrés à l'étude de la notion de morphisme
\emph{universellement ouvert} (2.4.2). On a déjà vu (2.4.6) qu'un morphisme
\emph{plat} et localement de présentation finie est universellement ouvert, la
réciproque étant inexacte. Au \S~14, on examine d'abord les propriétés des
dimensions des fibres d'un morphisme universellement ouvert $f:X\to Y$~;
lorsque $X$ et $Y$ sont localement noethériens, $f$ se comporte à cet égard
(14.2.1) comme un morphisme plat (cf. 6.1.2), et est en particulier
\emph{équidimensionnel} lorsqu'il est localement de type fini et dominant et que
$X$ est irréductible (14.2.2). Inversement, un morphisme équidimensionnel
localement de type fini $f:X\to Y$ est universellement ouvert lorsqu'on suppose
en outre que $Y$ est \emph{géométriquement unibranche}, et en particulier lorsque
$Y$ est \emph{normal} (critère de Chevalley, 14.4.4). Nous montrons aussi que les
morphismes universellement ouverts localement de type fini $f:X\to Y$ (lorsque
$X$ et $Y$ sont
\oldpage[IV-3]{200}
localement noethériens et $Y$ irréductible) admettent «~suffisamment~» de
\emph{quasi-sections}, i.e. au voisinage d'un point fermé $x$ d'une fibre
$f^{-1}(y)$, il y a un sous-préschéma fermé $X'$ de $X$ contenant $x$ et tel que
la restriction $X'\to Y$ de $f$ soit un morphisme \emph{quasi-fini} (donc à fibres
\emph{discrètes}) et dominant (14.5.3).

Au \S~15 on étudie diverses propriétés des fibres des morphismes universellement
ouverts $f:X\to Y$, localement de type fini, notamment lorsque $X$ et $Y$ sont
localement noethériens. On obtient ainsi en particulier un critère pour qu'un
point $x\in X$ n'appartienne qu'à \emph{une seule} composante irréductible de $X$,
en termes de propriétés de \emph{la fibre} $f^{-1}(f(x))$ de $x$~: il suffit que
$Y$ soit géométriquement unibranche (par exemple normal) au point $f(x)$ et que
$f^{-1}(f(x))$ soit géométriquement ponctuellement intègre au point $x$ (15.3.3)~;
si de plus $Y$ est localement intègre au point $f(x)$, $f$ est \emph{plat} au point
$x$ et $X$ localement intègre au point $x$. On étudie aussi la variation du
\emph{nombre géométrique de composantes connexes} d'une fibre $f^{-1}(y)$~; par
exemple, si $f$ est universellement ouvert et \emph{propre}, et les fibres de $f$
géométriquement réduites, ce nombre est \emph{localement constant} (15.5.7).
Enfin, lorsque $f$ admet une \emph{section} $g$ (ce qui sera le cas lorsque $X$ est
un $Y$-\emph{schéma en groupes}) et que pour tout $y\in Y$, on désigne par $X_y^0$
la composante connexe de la fibre $X_y=f^{-1}(y)$ au point $g(y)$ («~composante
neutre~» dans le cas des groupes), on étudie la réunion $X^0$ des $X_y^0$ pour
$y\in Y$, et on montre (15.6.4) que si $f$ est universellement ouvert et les fibres
$X_y$ géométriquement réduites, alors $X^0$ est un ensemble \emph{ouvert} dans $X$.

\subsection{Morphismes ouverts}
\label{subsection:IV.14.1-fr}

\begin{env}[14.1.1]
\label{IV.14.1.1-fr}
Rappelons (1.10.2) qu'une application continue $\psi:X\to Y$ est dite
\emph{ouverte en un point} $x\in X$ si l'image par $\psi$ de \emph{tout} voisinage
de $x$ dans $X$ est un voisinage de $\psi(x)$ dans $Y$.

On notera que cela \emph{n'implique pas} qu'il existe un système fondamental de
voisinages de $x$ dont les images soient \emph{ouvertes} dans $Y$.
\end{env}

\begin{proposition}[14.1.2]
\label{IV.14.1.2-fr}
Soient $X$, $Y$ deux espaces topologiques, $\psi:X\to Y$ une application
continue, $x$ un point de $X$, $y=\psi(x)$.
\begin{enumerate}[label=(\roman*)]
  \item Si $\psi$ est ouverte en $x$, alors pour toute partie $Y'$ de $Y$ contenant
  $\psi(x)$ la restriction $\psi^{-1}(Y')\to Y'$ de $\psi$ à $Y'$ est ouverte au
  point $x$.
  \item Supposons que $Y$ soit réunion d'une famille localement finie de parties
  fermées $(Y_i)$ et que pour tout $i$ tel que $\psi(x)\in Y_i$, la restriction
  $\psi^{-1}(Y_i)\to Y_i$ de $\psi$ soit ouverte au point $x$~; alors $\psi$ est
  ouverte au point $x$.
  \item Soient $\gamma:X'\to X$ une application continue, $x'$ un point de $X'$~;
  si l'application composée $\psi\circ\gamma:X'\to Y$ est ouverte au point $x'$,
  $\psi$ est ouverte au point $\gamma(x')$.
\end{enumerate}
\end{proposition}

\begin{proof}
Si $U$ est un voisinage de $x$, on a
$\psi(U\cap\psi^{-1}(Y'))=\psi(U)\cap Y'$~; d'où aussitôt (i). Pour prouver (ii),
notons qu'il y a un voisinage $W$ de $\psi(x)$ dans $Y$ ne rencontrant que les
parties fermées $Y_i$ en nombre \emph{fini} qui \emph{contiennent} $\psi(x)$~; donc
$W$ est réunion des $W\cap Y_i$ pour ces indices. Or, si $U$ est un voisinage de
$x$ tel que $\psi(U)\cap Y_i$ soit un voisinage de $\psi(x)$ dans $Y_i$, il existe
un voisinage $V\subset W$ de $\psi(x)$ dans $Y$ tel que pour tous les $i$ tels que
$\psi(x)\in Y_i$,
\oldpage[IV-3]{201}
on ait $V\cap Y_i\subset\psi(U)\cap Y_i$ et comme la réunion des $V\cap Y_i$ pour
ces indices est $V$, on a $V\subset\psi(U)$, donc $\psi(U)$ est un voisinage de
$\psi(x)$ dans $Y$. L'assertion (iii) est triviale.
\end{proof}

\begin{remarks}[14.1.3]
\label{IV.14.1.3-fr}
\begin{enumerate}[label=(\roman*)]
  \item L'ensemble $Z$ des points $x\in X$ où un morphisme est ouvert
  \emph{n'est pas nécessairement ouvert}. Par exemple, soient $K$ un corps, $A$
  l'anneau de polynômes $K[S,T]$, $V$ le plan affine $\operatorname{Spec}(A)$,
  $X$ le sous-préschéma fermé de $V$ «~réunion de la droite $X_1$ définie par
  $T=0$ et de la droite $X_2$ définie par $S=0$~», c'est-à-dire
  $\operatorname{Spec}(A/\mathfrak a)$, où $\mathfrak a=AST$~; prenons
  $Y=X_2=\operatorname{Spec}(A/AS)$ et pour $f$ la projection correspondant à
  l'injection canonique $K[T]\to A/\mathfrak a$~; alors on a $Z=X_2$, qui n'est
  pas ouvert dans $X$.

  \item Soient $X$, $Y$ deux préschémas noethériens irréductibles de points
  génériques $\xi$, $\eta$ respectivement, $f:X\to Y$ un morphisme
  \emph{dominant localement de type fini}~; alors $f$ \emph{est ouvert au point
  $\xi$}. En effet, on peut évidemment se borner au cas où $X$ et $Y$ sont réduits
  (donc intègres) (1.5.4) et affines~; en vertu du théorème de platitude générique
  (6.9.1), il existe un ouvert non vide $V$ de $Y$ tel que la restriction
  $f^{-1}(V)\to V$ de $f$ soit un morphisme \emph{plat}, et on conclut de (2.4.6)
  que cette restriction est un morphisme ouvert. Toutefois, comme il y a des
  morphismes dominants de type fini $f:X\to Y$ (où $X$ et $Y$ sont
  \emph{irréductibles}) qui ne sont pas ouverts (voir par exemple (\textbf{II},
  8.1)), l'ensemble des points où un tel morphisme est ouvert n'est pas
  nécessairement \emph{fermé} dans $X$.

  \item Nous ignorons si, lorsque $X$ et $Y$ sont localement noethériens et
  $f:X\to Y$ un morphisme localement de type fini, l'ensemble des points de $X$
  où $f$ est ouvert est ou non \emph{localement constructible}.
\end{enumerate}
\end{remarks}

\begin{proposition}[14.1.4]
\label{IV.14.1.4-fr}
Soient $X$, $Y$ deux espaces topologiques, $\psi:X\to Y$ une application
continue. Pour tout $y\in Y$, l'ensemble des $x\in\psi^{-1}(y)$ où $\psi$ est
ouvert est une partie fermée de $\psi^{-1}(y)$.
\end{proposition}

\begin{proof}
En effet, supposons que $\psi$ ne soit pas ouvert en un point
$x\in\psi^{-1}(y)$~; il existe alors un voisinage ouvert $V$ de $x$ dans $X$ tel
que $\psi(V)$ ne soit pas un voisinage de $y$~; il en résulte que pour tout
$x'\in V\cap\psi^{-1}(y)$, $\psi$ n'est pas ouvert au point $x'$.
\end{proof}

\begin{remark}[14.1.5]
\label{IV.14.1.5-fr}
Même si $X$ et $Y$ sont localement noethériens et $f$ un morphisme \emph{fini},
il peut se faire que $f$ soit ouvert en tous les points d'une fibre $f^{-1}(y)$,
sans qu'il existe un \emph{voisinage} de $f^{-1}(y)$ en tous les points duquel $f$
soit ouvert. Soient par exemple $K$ un corps, $A$ l'anneau de polynômes
$K[T_1,T_2,T_3]$, $V=\operatorname{Spec}(A)$ l'espace affine à 3 dimensions sur
$K$, $X=\operatorname{Spec}(A/\mathfrak a)$, où $\mathfrak a=\mathfrak b\mathfrak c$,
avec $\mathfrak b=(T_3)$ et $\mathfrak c=(T_1)+(T_2-T_3)$ dans $A$, de sorte que
$X$ est réunion du plan $X_1=\operatorname{Spec}(A/\mathfrak b)$ («~plan
d'équation $T_3=0$~») et de la droite $X_2=\operatorname{Spec}(A/\mathfrak c)$
(«~droite d'équations $T_1=0$, $T_2=T_3$~») qui sont ses composantes
irréductibles. Prenons $Y=X_1$ et soit $f:X\to Y$ la projection correspondant à
l'injection canonique $K[T_1,T_2]\to A/\mathfrak a$~; si $y$ est le point commun
à $X_1$ et $X_2$, $f^{-1}(y)$ est réduit à $y$ et $f$ est ouvert en ce point mais
n'est ouvert en aucun point de $X_2$ dans un voisinage de $y$ et distinct de $y$.
\end{remark}

\begin{env}[14.1.6]
\label{IV.14.1.6-fr}
Dans ce qui suit, le rôle essentiel sera joué par le critère (1.10.3)
caractérisant les morphismes \emph{localement de présentation finie} $f:X\to Y$
qui sont ouverts en un point $x$ par le «~relèvement des générisations~»~: pour
toute générisation $y'$ de $y=f(x)$, il existe $x'\in X$, générisation de $x$, tel
que $y'=f(x')$.
\end{env}

\oldpage[IV-3]{202}
\subsection{Morphismes ouverts et formule des dimensions}
\label{subsection:IV.14.2-fr}

\begin{theorem}[14.2.1]
\label{IV.14.2.1-fr}
Soient $X$, $Y$ deux préschémas localement noethériens, $f:X\to Y$ un morphisme,
$x$ un point de $X$, $y=f(x)$. On suppose que $f$ est ouvert aux points génériques
des composantes irréductibles de $f^{-1}(y)$ contenant $x$. Alors on a la relation
\begin{equation}
\tag{14.2.1.1}
\label{IV.14.2.1.1-fr}
  \dim(\mathcal O_x)=\dim(\mathcal O_y)+
  \dim(\mathcal O_x\otimes_{\mathcal O_y}k(y)).
\end{equation}
\end{theorem}

\begin{proof}
Soit $y'$ une générisation quelconque de $y$ distincte de $y$, et considérons le
sous-préschéma fermé réduit $Y'$ de $Y$ ayant pour espace sous-jacent
$\overline{\{y'\}}$~; alors aucune composante irréductible $X'$ de $f^{-1}(Y')$
contenant $x$ ne peut être contenue dans $f^{-1}(y)$~: en effet, si $x'$ est le
point générique de $X'$, on aurait $f(x')=y\ne y'$, et comme $x'$ est sa seule
générisation dans $f^{-1}(Y')$, cela contredirait l'hypothèse que $f$ est ouvert
au point $x'$, en vertu du fait que $a)$ implique $c)$ dans (1.10.3). On peut donc
appliquer (6.1.2) à l'homomorphisme local d'anneaux noethériens
$\mathcal O_y\to\mathcal O_x$, d'où la conclusion.
\end{proof}

\begin{corollary}[14.2.2]
\label{IV.14.2.2-fr}
Soient $Y$ un préschéma localement noethérien, $f:X\to Y$ un morphisme localement
de type fini, $x$ un point de $X$, $y=f(x)$. Supposons $f$ ouvert aux points
génériques des composantes irréductibles de $f^{-1}(y)$ contenant $x$. Alors $f$
est équidimensionnel au point $x$ dans chacun des deux cas suivants~:
\begin{enumerate}[label=(\roman*)]
  \item $X$ est irréductible et $f$ est dominant.
  \item Les anneaux $\mathcal O_x$ et $\mathcal O_y$ sont équidimensionnels.
\end{enumerate}
\end{corollary}

\begin{proof}
Pour (i), cela résulte de (14.2.1) et de (13.2.3). Pour (ii), cela résulte de
(14.2.1) et de (13.3.6).
\end{proof}

\begin{proposition}[14.2.3]
\label{IV.14.2.3-fr}
Soient $Y$ un préschéma irréductible localement noethérien, $\eta$ son point
générique, $f:X\to Y$ un morphisme localement de type fini, $y$ un point de
$f(X)$, $Z$ une composante irréductible de $f^{-1}(y)$ telle que $f$ soit ouvert
au point générique de $Z$. Alors $Z$ est contenue dans une composante irréductible
$X'$ de $X$ dominant $Y$ et telle que
\[
  \dim(Z)=\dim(X'\cap f^{-1}(y))=\dim(X'\cap f^{-1}(\eta)).
\]
\end{proposition}

\begin{proof}
Cela résulte en effet de (14.2.1) appliqué au point générique de $Z$, et de
(13.2.7).
\end{proof}

\begin{corollary}[14.2.4]
\label{IV.14.2.4-fr}
Avec les notations de (14.2.3), supposons que $f$ soit ouvert en tous les points
de $f^{-1}(y)$. Pour tout $z\in Y$, soient $E(z)$ l'ensemble des dimensions des
composantes irréductibles de $f^{-1}(z)$ et
$d(z)=\sup E(z)=\dim(f^{-1}(z))$. On a alors $E(y)\subset E(\eta)$, d'où
$d(y)\leq d(\eta)$.
\end{corollary}

\begin{proof}
Cela résulte aussitôt de (14.2.3).
\end{proof}

\begin{corollary}[14.2.5]
\label{IV.14.2.5-fr}
Soient $Y$ un préschéma localement noethérien, $f:X\to Y$ un morphisme propre,
$y\in Y$ un point tel que $f$ soit ouvert en tous les points de $f^{-1}(y)$. Alors
la fonction $z\mapsto\dim(f^{-1}(z))$ est constante dans un voisinage de $y$.
\end{corollary}

\begin{proof}
Soient $Y_i$ ($1\leq i\leq n$) les sous-préschémas fermés réduits de $Y$ ayant
pour espaces sous-jacents les composantes irréductibles de $Y$ contenant $y$~;
si $f_i:f^{-1}(Y_i)\to Y_i$ est la restriction de $f$, on sait que chacun des
$f_i$ est propre (\textbf{II}, 5.4.5) et ouvert en tous les points de $f^{-1}(y)$
(14.1.2). Comme la réunion des $Y_i$ est un voisinage de $y$ dans $Y$,
\oldpage[IV-3]{203}
on voit qu'on peut se borner à démontrer le corollaire lorsque $Y$ est
irréductible~; soit $\eta$ son point générique. Il résulte de (14.2.4) que
$d(y)\leq d(\eta)$~; d'autre part, comme $f$ est propre, on déduit de (13.1.5) que
$z\mapsto\dim(f^{-1}(z))$ est semi-continue supérieurement~; en particulier
$d(y)\geq d(\eta)$, donc $d(y)=d(\eta)$. De plus, il y a un voisinage $V$ de $y$
tel que $d(z)\leq d(y)$ pour tout $z\in V$~; mais comme $z$ est spécialisation de
$\eta$, on a d'autre part $d(z)\geq d(\eta)$, d'où finalement $d(z)=d(y)$ pour
$z\in V$.
\end{proof}

\begin{corollary}[14.2.6]
\label{IV.14.2.6-fr}
Soient $Y$ un préschéma localement noethérien, $f:X\to Y$ un morphisme ouvert de
type fini, $y$ un point de $Y$. Soient $y_i$ ($1\leq i\leq n$) les points
génériques des composantes irréductibles de $Y$ contenant $y$. Alors, avec les
notations de (14.2.4)~:
\begin{enumerate}[label=(\roman*)]
  \item Si pour $1\leq i\leq n$, on a $E(y)=E(y_i)$ (resp. $d(y)=d(y_i)$), la
  fonction $E$ (resp. $d$) est constante dans un voisinage de $y$.
  \item Il existe un voisinage ouvert $U$ de $f^{-1}(y)$ tel que la fonction
  $z\mapsto\dim(U\cap f^{-1}(z))$ soit constante dans un voisinage de $y$.
\end{enumerate}
\end{corollary}

\begin{proof}
\begin{enumerate}[label=(\roman*)]
  \item Le même raisonnement que dans (14.2.5) montre qu'on peut se borner à
  considérer le cas où $Y$ est irréductible, de point générique $\eta$. Si $y'$
  est une générisation quelconque de $y$, on a alors (en appliquant (14.2.4) à la
  restriction $f^{-1}(Y')\to Y'$ de $f$, où $Y'=\overline{\{y'\}}$, et utilisant
  (14.1.2)) $E(y)\subset E(y')\subset E(\eta)$ et
  $d(y)\leq d(y')\leq d(\eta)$~; cela montre que la relation $E(y)=E(\eta)$
  (resp. $d(y)=d(\eta)$) entraîne $E(y)=E(y')$ (resp. $d(y)=d(y')$) pour toute
  générisation $y'$ de $y$. Or, en vertu de (9.5.5), les fonctions $E$ et $d$ sont
  localement constructibles, donc il résulte de ($\mathbf{0}_{\mathrm{III}}$,
  9.2.5) appliqué à l'ensemble des $z$ tels que $E(z)=E(y)$ (resp.
  $d(z)=d(y)$) que cet ensemble est un voisinage de $y$.

  \item Pour chacun des $y_i$, $f^{-1}(y_i)$ est un espace noethérien puisque $f$
  est de type fini~; soient $x_{ij}$ ($1\leq j\leq m_i$) les points génériques de
  ses composantes irréductibles, et posons $X'_{ij}=\overline{\{x_{ij}\}}$~;
  montrons que le complémentaire $U$ dans $X$ de la réunion des $X'_{ij}$ qui ne
  rencontrent pas $f^{-1}(y)$ (qui est évidemment un voisinage ouvert de
  $f^{-1}(y)$) répond à la question. En effet, la restriction $f|U:U\to Y$ de $f$
  est un morphisme ouvert de type fini (I, 6.3.5)~; pour tout couple $(i,j)$ tel
  que $x_{ij}\in U$, $x_{ij}$ est point maximal de $U\cap f^{-1}(y_i)$, et il
  résulte de (13.1.1) appliqué à la restriction
  $X'_{ij}\to Y_i=\overline{\{y_i\}}$ de $f$ (compte tenu de (I, 5.2.2) et de
  (1.5.4)) que toutes les composantes irréductibles de
  $X'_{ij}\cap f^{-1}(y)$ ont une dimension
  $\geq\dim(X'_{ij}\cap f^{-1}(y_i))$. Compte tenu de (14.2.3) appliqué à la
  restriction $f^{-1}(Y_i)\to Y_i$ de $f$, qui est un morphisme ouvert (14.1.2),
  $f^{-1}(y)$ est, pour chaque $i$, la réunion des composantes irréductibles de
  $X'_{ij}\cap f^{-1}(y)$ ($1\leq j\leq m_i$), donc
  $\dim(f^{-1}(y))\geq\dim(U\cap f^{-1}(y_i))$~; mais comme $f|U$ est ouvert, on a
  aussi $\dim(f^{-1}(y))\leq\dim(f^{-1}(y_i)\cap U)$ en vertu de (14.2.4)~; on voit
  donc qu'on peut appliquer (i) à $f|U$, d'où la conclusion.
\end{enumerate}
\end{proof}

\begin{remark}[14.2.7]
\label{IV.14.2.7-fr}
L'exemple (13.2.12, (iii)) montre que sous les hypothèses de (14.2.5) ou
(14.2.6, (ii)), on ne peut, dans la conclusion, remplacer la fonction
$z\mapsto\dim(f^{-1}(z))$ par la fonction $z\mapsto E(z)$~; en effet, dans cet
exemple, $f$ est propre et plat (donc universellement ouvert (2.4.6)) et tout
voisinage de l'unique point de $f^{-1}(y_0)$ contient les points génériques des
composantes irréductibles de $f^{-1}(\eta)$.
\end{remark}

\oldpage[IV-3]{204}
\subsection{Morphismes universellement ouverts}
\label{subsection:IV.14.3-fr}

\begin{env}[14.3.1]
\label{IV.14.3.1-fr}
Rappelons (2.4.2) que dire qu'un morphisme $f:X\to Y$ est universellement ouvert
signifie que pour tout morphisme $g:Y'\to Y$, $f_{(Y')}:X_{(Y')}\to Y'$ est
ouvert~; il suffit d'ailleurs qu'il en soit ainsi lorsque
$Y'=Y[T_1,\ldots,T_e]$, pour tout $e$ (8.10.2) (et si $Y$ est localement
noethérien, il suffit donc qu'il en soit ainsi pour tout $Y'$ localement
noethérien).
\end{env}

\begin{proposition}[14.3.2]
\label{IV.14.3.2-fr}
Soit $f:X\to Y$ un morphisme de préschémas. Si $f$ est universellement ouvert,
alors, pour tout morphisme $g:Y'\to Y$, où $Y'$ est irréductible, l'image par
$f'=f_{(Y')}:X_{(Y')}\to Y'$ de toute composante irréductible de $X_{(Y')}$ est
dense dans $Y'$. Réciproquement, si cette condition est vérifiée pour tout $Y'$
irréductible et tout morphisme de type fini $g:Y'\to Y$, et si de plus $f$ est
localement de présentation finie, alors $f$ est universellement ouvert.
\end{proposition}

\begin{proof}
En effet, il résulte de (1.10.4) que si $f$ est universellement ouvert, il vérifie
la condition de l'énoncé. Inversement, supposons que $f$ soit localement de
présentation finie, et montrons que pour tout entier $e$, si l'on pose
$Y''=Y[T_1,\ldots,T_e]$, $f_{(Y'')}$ est ouvert. En effet, soit $Y'$ un
sous-préschéma fermé de $Y''$ ayant pour espace sous-jacent une partie fermée
irréductible de $Y''$~; le morphisme composé $Y'\to Y''\to Y$ est de type fini,
donc toute composante irréductible de $X_{(Y')}$ domine $Y'$ par hypothèse~; on
déduit donc de (1.10.4) que $f_{(Y'')}$ est un morphisme ouvert.
\end{proof}

Cette proposition montre que la définition (\textbf{III}, 4.3.9) coïncide dans
le cas considéré avec la définition générale des morphismes universellement
ouverts donnée dans (2.4.2).

A la notion de morphisme ouvert en un point (14.1.1) correspond de même la
suivante~:

\begin{definition}[14.3.3]
\label{IV.14.3.3-fr}
Soient $f:X\to Y$ un morphisme de préschémas, $x$ un point de $X$. On dit que $f$
est \emph{universellement ouvert au point $x$} si, pour tout morphisme $g:Y'\to Y$,
en posant $X'=X\times_Y Y'$, le morphisme $f'=f_{(Y')}:X'\to Y'$ est ouvert en
tout point $x'$ de $X'$ dont la projection dans $X$ est $x$.
\end{definition}

\begin{remarks}[14.3.3.1]
\label{IV.14.3.3.1-fr}
\begin{enumerate}[label=(\roman*)]
  \item Le raisonnement de (8.10.1) montre (avec les mêmes notations) que si $x$
  est un point de $X$, $x_\lambda$ sa projection dans $X_\lambda$, alors, si
  $f_\mu$ est ouvert au point $x_\mu$ pour tout $\mu\geq\lambda$, $f$ est ouvert
  au point $x$~; il suffit de se borner aux ouverts $U$ de $X$ contenant $x$ et de
  remarquer que l'hypothèse implique que $f_\mu(U_\mu)$ est un voisinage de
  $f_\mu(x_\mu)$, donc $f(v_\mu^{-1}(U_\mu))$ est un voisinage de $f(x)$. On en
  déduit que l'énoncé (8.10.2) est encore exact quand on y remplace
  «~universellement ouvert~» par «~universellement ouvert au point $x$~», et
  «~morphisme ouvert~» par «~morphisme ouvert en tout point $x_n$ de $X_n$ dont
  la projection dans $X$ est $x$~»~: il suffit dans la démonstration de se limiter
  aux ouverts $V$ contenant un $x_n$.

  \item Le résultat de (14.1.4) est encore valable pour un morphisme $f:X\to Y$,
  en y remplaçant «~ouvert~» par «~universellement ouvert~». En effet, supposons
  que $f$ ne soit pas universellement ouvert en un point $x\in f^{-1}(y)$~; il y a
  par suite un morphisme $Y'\to Y$ et un point $x'\in X'$ se projetant en $x$, tels
  que $f'$ ne soit pas ouvert au point $x'$.
\end{enumerate}
\end{remarks}

\oldpage[IV-3]{205}

Or, si $y'=f'(x')$, la projection $f'^{-1}(y')\to f^{-1}(y)$ est un morphisme
ouvert (2.4.10), et il y a, en vertu de (14.1.4), un voisinage $V'$ de $x'$ dans
$f'^{-1}(y')$ où $f'$ n'est pas ouvert~; donc $f$ n'est pas universellement ouvert
aux points de l'image de $V'$ dans $f^{-1}(y)$, qui est un voisinage de $x$ dans
$f^{-1}(y)$.

\begin{proposition}[14.3.4]
\label{IV.14.3.4-fr}
\begin{enumerate}[label=(\roman*)]
  \item Soient $f:X\to Y$, $g:Y\to Z$ deux morphismes, $x$ un point de $X$,
  $y=f(x)$. Si $f$ est universellement ouvert au point $x$ et $g$ universellement
  ouvert au point $y$, alors $g\circ f$ est universellement ouvert au point $x$.
  Inversement, si $g\circ f$ est universellement ouvert au point $x$, $g$ est
  universellement ouvert au point $y$.

  \item Si $f:X\to Y$ est un $S$-morphisme universellement ouvert au point
  $x\in X$, alors, pour tout changement de base $S'\to S$,
  $f_{(S')}:X_{(S')}\to Y_{(S')}$ est universellement ouvert en tout point de
  $X_{(S')}$ au-dessus de $x$.

  \item Pour que $f:X\to Y$ soit universellement ouvert au point $x\in X$, il
  faut et il suffit que $f_{\mathrm{red}}$ le soit.

  \item Soient $f:X\to Y$ un morphisme localement de présentation finie, $x$ un
  point de $X$, $y=f(x)$~; posons $Y_1=\operatorname{Spec}(\mathcal O_y)$,
  $X_1=X\times_Y Y_1$, $f_1=f_{(Y_1)}$. Pour que $f$ soit universellement ouvert
  au point $x$, il faut et il suffit que $f_1$ le soit (on rappelle (I, 3.6.5) que
  $X_1$ est canoniquement identifié à un sous-espace de $X$).
\end{enumerate}
\end{proposition}

\begin{proof}
En effet (ii) est une conséquence évidente de la définition (14.3.3)~; il résulte
aussi de la définition que pour démontrer l'assertion (i), il suffit de le faire
lorsqu'on y supprime partout le mot «~universellement~», et cela résulte alors de
(14.1.2). L'assertion (iii) résulte de (ii) et de ce que le morphisme canonique
$X_{\mathrm{red}}\to X$ est surjectif. Enfin, la condition (iv) est trivialement
nécessaire. D'autre part, si elle est vérifiée, et si $g:Y'\to Y$ est un morphisme
quelconque, $X'=X_{(Y')}$, $f'=f_{(Y')}$, $x'$ un point de $X'$ au-dessus de $x$,
alors $f'$ est localement de présentation finie, et pour voir qu'il est ouvert au
point $x'$, il suffit d'appliquer le critère (1.10.3, $b$). Soient
$y'=f'(x')$, $Y'_1=\operatorname{Spec}(\mathcal O_{y'})$,
$X'_1=X'\times_{Y'}Y'_1$, $f'_1=f_{(Y'_1)}$. Comme $g(y')=y$, le morphisme
composé $Y'_1\to Y'\to Y$ se factorise en $Y'_1\to Y_1\to Y$ (I, 2.4.4), donc
$X'_1=X_1\times_{Y_1}Y'_1$ et $f'_1=(f_1)_{(Y'_1)}$~; la conclusion résulte alors
aussitôt de l'hypothèse que $f'_1$ est ouvert au point $x'$ et de (1.10.3).
\end{proof}

\begin{proposition}[14.3.5]
\label{IV.14.3.5-fr}
Soient $X$, $Y$ deux préschémas, $f:X\to Y$ un morphisme, $x$ un point de $X$.
Soit $(Y_i)_{1\leq i\leq n}$ une famille localement finie de sous-préschémas
fermés de $Y$ telle que l'espace $Y$ soit réunion des $Y_i$, et supposons que pour
tout $i$ tel que $f(x)\in Y_i$, la restriction $f^{-1}(Y_i)\to Y_i$ de $f$ soit un
morphisme universellement ouvert au point $x$~; alors $f$ est universellement
ouvert au point $x$.
\end{proposition}

\begin{proof}
Compte tenu de la définition, cela résulte de la proposition analogue (14.1.2,
(ii)) pour les morphismes ouverts en un point.
\end{proof}

\begin{proposition}[14.3.6]
\label{IV.14.3.6-fr}
Soient $Y$ un préschéma localement noethérien, $f:X\to Y$ un morphisme localement
de type fini. Pour que $f$ soit universellement ouvert en un point $x\in X$, il
faut et il suffit que la condition suivante soit vérifiée~: pour tout morphisme
$g:Y'\to Y$, où $Y'=\operatorname{Spec}(A)$ est le spectre d'un anneau de
valuation discrète, tel que l'image $g(y')$ du point fermé $y'$ de $Y'$ soit égale
à $y=f(x)$, et pour tout point $x'\in X'=X\times_Y Y'$ dont les projections sur
$X$ et $Y'$ sont $x$ et $y'$ respectivement, il existe une générisation $z'$ de
$x'$ dans $X'$ dont la projection dans $Y'$ est le point générique de $Y'$
(autrement dit, il existe une composante irréductible de $X'$ contenant $x'$ et
dominant $Y'$).

\oldpage[IV-3]{206}

De plus, on peut, dans la condition précédente, se borner au cas où $A$ est
complet, a un corps résiduel algébriquement clos, et où $x'$ est rationnel sur
$k(y')$.
\end{proposition}

\begin{proof}
Si $Y'$ est comme dans l'énoncé, la nécessité de la condition résulte de ce que
$f'$ doit être ouvert au point $x'$, et du critère (1.10.3). Pour voir que la
condition est suffisante, considérons un morphisme de type fini $g:Y''\to U$, et
soient $X''=X\times_Y Y''$, $f''=f_{(Y'')}:X''\to Y''$, et $x''$ un point de $X''$
au-dessus de $x$. Posons $y''=f''(x'')$, et soit $t$ une générisation de $y''$
dans $Y''$, distincte de $y''$. Puisque $Y''$ est localement noethérien, il
résulte de (\textbf{II}, 7.1.9) et ($\mathbf{0}_{\mathrm{III}}$, 10.3.1) qu'il
existe un schéma $Y'=\operatorname{Spec}(A)$, où $A$ est un anneau de valuation
discrète complet, dont le corps résiduel est une clôture algébrique de $k(x'')$,
et un morphisme $g:Y'\to Y''$ tels que, si $s$ et $y'$ sont le point générique et
le point fermé de $Y'$, on ait $g(s)=t$ et $g(y')=y''$. Il y a alors un point $x'$
de $X'=X\times_Y Y'=X''\times_{Y''}Y'$ dont les projections dans $X''$ et $Y'$
soient $x''$ et $y'$ respectivement, et qui soit rationnel sur $k(y')$ (I, 3.4.9).
L'hypothèse entraîne qu'il y a une générisation $z'$ de $x'$ dans $X'$ dont la
projection dans $Y'$ soit $s$~; si $z''$ est la projection de $z'$ dans $X''$,
$z''$ est une générisation de $x''$ et sa projection dans $Y''$ est $t$~; on
conclut donc de (1.10.3) que $f''$ est ouvert au point $x''$, donc que $f$ est
universellement ouvert au point $x$ (14.3.3.1, (i)).
\end{proof}

\begin{corollary}[14.3.7]
\label{IV.14.3.7-fr}
Les notations étant celles de (14.3.6)~:
\begin{enumerate}[label=(\roman*)]
  \item Étant donné un point $y\in Y$, pour que $f$ soit universellement ouvert en
  tous les points de $f^{-1}(y)$, il faut et il suffit que pour tout morphisme
  $g:Y'\to Y$, où $Y'$ est le spectre d'un anneau de valuation discrète, et où
  l'image $g(y')$ du point fermé $y'$ de $Y'$ est $y$, toute composante
  irréductible de $X'=X\times_Y Y'$ domine $Y'$.
  \item Pour que $f$ soit universellement ouvert, il faut et il suffit que pour
  tout morphisme $g:Y'\to Y$, où $Y'$ est le spectre d'un anneau de valuation
  discrète, toute composante irréductible de $X'=X\times_Y Y'$ domine $Y'$.
\end{enumerate}
\end{corollary}

\begin{proof}
Il est clair qu'il suffit de prouver (i)~; la nécessité de (i) résulte de (14.3.2)
et sa suffisance de (14.3.6).
\end{proof}

\begin{proposition}[14.3.8]
\label{IV.14.3.8-fr}
Soient $Y$ un préschéma localement noethérien, irréductible, régulier et de
dimension 1 (par exemple le spectre d'un anneau de Dedekind), $f:X\to Y$ un
morphisme localement de type fini, $y$ un point de $Y$. Les conditions suivantes
sont équivalentes~:
\begin{enumerate}[label=\emph{\alph*)}]
  \item $f_{\mathrm{red}}$ est plat en tout point de $f^{-1}(y)$.
  \item $f$ est universellement ouvert dans un voisinage de $f^{-1}(y)$.
  \item $f$ est ouvert dans un voisinage de $f^{-1}(y)$.
  \item Toute composante irréductible de $X$ rencontrant $f^{-1}(y)$ domine $Y$.
\end{enumerate}
\end{proposition}

\begin{proof}
Comme $f_{\mathrm{red}}$ est localement de type fini (1.3.4), $a)$ entraîne que
$f_{\mathrm{red}}$ est plat dans un voisinage de $f^{-1}(y)$ (11.1.1), et il
suffit d'appliquer (2.4.6) dans un tel voisinage pour voir que $a)$ entraîne $b)$.
L'implication $b)\Rightarrow c)$ est triviale, et l'implication
$c)\Rightarrow d)$ résulte de (1.10.4) appliqué à un voisinage de $f^{-1}(y)$.
Reste à voir que $d)$ entraîne $a)$. On peut évidemment, en vertu de (1.3.4), se
borner au cas où $X$ est réduit. La question étant en outre locale sur $X$ et sur
$Y$, on peut supposer $Y=\operatorname{Spec}(A)$ et
$X=\operatorname{Spec}(B)$
\oldpage[IV-3]{207}
affines~; si $X_i$ ($1\leq i\leq n$) sont les sous-préschémas fermés (intègres)
de $X$ ayant pour espaces sous-jacents les composantes irréductibles de $X$,
alors, pour tout $x\in X$, $\mathcal O_{X,x}$, étant réduit, est un sous-anneau du
composé direct des $\mathcal O_{X_i,x}$~; si $y=f(x)$, il suffira de montrer que
chacun des $\mathcal O_{X_i,x}$ est un $\mathcal O_y$-module \emph{sans torsion},
car il en sera alors de même de $\mathcal O_{X,x}$~; comme par hypothèse
$\mathcal O_y$ est un anneau local régulier de dimension 1, c'est-à-dire un anneau
de valuation discrète (\textbf{II}, 7.1.6), il résultera alors de
($\mathbf{0}_{\mathrm I}$, 6.3.4) que $\mathcal O_{X,x}$ est un
$\mathcal O_y$-module \emph{plat}. Mais si $X_i=\operatorname{Spec}(B_i)$, où
$B_i$ est un anneau intègre, l'hypothèse $d)$ entraîne que l'homomorphisme
$A\to B_i$ est \emph{injectif} (I, 1.2.7)~; donc $B_i$ est un $A$-module sans
torsion, et \emph{a fortiori} $\mathcal O_{X_i,x}$ est un $\mathcal O_y$-module
sans torsion.
\end{proof}

\begin{remarks}[14.3.9]
\label{IV.14.3.9-fr}
\begin{enumerate}[label=(\roman*)]
  \item Dans l'énoncé de (14.3.8), on ne peut se dispenser de l'hypothèse que $Y$
  est \emph{régulier}. Avec les notations de (11.7.5), prenons en effet
  $Y=\operatorname{Spec}(A)$, $X=\operatorname{Spec}(\overline A)$, de sorte que
  $f:X\to Y$ est un morphisme \emph{fini surjectif}~; comme $A$ est un anneau
  local intègre de dimension 1, ainsi que $\overline A$, il résulte aussitôt de
  (1.10.4) que le morphisme $f$ est \emph{ouvert}. Cependant $f$ n'est pas
  universellement ouvert (ni \emph{a fortiori} plat), comme le montre (11.7.5).
  On aurait un exemple analogue en prenant pour $Y$ le schéma local au point
  double d'une courbe algébrique ayant un «~point double ordinaire~» et pour $X$
  le normalisé de $Y$.

  \item L'exemple de (14.1.3, (i)) montre que l'ensemble des points de $X$ où un
  morphisme est universellement ouvert n'est pas nécessairement ouvert, le
  morphisme $f$ dans cet exemple étant universellement ouvert en tous les points
  où il est ouvert. L'exemple $f:X\to Y$ vu ci-dessus en (i) montre de même que
  l'ensemble des points où un morphisme est universellement ouvert n'est pas
  nécessairement fermé, car il est immédiat qu'en tous les points de $X$ sauf un
  le morphisme $f$ est universellement ouvert (c'est même un isomorphisme local).
  Il serait intéressant de savoir si l'ensemble des points où un morphisme est
  universellement ouvert est localement constructible.
\end{enumerate}
\end{remarks}

Les deux propositions suivantes nous ont été signalées par M.~Artin~:

\begin{lemma}[14.3.10]
\label{IV.14.3.10-fr}
Soient $A$ un anneau de valuation (non nécessairement discrète), $k$ son corps
résiduel, $K$ son corps des fractions, et posons $S=\operatorname{Spec}(A)$. Soient
$X$ un préschéma irréductible, $f:X\to S$ un morphisme dominant de type fini~; soit
$X_0=X\otimes_A k$ (resp. $X_1=X\otimes_A K$) la fibre de $f$ au point fermé
(resp. au point générique) de $S$. Alors, si $X_0\ne\varnothing$, on a
$\dim(X_0)=\dim(X_1)$.
\end{lemma}

On peut se borner au cas où $X$ est affine, en remplaçant au besoin $X$ par un
ouvert affine contenant un point générique d'une composante irréductible de $X_0$,
de dimension maximale, et utilisant (4.1.1.3). Soit $n=\dim(X_0)\geq 0$~; il résulte
de (13.3.1.1) qu'il existe un voisinage $U$ de $X_0$ dans $X$ et un $S$-morphisme
quasi-fini
\[
  g:U\to S[T_1,\ldots,T_n]=Z
\]
tel que le morphisme restriction
$g_0:U_0=U\cap X_0\to Z_0=\operatorname{Spec}(k[T_1,\ldots,T_n])$ soit fini et
\emph{surjectif}. Par changement de base $\operatorname{Spec}(K)\to S$ et restriction
à l'ouvert $U_1=U\cap X_1$ de $X_1$, on déduit de $g$ un morphisme quasi-fini
$g_1:U_1\to Z_1=\operatorname{Spec}(K[T_1,\ldots,T_n])$. Comme $U_1$ est dense dans
$X_1$, on a $\dim(U_1)=\dim(X_1)$ (4.1.1.3)~; la proposition sera établie, en vertu
de (4.1.2), si nous prouvons que le morphisme $g_1$ est \emph{dominant}. Supposons
le contraire~; il existerait alors un polynôme $F_1\ne 0$ dans
$K[T_1,\ldots,T_n]$ tel que $g_1(U_1)\subset V(F_1)$. Si $\omega$ est une valuation
sur $K$ associée à $A$, et si $(a_\alpha)$ est la famille des coefficients de $F_1$,
on peut, après multiplication de $F_1$ par un élément $\ne 0$ de $K$, supposer que
l'on a $\inf_\alpha(\omega(a_\alpha))=0$~; autrement dit, $F_1$ provient d'un
polynôme $F\in A[T_1,\ldots,T_n]$ (avec lequel il s'identifie) tel que l'image $F_0$
de $F$ dans $k[T_1,\ldots,T_n]$ soit \emph{non nul}. Considérons alors dans $Z$
l'ensemble fermé $V(F)$~; on a $V(F)\cap Z_1=V(F_1)$, et comme $U_1$ est dense
dans $U$ (puisqu'il contient le point générique de $X$), $g(U)\subset V(F)$~; en
particulier, on aurait $g_0(U_0)\subset V(F)\cap Z_0=V(F_0)$~; mais puisque
$F_0\ne 0$, $V(F_0)$ est une partie fermée de $Z_0$ \emph{distincte de $Z_0$} et
l'on aboutit à une contradiction. C.Q.F.D.

\oldpage[IV-3]{208}

\begin{proposition}[14.3.11]
\label{IV.14.3.11-fr}
Soient $f:X\to S$ un morphisme de type fini, $(g_\lambda)_{\lambda\in L}$ une
famille de morphismes universellement ouverts $g_\lambda:Y_\lambda\to S$, et pour
tout $\lambda\in L$, soit $u_\lambda:Y_\lambda\to X$ un $S$-morphisme. Pour tout
$s\in S$, posons $X_s=X\times_S\operatorname{Spec}(k(s))$,
$(Y_\lambda)_s=Y_\lambda\times_S\operatorname{Spec}(k(s))$, et soit
$(u_\lambda)_s:(Y_\lambda)_s\to X_s$ le morphisme $u_\lambda\times 1$ déduit de
$u_\lambda$ par changement de base. Soit $Z(s)$ l'adhérence dans $X_s$ de la
réunion des ensembles $(u_\lambda)_s((Y_\lambda)_s)$, et posons
$d(s)=\dim(Z(s))$. Alors, pour toute générisation $s'$ d'un point $s\in S$, on a
$d(s)\leq d(s')$.
\end{proposition}

\begin{proof}
On sait (\textbf{II}, 7.1.4) qu'il existe un anneau de valuation $A'$ et un
morphisme $h:S'=\operatorname{Spec}(A')\to S$ tels que si $a$ (resp. $b$) est le
point fermé (resp. le point générique) de $S'$, on ait $h(a)=s$, $h(b)=s'$. En
outre, le morphisme projection
$p:X_s\otimes_{k(s)}k(a)\to X_s$ est surjectif et ouvert (2.4.10), donc fait de
$X_s$ un espace quotient de $X_s\otimes_{k(s)}k(a)$ par une relation d'équivalence
ouverte~; pour toute partie $M$ de $X_s$, $p^{-1}(\overline M)$ est donc égal à
l'adhérence $\overline{p^{-1}(M)}$ (Bourbaki, \emph{Top. gén.}, chap.~I\textsuperscript{er},
4\textsuperscript{e} éd., \S~5, n\textsuperscript{o}~3, prop.~7)~; on raisonne de
même pour $X_{s'}$, et compte tenu de (I, 3.4.8), de (4.2.7) et du fait que les
$g_\lambda$ sont universellement ouverts, on voit qu'on peut se ramener à prouver
la proposition sur la situation obtenue après changement de base $S'\to S$.
Supposons donc $S'=S$, $s$ étant le point fermé et $s'$ le point générique de $S$.
L'hypothèse que $g_\lambda$ est ouvert entraîne que toute composante irréductible
de $Y_\lambda$ domine $S$ (1.10.4), donc que son point générique est un point
maximal de $(Y_\lambda)_{s'}$~; si $Z$ désigne l'adhérence dans $X$ de $Z(s')$, on
a donc $u_\lambda(Y_\lambda)\subset Z$, et par suite
$(u_\lambda)_s((Y_\lambda)_s)\subset Z_s=Z\cap X_s$~; autrement dit, on a
$Z(s)\subset Z_s$, d'où $\dim(Z(s))\leq\dim(Z_s)$. Mais en appliquant (14.3.10) à
un sous-préschéma réduit de $X$ ayant pour espace sous-jacent une composante
irréductible de $Z$, on obtient $\dim(Z_s)\leq\dim(Z_{s'})$ (l'inégalité provenant
de ce qu'il peut y avoir des composantes irréductibles de $Z$ ne rencontrant pas
$X_s$). D'autre part, comme $Z(s')$ est par définition fermé dans $X_{s'}$, on a
$Z_{s'}=Z(s')$, ce qui achève la démonstration.
\end{proof}

\begin{remark}[14.3.12]
\label{IV.14.3.12-fr}
Le cas envisagé par M.~Artin était celui où $Y_\lambda=S$ pour tout $\lambda$,
autrement dit le cas où $(u_\lambda)$ est une famille de $S$-sections de $X$. Un
autre cas utile est celui où la famille $(u_\lambda)$ est réduite à un seul
élément~; on peut d'ailleurs toujours se ramener à ce cas en considérant le
préschéma $Y$ somme des $Y_\lambda$ et les morphismes $g:Y\to S$ et $u:Y\to X$
dont les restrictions à chaque $Y_\lambda$ sont respectivement $g_\lambda$ et
$u_\lambda$.
\end{remark}

\begin{proposition}[14.3.13]
\label{IV.14.3.13-fr}
Soient $f:X\to Y$ un morphisme localement de type fini, $y$ un point de $Y$, $x$
un point maximal de la fibre $X_y=X\times_Y\operatorname{Spec}(k(y))$.
Considérons les conditions suivantes~:
\begin{enumerate}[label=\emph{\alph*)}]
  \item $f$ est universellement ouvert au point $x$ (ou encore en tout point de la
  composante irréductible de $X_y$ de point générique $x$ (14.3.3.1, (ii))).
  \item Pour toute composante irréductible $Y_0$ de $Y$ contenant $y$, il existe
  une composante irréductible $Z$ de $X$ contenant $x$, dominant $Y_0$ et telle que
  $\dim_x(X_y)=\dim_x(Z\cap X_y)=\dim(Z\cap X_\eta)$, où $\eta$ est le point
  générique de $Y_0$ (ce qui entraîne que $Z$ est équidimensionnelle sur $Y_0$ au
  point $x$ (13.2.2)).
\end{enumerate}

\oldpage[IV-3]{209}

\begin{enumerate}[label=\emph{\alph*$'$)}]
  \item[b$'$)] Pour tout voisinage ouvert $U$ de $x$ dans $X$ et toute
  générisation $y'$ de $y$, on a
  $\dim(U\cap X_{y'})\geq\dim_x(U\cap X_y)$.
\end{enumerate}
Alors on a les implications $a)\Rightarrow b)\Leftrightarrow b')$.
\end{proposition}

\begin{proof}
Pour montrer que $b)$ implique $b')$, il suffit de remarquer que $y'$ appartient à
une composante irréductible $Y_0$ de $Y$ contenant $y$, de point générique $\eta$~;
prenant $Z$ comme dans $b)$ et notant que le point générique de $Z$ (qui est aussi
celui de $Z\cap X_\eta$ ($\mathbf{0}_{\mathrm I}$, 2.1.8)) est contenu dans $U$, on
a $\dim(U\cap X_{y'})\geq\dim(U\cap Z\cap X_{y'})$, et, en vertu de (13.1.6),
$\dim(U\cap Z\cap X_{y'})\geq\dim(Z\cap X_\eta)$~; d'où l'assertion, puisque
$\dim_x(U\cap X_y)=\dim_x(X_y)$ (4.1.1.3).

Pour prouver que $b')$ implique $b)$, on peut d'abord remplacer $X$ par
$f^{-1}(Y_0)$, donc supposer $Y_0=Y$~; on peut se borner au cas où $X$ et $Y$ sont
affines, puisque $\dim_x(U\cap X_y)=\dim_x(X_y)$ (4.1.1.3). Les composantes
irréductibles $W_i$ du préschéma noethérien $X_\eta$ sont alors en nombre fini, et
le complémentaire de la réunion des $\overline W_i$ (adhérences dans $X$) qui ne
contiennent pas $x$ est un voisinage ouvert $V$ de $x$. En remplaçant $X$ par $V$,
on peut donc supposer que $x\in\overline W_i$ pour tout $i$ (l'hypothèse $b')$
entraîne $\dim(V\cap X_\eta)\geq 0$, donc $x$ appartient à un des $\overline W_i$
pour un $i$ au moins). En outre, les $\overline W_i$ sont exactement les
composantes irréductibles de $X$ qui dominent $Y$ ($\mathbf{0}_{\mathrm I}$,
2.1.8)~; cela étant, si l'on avait, pour chacune de ces composantes $Z$,
$\dim(Z\cap X_\eta)<\dim_x(X_y)$, on en conclurait
$\dim(X_\eta)<\dim_x(X_y)$, contrairement à l'hypothèse $b')$. La relation
$\dim_x(Z\cap X_y)=\dim(Z\cap X_\eta)$ résulte alors de (13.1.6).

Reste à démontrer que $a)$ entraîne $b')$. Tenant compte de (\textbf{II}, 7.1.4)
et de l'invariance des hypothèses et de la conclusion par changement de base (en
vertu de (4.2.7)), on peut se borner au cas où $Y$ est un spectre d'anneau de
valuation, de point fermé $y$ et de point générique $y'$, et où $U=X$.
L'hypothèse que $f$ est ouvert au point $x$ entraîne qu'il existe une composante
irréductible $Z$ de $X$ contenant $x$ et dominant $Y$ (1.10.3). Appliquant
(14.3.10) à un voisinage de $x$ dans $Z$, on en conclut que
$\dim(Z\cap X_{y'})=\dim_x(Z\cap X_y)$~; mais comme $x$ est maximal dans $X_y$,
$Z\cap X_y$ contient la composante irréductible de $X_y$ de point générique $x$,
donc $\dim_x(Z\cap X_y)=\dim_x(X_y)$~; d'autre part, on a
$\dim(X_{y'})\geq\dim(Z\cap X_{y'})$, ce qui achève de prouver $b')$.
\end{proof}

\begin{remark}[14.3.14]
\label{IV.14.3.14-fr}
Nous ignorons si dans (14.3.13), la conclusion subsiste lorsqu'on remplace
l'hypothèse $a)$ par l'hypothèse plus faible que $f$ est ouvert au point $x$. On
peut montrer aisément qu'il suffirait de traiter le cas où $Y$ est spectre d'un
anneau local intègre, dont le point générique est isolé, et où $X$ est un
sous-préschéma fermé du fibré vectoriel $Y[T]$.
\end{remark}

\subsection{Le critère de Chevalley pour les morphismes universellement ouverts}
\label{subsection:IV.14.4-fr}

\begin{theorem}[14.4.1]
\label{IV.14.4.1-fr}
Soient $f:X\to Y$ un morphisme localement de type fini, $y$ un point de $Y$, $x$
un point maximal de la fibre $X_y=X\times_Y\operatorname{Spec}(k(y))$. Supposons
$y$ géométriquement unibranche. Alors les conditions suivantes sont équivalentes~:
\begin{enumerate}[label=\emph{\alph*)}]
  \item $f$ est universellement ouvert au point $x$ (ou encore en tout point de la
  composante irréductible de $X_y$ de point générique $x$ (14.3.3.1, (iii))).

  \oldpage[IV-3]{210}

  \item Si $Y_0$ est l'unique composante irréductible de $Y$ contenant $y$, $\eta$
  son point générique, il existe une composante irréductible $Z$ de $X$ contenant
  $x$, dominant $Y_0$ et telle que
  $\dim_x(Z\cap X_y)=\dim(Z\cap X_\eta)$ (ce qui signifie que $Z$ est
  équidimensionnelle sur $Y_0$ au point $x$ (13.2.2)).
  \item[b$'$)] Pour tout voisinage ouvert $U$ de $x$ dans $X$ et toute
  générisation $y'$ de $y$, on a
  $\dim(U\cap X_{y'})\geq\dim_x(U\cap X_y)$.
\end{enumerate}
Si de plus $Y$ est localement noethérien, ces conditions sont aussi équivalentes à
la suivante~:
\begin{enumerate}[label=\emph{\alph*)}]
  \item[c)] $f$ est ouvert au point $x$.
\end{enumerate}
\end{theorem}

Notons d'abord que puisque $(\mathcal O_y)_{\mathrm{red}}$ est intègre, $y$
n'appartient qu'à une seule composante irréductible $Y_0$ de $Y$. Le fait que $b)$
et $b')$ sont équivalentes et que $a)$ implique $b')$ résulte de (14.3.13)~;
d'autre part, si $Y$ est localement noethérien, on a vu dans (14.2.3) que $c)$
implique $b)$. Il reste donc à montrer que lorsque $y$ est géométriquement
unibranche, $b)$ entraîne $a)$.

\begin{lemma}[14.4.1.1]
\label{IV.14.4.1.1-fr}
Soient $Y$ un préschéma, $Y'=Y[T_1,\ldots,T_n]$, $y'$ un point de $Y'$, $y$ son
image dans $Y$. Si $Y$ est géométriquement unibranche au point $y$, alors $Y'$ est
géométriquement unibranche au point $y'$.
\end{lemma}

En effet, comme pour tout $y\in Y$,
$\operatorname{Spec}(k(y)[T_1,\ldots,T_n])$ est géométriquement régulier sur $k(y)$
($\mathbf{0}_{\mathrm{IV}}$, 17.3.7), le morphisme structural $Y'\to Y$ est lisse
(6.8.1), et il suffit d'appliquer (11.3.14).

\begin{lemma}[14.4.1.2]
\label{IV.14.4.1.2-fr}
Soient $A$ un anneau local intègre unibranche, $B$ un anneau intègre contenant $A$
et entier sur $A$, $\mathfrak n$ un idéal premier de $B$ au-dessus de l'idéal
maximal $\mathfrak m$ de $A$. Alors le morphisme
$\operatorname{Spec}(B_{\mathfrak n})\to\operatorname{Spec}(A)$ est surjectif,
autrement dit, pour tout idéal premier $\mathfrak p$ de $A$, il existe un idéal
premier $\mathfrak q$ de $B$ tel que $\mathfrak q\subset\mathfrak n$ et
$\mathfrak q\cap A=\mathfrak p$.
\end{lemma}

Soient $K$ (resp. $L$) le corps des fractions de $A$ (resp. $B$), $A'$ la clôture
intégrale de $A$, $B'$ le sous-anneau de $L$ engendré par $A'$ et $B$, de sorte
qu'on a un diagramme commutatif d'injections canoniques
\[
  \xymatrix{
    B \ar[r] & B' \\
    A \ar[u] \ar[r] & A' \ar[u]
  }
\]
Comme $B'$ est entier sur $B$, il existe un idéal premier $\mathfrak n'$ de $B'$
tel que $\mathfrak n'\cap B=\mathfrak n$ (Bourbaki, \emph{Alg. comm.}, chap.~V,
\S~2, n\textsuperscript{o}~1, th.~1), et (pour la même raison)
$\operatorname{Spec}(A')\to\operatorname{Spec}(A)$ est surjectif. D'autre part,
comme $A$ est unibranche, $A'$ est un anneau \emph{local}, donc
$\mathfrak n'\cap A'$, qui est au-dessus de l'idéal maximal $\mathfrak m$ de $A$,
est nécessairement égal à l'unique idéal maximal $\mathfrak m'$ de $A'$. En vertu
du second théorème de Cohen--Seidenberg (\emph{loc. cit.}, \S~2,
n\textsuperscript{o}~4, th.~3) le morphisme
$\operatorname{Spec}(B'_{\mathfrak n'})\to\operatorname{Spec}(A')$ est surjectif,
donc il en est de même du composé
$\operatorname{Spec}(B'_{\mathfrak n'})\to\operatorname{Spec}(A')\to
\operatorname{Spec}(A)$~; mais ce morphisme est aussi le composé
$\operatorname{Spec}(B'_{\mathfrak n'})\to\operatorname{Spec}(B_{\mathfrak n})
\to\operatorname{Spec}(A)$, donc le morphisme
$\operatorname{Spec}(B_{\mathfrak n})\to\operatorname{Spec}(A)$ est surjectif.

Ces lemmes étant établis, revenons à la démonstration de l'implication $b)\Rightarrow
a)$ dans (14.4.1). En vertu de (14.3.3.1, (i)), il suffit de prouver que, pour tout
entier $n\geq 0$ et tout point $x'$ de $X'=X[T_1,\ldots,T_n]$ au-dessus de $x$, le
morphisme

\oldpage[IV-3]{211}

$f':X'\to Y'=Y[T_1,\ldots,T_n]$, déduit de $f$ par changement de base, est ouvert
au point $x'$. Compte tenu du lemme (14.4.1.1), de (2.3.4) et (4.2.7), on est donc
ramené à prouver que $f$ est \emph{ouvert} au point $x$~: en outre, il suffit
évidemment (14.1.2, (iii)) de montrer que la restriction de $f$ à un
sous-préschéma fermé de $X$ ayant $Z$ pour espace sous-jacent est ouverte au point
$x$, si bien qu'on peut se borner au cas où $X=Z$ est irréductible. Quitte à
remplacer $X$ par un voisinage ouvert $V$ de $x$ tel que $V\cap X_y$ soit
irréductible, on peut supposer, en vertu de (13.3.1), que le morphisme $f$ se
factorise en
\[
  X\xrightarrow{g}Y''=Y[T_1,\ldots,T_m]\to Y,
\]
où $g$ est quasi-fini, dominant et localement de type fini. Comme le morphisme
structural $Y''\to Y$ est ouvert (2.4.6), on est ramené à prouver que $g:X\to Y''$
est ouvert au point $x$. D'ailleurs, en vertu de (14.4.1.1), $g(x)$ est un point
géométriquement unibranche de $Y''$. On est donc ramené à prouver le lemme suivant~:

\begin{lemma}[14.4.1.3]
\label{IV.14.4.1.3-fr}
Soient $X$, $Y$ deux préschémas irréductibles, $f:X\to Y$ un morphisme localement
quasi-fini et dominant. Si $x\in X$ est tel que $y=f(x)$ soit unibranche sur $Y$,
alors $f$ est ouvert au point $x$.
\end{lemma}

Il suffit de prouver que
$f(\operatorname{Spec}(\mathcal O_{X,x}))=\operatorname{Spec}(\mathcal O_{Y,y})$
(1.10.3). On peut donc se borner, par le changement de base
$\operatorname{Spec}(\mathcal O_{Y,y})\to Y$, au cas où $Y=\operatorname{Spec}(A)$,
où $A$ est un anneau local et $y$ est le point fermé de $Y$ (tenant compte de
(I, 3.6.5) et ($\mathbf{0}_{\mathrm I}$, 2.1.8), qui prouvent que
$X\times_Y\operatorname{Spec}(\mathcal O_{Y,y})$ est irréductible)~; remplaçant $f$
par $f_{\mathrm{red}}$, on peut supposer $X$ et $Y$ réduits, donc intègres.
Remplaçant éventuellement $X$ et $Y$ par des voisinages affines de $x$ et $y$
respectivement, on peut supposer (8.12.9) que le morphisme $f$ se factorise en
$X\xrightarrow{j}X_1\xrightarrow{g}Y$, où $j$ est une immersion ouverte et $g$ un
morphisme \emph{fini} (évidemment dominant)~; comme $X$ et $X_1$ sont affines, $j$
est affine, donc séparé et quasi-compact, et par suite se factorise en
$X\xrightarrow{u}X_2\xrightarrow{h}X_1$, où $X_2$ est l'image fermée de $X$ par
$j$, $h$ l'injection canonique et $u$ une immersion ouverte (I, 9.5.3). En d'autres
termes, on peut supposer que $X_1$ est intègre, ou encore de la forme
$X_1=\operatorname{Spec}(B)$, où $B$ est une $A$-algèbre intègre et finie, contenant
$A$ puisque $g$ est dominant. Si $\mathfrak n$ est l'idéal premier de $B$
correspondant au point $x$, l'hypothèse que $A$ est unibranche implique alors
(14.4.1.2) que le morphisme
$\operatorname{Spec}(B_{\mathfrak n})\to\operatorname{Spec}(A)$ est surjectif,
c'est-à-dire que
$\operatorname{Spec}(\mathcal O_{X,x})\to\operatorname{Spec}(\mathcal O_{Y,y})$
est surjectif. C.Q.F.D.

\begin{corollary}[14.4.2]
\label{IV.14.4.2-fr}
Soient $f:X\to Y$ un morphisme localement de type fini, $y$ un point
géométriquement unibranche de $Y$, $\eta$ le point générique de l'unique composante
irréductible $Y_0$ de $Y$ contenant $y$. Les conditions suivantes sont
équivalentes~:
\begin{enumerate}[label=\emph{\alph*)}]
  \item $f$ est universellement ouvert en tous les points de $X_y$ (ou, ce qui
  revient au même (14.3.3.1, (ii)), aux points maximaux de $X_y$).
  \item Pour tout $x\in X_y$, il existe une composante irréductible $Z$ de $X$
  contenant $x$ et équidimensionnelle sur $Y$ au point $x$ (13.2.2).
  \item[b$'$)] Pour tout $x\in X_y$ et tout voisinage ouvert $U$ de $x$ dans $X$,
  on a $\dim(U\cap X_\eta)\geq\dim_x(U\cap X_y)$.
  \item[b$''$)] Pour tout ouvert $U$ de $X$, on a
  $\dim(U\cap X_\eta)\geq\dim(U\cap X_y)$.
\end{enumerate}

\oldpage[IV-3]{212}

Lorsque de plus $Y$ est localement noethérien, ces conditions sont encore
équivalentes à la suivante~:
\begin{enumerate}[label=\emph{\alph*)}]
  \item[c)] $f$ est ouvert en tous les points de $X_y$ (ou, ce qui revient au même
  (14.3.3.1, (ii)), aux points maximaux de $X_y$).
\end{enumerate}
\end{corollary}

L'équivalence de $a)$ et $c)$ lorsque $Y$ est localement noethérien résulte de
(14.4.1)~; les conditions $b)$ ou $b')$, appliquées aux \emph{points maximaux} de
$X_y$, entraînent $a)$ en vertu également de (14.3.3.1, (ii))~; enfin, $b')$ et
$b'')$ sont équivalentes, car
\[
  \dim(U\cap X_y)=\sup_{x\in X_y}(\dim_x(U\cap X_y)).
\]
Reste à voir que la condition $a)$ entraîne $b)$ et $b')$ en \emph{tout} point
$x\in X_y$. Posons $d=\dim_x(X_y)$, et soit $x'$ le point générique d'une
composante irréductible de $X_y$ contenant $x$ et de dimension $d$. En vertu de
$a)$ et de (14.4.1), il y a une composante irréductible $Z$ de $X$ contenant $x'$
et équidimensionnelle sur $Y$ au point $x'$, donc telle que
$\dim_{x'}(Z\cap X_y)=\dim(Z\cap X_\eta)$. Mais par construction
$\dim_{x'}(Z\cap X_y)=\dim_{x'}(X_y)=d$, et
$\dim_x(Z\cap X_y)\leq\dim(Z\cap X_y)=d$~; compte tenu de (13.1.6), cela prouve
que $Z$ est équidimensionnel sur $Y$ au point $x$~; donc $a)$ entraîne $b)$. De
plus, on a $\dim(X_\eta)\geq\dim(Z\cap X_\eta)=d=\dim_x(X_y)$. Remplaçant $X$ par
un voisinage ouvert $U$ de $x$, on voit ainsi que $a)$ entraîne $b')$. C.Q.F.D.

\begin{corollary}[14.4.3]
\label{IV.14.4.3-fr}
Soient $Y$ un préschéma géométriquement unibranche, $f:X\to Y$ un morphisme
localement de type fini. Les conditions suivantes sont équivalentes~:
\begin{enumerate}[label=\emph{\alph*)}]
  \item $f$ est universellement ouvert.
  \item Pour tout ouvert $U$ de $X$, tout $y\in Y$ et toute générisation $y'$ de
  $y$, on a $\dim(U\cap X_{y'})\geq\dim(U\cap X_y)$.
\end{enumerate}
Si de plus $Y$ est localement noethérien, ces conditions sont aussi équivalentes à
la suivante~:
\begin{enumerate}[label=\emph{\alph*)}]
  \item[c)] $f$ est ouvert.
\end{enumerate}
\end{corollary}

\begin{corollary}[14.4.4]
\label{IV.14.4.4-fr}
(critère de Chevalley). Soit $f:X\to Y$ un morphisme localement de type fini.
\begin{enumerate}[label=(\roman*)]
  \item Si $f$ est équidimensionnel en un point $x\in X$ (13.3.2) et si
  $y=f(x)$ est un point géométriquement unibranche de $Y$, $f$ est
  universellement ouvert au point $x$.
  \item Si $Y$ est géométriquement unibranche, $f$ est universellement ouvert en
  tous les points de $X$ où $f$ est équidimensionnel, et l'ensemble de ces points
  est ouvert dans $X$. En particulier, si $f$ est équidimensionnel, il est
  universellement ouvert.
\end{enumerate}
\end{corollary}

L'assertion (ii) est conséquence triviale de (i), puisqu'on sait déjà que l'ensemble
des points où $f$ est équidimensionnel est ouvert (13.3.2). Quant à l'assertion (i),
elle résulte de ce que l'hypothèse implique que la condition $b)$ de (14.4.1) est
vérifiée au point générique d'une composante irréductible de $X_y$ contenant $x$,
compte tenu de (13.3.1)~; il suffit donc d'appliquer (14.4.1).

\begin{remark}[14.4.5]
\label{IV.14.4.5-fr}
On peut prouver que si $Y$ est localement noethérien, et si toutes les générisations
de $y=f(x)$ sont des points géométriquement unibranches de $Y$ (cf. (6.15.2)),
alors, si $f$ est équidimensionnel au point $x$, il est universellement ouvert
\emph{dans un voisinage de $x$}.
\end{remark}

\oldpage[IV-3]{213}

\begin{corollary}[14.4.6]
\label{IV.14.4.6-fr}
Soient $Y$ un préschéma localement noethérien, $f:X\to Y$ un morphisme localement
de type fini. Soit $y$ un point géométriquement unibranche de $Y$, et supposons en
outre que pour tout $x\in f^{-1}(y)$, l'anneau $\mathcal O_{X,x}$ soit
équidimensionnel. Alors les conditions suivantes sont équivalentes~:
\begin{enumerate}[label=\emph{\alph*)}]
  \item $f$ est équidimensionnel (13.3.2) en tous les points de $f^{-1}(y)$.
  \item $f$ est ouvert en tous les points de $f^{-1}(y)$.
  \item $f$ est universellement ouvert en tous les points de $f^{-1}(y)$.
\end{enumerate}
\end{corollary}

En effet, $c)$ implique trivialement $b)$, et $a)$ implique $c)$ en vertu de
(14.4.4)~; enfin, en raison des hypothèses sur $\mathcal O_y$ et $\mathcal O_x$,
$b)$ implique $a)$ par (14.2.2). Plus généralement~:

\begin{proposition}[14.4.7]
\label{IV.14.4.7-fr}
Soient $Y$ un préschéma localement noethérien, $f:X\to Y$ un morphisme localement
de type fini, $x$ un point de $X$, tel que $y=f(x)$ soit géométriquement unibranche.
Les conditions suivantes sont équivalentes~:
\begin{enumerate}[label=\emph{\alph*)}]
  \item $f$ est équidimensionnel (13.3.2) au point $x$.
  \item L'anneau $\mathcal O_x$ est équidimensionnel et $f$ est ouvert aux points
  génériques des composantes irréductibles de $f^{-1}(y)$ contenant $x$ (donc
  aussi en tout point d'une telle composante).
  \item L'anneau $\mathcal O_x$ est équidimensionnel, et $f$ est universellement
  ouvert aux points génériques des composantes irréductibles de $f^{-1}(y)$
  contenant $x$ (donc aussi en tout point d'une telle composante).
\end{enumerate}
En outre, lorsque ces conditions sont satisfaites, pour tout sous-préschéma fermé
réduit $X_i$ de $X$ ayant pour espace sous-jacent une composante irréductible de $X$
contenant $x$, la restriction $f_i:X_i\to Y$ de $f$ est un morphisme
équidimensionnel au point $x$, et universellement ouvert en tous les points des
composantes irréductibles de $f^{-1}(y)\cap X_i$ qui contiennent $x$.
\end{proposition}

La condition $a)$ implique que $f$ est équidimensionnel en tous les points
génériques des composantes irréductibles de $f^{-1}(y)$ contenant $x$ (13.3.1), et
par suite (14.4.4) universellement ouvert en ces points~; le même raisonnement
appliqué à chaque $f_i$ (compte tenu de (13.3.3)) prouve la dernière assertion de
la proposition, compte tenu de (14.3.3.1, (iii)). En outre, d'après (14.2.1), on a
les relations
\begin{equation}
\tag{14.4.7.1}
\label{IV.14.4.7.1-fr}
  \dim(\mathcal O_{X_i,x})=\dim(\mathcal O_y)+
  \dim(\mathcal O_{X_i,x}\otimes_{\mathcal O_y}k(y))
\end{equation}
\begin{equation}
\tag{14.4.7.2}
\label{IV.14.4.7.2-fr}
  \dim(\mathcal O_{X,x})=\dim(\mathcal O_y)+
  \dim(\mathcal O_{X,x}\otimes_{\mathcal O_y}k(y))
\end{equation}
et puisque $f$ est équidimensionnel au point $x$, il résulte de (13.3.1) que l'on
a $\dim_x(X_i\cap f^{-1}(y))=\dim_x(f^{-1}(y))$, donc (5.2.3)
\[
  \dim(\mathcal O_{X_i,x}\otimes_{\mathcal O_y}k(y))=
  \dim(\mathcal O_{X,x}\otimes_{\mathcal O_y}k(y)).
\]
On conclut donc que $\dim(\mathcal O_{X_i,x})=\dim(\mathcal O_{X,x})$ pour tout
$i$, autrement dit $\mathcal O_{X,x}$ est équidimensionnel, et cela achève de
montrer que $a)$ entraîne $c)$. Il est clair que $c)$ entraîne $b)$~; enfin, $b)$
entraîne la relation (14.4.7.2) en vertu de (14.2.1)~; il résulte alors de (13.3.6)
que $b)$ entraîne $a)$.

\begin{proposition}[14.4.8]
\label{IV.14.4.8-fr}
Soient $Y$ un préschéma noethérien, $f:X\to Y$ un morphisme localement de type
fini, $y$ un point de $Y$, $x$ un point maximal de $X_y=f^{-1}(y)$. Les conditions
suivantes sont équivalentes~:
\begin{enumerate}[label=\emph{\alph*)}]
  \item Le morphisme $f$ est universellement ouvert au point $x$, autrement dit,
  pour tout changement de base $g:Y'\to Y$, on a la propriété $P(Y')$~: pour tout
  point $x'$ de $X'=X\times_Y Y'$ au-dessus de $x$, le morphisme
  $f'=f_{(Y')}:X'\to Y'$ est ouvert au point $x'$.

  \oldpage[IV-3]{214}

  \item[a$'$)] La propriété $P(Y')$ est vraie pour tout morphisme fini
  $g:Y'\to Y$.
  \item[a$''$)] La propriété $P(Y'')$ est vraie pour le normalisé $Y''$ de
  $Y_{\mathrm{red}}$ (\textbf{II}, 6.3.8).
  \item Pour tout point $x''$ de $X''=X\times_Y Y''$ au-dessus de $x$, il existe
  une composante irréductible $Z''$ de $X''$ contenant $x''$ et
  équidimensionnelle sur $Y''$ au point $x''$.
\end{enumerate}
\end{proposition}

Il est trivial que $a)$ implique $a')$. Pour montrer que $a')$ implique $a'')$
notons que l'on peut écrire $Y''=\operatorname{Spec}(\mathcal B)$, où $\mathcal B$
est une $\mathcal O_Y$-Algèbre quasi-cohérente entière sur $\mathcal O_Y$~; comme
$Y$ est noethérien, $\mathcal B$ est limite inductive de ses sous-$\mathcal O_Y$-
Algèbres $\mathcal B_\lambda$ qui sont quasi-cohérentes et de type fini (I, 9.6.6)~;
mais alors les $\mathcal B_\lambda$ sont des $\mathcal O_Y$-Algèbres \emph{finies}
(\textbf{II}, 6.1.2)~; on peut donc écrire
$Y''=\varprojlim Y'_\lambda$, où
$Y'_\lambda=\operatorname{Spec}(\mathcal B_\lambda)$, d'où
$X''=X\times_Y Y''=\varprojlim X'_\lambda$, avec
$X'_\lambda=X\times_Y Y'_\lambda$. En vertu de $a')$, les morphismes
$f'_\lambda:X'_\lambda\to Y'_\lambda$ sont ouverts en tous les points de
$X'_\lambda$ au-dessus de $x$~; on en conclut que $f''$ est ouvert en tous les
points de $X''$ au-dessus de $x$, par (8.10.1) et (14.3.3.1, (i)).

Comme le préschéma $Y''$ est \emph{normal} par définition, le fait que $b)$ entraîne
$a'')$ résulte de (14.4.4) appliqué à la composante irréductible
équidimensionnelle de l'énoncé et à la restriction de $f''$ à cette composante. Il
reste donc à montrer que $a'')$ entraîne $a)$ et $b)$. Compte tenu de (1.10.3), on
peut se borner au cas où $Y=\operatorname{Spec}(\mathcal O_y)$, en notant que le
morphisme canonique $\operatorname{Spec}(\mathcal O_y)\to Y$ est universellement
bicontinu (I, 3.6.5), et d'autre part que
$Y''\times_Y\operatorname{Spec}(\mathcal O_y)$ est le normalisé de
$(\operatorname{Spec}(\mathcal O_y))_{\mathrm{red}}$ comme il résulte de la
permutabilité des opérations de fermeture intégrale et de localisation (Bourbaki,
\emph{Alg. comm.}, chap.~V, \S~1, n\textsuperscript{o}~5, prop.~16). Supposant donc
$Y=\operatorname{Spec}(A)$, où $A$ est un anneau local noethérien, et $y$ le point
fermé de $Y$, on sait ($\mathbf{0}_{\mathrm{IV}}$, 23.2.5) qu'il existe une
factorisation
\[
  Y''\xrightarrow{u}Y_1\xrightarrow{v}Y
\]
du morphisme structural, telle que $v$ soit un morphisme \emph{fini surjectif}, $u$
un morphisme entier, \emph{radiciel et dominant}, donc surjectif (\textbf{II},
6.1.10) et par suite un \emph{homéomorphisme universel} (2.4.5). Si l'on pose
$X_1=X\times_Y Y_1$, la projection $X''\to X_1$ est donc un homéomorphisme, et
l'hypothèse $a'')$ entraîne par suite que $f_1=f_{(Y_1)}:X_1\to Y_1$ est ouvert en
tous les points de $X_1$ au-dessus de $x$. En outre, cela montre que pour prouver
la propriété $b)$, il suffit de prouver la même propriété où l'on remplace $Y''$,
$X''$ et $x''$ par $Y_1$, $X_1$ et un point $x_1$ de $X_1$ au-dessus de $x$~: Mais
$Y_1$ est noethérien et en outre il est géométriquement unibranche puisque $Y''$
est normal et $u$ radiciel (6.15.1)~; la propriété à prouver résulte donc de
(14.4.1). Il reste à montrer que $f$ est universellement ouvert au point $x$, ce
qui résultera du lemme suivant~:

\begin{lemma}[14.4.8.1]
\label{IV.14.4.8.1-fr}
Soit $v:Y_1\to Y$ un morphisme fermé (resp. universellement fermé) et surjectif.
Pour qu'un morphisme $f:X\to Y$ soit ouvert (resp. universellement ouvert) en un
point $x\in X$, il suffit que $f_1=f_{(Y_1)}:X_1=X\times_Y Y_1\to Y_1$ soit ouvert
(resp. universellement ouvert) en tous les points de $X_1$ au-dessus de $x$.
\end{lemma}

La seconde assertion résulte trivialement de la première et du fait que pour tout
changement de base $Y'\to Y$, le morphisme
$v_{(Y')}:Y_1\times_Y Y'\to Y'$ est encore surjectif et est fermé si $v$ est
universellement fermé. Pour prouver la première assertion, considérons

\oldpage[IV-3]{215}

un voisinage ouvert $U$ de $x$ dans $X$~; comme $v$ est fermé et surjectif, pour
que $f(U)$ soit un voisinage de $y=f(x)$, il faut et il suffit que
$v^{-1}(f(U))$ soit un voisinage de $v^{-1}(y)$. Mais si $p:X_1\to X$ est la
projection canonique, on a $v^{-1}(f(U))=f_1(p^{-1}(U))$ (I, 3.4.8), et
l'hypothèse entraîne que $f_1(p^{-1}(U))$ est un voisinage de
$f_1(p^{-1}(x))=v^{-1}(y)$ (I, 3.4.8).

\begin{corollary}[14.4.9]
\label{IV.14.4.9-fr}
Soient $Y$ un préschéma noethérien, $f:X\to Y$ un morphisme localement de type
fini. Les conditions suivantes sont équivalentes~:
\begin{enumerate}[label=\emph{\alph*)}]
  \item $f$ est universellement ouvert, autrement dit, pour tout changement de
  base $Y'\to Y$, le morphisme $f_{(Y')}:X\times_Y Y'\to Y'$ est ouvert.
  \item[a$'$)] Pour tout morphisme fini $Y_1\to Y$, $f_{(Y_1)}$ est ouvert.
  \item[a$''$)] Si $Y''$ est le normalisé de $Y_{\mathrm{red}}$, $f_{(Y'')}$ est
  ouvert.
  \item Pour tout point $x''$ de $X''=X\times_Y Y''$, il existe une composante
  irréductible $Z''$ de $X''$ contenant $x''$ et équidimensionnelle sur $Y''$ au
  point $x''$ (cf. (14.4.10, (ii))).
\end{enumerate}
\end{corollary}

Cela résulte aussitôt de (14.4.8) et (14.1.4).

\begin{remarks}[14.4.10]
\label{IV.14.4.10-fr}
\begin{enumerate}[label=(\roman*)]
  \item L'équivalence des conditions $a)$ et $b)$ dans (14.4.8) (resp. (14.4.9))
  est encore valable pour un préschéma $Y$ quelconque et un morphisme $f$
  localement de type fini. En effet, $a)$ entraîne $b)$ en vertu de (14.4.1)~;
  inversement, $b)$ entraîne que $f''$ est universellement ouvert aux points de
  $X''$ au-dessus de $x$ (resp. en tout point de $X''$) en vertu de (14.4.1), et
  on en conclut la propriété $a)$ en appliquant le lemme (14.4.8.1) au morphisme
  entier et surjectif $Y''\to Y$.

  Il se peut que, dans (14.4.1), pour l'équivalence de $a)$ et $c)$, l'hypothèse
  supplémentaire que $Y$ est noethérien soit superflue (cf. (14.3.14)). S'il en
  est ainsi, les hypothèses noethériennes sont également superflues dans (14.4.2),
  (14.4.3), (14.4.8) et (14.4.9).

  \item On peut donner des exemples de morphismes $f:X\to Y$ ayant les propriétés
  suivantes~: $Y$ est noethérien, régulier et de dimension 2, $f$ est
  \emph{universellement ouvert} et de type fini, $X$ a deux composantes
  irréductibles $X_1$, $X_2$, mais la restriction $X_1\to Y$ de $f$ à l'une
  d'elles \emph{n'est pas un morphisme ouvert}. Le principe de la construction
  s'appuie sur la méthode générale de «~recollement~» qui sera exposée au chap.~V,
  et ne peut donc qu'être esquissé ici. On part d'un point fermé $y$ de $Y$, et on
  considère le $Y$-schéma $Y_1$ obtenu en faisant \emph{éclater} $y$
  (\textbf{II}, 8.1.3)~; si $f_1:Y_1\to Y$ est le morphisme structural, on sait que
  la restriction de $f_1$ à $f_1^{-1}(Y-\{y\})$ est un \emph{isomorphisme} sur
  $Y-\{y\}$ (\emph{loc. cit.}), tandis que la fibre $f_1^{-1}(y)$ est isomorphe à
  $\operatorname{Proj}(S)$, où
  $S=\bigoplus_{k=0}^{\infty}\mathfrak m_y^k/\mathfrak m_y^{k+1}$
  (\textbf{II}, 3.5.3), c'est-à-dire ici à $\mathbf P^1_{k(y)}$~; il résulte de
  (14.4.1) que $f_1$ n'est pas ouvert au point générique de $f_1^{-1}(y)$. D'autre
  part, posons $Y_2=\mathbf P^1_Y$, et soit $f_2:Y_2\to Y$ le morphisme structural~;
  il résulte de (\textbf{II}, 8.4.4) que $f_2$ est plat, donc universellement ouvert
  (2.4.6)~; en outre (\textbf{II}, 3.5.3), $f_2^{-1}(y)$ est isomorphe à
  $\mathbf P^1_{k(y)}$~; il suffit alors de «~recoller~» $Y_1$ et $Y_2$ le long des
  fibres isomorphes $f_1^{-1}(y)$ et $f_2^{-1}(y)$, ce qui donne un morphisme
  $f:X\to Y$ où les composantes irréductibles $X_1$, $X_2$ de $X$ s'identifient
  canoniquement à $Y_1$ et $Y_2$ respectivement, et les restrictions de $f$ à ces
  composantes à $f_1$ et $f_2$.

  Rappelons toutefois (12.1.1.5) que si $Y$ est localement noethérien, $f:X\to Y$
  de type fini et \emph{plat}, alors toute composante irréductible de $X$ est
  \emph{équidimensionnelle} sur $Y$ (et par suite la restriction de $f$ à une telle
  composante est \emph{universellement ouverte} si tous les points de $Y$ sont
  \emph{géométriquement unibranches}).

  \item Rappelons (12.1.2, (i)) qu'il y a des morphismes $f:X\to Y$ ayant les
  propriétés suivantes~: $Y$ est noethérien (non géométriquement unibranche), $f$
  est fini et plat (et même étale (17.6.3)), mais la restriction de $f$ à une
  composante irréductible de $X$ \emph{n'est pas un morphisme ouvert} (bien que $f$
  lui-même le soit d'après (2.4.6)).

  \item Le critère de Chevalley (14.4.4) explique l'importance de la notion de
  morphisme universellement ouvert. Cette notion permet en effet, dans de nombreux
  résultats, plus ou moins classiques, de remplacer une hypothèse de \emph{normalité}
  par l'hypothèse qu'un certain morphisme est universellement ouvert~; l'énoncé plus
  général

  \oldpage[IV-3]{216}

  ainsi obtenu s'appliquera en particulier à des morphismes \emph{plats} (2.4.6),
  dont l'importance en géométrie algébrique va croissant. On peut considérer que
  les énoncés faisant intervenir l'hypothèse qu'un morphisme est universellement
  ouvert sont des généralisations communes d'énoncés faisant intervenir une
  hypothèse de normalité et d'énoncés faisant intervenir une hypothèse de platitude.
\end{enumerate}
\end{remarks}

\subsection{Morphismes universellement ouverts et quasi-sections}
\label{subsection:IV.14.5-fr}

\begin{lemma}[14.5.1]
\label{IV.14.5.1-fr}
Soient $Y$ un préschéma irréductible localement noethérien, $f:X\to Y$ un
morphisme localement de type fini, $x$ un point de $X$, tel que $f$ soit
équidimensionnel (13.3.2) au point $x$~; on pose $y=f(x)$,
$e=\dim_x(f^{-1}(y))$. Soit $X'$ une partie fermée irréductible de $X$ contenant
$x$, $n$ un entier tel que l'on ait $\operatorname{codim}(X',X)\leq n$ et
$\dim_x(X'\cap f^{-1}(y))\leq e-n$. Alors on a nécessairement
$\operatorname{codim}(X',X)=n$, $\dim_x(X'\cap f^{-1}(y))=e-n$, et la restriction $X'\to Y$
de $f$ est un morphisme équidimensionnel en $x$ (et a fortiori dominant).
\end{lemma}

La question étant locale sur $X$, on peut supposer que $f$ est de type fini et
\emph{équidimensionnel}, de sorte que pour \emph{tout} $z\in Y$, \emph{toutes} les
composantes irréductibles de $f^{-1}(z)$ sont de dimension $e$ (13.3.1, $a''$).
Soient $x'$ le point générique de $X'$, $y'=f(x')$,
$Y'=\overline{\{y'\}}=\overline{f(X')}$ et posons $Z=f^{-1}(Y')$. En vertu de
($\mathbf 0$, 14.2.2), on a
\begin{equation}
  \operatorname{codim}(X',Z)\leq n.
  \tag{14.5.1.1}\label{IV.14.5.1.1-fr}
\end{equation}
Et si les deux membres sont égaux, on a nécessairement $\operatorname{codim}(X',X)=n$ et $Z$
contient une composante irréductible de $X$, ce qui entraîne (13.3.1) que
$f(X')$ est dense dans $Y$ et par suite $Z=X$~; donc l'égalité
$\operatorname{codim}(X',Z)=n$ équivaut à la conjonction de l'égalité $\operatorname{codim}(X',X)=n$ et de
la relation $Z=X$.

D'autre part, en raisonnant dans les préschémas réduits de $Y$ et $X$ ayant $Y'$
et $Z$ respectivement pour espaces sous-jacents, on déduit de (5.1.2) et de
(\textbf{I}, 3.6.5) que l'on a
\begin{equation}
  \operatorname{codim}(X'\cap f^{-1}(y'),f^{-1}(y'))=\operatorname{codim}(X',Z).
  \tag{14.5.1.2}\label{IV.14.5.1.2-fr}
\end{equation}
En vertu de l'hypothèse, $f^{-1}(y')$ est biéquidimensionnel (5.2.1) et de
dimension $e$, donc ($\mathbf 0$, 14.3.5), on a, en vertu de (14.5.1.2) et
(14.5.1.1)
\begin{equation}
  e'=\dim(X'\cap f^{-1}(y'))=e-\operatorname{codim}(X',Z)\geq e-n,
  \tag{14.5.1.3}\label{IV.14.5.1.3-fr}
\end{equation}
l'égalité ayant lieu si et seulement si $\operatorname{codim}(X',X)=n$ et si $f(X')$ est dense
dans $Y$. Enfin, (13.1.1), on a $\dim_x(X'\cap f^{-1}(y))\geq e'$, d'où, par
(14.5.1.3),
\begin{equation}
  \dim_x(X'\cap f^{-1}(y))\geq e'\geq e-n.
  \tag{14.5.1.4}\label{IV.14.5.1.4-fr}
\end{equation}
Or, par hypothèse, on a aussi $\dim_x(X'\cap f^{-1}(y))\leq e-n$, d'où les
conclusions de la proposition.

\begin{corollary}[14.5.2]
\label{IV.14.5.2-fr}
Les hypothèses sur $f$, $X$, $Y$, $x$ étant celles de (14.5.1), supposons en
outre que $x$ ne soit pas point maximal de $f^{-1}(y)$. Alors il existe un
voisinage ouvert affine $U$ de $x$ dans $X$, et une section
$g\in\Gamma(U,\mathcal O_X)$ telle que l'ensemble $X'$ des $x'\in U$ tels que
$g(x')=0$ contienne $x$ et ne contienne aucun point maximal de $f^{-1}(y)$. Pour
tout $g$ ayant ces propriétés, $X'$ est équidimensionnel sur $Y$ au point $x$, et
l'on a
\begin{equation}
  \dim_x(X'\cap f^{-1}(y))=e-1 \qquad\text{et}\qquad \operatorname{codim}(X',X)=1.
  \tag{14.5.2.1}\label{IV.14.5.2.1-fr}
\end{equation}
\end{corollary}

\oldpage[IV-3]{217}

On peut se borner au cas où $X=U$ est un voisinage ouvert affine de $x$ tel que
toutes les composantes irréductibles de $f^{-1}(y)$ contiennent $x$. Ces
composantes correspondent aux idéaux premiers minimaux de
$\mathcal O_x/\mathfrak m_y\mathcal O_x$, et par hypothèse ces idéaux sont
distincts de $\mathfrak m_x/\mathfrak m_y\mathcal O_x$ (Bourbaki,
\emph{Alg. comm.}, chap.~II, \S~1, n\textsuperscript{o}~1, prop.~2)~; pour obtenir
un $g\in\Gamma(U,\mathcal O_X)$ vérifiant les conditions de l'énoncé, il suffit de
prendre $g\in\mathfrak j_x$ tel que l'image de $g$ dans $\mathfrak m_x$
n'appartienne à aucun des idéaux premiers précédents. D'ailleurs, on a
$\operatorname{codim}(X',X)\leq 1$ (5.1.8), et comme $X'$ ne contient aucune des composantes
irréductibles de $f^{-1}(y)$ et que celles-ci sont de dimension $e$, on a
($\mathbf 0$, 14.2.2.2) $\dim_x(X'\cap f^{-1}(y))\leq e-1$. Il suffit alors
d'appliquer (14.5.1).

\begin{proposition}[14.5.3]
\label{IV.14.5.3-fr}
Soient $Y$ un préschéma irréductible localement noethérien, $f:X\to Y$ un
morphisme localement de type fini, $x$ un point de $X$. On suppose que $f$ est
équidimensionnel au point $x$ et que $x$ est fermé dans $f^{-1}(f(x))$. Alors il
existe une partie irréductible $X'$ de $X$, localement fermée dans $X$, contenant
$x$ et telle que la restriction $X'\to Y$ de $f$ (où $X'$ est le sous-préschéma
réduit de $X$ ayant $X'$ pour espace sous-jacent) soit un morphisme quasi-fini et
dominant.
\end{proposition}

En effet, avec les notations de (14.5.2), l'hypothèse que $x$ est fermé dans
$f^{-1}(y)$ entraîne que $x$ n'est pas point maximal de $X'\cap f^{-1}(y)$ tant
que $e-1\geq 1$. Il suffit donc d'appliquer (14.5.2) en raisonnant par récurrence
descendante sur $e=\dim_x(f^{-1}(f(x)))$ jusqu'à ce que l'on arrive à $e=1$~;
l'application de (14.5.2) dans ce dernier cas donne un $X'$ tel que
$X'\cap f^{-1}(f(x))$ soit noethérien et de dimension 0, donc fini et discret~;
comme alors $X'$ est équidimensionnel sur $Y$, $X'\cap f^{-1}(f(x'))$ est de
dimension 0 pour tout $x'\in X'$, ce qui entraîne que la restriction $X'\to Y$ de
$f$ est un morphisme quasi-fini (\textbf{II}, 6.2.2).

\begin{corollary}[14.5.4]
\label{IV.14.5.4-fr}
Sous les hypothèses de (14.5.3), supposons en outre que $Y=\operatorname{Spec}(A)$, où $A$ est
un anneau local noethérien intègre et complet, et que $y=f(x)$ soit l'unique point
fermé de $Y$. Alors, il existe un anneau local intègre $A'$, contenant $A$, qui est
une $A$-algèbre finie et possède la propriété suivante~: si l'on pose
$Y'=\operatorname{Spec}(A')$ et $X'=X\times_Y Y'$, il existe une $Y'$-section $h:Y'\to X'$ telle
que le morphisme composé $Y'\xrightarrow{h}X'\to X$ soit une immersion dont
l'image contient $x$.
\end{corollary}

Remplaçant au besoin $X$ par un sous-préschéma réduit irréductible de $X$, on peut,
en vertu de (14.5.3), se borner au cas où le morphisme $f$ est déjà quasi-fini et
dominant. Utilisant (\textbf{II}, 6.2.5), on en déduit que $\mathcal O_x=A'$ est un
anneau intègre qui est une $A$-algèbre finie, et que $X$ est somme disjointe du
sous-préschéma fermé $Y'=\operatorname{Spec}(A')$ et d'un sous-préschéma $Y''$~; le schéma $Y'$
répond à la question, le morphisme composé $Y'\to X'\to X$ n'étant autre que le
morphisme canonique $\operatorname{Spec}(\mathcal O_x)\to X$.

\begin{remark}[14.5.5]
\label{IV.14.5.5-fr}
Si on n'exige pas que dans l'énoncé de (14.5.4), le morphisme $Y'\to X$ soit une
immersion, on peut supposer que $A'$ est en outre \emph{intégralement clos}~: il
suffit en effet de remplacer $A'$ par sa \emph{clôture intégrale} $A_1$, car on
sait ($\mathbf 0$, 23.1.5) que $A_1$ est un $A$-module de type fini.
\end{remark}

\begin{proposition}[14.5.6]
\label{IV.14.5.6-fr}
Soient $Y$ un préschéma localement noethérien, irréductible, régulier et de
dimension 1, $f:X\to Y$ un morphisme localement de type fini, $y$ un point de $Y$.
Les conditions suivantes sont équivalentes~:
\begin{enumerate}[label=\emph{\alph*)}]
  \item $f_{\mathrm{red}}$ est plat en tout point de $f^{-1}(y)$.

  \oldpage[IV-3]{218}

  \item $f$ est universellement ouvert dans un voisinage de $f^{-1}(y)$.
  \item $f$ est ouvert dans un voisinage de $f^{-1}(y)$.
  \item Toute composante irréductible de $X$ rencontrant $f^{-1}(y)$ domine $Y$.
  \item Pour tout point $x\in f^{-1}(y)$, fermé dans $f^{-1}(y)$, et toute
  composante irréductible $X_i$ de $X$ contenant $x$, il existe une partie
  irréductible $X'$ de $X$, localement fermée dans $X$, contenant $x$, contenue
  dans $X_i$, et telle que la restriction $X'\to Y$ de $f$ soit un morphisme
  quasi-fini et dominant.
\end{enumerate}
\end{proposition}

L'équivalence de $a)$, $b)$, $c)$ et $d)$ a déjà été démontrée (14.3.8). Il est
clair que $e)$ entraîne $d)$, car pour toute composante irréductible $X_i$ de $X$
rencontrant $f^{-1}(y)$, il existe dans $X_i\cap f^{-1}(y)$ un point fermé dans cet
espace (5.1.11). Enfin, pour prouver que $c)$ entraîne $e)$, on peut se borner au
cas où $f$ est un morphisme de type fini~; considérons un point fermé $x$ de
$f^{-1}(y)$ et soit $X_i$ une composante irréductible de $X$ contenant $x$. Comme
la restriction $f_i:X_i\to Y$ de $f$ à $X_i$ est un morphisme dominant, il résulte
de l'équivalence de $c)$ et $d)$ pour $f_i$ que ce morphisme est ouvert aux points
génériques de $X_i\cap f^{-1}(y)$. Il résulte donc de (14.2.2) que $f_i$ est
équidimensionnel au point $x$, et on conclut alors à l'aide de (14.5.3).

\begin{remark}[14.5.7]
\label{IV.14.5.7-fr}
Si, dans l'énoncé de (14.5.5), on suppose que $Y=\operatorname{Spec}(A)$, où $A$ est un anneau
de valuation discrète complet, et que $y$ est le point fermé de $Y$, on peut en
outre supposer que $X'=\operatorname{Spec}(A')$, où $A'$ est un anneau de valuation discrète qui
est une $A$-algèbre finie, comme le montre la démonstration de (14.5.4) et le fait
qu'un anneau local intègre régulier et de dimension 1 est un anneau de valuation
discrète (\textbf{II}, 7.1.6).
\end{remark}

\begin{proposition}[14.5.8]
\label{IV.14.5.8-fr}
Soient $Y$ un préschéma localement noethérien, $f:X\to Y$ un morphisme localement
de type fini, $y$ un point de $Y$.

Pour tout $Y$-préschéma $Y'$, posons $X(Y')=\operatorname{Hom}_Y(Y',X)$. Soit $x$ un point de
$f^{-1}(y)$~; afin que $f$ soit universellement ouvert au point $x$, il faut et il
suffit que la condition suivante soit vérifiée~:

Pour tout anneau de valuation discrète complet $A$, de corps résiduel
algébriquement clos, tout morphisme $g:Y'=\operatorname{Spec}(A)\to Y$ tel que l'image par $g$ du
point fermé $y'$ de $Y'$ soit égal à $y$, et tout élément
$u_0\in X(\operatorname{Spec}(k(y')))$ tel que $u_0(y')=x$, il existe un anneau de valuation
discrète $B$, un homomorphisme local $A\to B$ faisant de $B$ une $A$-algèbre finie,
et, si l'on pose $Z'=\operatorname{Spec}(B)$, un élément $u\in X(Z')$ tels que, si $z'$ est le
point fermé de $Z'$, le diagramme
\[
  \xymatrix{
    \operatorname{Spec}(k(z')) \ar[r] \ar[d]^{\ell} & Z' \ar[d]^{u} \\
    \operatorname{Spec}(k(y')) \ar[r]_{u_0} & X
  }
\]
soit commutatif.
\end{proposition}

\oldpage[IV-3]{219}

Notons que si $A$ et $B$ vérifient les conditions de l'énoncé, $B$ est un anneau
de valuation discrète complet (Bourbaki, \emph{Alg. comm.}, chap.~III, \S~3,
n\textsuperscript{o}~3, prop.~7 et chap.~IV, \S~2, n\textsuperscript{o}~2,
cor.~3 de la prop.~9) de corps résiduel isomorphe à celui de $A$, donc
algébriquement clos, et puisque $u(z')=x$, il y a dans $X''=X\times_Y Z'$ un point
$x''$, dont les projections dans $X$ et $Z'$ sont $x$ et $z'$ et qui est rationnel
sur $k(z')$~; en outre, puisqu'il existe une $Z'$-section $v$ de $X''$ telle que
$v(z')=x''$, l'image par $v$ du point générique $s$ de $Z'$ est une générisation
$t$ de $x''$ dont la projection dans $Z'$ est $s$~; en appliquant (14.3.6), on voit
que la condition de l'énoncé est \emph{suffisante}. Prouvons maintenant qu'elle est
\emph{nécessaire}. Posons $X'=X\times_Y Y'$, $f'=f_{(Y')}:X'\to Y'$. Il y a par
hypothèse un point $x'\in X'$ au-dessus de $x$ et de $y'$ et rationnel sur $k(y')$
(\textbf{I}, 3.3.14), donc fermé dans $f'^{-1}(y')$. En vertu de (14.3.6), il y a
une composante irréductible $T'$ de $X'$ contenant $x'$ et dominant $Y'$, donc
((14.3.8) et (14.3.13)) la restriction $T'\to Y'$ de $f'$ est équidimensionnelle au
point $x$. On déduit donc de (14.5.5) qu'il y a une $A$-algèbre finie $B$ qui est
un anneau local intègre et intégralement clos et domine $A$ (donc est un anneau de
valuation discrète) et, en posant $Z'=\operatorname{Spec}(B)$ et
$X''=X\times_Y Z'=X'\times_{Y'}Z'$, une $Z'$-section $v:Z'\to X''$ telle que
l'image du point fermé $z'$ de $Z'$ par le morphisme composé $Z'\to X''\to X'$ soit
$x'$~; le morphisme composé $u:Z'\to X''\to X$ répond donc à la question.

\begin{proposition}[14.5.9]
\label{IV.14.5.9-fr}
Soient $Y$ un préschéma localement noethérien irréductible, $f:X\to Y$ un
morphisme localement de type fini, $y$ un point géométriquement unibranche de $Y$.
Alors les conditions suivantes sont équivalentes~:
\begin{enumerate}[label=\emph{\alph*)}]
  \item $f$ est universellement ouvert en tout point de $f^{-1}(y)$.
  \item $f$ est ouvert en tout point de $f^{-1}(y)$.
  \item Pour toute composante irréductible $Z$ de $f^{-1}(y)$, de point générique
  $z$, il existe une composante irréductible $X_i$ de $X$ contenant $z$ et
  équidimensionnelle sur $Y$ au point $z$.
  \item Pour tout point fermé $x$ de $f^{-1}(y)$, il existe une partie irréductible
  $X'$ de $X$, localement fermée dans $X$, contenant $x$, et telle que la
  restriction $X'\to Y$ de $f$ soit un morphisme quasi-fini et dominant.
\end{enumerate}
\end{proposition}

L'équivalence de $a)$, $b)$ et $c)$ a déjà été démontrée (14.4.2). Pour prouver
que $a)$ entraîne $d)$, notons qu'en vertu de (14.4.2), $a)$ entraîne qu'il existe
une composante irréductible $Z$ de $X$ contenant $x$ et équidimensionnelle sur $Y$
au point $x$~; l'existence de $X'$ provient alors de (14.5.3) appliqué à la
restriction $Z\to Y$ de $f$. Inversement, supposons $d)$ vérifiée~; en vertu du
critère de Chevalley (14.4.4), la restriction $X'\to Y$ de $f$ est un morphisme
universellement ouvert au point $x$, et a fortiori $f$ est ouvert au point $x$~;
$f$ est donc ouvert en tous les points fermés de $X_y=f^{-1}(y)$. Mais $X_y$ est
un $k(y)$-préschéma localement de type fini, donc un préschéma de Jacobson~;
l'ensemble des points fermés de $X_y$ est donc \emph{dense} dans $X_y$ (10.3.1),
et il résulte de (14.1.4) que $f$ est ouvert en \emph{tous} les points de $X_y$, ce
qui achève de prouver que $d)$ entraîne $b)$.

Le résultat suivant a été mis en évidence par D.~Mumford~:

\begin{proposition}[14.5.10]
\label{IV.14.5.10-fr}
Soient $Y$ un préschéma noethérien, $f:X\to Y$ un morphisme universellement ouvert,
surjectif et localement de type fini. Alors il existe un morphisme fini surjectif
$g:Y'\to Y$ tel que, si l'on pose $X'=X\times_Y Y'$ et
$f'=f_{(Y')}:X'\to Y'$, tout point $y'\in Y'$ admette un voisinage ouvert $U'$ tel
qu'il existe une $U'$-section de $f'^{-1}(U')$.
\end{proposition}

Nous démontrerons la proposition en plusieurs étapes.

\oldpage[IV-3]{220}

\emph{I) Réduction au cas où $Y$ est intègre.} -- Si l'on a prouvé la proposition
pour chacun des sous-préschémas réduits $Y_i$ ayant pour espace sous-jacent une
composante irréductible de $Y$, et pour les images réciproques $f^{-1}(Y_i)$, il est
clair que le préschéma $Y'$ \emph{somme} des $Y'_i$ correspondants répondra à la
question. On peut donc supposer $Y$ \emph{intègre} et on va dans ce qui suit se
borner à ce cas. Alors, dans la conclusion, on pourra aussi prendre $Y'$ intègre
(quitte à le remplacer par une composante irréductible convenable).

\emph{II) Caractère local sur $Y$.} -- Nous allons montrer que si l'on peut
recouvrir $Y$ par des ouverts $U_j$ en nombre fini, tels que, pour tout $j$, la
conclusion de la proposition soit vraie pour le morphisme $f^{-1}(U_j)\to U_j$,
restriction de $f$, alors la conclusion est vraie également pour $f$. En effet, on
peut évidemment supposer les $U_j$ affines, de sorte que $U_j=\operatorname{Spec}(A_j)$, où
$A_j$ est un anneau intègre noethérien dont le corps des fractions $K=R(Y)$ est le
corps des fonctions rationnelles sur $Y$. Pour tout $j$, il y a par hypothèse une
$A_j$-algèbre finie intègre $A'_j$ telle que l'homomorphisme $A_j\to A'_j$ soit
injectif (\textbf{I}, 1.2.7) et que le morphisme correspondant
$g_j:U'_j=\operatorname{Spec}(A'_j)\to\operatorname{Spec}(A_j)=U_j$ vérifie les conditions de la proposition
(pour $U_j$ et $f^{-1}(U_j)$). Soit alors $K'$ une extension finie de $K$ contenant
les corps des fractions de tous les $A'_j$ (qui sont des extensions finies de $K$).
Considérons le normalisé $Y''$ de $Y$ dans $K'$ (\textbf{II}, 6.3.8), qui est de la
forme $\operatorname{Spec}(\mathcal B)$, où $\mathcal B$ est une $\mathcal O_Y$-Algèbre entière
quasi-cohérente, fermeture intégrale de $\mathcal O_Y$ dans $K'$ (\textbf{II},
6.3.4). Ces définitions prouvent que pour tout $j$,
\[
  A'_j=\Gamma(U'_j,\mathcal O_{U'_j})
      =\Gamma(U_j,(g_j)_*(\mathcal O_{U'_j}))
\]
s'identifie à une sous-$A_j$-algèbre finie de $\Gamma(U_j,\mathcal B)$~; autrement
dit, $(g_j)_*(\mathcal O_{U'_j})=\mathcal A'_j$ est une
$\mathcal O_{U_j}$-Algèbre cohérente, sous-Algèbre de $\mathcal B|U_j$. Il existe
donc un sous-$\mathcal O_Y$-Module $\mathcal C'_j$ cohérent de $\mathcal B$ tel que
$\mathcal C'_j|U_j=\mathcal A'_j$ (\textbf{I}, 9.4.7). Si l'on pose
$\mathcal C'=\sum_j\mathcal C'_j$, la sous-$\mathcal O_Y$-Algèbre $\mathcal B'$ de
$\mathcal B$ engendrée par $\mathcal C'$ est cohérente puisque $\mathcal B$ est une
$\mathcal O_Y$-Algèbre entière. Montrons alors que $Y'=\operatorname{Spec}(\mathcal B')$ répond à
la question. En effet, il est clair que le morphisme $g:Y'\to Y$ est fini surjectif
et que $Y'$ est intègre~; en outre, pour tout $j$, $\Gamma(U_j,\mathcal B')$ est une
$A'_j$-algèbre finie, autrement dit le morphisme $g^{-1}(U_j)\to U_j$, restriction
de $g$, se factorise en $g^{-1}(U_j)\to U'_j\to U_j$, et comme l'existence locale
de sections est stable par changement de base, cela établit notre assertion, tout
$y\in Y$ appartenant à un $U_j$.

\emph{III) Réduction au cas où $Y$ est intègre, local et géométriquement
unibranche.} -- Supposons d'abord que la proposition ait été prouvée lorsque $Y$
est intègre et \emph{local} (avec $Y'$ intègre), et montrons qu'elle est valable
lorsque $Y$ est intègre (noethérien) affine quelconque. En effet, en vertu de la
réduction II), il suffit de prouver que pour tout point $y\in Y$, la proposition est
vraie pour un voisinage ouvert affine $V$ de $y$ dans $Y$. Soient $Y=\operatorname{Spec}(A)$,
$Y_1=\operatorname{Spec}(A_{\mathfrak p})$, où $\mathfrak p=\mathfrak j_y$, et posons
$X_1=X\times_Y Y_1$~; par hypothèse, il existe un morphisme fini surjectif
$g_1:Y'_1\to Y_1$, où $Y'_1=\operatorname{Spec}(B_1)$, $B_1$ étant une
$A_{\mathfrak p}$-algèbre intègre finie, donc un anneau semi-local, tel que $g_1$
vérifie les conditions de l'énoncé pour $Y_1$ et $X_1$. Si $y'_j$
($1\leq j\leq r$) sont les points fermés de $Y'_1$, il y a donc un recouvrement de
$Y'_1$ par des ouverts $U'_j$ tels que $y'_j\in U'_j$ et qu'il existe une
$U'_j$-section $h'_j$ de $X\times_Y U'_j$ ($1\leq j\leq r$). Le
$A_{\mathfrak p}$-module $B_1$ admet un système fini de générateurs de la forme
$z_i/s$ (avec $s\in A-\mathfrak p$,

\oldpage[IV-3]{221}

$z_i$ entiers sur $A$), que l'on peut supposer (en multipliant au besoin $s$ par un
élément de $A$) être des éléments du corps des fractions de $B_1$, entiers sur
$A_s$, de sorte que si $V$ est l'ouvert affine $D(s)=\operatorname{Spec}(A_s)\subset Y$, $Y'_1$
s'identifie à $Y_1\times_V V'$, où $V'$ est le spectre de la $A_s$-algèbre finie
engendrée par les $z_i/s$~; $g:V'\to V$ est donc un morphisme fini surjectif et
$g_1=g_{(Y_1)}$. De plus, appliquant la méthode de (8.1.2, $a)$), on peut supposer
que chacun des $U'_j$ est l'image réciproque par $p:Y'_1\to V'$ d'un ouvert $W'_j$
de $V'$, tel que les $W'_j$ recouvrent $V'$ (8.3.11), et que chacune des sections
$h'_j$ est de la forme $(v'_j)_{(Y_1)}$, où $v'_j$ est une $W'_j$-section de
$X\times_Y W'_j$ (8.8.2, (i)). On est donc bien ramené à prouver la proposition
lorsque $Y=\operatorname{Spec}(A)$, $A$ étant un anneau local intègre et noethérien. On sait alors
qu'il existe une $A$-algèbre intègre finie $B$, ayant même corps des fractions que
$A$, telle que $A\subset B$ et que $\operatorname{Spec}(B)$ soit géométriquement unibranche
(($\mathbf 0$, 23.2.5) et (6.15.5))~; comme le morphisme
$\operatorname{Spec}(B)\to\operatorname{Spec}(A)$ est surjectif, on peut évidemment remplacer $Y$ par
$\operatorname{Spec}(B)$ et $X$ par $X\otimes_A B$ pour prouver la proposition. Raisonnant comme
au début de la réduction III), on peut donc supposer $A$ local, intègre et
$Y=\operatorname{Spec}(A)$ géométriquement unibranche.

\emph{IV) Réduction au cas où $X$ est intègre, affine, et $f$ quasi-fini,
surjectif, birationnel et universellement ouvert.} -- Supposons donc $Y=\operatorname{Spec}(A)$
intègre, local et géométriquement unibranche. Il existe alors un sous-préschéma
irréductible $X_0$ de $X$ tel que la restriction $f_0:X_0\to Y$ de $f$ soit un
morphisme quasi-fini et dominant et que $f_0(X_0)$ contienne le point fermé $y$ de
$Y$ (14.5.9)~; comme $Y$ est géométriquement unibranche, il résulte de (14.4.1) que
$f_0$ est encore \emph{universellement ouvert}. Comme en outre on peut supposer
$X_0$ réduit, donc intègre, on voit qu'on peut, en remplaçant $X$ par $X_0$,
supposer que $X$ est \emph{intègre} et $f$ \emph{quasi-fini et dominant}, et tel que
$f(X)$ contienne $y$~; comme $f$ est \emph{ouvert} et que tout ouvert de $Y$
contenant $y$ est égal à $Y$, $f$ est surjectif.

Soient $\xi$, $\eta$ les points génériques de $X$ et $Y$ respectivement~;
$k(\xi)$ est donc une extension \emph{finie} de $K=k(\eta)$. Par suite (4.6.8), il
y a une extension \emph{finie} $K'$ de $K$ telle que
$(\operatorname{Spec}(k(\xi)\otimes_K K'))_{\mathrm{red}}$ soit géométriquement réduit sur $K$ et
que ses composantes irréductibles soient géométriquement irréductibles~; il en
résulte ((4.5.9) et (4.6.1)) que les corps résiduels de
$\operatorname{Spec}(k(\xi)\otimes_K K')$ sont des extensions finies, primaires et séparables de
$K'$, donc sont \emph{égaux à} $K'$. Appliquant de nouveau ($\mathbf 0$, 23.2.5) et
(6.15.5), il existe une $A$-algèbre intègre \emph{finie} $B$, ayant $K'$ pour
corps des fractions, contenant $A$ et telle que $\operatorname{Spec}(B)$ soit géométriquement
unibranche~; on peut donc de nouveau remplacer $Y$ par $\operatorname{Spec}(B)$ et $X$ par
$(X\otimes_A B)_{\mathrm{red}}$ pour prouver la proposition~; la réduction III)
permet alors de supposer encore $A$ local, intègre et $Y=\operatorname{Spec}(A)$ géométriquement
unibranche, de point fermé $y$. En outre, $f$ est alors un morphisme \emph{quasi-fini,
surjectif et universellement ouvert}~; chacune des composantes irréductibles $X_i$
de $X$ ($1\leq i\leq m$) domine $Y$ (1.10.4) et est \emph{birationnelle} sur $Y$.
Soit alors $x$ un point de $f^{-1}(y)$ et soit $X_i$ une composante irréductible de
$X$ contenant $x$~; désignons encore par $X_i$ le sous-préschéma réduit (donc
intègre) de $X$ ayant $X_i$ pour espace sous-jacent, et soit $f_i:X_i\to Y$ la
restriction de $f$. Appliquant de nouveau (14.4.1) au morphisme quasi-fini et
dominant $f_i$, on voit que $f_i$ est universellement ouvert~; si $U$ est un ouvert
affine de $X_i$ contenant $x$, $f_i(U)$ est donc

\oldpage[IV-3]{222}

ouvert dans $Y$ et contient $y$, donc est égal à $Y$. On peut ainsi remplacer $X$
par $U$ pour prouver la proposition.

\emph{V) Fin de la démonstration.} -- Nous supposons donc $Y$ local, intègre et
géométriquement unibranche, $X$ intègre et affine, $f$ quasi-fini, surjectif,
birationnel et universellement ouvert~; il en résulte que $f$ est automatiquement
\emph{séparé}. En vertu du Main theorem (8.12.6), il existe une factorisation
$X\xrightarrow{j}Z\xrightarrow{u}Y$, où $j$ est une immersion ouverte et $u$ un
morphisme \emph{fini}. Remplaçant en outre $Z$ par l'image fermée de $j$
(\textbf{I}, 9.5.10), on peut supposer que $Z$ est \emph{intègre}~; comme $u$ est
fini et birationnel et $Y$ géométriquement unibranche, il résulte de
(\textbf{III}, 4.3.5 et 4.3.4) que $u$ (et par suite $f$) est un morphisme
\emph{radiciel}~; comme $f$ est universellement ouvert et surjectif, c'est donc un
\emph{homéomorphisme universel}. Par suite (2.4.5, (ii)) $f$ est un morphisme
\emph{fini}. Mais alors on répond aux conditions de l'énoncé avec $Y'$ \emph{intègre}
en prenant simplement $Y'=X$ et $g=f$, la $Y'$-section étant le morphisme diagonal
$\Delta_f$. C.Q.F.D.

Ce résultat peut être utilisé pour développer des critères de «~descente~» de
diverses propriétés par les morphismes universellement ouverts et surjectifs.
Signalons en particulier le critère suivant, dû à D.~Mumford~:

\begin{corollary}[14.5.11]
\label{IV.14.5.11-fr}
Soient $Y$ un préschéma localement noethérien, $f:X\to Y$ un morphisme localement
de type fini, $g:Y'\to Y$ un morphisme localement de type fini, universellement
ouvert et surjectif~; posons $X'=X\times_Y Y'$, $f'=f_{(Y')}:X'\to Y'$. Alors,
pour que $f$ soit affine, il faut et il suffit que $f'$ le soit.
\end{corollary}

On doit seulement prouver que la condition est suffisante. La question étant locale
sur $Y$, on peut supposer $Y$ affine, donc noethérien. En vertu de (14.5.10), il
existe un morphisme fini surjectif $h:Y_1\to Y$ tel que, si l'on pose
$Y'_1=Y'\times_Y Y_1$, et $g'=g_{(Y_1)}:Y'_1\to Y_1$, tout point de $Y_1$ admette
un voisinage ouvert $U_1$ tel qu'il existe une $U_1$-section de $g'^{-1}(U_1)$. Si
l'on pose $X_1=X\times_Y Y_1$, la projection canonique $p:X_1\to X$ est un
morphisme fini surjectif, donc, si l'on prouve que $X_1$ est un schéma affine, il en
résultera d'abord que $X$ est quasi-compact, donc noethérien, puis que $X$ est affine
en vertu du théorème de Chevalley (\textbf{II}, 6.7.1). Il suffit par suite de
prouver que le morphisme $f_1=f_{(Y_1)}:X_1\to Y_1$ est affine. Or, si l'on pose
\[
  X'_1=X_1\times_{Y_1}Y'_1=X'\times_{Y'}Y'_1
  \qquad\text{et}\qquad
  f'_1=(f_1)_{(Y'_1)}=f'_{(Y'_1)}:X'_1\to Y'_1,
\]
$f'_1$ est affine en vertu de l'hypothèse. On est donc ramené à prouver le corollaire
en y remplaçant $Y$, $X$, $f$ et $Y'$ par $Y_1$, $X_1$, $f_1$ et $Y'_1$,
autrement dit, il suffit de prouver que $f$ est affine lorsqu'on fait de plus, dans
l'énoncé de (14.5.11), l'hypothèse que tout point de $Y$ admet un voisinage ouvert
$U$ tel qu'il existe une $U$-section de $g^{-1}(U)$. La question étant locale sur
$Y$, on peut même supposer qu'il existe une $Y$-section $s$ de $Y'$. Or, on a le
lemme élémentaire suivant (valable dans toute catégorie admettant des produits
fibrés)~:

\begin{lemma}[14.5.11.1]
\label{IV.14.5.11.1-fr}
Soient $f:X\to S$, $g:Y\to S$ deux morphismes, $p_1:X\times_S Y\to X$,
$p_2:X\times_S Y\to Y$ les projections canoniques. Si $s:S\to Y$ est une section
de $g$, alors $s'=(1,s\circ f)_S$ est une section de $p_1$ et $X$, muni des
morphismes $f:X\to S$ et $s':X\to X\times_S Y$, s'identifie au produit des
$Y$-préschémas $S$ et $X\times_S Y$ pour les morphismes $s:S\to Y$ et
$p_2:X\times_S Y\to Y$.
\end{lemma}

\oldpage[IV-3]{223}

\begin{proof}
C'est un cas particulier de (\textbf{I}, 3.3.11), où l'on remplace le diagramme par
\[
  \xymatrix{
    X \ar[r]^{p_1} \ar[d]_{f} & X\times_S Y \ar[r]^{s'} \ar[d]^{p_2} & X \ar[d]^{f} \\
    S \ar[r]_{s} & Y \ar[r]_{g} & S
  }
\]
\end{proof}

Appliquant ce lemme en y remplaçant $S$, $Y$ par $Y$, $Y'$, on voit que l'on peut
écrire $f=(f')_{(Y)}$ pour le changement de base $s:Y\to Y'$, donc $f$ est affine
puisque $f'$ l'est. C.Q.F.D.

Une variante de ce critère est la suivante~:

\begin{corollary}[14.5.12]
\label{IV.14.5.12-fr}
Avec les notations et hypothèses générales de (14.5.11), soit $\mathcal L$ un
$\mathcal O_X$-Module inversible. Supposons qu'il existe une partie fermée $Z$ de
$X$, propre sur $Y$ (\textbf{II}, 5.4.10) et telle que le préschéma induit par $X$
sur l'ouvert $X-Z$ soit normal. Alors, pour que $\mathcal L$ soit ample relativement
à $f$, il faut et il suffit que
$\mathcal L'=\mathcal L\otimes_{\mathcal O_Y}\mathcal O_{Y'}$ soit ample
relativement à $f'$.
\end{corollary}

Gardons les notations de la démonstration de (14.5.11). Posons
$\mathcal L_1=\mathcal L\otimes_{\mathcal O_Y}\mathcal O_{Y_1}$~; il suffira de
prouver que $\mathcal L_1$ est ample relativement à $f\circ p$, en vertu de
(\textbf{III}, 2.6.2)~; mais $f\circ p=h\circ f_1$, et, compte tenu de
(\textbf{II}, 4.6.13, (v)), il suffira de prouver que $\mathcal L_1$ est ample
relativement à $f_1$. La question étant locale sur $Y_1$, on peut de nouveau
supposer que $g'$ admet une section $s$~; si
$\mathcal L'_1=\mathcal L'\otimes_{\mathcal O_{Y'}}\mathcal O_{Y'_1}$, on peut
alors écrire, en vertu du lemme (14.5.11.1),
$\mathcal L_1=\mathcal L'_1\otimes_{\mathcal O_{Y'_1}}\mathcal O_{Y_1}$ pour le
changement de base $s:Y_1\to Y'_1$~; la conclusion résulte donc de deux
applications de (\textbf{II}, 4.6.13, (iii)).

% IV3_B31_P223_END_BEFORE_SECTION_15

\section{Étude des fibres d'un morphisme universellement ouvert}
\label{section:IV.15-fr}

\subsection{Multiplicités des fibres d'un morphisme universellement ouvert}
\label{subsection:IV.15.1-fr}

\begin{proposition}[15.1.1]
\label{IV.15.1.1-fr}
\oldpage[IV-3]{223}
Soient $Y$ un préschéma irréductible localement noethérien, de point générique
$\eta$, $f:X\to Y$ un morphisme localement de type fini, $y$ un point de $Y$,
$\mathcal F$ un $\mathcal O_X$-Module cohérent~; on pose
$\mathcal F_y=\mathcal F\otimes_{\mathcal O_Y}\mathbf{k}(y)$ (faisceau de modules
sur la fibre $f^{-1}(y)$). Soit $z$ le point générique d'une composante
irréductible de $\operatorname{Supp}(\mathcal F_y)$, et soient $X_i$
($1\leqslant i\leqslant n$) ceux des sous-préschémas fermés réduits de $X$ dont
l'espace sous-jacent est une composante irréductible de
$\operatorname{Supp}(\mathcal F)$ contenant $z$, et qui sont tels que la
restriction $X_i\to Y$ de $f$ soit un morphisme universellement ouvert au point
$z$~; soit $z_i$ le point générique de $X_i$ ($1\leqslant i\leqslant n$). Si
$\lambda_x(\mathcal F_y)$ (resp. $\lambda_t(\mathcal F_\eta)$) désigne la longueur
géométrique de $\mathcal F_y$ en un point $x\in f^{-1}(y)$ (resp. celle de
$\mathcal F_\eta$ en un point $t\in f^{-1}(\eta)$) (4.7.5), on a
\[
\tag{15.1.1.1}
\label{IV.15.1.1.1-fr}
  \lambda_z(\mathcal F_y)\geqslant
  \sum_{i=1}^{n}\lambda_{z_i}(\mathcal F_\eta).
\]
Si de plus on suppose l'anneau $\mathcal O_y$ régulier, on a
\[
\tag{15.1.1.2}
\label{IV.15.1.1.2-fr}
  \operatorname{long}((\mathcal F_y)_z)\geqslant
  \sum_{i=1}^{n}\operatorname{long}((\mathcal F_\eta)_{z_i}).
\]
\end{proposition}

\begin{proof}
\oldpage[IV-3]{224}
Les fibres $f^{-1}(y')$ de $f$ (pour tout $y'\in Y$) ne changeant pas lorsque
l'on remplace $f$ par $f_{\mathrm{red}}$, on peut supposer que $X$ et $Y$ sont
réduits (2.4.3, (vi)), donc $Y$ intègre~; on posera
$K=\mathbf{k}(\eta)$. On peut en outre remplacer $f$ par la restriction
$\operatorname{Supp}(\mathcal F)\to Y$ de $f$ au sous-préschéma fermé réduit de
$X$ ayant $\operatorname{Supp}(\mathcal F)$ comme espace sous-jacent, donc se
borner au cas où $X=\operatorname{Supp}(\mathcal F)$. Enfin, si $y=\eta$, il n'y
a qu'une seule composante irréductible $X_i$ de $X$ contenant $z$
($\mathbf{0}_{\mathrm I}$, 2.1.8) et on a $z_i=z$, ce qui rend (15.1.1.1) et
(15.1.1.2) triviales (sans hypothèses sur $f$). On peut donc se borner au cas
$y\ne\eta$~; il résulte alors de l'hypothèse et de (1.10.3) que les $z_i$
appartiennent à $f^{-1}(\eta)$, les seconds membres de (15.1.1.1) et
(15.1.1.2) étant donc définis. Soit $Z_i$ le sous-préschéma fermé réduit de
$f^{-1}(\eta)$ ayant pour espace sous-jacent $X_i\cap f^{-1}(\eta)$~; on sait
(4.6.6) qu'il existe une extension finie radicielle $K'$ de $K$ telle que, pour
tout $i$, le $K'$-préschéma $(Z_i\otimes_K K')_{\mathrm{red}}$ soit
géométriquement réduit sur $K'$. En outre, le morphisme projection
$f^{-1}(\eta)\otimes_K K'\to f^{-1}(\eta)$ est fini, dominant et radiciel
(\textbf{I}, 3.5.7), donc est un homéomorphisme (2.4.5). Si $z_i'$ est l'unique
point de $f^{-1}(\eta)\otimes_K K'$ au-dessus de $z_i$, il résulte de (4.7.9) et
(4.7.5) que l'on a
\[
\tag{15.1.1.3}
\label{IV.15.1.1.3-fr}
  \lambda_{z_i}(\mathcal F_\eta)
  =\lambda_{z_i'}(\mathcal F_\eta\otimes_K K')
  =\operatorname{long}((\mathcal F_\eta\otimes_K K')_{z_i'}).
\]

Soit $V$ un anneau de valuation discrète, de corps de fractions $K'$, dominant
$\mathcal O_y$ (\textbf{II}, 7.1.7)~; posons $Y'=\operatorname{Spec}(V)$,
$X'=X\times_Y Y'$, $\mathcal F'=\mathcal F\otimes_Y Y'$,
$f'=f_{(Y')}:X'\to Y'$~; si $g:Y'\to Y$ est le morphisme structural, $\eta'$ le
point générique de $Y'$ et $y'$ son point fermé, on a $g(y')=y$, $g(\eta')=\eta$~;
le morphisme $\operatorname{Spec}(\mathbf{k}(y'))\to
\operatorname{Spec}(\mathbf{k}(y))$ étant fidèlement plat, il en est de même de
la projection $g':f'^{-1}(y')\to f^{-1}(y)$ (2.2.13), donc (2.3.4), il y a un
point générique $z'$ d'une composante irréductible de $f'^{-1}(y')$ tel que
$g'(z')=z$. Par construction, on a $\mathbf{k}(\eta')=K'$, donc
$f'^{-1}(\eta')=f^{-1}(\eta)\otimes_K K'$~; d'autre part, si l'on pose
$X_i'=X_i\times_Y Y'$, la restriction $X_i'\to Y'$ de $f'$ est un morphisme
universellement ouvert au point $z'$, donc (1.10.3) il existe un point
$z_i''\in f'^{-1}(\eta')$ qui est une générisation de $z'$, et l'on peut
évidemment supposer que $z_i''$ est point générique d'une composante
irréductible de $X_i'\cap f'^{-1}(\eta')$~; comme la projection de
$X_i'\cap f'^{-1}(\eta')$ dans $f^{-1}(\eta)$ est $Z_i$, on a $z_i''=z_i'$.
En vertu de (4.7.9) et (4.7.5), on a
\[
\tag{15.1.1.4}
\label{IV.15.1.1.4-fr}
  \lambda_z(\mathcal F_y)=\lambda_{z'}(\mathcal F'_{y'})
  \geqslant\operatorname{long}((\mathcal F'_{y'})_{z'})
\]
et il résulte donc de cette inégalité et de (15.1.1.3) qu'il suffit, pour établir
(15.1.1.1), de démontrer l'inégalité
\[
\tag{15.1.1.5}
\label{IV.15.1.1.5-fr}
  \operatorname{long}((\mathcal F'_{y'})_{z'})\geqslant
  \sum_{i=1}^{n}\operatorname{long}((\mathcal F'_{\eta'})_{z_i'}).
\]

Considérons maintenant la seconde inégalité (15.1.1.2)~; nous utiliserons le
lemme suivant~:

\begin{lemma}[15.1.1.6]
\label{IV.15.1.1.6-fr}
Soient $A$ un anneau local régulier qui n'est pas un corps, $K$ son corps des
fractions, $k$ son corps résiduel. Il existe un anneau de valuation discrète $V$
dominant $A$, ayant $K$ pour corps des fractions et dont le corps résiduel est
une extension transcendante pure de $k$.
\end{lemma}

\oldpage[IV-3]{225}
Posons $S=\operatorname{Spec}(A)$, et considérons le $S$-schéma $P$ obtenu en
faisant éclater le point fermé $s$ de $S$ (\textbf{II}, 8.1.3)~; si $\mathfrak m$
est l'idéal maximal de $A$, on a donc par définition
$P=\operatorname{Proj}(B)$, où $B$ est la $A$-algèbre graduée
$\bigoplus_{n\geqslant0}\mathfrak m^n$. Si $h:P\to S$ est le morphisme
structural, la fibre $h^{-1}(s)$ est isomorphe à
$\operatorname{Proj}(B\otimes_A k)$ (\textbf{II}, 2.8.10), et par définition
$B\otimes_A k$ est la $k$-algèbre graduée
$\bigoplus_{n\geqslant0}\mathfrak m^n/\mathfrak m^{n+1}$~; mais comme $A$ est
régulier, cette algèbre est isomorphe à une algèbre de polynômes
$k[t_1,\ldots,t_e]$ ($t_i$ indéterminées) ($\mathbf{0}_{\mathrm{IV}}$, 17.1.1)
et par suite $h^{-1}(s)$ est isomorphe à $\mathbf P_k^{e-1}$. Soit $\zeta$ le
point générique de $h^{-1}(s)$~; si $t_i'$ ($1\leqslant i\leqslant e$) sont les
éléments de $\mathfrak m$ dont les classes mod. $\mathfrak m^2$ sont les $t_i$,
$B_{(t_e')}$ est l'anneau d'un voisinage ouvert affine de $\zeta$ dans $P$, et
l'on a $\mathbf{k}(\zeta)=k(t_1,\ldots,t_{e-1})$~; d'autre part, comme $A$ est
intégralement clos, on vérifie aisément qu'il en est de même de $B$ (Bourbaki,
\emph{Alg. comm.}, chap.~V, \S~1, n\textsuperscript{o}~8, cor.~1 de la prop.~21),
donc $B_{(t_e')}$ est aussi intégralement clos, et il en est par suite de même de
l'anneau local $\mathcal O_{P,\zeta}$. Enfin, comme $A=\mathcal O_s$ est
régulier, donc un anneau universellement caténaire (5.6.4), on a, en vertu de
(5.6.5),
$\dim(\mathcal O_{P,\zeta})=\dim(\mathcal O_s)-(e-1)$, puisque $h$ est un
morphisme birationnel, que $\dim(h^{-1}(s))=e-1$ et que l'anneau local de la
fibre $h^{-1}(s)$ en son point générique $\zeta$ est de dimension $0$. Mais
$\dim(\mathcal O_s)=e$, donc $\dim(\mathcal O_{P,\zeta})=1$, et
$\mathcal O_{P,\zeta}=V$ est un anneau de valuation discrète, étant noethérien,
de dimension $1$ et intégralement clos (\textbf{II}, 7.1.6)~; il répond
évidemment à la question.

Ce lemme étant établi, on peut reprendre le raisonnement fait pour l'inégalité
(15.1.1.1) avec $Y'=\operatorname{Spec}(V)$~; cette fois $\mathbf{k}(y')$ est
une extension \emph{séparable} de $\mathbf{k}(y)$, et par suite (4.7.9), on a
$\operatorname{long}((\mathcal F_y)_z)=
\operatorname{long}((\mathcal F'_{y'})_{z'})$~; comme par ailleurs
$f'^{-1}(\eta')=f^{-1}(\eta)$ et $\mathcal F'_{\eta'}=\mathcal F_\eta$, on est
encore ramené à prouver l'inégalité (15.1.1.5). Or, si $t$ est une uniformisante
de $\mathcal O_{y'}$, $f'^{-1}(y')$ est le sous-préschéma fermé de $X'$ défini
par l'idéal $t\mathcal O_{X'}$, et on a
$\mathcal F'_{y'}=\mathcal F'/t\mathcal F'$~; d'autre part, on vérifie aussitôt
(\textbf{I}, 9.1.12) que
$(\mathcal F'_{\eta'})_{z_i'}=\mathcal F'_{z_i'}$~; l'inégalité à démontrer
(15.1.1.5) n'est autre alors qu'un cas particulier de (3.4.1.1).
\end{proof}

\begin{corollary}[15.1.2]
\label{IV.15.1.2-fr}
Les hypothèses étant celles de (15.1.1), supposons de plus que
$\lambda_z(\mathcal F_y)=1$ (resp. que $\mathcal O_y$ soit régulier et
$\operatorname{long}((\mathcal F_y)_z)=1$). Alors il existe au plus une composante
irréductible $X_i$ de $\operatorname{Supp}(\mathcal F)$ contenant $z$ et telle que
la restriction $X_i\to Y$ de $f$ soit universellement ouverte au point $z$. De
plus, si $z_i$ est le point générique de cette composante, on a
$\lambda_{z_i}(\mathcal F_\eta)=1$ (resp.
$\operatorname{long}((\mathcal F_\eta)_{z_i})=1$).
\end{corollary}

Cela résulte aussitôt de (15.1.1) en notant que l'on a nécessairement, par
définition de $X_i$, $(\mathcal F_\eta)_{z_i}\ne0$ (\textbf{I}, 9.1.13).

\begin{corollary}[15.1.3]
\label{IV.15.1.3-fr}
Soient $Y$ un préschéma irréductible localement noethérien de point générique
$\eta$, $f:X\to Y$ un morphisme localement de type fini, $\mathcal F$ un
$\mathcal O_X$-Module cohérent de support $X$, $x$ un point de $X$, $y=f(x)$. On
suppose que~: $1^\circ$ $x$ n'appartient qu'à une seule composante irréductible
$Z$ de $f^{-1}(y)$~; $2^\circ$
$\mathcal F_y=\mathcal F\otimes_{\mathcal O_Y}\mathbf{k}(y)$ est géométriquement
réduit sur $\mathbf{k}(y)$, autrement dit, si $z$ est le point générique de $Z$,
$\lambda_z(\mathcal F_y)=1$ (resp. $\mathcal O_y$ est régulier et
$\operatorname{long}((\mathcal F_y)_z)=1$). Alors, il existe au plus une
composante irréductible $X'$ de $X$ contenant $x$, telle que
$\dim_x(X'\cap f^{-1}(y))=\dim_x(f^{-1}(y))$ et que la restriction $X'\to Y$ de
$f$ soit universellement ouverte aux points génériques des composantes
irréductibles de $X'\cap f^{-1}(y)$ contenant $x$. De plus, s'il existe une telle
composante $X'$ et si $z'$ est son point générique, on a
$\lambda_{z'}(\mathcal F_\eta)=1$ (resp.
$\operatorname{long}((\mathcal F_\eta)_{z'})=1$).
\end{corollary}

\oldpage[IV-3]{226}
En effet, une composante irréductible $X'$ de $X$ contenant $x$ et telle que
$\dim_x(X'\cap f^{-1}(y))=\dim_x(f^{-1}(y))$ contient nécessairement $Z$,
puisqu'une composante irréductible de $X'\cap f^{-1}(y)$ contenant $x$ doit être
de même dimension que la seule composante irréductible $Z$ de $f^{-1}(y)$
contenant $x$. On peut alors appliquer (15.1.2) à $z$.

\begin{remarks}[15.1.4]
\label{IV.15.1.4-fr}
\begin{enumerate}[label=(\roman*)]
  \item Dans l'énoncé de (15.1.1), on ne peut remplacer «~universellement
  ouvert~» par «~équidimensionnel~», comme le montre l'exemple (14.4.10, (ii)) où
  l'on fait $\mathcal F=\mathcal O_X$~; les fibres de $f$ sont alors des schémas
  artiniens réduits, donc (avec les notations introduites \emph{loc. cit.}) on a
  $\lambda_{x'}(\mathcal F_y)=1$, mais il y a deux composantes irréductibles
  $X_1$, $X_2$ de $X$ contenant $x'$, et le second membre de (15.1.1.1) est donc
  égal à $2$, bien que la restriction de $f$ à chacune des composantes $X_1$,
  $X_2$ soit un morphisme équidimensionnel en tout point. Ce même exemple montre
  aussi que dans (15.1.2) et (15.1.3), on ne peut remplacer l'hypothèse
  «~universellement ouvert~» par «~équidimensionnel~».
  \item Il est vraisemblable que pour la validité de l'inégalité (15.1.1.2), on
  ne peut supprimer l'hypothèse que $\mathcal O_y$ est \emph{régulier}.
\end{enumerate}
\end{remarks}

\subsection{Platitude des morphismes universellement ouverts à fibres géométriquement réduites}
\label{subsection:IV.15.2-fr}

\begin{env}[15.2.1]
\label{IV.15.2.1-fr}
On a vu (2.4.6) qu'un morphisme \emph{plat} et localement de présentation finie
est \emph{universellement ouvert}. On a une réciproque partielle de ce résultat
moyennant des hypothèses supplémentaires~:
\end{env}

\begin{theorem}[15.2.2]
\label{IV.15.2.2-fr}
Soient $Y$ un préschéma localement noethérien, $f:X\to Y$ un morphisme localement
de type fini, $\mathcal F$ un $\mathcal O_X$-Module cohérent, $x$ un point de
$X$, $y=f(x)$. On suppose vérifiées les conditions suivantes~:
\begin{enumerate}[label=(\roman*)]
  \item $\operatorname{Supp}(\mathcal F)=X$.
  \item $f$ est universellement ouvert aux points génériques des composantes
  irréductibles de $f^{-1}(y)$ contenant $x$.
  \item $\mathcal F_y=\mathcal F\otimes_{\mathcal O_Y}\mathbf{k}(y)$ est
  géométriquement réduit au point $x$ (4.6.22).
  \item L'anneau $\mathcal O_y$ est réduit.
\end{enumerate}
Alors $\mathcal F$ est $f$-plat au point $x$.
\end{theorem}

\begin{proof}
Comme $\mathcal O_y$ est réduit et noethérien, nous allons appliquer le critère
valuatif de platitude (11.8.1). Considérons donc un homomorphisme local
$\mathcal O_y\to A'$, où $A'$ est un anneau de valuation discrète~; posons
$Y'=\operatorname{Spec}(A')$, $X'=X\times_Y Y'$, $f'=f_{(Y')}:X'\to Y'$, et
soit $x'$ un point de $X'$ dont les projections dans $X$ et $Y'$ sont
respectivement $x$ et le point fermé $y'$ de $Y'$~; si l'on pose
$\mathcal F'=\mathcal F\otimes_Y Y'$, on a
$\operatorname{Supp}(\mathcal F')=X'$ par (\textbf{I}, 9.1.13)~; tout point
générique $z'$ d'une composante irréductible de $f'^{-1}(y')$ contenant $x'$ a
pour projection dans $f^{-1}(y)$ un point générique $z$ d'une composante
irréductible de $f^{-1}(y)$ contenant $x$, en vertu de (4.2.6)~; donc $f'$ est
universellement ouvert au point $z'$. Enfin, il résulte de (4.7.11) que
$\mathcal F'_{y'}$ est géométriquement réduit au point $x'$. On voit donc qu'on
est ramené à démontrer le

\oldpage[IV-3]{227}
\begin{lemma}[15.2.2.1]
\label{IV.15.2.2.1-fr}
Soient $Y=\operatorname{Spec}(A)$ le spectre d'un anneau de valuation discrète,
$y$ son point fermé, $f:X\to Y$ un morphisme localement de type fini, $x$ un
point de $f^{-1}(y)$, $\mathcal F$ un $\mathcal O_X$-Module cohérent. On suppose
vérifiées les conditions suivantes~:
\begin{enumerate}[label=(\roman*)]
  \item $\operatorname{Supp}(\mathcal F)=X$.
  \item $f$ est ouvert aux points génériques des composantes irréductibles de
  $f^{-1}(y)$ contenant $x$.
  \item $\mathcal F_y=\mathcal F\otimes_{\mathcal O_y}\mathbf{k}(y)$ est réduit
  au point $x$ (3.2.2).
\end{enumerate}
Alors $\mathcal F$ est $f$-plat au point $x$.
\end{lemma}

Désignons en effet par $t$ une uniformisante de $A$, posons $B=\mathcal O_x$ et
$\mathcal F_x=M$~; la condition (i) implique que
$\operatorname{Supp}(M)=\operatorname{Spec}(B)$ (\textbf{I}, 9.1.13). La
condition (ii) implique, en vertu de (14.3.2), que toute composante irréductible
de $X$ contenant $x$ domine $Y$~; en d'autres termes, l'image réciproque dans $A$
de tout idéal premier minimal $\mathfrak p_i$ de $B$ est $0$, ce qui signifie
encore que l'on a $t\notin\mathfrak p_i$ pour tout $i$. Enfin, (iii) signifie que
$M/tM$ est un $(B/tB)$-module \emph{réduit} (3.2.2). Il s'agit de montrer que
sous ces conditions $M$ est un $A$-module \emph{sans torsion}
($\mathbf{0}_{\mathrm I}$, 6.3.4), et pour cela il suffit évidemment de montrer
que $t$ est un élément $M$-\emph{régulier}~; mais cela résulte de (3.4.6.1).
\end{proof}

\begin{corollary}[15.2.3]
\label{IV.15.2.3-fr}
Soient $Y$ un préschéma localement noethérien, $f:X\to Y$ un morphisme localement
de type fini, $x$ un point de $X$, $y=f(x)$. On suppose vérifiées les conditions
suivantes~:
\begin{enumerate}[label=(\roman*)]
  \item $f$ est universellement ouvert aux points génériques des composantes
  irréductibles de $f^{-1}(y)$ contenant $x$.
  \item $f^{-1}(y)$ est géométriquement réduit (sur $\mathbf{k}(y)$) au point
  $x$ (4.6.9).
  \item L'anneau $\mathcal O_y$ est réduit.
\end{enumerate}
Sous ces conditions $f$ est plat au point $x$.
\end{corollary}

Il suffit d'appliquer (15.2.2) à $\mathcal F=\mathcal O_X$ (4.6.22).

\begin{remarks}[15.2.4]
\label{IV.15.2.4-fr}
\begin{enumerate}[label=(\roman*)]
  \item Les trois premières des conditions de (15.2.2) ne changent pas lorsqu'on
  remplace $f$ par $f_{\mathrm{red}}$~; mais si $\mathfrak N$ est le nilradical
  d'un anneau local noethérien $A$, $A/\mathfrak N$ est plat sur $A/\mathfrak N$
  mais non sur $A$ lui-même lorsque $\mathfrak N\ne0$ (Bourbaki,
  \emph{Alg. comm.}, chap.~II, \S~3, n\textsuperscript{o}~2, cor.~2 de la
  prop.~5). On voit donc qu'on ne peut dans (15.2.2) supprimer la condition (iv).

  \item Il résulte de (2.4.6) que si la conclusion de (15.2.2) est vraie, ainsi
  que l'hypothèse (i), $f$ est universellement ouvert dans un voisinage de $x$,
  et en particulier aux points génériques des composantes irréductibles de
  $f^{-1}(y)$ contenant $x$. Même lorsque (i), (iii) et (iv) sont vérifiées, on
  ne peut remplacer (ii) par l'hypothèse plus faible que $f$ est universellement
  ouvert \emph{au point $x$ seulement}. C'est ce que montre l'exemple
  (14.1.3, (i)), où $f$ est évidemment universellement ouvert en tout point de
  $X_2$, donc en particulier au point $x$, intersection de $X_1$ et $X_2$~; les
  conditions (ii) et (iii) de (15.2.3) sont aussi vérifiées par cet exemple.
  Toutefois, $\mathcal O_x$ n'est pas un $\mathcal O_y$-module sans torsion,
  $\mathcal O_y$ s'identifiant à un sous-anneau de $\mathcal O_x$ qui contient
  des diviseurs de $0$ dans $\mathcal O_x$~; donc $f$ n'est pas plat au point $x$.

  \item Dans (15.2.3), la conclusion n'est plus nécessairement valable lorsqu'on
  remplace l'hypothèse (ii) par l'hypothèse plus faible que $f^{-1}(y)$,
  considéré comme $\mathbf{k}(y)$-préschéma, n'a pas de cycles premiers associés
  immergés. Un exemple est fourni en prenant pour $Y$ une courbe ayant un point
  de rebroussement $y$, par exemple
  $Y=\operatorname{Spec}(K[S,T]/(S^3-T^2))$, $K$ étant un corps algébriquement
  clos, et pour $X$ sa normalisée (\textbf{II}, 6.3.8). Si $s$ et $t$ sont les
  classes de $S$ et $T$ dans $A=K[S,T]/(S^3-T^2)$, on vérifie aussitôt que
  $X=\operatorname{Spec}(B)$, où $B=K[s,t,u]$, $u$ étant l'élément $t/s$ du corps
  des fractions de $A$, et comme $s=u^2$, $t=u^3$, on a aussi $B=K[u]$,
  isomorphe à l'anneau de polynômes à une indéterminée sur $K$. Le seul point $x$
  de $X$ au-dessus du point $y$ correspondant à l'idéal maximal $(s)+(t)$
  correspond à l'idéal maximal $(u)$~; mais comme la classe $\bar u$ de $u$ dans
  $\mathcal O_x/\mathfrak m_y\mathcal O_x$ est telle que $\bar u^2=0$,
  $f^{-1}(y)$ n'est pas un $\mathbf{k}(y)$-préschéma réduit. Ici $f$ est un
  morphisme fini, surjectif et radiciel, donc un homéomorphisme universel (2.4.5),
  mais n'est pas plat au point $x$, car si $\mathcal O_x$ était un
  $\mathcal O_y$-module plat, il serait un $\mathcal O_y$-module
  \oldpage[IV-3]{228}
  libre de rang $1$ engendré par l'élément $1$ de $\mathcal O_x$ (Bourbaki,
  \emph{Alg. comm.}, chap.~II, \S~3, n\textsuperscript{o}~2, prop.~5), ce qui est
  absurde.

  \item Montrons enfin que dans (15.2.2) ou (15.2.3), on ne peut remplacer
  «~géométriquement réduit~» par «~réduit~». Nous allons en effet définir deux
  anneaux $A$, $A'$ ayant les propriétés suivantes~:
  \begin{enumerate}[label=\arabic*)]
    \item $A$ est un anneau local noethérien d'idéal maximal $\mathfrak m$,
    intègre, complet, de dimension $1$ et géométriquement unibranche.
    \item $A'$ est la clôture intégrale de $A$, une $A$-algèbre finie dont
    l'idéal maximal est $\mathfrak mA'$, et le corps résiduel $k'$ une extension
    radicielle finie non triviale du corps résiduel $k=A/\mathfrak m$ de $A$.
  \end{enumerate}
  On peut alors prendre $Y=\operatorname{Spec}(A)$,
  $X=\operatorname{Spec}(A')$, $\mathcal F=\mathcal O_X$, $y$ et $x$ étant les
  points fermés de $Y$ et $X$ respectivement~; les hypothèses (i) et (iv) de
  (15.2.2) sont trivialement vérifiées~; comme $A$ est de dimension $1$,
  $f:X\to Y$ est radiciel, fini et surjectif, donc un homéomorphisme universel
  (2.4.5), ce qui prouve l'hypothèse (i) de (15.2.3). Enfin, il est clair que
  $f^{-1}(y)$ est réduit. Cependant $f$ n'est pas plat au point $x$, car si $A'$
  était un $A$-module plat, ce serait un $A$-module libre de rang $1$ (puisqu'il
  a même corps des fractions que $A$) engendré par l'élément $1$ de $A$
  (Bourbaki, \emph{Alg. comm.}, chap.~II, \S~3, n\textsuperscript{o}~2,
  prop.~5), ce qui est absurde.

  Pour construire les anneaux $A$ et $A'$, partons d'un corps imparfait $k$ de
  caractéristique $p>0$~; soient $K_0=k((T))$ le corps des séries formelles sur
  $k$, $A_0$ l'anneau de valuation pour $K_0$ correspondant à la valuation
  discrète égale à l'\emph{ordre} des séries formelles~; le corps résiduel de
  $A_0$ est donc $k$. Soit $\alpha$ un élément de $k$ qui n'est pas puissance
  $p$-ième, et prenons $k'=k(\alpha^{1/p})$~; $K=k'((T))$ est alors une extension
  radicielle finie de $K_0$, et $A'=A_0\otimes_k k'$ l'anneau de valuation
  discrète pour $K$ correspondant à l'ordre des séries formelles sur $K$. Si
  $\mathfrak m'$ est l'idéal maximal de $A'$, on répondra aux conditions 1) et 2)
  en prenant $A=A_0+\mathfrak m'$~: en effet, comme $A_0$ est noethérien et $A'$
  un $A_0$-module de type fini, $A$ est une $A_0$-algèbre finie, donc un anneau
  noethérien~; comme $A'$ est fini sur $A$, $\mathfrak m'$ est le seul idéal
  maximal de $A$, qui est donc un anneau local complet, évidemment de dimension
  $1$ (étant fini sur $A_0$) et géométriquement unibranche puisque $A'$ est
  intégralement clos et a même corps de fractions $K$ que $A$.

  \item La démonstration de (15.2.2) montre toutefois que l'on peut affaiblir la
  condition (iii) en y remplaçant «~géométriquement réduit~» par «~réduit~»
  lorsqu'on peut appliquer le critère valuatif de platitude (11.8.1) en utilisant
  seulement des anneaux de valuation discrète $A'$ dont le corps résiduel $k'$
  est séparable sur $\mathbf{k}(y)$~; en effet, si $x'$ et $z'$ ont les mêmes
  significations que dans (15.2.2), les multiplicités radicielles
  $\mu_i(\mathbf{k}(z')\mid k')$ et
  $\mu_i(\mathbf{k}(z)\mid\mathbf{k}(y))$ sont les mêmes, en vertu de (4.7.3),
  donc il résulte de (4.7.9) que $\operatorname{long}((\mathcal F_y)_z)$ et
  $\operatorname{long}((\mathcal F'_{y'})_{z'})$ sont égales, et l'on est encore
  ramené au cas où $Y$ est le spectre d'un anneau de valuation discrète et $y$
  son point fermé. Par exemple, on pourra obtenir ce résultat plus fort lorsque
  $\mathcal O_y$ est unibranche et dominé par un anneau de valuation discrète
  dont le corps résiduel est extension séparable de $\mathbf{k}(y)$, en vertu de
  (11.6.4). C'est le cas lorsque $\mathcal O_y$ est un anneau régulier, d'après
  (15.1.1.6). Mais, comme nous l'a signalé Hironaka, il y a des anneaux locaux
  intégralement clos noethériens, de dimension $2$ (provenant de schémas
  algébriques sur des corps imparfaits) qui ne satisfont pas à la condition
  précédente.
\end{enumerate}
\end{remarks}

\subsection{Applications : critères de réduction et d'irréductibilité}
\label{subsection:IV.15.3-fr}

\begin{proposition}[15.3.1]
\label{IV.15.3.1-fr}
\oldpage[IV-3]{228}
Soient $Y$ un préschéma localement noethérien, $f:X\to Y$ un morphisme localement
de type fini, $\mathcal F$ un $\mathcal O_X$-Module cohérent, $x$ un point de
$X$, $y=f(x)$. Supposons vérifiées les conditions (i), (ii) et (iii) de
(15.2.2). Alors~:
\begin{enumerate}[label=(\roman*)]
  \item Il existe un voisinage $U$ de $x$ dans $X$ tel que $f|U$ soit
  universellement ouvert et que, pour tout $x'\in U$,
  $\mathcal F_{f(x')}$ soit géométriquement réduit sur
  $\mathbf{k}(f(x'))$ au point $x'$.
  \item Si de plus $Y$ est réduit au point $y$, il existe un voisinage
  $U_1\subset U$ de $x$ tel que $\mathcal F|U_1$ soit $f$-plat et réduit.
\end{enumerate}
\end{proposition}

\begin{proof}
\oldpage[IV-3]{229}
Compte tenu de (\textbf{I}, 5.1.8), (4.2.7) et (4.7.11), les hypothèses ne
changent pas, ni la conclusion (i) de l'énoncé, si l'on remplace $Y$ par
$Y_{\mathrm{red}}$ et $X$ par $X\times_Y Y_{\mathrm{red}}$, et l'on peut donc
supposer $Y$ \emph{réduit}. Il résulte alors de (15.2.2) que $\mathcal F$ est
$f$-plat au point $x$, donc aussi (11.1.1) dans un voisinage de $x$, et
\emph{a fortiori} $f$ est universellement ouvert dans ce voisinage (2.4.6). Le
fait que $\mathcal F_{f(x')}$ soit alors géométriquement réduit au point $x'$
pour tous les points d'un voisinage de $x$ résulte de (12.1.1, (vii)).
L'assertion (ii) résulte donc de ce qui précède et de (3.3.4).
\end{proof}

\begin{proposition}[15.3.2]
\label{IV.15.3.2-fr}
Les notations étant celles de (15.3.1), supposons que les conditions (i), (ii),
(iii) de (15.2.2) soient vérifiées, et de plus que $f^{-1}(y)$ soit
équidimensionnel au point $x$. Alors $f$ est équidimensionnel au point $x$.
\end{proposition}

\begin{proof}
On a, pour tout $y'\in Y$, $\operatorname{Supp}(\mathcal F_{y'})=f^{-1}(y')$.
Par hypothèse $\mathcal F_y$ est réduit (donc vérifie $(S_1)$) et est
équidimensionnel au point $x$. En vertu de (12.1.1, (ii)), on peut donc supposer
que $\mathcal F_{f(x')}$ est équidimensionnel et de dimension constante
$e=\dim_x(f^{-1}(y))$ au point $x'$ pour tous les $x'\in U$, ce qui signifie que
$f^{-1}(f(x'))$ est équidimensionnel et de dimension $e$ au point $x'$. Comme
par ailleurs $\mathcal F|U$ est $f$-plat, il résulte de (2.3.4) et de
(13.3.1, $a$)) que $f$ est équidimensionnel au point $x$.
\end{proof}

\begin{corollary}[15.3.3]
\label{IV.15.3.3-fr}
Les notations étant celles de (15.3.1), supposons vérifiées les conditions
suivantes~:
\begin{enumerate}[label=(\roman*)]
  \item $\operatorname{Supp}(\mathcal F)=X$.
  \item $f$ est universellement ouvert aux points génériques des composantes
  irréductibles de $f^{-1}(y)$ contenant $x$.
  \item $\mathcal F_y$ est géométriquement ponctuellement intègre sur
  $\mathbf{k}(y)$ au point $x$ (4.6.22).
  \item $Y$ est géométriquement unibranche au point $y$.
\end{enumerate}
Alors $x$ n'appartient qu'à une seule composante irréductible de $X$.

Si de plus $Y$ est intègre au point $y$ et $\mathcal F=\mathcal O_X$, alors $X$
est intègre au point $x$ et $f$ est plat au point $x$.
\end{corollary}

\begin{proof}
Notons d'abord que (i) et (iii) entraînent que $x$ n'appartient qu'à une seule
composante irréductible $Z$ de $f^{-1}(y)$. Il résulte donc de (14.4.7) et de
(15.3.2) que pour toute composante irréductible $X_i$ de $X$ contenant $x$
(donc $Z$), la restriction $f_i$ de $f$ à $X_i$ est un morphisme
équidimensionnel au point $x$, et universellement ouvert au point générique de
$Z$. La première assertion de l'énoncé résulte alors de (15.1.3) et la seconde
de (15.3.1).
\end{proof}

\begin{remarks}[15.3.4]
\label{IV.15.3.4-fr}
\begin{enumerate}[label=(\roman*)]
  \item L'exemple (14.4.10, (ii)) montre que, dans (15.3.3), on ne peut
  supprimer l'hypothèse que $Y$ est géométriquement unibranche au point $y$.
  \item La remarque (15.2.4, (vi)) montre que la conclusion de (15.3.3) est
  encore vraie sous les hypothèses suivantes~: $\mathcal O_y$ est régulier,
  $\operatorname{Supp}(\mathcal F)=X$ et $\mathcal F_y$ est intègre au point
  $x$. Nous ignorons si dans cet énoncé, on peut remplacer l'hypothèse que
  $\mathcal O_y$ est régulier par celle qu'il est géométriquement unibranche, ou
  même intègre et intégralement clos.
\end{enumerate}
\end{remarks}

\subsection{Compléments sur les morphismes de Cohen-Macaulay}
\label{subsection:IV.15.4-fr}

\begin{proposition}[15.4.1]
\label{IV.15.4.1-fr}
Soient $Y$ un préschéma localement noethérien, $f:X\to Y$ un morphisme localement
de type fini, $\mathcal F$ un $\mathcal O_X$-Module cohérent tel que
$\operatorname{Supp}(\mathcal F)=X$, $x$ un point de $X$, $y=f(x)$. On pose,
pour tout $z\in X$, $X_{f(z)}=f^{-1}(f(z))$,
$\mathcal F_{f(z)}=\mathcal F\otimes_{\mathcal O_{f(z)}}\mathbf{k}(f(z))$.
Supposons que $\mathcal F$ soit $f$-plat au point $x$ et que $\mathcal F_y$ soit
un $\mathcal O_{X_y}$-Module de Cohen-Macaulay au point $x$. Alors $f$ est
\oldpage[IV-3]{230}
équidimensionnel au point $x$.

En effet ((11.1.1) et (12.1.1, (vi))), il existe un voisinage $U$ de $x$ tel que
$\mathcal F|U$ soit $f$-plat et que pour tout $x'\in U$,
$\mathcal F_{f(x')}$ soit un $\mathcal O_{X_{f(x')}}$-Module de Cohen-Macaulay
au point $x'$~; cette dernière propriété entraîne que $\mathcal F_{f(x')}$ est
équidimensionnel au point $x'$ et que $x'$ n'appartient à aucun cycle premier
immergé associé à $\mathcal F_{f(x')}$ ($\mathbf{0}_{\mathrm{IV}}$, 16.5.4). En
outre, en vertu de (12.1.1, (ii)), on peut supposer que les dimensions des
composantes irréductibles de
$\operatorname{Supp}(\mathcal F_{f(x')})|(U\cap X_{f(x')})$ aient une (même)
valeur indépendante de $x'$. Comme
$\operatorname{Supp}(\mathcal F_{f(x')})=X_{f(x')}$ par hypothèse, il résulte de
(2.3.4) et de (13.3.1, $a$)) que $f$ est équidimensionnel au point $x$.
\end{proposition}

\begin{proposition}[15.4.2]
\label{IV.15.4.2-fr}
Soient $Y$ un préschéma localement noethérien, $f:X\to Y$ un morphisme localement
de type fini, $\mathcal F$ un $\mathcal O_X$-Module cohérent tel que
$\operatorname{Supp}(\mathcal F)=X$, $x$ un point de $X$, $y=f(x)$. On suppose
que $\mathcal O_y$ est régulier et que $\mathcal F$ est un $\mathcal O_X$-Module
de Cohen-Macaulay au point $x$. Alors (avec les notations de (15.4.1)) les
conditions suivantes sont équivalentes~:
\begin{enumerate}[label=\emph{\alph*})]
  \item $\mathcal F$ est $f$-plat au point $x$ et $\mathcal F_y$ est un
  $\mathcal O_{X_y}$-Module de Cohen-Macaulay au point $x$.
  \item $\mathcal F$ est $f$-plat au point $x$.
  \item $f$ est universellement ouvert dans un voisinage de $x$.
  \item $f$ est ouvert aux points génériques des composantes irréductibles de
  $X_y$ contenant $x$.
  \item $\dim(\mathcal F_x)=\dim(\mathcal O_y)+\dim((\mathcal F_y)_x)$.
  \item[\emph{e}$'$)]
  $\dim(\mathcal O_{X,x})=\dim(\mathcal O_y)+
  \dim(\mathcal O_{X,x}\otimes_{\mathcal O_y}\mathbf{k}(y))$.
\end{enumerate}
Comme $\operatorname{Supp}(\mathcal F)=X$, les conditions $e)$ et $e')$ sont
équivalentes par définition (5.1.12). La condition $a)$ implique trivialement
$b)$~; $b)$ entraîne que $\mathcal F$ est $f$-plat aux points d'un voisinage de
$x$ (11.1.1), donc $b)$ entraîne $c)$ (2.4.6)~; $c)$ entraîne trivialement $d)$,
et $d)$ entraîne $e')$ par (14.2.1). Enfin, puisque $\mathcal O_y$ est régulier
et que $\mathcal F_x$ est un $\mathcal O_{X,x}$-module de Cohen-Macaulay, il
résulte de (6.1.4) que $e)$ entraîne $a)$.
\end{proposition}

\begin{proposition}[15.4.3]
\label{IV.15.4.3-fr}
Soient $f:X\to Y$ un morphisme localement de présentation finie, $\mathcal F$ un
$\mathcal O_X$-Module de présentation finie, tel que
$\operatorname{Supp}(\mathcal F)=X$, $x$ un point de $X$, $y=f(x)$. Alors (avec
les notations de (15.4.1)) les conditions suivantes sont équivalentes~:
\begin{enumerate}[label=\emph{\alph*})]
  \item $\mathcal F$ est $f$-plat au point $x$ et $\mathcal F_y$ est un
  $\mathcal O_{X_y}$-Module de Cohen-Macaulay au point $x$.
  \item Il existe un voisinage ouvert $U$ de $x$ dans $X$, et un $Y$-morphisme
  quasi-fini $g:U\to Y'=Y[T_1,\ldots,T_n]$ ($T_i$ indéterminées), tels que
  $\mathcal F|U$ soit $g$-plat au point $x$.
\end{enumerate}
\end{proposition}

\begin{proof}
Montrons d'abord que $b)$ entraîne $a)$. Posons $h:Y'\to Y$~; le
$\mathbf{k}(y)$-préschéma $Y_y'=h^{-1}(y)$ est régulier
($\mathbf{0}_{\mathrm{IV}}$, 17.3.7)~; soient $A$ l'anneau local de ce préschéma
au point $y'=g(x)$, $k$ son corps résiduel, et soit $B$ l'anneau local de $X_y$
au point $x$~; posons d'autre part $(\mathcal F_y)_x=M$, qui est un $B$-module
de type fini. L'hypothèse que le morphisme $g$ est quasi-fini entraîne qu'il en
est de même de $g_y:X_y\to Y_y'$ (\textbf{II}, 6.2.4), donc $M\otimes_A k$ est
un $(B\otimes_A k)$-module de longueur finie (\textbf{II}, 6.2.2)~; comme par
hypothèse $M$ est un $A$-module plat et $A$ un anneau de Cohen-Macaulay, $M$ est
un $B$-module de Cohen-Macaulay (6.3.3). Le fait que $\mathcal F$ est $f$-plat
au point $x$ résulte de ce qu'il est $g$-plat et de ce que le morphisme $h$ est
plat (2.1.6).

Montrons maintenant que $a)$ entraîne $b)$. La question étant locale sur $X$ et
$Y$, on peut supposer $X$ et $Y$ affines~; utilisant (8.9.1) et (11.2.7), on se
ramène au cas où $X$ et $Y$ sont noethériens, compte tenu de (6.7.1).
L'hypothèse entraîne que $f$ est
\oldpage[IV-3]{231}
équidimensionnel au point $x$ (15.4.1), d'où l'existence d'un morphisme
quasi-fini $g:U\to Y'$ tel que toute composante irréductible de $U$ domine $Y'$
(13.3.1, $b$)) et que $g_y:U_y\to Y_y'$ soit fini et surjectif (13.3.1.1).
D'autre part, pour prouver que $\mathcal F|U$ est $g$-plat au point $x$, il
suffit, puisque $h:Y'\to Y$ est plat, de montrer (en vertu du critère de
platitude par fibres (11.3.10)) que $\mathcal F_y$ est $g_y$-plat au point $x$.
Mais par hypothèse $\mathcal F_y$ est un $\mathcal O_{X_y}$-Module de
Cohen-Macaulay au point $x$ et $Y_y'$ un préschéma \emph{régulier}~; d'autre
part, comme $g_y$ est fini et surjectif, il vérifie la condition $e')$ de
(15.4.2) en vertu de (5.6.10)~; il suffit donc pour conclure d'appliquer (15.4.2)
à $g_y$ et à $\mathcal F_y$.
\end{proof}

% IV3_B32_P231_END_BEFORE_SUBSECTION_15.5

\subsection{Rang séparable des fibres d'un morphisme quasi-fini et universellement ouvert. Application aux composantes connexes géométriques des fibres d'un morphisme propre}
\label{subsection:IV.15.5-fr}

\begin{proposition}[15.5.1]
\label{IV.15.5.1-fr}
Soient $Y$ un préschéma localement noethérien, $f:X\to Y$ un morphisme séparé et quasi-fini. Pour tout $z\in Y$, soit $n(z)$ le nombre géométrique des points de $f^{-1}(z)$ (ou rang séparable de $f^{-1}(z)$ sur $\mathbf{k}(z)$, cf. \textbf{I}, 6.4.8). On a les propriétés suivantes~:
\begin{enumerate}[label=\emph{(\roman*)}]
  \item La fonction $z\mapsto n(z)$ est semi-continue inférieurement en tout point $y\in Y$ tel que $f$ soit universellement ouvert aux points de $f^{-1}(y)$.
  \item Soit $y$ un point de $Y$ tel que $f$ soit universellement ouvert aux points de $f^{-1}(y)$~; si la fonction $z\mapsto n(z)$ est constante dans un voisinage de $y$, alors $f$ est propre au point $y$ (15.7.1).
  \item On suppose que $f$ soit universellement ouvert et que tout point de $Y$ soit géométriquement unibranche~; alors, si la fonction $z\mapsto n(z)$ est localement constante dans $Y$, les composantes irréductibles de $X$ sont deux à deux disjointes.
\end{enumerate}
\end{proposition}

\begin{proof}
(i) Notons que puisque $f$ est quasi-fini, le nombre $n(z)$ est le nombre géométrique des composantes connexes de $f^{-1}(z)$~; on sait que la fonction $z\mapsto n(z)$ est localement constructible (9.7.9)~; compte tenu de ($\mathbf{0}_{\mathrm{III}}$, 9.3.4), on est ramené à prouver le

\begin{corollary}[15.5.2]
\label{IV.15.5.2-fr}
Sous les hypothèses générales de (15.5.1), soit $y$ un point de $Y$ tel que $f$ soit universellement ouvert aux points de $f^{-1}(y)$. Alors, si $y'$ est une générisation de $y$, on a
\begin{equation}
\tag{15.5.2.1}
\label{IV.15.5.2.1-fr}
  n(y)\leqslant n(y').
\end{equation}
\end{corollary}

Pour tout changement de base $g:Z\to Y$, $f_{(Z)}:X_{(Z)}\to Z$ est séparé, quasi-fini et universellement ouvert aux points de $X_{(Z)}$ se projetant dans $f^{-1}(y)$~; en outre, si $z,z'$ sont deux points de $Z$ tels que $g(z)=y$, $g(z')=y'$ et que $z'$ soit une générisation de $z$, on a $n(z)=n(y)$ et $n(z')=n(y')$ (I, 6.4.12)~; il suffit donc de démontrer (15.5.2.1) pour un $g$ convenable. On peut évidemment supposer $y\ne y'$. Nous utiliserons le lemme suivant~:

\begin{lemma}[15.5.2.2]
\label{IV.15.5.2.2-fr}
Sous les hypothèses générales de (15.5.1), si $y$ est un point de $Y$, $y'$ une générisation de $y$ distincte de $y$, il existe un spectre d'anneau de valuation discrète $Z$, de point fermé $z$ et de point générique $z'$, et un morphisme $g:Z\to Y$ tels que~:
\begin{enumerate}[label=\emph{\arabic*\textsuperscript{\textdegree}}]
  \item $g(z)=y$, $g(z')=y'$~;
  \item si l'on pose $f'=f_{(Z)}$, $n(y)$ est le nombre de points de $f'^{-1}(z)$ et $n(y')$ est le nombre de points de $f'^{-1}(z')$.
\end{enumerate}
\end{lemma}

En effet, il existe un anneau de valuation discrète $A_1$ et un morphisme $g_1$ de $Z_1=\operatorname{Spec}(A_1)$ dans $Y$ tels que, si $z_1$ et $z'_1$ sont le point fermé et le point générique de $Z_1$,
\oldpage[IV-3]{232}
on ait $g_1(z_1)=y$ et $g_1(z'_1)=y'$ (\textbf{II}, 7.1.9)~; on peut donc déjà supposer que $Y$ est le spectre d'un anneau de valuation discrète $A$, $y$ son point fermé et $y'$ son point générique. Pour tout $x\in f^{-1}(y)$, $\mathbf{k}(x)$ est une extension finie de $\mathbf{k}(y)$ par hypothèse (\textbf{II}, 6.2.2)~; on déduit de (4.5.11) qu'il existe une extension algébrique finie $k$ de $\mathbf{k}(y)$ telle que le nombre de points de $f^{-1}(y)\otimes_{\mathbf{k}(y)}k$ soit égal à $n(y)$. Comme il y a un anneau de valuation discrète $A_2$ dominant $A$ et dont le corps résiduel est fini sur $\mathbf{k}(y)$ et contient $k$ (\textbf{II}, 7.1.2), on peut en second lieu supposer que $n(y)$ est le nombre de points de $f^{-1}(y)$. Enfin, le même raisonnement montre qu'il y a une extension finie $K$ de $\mathbf{k}(y')$ telle que le nombre de points de $f^{-1}(y')\otimes_{\mathbf{k}(y')}K$ soit égal à $n(y')$~; si $w$ est une valuation de $\mathbf{k}(y')$ associée à $A$, il y a une valuation discrète de $K$ prolongeant $w$, et l'anneau $A_3$ de cette valuation domine $A$, a $K$ pour corps des fractions, et répond aux conditions du lemme.

On peut donc désormais supposer que $Y$ est le spectre d'un anneau de valuation discrète, $y$ son point fermé, $y'$ son point générique, et que $n(y)$ et $n(y')$ sont les nombres de points de $f^{-1}(y)$ et $f^{-1}(y')$, de sorte que pour tout $x\in f^{-1}(y)$ (resp. tout $x'\in f^{-1}(y')$) on a $\mathbf{k}(x)=\mathbf{k}(y)$ (resp. $\mathbf{k}(x')=\mathbf{k}(y')$).

Désignons alors par $X_i$ les sous-préschémas réduits de $X$ ayant pour espaces sous-jacents les composantes irréductibles de $X$ ($1\leqslant i\leqslant m$). En vertu de (1.10.4), la restriction $X_i\to Y$ de $f$ à tout $X_i$ est un morphisme dominant, et comme $f^{-1}(y')$ est discret, chacune des intersections $X_i\cap f^{-1}(y')$ est réduite au point générique de $X_i$~; d'où (en désignant par $n_i(z)$ le nombre géométrique de points de $X_i\cap f^{-1}(z)$ pour tout $z\in Y$), $n_i(y')=1$ pour tout $i$, et $n(y')=\sum_{i=1}^{m}n_i(y')=m$. Par ailleurs, $f^{-1}(y)$ est réunion des $X_i\cap f^{-1}(y)$, d'où
\begin{equation}
\tag{15.5.2.3}
\label{IV.15.5.2.3-fr}
  n(y)\leqslant\sum_{i=1}^{m}n_i(y),
\end{equation}
et, pour prouver (15.5.2.1), il suffit d'établir que $n_i(y)\leqslant 1$ pour tout $i$. Autrement dit, on peut supposer $X$ intègre, et le morphisme $f$ birationnel. Mais alors, pour tout $x\in f^{-1}(y)$, on a $\mathcal O_y\subset\mathcal O_x\subset K$, où $K=\mathbf{k}(y')$ est le corps des fractions de $\mathcal O_y$. Comme $\mathcal O_y$ est un anneau de valuation et que $\mathcal O_x$ domine $\mathcal O_y$, on a nécessairement $\mathcal O_y=\mathcal O_x$ (\textbf{II}, 7.1.1). Il existe alors une $Y$-section $h:Y\to X$ telle que $h(y)=x$ (\textbf{I}, 2.4.4)~; comme $Y$ est réduit et que $X$ est séparé sur $Y$, une telle section est unique (\textbf{I}, 7.2.2), donc l'ensemble $f^{-1}(y)$ ne contient qu'un seul point, ce qui achève de prouver (i).

(ii) Comme au début de (i), on est ramené à montrer que si les deux membres de (15.5.2.1) sont égaux pour toute générisation $y'$ de $y$, $f$ est propre au point $y$. Grâce au critère de propreté locale (15.7.5) et aux remarques du début on peut encore supposer que $Y=\operatorname{Spec}(A)$ est le spectre d'un anneau de valuation discrète $A$, $y$ son point fermé et $y'$ son point générique, et il s'agit alors de montrer que $f$ est propre. Utilisons (15.5.2.2) et notons que si $A'$ est un anneau de valuation discrète dominant $A$, $A'$ est un $A$-module sans torsion, donc le morphisme $\operatorname{Spec}(A')\to\operatorname{Spec}(A)$ est fidèlement plat et quasi-compact, et il revient donc au même de dire que $f$ est propre ou que le morphisme déduit de $f$ par le changement de base $\operatorname{Spec}(A')\to\operatorname{Spec}(A)$ est
\oldpage[IV-3]{233}
propre (2.7.1, (vii))~; on peut donc supposer que $n(y)$ et $n(y')$ sont les nombres de points des fibres $f^{-1}(y)$ et $f'^{-1}(y')$ respectivement. Avec les notations de (i), on a alors par (15.5.2.3) et (15.5.2.1)
\begin{equation}
\tag{15.5.2.4}
\label{IV.15.5.2.4-fr}
  n(y)\leqslant\sum_i n_i(y)\leqslant\sum_i n_i(y')=n(y').
\end{equation}
et, pour que les membres extrêmes soient égaux, il faut que $n_i(y)=n_i(y')=1$ pour tout $i$. Comme on l'a vu dans (i), si $n_i(y)=1$ pour tout $i$, $X_i\cap f^{-1}(y)$ n'est pas vide et la restriction $X_i\to Y$ de $f$ est un isomorphisme~; comme les $X_i$ sont des sous-préschémas fermés de $X$, cela entraîne que $f$ est propre (\textbf{II}, 5.4.5).

(iii) On peut se borner au cas où $Y$ est affine et réduit, et comme $Y$ est par hypothèse localement intègre (\textbf{I}, 5.1.4), on peut supposer $Y$ intègre, et la fonction $y\mapsto n(y)$ constante dans $Y$. Soient encore $X_i$ ($1\leqslant i\leqslant m$) les composantes irréductibles de $X$. Il résulte de (1.10.4) que si $y'$ est le point générique de $Y$, chacun des $X_i\cap f^{-1}(y')$ est réduit à un seul point~; la restriction $f_i:X_i\to Y$ de $f$ étant un morphisme quasi-fini et dominant et tout point de $Y$ étant géométriquement unibranche, il résulte du critère de Chevalley (14.4.4) que $f_i$ est universellement ouvert. Si alors $y$ est un point quelconque de $Y$, on a $n_i(y)\leqslant n_i(y')$~; d'autre part, comme les $X_i\cap f^{-1}(y)$ sont deux à deux disjoints, on a $\sum_i n_i(y')=n(y')$~; on a donc encore les relations (15.5.2.4). Mais par hypothèse les membres extrêmes sont égaux, donc $n(y)=\sum_i n_i(y)$, ce qui implique que les $X_i\cap f^{-1}(y)$ sont deux à deux disjoints, et achève la démonstration.
\end{proof}

Combinant cette proposition avec le théorème de connexion de Zariski et ses conséquences (\textbf{III}, 4.3), on obtient les résultats suivants qui complètent (\textbf{III}, 4.3.7 et 4.3.10)~:

\begin{proposition}[15.5.3]
\label{IV.15.5.3-fr}
Soient $Y$ un préschéma irréductible localement noethérien, $f:X\to Y$ un morphisme propre~; pour tout $y\in Y$, soit $n(y)$ le nombre géométrique de composantes connexes (4.5.2) de $f^{-1}(y)$. On suppose que la restriction de $f$ à toute composante irréductible de $X$ soit un morphisme dominant dans $Y$. Alors, si $y_0$ est un point géométriquement unibranche de $Y$, la fonction $y\mapsto n(y)$ est semi-continue inférieurement au point $y_0$.
\end{proposition}

\begin{proof}
Considérons la factorisation de Stein $f=g\circ f'$ de $f$ (\textbf{III}, 4.3.3), où $g:Y'\to Y$ est un morphisme fini et $f':X\to Y$ est un morphisme surjectif dont les fibres sont géométriquement connexes (\textbf{III}, 4.3.4). Puisque $f'$ est surjectif, l'image réciproque par $f'$ de toute composante irréductible de $Y'$ contient une composante irréductible de $X$ au moins~; donc chacune des composantes irréductibles de $Y'$ domine $Y$. On peut donc appliquer le critère de Chevalley (14.4.4) aux restrictions de $g$ à chacune de ces composantes irréductibles, et on en conclut que $g$ est universellement ouvert en tout point de $g^{-1}(y_0)$. La conclusion résulte donc de (15.5.2) et du fait que le nombre géométrique des composantes connexes de $f^{-1}(y)$ est égal au nombre géométrique des points de $g^{-1}(y)$ (\textbf{III}, 4.3.4).
\end{proof}

\begin{corollary}[15.5.4]
\label{IV.15.5.4-fr}
Soient $Y$ un préschéma localement noethérien, $f:X\to Y$ un morphisme propre~; soit $y$ un point de $Y$ tel que $f$ soit universellement ouvert aux points de $f^{-1}(y)$~;
\oldpage[IV-3]{234}
alors, avec les notations de (15.5.3), la fonction $z\mapsto n(z)$ est semi-continue inférieurement au point $y$.
\end{corollary}

\begin{proof}
Avec les notations de (15.5.3), notons que dans la factorisation de Stein $f=g\circ f'$, $f'$ est surjectif~; comme $f$ est universellement ouvert aux points de $f^{-1}(y)$, $g$ est universellement ouvert aux points de $g^{-1}(y)$ (14.3.4, (i))~; il suffit alors d'appliquer à $g$ la proposition (15.5.1), (i) en tenant compte de (\textbf{III}, 4.3.4).
\end{proof}

\begin{remarks}[15.5.5]
\label{IV.15.5.5-fr}
\begin{enumerate}[label=\emph{(\roman*)}]
  \item Même si $Y$ est intègre et normal, $X$ intègre, $f$ fini et surjectif (donc universellement ouvert en vertu du critère de Chevalley (14.4.4)), la fonction $y\mapsto n(y)$ n'est pas nécessairement localement constante. On en a un exemple en prenant $Y=\operatorname{Spec}(k[T])$ où $k$ est algébriquement clos, et $X=\operatorname{Spec}(A)$, où $A=k[S,T]/(S^2+T^2-1)$~; aux points $y',y''$ de $Y$ correspondant aux idéaux maximaux $(T-1)$ et $(T+1)$, on a $n(y')=n(y'')=1$, mais $n(y)=2$ en tous les autres points de $Y$. On va donner ci-dessous une condition supplémentaire assurant que $y\mapsto n(y)$ est localement constante (15.5.7).
  \item L'exemple (14.4.10, (ii)) montre que dans (15.5.1, (iii)), on ne peut supprimer l'hypothèse que les points de $Y$ sont géométriquement unibranches.
  \item L'exemple (11.7.5) montre que la conclusion de (15.5.1, (i)) n'est plus valable si on suppose seulement que le morphisme $f$ est ouvert (même si, comme c'est le cas dans l'exemple cité, $f$ est fini et surjectif).
  \item Enfin, dans (15.5.1), on ne peut se dispenser de l'hypothèse que le morphisme $f$ est séparé, comme le montre l'exemple où $X$ est la droite affine dont on a «~dédoublé un point~» (\textbf{I}, 5.5.11), et $Y$ la droite affine~: ici $f$ est un isomorphisme local (donc est plat) et est quasi-fini, mais $z\mapsto n(z)$ n'est pas semi-continue inférieurement.
\end{enumerate}
\end{remarks}

\begin{lemma}[15.5.6]
\label{IV.15.5.6-fr}
Soit $Y$ le spectre d'un anneau de valuation discrète, $a$ son point fermé, $b$ son point générique. Soit $f:X\to Y$ un morphisme localement de type fini et ouvert. Supposons que $X$ soit connexe et que la fibre $f^{-1}(a)$ soit un préschéma réduit. Alors $f^{-1}(b)$ est connexe.
\end{lemma}

\begin{proof}
Nous utiliserons le lemme purement topologique suivant (cas particulier d'un résultat plus général du chap. III, 3\textsuperscript{e} Partie)~:
\begin{lemma}[15.5.6.1]
\label{IV.15.5.6.1-fr}
Soit $X$ un espace localement noethérien, dont toute partie fermée irréductible admet un point générique. Soit $Z$ une partie fermée rare de $X$ ayant la propriété suivante~: pour tout $x\in Z$, si l'on désigne par $T_x$ l'ensemble des générisations de $x$ dans $X$, alors $T_x-\{x\}$ est connexe. Dans ces conditions, si $X$ est connexe, $X-Z$ est connexe.
\end{lemma}

\begin{proof}
Donnons-en une démonstration indépendante. Raisonnant par l'absurde, supposons que l'ouvert $U=X-Z$, partout dense dans $X$, soit réunion de deux ouverts non vides sans point commun $U'$ et $U''$, et désignons par $X'$ et $X''$ les adhérences dans $X$ de $U'$ et $U''$. Comme $X=\overline{U}=X'\cup X''$ et que $X$ est connexe, on a $X'\cap X''\ne\varnothing$, et il est clair que $X'\cap X''\subset Z$. Soit $x$ un point générique d'une composante irréductible de $X'\cap X''$. Comme $X$ est localement noethérien, il y a un voisinage ouvert $V$ de $x$ dans $X$ tel que $V\cap U'$ et $V\cap U''$ n'aient qu'un nombre fini de composantes irréductibles~; comme $x$ est adhérent à $U'$ et à $U''$, il est nécessairement dans l'adhérence d'une composante irréductible $Z'$ de $V\cap U'$ et dans l'adhérence d'une composante irréductible $Z''$ de $V\cap U''$. Mais $Z'$ (resp. $Z''$) est fermé et irréductible dans $X$, donc admet un point générique $z'$ (resp. $z''$), et on a nécessairement $z'\in U'$ et $z''\in U''$. Ceci prouve que les intersections de $T_x-\{x\}$ avec $U'$ et $U''$ sont non vides. Mais par définition $(X'\cap X'')\cap(T_x-\{x\})=\varnothing$~; donc les intersections de $X'$ et de $X''$ avec $T_x-\{x\}$ sont deux fermés non vides disjoints dans $T_x-\{x\}$, dont la réunion est $T_x-\{x\}$~; or cela contredit l'hypothèse que $T_x-\{x\}$ est connexe. C.Q.F.D.
\end{proof}
\oldpage[IV-3]{235}
Pour prouver (15.5.6), nous appliquerons le lemme (15.5.6.1) à $X$ et à sa partie fermée $Z=f^{-1}(a)$ qui est rare puisque $f$ est ouvert. Notons que l'hypothèse, compte tenu de (15.2.2.1), implique que $f$ est plat aux points de $f^{-1}(a)$. Pour un tel point $x$, $\mathcal O_x$ est donc un $\mathcal O_a$-module sans torsion ($\mathbf{0}_{\mathrm I}$, 6.3.4), et en particulier, si $t$ est une uniformisante de $\mathcal O_a$, $t$ est $\mathcal O_x$-régulier. En outre $\mathcal O_x/t\mathcal O_x$ est réduit par hypothèse~; s'il est de profondeur $0$, c'est donc un corps ($\mathbf{0}$, 16.4.7), et comme $t\mathcal O_x$ est alors l'idéal maximal de $\mathcal O_x$, $\mathcal O_x$ est un anneau de valuation discrète ($\mathbf{0}$, 17.1.4), donc $\operatorname{Spec}(\mathcal O_x)-\{x\}$ est réduit à un point, et a fortiori est connexe. Si au contraire $\operatorname{prof}(\mathcal O_x/t\mathcal O_x)\geqslant 1$, on a $\operatorname{prof}(\mathcal O_x)\geqslant 2$ puisque $t$ est $\mathcal O_x$-régulier ($\mathbf{0}$, 16.4.6)~; il résulte alors du théorème de Hartshorne (5.10.7) que $\operatorname{Spec}(\mathcal O_x)-\{x\}$ est connexe, ce qui termine la démonstration.
\end{proof}

\begin{proposition}[15.5.7]
\label{IV.15.5.7-fr}
Soient $Y$ un préschéma localement noethérien, $f:X\to Y$ un morphisme propre. Pour tout $z\in Y$, soit $n(z)$ le nombre géométrique de composantes connexes de $f^{-1}(z)$. Soit $y$ un point de $Y$ tel que $f$ soit universellement ouvert aux points de $f^{-1}(y)$ et que $f^{-1}(y)$ soit géométriquement réduit sur $\mathbf{k}(y)$ (4.6.2). Alors la fonction $z\mapsto n(z)$ est constante dans un voisinage de $y$.
\end{proposition}

\begin{proof}
On sait déjà (15.5.4) que sous les conditions de l'énoncé, $z\mapsto n(z)$ est semi-continue inférieurement au point $y$~; tout revient à voir qu'elle est aussi semi-continue supérieurement, et par le même raisonnement qu'au début de la démonstration de (15.5.1, (i)), il suffit de montrer que si $y'$ est une générisation de $y$, on a
\begin{equation}
\tag{15.5.7.1}
\label{IV.15.5.7.1-fr}
  n(y')\leqslant n(y).
\end{equation}

Utilisant le raisonnement du début de la démonstration de (15.5.2) et le lemme (15.5.2.2) appliqué au morphisme fini $g$ de la factorisation de Stein $f=g\circ f'$ de $f$ (en se souvenant que le nombre géométrique de points de $g^{-1}(y)$ est égal à celui des composantes connexes de $f^{-1}(y)$ ($\mathbf{III}$, 4.3.4)), on se ramène au cas où $n(y)$ et $n(y')$ sont respectivement le nombre de composantes connexes de $f^{-1}(y)$ et $f^{-1}(y')$, et où $Y$ est le spectre d'un anneau de valuation discrète, $y$ le point fermé et $y'$ le point générique de $Y$. On peut en outre remplacer $X$ par une quelconque de ses composantes connexes, ces dernières étant ouvertes et fermées dans $X$, et propres sur $Y$. Supposons donc $X$ connexe~; l'hypothèse que $f$ est ouvert aux points de $f^{-1}(y)$ entraîne que si $f^{-1}(y)\ne\varnothing$, on a aussi $f^{-1}(y')\ne\varnothing$ ($\mathbf{0}_{\mathrm I}$, 1.10.3)~; l'hypothèse que $f$ est fermé entraîne que si $f^{-1}(y')\ne\varnothing$, on a aussi $f^{-1}(y)\ne\varnothing$ (\textbf{II}, 7.2.1). Donc, si $X\ne\varnothing$ (ce qu'on peut évidemment supposer, la proposition étant triviale dans le cas contraire), $f^{-1}(y)$ et $f^{-1}(y')$ sont tous deux non vides. En outre, l'hypothèse entraîne que le préschéma $f^{-1}(y)$ est réduit (4.6.1)~; comme $f$ est ouvert aux points de $f^{-1}(y)$, le lemme (15.5.6) implique que $f^{-1}(y')$ est connexe, autrement dit $n(y')=1$. Comme $f^{-1}(y)$ n'est pas vide, cela prouve (15.5.7.1).
\end{proof}

\begin{remark}[15.5.8]
\label{IV.15.5.8-fr}
La relation (15.5.7.1) n'est plus nécessairement valable si $f$ n'est pas supposé propre, même s'il vérifie les autres hypothèses de (15.5.7). Avec les notations de (15.5.5, (i)), il suffit pour le voir de considérer la restriction de $f$ à $X-\{x'\}$, où $x'$ est un des deux points de $X$ au-dessus d'un point $y$ distinct de $y'$ et $y''$~; on a alors $n(y)=1$, tandis que si $\eta$ est le point générique de $Y$, $n(\eta)=2$.
\end{remark}

\oldpage[IV-3]{236}
\begin{proposition}[15.5.9]
\label{IV.15.5.9-fr}
Soit $f:X\to Y$ un morphisme plat et localement de présentation finie~; pour tout $y\in Y$, soit $n(y)$ le nombre géométrique de composantes connexes de $f^{-1}(y)$.
\begin{enumerate}[label=\emph{(\roman*)}]
  \item Si $f$ est séparé et quasi-fini, la fonction $y\mapsto n(y)$ (égale alors au nombre géométrique des points de $f^{-1}(y)$) est semi-continue inférieurement dans $Y$~; si elle est constante dans un voisinage de $y_0$, $f$ est propre au point $y_0$.
  \item Si $f$ est propre, la fonction $y\mapsto n(y)$ est semi-continue inférieurement dans $Y$~; si de plus $f^{-1}(y_0)$ est géométriquement réduit sur $\mathbf{k}(y_0)$, $n(y)$ est constante dans un voisinage de $y_0$.
\end{enumerate}
\end{proposition}

\begin{proof}
Dans tous les cas, les questions sont locales sur $Y$, donc on peut supposer $Y$ affine et $f$ de présentation finie~; utilisant (8.9.1), (8.10.5) et (11.2.7), on se ramène au cas où $Y$ est noethérien (en utilisant en outre (8.2.11) pour la seconde assertion de (i)). Comme de plus $f$ est alors universellement ouvert (2.4.6), il suffit d'appliquer (15.5.1), (15.5.4) et (15.5.7) pour conclure.
\end{proof}

\begin{remark}[15.5.10]
\label{IV.15.5.10-fr}
On notera qu'on a ainsi obtenu une autre démonstration de (12.2.4, (vi)). Inversement, on peut déduire (15.5.7) de (12.2.4, (vi))~: en effet, on peut se borner au cas où $Y$ est réduit (en remplaçant $f$ par $f_{\mathrm{red}}$), et il résulte alors des hypothèses de (15.5.7) et de (15.2.3) que $f$ est plat dans un voisinage ouvert de $f^{-1}(y)$~; comme de plus $f$ est propre, on peut supposer que ce voisinage est de la forme $f^{-1}(U)$, où $U$ est un voisinage ouvert de $y$ dans $Y$. On conclut alors par (12.2.4, (vi)). La démonstration de (15.5.7) donnée plus haut a l'avantage de mettre en évidence le résultat (15.5.6), qui a un intérêt indépendant.
\end{remark}

% IV3_B33_P236_END_BEFORE_SUBSECTION_15.6

\subsection{Composantes connexes des fibres le long d'une section}
\label{subsection:IV.15.6-fr}

\begin{proposition}[15.6.1]
\label{IV.15.6.1-fr}
Soient $Y$ un préschéma localement noethérien, $f:X\to Y$ un morphisme localement de type fini, $g:Y\to X$ une $Y$-section de $X$ (\textbf{I}, 2.5.5). Pour tout $y\in Y$, désignons par $X_y^0$ la composante connexe de $f^{-1}(y)$ contenant le point $g(y)$. Soient $y$ un point de $Y$, $y'$ une générisation de $y$ dans $Y$. Alors :
\begin{enumerate}[label=\emph{(\roman*)}]
  \item S'il existe une composante irréductible $Z$ de $X_{y'}^0$, de dimension égale à $\dim(X_{y'}^0)$ et contenant $g(y')$ (ce qui sera le cas si $X_{y'}^0$ est irréductible), on a
\begin{equation}
\tag{15.6.1.1}
\label{IV.15.6.1.1-fr}
  \dim(X_y^0)\geqslant\dim(X_{y'}^0).
\end{equation}
  \item Si de plus $X_y^0$ est irréductible et si les deux membres de (15.6.1.1) sont égaux, alors $X_y^0\subset\overline{X_{y'}^0}$ (adhérence dans $X$).
\end{enumerate}
\end{proposition}

\begin{proof}
(i) Soit $\overline Z$ l'adhérence de $Z$ dans $X$~; comme $y$ est adhérent à $y'$ et $g$ continue, on a $g(y)\in\overline Z$. Comme $Z=\overline Z\cap f^{-1}(y')$, il résulte du théorème de semi-continuité de Chevalley (13.1.3) appliqué à la restriction de $f$ au sous-préschéma fermé réduit de $X$ ayant $\overline Z$ pour espace sous-jacent, que l'on a
\[
  \dim_{g(y)}(\overline Z\cap f^{-1}(y))\geqslant\dim(Z)~;
\]
mais les composantes irréductibles de $\overline Z\cap f^{-1}(y)$ contenant $g(y)$ sont évidemment contenues dans $X_y^0$, d'où \emph{a fortiori} l'inégalité (15.6.1.1).

(ii) Le raisonnement de (i) montre en outre que si les deux membres de (15.6.1.1) sont égaux, on a nécessairement
\[
  \dim_{g(y)}(\overline Z\cap f^{-1}(y))=\dim_{g(y)}(X_y^0)~;
\]
si $X_y^0$ est irréductible, cela entraîne $\overline Z\cap f^{-1}(y)=X_y^0$, donc $X_y^0$ est contenu dans $\overline Z$, et \emph{a fortiori} dans $\overline{X_{y'}^0}$.
\end{proof}

\oldpage[IV-3]{237}
\begin{corollary}[15.6.2]
\label{IV.15.6.2-fr}
Avec les notations de (15.6.1), supposons en outre que pour tout $y\in Y$, $X_y^0$ soit irréductible. Alors la fonction $z\mapsto\dim(X_z^0)$ est semi-continue supérieurement dans $Y$. Si en outre cette fonction est constante dans un voisinage d'un point $y\in Y$, alors on a $X_y^0\subset\overline{X_{y'}^0}$ pour toute générisation $y'$ de $y$.
\end{corollary}

\begin{proof}
Soit $y$ un point de $Y$, et soit $U$ un voisinage ouvert affine de $g(y)$ dans $X$~; alors il y a un voisinage ouvert $V$ de $y$ dans $Y$ tel que $g(V)\subset U$, et pour tout $z\in V$, $U\cap X_z^0$ est dense dans $X_z^0$, donc $\dim(X_z^0)=\dim(U\cap X_z^0)$ (4.1.1.3). Remplaçant $X$ par $U$, on peut donc supposer que $f$ est un morphisme de type fini. On sait alors (9.7.10) que la fonction $z\mapsto\dim(X_z^0)$ est localement constructible, et les assertions du corollaire résultent donc de (15.6.1) et de ($\mathbf{0}_{\mathrm{III}}$, 9.3.4).
\end{proof}

\begin{proposition}[15.6.3]
\label{IV.15.6.3-fr}
Avec les notations de (15.6.1), supposons que pour tout $y\in Y$, $X_y^0$ soit géométriquement irréductible (4.5.2). Alors, pour tout $y\in Y$ tel que la fonction $z\mapsto\dim(X_z^0)$ soit constante dans un voisinage de $y$, $f$ est universellement ouvert aux points de $X_y^0$.
\end{proposition}

\begin{proof}
Appliquons le critère (14.3.7). Soit donc $h:Y'\to Y$ un morphisme, $Y'$ étant le spectre d'un anneau de valuation discrète, dont nous désignerons par $t$, $t'$ le point fermé et le point générique respectivement, et supposons que $h(t)=y$, de sorte que $y'=h(t')$ est une générisation de $y$. Posons $X'=X\times_Y Y'$, $f'=f_{(Y')}:X'\to Y'$, $g'=g_{(Y')}:Y'\to X'$. Avec les mêmes notations que dans (15.6.1), l'hypothèse faite sur les $X_y^0$ implique que
\[
  X_t^{\prime 0}=X_y^0\otimes_{\mathbf{k}(y)}\mathbf{k}(t)
  \quad\text{et}\quad
  X_{t'}^{\prime 0}=X_{y'}^0\otimes_{\mathbf{k}(y')}\mathbf{k}(t'),
\]
et que $X_t^{\prime 0}$ et $X_{t'}^{\prime 0}$ sont irréductibles (4.4.1)~; comme $\dim(X_y^0)=\dim(X_t^{\prime 0})$ et $\dim(X_{y'}^0)=\dim(X_{t'}^{\prime 0})$ (4.1.4), on voit qu'on est ramené à prouver la proposition pour $f'$ et $X_t^{\prime 0}$, autrement dit on peut se borner au cas où $Y$ est spectre d'un anneau de valuation discrète, $y$ son point fermé et $y'$ son point générique~; si $x'$ est le point générique de $X_{y'}^0$, il résulte alors de (15.6.2) que tout point de $X_y^0$ est spécialisation de $x'$, d'où la conclusion.
\end{proof}

\begin{proposition}[15.6.4]
\label{IV.15.6.4-fr}
Sous les hypothèses générales de (15.6.1), soit $X^0$ la réunion des $X_y^0$ lorsque $y$ parcourt $Y$. Si $y\in Y$ est tel que $X_y^0$ soit géométriquement réduit sur $\mathbf{k}(y)$ (4.6.2) et si $f$ est universellement ouvert en tout point de $X_y^0$, alors $X^0$ est un voisinage de $X_y^0$ dans $X$. En particulier, si $f$ est universellement ouvert et si, pour tout $y\in Y$, $X_y^0$ est géométriquement réduit sur $\mathbf{k}(y)$, $X^0$ est ouvert dans $X$.
\end{proposition}

\begin{proof}
Montrons d'abord qu'on peut se ramener au cas où $f$ est un morphisme \emph{de type fini}. Soit $x$ un point de $X_y^0$~; il résulte de (5.10.8.1) (avec $\mathfrak F=\{\varnothing\}$) qu'il y a une suite finie $(Z_i)_{0\leqslant i\leqslant n}$ de composantes irréductibles de $f^{-1}(y)$ telle que $g(y)\in Z_0$, $x\in Z_n$ et $Z_i\cap Z_{i+1}\ne\varnothing$ pour $0\leqslant i\leqslant n-1$~; les $Z_i$ étant irréductibles, il y a une suite finie $(U_i)$ ($1\leqslant i\leqslant n$) d'ouverts affines dans $f^{-1}(y)$ tels que $g(y)\in U_1$, $x\in U_n$, $U_i$ étant un voisinage affine d'un point $x_i$ de $Z_i\cap Z_{i-1}$ pour $1\leqslant i\leqslant n$. Il y a pour chaque $i$ un ouvert quasi-compact $W_i$ dans $X$ tel que $U_i=f^{-1}(y)\cap W_i$~; soit $W$ l'ouvert quasi-compact de $X$ réunion des $W_i$. Comme $g$ est continue, il y a un voisinage ouvert affine $V$ de $y$ dans $Y$ tel que $g(V)\subset W$~; si $X'=W\cap f^{-1}(V)$, la restriction de $f$ à $X'$ est de type fini~; en outre, si, pour tout $z\in V$, $X_z^{\prime 0}$ est la composante connexe de $X'\cap f^{-1}(z)$ contenant $g(z)$, on a $X_z^{\prime 0}\subset X_z^0$, et $x$ appartient à $X_y^{\prime 0}$. En effet, chacun des $U_i$ contient les deux ensembles irréductibles $U_i\cap Z_{i-1}$ et $U_i\cap Z_i$ qui se rencontrent, et les ensembles irréductibles $U_i\cap Z_{i-1}$ et $U_{i-1}\cap Z_{i-1}$

\oldpage[IV-3]{238}
se rencontrent pour $2\leqslant i\leqslant n$~; la réunion des $U_i\cap Z_{i-1}$ et des $U_i\cap Z_i$ pour $1\leqslant i\leqslant n$ est donc connexe et contient $x$ et $g(y)$. Comme $f|X'$ est universellement ouvert aux points de $X_y^{\prime 0}$, cela achève de prouver notre assertion, car si la proposition est prouvée lorsque $f$ est de type fini, la réunion $X^{\prime 0}$ des $X_z^{\prime 0}$ pour $z\in V$ sera un voisinage de $x$, et il en sera de même de $X^0$.

Supposons donc $f$ de type fini. Alors $X^0$ est localement constructible (9.7.12)~; il suffit par suite de montrer que $X^0$ contient toute générisation $x'$ de $x$ ($\mathbf{0}_{\mathrm{III}}$, 9.2.5). Posons $y'=f(x')$~; il existe un anneau de valuation discrète $A$ et un morphisme $u$ de $Y'=\operatorname{Spec}(A)$ dans $X$ tels que, si $z$ est le point fermé et $z'$ le point générique de $Y'$, on ait $u(z)=x$ et $u(z')=x'$ (\textbf{II}, 7.1.9)~; si $h=f\circ u$, on a donc $h(z)=y$ et $h(z')=y'$. Posons $X'=X\times_Y Y'$, $f'=f_{(Y')}:X'\to Y'$, qui est un morphisme de type fini universellement ouvert aux points de $f'^{-1}(z)$, et $g'=g_{(Y')}$, qui est une $Y'$-section de $X'$. Si $p_y:f'^{-1}(z)\to f^{-1}(y)$ (resp. $p_{y'}:f'^{-1}(z')\to f^{-1}(y')$) est la projection canonique, $p_y^{-1}(X_y^0)$ (resp. $p_{y'}^{-1}(X_{y'}^0)$) est la composante connexe de $f'^{-1}(z)$ (resp. $f'^{-1}(z')$) contenant $g'(z)$ (resp. $g'(z')$), car les $X_t^0$ sont géométriquement connexes (4.5.13) et il suffit d'appliquer (4.4.1). De plus, le morphisme $u$ correspond à une $Y'$-section $v$ de $X'$ telle que $p_y(v(y))=x$, $p_{y'}(v(y'))=x'$, de sorte que $v(y')$ est une générisation de $v(y)$, et il suffira donc de prouver que $v(y')$ appartient à $p_{y'}^{-1}(X_{y'}^0)$. Enfin, il est clair que $p_y^{-1}(X_y^0)$ est géométriquement réduit sur $\mathbf{k}(z)$.

En d'autres termes, on est ramené au cas où $Y$ est spectre d'un anneau de valuation discrète, $y$ son point fermé, $y'$ son point générique. Comme $f^{-1}(y)-X_y^0$ est alors fermé dans $X$, on peut, en remplaçant $X$ par l'ouvert $X-(f^{-1}(y)-X_y^0)$, supposer que $X_y^0=f^{-1}(y)$~; de même, on peut remplacer $X$ par l'ouvert, composante connexe de $X$, qui contient $g(Y)$, cette composante connexe contenant évidemment $X^0$~; autrement dit, on peut supposer $X$ connexe. La proposition sera alors établie si l'on montre que $X^0=X$, ou encore que $X_{y'}$ est connexe~; mais $f^{-1}(y')$ est géométriquement réduit sur $\mathbf{k}(y')$, et \emph{a fortiori} réduit~; comme en outre $f$ est ouvert aux points de $f^{-1}(y)$, on peut appliquer le lemme (15.5.6), d'où la conclusion.
\end{proof}

\begin{corollary}[15.6.5]
\label{IV.15.6.5-fr}
Soient $f:X\to Y$ un morphisme plat de présentation finie, $g$ une $Y$-section de $X$~; pour tout $y\in Y$, soit $X_y^0$ la composante connexe de $f^{-1}(y)$ contenant $g(y)$, et soit $X^0$ la réunion des $X_y^0$ quand $y$ parcourt $Y$. Alors, pour tout $y\in Y$ tel que $X_y^0$ soit géométriquement réduit sur $\mathbf{k}(y)$, $X^0$ est un voisinage de $X_y^0$ dans $X$. En particulier, si $f$ est réduit (6.8.1), $X^0$ est ouvert dans $X$.
\end{corollary}

\begin{proof}
La question étant locale sur $Y$, on peut supposer $Y=\operatorname{Spec}(A)$ affine. Il existe donc un sous-anneau noethérien $A_0$ et un morphisme plat de type fini $f_0:X_0\to Y_0=\operatorname{Spec}(A_0)$ tels que $X=X_0\times_{Y_0}Y$ et $f=(f_0)_{(Y)}$ (8.9.1 et 11.2.7)~; en outre, on peut supposer qu'il existe une $Y_0$-section $g_0$ de $X_0$ telle que $g=(g_0)_{(Y)}$ (8.8.2). Pour tout $y\in Y$, soit $y_0$ sa projection dans $Y_0$~; il résulte de (4.5.13) que $(X_0)_{y_0}^0$ est géométriquement connexe, donc, si $p:f^{-1}(y)\to f_0^{-1}(y_0)$ est la projection canonique, on a $X_y^0=p^{-1}((X_0)_{y_0}^0)$ (4.4.1)~; de plus, l'hypothèse que $X_y^0$ est géométriquement réduit sur $\mathbf{k}(y)$ entraîne que $(X_0)_{y_0}^0$ est géométriquement réduit sur $\mathbf{k}(y_0)$ (4.6.10). On est ainsi ramené au cas où on suppose

\oldpage[IV-3]{239}
en outre $Y$ noethérien~; comme $f$ est universellement ouvert (11.1.1), il suffit alors d'appliquer (15.6.4).
\end{proof}

\begin{proposition}[15.6.6]
\label{IV.15.6.6-fr}
Soit $f:X\to Y$ un morphisme tel que $Y$ soit localement noethérien et $f$ localement de type fini, ou que $f$ soit de présentation finie. Soit $g$ une $Y$-section de $X$~; pour tout $y\in Y$, soit $X_y^0$ la composante connexe de $f^{-1}(y)$ contenant $g(y)$~; enfin, soit $X^0$ la réunion des $X_y^0$ lorsque $y$ parcourt $Y$.

Supposons que, pour tout $y\in Y$, $X_y^0$ soit géométriquement irréductible. Alors :
\begin{enumerate}[label=\emph{(\roman*)}]
  \item La fonction $y\mapsto\dim(X_y^0)$ est semi-continue supérieurement dans $Y$.
  \item Soient $x_0\in X^0$, $y_0=f(x_0)$. Si la fonction $y\mapsto\dim(X_y^0)$ est constante dans un voisinage de $y_0$, il existe un voisinage ouvert $U$ de $y_0$ tel que $f$ soit universellement ouvert aux points de $X^0\cap f^{-1}(U)$. Si de plus $Y$ est réduit et $f^{-1}(y_0)$ géométriquement réduit sur $\mathbf{k}(y_0)$ au point $x_0$, $f$ est plat au point $x_0$.
  \item Inversement, supposons que $f$ soit universellement ouvert au point générique de $X_{y_0}^0$ et en outre que l'une des conditions suivantes soit vérifiée :
  \begin{enumerate}
    \item[$\alpha$)] La fibre $f^{-1}(y_0)$ est géométriquement réduite sur $\mathbf{k}(y_0)$ au point $g(y_0)$ et l'anneau $\mathcal O_{Y,y_0}$ est noethérien.
    \item[$\beta$)] Pour tout point générique $y'$ d'une composante irréductible de $Y$ contenant $y_0$, on a $\dim(f^{-1}(y'))=\dim(X_{y'}^0)$.
  \end{enumerate}
  Alors la fonction $y\mapsto\dim(X_y^0)$ est constante dans un voisinage de $y_0$.
  \item L'ensemble $W$ des points $x\in X^0$ tels que la fonction $y\mapsto\dim(X_y^0)$ soit constante dans un voisinage de $f(x)$ et que $X_{f(x)}^0$ soit géométriquement réduit sur $\mathbf{k}(f(x))$, est ouvert dans $X$.
\end{enumerate}
\end{proposition}

\begin{proof}
Les questions étant locales sur $Y$, on peut supposer que $Y=\operatorname{Spec}(A)$ est affine. Lorsque $f$ est de présentation finie, il existe donc un sous-anneau noethérien $A_1$ de $A$ et un morphisme de type fini $f_1:X_1\to Y_1=\operatorname{Spec}(A_1)$ tels que $X=X_1\times_{Y_1}Y$ et $f=(f_1)_{(Y)}$ (8.9.1)~; en outre, on peut supposer qu'il existe une $Y_1$-section $g_1$ de $X_1$ telle que $g=(g_1)_{(Y)}$ (8.8.2). Pour tout $y\in Y$, soit $y_1$ sa projection dans $Y_1$~; il résulte de (4.5.13) que $(X_1)_{y_1}^0$ est géométriquement connexe, donc, si $p:f^{-1}(y)\to f_1^{-1}(y_1)$ est la projection canonique, on a $X_y^0=p^{-1}((X_1)_{y_1}^0)$ (4.4.1). Notons d'autre part que l'on peut, en vertu de (9.7.11) et (9.7.12), appliquer (9.3.3) à la propriété d'être géométriquement irréductible~; on peut donc supposer $A_1$ choisi de telle sorte que $(X_1)_{y_1}^0$ soit géométriquement irréductible pour tout $y_1\in Y_1$. On a $\dim(X_y^0)=\dim((X_1)_{y_1}^0)$ (4.1.4), et en appliquant de nouveau (9.3.3) à la propriété d'avoir une dimension donnée (et utilisant (9.5.5) et (9.7.12)), on peut supposer (en restreignant au besoin $Y$ à un voisinage de $y_0$) que si $\dim(X_y^0)$ est constante dans $Y$, alors $\dim((X_1)_{y_1}^0)$ est constante dans $Y_1$. Si $Y$ est réduit, il en est de même de $Y_1$ puisque $A_1\subset A$~; d'autre part, si $f^{-1}(y_0)$ est géométriquement réduit au point $x_0$, $f_1^{-1}(y_{01})$ l'est au point $x_{01}$ ($x_{01}$ et $y_{01}$ étant les projections respectives dans $X_1$ et $Y_1$ de $x_0$ et $y_0$) (4.6.10).

Ces remarques montrent que pour démontrer (i) et (ii), on peut se borner au cas où $Y$ est noethérien et $f$ localement de type fini. Mais dans ce cas, l'assertion (i) résulte de (15.6.2) et la première assertion de (ii) de (15.6.3). D'autre part, si $Y$ est réduit et $f^{-1}(y_0)$ géométriquement réduit sur $\mathbf{k}(y_0)$ au point $x_0$ ($Y$ étant toujours supposé

\oldpage[IV-3]{240}
noethérien), comme $X_y^0$ est la seule composante irréductible de $f^{-1}(y)$ contenant $x$, on déduit de (15.6.3) et (15.2.3) que si $\dim(X_y^0)$ est constante dans un voisinage de $y_0$, $f$ est plat au point $x_0$.

Pour prouver (iii), remarquons d'abord que l'ensemble $X^0$ est localement constructible dans $X$ (9.7.12), donc, par (9.5.5), l'ensemble $E$ des $y\in Y$ tels que $\dim(X_y^0)=\dim(X_{y_0}^0)$ est localement constructible dans $Y$. Appliquons alors (1.10.1) à $E$~: il suffit de prouver (compte tenu de (i)) que $\dim(X_{y'}^0)=\dim(X_{y_0}^0)$ pour le point générique $y'$ d'une composante irréductible de $Y$ contenant $y_0$~; ceci permet déjà de supposer que $Y=\operatorname{Spec}(\mathcal O_{Y,y_0})$.

Si on est dans le cas $\alpha$), on se ramène aussitôt, en vertu de (\textbf{II}, 7.1.9), et utilisant (4.5.13) et (4.4.1) comme dans (15.6.4), au cas où $Y$ est spectre d'un anneau de valuation discrète, $y_0$ son point fermé et $y'$ son point générique. Mais alors les hypothèses entraînent en vertu de (15.2.2.1), que $f$ est \emph{plat} au point $g(y_0)$, donc aussi dans un voisinage $V$ de ce point dans $X$ (11.1.1). Pour démontrer notre assertion, on peut remplacer $X$ par $V$, car $V\cap X_{y'}^0$ n'est pas vide en vertu de (2.3.4), donc est un ouvert partout dense dans $X_{y'}^0$ et a par suite même dimension (4.1.1.3)~; d'ailleurs, comme $X_{y'}^0$ est aussi la composante connexe de $f^{-1}(y')$ contenant $g(y')$, on peut supposer $V$ pris tel que $V$ ne rencontre aucune autre composante irréductible de $f^{-1}(y_0)$ ni de $f^{-1}(y')$, autrement dit on peut supposer que $X^0=X$~; mais alors la conclusion résulte de (12.1.1, (i)), puisque par hypothèse $X_{y_0}^0$ est intègre.

Supposons maintenant qu'on soit dans le cas $\beta$). Appliquons cette fois (\textbf{II}, 7.1.4) de la même manière que (\textbf{II}, 7.1.9) dans le cas $\alpha$)~: on est alors ramené au cas où $Y$ est spectre d'un anneau de valuation (non nécessairement discrète), $y_0$ son point fermé et $y'$ son point générique. En vertu de (14.3.13), il existe une composante irréductible $Z$ de $X$ contenant $X_{y_0}^0$ et dominant $Y$, et telle que $\dim(Z\cap f^{-1}(y'))=\dim(X_{y_0}^0)$~; mais l'hypothèse $\beta$) et le fait que $\dim(Z\cap f^{-1}(y'))\leqslant\dim(f^{-1}(y'))$ montrent que $\dim(X_{y_0}^0)\leqslant\dim(X_{y'}^0)$, d'où $\dim(X_{y_0}^0)=\dim(X_{y'}^0)$ en vertu de (i).

Reste à prouver (iv). Notons d'abord que les ensembles envisagés ne changent pas lorsqu'on remplace $f$ par $f_{(Y_{\mathrm{red}})}:X\times_Y Y_{\mathrm{red}}\to Y_{\mathrm{red}}$, la projection $X\times_Y Y_{\mathrm{red}}\to X$ étant un homéomorphisme. Autrement dit, on peut supposer $Y$ réduit, et alors il résulte de (ii) que $f$ est \emph{plat} aux points de $W$. Considérons un point $x_0$ de $W$ et prouvons que $W$ est un voisinage de $x_0$~; procédant comme au début de la démonstration, et utilisant aussi (11.2.7), on peut de plus supposer $Y$ noethérien~; il résulte alors de (ii) et de (15.6.4) que $X^0$ est un voisinage de $x_0$ dans $X$, et de (12.1.1, (vii)) que $W$ est aussi un voisinage de $x_0$ dans $X$.
\end{proof}

\begin{corollary}[15.6.7]
\label{IV.15.6.7-fr}
Supposons vérifiées les conditions préliminaires de (15.6.6) sur $f$, et supposons en outre que pour tout $y\in Y$, $X_y^0$ soit géométriquement intègre sur $\mathbf{k}(y)$. Alors les conditions suivantes sont équivalentes :
\begin{enumerate}[label=\emph{\alph*)}]
  \item La fonction $y\mapsto\dim(X_y^0)$ est localement constante dans $Y$.
  \item Le morphisme $f_{(Y_{\mathrm{red}})}:X\times_Y Y_{\mathrm{red}}\to Y_{\mathrm{red}}$ déduit de $f$ par changement de base, est plat aux points de $X^0$.
\end{enumerate}

\oldpage[IV-3]{241}
En outre, ces conditions entraînent que $X^0$ est ouvert dans $X$. Enfin, lorsque $Y$ est localement noethérien, a) et b) sont aussi équivalentes à
\begin{enumerate}
  \item[\rm b$'$)] $f$ est universellement ouvert aux points de $X^0$.
\end{enumerate}
\end{corollary}

\begin{proof}
Le fait que a) entraîne b) résulte de (15.6.6, (ii)), ainsi que le fait que $X^0$ est alors ouvert dans $X$. Si b) est vérifiée, on peut se borner au cas où $Y$ est réduit et $f$ plat~; alors, on se ramène, comme au début de (15.6.6), et en utilisant en outre (11.2.7), au cas où $Y$ est noethérien, cas où la conclusion résulte de (15.6.6, (iii), cas $\alpha$)). L'équivalence de b) et b$'$) a déjà été prouvée lorsque $Y$ est localement noethérien (15.2.2.1), compte tenu de ce que le morphisme $X\times_Y Y_{\mathrm{red}}\to X$ est un homéomorphisme universel.
\end{proof}

\begin{proposition}[15.6.8]
\label{IV.15.6.8-fr}
Soient $f:X\to Y$ un morphisme séparé de présentation finie, $g$ une $Y$-section de $X$~; pour tout $y\in Y$, soit $X_y^0$ la composante connexe de $f^{-1}(y)$ qui contient $g(y)$, et soit $X^0$ la réunion des $X_y^0$ quand $y$ parcourt $Y$. Alors, si $y\in Y$ est tel que $X_y^0$ soit propre sur $\mathbf{k}(y)$, $X^0$ est propre sur $Y$ au point $y$ (i.e. (15.7.1), il existe un voisinage ouvert $U$ de $y$ dans $Y$ tel que $X^0\cap f^{-1}(U)$ soit fermé dans $f^{-1}(U)$ et propre sur $U$).
\end{proposition}

\begin{proof}
Procédant comme au début de (15.6.6) (où on utilise (8.10.5, (v)), et où on remplace (9.7.11) par (9.6.7)), on se ramène au cas où $Y$ est noethérien et $f$ de type fini. On sait alors (9.7.12) que $X^0$ est localement constructible dans $X$~; d'autre part, les $X_z^0$ sont géométriquement connexes (pour tout $z\in Y$) en vertu de (4.5.13)~; enfin, comme $g(Y)\subset X^0$, $X^0$ est universellement submersif sur $Y$ (15.7.8). Il suffit alors d'appliquer à $X^0$ le critère (15.7.9).
\end{proof}

\begin{corollary}[15.6.9]
\label{IV.15.6.9-fr}
Soit $f:X\to Y$ un morphisme séparé tel que $Y$ soit localement noethérien et $f$ localement de type fini, ou que $f$ soit de présentation finie. Soit $g$ une $Y$-section de $X$, et pour tout $y\in Y$, soit $X_y^0$ la composante connexe de $f^{-1}(y)$ contenant $g(y)$. Soit $y$ un point de $Y$ tel que $X_y^0$ soit propre sur $\mathbf{k}(y)$. Alors, pour toute générisation $y'$ de $y$, on a $\dim(X_y^0)\geqslant\dim(X_{y'}^0)$.
\end{corollary}

\begin{proof}
Supposons d'abord $f$ de présentation finie~; alors en remplaçant au besoin $Y$ par un voisinage de $y$, on peut supposer que $f|X^0$ est \emph{propre} (15.6.8), et il suffit d'appliquer (13.1.5) à cette restriction. Si $Y$ est localement noethérien et $f$ localement de type fini, $X$ est localement noethérien~; en outre l'hypothèse que $X_y^0$ soit propre sur $\mathbf{k}(y)$ entraîne que $X_y^0$ est noethérien~; il existe donc un ouvert noethérien $V\subset X$ contenant $X_y^0$~; par suite il y a un voisinage ouvert $W$ de $y$ dans $Y$ tel que $g(W)\subset V$. En outre, si $z$ est point maximal de $X_{y'}^0$, on peut supposer que $z\in V$. La restriction de $f$ à $V\cap f^{-1}(W)$ est alors un morphisme de type fini, et $g|W$ une $W$-section~; en appliquant à ce morphisme le résultat déjà obtenu, on voit que si $Z$ est la composante irréductible de $X_{y'}^0$ de point générique $z$, on a $\dim(Z)\leqslant\dim(X_y^0)$. Comme $z$ est un point maximal quelconque de $X_{y'}^0$, on a bien $\dim(X_{y'}^0)\leqslant\dim(X_y^0)$.
\end{proof}

\begin{remarks}[15.6.10]
\label{IV.15.6.10-fr}
(i) L'hypothèse supplémentaire sur l'existence d'une composante irréductible de $X_{y'}^0$ contenant $g(y')$ et de dimension égale à $\dim(X_{y'}^0)$, faite dans (15.6.1), n'est pas superflue, comme le montre l'exemple suivant. Soient $A$ un anneau de valuation discrète, $\pi$ une uniformisante de $A$, $K$ le corps des fractions de $A$, $k$ son corps résiduel. Prenons $Y=\operatorname{Spec}(A)$, et désignons par $y$ et $y'$ le point fermé et le point générique de $Y$. Considérons les deux $Y$-schémas suivants~: $X_1=\operatorname{Spec}(A[t])$, $X_2=\operatorname{Spec}(K[u,v])$, où $t,u,v$ sont des indéterminées (on notera que la projection de $X_2$ dans $Y$ est donc $\{y'\}$). Dans l'anneau $A[t]$, l'idéal principal $(\pi t-1)$ est \emph{maximal}, et correspond donc à un point fermé $x_1$ de $X_1$, qui se projette au point générique $y'$ de $Y$ et est tel que $\mathbf{k}(x_1)=K$.

\oldpage[IV-3]{242}
Considérons d'autre part le point fermé $x_2$ de $X_2$ correspondant à l'idéal maximal $\left(u-\frac{1}{\pi}\right)+(v)$ de $K[u,v]$~; on a aussi $\mathbf{k}(x_2)=K$. Comme on le verra au chap.~V, on peut donc \emph{recoller} $X_1$ et $X_2$ le long des deux sous-préschémas fermés $Z_1$, $Z_2$ de $X_1$, $X_2$ ayant respectivement $\{x_1\}$ et $\{x_2\}$ pour espaces sous-jacents, suivant l'unique $Y$-isomorphisme de $Z_1$ sur $Z_2$. Désignons par $X$ le $Y$-préschéma ainsi obtenu, par $x$ le point de $X$ image de $x_1$ et $x_2$. La «~section nulle~» $g$ de $X_1$ (\textbf{II}, 1.7.9) est encore alors une $Y$-section de $X$~; $f^{-1}(y)$ est égale à $(X_1)_y$, tandis que $f^{-1}(y')$ a deux composantes irréductibles, respectivement isomorphes à $(X_1)_{y'}^0$ et $(X_2)_{y'}^0$, ayant un unique point commun $x$, et de dimensions respectives 1 et 2. Voir cependant la prop.~(15.6.9).

(ii) Considérons le morphisme $f$ défini dans (12.2.3, (ii)), qui est propre et plat~; en outre, la restriction de $f$ à $X_1\cap f_0^{-1}(Y)$ est un isomorphisme, donc le morphisme réciproque $g:Y\to X$ est une $Y$-section de $X$. On a alors $\dim(X_{y_0})=1$ tandis que $\dim(X_y^0)=0$ pour $y\ne y_0$, bien que $X_y^0$ soit géométriquement irréductible pour tout $y\in Y$ (mais $X_{y_0}^0$ n'est pas réduit)~; on voit donc que dans (15.6.6, (iii)), on ne peut supprimer les hypothèses $\alpha$) et $\beta$). En outre, $X^0$ n'est pas un voisinage de $X_{y_0}^{0}$, ce qui prouve que dans (15.6.4), on ne peut se dispenser de l'hypothèse que $X_{y_0}^0$ est réduit.

(iii) Au chap.~VI, nous appliquerons les résultats précédents aux $Y$-préschémas \emph{en groupes} localement de type fini sur un préschéma localement noethérien $Y$. Si $G$ est un tel préschéma, il existe une $Y$-section canonique $g$, la «~section unité~», et l'on montre que $G_y^0$ est toujours géométriquement irréductible et que l'on a $\dim(G_y^0)=\dim(G_y)$, en posant $G_y=f^{-1}(y)$. Il résultera donc de (15.6.2) et (15.6.3) que la fonction $z\mapsto\dim(G_z)$ est semi-continue supérieurement, et que si elle est constante au voisinage d'un point $y$, $f$ est universellement ouvert aux points de $G_y^0$. Si de plus $G_y$ est géométriquement réduit sur $\mathbf{k}(y)$ en un de ses points, on montre que $G_y$ est \emph{lisse} (6.8.1) sur $\mathbf{k}(y)$ en tous ses points~; il résultera alors de (15.6.6) que si de plus $\mathcal O_y$ est réduit, alors $f$ est lisse en tous les points de $G_y^0$ (mais pas nécessairement en tous les points de $G_y$). Ces résultats s'appliqueront par exemple dans la théorie de schémas de Picard, où nous disposerons d'un théorème général assurant que, sous certaines conditions, la fonction $z\mapsto\dim(G_z)$ est localement constante. Remarquons que c'est en vue d'applications de cette nature que les énoncés tels que (15.6.1) sont donnés pour les morphismes \emph{localement} de type fini, et non seulement pour les morphismes de type fini.

On notera aussi que dans le cas d'un $Y$-préschéma en groupes $G$, l'hypothèse $\beta$) de (15.6.6) est toujours vérifiée.
\end{remarks}

% EGA IV-3 terminal continuation, printed pp. 242--255.
% Direct diplomatic French transcription from the NUMDAM authority.

\subsection{Appendice : Critères valuatifs de propreté locale}
\label{subsection:IV.15.7-fr}

Ce numéro donne des compléments au critère valuatif de propreté démontré dans (\textbf{II}, 7.3.10)~; il est indépendant du reste du \S~15.

\begin{definition}[15.7.1]
\label{IV.15.7.1-fr}
Soient $f:X\to Y$ un morphisme de préschémas, $y$ un point de $Y$. On dit que $f$ est propre au point $y$ s'il existe un voisinage ouvert $U$ de $y$ dans $Y$ tel que la restriction $f^{-1}(U)\to U$ de $f$ soit un morphisme propre. On dit qu'une partie $Z$ de $X$ est propre sur $Y$ au point $y$ s'il existe un voisinage ouvert $U$ de $y$ tel que $Z\cap f^{-1}(U)$ soit fermée dans $f^{-1}(U)$ et propre sur $U$ (\textbf{II}, 5.4.10) (ce qui revient à dire que pour tout sous-préschéma fermé de $X$ ayant $Z$ pour espace sous-jacent, la restriction $Z\to Y$ de $f$ à $Z$ est propre au point $y$).
\end{definition}

Soit $g:Y'\to Y$ un morphisme quelconque~; posons $X'=X\times_Y Y'$, $f'=f_{(Y')}$ et, si $p:X'\to X$ est la projection canonique, $Z'=p^{-1}(Z)$. Alors, si $Z$ est propre sur $Y$ au point $y$, $Z'$ est propre sur $Y'$ en tout point $y'$ au-dessus de $y$ (\textbf{II}, 5.4.2).

\begin{proposition}[15.7.2]
\label{IV.15.7.2-fr}
Soient $Y$ un préschéma intègre localement noethérien, $X$ un préschéma intègre, $f:X\to Y$ un morphisme séparé, dominant et de type fini, $y$ un point de $Y$. Les conditions suivantes sont équivalentes :
\begin{enumerate}[label=\emph{\alph*)}]
  \item $f$ est propre au point $y$.
  \item Si $Y_1=\operatorname{Spec}(\mathcal O_y)$, $X_1=X\times_Y Y_1$, le morphisme $f_1=f_{(Y_1)}:X_1\to Y_1$ est propre.
  \item Tout anneau de valuation discrète ayant $R(X)$ pour corps des fractions et dominant $\mathcal O$ domine un anneau local $\mathcal O_x$ de $X$ (auquel cas on a nécessairement $f(x)=y$).
\end{enumerate}
\end{proposition}

\oldpage[IV-3]{243}
\begin{proof}
Le fait que a) implique b) résulte de (\textbf{II}, 5.4.2), et l'implication b) $\Rightarrow$ c) résulte de (\textbf{II}, 7.3.10). Reste donc à montrer que c) entraîne a). La question étant locale sur $Y$, on peut supposer $Y$ affine, donc noethérien. En vertu du lemme de Chow (\textbf{II}, 5.6.1), il existe un préschéma intègre $X'$, un morphisme projectif $p:P\to Y$, une immersion ouverte dominante $j:X'\to P$, et un morphisme projectif et birationnel (donc surjectif) $g:X'\to X$ tels que le diagramme
\[
\xymatrix{
P\ar[d]_-{p} & X'\ar[l]_-{j}\ar[d]^-{g}\\
Y & X\ar[l]^-{f}
}
\]
soit commutatif. Comme $X'$ est intègre et $j$ dominant, $P$ est irréductible, et l'on peut, en remplaçant $P$ par $P_{\mathrm{red}}$, supposer $P$ intègre (\textbf{I}, 5.2.2 et \textbf{II}, 5.5.5, (vi)). Tout revient à prouver qu'il existe un voisinage ouvert $U$ de $y$ tel que $p^{-1}(U)\subset j(X')$, car alors la restriction de $j$ à
\[
j^{-1}\bigl(p^{-1}(U)\bigr)=g^{-1}\bigl(f^{-1}(U)\bigr)=V
\]
sera un isomorphisme sur $p^{-1}(U)$, donc la restriction $V\to U$ de $f\circ g$ sera propre, et comme la restriction $V\to f^{-1}(U)$ de $g$ est surjective, la restriction $f^{-1}(U)\to U$ de $f$ sera propre (\textbf{II}, 5.4.3). Comme $T=P-j(X')$ est fermé dans $P$, $p(T)$ est fermé dans $Y$, et il suffit de montrer que $y\notin p(T)$, ou encore que tout $z'\in P$ tel que $p(z')=y$ \emph{appartient à} $j(X')$. Or, le corps des fractions rationnelles $R(P)$ est égal à $R(X')$, donc à $R(X)$ par construction; comme on peut se borner au cas où $z'$ n'est pas le point générique de $P$, il existe un anneau de valuation discrète $A$ ayant $R(X)$ pour corps des fractions et dominant $\mathcal O_{z'}$ (\textbf{II}, 7.1.7); comme $p(z')=y$, $A$ domine aussi $\mathcal O_y$, donc par hypothèse il existe un $x\in X$ tel que $f(x)=y$ et que $A$ domine $\mathcal O_x$. Comme $g$ est propre, il existe $x'\in X'$ tel que $g(x')=x$ et que $A$ domine $\mathcal O_{x'}$ (\textbf{II}, 7.3.10); donc $z'$ et $j(x')$ sont deux points de $P$ dont les anneaux locaux $\mathcal O_z$ et $\mathcal O_{j(x')}=\mathcal O_{x'}$ sont apparentés (\textbf{I}, 8.1.4); comme $P$ est un schéma, cela entraîne $z'=j(x')$ (\textbf{I}, 8.2.2), ce qui achève la démonstration.
\end{proof}

\begin{remark}[15.7.3]
\label{IV.15.7.3-fr}
On peut dans (15.7.2) supprimer l'hypothèse que $Y$ est localement noethérien, en remplaçant dans la condition c) les anneaux de valuation discrète par des anneaux de valuation quelconque; la démonstration est alors inchangée, compte tenu de (\textbf{II}, 7.3.10).
\end{remark}

\begin{corollary}[15.7.4]
\label{IV.15.7.4-fr}
Soient $Y$ un préschéma noethérien, $f:X\to Y$ un morphisme séparé de type fini, $Z$ une partie de $X$ tel que les points maximaux de son adhérence $\overline Z$ appartiennent à $Z$ (ce qui est le cas lorsque $Z$ est fini, ou lorsque $Z$ est constructible ($\mathbf{0}_{\mathrm{III}}$, 9.2.2)). Soit $y$ un point de $Y$. Les conditions suivantes sont équivalentes :
\begin{enumerate}
  \item[\rm a)] L'ensemble $\overline Z$ est propre sur $Y$ au point $y$.
  \item[\rm b)] Si $Y_1=\operatorname{Spec}(\mathcal O_y)$, $X_1=X\times_Y Y_1$, et si $g_1:X_1\to X$ est la projection canonique et $Z_1=g_1^{-1}(Z)$, l'adhérence $\overline Z_1$ de $Z_1$ dans $X_1$ est propre sur $Y_1$.
  \item[\rm c)] Pour tout schéma $Y'=\operatorname{Spec}(A')$, où $A'$ est un anneau de valuation discrète, et tout morphisme $g:Y'\to Y$, tel que l'image $g(y')$ du point fermé $y'$ de $Y'$ soit $y$, on a la propriété suivante : posant $X'=X\times_Y Y'$, $f'=f_{(Y')}:X'\to Y'$, $g'=g_{(X)}:X'\to X$, $Z'=g'^{-1}(Z)$, et désignant par $\eta'$
\oldpage[IV-3]{244}
le point générique de $Y'$, alors, pour tout point $x'\in Z'\cap f'^{-1}(\eta')$ rationnel sur $\mathbf k(\eta')$, il existe une $Y'$-section de $X'$ contenant $x'$.
\end{enumerate}
\end{corollary}

\begin{proof}
Le fait que a) entraîne b) résulte de ce que $\overline Z_1\subset g_1^{-1}(\overline Z)$ et de ce que $g_1^{-1}(\overline Z)$ est propre sur $Y_1$ (15.7.1). Le fait que b) entraîne c) résulte de (\textbf{II}, 7.3.3). Reste donc à voir que c) implique a). Il suffit évidemment de prouver que toute composante irréductible de $\overline Z$ est propre sur $Y$ au point $y$, et l'hypothèse permet donc de se borner au cas où $Z$ est réduit à un seul point $z$. On peut évidemment aussi, compte tenu de (\textbf{I}, 5.2.2) et (\textbf{II}, 5.4.6), supposer que $X$ et $Y$ sont réduits, puis (par définition d'une partie propre (\textbf{II}, 5.4.10)) supposer que $\overline Z=X=\overline{\{z\}}$ et (appliquant encore (\textbf{I}, 5.2.2)) que $f$ est dominant, donc $X$ et $Y$ intègres et noethériens; il faut alors prouver que $f$ est propre au point $y$ de $Y$. Montrons pour cela que $f$ vérifie la condition c) de (15.7.2). Soit $A'$ un anneau de valuation discrète ayant $R(X)=\mathbf k(z)$ pour corps des fractions et dominant $\mathcal O_y$; avec les notations de c), on a donc $g(y')=y$, et $g(\eta')$ est le point générique $\eta=f(z)$ de $Y$. Comme par définition $\mathbf k(\eta')=\mathbf k(z)$, il existe un $\mathbf k(\eta)$-morphisme $\operatorname{Spec}(\mathbf k(\eta'))\to\operatorname{Spec}(\mathbf k(z))$ transformant $\eta'$ en $z$, donc une $\mathbf k(\eta')$-section de $f'^{-1}(\eta')$ (\textbf{I}, 3.3.14); si $z'$ est l'image de $\eta'$ par cette section, $z'$ est donc rationnel sur $\mathbf k(\eta')$ et $g'(z')=z$. On conclut donc de l'hypothèse c) qu'il existe une $Y'$-section $h'$ de $X'$ telle que $h'(\eta')=z'$; soit $h=g'\circ h':Y'\to X$ le $Y$-morphisme correspondant, et soit $x=h(y')$; il est clair que $f(x)=y$ et que $A'=\mathcal O_{y'}$ domine $\mathcal O_x$, ce qui achève la démonstration.
\end{proof}

\begin{corollary}[15.7.5]
\label{IV.15.7.5-fr}
Soient $Y$ un préschéma localement noethérien, $f:X\to Y$ un morphisme séparé de type fini, $y$ un point de $Y$. Les conditions suivantes sont équivalentes :
\begin{enumerate}
  \item[\rm a)] $f$ est propre au point $y$.
  \item[\rm b)] Si $Y_1=\operatorname{Spec}(\mathcal O_y)$, $X_1=X\times_Y Y_1$, le morphisme $f_1=f_{(Y_1)}:X_1\to Y_1$ est propre.
  \item[\rm c)] Pour tout schéma $Y'=\operatorname{Spec}(A')$, où $A'$ est un anneau de valuation discrète, et tout morphisme $g:Y'\to Y$, tel que l'image $g(y')$ du point fermé $y'$ de $Y'$ soit $y$, on a la propriété suivante : posant $X'=X\times_Y Y'$, $f'=f_{(Y')}:X'\to Y'$ et désignant par $\eta'$ le point générique de $Y'$, alors, pour tout $x'\in f'^{-1}(\eta')$, rationnel sur $\mathbf k(\eta')$, il existe une $Y'$-section de $X'$ contenant $x'$.
\end{enumerate}
\end{corollary}

\begin{corollary}[15.7.6]
\label{IV.15.7.6-fr}
Soient $Y$ un préschéma localement noethérien, $f:X\to Y$ un morphisme de type fini. Les conditions suivantes sont équivalentes :
\begin{enumerate}
  \item[\rm a)] Il existe un morphisme propre $g:X\to Z$ et une immersion $h:Z\to Y$ tels que $f=h\circ g$.
  \item[\rm b)] Le sous-espace $f(X)$ est localement fermé dans $Y$, et si $Z'$ est le sous-préschéma de $Y$ image fermée de $X$ par $f$ (\textbf{I}, 9.5.3 et 9.5.1), et $Z''$ le préschéma induit par $Z'$ sur l'ouvert $f(X)$ de $Z'$, de sorte que $f$ se factorise de façon unique en $f=j\circ g$ (où $j:Z''\to Y$ est l'injection canonique), alors $g:X\to Z''$ est propre.
  \item[\rm c)] Pour tout $y\in f(X)$, $f$ est propre au point $y$.
  \item[\rm d)] Pour tout schéma $Y'=\operatorname{Spec}(A')$, où $A'$ est un anneau de valuation discrète, et tout morphisme $g:Y'\to Y$ tel que $g(Y')\subset f(X)$, toute $Y'$-section rationnelle (\textbf{I}, 7.1.2) de $X'=X\times_Y Y'$ se prolonge de façon unique en une $Y'$-section de $X'$.
\end{enumerate}
\end{corollary}

\begin{proof}
Il est clair que b) entraîne a). Pour voir que a) entraîne c), remarquons d'abord que a) entraîne que $h(Z)$ est localement fermé dans $Y$ et $g(X)$ fermé dans $Z$, donc $f(X)$ est localement fermé dans $Y$; pour tout $y\in f(X)$, il y a par suite un voisinage ouvert $U$
\oldpage[IV-3]{245}
de $y$ dans $Y$ tel que $f(X)\cap U$ soit fermé dans $U$; alors la restriction $h^{-1}(U)\to U$ de $h$ est une immersion fermée, et la restriction $f^{-1}(U)=g^{-1}(h^{-1}(U))\to h^{-1}(U)$ de $g$ est propre, donc la restriction $f^{-1}(U)\to U$ de $f$ est propre (\textbf{II}, 5.4.2). Pour voir que c) entraîne b), notons d'abord que, pour tout $y\in f(X)$, il existe, en vertu de c), un voisinage ouvert $U$ de $y$ dans $Y$ tel que $f(X)\cap U$ soit fermé dans $U$, donc $f(X)$ est localement fermé, et par suite ouvert dans son adhérence. Comme cette dernière est l'espace sous-jacent de $Z'$, et que $f$ se factorise en $f=j_0\circ g_0$, où $j_0:Z'\to Y$ est l'injection canonique et $g_0$ est un morphisme de $X$ dans $Z'$, le fait que $g_0(X)=Z''$ (qui est ouvert dans $Z'$) entraîne la factorisation de l'énoncé de b); le fait que $g$ est propre résulte alors du caractère local (sur $Z''$) de cette propriété, et de (\textbf{II}, 5.4.3). Il est clair que c) implique d) en vertu de (15.7.5), l'unicité du prolongement d'une $Y'$-section rationnelle résultant de l'hypothèse que $f$ est séparé (\textbf{I}, 7.2.2). Prouvons enfin que d) entraîne c); en premier lieu, il résulte de d), compte tenu de (\textbf{II}, 7.2.3 et 7.2.4), que $f$ est séparé; on peut alors appliquer le critère (15.7.5, c)), qui montre que $f$ est propre en tout point de $f(X)$.
\end{proof}

\begin{remarks}[15.7.7]
\label{IV.15.7.7-fr}
(i) Dans (15.7.4, c)), (15.7.5, c)) et (15.7.6, d)), on peut se restreindre au cas où l'anneau de valuation discrète $A'$ est \emph{complet} et admet un corps résiduel \emph{algébriquement clos}. En effet, dans la démonstration de (15.7.4), si l'on considère un anneau de valuation discrète $A''$ complet dominant $A'$ et ayant un corps résiduel algébriquement clos ($\mathbf{0}_{\mathrm{III}}$, 10.3.1), l'hypothèse c) est alors vérifiée pour $Y''=\operatorname{Spec}(A'')$, $X''=X\times_Y Y''$ et $f''=f_{(Y'')}$; mais $X''=X'\times_{Y'}Y''$, et il résulte alors de la remarque (\textbf{II}, 7.3.9, (i)) que $f'$ est propre; par suite $Y'$, $X'$ et $f'$ vérifient aussi l'hypothèse c) (\textbf{II}, 7.3.8).

(ii) Compte tenu de (\textbf{II}, 7.3.8), la condition c) de (15.7.5) peut se remplacer par la condition que $f'$ est \emph{propre}.
\end{remarks}

\begin{env}[15.7.8]
\label{IV.15.7.8-fr}
On dit qu'un morphisme de préschémas $f:X\to Y$ est \emph{submersif} s'il est surjectif et si la topologie de $Y$ est égale au quotient de la topologie de $X$ par la relation d'équivalence définie par $f$. On dit que $f$ est \emph{universellement submersif} si pour tout changement de base $Y'\to Y$, le morphisme $f'=f_{(Y')}:X\times_Y Y'\to Y'$ est submersif. Tout morphisme surjectif \emph{universellement ouvert}, ou \emph{universellement fermé}, ou \emph{fidèlement plat et quasi-compact} est universellement submersif (2.3.12). Étant donné un morphisme $f:X\to Y$, on dit qu'un sous-ensemble $Z$ de $X$ est \emph{submersif sur $f(Z)$} si la topologie du sous-espace $f(Z)$ de $Y$ est quotient de la topologie du sous-espace $Z$ de $X$ par la relation d'équivalence définie par $f|Z$. On dit que $Z$ est \emph{universellement submersif sur $f(Z)$} si, pour tout changement de base $g:Y'\to Y$, l'image réciproque $Z'=p^{-1}(Z)$ de $Z$ par la projection $p:X\times_Y Y'\to Y'$ est un ensemble submersif sur $f'(Z')=g^{-1}(f(Z))$. On notera que si le morphisme $f$ admet une $Y$-section $h:Y\to X$, tout ensemble $Z$ contenant $h(Y)$ est \emph{universellement submersif sur $Y$}.
\end{env}

\begin{proposition}[15.7.9]
\label{IV.15.7.9-fr}
Soient $Y$ un préschéma localement noethérien, $y$ un point de $Y$, $f:X\to Y$ un morphisme séparé de type fini. Soit $Z$ une partie localement constructible de $X$ ayant les propriétés suivantes :
\begin{enumerate}
  \item[\rm (i)] Pour toute générisation $y'$ de $y$, $Z_{y'}=Z\cap f^{-1}(y')$ est une composante connexe de $f^{-1}(y')$ et est géométriquement connexe sur $\mathbf k(y')$ (4.5.2).
\oldpage[IV-3]{246}
  \item[\rm (ii)] $Z_y$ est propre sur $\operatorname{Spec}(\mathbf k(y))$.
  \item[\rm (iii)] $Z$ est universellement submersif sur $f(Z)$ (15.7.8).
\end{enumerate}
Alors $Z$ est propre sur $Y$ au point $y$ (autrement dit (15.7.1) il existe un voisinage ouvert $U$ de $y$ tel que $Z\cap f^{-1}(U)$ soit fermé dans $f^{-1}(U)$ et propre sur $U$).
\end{proposition}

\begin{proof}
Utilisant la méthode de (8.1.2, a)) et (8.10.5, (xii)), on est ramené à prouver que lorsque $Y=\operatorname{Spec}(A)$ est spectre d'un anneau \emph{local} noethérien, les hypothèses (i), (ii) et (iii) entraînent que $Z$ est \emph{propre} sur $Y$.

Notons d'abord que si $A'$ est un anneau local noethérien, $\varphi:A\to A'$ un homomorphisme local, $Y'=\operatorname{Spec}(A')$ et $g:Y'\to Y$ le morphisme correspondant à $\varphi$, l'image réciproque $Z'$ de $Z$ par la projection canonique $p:X\times_Y Y'\to X$ a vis-à-vis de $Y'$ les propriétés (i), (ii), (iii) de l'énoncé, compte tenu de (1.8.2); utilisant (2.7.1, (vii)), il revient donc au même de démontrer que $Z'$ est propre sur $Y'$, pourvu que $A'$ soit un $A$-module \emph{fidèlement plat}. Nous allons utiliser cette remarque en prenant $A'=\widehat A$ ($\mathbf{0}_{\mathrm{I}}$, 7.3.5), autrement dit nous pouvons nous borner au cas où $A$ est \emph{complet}. Comme $Z_y$ est une composante connexe de $f^{-1}(y)$, propre sur $\mathbf k(y)$, il résulte de (\textbf{III}, 5.5.2) que $X$ est \emph{somme} de deux sous-préschémas $X_0$, $X_1$ induits sur des parties ouvertes et fermées de $X$, telles que $X_0$ soit propre sur $Y$ et que $X_0\cap f^{-1}(y)=Z_y$. Posons $Z_0=X_0\cap Z$, $Z_1=X_1\cap Z$, de sorte que ces ensembles forment une partition de $Z$ en deux parties ouvertes et fermées dans $Z$. Comme les $Z_{y'}$ sont connexes, on a nécessairement $Z_{y'}\subset Z_0$ ou $Z_{y'}\subset Z_1$, donc $f(Z_0)\cap f(Z_1)=\varnothing$. Mais comme $Z_0$ et $Z_1$ sont saturés pour la relation d'équivalence définie par $f|Z$, il résulte de (iii) que $f(Z_0)$ et $f(Z_1)$ sont ouverts dans $f(Z)$. Mais $f(Z_0)$ contient $y$, et tout voisinage de $y$ dans $Y$ est égal à $Y$, donc $f(Z_0)=f(Z)$, et par suite $f(Z_1)=\varnothing$, donc $Z_1=\varnothing$ et $Z\subset X_0$.

On peut donc désormais supposer de plus que $f$ est \emph{propre}, et il reste à prouver que $Z$ est fermé dans $X$. Autrement dit, il suffira de prouver qu'une partie \emph{constructible} $Z$ d'un préschéma $X$ \emph{propre} sur $Y$ qui satisfait aux \emph{deux seules conditions} (i) et (iii), est \emph{fermée} dans $X$. En vertu de ($\mathbf{0}_{\mathrm{III}}$, 9.2.5), il suffit donc de prouver que pour tout $x'\in Z$ et toute spécialisation $x$ de $x'$ dans $X$, on a $x\in Z$. Soient $A_1$ un anneau de valuation discrète \emph{complet}, $u$ un morphisme de $Y_1=\operatorname{Spec}(A_1)$ dans $X$ tel que, si $y_1$ et $y_1'$ sont le point fermé et le point générique de $Y_1$, on ait $u(y_1)=x$, $u(y_1')=x'$ (\textbf{II}, 7.1.7); posons $g=f\circ u:Y_1\to Y$ et $X_1=X\times_Y Y_1$; il existe une $Y_1$-section $u_1:Y_1\to X_1$ de $f_1=f_{(Y_1)}$ telle que $u=p\circ u_1$, où $p:X_1\to X$ est la projection canonique. Si $x_1=u_1(y_1)$, $x_1'=u_1(y_1')$, on a donc $p(x_1)=x$ et $p(x_1')=x'$, et $x_1$ est une spécialisation de $x_1'$. En outre, on a $x_1'\in Z_1=p^{-1}(Z)$; comme nous avons remarqué que les conditions (i) et (iii) sont stables pour \emph{tout} changement de base, ainsi que la propriété d'être constructible, on voit qu'on est ramené à la situation suivante : $Y$ est le spectre d'un anneau de valuation discrète complet, $X$ est propre sur $Y$, $f(x)=y$ est le point fermé de $Y$, $y'=f(x')$ son point générique, il existe une $Y$-section $h$ de $X$ telle que $h(y)=x$, $h(y')=x'$, et enfin on a $x'\in Z$; il faut prouver que $x\in Z$. Or, $Z_y$ est une composante connexe de $f^{-1}(y)$, \emph{propre} sur $\operatorname{Spec}(\mathbf k(y))$ puisque $X$ est propre sur $Y$. Appliquant de nouveau (\textbf{III}, 5.5.2), on voit comme plus haut qu'il y a une partie ouverte et fermée $X'$ de $X$ telle que $X'\cap f^{-1}(y)=Z_y$; comme $Z_y$ est connexe,
\oldpage[IV-3]{247}
on peut supposer $X'$ connexe (en remplaçant au besoin $X'$ par celle de ses composantes connexes contenant $Z_y$). Mais alors le même raisonnement qu'au début prouve qu'on a nécessairement $Z\subset X'$; comme $h(Y)$ est connexe et contient $x'$, il est nécessairement contenu dans $X'$, donc $x=h(y)\in X'\cap f^{-1}(y)=Z_y$. C.Q.F.D.
\end{proof}

\begin{corollary}[15.7.10]
\label{IV.15.7.10-fr}
Soient $Y$ un préschéma localement noethérien, $f:X\to Y$ un morphisme séparé de type fini, universellement submersif (15.7.8) et tel que toutes les fibres $X_y=f^{-1}(y)$ soient géométriquement connexes. Alors, si $y\in Y$ est tel que $X_y$ soit propre sur $\mathbf k(y)$, $f$ est propre au point $y$.
\end{corollary}

\begin{proof}
Compte tenu de (15.7.5), on est ramené au cas où $Y=\operatorname{Spec}(\mathcal O_y)$ puisque les hypothèses sont stables par changement de base. Mais dans ce cas il suffit d'appliquer (15.7.9) à $Z=X$.
\end{proof}

\begin{corollary}[15.7.11]
\label{IV.15.7.11-fr}
Soit $f:X\to Y$ un morphisme séparé de présentation finie; supposons qu'il existe un morphisme surjectif de présentation finie $g:Y'\to Y$, qui soit en outre propre ou plat, et tel que si l'on pose $X'=X\times_Y Y'$, il existe une $Y'$-section de $X'$ (ou encore un $Y$-morphisme $Y'\to X$ (\textbf{I}, 3.3.14)). Supposons enfin que les fibres $X_y=f^{-1}(y)$ soient géométriquement connexes. Alors l'ensemble $U$ des $y\in Y$ tels que $X_y$ soit propre sur $\mathbf k(y)$ est ouvert dans $Y$, et la restriction $f^{-1}(U)\to U$ de $f$ est propre.
\end{corollary}

\begin{proof}
La question étant locale sur $Y$, on peut supposer $Y=\operatorname{Spec}(A)$ affine. Appliquant (8.9.1) à $f$ et à $g$, on voit qu'il existe un sous-anneau noethérien $A_0$ de $A$ et deux morphismes de type fini $f_0:X_0\to Y_0=\operatorname{Spec}(A_0)$, $g_0:Y_0'\to Y_0$ tels que $X=X_0\times_{Y_0}Y$, $Y'=Y_0'\times_{Y_0}Y$, $f=(f_0)_{(Y)}$ et $g=(g_0)_{(Y)}$; en outre, utilisant (8.10.5, (v), (vi) et (xii)) et (11.2.7), on peut supposer $f_0$ séparé, $g_0$ surjectif et propre (resp. plat) si $g$ est propre (resp. plat); utilisant (8.8.2), on peut de plus supposer qu'il existe une $Y_0'$-section de $X_0$. D'autre part, utilisant (9.7.7), on peut appliquer (9.3.3) à la propriété d'être géométriquement connexe, et on peut donc supposer $A_0$ choisi de sorte que les $f_0^{-1}(y_0)$ soient géométriquement connexes pour tout $y_0\in Y_0$. Enfin, utilisant (2.7.1, (vii)), il revient au même de dire que $f^{-1}(y)$ est propre sur $\mathbf k(y)$ ou que $f_0^{-1}(y_0)$ est propre sur $\mathbf k(y_0)$, où $y_0$ est la projection de $y$ dans $Y_0$. Ces remarques montrent qu'on est ramené à démontrer le corollaire lorsque $Y$ est \emph{noethérien}. Notons maintenant le

\begin{lemma}[15.7.11.1]
\label{IV.15.7.11.1-fr}
Soient $f:X\to Y$ un morphisme, $g:Y'\to Y$ un morphisme surjectif, qui est en outre propre ou plat et quasi-compact. Posons $X'=X\times_Y Y'$, $f'=f_{(Y')}:X'\to Y'$. Alors, si $f'$ est submersif (resp. universellement submersif), il en est de même de $f$.
\end{lemma}

\begin{proof}
On peut se borner au cas où $f'$ est submersif. En effet, supposons le lemme prouvé dans ce cas, et montrons que si $f'$ est universellement submersif, $f$ l'est aussi. Soit donc $Y_1\to Y$ un morphisme quelconque, et posons $Y_1'=Y'\times_Y Y_1$; si $X_1'=X'\times_{Y'}Y_1'$ et $f_1'=(f')_{(Y_1')}:X_1'\to Y_1'$, $f_1'$ est submersif par hypothèse; de plus le morphisme $g_1:Y_1'\to Y_1$ est surjectif, et propre (resp. plat et quasi-compact) si $g$ l'est. Donc, si l'on pose $X_1=X\times_Y Y_1$, $f_1=f_{(Y_1)}$, il résulte de l'hypothèse et du fait que $X_1'=X_1\times_{Y_1}Y_1'$, que $f_1$ est submersif, donc $f$ universellement submersif. Supposons donc seulement que $f'$ soit submersif. Il est clair tout d'abord que $f$ est surjectif. Soit $E$ une partie de $Y$ telle que $f^{-1}(E)$ soit fermé dans $X$; alors, si $p:X'\to X$ est la projection canonique,
\oldpage[IV-3]{248}
$p^{-1}(f^{-1}(E))=f'^{-1}(g^{-1}(E))$ est fermé dans $X'$, donc, puisque $f'$ est submersif, $g^{-1}(E)$ est fermé dans $Y'$; mais comme $g$ est aussi universellement submersif (15.7.8), $E$ est fermé dans $Y$, ce qui prouve le lemme.
\end{proof}

Cela étant, comme il existe par hypothèse une $Y'$-section de $X'$, $f'$ est universellement submersif (15.7.8), donc il en est de même de $f$ d'après le lemme (15.7.11.1). Il suffit alors d'appliquer (15.7.10), en remarquant que l'ensemble des $y\in Y$ tels que $f$ soit propre en $y$ est par définition ouvert dans $Y$.
\end{proof}

\begin{center}
\emph{(A suivre.)}
\end{center}

\clearpage
\oldpage[IV-3]{249}
\phantomsection
\addcontentsline{toc}{section}{Bibliographie (suite)}
\section*{BIBLIOGRAPHIE (\emph{suite})}

\begingroup
\small
\setlength{\parindent}{0pt}
\setlength{\parskip}{3pt}
[39] J.-P.~\textsc{Serre}, \emph{Corps locaux}, Paris (Hermann), 1962.

[40] J.-P.~\textsc{Serre}, \emph{Groupes proalgébriques}, \emph{Publ. Inst. Hautes Études Sci.}, n\textsuperscript{o} 7 (1960).

\endgroup

\clearpage
\oldpage[IV-3]{250}
\phantomsection
\addcontentsline{toc}{section}{Index des notations}
\section*{INDEX DES NOTATIONS}

\begingroup
\small
\setlength{\parindent}{0pt}
\setlength{\parskip}{2pt}

$\mathfrak C(X)$, $\mathfrak D_c(X)$, $\mathfrak F_c(X)$,
$\mathfrak{LF}_c(X)$ ($X$ préschéma) : 8.3.9.

$\mathfrak Q(\mathcal F)$ ($\mathcal F$ $\mathcal O_X$-Module) : 8.5.10.

$\mathfrak{Spr}(Y)$, $\mathfrak{Spro}(Y)$, $\mathfrak{Sprf}(Y)$
($Y$ préschéma) : 8.6.1.

$\mathbf{Pro}(\mathcal C)$ ($\mathcal C$ catégorie) : 8.13.3.

$X_s$, $\mathcal F_s$, $u_s$, $g_s$, $Z_s$, $h_s$
($X$ $S$-préschéma, $s\in S$, $\mathcal F$ $\mathcal O_X$-Module,
$u$ homomorphisme de $\mathcal O_X$-Modules, $g$ section de
$\mathcal O_X$-Module, $Z$ partie de $X$, $h:X\to Y$ $S$-morphisme) : 9.4.1.

$\mathfrak{Qc}(X)$, $\mathfrak{Lqc}(X)$, $\mathfrak O(X)$,
$\mathfrak F(X)$, $\mathfrak{Lf}(X)$, $\mathfrak{Lc}(X)$
($X$ espace topologique) : 10.1.1.

$\operatorname{prof}^{*}_{x}(\mathcal F)$,
$\operatorname{prof}^{*}_{Z}(\mathcal F)$,
$\operatorname{prof}^{*}(\mathcal F)$
($\mathcal F$ $\mathcal O_X$-Module) : 10.8.1.

$S(X)$ ($X$ préschéma de Jacobson) : 10.9.1.

$S(f)$ ($f$ morphisme de préschémas de Jacobson) : 10.9.2.

$\operatorname{Spm}(A)$, $D^m(h)$ ($A$ anneau de Jacobson, $h\in A$) : 10.9.3.

$\operatorname{Tor.dim}_B(M)$ ($M$ $B$-module) : 12.3.2.

\endgroup

\clearpage
\oldpage[IV-3]{251}
\phantomsection
\addcontentsline{toc}{section}{Index terminologique}
\section*{INDEX TERMINOLOGIQUE}

\begingroup
\footnotesize
\setlength{\parindent}{0pt}
\setlength{\parskip}{1.1pt}
\newcommand{\EGAIVthreeTerm}[2]{#1 : #2.\par}

\EGAIVthreeTerm{Critère de platitude par fibres}{11.3.10}
\EGAIVthreeTerm{Espace algébrique sur $k$, espace $k$-algébrique, espace
préalgébrique sur $k$, espace $k$-préalgébrique}{10.10.2}
\EGAIVthreeTerm{Espace de Jacobson}{10.3.1}
\EGAIVthreeTerm{Espace $k$-préalgébrique réduit}{10.10.4}
\EGAIVthreeTerm{Famille de morphismes schématiquement dominante}{11.10.2}
\EGAIVthreeTerm{Famille de morphismes universellement schématiquement dominante}{11.10.7}
\EGAIVthreeTerm{Famille de sous-préschémas schématiquement dense}{11.10.2}
\EGAIVthreeTerm{Famille séparante d'homomorphismes de faisceaux de modules}{11.9.1}
\EGAIVthreeTerm{Famille séparante d'homomorphismes de modules}{11.9.4}
\EGAIVthreeTerm{Famille universellement séparante d'homomorphismes de faisceaux de modules}{11.9.14}
\EGAIVthreeTerm{Géométriquement isomorphes ($k$-préschémas)}{9.1.4}
\EGAIVthreeTerm{Homomorphisme de Modules universellement injectif}{11.9.14 et 11.9.18}
\EGAIVthreeTerm{Lemme de Chow pour les morphismes de présentation finie}{8.10.5.1}
\EGAIVthreeTerm{« Main theorem » de Zariski pour les morphismes de présentation finie}{8.12.6}
\EGAIVthreeTerm{Morphisme équidimensionnel en un point, morphisme équidimensionnel}{13.2.2 et 13.3.2}
\EGAIVthreeTerm{Morphisme propre en un point}{15.7.1}
\EGAIVthreeTerm{Morphisme pseudo-fini}{8.12.3}
\EGAIVthreeTerm{Morphisme submersif, morphisme universellement submersif}{15.7.8}
\EGAIVthreeTerm{Morphisme universellement ouvert en un point}{14.3.3}
\EGAIVthreeTerm{Ouvert ultraaffine}{10.9.5}
\EGAIVthreeTerm{Partie finie saturée d'un $k$-préschéma algébrique}{9.8.9}
\EGAIVthreeTerm{Partie propre en un point (d'un $S$-préschéma)}{15.7.1}
\EGAIVthreeTerm{Partie quasi-constructible, partie localement quasi-constructible d'un espace topologique}{10.1.1}
\EGAIVthreeTerm{Partie très dense d'un espace topologique}{10.1.3}
\EGAIVthreeTerm{Polynôme géométriquement irréductible}{9.7.4}
\EGAIVthreeTerm{Préschéma de Jacobson}{10.4.1}
\EGAIVthreeTerm{Préschéma équidimensionnel sur un autre en un point, équidimensionnel sur un autre}{13.2.2 et 13.3.2}
\EGAIVthreeTerm{Préschéma essentiellement affine sur un autre}{8.13.4}
\EGAIVthreeTerm{Principe de l'extension finie}{9.1.1}
\EGAIVthreeTerm{Profondeur rectifiée}{10.8.1}
\EGAIVthreeTerm{Pro-objet, morphisme de pro-objets}{8.13.3}
\EGAIVthreeTerm{Pro-objet essentiellement affine}{8.13.4, 8.14.1}
\EGAIVthreeTerm{Propriété constructible, propriété ind-constructible}{9.2.1 et 9.2.2, (i) et (ii)}
\EGAIVthreeTerm{Quasi-homéomorphisme d'espaces topologiques}{10.2.2}
\EGAIVthreeTerm{Quasi-isomorphisme d'espaces annelés}{10.2.8}
\EGAIVthreeTerm{Saturée d'une partie finie d'un $k$-préschéma algébrique}{9.8.9}
\EGAIVthreeTerm{Sous-ensemble d'un $S$-préschéma submersif, universellement submersif sur son image}{15.7.8}
\EGAIVthreeTerm{Squelette primaire d'un Module, squelette virtuel}{9.8.9}
\EGAIVthreeTerm{Spectre maximal d'un anneau de Jacobson}{10.9.3}
\EGAIVthreeTerm{Type primaire d'un Module, type primaire géométrique d'un Module}{9.8.9}
\EGAIVthreeTerm{Ultrapréschéma, morphisme d'ultrapréschémas}{10.9.5}
\EGAIVthreeTerm{$R$-ultrapréschéma}{10.10.2}

\let\EGAIVthreeTerm\relax
\endgroup

\clearpage
% La page imprimée 252 est blanche.
\thispagestyle{empty}
\null

\clearpage
\oldpage[IV-3]{253}
\phantomsection
\addcontentsline{toc}{section}{Table des matières}
\section*{TABLE DES MATIÈRES}

\begingroup
\small
\setlength{\parindent}{0pt}
\setlength{\parskip}{1.5pt}
\newcommand{\EGAIVthreeContentsLine}[3]{\textbf{#1}\quad #2\unskip\nobreak\hspace{0.4em}\nobreak\dotfill\nobreak\mbox{#3}\par}

\EGAIVthreeContentsLine{CHAPITRE IV.}{\textbf{Étude locale des schémas et des morphismes de schémas} \emph{(suite)}}{5}
\EGAIVthreeContentsLine{\S~8.}{Limites projectives de préschémas}{5}
\EGAIVthreeContentsLine{8.1.}{Introduction}{5}
\EGAIVthreeContentsLine{8.2.}{Limites projectives de préschémas}{7}
\EGAIVthreeContentsLine{8.3.}{Parties constructibles dans une limite projective de préschémas}{12}
\EGAIVthreeContentsLine{8.4.}{Critères d'irréductibilité et de connexion pour les limites projectives de préschémas}{17}
\EGAIVthreeContentsLine{8.5.}{Modules de présentation finie sur une limite projective de préschémas}{19}
\EGAIVthreeContentsLine{8.6.}{Sous-préschémas de présentation finie d'une limite projective de préschémas}{25}
\EGAIVthreeContentsLine{8.7.}{Critères pour qu'une limite projective de préschémas soit un préschéma réduit (resp. intègre)}{27}
\EGAIVthreeContentsLine{8.8.}{Préschémas de présentation finie sur une limite projective de préschémas}{28}
\EGAIVthreeContentsLine{8.9.}{Premières applications à l'élimination des hypothèses noethériennes}{34}
\EGAIVthreeContentsLine{8.10.}{Propriétés de permanence des morphismes par passage à la limite projective}{36}
\EGAIVthreeContentsLine{8.11.}{Application aux morphismes quasi-finis}{41}
\EGAIVthreeContentsLine{8.12.}{Nouvelle démonstration et généralisation du « Main Theorem » de Zariski}{43}
\EGAIVthreeContentsLine{8.13.}{Traduction en termes de pro-objets}{49}
\EGAIVthreeContentsLine{8.14.}{Caractérisation d'un préschéma localement de présentation finie sur un autre, en termes du foncteur qu'il représente}{52}

\EGAIVthreeContentsLine{\S~9.}{Propriétés constructibles}{54}
\EGAIVthreeContentsLine{9.1.}{Le principe de l'extension finie}{54}
\EGAIVthreeContentsLine{9.2.}{Propriétés constructibles et ind-constructibles}{56}
\EGAIVthreeContentsLine{9.3.}{Propriétés constructibles de morphismes de préschémas algébriques}{60}

\clearpage
\oldpage[IV-3]{254}
\EGAIVthreeContentsLine{9.4.}{Constructibilité de certaines propriétés des Modules}{62}
\EGAIVthreeContentsLine{9.5.}{Constructibilité de propriétés topologiques}{67}
\EGAIVthreeContentsLine{9.6.}{Constructibilité de certaines propriétés des morphismes}{71}
\EGAIVthreeContentsLine{9.7.}{Constructibilité des propriétés de séparabilité, d'irréductibilité géométrique et de connexité géométrique}{76}
\EGAIVthreeContentsLine{9.8.}{Décomposition primaire au voisinage d'une fibre générique}{83}
\EGAIVthreeContentsLine{9.9.}{Constructibilité des propriétés locales des fibres}{88}

\EGAIVthreeContentsLine{\S~10.}{Préschémas de Jacobson}{95}
\EGAIVthreeContentsLine{10.1.}{Parties très denses d'un espace topologique}{95}
\EGAIVthreeContentsLine{10.2.}{Quasi-homéomorphismes}{97}
\EGAIVthreeContentsLine{10.3.}{Espaces de Jacobson}{101}
\EGAIVthreeContentsLine{10.4.}{Préschémas de Jacobson et anneaux de Jacobson}{101}
\EGAIVthreeContentsLine{10.5.}{Préschémas de Jacobson noethériens}{104}
\EGAIVthreeContentsLine{10.6.}{Dimension dans les préschémas de Jacobson}{107}
\EGAIVthreeContentsLine{10.7.}{Exemples et contre-exemples}{109}
\EGAIVthreeContentsLine{10.8.}{Profondeur rectifiée}{110}
\EGAIVthreeContentsLine{10.9.}{Spectres maximaux et ultrapréschémas}{112}
\EGAIVthreeContentsLine{10.10.}{Espaces algébriques de Serre}{114}

\EGAIVthreeContentsLine{\S~11.}{Propriétés topologiques des morphismes plats de présentation finie. Critères de platitude}{116}
\EGAIVthreeContentsLine{11.1.}{Ensembles de platitude (cas noethérien)}{117}
\EGAIVthreeContentsLine{11.2.}{Platitude d'une limite projective de préschémas}{119}
\EGAIVthreeContentsLine{11.3.}{Application à l'élimination d'hypothèses noethériennes}{132}
\EGAIVthreeContentsLine{11.4.}{Descente de la platitude par des morphismes quelconques : cas d'un préschéma de base artinien}{143}
\EGAIVthreeContentsLine{11.5.}{Descente de la platitude par des morphismes quelconques : cas général}{150}
\EGAIVthreeContentsLine{11.6.}{Descente de la platitude par des morphismes quelconques : cas d'un préschéma de base unibranche}{154}
\EGAIVthreeContentsLine{11.7.}{Contre-exemples}{157}
\EGAIVthreeContentsLine{11.8.}{Un critère valuatif de platitude}{159}
\EGAIVthreeContentsLine{11.9.}{Familles séparantes et universellement séparantes d'homomorphismes de faisceaux de modules}{160}
\EGAIVthreeContentsLine{11.10.}{Familles schématiquement dominantes de morphismes et familles schématiquement denses de sous-préschémas}{170}

\EGAIVthreeContentsLine{\S~12.}{Étude des fibres des morphismes plats de présentation finie}{173}
\EGAIVthreeContentsLine{12.0.}{Introduction}{173}
\EGAIVthreeContentsLine{12.1.}{Propriétés locales des fibres d'un morphisme plat localement de présentation finie}{174}

\clearpage
\oldpage[IV-3]{255}
\EGAIVthreeContentsLine{12.2.}{Propriétés locales et globales des fibres d'un morphisme propre, plat et de présentation finie}{179}
\EGAIVthreeContentsLine{12.3.}{Propriétés cohomologiques locales des fibres d'un morphisme plat et localement de présentation finie}{183}

\EGAIVthreeContentsLine{\S~13.}{Morphismes équidimensionnels}{187}
\EGAIVthreeContentsLine{13.1.}{Le théorème de semi-continuité de Chevalley}{188}
\EGAIVthreeContentsLine{13.2.}{Morphismes équidimensionnels : cas des morphismes dominants de préschémas irréductibles}{190}
\EGAIVthreeContentsLine{13.3.}{Morphismes équidimensionnels : cas général}{194}

\EGAIVthreeContentsLine{\S~14.}{Morphismes universellement ouverts}{199}
\EGAIVthreeContentsLine{14.1.}{Morphismes ouverts}{200}
\EGAIVthreeContentsLine{14.2.}{Morphismes ouverts et formule des dimensions}{202}
\EGAIVthreeContentsLine{14.3.}{Morphismes universellement ouverts}{204}
\EGAIVthreeContentsLine{14.4.}{Le critère de Chevalley pour les morphismes universellement ouverts}{209}
\EGAIVthreeContentsLine{14.5.}{Morphismes universellement ouverts et quasi-sections}{216}

\EGAIVthreeContentsLine{\S~15.}{Étude des fibres d'un morphisme universellement ouvert}{223}
\EGAIVthreeContentsLine{15.1.}{Multiplicités des fibres d'un morphisme universellement ouvert}{223}
\EGAIVthreeContentsLine{15.2.}{Platitude des morphismes universellement ouverts à fibres géométriquement réduites}{226}
\EGAIVthreeContentsLine{15.3.}{Application : critères de réduction et d'irréductibilité}{228}
\EGAIVthreeContentsLine{15.4.}{Compléments sur les morphismes de Cohen-Macaulay}{229}
\EGAIVthreeContentsLine{15.5.}{Rang séparable des fibres d'un morphisme quasi-fini et universellement ouvert. Application aux composantes connexes géométriques des fibres d'un morphisme propre}{231}
\EGAIVthreeContentsLine{15.6.}{Composantes connexes des fibres le long d'une section}{236}
\EGAIVthreeContentsLine{15.7.}{Appendice : Critères valuatifs de propreté locale}{242}
\EGAIVthreeContentsLine{}{BIBLIOGRAPHIE \emph{(suite)}}{249}
\EGAIVthreeContentsLine{}{INDEX DES NOTATIONS}{250}
\EGAIVthreeContentsLine{}{INDEX TERMINOLOGIQUE}{251}

\vspace{2em}
\begin{center}
\emph{Manuscrit reçu le 15 novembre 1964.}
\end{center}

\let\EGAIVthreeContentsLine\relax
\endgroup


\end{document}
